% METADATA DEFINITIONS (Hidden from PDF output)













\documentclass[11pt, a5paper, twoside]{article}


    \raggedbottom


\usepackage{fontspec}
\usepackage{polyglossia}
\usepackage{geometry}
\usepackage{fancyhdr}
\usepackage{tcolorbox}
\usepackage{ifthen}
\usepackage{tikz}
\usetikzlibrary{shadows}
\usepackage{xcolor}
\usepackage{microtype}
\usepackage{graphicx}
\usepackage{float}
\usepackage[object = vectorian]{pgfornament}
\usepackage{tocloft}


% --- Robust Purnaviram Line Break Control (LuaTeX) ---
\usepackage{luacode}
\begin{luacode*}
function fix_devanagari_punctuation_breaks(head)
  local GLYPH = node.id("glyph")
  local GLUE  = node.id("glue")
  local KERN  = node.id("kern")
  local PENALTY = node.id("penalty")

  for n in node.traverse_id(GLYPH, head) do
    -- 0x0964 = । (Purnaviram), 0x0965 = ॥ (Deerghapurnaviram)
    if n.char == 0x0964 or n.char == 0x0965 then
      
      -- 1. Prohibit break BEFORE punctuation
      local curr = n.prev
      while curr and (curr.id == GLUE or curr.id == KERN or curr.id == PENALTY) do
        if curr.id == GLUE or curr.id == KERN then
            local p = node.new(PENALTY)
            p.penalty = 10000 -- Block break
            head = node.insert_before(head, curr, p)
        elseif curr.id == PENALTY then
            curr.penalty = 10000 -- Override existing penalty
        end
        curr = curr.prev
      end
      
      -- 2. Encourage break AFTER punctuation
      local nxt = n.next
      if nxt and nxt.id == GLUE then
         local p = node.new(PENALTY)
         p.penalty = -500 -- Preferred break point
         head = node.insert_before(head, nxt, p)
      end
    end
  end
  return head
end
luatexbase.add_to_callback("pre_linebreak_filter", fix_devanagari_punctuation_breaks, "Fix Devanagari Breaks")
\end{luacode*}

% Support settings for the above
\sloppy
\emergencystretch=5cm
\hyphenpenalty=10000
\exhyphenpenalty=10000
\linepenalty=1000

% Layout settings for better line breaking
\sloppy
\emergencystretch=5cm


\geometry{
    top=1in,
    bottom=1in,
    left=0.5in,
    right=0.5in
}

% --- spacing configuration ---
\linespread{1.0}
\setlength{\parskip}{6pt}
\setlength{\parindent}{20pt}

% --- Devanagari Page Numbering Macros ---
\makeatletter
\def\devanagaridigits#1{\expandafter\@devanagaridigits\number#1@}
\def\@devanagaridigits#1{%
  \ifx#1@\else
    \ifcase#1 ०\or १\or २\or ३\or ४\or ५\or ६\or ७\or ८\or ९\fi
    \expandafter\@devanagaridigits
  \fi
}

\def\devanagarialph#1{\expandafter\@devanagarialph\number#1}
\def\@devanagarialph#1{%
  \ifcase#1\or अ\or आ\or इ\or ई\or उ\or ऊ\or ऋ\or ॠ\or ऌ\or ॡ\or ए\or ऐ\or ओ\or औ\or क\or ख\or ग\or घ\or ङ\or च\or छ\or ज\or झ\or ञ\or ट\or ठ\or ड\or ढ\or ण\or त\or थ\or द\or ध\or न\or प\or फ\or ब\or भ\or म\or य\or र\or ल\or व\or श\or ष\or स\or ह\or ळ\or क्ष\or ज्ञ\else\@ctrerr\fi
}
\makeatother

\newcommand{\ornamentalpage}{%
    %
        \smash{%
        \begin{tikzpicture}[baseline=(page.base)]
            \node[color=theme, inner sep=0pt, yshift=0pt] (ornl) {\pgfornament[width=0.3cm]{21}};
            \node[anchor=west, xshift=5pt, inner sep=0pt] (page) at (ornl.east) {\thepage};
            \node[anchor=west, color=theme, xshift=5pt, inner sep=0pt, yshift=0pt] (ornr) at (page.east) {\pgfornament[width=0.3cm]{23}};
        \end{tikzpicture}%
        }%
    %
}

\pagestyle{fancy}
\fancyhf{} 


    \renewcommand{\thesection}{\devanagaridigits{\value{section}}}
    \renewcommand{\thesubsection}{\thesection.\devanagaridigits{\value{subsection}}}
    \renewcommand{\thesubsubsection}{\thesubsection.\devanagaridigits{\value{subsubsection}}}


% -- Rules --

    \renewcommand{\headrulewidth}{0.4pt}



    \renewcommand{\footrulewidth}{0.4pt}


% -- Header/Footer Positioning --

    % Mirrored (Outer/Inner logic)
    % Header Left -> Left Side on Odd (Outer), Right Side on Even (Outer)
    
    % Header Right -> Right Side on Odd (Outer), Left Side on Even (Outer)
    \fancyhead[RO,LE]{\leftmark}
    % Header Center
    
    
    % Footer Left
    
    % Footer Right
    \fancyfoot[RO,LE]{वैदिकसाहित्याध्ययनम्}
    % Footer Center
    \fancyfoot[CO,CE]{\ornamentalpage}


% --- Theme Color Definition ---
\definecolor{theme}{named}{black}

% --- Vishesh Vishaya Suchi (Special Table of Contents) ---
\newcommand{\visheshcontents}[5]{%
    \cleardoublepage
    \begingroup
    
    % 1. Use specialkhanda for title (handles marks, TOC entry, and fancy style)
    \specialkhanda{#1}%
    
    % 2. Suppress standard ToC title and force fancy style
    \renewcommand{\contentsname}{}% 
    \renewcommand{\cftaftertoctitle}{\thispagestyle{fancy}}% 
    \tocloftpagestyle{fancy}%
    
    % 3. Style Mahakhanda (Part)
    \cftpagenumbersoff{part}% Hide page numbers for Mahakhanda
    \renewcommand{\cftpartfont}{\hfill\raisebox{#5}{\mbox{\pgfornament[width=#4]{#2}}}\quad\Large\bfseries\color{theme}}% Left Orn + Title
    \renewcommand{\cftpartafterpnum}{\quad\raisebox{#5}{\mbox{\pgfornament[width=#4]{#3}}}\hfill\null\par\vspace{1em}}% Right Orn + Gap Below
    \setlength{\cftbeforepartskip}{2.5em}%
    \setlength{\cftpartnumwidth}{0pt}% 

    % 4. Style Khanda (Section)
    \renewcommand{\cftsecfont}{\large\bfseries}%
    \renewcommand{\cftsecleader}{\cftdotfill{\cftdotsep}}%
    \setlength{\cftbeforesecskip}{0.5em}%

    \tableofcontents
    \endgroup
    \pagebreak
}

% --- Custom Segments (Khanda) ---
\newcommand{\mahakhanda}[1]{
    \cleardoublepage
    \thispagestyle{empty}
    \vspace*{\fill}
    \begin{center}
        \textbf{\Huge \color{theme} #1}
    \end{center}
    \vspace*{\fill}
    \addcontentsline{toc}{part}{#1}
    \setcounter{section}{0}
    \pagebreak
}

\newcommand{\specialmahakhanda}[1]{
    \cleardoublepage
    \thispagestyle{empty}
    \vspace*{\fill}
    \begin{center}
    \begin{tikzpicture}
        \node[inner sep=20pt, fill=theme!10, rounded corners=10pt, drop shadow] (title) {
            \parbox{0.8\textwidth}{\centering \textbf{\Huge \color{theme} #1}}
        };
    \end{tikzpicture}
    \end{center}
    \vspace*{\fill}
    \addcontentsline{toc}{part}{#1}
    \setcounter{section}{0}
    \pagebreak
}

\newcommand{\specialmahakhandatwo}[2]{
    \cleardoublepage
    \thispagestyle{empty}
    \vspace*{\fill}
    \begin{center}
    \begin{tikzpicture}
        \node[inner sep=20pt, fill=theme!10, rounded corners=10pt, drop shadow] (title) {
            \parbox{0.8\textwidth}{\centering \textbf{\Huge \color{theme} #1} \par \vspace{0.6cm} \textbf{\Large \color{theme} #2}}
        };
    \end{tikzpicture}
    \end{center}
    \vspace*{\fill}
    \addcontentsline{toc}{part}{#1}
    \setcounter{section}{0}
    \pagebreak
}

\newcommand{\khanda}[1]{
    \markboth{#1}{}
    \vspace{0.8cm}
    \noindent\makebox[\textwidth][c]{{\huge \color{theme} \textbf{#1}}} \par
    \addcontentsline{toc}{section}{#1}
    \vspace{0.4cm}
}

\newcommand{\upakhanda}[1]{
    \vspace{0.6cm}
    \noindent\makebox[\textwidth][r]{{\Large \color{theme} \textbf{#1}}} \par
    \addcontentsline{toc}{subsection}{#1}
    \vspace{0.3cm}
}

\newcommand{\upaupakhanda}[1]{
    \vspace{0.4cm}
    \noindent\makebox[\textwidth][l]{{\normalsize \color{theme} \textbf{#1}}} \par
    \addcontentsline{toc}{subsubsection}{#1}
    \vspace{0.2cm}
}

% --- Custom Paragraph ---
\newcommand{\customparagraph}[1]{
    \vspace{0.2cm}
    \noindent \textbf{#1} \par
    \vspace{0.1cm}
}

% --- Alignment and Spacing Macros ---
\newcommand{\vspacer}[1]{\vspace{#1 cm}}

% --- Shloka Environment ---
\newcounter{shlokaline}
\newenvironment{shloka}[1][alternating]{
    \par\vspace{0.5cm}
    \begingroup
    \parindent=0pt
    \def\shlokaalign{#1}
    \setcounter{shlokaline}{0}
    \ifthenelse{\equal{\shlokaalign}{alternating}}{
        \let\oldpar\par
        \def\par{\oldpar\stepcounter{shlokaline}}
        \everypar{\ifodd\value{shlokaline}\raggedleft\else\raggedright\fi}
    }{
        \ifthenelse{\equal{\shlokaalign}{left}}{\raggedright}{}
        \ifthenelse{\equal{\shlokaalign}{right}}{\raggedleft}{}
        \ifthenelse{\equal{\shlokaalign}{center}}{\centering}{}
    }
}{
    \par\endgroup
    \vspace{0.5cm}
}

% --- Vishesh Shloka Environment (with Ornament) ---
\newenvironment{visheshshloka}[1][alternating]{
    \begin{shloka}[#1]
}{
    \end{shloka}
    \vspace{-0.8cm}
    \begin{center}
        \color{theme}\pgfornament[width=0.2\textwidth]{88}
    \end{center}
    \vspace{-0.3cm}
}

% --- Vishesh Shloka Two Environment (with Custom Ornament) ---
\newenvironment{visheshshlokatwo}[2][alternating]{
    \def\visheshornament{#2}
    \begin{shloka}[#1]
}{
    \end{shloka}
    \vspace{-0.8cm}
    \begin{center}
        \color{theme}\pgfornament[width=0.2\textwidth]{\visheshornament}
    \end{center}
    \vspace{0.1cm}
}

% --- Image Macros ---
% --- Vishesh Butta (Repeated Ornaments) ---
\newcommand{\visheshbutta}[5]{%
    \par\vspace{0.4cm}
    \noindent\begin{minipage}{\textwidth}
        \ifthenelse{\equal{#1}{Center}}{\centering}{}%
        \ifthenelse{\equal{#1}{Left}}{\raggedright}{}%
        \ifthenelse{\equal{#1}{Right}}{\raggedleft}{}%
        \foreach \n in {1,...,#3}{%
            \pgfornament[width=#4]{#2}\hspace{#5}%
        }%
    \end{minipage}\par\vspace{0.4cm}
}

\newcommand{\customimage}[3]{
    \begin{figure}[H]
        \ifthenelse{\equal{#2}{Center}}{\centering}{}
        \ifthenelse{\equal{#2}{Left}}{\raggedright}{}
        \ifthenelse{\equal{#2}{Right}}{\raggedleft}{}
        \includegraphics[width=#3\textwidth]{#1}
    \end{figure}
}

% --- Special segments with TikZ gradient ---
\newcommand{\specialkhanda}[1]{
    \markboth{#1}{}
    \vspace{1.2cm}
    \noindent\makebox[\textwidth][c]{{\huge \color{theme} \textbf{#1}}} \par \nointerlineskip
    \addcontentsline{toc}{section}{#1}
    \noindent\begin{tikzpicture}[remember picture, overlay]
        
            \shade[left color=white, right color=theme!80!white] (0,-0.4) rectangle (0.5\textwidth, -0.2);
            \shade[left color=theme!80!white, right color=white] (0.5\textwidth,-0.4) rectangle (\textwidth, -0.2);
        
    \end{tikzpicture}
    \vspace{0.8cm}
}

\newcommand{\specialupakhanda}[1]{
    \vspace{0.8cm}
    \noindent\makebox[\textwidth][r]{{\Large \color{theme} \textbf{#1}}} \par \nointerlineskip
    \addcontentsline{toc}{subsection}{#1}
    \noindent\begin{tikzpicture}[remember picture, overlay]
        
            \shade[left color=white, right color=theme!60!white] (0,-0.3) rectangle (\textwidth, -0.15);
        
    \end{tikzpicture}
    \vspace{0.6cm}
}

\newcommand{\specialupaupakhanda}[1]{
    \vspace{0.4cm}
    \noindent\makebox[\textwidth][l]{{\normalsize \color{theme} \textbf{#1}}} \par \nointerlineskip
    \noindent\begin{tikzpicture}[remember picture, overlay]
        
            \shade[left color=theme!40!white, right color=white] (0,-0.2) rectangle (\textwidth, -0.1);
        
    \end{tikzpicture}
    \vspace{0.4cm}
    \addcontentsline{toc}{subsubsection}{#1}
}

% --- Font Configuration ---

    \setmainfont[Script=Devanagari, AutoFakeBold=1.8]{Tiro Devanagari Sanskrit}


\begin{document}
\sloppypar
\renewcommand{\thepage}{\devanagarialph{\value{page}}}

\specialkhanda{\textbf{किञ्चिद् मन्तव्यम्}}

नेपालसंस्कृतविश्वविद्यालयेन शास्त्रिस्तरीयपाठ्यक्रमस्य पुनर्निर्माणं कृतम् । तत्क्रमे साहित्यविषयस्य पाठ्यक्रममपि नवीकृतम् । तत्र तृतीयखण्डस्य षष्ठपत्रे वैदिकसाहित्याध्ययनम् इति शीर्षके वैदिकसाहित्यस्य सामान्यज्ञानार्जनाय वेदोपनिषद्ब्राह्मणग्रन्थतः काश्चन मुख्यविषयान्नुधृत्य पाठ्यांशवस्तु निर्धारितं विद्यते । अस्य पाठ्यांशस्य मूलग्रन्था बृहत्तरा विद्यन्ते । अतच्छात्रैस्तानि सर्वाणि पुस्तकानि पृथक्तया प्राप्य पठनाय काठिन्यमनुभूयते । तस्मात् सर्वान् ग्रन्थांशानेकत्रैव सङ्कलय्य सारल्येन सम्पादनमावश्यकमिति मया अनुभूतम् । समयक्रमेण विद्यापीठे समागतैः कुलसचिवपदमलङ्‌कुर्वद्भिर्माधवप्रसादाधिकारीवर्यैः, पूर्वशिक्षाध्यक्ष्यैः प्रा. डा. भीमप्रसादखतिवडावर्यैर्यत् प्रेरकं वचनं सङ्‌घुष्टम्, तेनाप्यहं प्रेरितो जातः । विषयेऽस्मिन् अस्माकं विश्वविद्यालयस्य साहित्यविषयसमितेरध्यक्षपदासीनेभ्यः प्राप्तविद्यावारिधिकेभ्यः प्राध्यापकपदमलङ्‌कुर्वद्भ्यो नुरापतिपोखरेलगुरुवरेभ्यो विषयवस्तुविशिष्टीकरणादिकार्येषु अविस्मणीया प्रेरणा प्राप्ता । प्रथमतः कार्येऽस्मिन् प्रेरकान् सर्वान् मान्यजनान् सादरं नौमि ।

अहं साहित्यविषयाध्ययनाध्यापनकार्ये संलग्नोऽस्मि । तस्माद् मम कृते वैदिकसाहित्यस्य पूर्णं ज्ञानं नासीत् । अतः पाठ्यक्रमानुसारं पुस्तकमिदं निर्मातुं मया मूलग्रन्थानां सहायकग्रन्थानाञ्च अध्ययनं कृतम् । अस्मिन् कार्ये मम विद्यावारिध्युपाधिलाभाय कृतस्य शोधप्रबन्धस्य निर्देशकैर्नुरापतिपोखरेलगुरुवरैर्विहितेन पुस्तकादिसामग्रीसंप्रेषणेन प्राप्तं सौकर्यं शब्दवाच्यमेव न इति मन्ये । मम इदं प्रथमं पुस्तकसम्पादनकार्यम् । अतोऽत्र नैका त्रुटयः सञ्जाताः । तानपि सम्यक् विनिभाल्य परिमार्जनकार्यं नुरापतिगुरुभिरेव कृतम् । तस्माद् गुरुवरान् प्रति कृतकृत्यमनुभूय अनन्तं नतभावं ज्ञापयामि ।

विश्ववाङ्‌मये संस्कृतभाषा महनीया । इयं भाषा सामान्यलोकभाषा अतिरिच्य स्वकीयं वैशिष्ट्यं तनोति । इयञ्च भाषा लौकिकी वैदिकी चेत्युभयथा विभक्ता विद्यते । तयोर्वैदिकी भाषा कठिना अनुभूयते । तस्याः कृते प्राचीनैर्मुनिभिः कृतं भाष्यमेवावलोकनीयम् । लघुपुस्तकमिदं निर्मातुं मया सायणमुनिना कृतं सायणभाष्यं सहाय्येन स्वीकृतम् । सायणभाष्यानुसारमेव मया तत्र पर्यायशब्दाः समुल्लिखिताः । तदनुसारमेव संस्कृते नेपाल्यां च भावार्थप्रकाशनं कृतम् । प्रसङ्गेऽस्मिन् ऋषिरामरेग्मीति नाम्ना प्राध्यापकेन नेपालीभाषायां रचितं वैदिकसाहित्याध्ययनम् इति पुस्तकमपि सहायकसामग्रीत्वेन गृहीतम् । एवमेव अन्येऽपि संस्कृतहिन्दीनेपालीभाषासु प्रकाशिता ग्रन्था अपि सामग्रीत्वेन गृहीताः । पूर्ववर्तिनां येषां विदुषां वैदिकभावप्रकाशका ग्रन्था मया सामग्रीत्वेन स्वीकृतास्तान्प्रति हार्दं श्रद्धाभावं पुरस्करोमि ।

इतः पूर्वं मया केवलं लघुलेखादिकमेव रचितमासीत् । संस्कृतभाषायां पुस्तकाकारकृतिस्वियमेव प्रथमा। अतोऽत्रावसावधानवशात् त्रुटयो भवेयुः । ता विनिभाल्य विद्वांसः शोधनीयं प्रेरकञ्च वचनं सम्प्रेष्य ममोत्साहं वर्धयिष्यन्तीति मम बलवती आशा विद्यते । पुनश्च सर्वेभ्यो मान्यजनेभ्योऽग्रजेभ्यः शुभचिन्तकेभ्यो महानुभावेभ्यो शतशः नतभावं व्यहरामि । नमो नमः ।

\begin{center}

\textbf{विनीतः}

गुरुप्रसादकोइराला, विद्यावारिधिः

\textbf{सहप्राध्यापकः}

\textbf{धरानस्थस्य पिण्डेश्वरविद्यापीठस्य}

\end{center}

\visheshcontents{विषय सूची}{17}{18}{0.7cm}{5pt}

\setcounter{page}{1}

\specialmahakhandatwo{प्रथमः पाठः}{वैदिकसूक्ताध्ययनम्}

\specialkhanda{\textbf{वेदपरिचयः}}

‘शास्त्राणि विन्दन्ति प्राप्नुवन्ति अस्मिन्’ इति व्युत्पत्तौ  ‘विद्’ धातो ‘घञ्’ प्रत्ययेन ‘वेदः’ शब्दो निष्पद्यते । तत्र केचन विद्वासश्चतुर्भिर्धातुभिर्वेदशब्दं निष्पादयन्ति । पौरस्त्यज्ञानस्य मूलभूता रहस्यपूर्णाः सिद्धान्तास्तेषु वेदेषु प्रतिपादिताः सन्ति । भगवता मनुना प्रोक्तमपि विद्यते यत्- ‘वेदोऽखिलो धर्ममूलम्’ इति । अखिलब्रह्माण्डस्य समस्तधर्माणां मूलानि वेदा एव सन्ति । अद्य यावदुपलब्धेषु वाङ्मयग्रन्थेष्वाधुनिकैरपि वेदाः प्राचीनतमा मन्यन्ते । वेदाः कदा केन रचिता ? इति विषये विदुषां मतैक्यं नास्ति । विषयेऽस्मिन् आध्यात्मवादिनां भौतिकवादिनाञ्च पृथक् धारणा विद्यते । अध्यात्मवादिनः ‘वेदा अपौरुषेयाः, ध्यानस्थानामृषीणां मानसे प्रतिध्वनिरूपेण सञ्जाताः, पश्चाच्च लिपिबद्धाः कृताः, अत एव मन्त्रद्रष्टार ऋषयः’ इति वदन्ति । तद्यथा पराशरसंहितायाम्-

\begin{shloka}[center]

न कश्चिद् वेदकर्तास्ति वेदस्मर्ता चतुर्मुखः ।

ब्रह्माद्या ऋषयः सर्वे स्मारका न तु कारकाः॥

\end{shloka}

भौतिकवादिनः पाश्चात्त्यविद्वांसस्तु मुनिनिर्मिता वेदाः, तेषां रचनाकाल ईशवीवर्षारम्भात् प्राक् षट्‌सहस्रवर्षादारभ्य द्विसहस्रवर्षं यावद् इति मन्यन्ते । विषयेऽस्मिन् मतवैभिन्ये सत्यपि ‘सर्वप्राचीनो ग्रन्थो वेदः’ इत्यस्मिन् विषये न केषाञ्चिद् विमतिः ।

\renewcommand{\thepage}{\devanagaridigits{\value{page}}}

\customparagraph{\textbf{संहिता (वेदविभागः)}}

संग्रहार्थकेन ‘संहिता’ पदेन मन्त्राणां सङ्ग्रह इत्यर्थो ज्ञायते । वेदस्य प्रमुखाे विभाग एव ‘संहिता’ इति कथ्यते । ‘त्रयी’ इति शब्देन वेदानां संहितात्रितयत्वावस्थितिर्बुध्यते परं तास्तु ऋग्वेदसंहिता, यजुर्वेदसंहिता, सामवेदसंहिता, अथर्ववेदसंहिता चेति चतस्रो विद्यन्ते । मूलतः ‘त्रयी’ शब्दः संहितासङ्‌ख्यायाः कृते न, रचनाशैलीभेदस्य कृते प्रयुज्यते । तासु चतसृषु संहितासु गद्यं पद्यं गीतञ्चेति त्रिसृषु शैलीषु रचिता मन्त्राः प्राप्यन्ते । यत्र तत्र प्रसृतानां मन्त्राणामेकत्रीकरणमेव संहितीकरणमुच्यते । तच्च भगवता व्यासेन विहितम्, पैल-वैशम्पायन-जैमिनि-सुमन्तुभ्यः स्वशिष्येभ्यश्च समाम्नातम् । उपदेशक्रमे पैलाय ऋक्संहितायाः, वैशम्पायनाय यजुःसंहितायाः, जैमिनये सामसंहितायाः, सुमन्तवे चाथर्वसंहिताया उपदेेशो विहितः ।

वैदिकमन्त्राणां सङ्‍ग्रहरूपाः संहिताश्च्श्चतस्रः सन्ति । संहिताया विभाजनं मन्त्राणां कार्यमाधृत्य कृतमस्ति ।  यज्ञीयकार्यसम्पादनाय चत्वार ऋत्विजां गणा भवन्ति । तेषु प्रत्येकेषु गणेषु चत्वारः प्रमुखा ऋत्विजो निर्धारिता विद्यन्ते । तत्र होतृगणे होता, अध्वर्युगणे अध्वर्युः, उद्‌गातृगणे उद्गाता, ब्रह्मगणे च ब्रह्मा मुख्यो भवति । तेषु होता ऋक्संहितामन्त्रान्, अध्वर्युर्यजुस्संहितामन्त्रान्, उद्गाता  सामसंहितामन्त्रान्, ब्रह्मा चाथर्वसंहितामन्त्रान् पठति । अतो मन्त्राणां  भिन्न-भिन्नकार्यमनुलक्ष्य भगवता व्यासेन चतस्रः संहिताः सङ्गठिता इति ज्ञायते ।

\customparagraph{\textbf{ऋग्वेदसंहिता}}

होता मन्त्रान्नुच्चार्य देवान् आह्वयति स्तौति च । तस्य कृते स्तुतिपरका मन्त्राः ऋग्वेदसंहितायां सङ्‌गृहीताः सन्ति । चतसृषु संहितासु ऋक्संहिताया महत्त्वं सर्वाधिकं विद्यते । व्यासात् पैलाख्येन मुनिना श्रुतायामृक्संहितायां स्तुतिपरकाणां मन्त्राणां सङ्‌‌ग्रहोऽस्ति । ऋक्संहितायां बहुभिर्मुनिभिः सुललितैर्भावगभीरैः शब्दैर्विभिन्ना देवताः सादरं संस्तुताः सन्ति । अस्या उल्लिखितासु एकविंशतिशाखासु पञ्चशाखाः प्रसिद्धाः सन्ति । तासु शाकलः, शाङ्‌खायनीयः, आश्वलायनीयश्चेति तिस्रः शाखा एव समुपलब्धाः सन्ति । एतासु शाकलशाखाया एवाधुनाऽस्माकं विश्वविद्यालये पठनपाठनं भवति । पद्यमयी ऋक्संहिता अष्टकाध्यायवर्गभेदेन मण्डलानुवाकसूक्तभेदेन च द्विधा विभक्ता विद्यते । अष्टकक्रमतः ‘अष्टौ’ (८) अष्टकानि, ‘चतुःषष्टिः’ (६४) अध्यायाः, ‘षडधिकसहस्रद्वयमिता’श्च (२००६) वर्गाः सन्ति । बालखिल्यसूक्तानि विहाय पूर्णायामृक्संहितायां ‘दशसङ्‌ख्यकानि’ (१०) मण्डलानि, ‘पञ्चाशीतिः’(८५) अनुवाकाः, ‘सप्तदशोत्तरसहस्राणि’ (१०१७) सूक्तानि सन्ति। प्रत्यष्टकम् अष्टाध्यायाः सन्ति । अध्यायस्य अवान्तरविभागानां नाम ‘वर्गः’ अस्ति । वर्गः ऋचां समुदायस्य नामास्ति, परं वर्गेषु ऋचां सङ्‌ख्या निश्चिता नास्ति । दशसङ्‌ख्यायुतेषु मण्डलेषु सप्तदशोत्तरैकसहस्राणि सूक्तानि सन्ति। एतेभ्यः सूक्तेभ्योऽतिरिक्तानि एकादशसूक्तानि बालखिल्यनाम्ना प्रसिद्धानि सन्ति । ‘खिलः’ शब्दस्यार्थो भवति परिशिष्टो भागः । एतानि सूक्तानि ऋक्संहितायाः परिशिष्टे ‘बालखिल्य-’ नाम्ना सङ्‌कलिता विद्यन्ते । बालखिल्यसूक्तैः सह ऋग्वेदस्थसूक्तानां सङ्ख्या ‘१०२८’ भवति ।

\customparagraph{\textbf{यजुर्वेदसंहिता}}

इज्यते यज्ञं सम्पद्यतेऽनेन इति यजुः, तस्य संहिता सङ्‌ग्रहो यत्र सा यजुस्संहिता । यज्ञकर्मप्रयोज्या मन्त्रा यजुर्वेदसंहितायां न्यस्ताः सन्ति। गद्यमयी यजुस्संहिता कर्मकाण्डप्रधाना मन्यते । कर्मणैव यज्ञः सम्पद्यते । अस्याः संहितायाः पाठकोऽध्वर्युः । अध्वरं यज्ञं युनक्तीति- अध्वर्युः। अयं चतुर्षु ऋत्विग्गणेषु अध्वर्युगणस्य मुख्योऽस्ति । वेदव्यासाद् वैशम्पायनेन श्रुताया यजुर्वेदसंहितायाः कृष्णशुक्लभेदेन द्वौ भेदौ स्तः । कृष्णशुक्लयजुस्संहिताया विविधानां शाखानां सङ्केतः प्राप्यते परं पूर्णरूपेण कृष्णयजुर्वेदस्य तैत्तिरीय-मैत्रायणी-कठ-कपिष्ठलकठाख्याश्चतस्रः शाखाः समुपलभ्यन्ते । शुक्लयजुर्वेदस्य च माध्यन्दिनशाखा कण्वशाखा चेति द्वे शाखे उपलब्धे स्तः।  शुक्लयजुर्वेदस्य \textbf{माध्यन्दिनसंहितायाम्-} चत्वारिंशदध्यायाः १९७५ मन्त्राः सन्ति । \textbf{कण्वशाखायाम्-} चत्वारिंशदध्यायाः २०८६ मन्त्राश्च सन्ति । कृष्णयजुर्वेदस्य \textbf{तैत्तिरीयशाखायाम्-} सप्त काण्डानि, चतुश्चत्वारिंशत् प्रपाठकाः, ६५१ अनुवाकाः, २१९८ मन्त्राश्च सन्ति । \textbf{मैत्रायणीशाखायाम्-} चत्वारि काण्डानि, ११ प्रपाठकाः, ३१४४ मन्त्राः सन्ति । \textbf{कठसंहितायाम्}- ४० अध्यायाः, ८४३ अनुवाकाः, ३०९१ मन्त्राः सन्ति । \textbf{कपिष्ठलकठसंहितायाम्}- ६ अष्टकानि ४८ अध्यायाः सन्ति । कपिष्ठलकठसंहिता अपूर्णा एवोपलब्धा विद्यते ।

\customparagraph{\textbf{सामसंहिता}}

सामयति सान्त्वयति देवान् इति ‘सामन्’ शब्दस्यार्थः । शब्दोऽयं प्रसन्नतार्थे शान्त्यर्थे च प्रयुज्यते । अर्थाद् देवानां प्रसन्नतायै गीयमानाः ऋचः सामानि कथ्यन्ते । वेदव्यासाद् जैमिनिना मुनिना श्रुतायाः सामसंहिताया महत्त्वं ख्यापयन् भगवान् नारायणो गीतायामुक्तवान्- ‘वेदानां सामवेदोऽस्मि’ इति। संहितासु इयं सामसंहिता उपासनाकाण्डत्वेन प्रसिद्धा अस्ति । गानमय्या अस्याः संहितायाः ऋत्विग् ‘उद्गाता’ अस्ति। उद्गाता यज्ञस्थले सस्वरं मन्त्रगायनं करोति । तेन गेया मन्त्राः सामवेदसंहितायां विद्यन्ते। देवानां प्रसन्नतायै, प्राणाभिवृद्धये च सामगानमत्यावश्यकं मन्यते । सामवेदसंहितायाः  पूर्वार्चिकः उत्तरार्चिकश्चेति द्वौ भागौ स्तः । पूर्वार्चिकः- छन्दः, छन्द आर्चिकः, छन्दःसंहिता, योनिग्रन्थश्च कथ्यते । उत्तरार्चिकः- उत्तरा, उत्तरसंहिता इत्यादिनाम्ना ज्ञायते । पूर्वार्चिकभेदे षट् अध्यायाः, षडेव प्रपाठकाः सन्ति । पूर्वार्चिकवद् उत्तरार्चिकभेदस्यापि विभागाः सन्ति । सर्वान् संयोज्य पूर्वार्चिकभागे ६५० तथा चोत्तरार्चिकभागे १२२५ मन्त्राः सन्ति । सम्प्रति सामसंहितायाः कौथुमीया, राणायनीया, जैमिनीया चेति तिस्रः शाखाः समुपलभ्यन्ते ।

\customparagraph{\textbf{अथर्वसंहिता}}

चतसृषु संहितासु चरमा संहिता अथर्वसंहिता अस्ति । चतुर्षु ऋत्विक्षु ‘ब्रह्मा’ नामक ऋत्विक् यज्ञस्य प्रमुखाध्यक्षो भवति । एष ऋत्विग् ‘यज्ञकार्ये त्रुटयो मा भवन्तु’ इति विचार्य सावधानेन यज्ञस्थलस्य निरीक्षणं करोति । अयं सर्ववेदज्ञो भवति । तथाप्यस्य कृते आवश्यका मन्त्रा अथर्ववेदसंहितायां सङ्‌कलिताः सन्ति । ऋगादयस्त्रयो वेदा आमुष्मिकं फलं प्रयच्छन्ति । परमथर्वसंहिता तु आमुष्मिकेण सह ऐहलौकिकं फलमपि प्रयच्छति । अथर्वसंहितायां शान्तिलाभस्य पौप्टिककार्यस्य च सम्पादनाय विधानानि वर्णितानि सन्ति । वेदव्यासः सुमन्त- नामकं शिष्यमथर्ववेदमध्यापितवान् । अत्र याज्ञिकविधानविस्तरेण सह यज्ञकर्तॄणामन्येषाञ्च मनुष्याणां समस्यानां समाधानं वर्णितमस्ति । तस्माद् अथर्ववेदसंहिता वैदिकसंहितासु पृथग् वैशिष्ट्यमावहति । महाभाष्यकारेण पतञ्जलिना पस्पशाह्निके ‘नवधाऽथर्वणो वेदः’ इत्युक्तम्। परमधुना नवशाखासु पैप्पलादशौनकाख्ये द्व एव शाखे प्रकाशिते स्तः । पैप्पलादशाखायाः सामान्यैव चर्चा प्राप्यते । शौनकसंहितायाम्- २० काण्डानि, ३६ प्रपाठकाः, १११ अनुवाकाः, ७३६ सूक्तानि, ५९८७ मन्त्राः सन्ति ।

सामान्यतया यज्ञसम्पादनपरासु चतुसृषु संहितासु ऋग्वेदसंहितायां स्तुतिपरका मन्त्राः, यजुर्वेदसंहितयां यज्ञविधिविषयका मन्त्राः, सामसंहितायां गायनविषयभूता मन्त्राः, अथर्वसंहितायां विघ्नविनाशका मन्त्राः सङ‌्‌गृहीता विद्यन्ते ।

प्राधान्येन वेदो द्विविधो मन्यते- मन्त्ररूपः, ब्राह्मणरूपश्चेति । मन्त्रणां सङ्ग्रह एव ‘संहिता’ इति व्यवह्रियते । संहिताया एव व्याख्यात्मको ब्राह्मणग्रन्थः । वेदस्यायं ब्राह्मणभागो यागस्वरूपस्य ‘बोधकः’ इति प्रसिद्धः । ब्राह्मणग्नन्थोऽपि त्रिधा विभक्तो विद्यते, यत्- ब्राह्मणम्, आरण्यकम्, उपनिषद् चेति । तेषु यज्ञस्वरूपप्रतिपादका ब्राह्मणग्रन्था; सामान्यतया अरण्ये पठनीयाः यज्ञस्याध्यात्मिकस्वरूपविवेचका आरण्यकग्रन्थाः; ब्रह्मस्वरूपस्य बोधकाः, मोक्षसाधकाश्च उपनिषद्‌ग्रन्था इति । सामान्यतया ब्राह्मणभागो गृहस्थाश्रमावलम्बिनाम्, आरण्यकभागो वानप्रस्थाश्रमाश्रितानाम्, उपनिषद्भागश्च तुर्याश्रमसेविनाम् उपयोगी इति मन्यते ।

\customparagraph{\textbf{वेदाङ्गानि}}

वेदानां सूत्रात्मकत्वेन अर्थज्ञाने काठिन्यमनुभूयते । तन्निराकरणायैव वेदाङ्गानां प्रणयनम् । तस्माद् वेदस्य पू्र्णं ज्ञानं वेदाङ्गैर्जायते । ‘अङ्‌ग्यन्ते ज्ञायन्ते अमीभिः अङ्गानि’ इति अङ्गशब्दस्याशय उपकारको बोध्यते । तस्माद् वेदस्य यानि अङ्गानि सन्ति तानि सर्वाणि वेदस्य तत्त्विकज्ञाने उपकारकाणीति ज्ञायन्ते । वेदस्य वास्तविकज्ञानाय षडङ्गानि ज्ञातव्यानि । तानि च -

\begin{shloka}[center]

शिक्षा व्याकरणं छन्दो निरुक्तं ज्योतिषं तथा ।

कल्पश्चेति षडङ्गानि वेदस्याहुर्मनीषिणः॥

\end{shloka}

शिक्षा, व्याकरणम्, छन्दः, निरुक्तम्, ज्योतिषम्, कल्पश्चेति षड्‌ वेदाङ्गानि सन्ति । प्रतीकात्मकरूपेण वेदं शरीरत्वेन परिकल्प्य वेदशरीरस्य षडङ्गानां कार्यमपि विद्वद्भिः समुल्लिखितं प्राप्यते । यथा पाणिनीयशिक्षायाम्-

\begin{shloka}[center]

छन्दः पादौ तु वेदस्य हस्तौ कल्पोऽथ पठ्यते ।

ज्योतिषामयनं चक्षुर्निरुक्तं श्रोत्रमुच्यते ॥

शिक्षा घ्राणं तु वेदस्य मुखं व्याकरणं स्मृतम् ।

तस्मात्साङ्गमधीयानो ब्रह्मलोके महीयते ॥

\end{shloka}

वेदस्य छन्दः पादौ स्तः । कल्पो हस्तौ स्तः । ज्योतिषशास्त्रं नयने स्तः । निरुक्तं कर्णौ स्तः। शिक्षा नाशिका उल्लिखिता । व्याकरणशास्त्रं मुखत्वेन परिकल्पितं विद्यते । अत्र विचारणीयं विद्यते यत्- अस्माकं शरीरे पादादिषु कस्यचिदेकस्याङ्गस्य अभावो भवति चेज्जीवने महत् काठिन्यमनुभूयते । सर्वेषामङ्गानामन्यूनानतिरिक्तं महत्त्वं विद्यते । एवमेव वेदशरीरेऽपि शिक्षादिषु षडङ्गेषु एकेनापि विना तस्य ज्ञानजीवनं कष्टावृतं जायते । अतो वैदिकसाहित्यज्ञानाय शिक्षाव्याकरणादीनां षण्णां वेदाङ्गानां महत्त्वं समानमेव संलक्ष्यते ।

\customparagraph{\textbf{शिक्षा (उच्चारणविधिः)}}

वेदाङ्गेषु शिक्षायाः प्रमुखं स्थानं विद्यते । पाणिनीयशिक्षायां ‘शिक्षा’ वेदस्य घ्राणम् (नासिका) समुल्लिखितं विद्यते । शिक्षायाः सामान्यार्थो भवति यत्- वर्णोच्चारणस्य सम्यक् ज्ञानकरणम् । मन्त्रोच्चारणे वर्णोच्चारणस्य शिक्षाप्रदानमेव शिक्षाया उद्देश्यं भवति । कस्य वर्णस्योच्चारणं कथं करणीयम् ? कः प्रयत्नो विधेयः ? वर्णानां विभाजनं कथं भवति ? कति स्थानानि सन्ति ? मन्त्रपाठक्षणे श्वासवायुर्वर्णरूपेण कथं परिवर्तते ? इत्यादिविषया एव शिक्षाख्ये वेदाङ्गशास्त्रे प्रतिपादिता भवन्ति । एतेषां शिक्षाप्रतिपादितानां विषयाणामध्ययनं विना वेदमन्त्राणां पाठकार्यम् अर्थज्ञानञ्चासम्भवप्रायं ज्ञायते । सम्प्रति त्रिंशत्सङ्‌ख्यकाः शिक्षाग्नन्था उपलभ्यन्ते । तेषुु पाणिनीयशिक्षा प्रसिद्धतरा विद्यते । अत्र प्रातिशाख्यग्रन्था अपि स्मरणीयाः सन्ति । प्रतिवेदं तेऽपि ग्रन्थाः पृथक् विद्यन्ते ।

\customparagraph{\textbf{व्याकरणम् (शब्दव्युत्पत्तिः)}}

‘मुखं व्याकरणं स्मृतम्’ इति पाणिनीयशिक्षायां समुल्लिखितं विद्यते । व्याकरणं वेदस्य मुखं मन्यते । यथा मुखं विना भोजनादिव्यवहाराभावेन शरीरस्य पोषणादिकं न भवति तथैव वेदस्यापि व्याकरणात्मकं मुखं विना महत्त्वादिकं नैव वर्धते । व्याकरणशास्त्रमन्तरा शब्दानां सम्यक् ज्ञानं न जायते । वेदे स्थितानां कठिनशब्दानां व्युत्पत्तये व्याकरणशास्त्रस्य महत्त्वमन्यूनं विद्यते । पाणिनीयव्याकरणं वेदाङ्गं मन्यते । पाणिनेरष्टाध्यायी, कात्यायनस्य वार्तिकम्, पतञ्जलेर्महाभाष्यं पाणिनीयव्याकरणनाम्ना मुनित्रयव्याकरणनाम्ना च प्रसिद्धम् ।

\customparagraph{\textbf{छन्दःशास्त्रम् (लयविधानम्)}}

‘छद्’ आच्छादने धातोः छन्दःशब्दो निष्पद्यते । छन्दो वेदस्य पादस्थानीयं मन्यते । यथा पादरूपे आधारे शरीरमवस्थितं भवति तथैव वेदशरीरमपि छन्दस आधारे स्थितं भवति । छन्दांसि ज्ञेयानि भवन्ति । तदाधारिता रचना श्रोतरि वक्तरि च श्रवणानन्दमाविष्करोति । वेदमन्त्रा अपि छन्दसः कारणेन श्रवणीयाः सन्ति । तस्मात् छन्दोबद्धानां वेदमन्त्राणां सम्यक्तया गायनाय छन्दःशास्त्रस्याध्ययनमावश्यकमिति । महामुनेः पिङ्गलस्य छन्दशास्त्रं प्रसिद्धम् । तत्र वैदिकानि लौकिकानि च छन्दांसि समुल्लिखितानि सन्ति ।

\customparagraph{\textbf{निरुक्तम् (वैदिकशब्दनिर्वचनम्)}}

‘निरुक्तं श्रोत्रमुच्यते’ इति वेदस्य श्रवणेन्द्रियरूपेण निरुक्तं मन्यते । निरुच्यते निःशेषेण उपदिश्यते इति निरुक्तम् । निरुक्तशब्दस्य सामान्यार्थो भवति- निर्वचनं व्युत्पत्तिश्च । शब्दस्योत्पत्तिविकासप्रभृतिविषया निरुक्ते समाविश्यते। कस्यापि शब्दस्य प्रकृतिप्रत्यययोः स्पष्टीकरणम्, समानार्थकनानार्थकादिशब्दानां विवेचनं निरुक्तं विना असम्भवमेव । वैदिकशब्दा अपि निर्वचनेन विना ज्ञातुमशक्या एव । तस्मात् वेेदस्य सम्यगध्ययनाय निरुक्तशास्त्रस्यापि महत्त्वं समधिकं विद्यते । यास्कमुनेर्निरुक्तं प्रसिद्धं विद्यते । तत्र चतुर्दशाध्यायाः सन्ति ।

\customparagraph{\textbf{कल्पः (कर्मकाण्डाचारविधिः)}}

‘कल्पः’ वेदस्य हस्तौ मन्यते । यथा हस्तौ शरीरं विविधकर्मसु प्रेरयतः, तथैव यज्ञयागादिवैदिककृत्यानां विधानं कल्पग्रन्थेषु विद्यते । यज्ञसम्बद्धानां समस्तानां विधीनां समर्थनं निर्देशनं प्रतिपादनञ्च कल्पशास्त्रे लभ्यते । अतो वेदाध्ययने कल्पशास्त्रस्य महती उपयोगिता विद्यते । कल्पसूत्रमपि चतुर्धा विभक्तमस्ति- श्रौतसूत्रम्, गृह्यसूत्रम्, धर्मसूत्रम्, शुल्वसूत्रम् (परिशिष्टश्च)। श्रौतकर्मप्रतिपादकं श्रौतसूत्रम्, गृहस्थाश्रमसम्बद्धकर्मप्रतिपादकं गृह्यसूत्रम्, वर्णाश्रमादि- धर्मप्रतिपादकं धर्मसूत्रम्, यज्ञवेदीनिर्माणादिसम्बद्धम् शुल्वसूत्रम् । श्रौतगृह्य- सूत्रापेक्षितविषयप्रतिपादकः परिशिष्टोऽपि कल्पसूत्रान्तर्गत एव । केचित् शुल्वसूत्रम्, परिशिष्टञ्च कल्पासूत्रावयवं न मन्यन्ते । पुनश्च प्रत्येकानां वेदानां पृथक् पृथक्  कल्पसूत्राणि विद्यन्ते ।

\customparagraph{\textbf{ज्योतिषम् (यज्ञसमयनिरूपणम्)}}

वेदाङ्गेषु ज्योतिषशास्त्रस्य महत्त्वपूर्णं स्थानं विद्यते । पाणिनिः कथयति यत्- ‘ज्यौतिषामयनं चक्षुः’ इति । यज्ञयागादीनां सम्यग् रूपेण सम्पादनाय समयविशेषस्य ज्ञानमपेक्ष्यते। यज्ञस्य साफल्याय केवलमुचितविधानस्यैव नापितु उचितसमयस्यापि आवश्यकता भवति । मङ्गलमुहूर्ते क्रियमाणा यज्ञादिक्रिया एव फलान्विता भवन्ति नान्यदा । मुहूर्तादिसमयज्ञानं ज्यौतिषशास्त्रसम्बद्धम् । तस्माद् वेदज्ञानलाभाय ज्यौतिषशास्त्रमपि ज्ञातव्यं भवति । वेदानां पृथक् पृथक् ज्योतिषशास्त्रमासीत् । ऋग्वेदस्य आर्चज्योतिषम्, यजुर्वेदस्य याजुषज्योतिषम् अथर्ववेदस्य आथर्वणज्योतिषमिति । सम्प्रति लगधप्रणीता ऋगादीनां ज्योतिषग्रन्थाः प्राप्यन्ते, परं सामवेदस्य न कोऽपि ज्योतिषग्रन्थ उपलभ्यते।

संस्कृतवाङ‌मयस्य गौरवभूताद् वेदाद् गभीरज्ञानावाप्तये तन्मर्मबोधाय च तस्य अङ्गभूतानि षट्‍शास्त्राणि ज्ञातव्यानि भवन्ति । इमानि षट्‌शास्त्राणि वेदमन्त्राणां सम्यग् रहस्यबोधनायैव निर्मितानि सन्ति ।

\customparagraph{\textbf{प्राचीनतमाया ऋक्संहिताया वैशिष्ट्यम्}}

ऋच्यते स्तूयते यया सा ‘ऋक्’ इति ‘ऋच स्तुतौ’ इति धातोः क्विप्प्रत्ययेन निष्पन्नोऽयं स्तवनपरकः ‘ऋक्’ शव्दः । ऋचां समूहः ‘ऋग्वेदः’ कथ्यते । प्राचीनकाले वेद एक एवासीत् । भगवता कृष्णद्वैपायनव्यासेन द्वापरयुगस्यान्तेे मानवबुद्ध्या ग्रहणयोग्यमध्ययनसौकर्यं चाभिलक्ष्य एक एव वेद ऋग्यजुसामाथर्वरूपेण चतुर्धा विभक्तः । सम्प्रति ऋग्वेदसंहितापदेन सामान्यतया शाकलसंहिता गृह्यते । अस्यां च संहितायामष्टकवर्णीकरणम्, मण्डलवर्गीकरणमिति द्विविधं वर्गीकरणं विद्यते । अष्टकवर्गीकरणानुसारेण ऋग्वेदवेदस्य अष्टकः, अध्यायः, वर्गः, इति ऋचां क्रमविभागो विद्यते । तेष्वष्टसङ्‌ख्यका (८) अष्टकाः, प्रत्येकेष्वष्टकेषु अष्टाध्यायानां गणनया चतुष्षष्टिः (६४) अध्यायाः सन्ति । षडधिकद्विसहस्रं (२००६ ) वर्गाः (बालखिल्यसूक्तस्य अष्टादश वर्गान् सङ्‌गृह्य चतुर्विंशत्यधिकद्विसहस्रं (२०२४) वर्गाः सन्ति । ऋग्वेदे  मण्डलक्रमेण दशमण्डलानि, पञ्चाशीति (८५) अनुवाकाः, सप्तदशाधिकानि सहस्राणि (१०१७) सूक्तानि च विद्यन्ते । द्विसप्तत्युत्तरचतुश्शताधिकैकायुतं (१०४७२) (बालखिल्यसूक्तस्य अशीतिसङ्‌ख्यायोगेन द्विपञ्चाशदुत्तरपञ्चशताधिकैकायुतम्- १०५५२) ऋचो विद्यन्ते । ऋग्वेदस्य एकविंशतिः शाखा आसन्, परमधुना तिस्रः शाखा एव लभ्यन्ते । शाकलशाखा, आश्वलायनशाखा, शाङ्‌खायनशाखा चेति । सर्वेषां वैदिकलौकिकसाहित्यानाम् आधारभूता, प्राचीततमा, विविधासु शाखासु प्रसृता, ज्ञानराशिसंवलिता, ऋग्वेदसंहिता पौरस्त्यवाङ्मयस्यैव नापितु विश्ववाङ्मयस्य कृतेऽपि महत्त्वाधायिका मन्यते ।

\customparagraph{\textbf{साहित्यिकदृष्ट्या ऋग्वेदस्य विश्लेषणम्}}

पौरस्त्यसाहित्यं वैदिकलौकिकयोर्द्वयोर्भागयोर्विभक्तमस्ति । द्वयोर्मध्ये वैदिकसाहित्यं प्राचीनतमं विद्यते । वैदिकसाहित्ये वेदाः चतुर्धा विभज्यन्ते- ऋक्, यजुस्, सामस्, अथर्वणञ्चेति । तेषां प्रत्येकानां पुनश्चतुर्धा विभागः संहिताब्राह्मणारण्यकोपनिषद्‌भेदात् । एते चत्वारो वेदा ब्रह्मचर्येण, गार्ह्यस्थ्येन, आरण्यकेन, संन्यासेन आश्रमेण सम्बद्धाः । एषु सर्वेषु ऋग्वेदः प्राचीनतमः । तत्र दशमण्डलानि सबालखिल्यम् अष्टाविंशत्यधिकसहस्राणि सूक्तानि सन्ति । सृष्टेः प्राक्कालस्य तस्य अव्यवहितपरकालस्य च यज्ज्ञानम् ऋषिभिः प्रत्यक्षीकृतं तस्य अवधारणा अतिगूढतया सूत्ररूपेण एषु मन्त्रेषु उपलभ्यते । तत्र लौकिकसाहित्ये लोकसमाजे चावश्यकीयानि तत्त्वानि निबद्धानि सन्ति ।

ऋग्वेदस्य संवादसूक्तानि काव्यमयत्वादत्यन्तमेव हृद्यानि । तेषु पुरुरवोर्वशीसंवादयमयमीसंवादप्रभृतयः प्रसिद्धाः सन्ति । अत्र सामाजिकरीतिसम्बद्धानि सूक्तानि लभ्यन्ते, येषां सङ्ख्या प्रायो विंशतिर्विद्यते । एवमेव शकुनसूक्तम्, विवाहसूक्तम्, पितृृसूक्तम्, रहस्यसूक्तमित्यादीनि च प्रसिद्धानि । तत्रत्यानि नासदीयपुरुषहिरण्यगर्भसूक्तानि दार्शनिकविषयसम्बद्धानि ।  विद्वद्भिरेतान्येव सूक्तान्यङ्गीकृत्य रसमयानि हृद्यानि विक्रमोर्वशीयमित्यादीनि काव्यानि प्रणीतानि ।

संस्कृतकाव्यशास्त्रेऽधुना विकसितानि रसध्वन्यलङ्कारादीनि यानि काव्यतत्त्वानि सन्ति, प्रायशः तेषां सर्वेषां बीजम् ऋग्वेदसूक्तेषु लभ्यते । तस्मात् काव्यरचनाविधिनिर्माणस्य काव्यप्रणयनस्य च प्रेरकत्वेन ऋग्वेदसूक्तानां महत्त्वमसामान्यं विद्यते ।

\specialkhanda{\textbf{ऋग्वेदीयं नासदीयसूक्तम्}}

\customparagraph{\textbf{परिचयः}}

ऋग्वेदसंहिताया दशमे मण्डले एकनवत्युत्तरशतसङ्ख्यकानि सूक्तानि सन्ति । तेष्वेकोनत्रिंशदुत्तरशतसङ्‍ख्यकं सूक्तं ‘नासदीयसूक्तम्’ इति ज्ञायते । अत्र सप्त ऋचः सन्ति । अस्य सूक्तस्य परमेष्ठी प्रजापतिर्ऋषिः, सृिष्टिस्थितप्रलयकर्ता परमात्मा देवता, त्रिष्टुप्च्छन्दश्चास्ति । सूक्तेऽस्मिन् सृष्ट्यारम्भात् प्राग् जगतः कीदृशं स्वरूपमासीदिति चिन्तनमस्ति । यथा चास्य सूक्तस्य नाम्नैव ज्ञायते यत्- न+असत् इति नासत्, तदीयं सूक्तं नासदीयसूक्तम् । सृष्टेः पूर्वं जगत्कीदृशमासीदिति जिज्ञासया प्रवृत्तं सूक्तमिदं पौरस्त्यदार्शनिकविचारधारायाः प्रवाहरूपो मन्यते । सूक्तेऽस्मिन् ब्राह्मी सृष्टिरपि निर्दिष्टा विद्यते । नासदासीन्नो सदासीत्तदानीम् इत्यारभ्य प्रस्तुतं नासदीयसूक्तं सृष्टिविज्ञानसूक्तमित्यपरनाम्नापि विख्यातमस्ति ।

\customparagraph{\textbf{वैशिष्ट्यम्}}

प्राणिषु समुत्कृष्ट विवेकसमन्वितोमानवः । एष स्वभाविकरूपेण जिज्ञासुर्भवति,  तस्माद् वयं सर्वे जिज्ञासवः । तदनुसारमस्माकं सर्वेषां हृदि ज्ञातुमिच्छा भवत्येव यत्- यत्र वयं वसामः, तस्य स्थानस्य किं महत्त्वम् ? तस्यैतिहासिकी पृष्ठभूमिश्च का वर्तते ? अस्माकं पूर्वजाः कथमत्रागत्य निवासं चक्रुः ? अस्या आवासभूमेः किं वैशिष्ट्यमिति ? अस्माकं पूर्वजानामागमनात् प्रागस्यां भूमौ किमासीत् ? एतेषां समेषां प्रश्नानामुत्तराणि अस्माभिरन्वेषणीयान्येव भवन्ति । एवमेव साध्याः देवाः सृष्ट्युत्पत्तेः प्राक्कालीनां प्रकृतेः स्थितिमन्वेषयन्तो नैकधा चिन्तितवन्तः । तेषां चिन्तनस्य यन्निष्कर्षमासीत्तदेव नासदीयसूक्तरूपेणास्माकं समक्षं विद्यमानमस्तीति ज्ञायते । अतो विश्वव्यापिनमीश्वरमनुज्ञातुम्, तत्सृष्टस्यास्य ब्रह्माण्डस्य स्वरूपमनुमातुञ्च अस्य सूक्तस्य महत्त्वं सार्वत्रिकं सार्वकालिकं सार्वजनीनञ्चास्ति ।

\visheshbutta{Center}{4}{3}{0.3cm}{0.2cm}

\begin{visheshshloka}[center]

\textbf{नासदासीन्नो सदासीत्तदानीं नासीद्रजो नो व्योमा परो यत् ।}

\textbf{किमावरीवः कुह कस्य शर्मन्नम्भः किमासीद्‍गहनं गभीरम् ॥ १ ॥}

\end{visheshshloka}

\customparagraph{\textbf{अन्वयः}}

तदानीम् असत् न आसीत्, सत् नो आसीत्, रजः न आसीत्, व्योम न (आसीत्) यत् परः किम् आवरीवः ? कुह ? कस्य शर्मन् ? किम् गहनम् गभीरम् अम्भ आसीत् ?

\customparagraph{\textbf{पदार्थः}}

तदानीम्‌‌= तस्मिन् काले, सृष्टेः पूर्ववर्तिनि समये इत्यर्थः । असत्= नामरूपविहीनं तत्त्वम्, न आसीत्= नैव अवस्थितमासीदित्यर्थः (परमतत्त्वस्य शशविषाणस्यैव सर्वथाभावो नासीत्)। सत्= नामरूपयुक्तं किमपि रजः= पृथ्व्यादयो लोकाः व्योमा= अन्तरिक्षम् न आसीत्, परः= आकाशादुपर्यपि यत्= यदिदानीमस्ति तत् किम्= कीदृशं तत्त्वम्, आवरीवः= आवरकमासीत् ? कुह= कुत्र देशे कस्य= भोक्तुः जीवस्य वा शर्मन्= सुखे संरक्षणे वा किम्= कति गहनम्= पारावाररहितम्, गभीरम्= गम्भीरम्, दुष्प्रवेश्यम् अम्भः= जलम् आसीद्, वा किमपि नासीत् ।

\customparagraph{\textbf{भावार्थः}}

तदानीं सृष्ट्युत्पत्तेः प्राक् प्रलयावस्थायाम् असद्भावयुक्तं सद्भावात्मकञ्च तत्त्वं नासीत् । तथा च सत्तात्मकं तत्त्वमपि नासीत् । पृथिवीमारभ्य पातालपर्यन्ता ये रजोलोकविशिष्टाः सन्ति तेऽपि नासन् । नान्तरिक्षमासीत् । अन्तरिक्षात्परश्चापि यदिदानीमनुज्ञायमानमस्ति तत्सर्वं किमपि नासीत्। परमिदं तु ज्ञायते स्म, यदिदं समस्तं शून्यमयं जगत् केनापि आवरणेन समावृतमासीत्, तत् आवरकं तत्त्वं किमासीत् ? कुत्रासीत् ? कस्य संरक्षणे आसीत् ? एतत् सर्वं पाषाणे तैलमिव निरुपाख्यमासीत् । तदानीम् दुष्प्रवेश्यम् अत्यन्तं गभीरञ्च जलं यदिदानीमस्ति तत्किमासीत्  इत्यपि न ज्ञायते स्म । अर्थात् तदानीं नैतत् किमप्यासीदित्यर्थः ।

\customparagraph{\textbf{नेपाली भावार्थ}}

यस जगत्‌को सृष्टिभन्दा पहिले असत् अर्थात् नामरूपविहीन तत्त्व थिएन, सत् अर्थात् नामरूपधारी स्थावरजङ्गमादि केही पनि थिएन । त्यस समयमा स्वर्ग, मर्त्य, पातालजस्ता लोक पनि थिएनन् । अन्तरिक्ष र अन्तरिक्षभन्दा पर पनि केही थिएन । त्यस वेला कहाँ ? कसका आश्रयमा ? केले केलाई ढाकेको थियो ? गहन र गम्भीर पानी थियो वा थिएन केही निश्चित थिएन।

\begin{visheshshloka}[center]

\textbf{न मृत्युरासीदमृतं न तर्हि न रात्र्या अह्न आसीत् प्रकेतः ।}

\textbf{आनीदवातं स्वधया तदेकं तस्माद्धान्यन्न परः किञ्चनास ॥ २॥}

\end{visheshshloka}

\customparagraph{\textbf{अन्वयः}}

तर्हि मृत्युः न आसीत् । अमृतं न (आसीत्) । रात्र्या अह्नः प्रकेतः न आसीत् । तत् आनीत् अवातम् स्वधया एकम् । ह तस्मात् अन्यत् किञ्चन न आस न परः ।

\customparagraph{\textbf{पदार्थः}}

तर्हि= तदानीम्, मृत्युः= मरणधर्मविशेषः, न= नैव, आसीत्= स्थितः । अमृतम्= अमरणधर्मविशेषः, न= नैव आसीत् । रात्र्याः= रात्रिकालस्य, अह्नः= दिनस्य, प्रकेतः= ज्ञानम्, न= नैव, आसीत्= वर्तते स्म । तत्= ब्रह्मतत्त्वमेव, आनीत्= प्राणेन संयुतम्, अवातम्= क्रियाशून्यम्, स्वधया= स्वमायया, एकम्= अविभक्तं सद् एकरूपे विद्यमानमासीत् । ह= खलु तस्मात्= पूर्वोक्ताद् मायासहिताद् ब्रह्मणः, अन्यत्= भिन्नः, किञ्चन=  किमपि वस्तु भूतभौतिकात्मकं जगत्, न= नैव आस= आसीत् । न= नैव परः= परस्तात् सृष्टेरूर्ध्वम् आसीत् । वर्त्तमानमिदं जगत्तदानीं शून्यावस्थायामासीत्  ।

\customparagraph{\textbf{भावार्थः}}

तस्यां प्रलयावस्थायां प्राणिनां मृत्युः अमरणादिकं वा किमपि नासीत् । न च कश्चित् कालो विद्यते स्म । कालस्य परिज्ञानमहोरात्राभ्यां भवति, तद्हेतुभूतयोः सूर्याचन्द्रमसोरभावाद् मासर्तुसंवत्सरादिकालस्याभाव आसीत् । तद् ब्रह्म एव एकः प्राणिवत् प्राणनकर्मणायुक्तः क्रियाहीनः स्वधया (स्वस्मिन् धीयते इति स्वधा माया तया सह) स्वकीयया मायया भेदरहितोऽद्वितीयः सन् विद्यमानमासीत् । तस्माद् मायासहिताद् ब्रह्मणो भिन्नं जगत्यस्मिन् न किमपि आसीत् । न च तस्मात् परस्तात् किमपि आसीदित्यर्थः ।

\customparagraph{\textbf{नेपाली भावार्थ}}

सृष्टिभन्दा पहिले न मृत्यु थियो न त अमृत वा जीवन नै थियो । रात र दिनको कुनै सङ्केत थिएन अर्थात् सूर्य र चन्द्र पनि थिएनन् । त्यसवेला एउटा त्यही तत्त्व आफूमै आधृत भएर क्रियाविहीन तर जीवित थियो । त्यसभन्दा पर अन्य केही पनि थिएन ।

\begin{visheshshloka}[center]

\textbf{तम आसीत्तमसा गूढमग्रेऽप्रकेतं सलिलं सर्वमा इदम् ।}

\textbf{तुच्छ्येनाभ्वपिहितं यदासीत्तपसस्तन्महिनाजायतैकम् ॥ ३॥}

\end{visheshshloka}

\customparagraph{\textbf{अन्वयः}}

अग्रे तमसा गूढम् तमः आसीत् । अप्रकेतम् इदम् सर्वम् सलिलम् आः । यत् आभु तुच्छ्येन अपिहितम् आसीत् । तत् एकम् तपसः महिना अजायत ।

\customparagraph{\textbf{पदार्थः}}

अग्रे= सृष्टिरचनात्‌ प्राक्, तमसा= अन्धकारेण गूढम्= आच्छादितम्, तमः= माया (सृष्टिमूलकारणम्), आसीत्= विद्यमानमासीत् । अप्रकेतम्= अप्रकाशितम्, इदं= जगत्, सर्वम्= सम्पूर्णम्, सलिलम्= सलिलमयम्, आः= समन्तादासीत् । यत्= जगत्, आभु= व्यापकेन तुच्छ्येन= अज्ञानेन, अपिहितम्= आच्छादितमासीत् । तत्= जगत्, एकम्= एकीभूतं सत्, तपसः= तपस्यायाः (सङ्कल्परूपं तपस्तस्य) महिना= माहात्म्येन, अजायत= प्रादुरभूत् ।

\customparagraph{\textbf{भावार्थः}}

सृष्टेः प्राक् सर्वं जगद् अज्ञानरूपेणान्धकारेणावृतमासीत् । सर्वत्र तमसा प्रकीर्णम् अप्रकेतं नामरूपाद्युपाधिहीनं शून्यमिव अप्रज्ञायमानं सलिलं सर्वत्र व्याप्तमासीत् । तदा तुच्छ्येन शून्यसदृशेन ‘आभुना’ आसमन्तात् भवतीति आभु व्याप्तरूपेणाज्ञानेन आच्छादितमासीत्। तद् एकीभूतम् अविभक्तं परमतत्त्वस्य सङ्कल्पमयेन मायासामर्थ्येन किमपि तत्त्वं (मनः) सञ्जायत ।

\customparagraph{\textbf{नेपाली भावार्थ}}

सृष्टिभन्दा पहिले यो ब्रह्माण्ड अन्धकारले (मायाले) ढाकिएको अवस्थामा सर्वत्र अन्धकार (अज्ञान) मात्र थियो । यो सम्पूर्ण जगत्‌मा अज्ञानमय जल व्याप्त थियो । सर्वत्र व्याप्त जगत् तुच्छ अज्ञानले ढाकिएको थियो । त्यसैवेला ब्रह्मको इच्छाशक्तिको महिमाले विशेष तत्त्व (मन) उत्पन्न भयो ।

\begin{visheshshloka}[center]

\textbf{कामस्तदग्रे समवर्तताधि मनसो रेतः प्रथमं यदासीत् ।}

\textbf{सतो बन्धुमसति निरविन्दन् हृदि प्रतीष्या कवयो मनीषा ॥४॥}

\end{visheshshloka}

\customparagraph{\textbf{अन्वयः}}

अग्रे कामः तत् अधि समवर्तत यत् मनसः प्रथमम् रेतः आसीत्, कवयः मनीषा हृदि प्रतीष्य असति सतो बन्धुम् निरविन्दन् ।

\customparagraph{\textbf{पदार्थः}}

अग्रे= सृष्टेः प्राक्, कामः= सङ्कल्पः (जगतः स्रष्टुमिच्छा), तत्= परतत्त्वम्, अधि= अधिकृत्य, समवर्तत= सम्यक् रूपेण अजायत । यत्= तथ्यम्, मनसः= अन्तःकरणस्य प्रथमम्=अग्रम्, रेतः= विकाररूपं सृष्टेर्बीजभूतम् आसीत् । कवयः= गूढार्थवेत्तारो मनीषिणः, मनिषा= विलक्षणप्रज्ञया हृदि= अन्तःस्करणे प्रतीष्य= अनुसन्धाय असति= अस्थिते नामरूपादिविहीने सतो व्याकृतस्य= नामरूपादियुक्तस्य बन्धुम्= सम्बन्धम्, निरविन्दन्= प्राप्तवन्तः ।

\customparagraph{\textbf{भावार्थः}}

अस्य जगतः सृष्टेः प्रागवस्थायां परमब्रह्मणो मनसि ब्रह्माण्डस्य निर्माणे इच्छा सम्यगजायत। कामरूपं यत्तथ्यं मनसोऽन्तःकरणस्य ज्ञानस्य प्रथमं रेताख्यं कार्यरूपो विकार आसीत् । एतत् तथ्यं तत्त्ववेत्तारो मनीषिणः स्वप्रज्ञया हृदि गभीररूपेण विचार्य नामरूपादिविहीने अव्याकृते नामरूपादियुक्तस्य व्याकृतस्य सम्बन्धं सम्यग् रूपेण ज्ञातवन्तः ।

\customparagraph{\textbf{नेपाली भावार्थ}}

त्यसवेला त्यस परमात्मतत्त्वको सङ्कल्पसिर्जित मनमा ब्रह्माण्डको निर्माण गर्ने कामनाको प्रादुर्भाव भयो । जसमा मनको विकार रूप शक्ति थियो । तत्त्ववेत्ता महर्षिहरूले आफ्नै प्रज्ञाले हृदयमा विचार गरेर नामरूपादिविहीन अवस्थासँग नामरूपादिसहितको जगत्‌को सम्बन्ध थाहा पाए ।

\begin{visheshshloka}[center]

\textbf{तिरश्चिनो विततो रश्मिरेषामधः स्विदासी३ दुपरि स्विदासी३त् ।}

\textbf{रेतोधा आसन्महिमान आसन् स्वधा अवस्तात्प्रयतिः परस्तात् ॥५॥}

\end{visheshshloka}

\customparagraph{\textbf{अन्वयः}}

एषां रश्मिः तिरश्चिनः विततः स्वित् अधः आसीत्, स्वित् उपरि आसीत्, रेतोधाः आसन्, महिमानः आसन्, स्वधा अवस्तात्, प्रयतिः परस्तात् ।

\customparagraph{\textbf{पदार्थः}}

एषाम्= अविद्याकामकर्मणाम्, वियदादिभूतजातानि सृजताम्, रश्मिः= सूत्रम् जालं वा, तिरश्चिनः= तिर्यगवस्थितः, विततः= मध्ये विस्तृतः स्वित्= किम् अधः= अधस्तात् आसीत्= अवस्थितमासीत् । स्वित्= आहोस्वित् उपरि= उपरिष्टात् आसीत्= किमासीत् ? रेतोधाः= बीजभूतस्य कर्मणो विधातारः कर्तारो जीवाः आसन्= स्थिता आसन् । महिमानः= महान्तो वियदादयो भोग्याः आसन्= विद्यन्ते स्म । स्वधाः= भोग्यप्रपञ्चः अवस्तात्= अवरो निकृष्टः आसीत् । प्रयतिः= प्रयतिता भोक्ता परस्तात्= पर उत्कृष्ट आसीत् ।

\customparagraph{\textbf{भावार्थः}}

अविद्याकामकर्मणां चेतनाचेतनानां सृजतां रश्मिः सूर्यरश्मिसदृशः सर्वं जगत् व्याप्नुवन् यः कार्यवर्गः प्रसृतः आसीत् । स कार्यवर्गः प्रथमतः किं तिर्यगवस्थितः, किं मध्ये स्थितः आसीत् किं वा अधः अधस्तात् आसीत् । आहोस्विद् उपरि उपरिष्टात् किमासीत् । एषु कार्येषु मध्ये केचिद्भावा रेतोधा बीजभूतस्य कर्मणो विधातारः कर्तारो भोक्तारश्च जीवा आसन् । तत्र च भोक्तृभोग्ययोर्मध्ये भोग्यप्रपञ्चो निकृष्ट आसीत् । भोक्ता उत्कृष्ट आसीत् ।

\customparagraph{\textbf{नेपाली भावार्थ}}

विद्वान्‌हरूले पत्ता लगाएको दृश्य पदार्थ प्रकाशका रूपमा थियो । त्यो तेर्सो फैलिएको थियो कि बीचमा फैलिएको थियो कि ? तलतिर फैलिएको थियो वा माथि फैलिएको थियो निश्चय गर्न नसकिने थियो । एकैपटक सर्वत्र फैलिएको देखियो । त्यससँगै कर्मका कर्ता र भोक्ताहरू पनि उत्पन्न भएका थिए । साथै विभिन्न प्रकारका भोग्यपदार्थहरू पनि थिए । जसमध्ये भोग्यपदार्थ केही निकृष्ट थिए । प्रयत्न गर्ने भोक्ता चाहिँ माथिल्लो तहको थियो ।

\begin{visheshshloka}[center]

\textbf{को अद्धा वेद क इह प्रवोचत् कुत आजाता कुत इयं विसृष्टिः ।}

\textbf{अर्वाग् देवा अस्य विसर्जनेनाथा को वेद यत आ बभूव ॥६ ॥}

\end{visheshshloka}

\customparagraph{\textbf{अन्वयः}}

कः अद्धा वेद, कः इह प्रवोचत् कुत इयम् विसृष्टिः कुतः आजाता, देवाः अस्य विसर्जनेन अर्वाक्, अथ कः वेद यतः आबभूव ।

\customparagraph{\textbf{पदार्थः}}

कः= पुरुषः अद्धा= वास्तविकरूपेण वेद= जानाति कः= कः पुरुषः इह= अस्मिन् लोके प्रवोचत्= प्रब्रूयात्, इयम्= दृश्यमाना विसृष्टिः= बहुप्रकारा सृष्टिः कुतः= कस्मादुपादानकारणात् कुतः= कस्माच्च निमित्तोपादानकारणात् आ= समन्तात् जाता= प्रादुर्भूता, देवाः= स्वर्गस्था इन्द्रादयो देवा ये सर्वज्ञास्ते अस्य= जगतः विसर्जनेन= विविधं यद्भौतिकं सर्जनम्, तेन अवाक्= अर्वाचीनाः सन्ति । अथ= तस्मात् देवा अपि न जानन्ति । तद्व्यतिरक्तः कः= को नाम मनुष्यादिसचेतः वेद= जानाति यतः= यस्मात् कारणाद् आबभूव= कृत्स्नं जगद् अजायत । अर्थाद् जगतः आदिसमयं परब्रह्मव्यतिरिक्तः कोऽपि न जानाति ।

\customparagraph{\textbf{भावार्थः}}

को विवेकशीलः पारमार्थ्येन इमां सृष्टेः प्रागवस्थां जानाति ? एवमेव को वा पुरुषो लोकेऽस्मिन् तद्वक्तुं समर्थः स्याद्यद् इयं दृश्यमाना भूतभौतिकभोक्तृभोग्यादिरूपेण बहुप्रकारा सृष्टिः कस्मात् उपादानकारणात् कस्माच्च निमित्तोपादानकारणात् प्रादुर्भूता इति । एतदुभयं सम्यक् को वेद ? स्वर्गवासिनो वासवादिदेवाश्चापि अस्य जगतो वियदादिभूतोत्पत्तेरनन्तरं जाताः । तादृशास्ते स्वोत्पत्तेः पूर्वकाले प्रादुर्भूतां सृष्टिं कथं जानीयुः । अज्ञानवन्तस्ते कथं प्रब्रूयुः । तस्माद् देवा अपि न जानन्ति । तद्व्यतिरिक्तः को नाम मनुष्यादिर्वेद, यस्मात्कारणात् कृत्स्नं जगदजायत, अर्थात् परब्रह्मव्यतिरिक्तः कोऽपि न जानातीत्यर्थः ।

\customparagraph{\textbf{नेपाली भावार्थ}}

यस बारेमा कसले यथार्थ रूपमा जानेको छ ? कसले यथार्थ रूपमा बताउला र !  यो सृष्टि के र कसबाट भयो ? देवताहरू पनि यस सृष्टिको उत्पत्ति भएपछिका हुन् । अर्थात् त्यसवेला देवताको अस्तित्व नभएकाले देवताहरू पनि सृष्टिका बारेमा जान्दैनन् । त्यसकारण यो सृष्टि केबाट कसरी भयो भन्ने कसले जान्न सक्छ र ! अर्थात् त्यो परमेश्वरबाहेक कोही पनि ज्ञाता छैनन्।

\begin{visheshshloka}[center]

\textbf{इयं विसृष्टिर्यत आबभूव यदि वा दधे यदि वा न ।}

\textbf{यो अस्याध्यक्षः परमे व्योमन्त्सो अङ्ग वेद यदि वा न वेद ॥७॥}

\end{visheshshloka}

\customparagraph{\textbf{अन्वयः}}

इयम् विसृष्टिः यतः आबभूव यदि वा दधे यदि वा न । अस्य यः अध्यक्षः परमे व्योमन् । अङ्ग सः वेद यदि वा न वेद ।

\customparagraph{\textbf{पदार्थः}}

इयं= दृश्यमाना विसृष्टिः= चराचरविशिष्टा विचित्रा सृष्टिः यतः= यस्मान्निमित्तोपादानभूतात्परमात्मन आ= समन्तात् बभूव= समुत्पन्ना यदि= सोऽपि परमात्मा यदीमाम् वा= उत दधे= धारयति । यदि= अथ वा= उत न= न धारयति तर्हि कोऽन्यः धर्तुं शक्नुयात् ? अस्य= जगतः यः= परमात्मा अध्यक्षः= स्वामी ईश्वरो वा सः परमे= उत्कृष्टे व्योमन्= आनन्दस्वरूपे आकाशे प्रतिष्ठितः, अङ्ग= हे प्रियाः श्रोतारः ! सः= सर्वज्ञ ईश्वर एव ताम् वेद= जानाति यदि सोऽपि न जानाति वा= तर्हि न कश्चिदन्यस्तां सृष्टिं वेद= जानीयात् ।

\customparagraph{\textbf{भावार्थः}}

इयं दृश्यमाना विचित्रा सृष्टिः, यस्मात् निमित्तोपादानकारणाभ्यां परमात्मनः सकाशात् समुत्पन्ना तस्मात् सर्वजगतो हेतुभूतः परमात्मा एव इमां सृष्टिं धारयन्नस्ति । तद्व्यतिरिक्तः कोऽपि अन्योऽस्या धारको नास्ति । सः परमात्मा सृष्टिमिमां धारयेन्न वा धारयेत् अन्यस्तु न कश्चिद्धर्तुमिमां शक्नुयात् । यतः परमात्मा अस्य जगतः स्वामी ईश्वरो वा अस्ति सः परमुत्कृष्टे स्थाने आनन्दस्वरूपे आकाशे आकाशवन्निर्मले स्वप्रकाशे प्रतिष्ठितोऽस्ति । हे प्रियश्रोतारः ! एष सुखस्वरूपपरमात्मैव सृष्टिं विजानाति । यदि सर्वज्ञानमयस्वरूपः सोऽपि न जानाति चेत् तद्व्यतिरिक्तो न कश्चिदपि इमां ज्ञातुं समर्थः स्यादित्यर्थः ।

\customparagraph{\textbf{नेपाली भावार्थ}}

यो सृष्टि जसबाट भयो तिनै परब्रह्मले यसलाई धारण गरेका छन् । अरूले त कसैले धारण गर्न सक्दैनन् नै । यस सृष्टिका अध्यक्ष परमात्मा, जो अन्तरिक्षमा रहन्छन्, तिनैले यो सिर्जित जगत्‌का बारेमा जान्न सक्छन्, यदि तिनले पनि जान्दैनन् भने कसैले पनि जान्दैनन् ।

\customparagraph{\textbf{सारः}}

सूक्तेऽस्मिन् सप्त ऋचः सन्ति । प्रथमायामृचि विद्यमानं नासदित्यंशमादाय नासदीयसूक्तमित्यस्य सूक्तस्य नामकरणं विहितमस्ति । सूक्तानुसारं प्रलयावस्थायाम् अस्मिञ्जगति सदसत् किमपि नासीत्, न रजोलोका आसन्, न व्योम एव । व्योमात्परेऽन्तरिक्षलोकाश्चापि नासन् । अस्य जगतः आवरणं किमासीदित्यपि न कश्चिद्वक्तुं समर्थ आसीत् । अम्भः किमासीत् ? आसीत् चेत् तत् कति गम्भीरं गहनञ्चासीदित्यपि न स्पष्टमासीत् । तदानीं मृत्युर्नासीत्, अमृतमपि नासीत्, सूर्याचन्द्रमसोरभावाद् अहोरात्र्योरप्यभाव आसीत्, एकलः स परमात्मैवासीत्, तद्भिन्नो न कश्चिदासीत् । सर्वत्रान्धकारो व्याप्त आसीत्, सर्वत्र जलमेव व्याप्तमासीत् नान्यत् किमपि चिह्नम् अवलोक्यते स्म । सर्वमिदं जगत् आज्ञानेन समन्तात् पिहितमासीत् । तस्य परमेश्वरस्य सङ्कल्परूपतपसः सकाशादिदं जगदुत्पन्नमिति । सोऽकामयत एकोऽहं बहु स्यां प्रजायेय इति। तस्य हृदि सिसृक्षा जागृता। तस्यैवैतत्फलं यदधुनेदं विविधापन्नं जगत् विद्यमानमस्ति । सृष्टिमिमां तं परमात्मानं विहाय नान्यः कश्चिज्जानाति यतस्तु सर्वेऽपि स्वर्गस्था देवाः सृष्टेः पश्चादेवोत्पन्नाः । तस्मात् यदि सोऽपि न ज्ञातुं शक्नुयात् चेत्तर्हि न कोऽप्यन्यः ज्ञातुमर्हति । इत्थं सूक्तेऽस्मिन् सृष्टिपूर्वावस्थाया ब्राह्मसृष्टेश्च वर्णनं विद्यते ।

\begin{visheshshlokatwo}[center]{84}

\textbf{इति सायणभाष्याश्रितं नासदीयसूक्तम्}

\end{visheshshlokatwo}

\pagebreak

\specialkhanda{\textbf{ऋग्वेदीयम् उषासूक्तम्}}

\customparagraph{\textbf{परिचयः}}

ओषति नाशयत्यन्धकारमिति उषा । सुष्ठु = शोभनम्, उक्तम्= कथनमिति सूक्तम् । अनयोः शब्दयोर्योगेन ‘उषासूक्तम्’ इति शब्दो निष्पद्यते । यत्र उषादेव्याः प्रशस्तिर्वर्णिता विद्यते, तदेव ‘उषासूक्तम्’ कथ्यते। वैदिकसाहित्यस्याधारभूताश्चतस्रो वैदिकसंहिताः । ताषु पृथक् देवतानां नैकानि सूक्तानि सन्ति । तत्र ऋक्संहितायाम् अष्टाविंशत्यधिकसहस्रसङ्‌ख्यकानि सूक्तानि प्राप्यन्ते । तत्र प्रसिद्धेषु सूक्तेषु उषासूक्तस्य नाम प्राक् समागच्छति । सूर्योदयात् प्राग् यः स्वर्णाभो ज्योतिःपुञ्जो दृश्यते सा एव उषा । तामेव मानवीकृत्य आराध्यां देवीं वा परिकल्प्य विहिता तस्याः स्तुतिरेव उषासूक्तमिति प्रसिद्धिरस्ति । सूक्तमिदं प्रकृतिवर्णनात्मकानां लौकिककाव्यानामाधारभूतं स्वीक्रियते । उषासूक्तौ प्रकृतिवर्णनस्याकर्षकं सौन्दर्यं लौकिककाव्यानां कृते प्रेरकत्वेन चकास्ति ।

\customparagraph{\textbf{प्राकृतिकवर्णनवैशिष्ट्यम्}}

उषासूक्तम् ऋग्वेदादुद्धृतमस्ति । प्रातःकाले अरुणोदयात् प्राक् या ज्योत्स्ना समायाति सैवोषा । सा सर्वान् प्राणिनः स्व-स्व-कार्येषु प्रेरयति । अतोऽत्र सा प्राणिनां कृते प्राथमिकी शिक्षिका एवेति वर्णिता विद्यते । प्रकृतेर्ज्योत्स्नायुतः प्रारम्भिकः समय एव उषा इति कथ्यते ।

उषाकालः प्रकृतेर्महत्त्वाधायकः समयोऽस्ति । एष कालः प्राणिनः कृते दैनिककार्यसम्पादनाय प्रेरको ज्योत्स्नापूर्णः प्रारम्भिकः समयो विद्यते । अस्यैव समयस्य नाम ‘उषा’ । इयं दिव्यशक्तियुता देवी इति उषासूक्ते सम्बोधिता विद्यते । अस्याः प्रकाशोऽन्येभ्यः प्रकाशेभ्यो विलक्षणो विद्यते । सूर्यस्य प्रखरः प्रकाशो जनैः सादरं ग्रहीतुं न शक्यते, चन्द्रस्य प्रकाशः कार्यसम्पादनाय पर्याप्तो न भवति परमुषायाः समुज्ज्वलः शीतत्त्वगुणयुतः प्रकाशः सर्वेषु प्राणिषु प्रतिदिनं शारीरिकं नावीन्यं वितरति । तदा प्रचलितः शीतलवायुर्जनानाम् उत्साहं वर्धयति । इयमस्माकं कर्मावरोधभूतमन्धकारजालमुत्सार्य अस्मान् कर्मपरम्परायां प्रेरयितुं शयनादुत्थापयति । अस्या एव प्रेरणया वयं यज्ञादिकर्मणि संशक्ता भवामः । यज्ञादिकर्मभिर्देवाः सन्तुष्टा भवन्ति, अत इयं देवानपि उपकरोति । अस्या एव कारणेन वनस्पतयोऽपि ओषधगुणसम्पन्ना भवन्ति । वनस्पतीनामन्नादिप्रदातृत्वेन इयम् अन्नवती वा धनवतीत्यपि कथ्यते । एवमेव फलरूपेण उपभोग्यवस्तुनः प्रदायकान् वनस्पतीन् प्रवर्धयति । इयं प्राणिनः पोष्यपदार्थानां प्रदायिका इत्यस्मिन् सूक्ते वर्णितं विद्यते यत्-

\begin{shloka}[center]

आवहन्ती पोष्या वार्याणि चित्रं केतुं कृणुते चेकिताना ।

ईयुषीणामुपमा शाश्वतीनां विभातीनां प्रथमोषा व्यश्वैत् ॥

\end{shloka}

एवम् प्रकारेण उषासूक्ते वर्णिता इयं प्रकृतिदेवी प्राणिनः कृते आदरणीया विद्यते ।

उषासूक्ते उषा प्रकृत्या देवी इति वर्णिता विद्यते । अत्र उषायाः कारणेन प्राणिन एव नापितु देवा अपि सन्तुष्टा भवन्तीति संवर्ण्यते । जगतः सर्वान् पदार्थान् मूलरूपेण स्वस्वकर्मणि प्रेरिकायाः प्रकृतिदेव्या उषायाः प्राकृतिकं वर्णनं मनोहरं विद्यते ।

\begin{visheshshloka}[center]

\textbf{इदं श्रेष्ठं ज्योतिषां ज्योतिरागाच्चित्रः प्रकेतो अजनिष्ट विभ्वा ।}

\textbf{यथा प्रसूता सवितुः सवायँ एवा रात्र्युषसे योनिमारैक् ॥ १॥}

\end{visheshshloka}

\customparagraph{\textbf{अन्वयः}}

ज्योतिषाम् इदं  श्रेष्ठम् ज्योतिः आ आगात् । चित्रः प्रकेतः विभ्वा अजनिष्ट । यथा रात्री सवितुः प्रसूता एव उषसे सवाय योनिम् आरैक् ।

\customparagraph{\textbf{शब्दार्थः}}

\textbf{ज्योतिषाम्}= ग्रहनक्षत्रादीनां प्रकाशमानानाम् मध्ये \textbf{इदम्}= उषाख्यम्,  \textbf{श्रेष्ठम्}= उत्कृष्टम्, \textbf{ज्योतिः=} प्रकाशपुञ्जः, \textbf{आगात्}= पूर्वस्यां दिशि समागमत् । \textbf{चित्रः}= महनीयः, \textbf{प्रकेतः}= अन्धकारावृतपदार्थस्य प्रज्ञापकः, \textbf{विभ्वा}= विभुर्व्याप्तः सन् भानुः \textbf{अजनिष्ट}= प्रादुरभूत्। \textbf{यथा रात्री}= रजनिः स्वयं \textbf{सवितुः}= सूर्यसकाशात् \textbf{प्रसूता}= उत्पन्ना । सूर्योऽस्तं गच्छति अनन्तरं रात्रिः समागच्छतीत्यर्थः । एवमेव \textbf{रात्रिः}= रजनिरपि उषसे \textbf{सवाय}= उषाया उत्पत्तये \textbf{योनिं}= स्थानम् \textbf{आरैक्}= कल्पितवती ।

\customparagraph{\textbf{भावार्थः}}

ज्योतिर्भेदेषु श्रेष्ठा एषा उषा समागता । अस्यामागतायां सत्यामेव ज्योतिषां मुख्यः तेजसां पतिः सूर्यः प्रकाशितः । यथा रात्रिः सूर्यात् समुत्पन्ना भवति सूर्यगमनानन्तरमुदेतीत्यर्थः तथैव रात्रिः उषाकालस्योत्पत्तये स्थानं दधाति रात्रेरनन्तरमुषा समागच्छतीत्यर्थः ।

\customparagraph{\textbf{नेपाली भावार्थ}}

धेरै प्रकारका ज्योतिहरूमध्ये उषाको ज्योति सर्वश्रेष्ठ छ । यो ज्योति पूर्वमा उदय भयो । उषा उदित भएपछि उज्यालो तेज छर्दै सूर्य उदायो । रात्रि सूर्यबाट जन्मिन्छिन् अर्थात् सूर्य अस्ताएपछि मात्र रात्रिको आगमन हुन्छ । त्यस्तै उषालाई पनि उदयका लागि रात्रिले ठाउँ दिन्छिन्। अर्थात् रात्रिको समाप्तिपछि उषाको उदय हुन्छ ।

\begin{visheshshloka}[center]

\textbf{रुषद्‌वत्सा रुशती श्वेत्यागादारैगु कृष्णा सदनान्यस्या ।}

\textbf{समानबन्धू अमृते अनूची द्यावा वर्णं चरत अमिनाने ॥२॥}

\end{visheshshloka}

\customparagraph{\textbf{अन्वयः}}

रुशती श्वेत्या रुशद्‌वत्सा आ अगात् । अस्याः कृष्णा सदनानि आरैक् उ समानबन्धू अमृते अनूची वर्णम् अमिनाने द्यावा चरतः ।

\customparagraph{\textbf{पदार्थः}}

\textbf{रुशती}= दीप्ता उज्ज्वलवर्णयुता \textbf{श्वेत्या}= उषा \textbf{रुशद्‌वत्सा}= रुशन् दीप्तिमान् सूर्यो वत्सो यस्या सा, यथा मातुरनु वत्सो धावति तथैव उषसोऽनन्तरं सूर्य उदेति तस्माद् रुशद्‌वत्सा उषेत्यर्थः । तादृशी उषा \textbf{आगात्}= आगतवती । \textbf{अस्याः}= उषायाः \textbf{कृष्णा}= कृष्णवर्णयुता रात्रिः, \textbf{सदनानि}= स्थानानि \textbf{आरैक्}= कल्पितवती । \textbf{उ} इति पादपूरणमात्रम् । एते रात्र्युषे, \textbf{समानबन्धू}= स्वस्वास्तित्वलाभाय सूर्यसम्बद्धे इत्यर्थः। \textbf{अमृते}= मरणरहिते (नित्यत्वात्) \textbf{अनुची}= सूर्यगत्यनुसारेण सञ्चलन्त्यौ, \textbf{वर्णम्}= सर्वेषां लोकानां रूपम्, \textbf{अमिनाने}= जरयन्त्यौ क्षीणतां नयन्त्यौ, अथवा परस्परं स्वकीयं रूपं हिंसन्त्यौ, रात्र्युषे \textbf{द्यावा}= नभसा अन्तरिक्षमार्गेण \textbf{चरतः}= प्रतिदिवसं गच्छतः, अनुदिनं भ्रमतः ।

\customparagraph{\textbf{भावार्थः}}

प्रकाशवती श्वेतवर्णा सूर्यरूपेण वत्सेन युता अत एव वत्सवती उषा समागता । आगच्छन्त्यै उषायै कृष्णवर्णयुता रात्रिः स्थानं दत्तवती । सूर्येण सह समानसम्बन्धयुते विनाशरहिते सूर्यगत्यनुसारेण सञ्चलन्त्यौ लोकानां स्वकीयं वा रूपं क्षीणतां नयन्त्यौ रात्र्युषे अन्तरिक्षमार्गेण निरन्तरं भ्रमतः ।

\customparagraph{\textbf{नेपाली भावार्थ}}

दिव्य प्रकाशका कारण उज्ज्वल वर्ण भएकी सूर्यलाई बाछोझैँ पछि लगाएर उषा उदाइन् । उदाउँदै गरेकी उषालाई कालो वर्ण भएकी रात्रिले ठाउँ दिइन् । सूर्यद्वारा नै अस्तित्व बनाएकाले दुबै सूर्यका साथमा समान सम्बन्ध भएका  विनाशरहित लोकको रूपलाई क्षीण बनाउँदै अथवा आफ्नो रूपलाई क्षीण बनाउँदै रात्रि र उषा आकाश मार्गबाट निरन्तर भ्रमण गरिरहन्छन् ।

\begin{visheshshloka}[center]

\textbf{समानो अध्वा स्वस्रोरनन्तस्तमन्यान्या चरतो देवशिष्टे ।}

\textbf{न मेथेते न तस्थतुः सुमेके नक्तोषासा समनसा विरूपे ॥३॥}

\end{visheshshloka}

\customparagraph{\textbf{अन्वयः}}

स्वस्रो अनन्तः अध्वा समानः । तं देवशिष्टे अन्यान्या चरतः । सुमेके नक्तोषासा विरूपे (अपि) समनसा न मेथेते । (तथा) न तस्थतुः ।

\customparagraph{\textbf{पदार्थः}}

\textbf{स्वस्रोः}= भगिन्योः  रात्र्युषयोः, \textbf{अनन्तः}= अवसानरहितः, \textbf{अध्वा}= गमनागमनाय साधनभूतो मार्गः, \textbf{समानः}= एक एव आकाश इति । \textbf{तम्}= मार्गम्, \textbf{देवशिष्टे}= सूर्येण निर्दिष्टे सत्यौ \textbf{अन्यान्या}= एकैका रात्रिः- अनु- उषा इति क्रमेण \textbf{चरतः}= गच्छतः । \textbf{सुमेके}= शोभनमेहने (सुप्रजनने) \textbf{नक्तोषासा}= रात्रिरुषा च \textbf{विरूपे}= रात्रिरन्धकार उषाप्रकाश इति विपरीतलक्षणयुते अपि \textbf{समनसा}= सङ्गतचित्तगुणयुता सत्यौ \textbf{न मेथेते}= परस्परं नाक्रोशतः, न \textbf{तस्थतुः}= न विश्रामं भजतः ।

\customparagraph{\textbf{भावार्थः}}

भगिनीगुणयुतयो रात्र्युषयोर्विरामरहित एक एव आकाशमार्गः । तस्मिन् मार्गे सूर्येण अनुशिष्टे सत्यौ ते क्रमेण भ्रमन्तः । सर्वेषां वनस्पत्यादीनां फलोत्पादनरूपे शोभनप्रजननशक्तियुते रात्र्युषे स्तः । ‘रात्रेरन्धकार उषायाः प्रकाशः’ इति विपरीतगुणयुते अपि ते ऐक्यभावं कृत्वा परस्परकार्ये बाधां न कुरुतः, अविश्रान्ततया केवलं स्वकार्ये संलग्ने भवत इति ।

\customparagraph{\textbf{नेपाली भावार्थ}}

दिदीबहिनीका रूपमा रहेका रात्रि र उषाको निरन्तर अविराम गतिले हिँड्ने आकाशमार्ग एउटै छ । त्यस मार्गमा सूर्यद्वारा निर्दिष्ट भएर दुबै क्रमैले हिँड्छन् । जगत्‌का वनस्पतिमा राम्रो उत्पादन शक्ति दिने रात्रि र उषा दुबैका अन्धकार र प्रकाशजस्ता विपरीत गुण छन् तापनि ती दुबै एकआपसका कार्यमा कहिल्यै बाधा नपुर्‍याई केवल आफ्नो कर्तव्यपथमा सधैँ भ्रमणशील रहन्छन् ।

\begin{visheshshloka}[center]

\textbf{भास्वती नेत्री सूनृतानामचेति चित्रा वि दुरो न आवः ।}

\textbf{प्रार्प्य जगद्व्यु नो रायो अख्यदुषा अजीगर्भुवनानि विश्वा ॥४॥}

\end{visheshshloka}

\customparagraph{\textbf{अन्वयः}}

भास्वती सूनृतानां नेत्री अचेति । चित्रा नः दुरो वि आवः । जगत् प्रार्प्य नः रायः वि अख्यत् उ । उषा विश्वा भुवनानि अजीगः ।

\customparagraph{\textbf{पदार्थः}}

\textbf{भास्वती}= विशिष्टप्रकाशनयुता \textbf{सूनृतानाम्}= वाचाम् \textbf{नेत्री}= उत्पादयित्री, उषायाः प्रादुर्भावपश्चादेव प्राणिनः शब्दं कुर्वन्तीत्यर्थः । एतादृशीमुषाम्, \textbf{अचेति}= अस्माभिरज्ञायि, अर्थात् वयं ज्ञातवन्तः । \textbf{चित्रा}= विचित्रसौन्दर्ययुता ज्ञाता सा \textbf{नः}= अस्माकम्  तमसा तिरोहितानि \textbf{दुरः}= द्वाराणि \textbf{वि आव}= विशेषेण आवृणोत्, उद्‌घाटयति । \textbf{जगत्}= सर्वं भुवनम्, \textbf{प्रार्प्य}= प्रकाशं गमयित्वा \textbf{नः}= अस्माकम् \textbf{रायः}= धनानि \textbf{वि अख्यत्}= विशेषरूपेण प्रकाशनयुतान्यकरोत् । \textbf{उ}= पादपूरणमात्रम् । एतादृशी \textbf{उषा}=देवी \textbf{विश्वा‍‍‍ भुवनानि}= सर्वाणि भुवनानि \textbf{अजीगः}= उद्गिरति, तमसा तिरोहितत्त्वेन अदृश्यानि भुवनानि उषा स्वमुखाद् निर्गमयति प्रकाशयतीत्यर्थः ।

\customparagraph{\textbf{भावार्थः}}

विशिष्टप्रकाशवती उषा प्राणिनः शब्दोच्चारणार्थं प्रेरयति । एतादृशीम् उषां वयं ज्ञातवन्तः । विचित्रसौन्दर्ययुता उषा अस्माकं द्वाराणि उद्‌घाटयति । सर्वाणि जगन्ति प्रकाशमयानि कृत्वा अस्माकं विविधस्वरूपात्मकं धनं प्रकाशयति । तस्माद् उषा अदृश्यानि सर्वाणि जगन्ति उद्गिरतीव प्रतीयते । उषा एव सर्वाणि जगन्ति नेत्रप्रत्यक्षीकरोतीत्याशयः ।

\customparagraph{\textbf{नेपाली भावार्थ}}

विशेष किसिमको प्रकाश छर्दै उषा दिनको प्रारम्भमा प्राणीहरूलाई शब्द उच्चारणका लागि प्रेरणा दिन्छिन् । उषाकाल भएपछि प्राणीहरू बोल्न थाल्छन् । त्यस्ती उषालाई हामीले राम्ररी चिनेका छौँ । प्रकाशको सुनौलो सौन्दर्य लिएर ती सबेरै हाम्रो ढोका खोल्न आइपुग्छिन् । सबैतिर उज्ज्यालो पारेर हाम्रा गाईवस्तु, खेत, बारी आदि चलअचल सम्पत्तिहरू प्रकाशित गरिदिन्छन् । यस्तो लाग्छ कि अन्धकारले अदृश्य पृथ्वीलाई उषा मुखबाट उगेल्छिन्, अर्थात् प्रकाश छरेर दर्शनीय बनाउँछिन् ।

\begin{visheshshloka}[center]

\textbf{जिह्मश्ये ३ चरितवे मघोन्याभोगय इष्टये राय उ त्वम् ।}

\textbf{दभ्रं पश्यद्भ्य उर्विया विचक्ष उषा अजीगर्भुवनानि विश्वा ॥५॥}

\end{visheshshloka}

\customparagraph{\textbf{अन्वयः}}

मघोनी जिह्मश्ये चरितवे त्वम् आभोगये (त्वं) इष्टये (त्वं) राये (व्युच्छन्ती भवति) उ दभ्रं पश्यद्भ्यः विचक्षे । उर्विया उषाः विश्वा भुवनानि अजीगः ।

\customparagraph{\textbf{शब्दार्थः}}

\textbf{मघोनी}= धनसम्पन्ना उषा, \textbf{जिह्मश्ये}= वक्राकारेण शयिताय पुरुषाय, \textbf{चरितवे}= उत्थाय स्वेप्सितं कार्यं प्रति गन्तुं व्युच्छन्ती भवति । त्वम्- \textbf{आभोगाय}= सुखादिभोगाय, त्वम्- \textbf{इष्टये}= यज्ञाय, त्वम्- \textbf{राये}= धनार्जनाय, व्युच्छन्ती भवति । \textbf{दभ्रम्}= स्वल्पम्, \textbf{पश्यद्भ्यः}= तिमिरावृतत्त्वेन ईषदेव द्रष्टुं पारयद्भ्यो मनुष्येभ्यः, \textbf{विचक्षे}= विशिष्टप्रकाशाय प्रेरयति । \textbf{उर्विया}= उर्वी विस्तीर्णा \textbf{उषाः} = देवी \textbf{विश्वा भुवनानि}= समस्तानि जगन्ति, \textbf{अजीगः}= अन्धकारेण तिरोहितानि भुवनानि उद्गीर्णानीव करोति, प्रकाशेन दृष्टिपथं नयतीत्यर्थः।

\customparagraph{\textbf{भावार्थः}}

प्रकाशधनसम्पन्ना उषा रात्रौ वक्राकारेण शयिताय पुरुषाय शयनादुत्थातुं प्रेरयति । त्वं भोगकार्ये, त्वं यज्ञकार्ये, त्वं धनार्जनकार्ये च सन्नद्धो भव इत्याशयेन सर्वान् जनान् स्वस्वकार्येषु प्रेरयति । रात्रिसमये ईषदेव दर्शनसमर्थेभ्यो मनुष्येभ्यो विशिष्टप्रकाशं वितीर्य अन्धकारेण निगलितानि भुवनानि उषया उद्गीर्णानीव दृश्यन्ते ।

\customparagraph{\textbf{नेपाली भावार्थ}}

प्रकाशरूपी सम्पत्तिले पूर्ण भएकी उषाले राती सुतेका मानिसलाई जगाएर विभिन्न काम प्रति प्रेरित गर्दछिन् । रातको अन्धकारले निलेझैं अदृश्य रहेका सम्पूर्ण भुवनलाई मानौं उगेलेझैं गरी प्रकट गराउँछिन् ।

\begin{visheshshloka}[center]

\textbf{क्षत्राय त्वं श्रवसे त्वं महीयै इष्टये त्वमर्थमिव त्वमित्यै ।}

\textbf{विसदृशा जीविताभिप्रचक्ष उषा अजीगर्भुवनानि विश्वा ॥६॥}

\end{visheshshloka}

\customparagraph{\textbf{अन्वयः}}

उषाः क्षत्राय त्वम्, श्रवसे त्वम्, महीयै इष्टये त्वम्, अर्थमिव इत्यै त्वम्, (व्युच्छन्ती) विसदृशा जीविता अभिप्रचक्षे । उषाः विश्वा भुवनानि अजीगः ।

\customparagraph{\textbf{शब्दार्थः}}

\textbf{उषाः}= देवी \textbf{त्वम्}= एकं व्यक्तिविशेषम्,  \textbf{क्षत्राय}= धनाय \textbf{त्वम्}= एकम्, \textbf{श्रवसे}= अन्नाय, \textbf{त्वम्}= एकम्, \textbf{महीयै}=महत्यै \textbf{इष्टये}= अग्निष्टोमादिमहायज्ञाय, \textbf{त्वम्}= एकम् \textbf{अर्थमिव}= अपेक्षितार्थं प्रति, \textbf{इत्यै}= गमनाय व्युच्छन्ती, उषा \textbf{विशदृशा}= विविधप्रकारकान्, \textbf{जीविता}= प्राणिजीवनभूतान्, अन्नादिभोग्यपदार्थानित्यर्थः \textbf{अभिप्रचक्षे}= आभिमुख्येन प्रकाशयितुं व्युच्छन्ती । एतादृशी उषाः, \textbf{विश्वा भुवनानि}= सर्वाणि जगन्ति \textbf{अजीगः}= उद्गीरति प्रकाशयतीत्यर्थः।

\customparagraph{\textbf{भावार्थः}}

उषा देवी केषाञ्चित् धनम्, केषाञ्चित् अन्नम्, केषाञ्चित् यज्ञीयकार्यम्, केषाञ्चित् अपेक्षितं कार्यं प्रति व्युच्छन्ती भवति । सा प्रणिनां जीवनरक्षणाय विविधानि जीविकासम्बद्धानि कार्याणि प्रकाशयितुं व्युच्छन्ती भवति । एतादृशी उषा सर्वाणि भुवनानि अन्धकारात् प्रकाशेन उद्गीर्णानीव प्रकाशमयानि करोति ।

\customparagraph{\textbf{नेपाली भावार्थ}}

उषा देवी कसैलाई धनका लागि कसैलाई अन्नका लागि कसैलाई यज्ञविधानका लागि कसैलाई ईप्सित कर्म गर्नका लागि उदाउँछिन् अर्थात् उषा उदय भएपछि मानिसहरू आ-आफ्नो कर्ममा लाग्दछन् । ती उषाले मानिसका लागि विभिन्न प्रकारका जीविकाका माध्यमलाई प्रकाशित गर्दछिन् । यस्ती उषाले अन्धकारमा डुबेको भुवनलाई दर्शनीय बनाउँछिन् ।

\begin{visheshshloka}[center]

\textbf{एषा दिवो दुहिता प्रत्यदर्शि व्युच्छन्ती युवतिः शुक्रवासाः ।}

\textbf{विश्वस्येशाना पार्थिवस्य वस्व उषो अद्येह सुभगे व्युच्छ ॥७॥}

\end{visheshshloka}

\customparagraph{\textbf{अन्वयः}}

दिवो दुहिता एषा युवतिः शुक्रवासाः विश्वस्य पार्थिवस्य वस्वः ईशाना व्युच्छन्ती (उषा) प्रति अदर्शि । हे सुभगे ! उषः ! अद्य इह व्युच्छ।

\customparagraph{\textbf{शब्दार्थः}}

\textbf{दिवः}= व्योम्नः \textbf{दुहिता}= दुहितृरूपा \textbf{एषा}= अद्यतनी उषा \textbf{युवतिः}= नित्ययौवनयुता \textbf{शुक्रवासाः}= श्वेतवस्त्रयुता \textbf{विश्वस्य}= पृथिव्या \textbf{वस्वः}= धनस्य \textbf{ईशाना}= ईश्वरी \textbf{व्युच्छन्ती}= तमो वर्जयन्ती \textbf{प्रत्यदर्शि}= सर्वैः प्राणिभिर्दृश्यमाना जाता । \textbf{हे सुभगे} != शोभनधने उषः ! तादृशी त्वम्, \textbf{अद्य}= अस्मिन् काले \textbf{इह}= अस्मिन् देवयजनदेशे \textbf{व्युुच्छ}= तमांसि विवासय, दूरीकुरु ।

\customparagraph{\textbf{भावार्थः}}

आकाशस्य देवी सर्वदा यौवनसौन्दर्ययुता श्वेतवस्त्रधारिणी जगतो धनस्य स्वामिनीपुत्री उषा अन्धकारमपसारयन्ती सर्वेषां प्राणिनां दृष्टिपथि समागता । सौभाग्यशालिनि हे उषे !  अद्य अस्मिन् यज्ञस्थलेऽपि अन्धकारमपसारय ।

\customparagraph{\textbf{नेपाली भावार्थ}}

आकाशकी छोरी, सर्वदा जबानीको सौन्दर्यले युक्त भएकी, सेतो कपडाले लगाएकी, जगत्‌मा रहेको सम्पूर्ण सम्पत्तिकी स्वामिनी उषा देवी अन्धकारलाई हटाउँदै सबै प्राणीका अगाडि उदाइन् । त्यस्ती सौभाग्यशाली हे उषा ! यस यज्ञभूमिको पनि अन्धकारलाई हटाऊ ।

\begin{visheshshloka}[center]

\textbf{परायतीनामन्वेति पाथ आयतीनां प्रथमा शश्वतीनाम् ।}

\textbf{व्युच्छन्ती जीवमुदीरयन्त्युषा मृतं कं चन बोधयन्ती ॥ ८॥}

\end{visheshshloka}

\customparagraph{\textbf{अन्वयः}}

परायतीनां पाथः अनु एति । आयतीनां शश्वतीनां प्रथमा व्युच्छन्ती जीवम् उदीरयन्ती उषाः मृतं कं चन बोधयन्ती (आगात्) ।

\customparagraph{\textbf{शब्दार्थः}}

\textbf{परायतीनाम्}= परागच्छन्तीनाम्, अतीतानामुषसामित्यर्थः । \textbf{पाथः}= अन्तरिक्षैकलक्षणं स्थानम्,  अद्यतनी उषा \textbf{अनु एति}= अन्वेति, अनुगच्छति। \textbf{आयतीनाम्}= आगच्छन्तीनाम्, \textbf{शश्वतीनाम्}= बह्वीनामुषसाम् \textbf{प्रथमा}= आद्या, \textbf{व्युच्छन्ती}= तमो वर्जयन्ती \textbf{जीवम्}= जीवात्मानम्, \textbf{उदीरयन्ती}= शयनात् उत्थातुं प्रेरयन्ती \textbf{उषाः =} देवी \textbf{मृतम्}= श्वापसमये प्रलीनेन्द्रियत्वात् मृतमिव सन्तम् \textbf{कं चन}= कमपि जनम्, \textbf{बोधयन्ती}= सचेतनं कुर्वती, समागतवती ।

\customparagraph{\textbf{भावार्थः}}

अद्यतनी एषा उषा प्राग्जातानामुषसाम् अन्तरिक्षमार्गेऽनुगच्छति । आगामिनीषु बह्वीषु उषस्सु प्रथमा इयम् अद्यतनी उषा अन्धकारमपसारयन्ती प्राणिनः शयनादुत्थापयन्ती सती समागतवती ।

\customparagraph{\textbf{नेपाली भावार्थ}}

अहिलेकी यी उषा पहिलेका उषाहरू हिँडेको आकाश मार्गमा पछिपछि हिँड्छिन् । पछि आउने अरू उषाहरूभन्दा पहिली, अहिलेकी यी उषा अन्धकारलाई हटाउँदै, इन्द्रियहरू निष्क्रिय हुनाले मरेजस्तै भएर सुतेका मानिसमा चेतना भर्दै आएकी छिन् ।

\begin{visheshshloka}[center]

\textbf{उषो यदग्निं समिधे चकर्थ वि यदावश्चक्षसा सूर्यस्य ।}

\textbf{यन्मानुषान् यक्षमाणाँ अजीगस्तद्देवेषु चक्रिषे भद्रमप्नः ॥९॥}

\end{visheshshloka}

\customparagraph{\textbf{अन्वयः}}

हे उषः ! (त्वम्) अग्निंं समिधे यत् चकर्थ । सूर्यस्य चक्षसा यत् वि आवः । मानुषान् यक्ष्यमाणान् यत् अजीगः । देवेषु भद्रं तत् अप्नः चकृषे । \textbf{शब्दार्थः}

हे उषः ! त्वम् \textbf{अग्निम्}= गार्हस्पत्यादिस्वरूपम् \textbf{समिधे}= प्रज्ज्वलनाय यत् \textbf{चकर्थः}= कृतवती, उषाकाले ह्यग्नयो हवनार्थमुपसमिध्यन्ते । \textbf{सूर्यस्य}= रवेः, \textbf{चक्षसा}= तेजसा यत् \textbf{वि आवः}= व्यावृणोः तमसा विश्लिष्टमकरोः । \textbf{मानुषान्=} मनोः पुत्रान् मनुष्यान् \textbf{यक्ष्यमाणान्=} यागं करिष्यतस्त्वं यत् \textbf{अजीगः}= प्राक् तमसा ग्रस्तान् प्रकाशेन उद्गीर्णानिवाकरोः । हे उषः ! \textbf{देवेषु}= इन्द्रादिदेवेषु त्वमेव \textbf{भद्रम्}= भजनीयम् \textbf{तत्}= एतत् विविधं \textbf{अप्नः}= कर्म \textbf{चकृषे}= कृतवती ।

\customparagraph{\textbf{भावार्थः}}

हे उषः ! त्वं गार्हस्पत्यादयो येऽग्नयः सन्ति तेषां प्रज्ज्वलनाय जनान् प्रेरयसि । सूर्यस्य किरणस्य माध्यमेन रात्रिजातं निविडम् अन्धकारं दूरीकरोसि । यज्ञकर्तॄन् जनान् शयनादुत्थाप्य प्रेरयसि । हे उषादेवि ! त्वं देवान् समुद्दिश्य विविधान् कार्यान् करोसि ।

\customparagraph{\textbf{नेपाली भावार्थ}}

हे उषा देवी !  तिमी सूर्यको किरणका  सहायताले अन्धकारलाई हटाउँछ्यौ । सुतेका यज्ञकर्तालाई सबेरै उठाउँछ्यौ । यज्ञकर्मका लागि अग्नि प्रज्ज्वलन गर्न प्रेरित गर्दछ्यौ । यसरी तिमी देवताका लागि आवश्यक विभिन्न कार्यहरू गर्दछ्यौ ।

\begin{visheshshloka}[center]

\textbf{कियात्या यत्समया भवाति या व्यूषुर्याश्च नूनं व्युच्छान् ।}

\textbf{अनु पूर्वाः कृषते वावशाना प्रदीध्याना जोषमन्याभिरेति ॥१०॥}

\end{visheshshloka}

\customparagraph{\textbf{अन्वयः}}

(उषाः) समया भवाति यत् कियति आ ? याः व्यूषुः नूनं याः च व्युच्छान् (तत्र) पूर्वा वावशानाः अनु कृषते । प्रदीध्यानाः अन्याभि जोषम् एति ।

\customparagraph{\textbf{शब्दार्थः}}

उषाः सम\textbf{या भवाति}= समीपस्था भवति । \textbf{यत्}= समयः स्थानञ्च \textbf{कियति आ}= कति भवति, आ इति प्रश्नार्थः । अर्थात् ज्ञातुं न शक्यते, एतेन उषसः अनन्तत्त्वं सूच्यते । पुरा \textbf{याः}= उषसो \textbf{व्यूषुः}= व्युष्टाः सञ्जाताः (प्रकाशिताः) \textbf{नूनम्=} अवश्यम् इतः परम् \textbf{याः च=} उषसो \textbf{व्युच्छान्}= व्युच्छन्ति, व्युष्टा भविष्यन्ति (प्रकाशयिष्यन्ति) । तत्र \textbf{पूर्वा}= अतीताः उषसो \textbf{वावशाना}= कामयमाना इदानीं वर्तमाना उषा \textbf{अनुकृषते}= अनुकल्पते । अतीता उषसो यथा प्रकाशमकुर्वन् तद्वत् अधुनातनी उषा अपि प्रकाशं करोतीत्यर्थः । एवमेव \textbf{प्रदीध्याना}= प्रकर्षेण दीप्यमाना, उषा \textbf{अन्याभिः}= आगमिनीभिः उषोभिः, \textbf{जोषम्}= सह \textbf{एति=} सङ्‌गच्छते । आगामिन्यः उषसः अपि एतदीयं प्रकाशमनुकुर्वन्तीत्यर्थः।

\customparagraph{\textbf{भावार्थः}}

इयम् उषा कति समयं यावत्, कति परिमितं स्थानं वा तिष्ठतीति न ज्ञायते, एतेन उषाया अनन्तत्त्वं बोध्यते । प्राग्‌जाता उषाः प्रकाशिता जाता, भविष्यत्काले जायमाना उषा अपि प्रकाशिता भविष्यन्ति । तत्र अतीता उषा यथा प्रकाशमनुकुर्वन् तथैव अधुना जाता उषा प्रकाशं करोति । एवमेव आगमिन्य उषा एतदीयं प्रकाशमनुकरिष्यन्तीति भावः ।

\customparagraph{\textbf{नेपाली भावार्थ}}

अहिले उदाएकी उषाको समय र स्थान कति हुन्छ ? बुझ्न गारो छ। यसबाट उषाको अनन्तता बुझिन्छ । आजभन्दा पहिले पनि उषा उदाएकी थिइन् । भविष्यमा पनि उषा उदाउनेछिन्। त्यस अवस्थामा पहिले उदाएकी उषाको अनुकरण अहिलेकी उषाले गरेकी छन् र पछि उदय हुने उषाले अहिलेकी उषाको अनुकरण गर्नेछिन् ।

\begin{visheshshloka}[center]

\textbf{ईयुष्टे ये पूर्वतरामपश्यन् व्युच्छन्तीमुषसं मर्त्यासः ।}

\textbf{अस्माभिरू नु प्रतिचक्ष्याभूदो ते यन्ति ये अपरीषु पश्यान् ॥११॥}

\end{visheshshloka}

\customparagraph{\textbf{अन्वयः}}

ये मर्त्यासः व्युच्छन्तीम् पूर्वतराम् उषसम् अपश्यन् ते ईयुः । अस्माभिः नु प्रतिचक्ष्या अभूत्। अपरीषु ये पश्यान् ते ऊ ओ यन्ति ।

\customparagraph{\textbf{शब्दार्थः}}

\textbf{ये मर्त्यासः}= मरणधर्मयुता मनुष्याः, \textbf{व्युच्छन्तीम्}= विवासयन्तीम् \textbf{पूर्वतराम्}= अतिशयेन पूृर्वाम्, प्राचीनामित्यर्थः, \textbf{उषसम्}= उषादेवीम् \textbf{अपश्यन्}= दृष्टवन्तः, \textbf{ते}= तादृशा मनुष्याः, \textbf{ईयुः}= गतवन्तः । एवमेव \textbf{अस्माभिः}= अस्मिन् समये सञ्जातैर्जनैरपि प्र\textbf{तिचक्ष्या=} प्रकर्षेण द्रष्टव्या \textbf{अभूत्=} सञ्जाता । एवमेव \textbf{अपरीषु}= भाविनीषु रात्रिषु  \textbf{ये}= मनुष्या एतामेव उषां \textbf{पश्यान्=} द्रक्ष्यन्तीत्यर्थः । \textbf{ते}= मनुष्या \textbf{यन्ति=} आयास्यन्ति एव । अत्र \textbf{उ ओ}= इति निपातद्वयम् । अवधारणार्थे प्रयुक्तं विद्यते । एवमुषसः आनन्त्यं द्योत्यते ।

\customparagraph{\textbf{भावार्थः}}

ये मनुष्याः प्रकाशवतीं प्राचीनाम् उषाम् अवलोकितवन्तः, ते पूर्वजा गताः । वर्तमाने समये सञ्जाता वयमपि प्रत्यक्षरूपेण प्रकाशवतीं उषाम् अवलोकयन्तः स्मः । एवमेव ये जना भविष्यति काले समागमिष्यन्ति, तेऽपि उषाम् अवलोकयिष्यन्ति । अतीतैर्वर्तमानैरागामिभिश्च सर्वैर्मनुष्यैरुषा दृश्यते नास्मदादिवत् कैश्चिदेव इति समस्तार्थः ।

\customparagraph{\textbf{नेपाली भावार्थ}}

हामीभन्दा पहिले जन्मिएका हाम्रा पूर्वजहरूले प्रकाशले युक्त भएकी उषालाई देखे ती गए। अहिले हामीले प्रकाशमय उषालाई देखिरहेका छौं । त्यस्तै भविष्यमा जन्मिने मानिसहरूले पनि उषालाई देख्नेछन् । अर्थात् सुनौलो प्रकाश लिएर प्रभात कालमा आउने उषा सार्वकालिक वा शाश्वत छिन् । यिनले सदावहार रात्रिको अन्धकार हटाएर सुनौलो बिहानी ल्याइरहनेछिन् ।

\begin{visheshshloka}[center]

\textbf{यावयद्‌द्वेषा ऋतपा ऋतेजाः सुम्नावरी सूनृता ईरयन्ती ।}

\textbf{सुमङ्गलीर्बिभ्रती देववीतिमिहाद्योषः श्रेष्ठतमा व्युच्छ ॥१२॥}

\end{visheshshloka}

\customparagraph{\textbf{अन्वयः}}

हे उषः ! यावयद्‌द्वेषाः ऋतपाः ऋतेजाः सूम्नावरी सूनृता ईरयन्ती सुमङ्गलीः देववीतिं बिभ्रती श्रेष्ठतमा (त्वम्) इह अद्य व्युच्छ ।

\customparagraph{\textbf{शब्दार्थः}}

हे उषः ! \textbf{यावद्‌द्वेषाः}= यावयन्ति अस्मत्तः पृथक्‌कृतानि द्वेषांसि द्वेषयुतानि राक्षसादीनि यया सा तथोक्ता, न ह्युषसि जातायां राक्षसादयोऽवतिष्ठन्ते तस्मात्ते निशाचराः । \textbf{ऋतपाः}= ऋतस्य सत्यस्य यज्ञस्य वा पालयित्री, \textbf{ऋतेजाः=} यज्ञार्थं प्रादुर्भूता, सत्यामुषायाम् अहनि यागा अनुष्ठीयन्ते, अतो यज्ञार्थं जातेत्युच्यते । \textbf{सुम्नावरी}= सुम्नमिति सुखनाम तद्वती, \textbf{सूनृता=} वाङ्‌नामैतत् । पशुपक्षिमृगादीनां वचांसि, \textbf{ईरयन्ती=} प्रेरयन्त्युत्पादयन्ती \textbf{सुमङ्गलीः}= सौमङ्गल्योपेता पत्या कदाचिदपि न वियुक्तेत्यर्थः । \textbf{देववीतिम्}= देवैः काम्यमानं यज्ञम्, \textbf{बिभ्रती}= धारयन्ती, \textbf{श्रेष्ठतमा}= उक्तेन प्रकारेण ऐश्वर्ययुता त्वम्, \textbf{इह}= अस्मिन् देवयजनदेशे, \textbf{अद्य}= अस्मिन् यागसमये, \textbf{व्युच्छ}= विवासय ।

\customparagraph{\textbf{भावार्थः}}

हे उषे ! रजन्या अन्धकारमपहृत्य अस्माभिः सह द्वेष्टॄन् राक्षसान् दूरीकरणे समर्था, यज्ञस्य पालयित्री, यज्ञार्थे समुदिता, शयनान्ते पशुपक्ष्यादिप्राणिनां शब्दकरणे प्रेरयित्री, मङ्गलमयी देवकार्यं बिभ्रती, स्वात्मना कृत्यैः कार्यैरतिप्रशस्ता त्वम् अत्र यज्ञस्थलमपि प्रकाशय ।

\customparagraph{\textbf{नेपाली भावार्थ}}

हे उषादेवी ! हाम्रा शत्रु राक्षसादि निशाचरलाई हामीबाट अगल गराउने, अर्थात् निशाचरको समय रात्रि हटाएर प्रभात काल ल्याउने, यज्ञको पालन वा संरक्षण गर्ने, यज्ञकै लागि उत्पन्न भएकी वा उदाएकी, पशुपक्षी आदि प्राणीलाई बिहानै बोल्न प्रेरणा दिने, मङ्गलप्रदान गर्ने, देवताहरूका लागि यज्ञलाई धारण गर्ने, उत्कृष्ट ऐश्वर्यशाली तिमी आज यस यज्ञस्थलमा पनि प्रकाशित होऊ ।

\begin{visheshshloka}[center]

\textbf{शश्वत्पुरोषा व्युवास देव्यथो अद्येदं व्यावो मघोनी ।}

\textbf{अथो व्युच्छादुत्तराँ अनु द्यू नजरामृता चरति स्वधाभिः ॥१३॥}

\end{visheshshloka}

\customparagraph{\textbf{अन्वयः}}

देवी उषाः पुरा शश्वत् व्युवास । अथो अद्य मघोनी इदं व्यावः । अथो उत्तरान् द्यून् अनु व्युच्छात् । अजरा अमृता स्वधाभिः चरति ।

\customparagraph{\textbf{शब्दार्थः}}

देवी उषाः \textbf{पूरा=} पूर्वस्मिन् काले \textbf{शश्वत्=} सततम् \textbf{व्युवास}= व्यौच्छत्, अन्धकारं दूरीकृतवतीत्यर्थः । \textbf{अथो}= अनन्तरम् \textbf{अद्य}= अस्मिन् काले \textbf{मघोनी}= धनवती उषा तमसा तिरोहितम् \textbf{इदं}= सर्वं जगत् \textbf{व्यावः}= विवासितम्, प्रकाशेन तमसा वियुक्तम् अकरोत् । \textbf{अथो}= अनन्तरम् \textbf{उत्तरान्}= ऊर्ध्वतरान् भाविनः \textbf{द्यून्}= दिवसान् अनुलक्ष्य आगामिष्वपि दिवसेषु \textbf{व्युच्छात्}= व्युच्छति, विवासयति दूरीकरोति ।  अतः कालत्रयव्यापिनी सा उषा, \textbf{अजरा}= वृद्धभावरहिता सर्वदैकरूपा इत्यर्थः । \textbf{अमरा}= मरणरहिता सति \textbf{स्वधाभिः}= आत्मीयैस्तेजोभिः सह \textbf{चरति}= भ्रमति ।

\customparagraph{\textbf{भावार्थः}}

दीप्तिमती उषा प्राचीनकाले रजन्या अन्धकारं दूरीचकार। अधुनाऽपि अन्धकारमपहृत्य जगत् प्रकाशयति । एवमेव आगामिष्वपि दिवसेषु अन्धकारं दूरीकरिष्यत्येव । एतादृशी कालत्रयव्यापिनी उषा अजरा अमरा च सती स्वकीयैः किरणैः सह निरवच्छिन्नरूपेण भ्रमति ।

\customparagraph{\textbf{नेपाली भावार्थ}}

शीतलप्रकाशकी धनी उषाले प्राचीन कालमा रात्रिबाट सिर्जित अन्धकार हटाइन् । अहिले पनि त्यस्तो अन्धकार हटाई नै रहेकी छिन् । एवम् भविष्यमा पनि रात्रिजनित अन्धकारलाई हटाउनेछिन् । यसरी भूत, वर्तमान र भविष्यमा फैलिएर रहेकी उषा अजर र अमर भई आफ्ना किरणका साथमा निरन्तर जगत्‌को भ्रमण गरिरहन्छिन् ।

\begin{visheshshloka}[center]

\textbf{व्य१ञ्जिभिर्दिव आतास्वद्यौदप कृष्णां निर्णिजं देव्यावः ।}

\textbf{प्रबोधयन्त्यरुणेभिरश्वैरोषा याति सुयुजा रथेन ॥१४॥}

\end{visheshshloka}

\customparagraph{\textbf{अन्वयः}}

दिवः आतासु अञ्जिभिः वि अद्यौत् । देवी कृष्णां निर्णिजम् अप आवः । अरुणेभिः अश्वैः सुयुजा रथेन प्रबोधयन्ती उषाः आ याति ।

\customparagraph{\textbf{शब्दार्थः}}

\textbf{दिवः}= नभसः सम्बन्धिनीषु \textbf{आतासु}= विस्तीर्णासु दिक्षु, उषा \textbf{अञ्जिभिः}= प्रकाशकैः किरणैः, वि \textbf{अद्यौत्}= विशेषेण प्रकाशते । सैषा \textbf{देवी}= उषा कृष्णां निर्णिजम्= रात्रिकृतं कृष्णं रूपम्, \textbf{अप आव}= अपावृणोत्, दूरीकृतवती । \textbf{अरुणेभिः}= अरुणैः लोहितवर्णैः \textbf{अश्वैः}= व्यापनशीलैः स्वकीयैः कीरणैस्तुरगैर्वा, \textbf{सुयुजा}= सम्यक् रूपेण योजितेन रथेन युक्ता, \textbf{प्रबोधयन्ती}= सुप्तान् सर्वान् प्राणिनः प्रबुद्धान् कुर्वती, \textbf{उषा}= देवी \textbf{आ याति}= आगच्छति ।

\customparagraph{\textbf{भावार्थः}}

अन्तरिक्षसम्बद्धासु विस्तृतासु सर्वासु दिक्षु  उषाः स्वकीयैः किरणैः प्रकाशते । तादृशी उषा कृष्णवर्णमन्धकारं दूरीकरोति । अरुणवर्णयुतैः किरणैरश्वैर्वा युक्तेन रथेन युता सा उषा निशान्ते प्राणिनः प्रबोधयन्ती समुदेति । निशान्ते उषा दिक्षु स्थितं कृष्णवर्णात्मकमन्धकारमपहृत्य प्राणिनः प्रबोध्य उदयं गच्छतीत्यर्थः ।

\customparagraph{\textbf{नेपाली भावार्थ}}

उषा देवी अन्तरिक्षका दिशाहरूलाई उज्यालो बनाउँदै प्रकाशित हुन्छिन् । रात्रिको कालो अन्धकारलाई हटाउँछिन् । अरुण रङका किरण वा घोडाले युक्त रथले सहित भएर ती उषा प्राणीहरूलाई ब्युँझाउँदै उदाउँछिन् ।

\begin{visheshshloka}[center]

\textbf{आवहन्ती पोष्या वार्याणि चित्रं केतुं कृणुते चेकिताना ।}

\textbf{ईयुषीणामुपमा शश्वतीनां विभातीनां प्रथमोषा व्यश्वैत् ॥१५॥}

\end{visheshshloka}

\customparagraph{\textbf{अन्वयः}}

पोष्या वार्याणि आवहन्ती चेकिताना उषा चित्रं केतुं कृणुते । ईयुषीणां शश्वतीनां उपमा विभातीनां प्रथमा वि अश्वैत् ।

\customparagraph{\textbf{शब्दार्थः}}

\textbf{पोष्या}= यावज्जीवं पोषणसमर्थानि \textbf{वार्याणि}= वरणीयानि धनानि \textbf{आवहन्ती}= अस्मभ्यमानयन्ती \textbf{चेकिताना}= सर्वं जनं प्रज्ञापयन्ती उषा \textbf{चित्रम्}= विचित्रमाश्चर्यजनकम् चायनीयम् वा \textbf{केतुम्}= प्रज्ञापकम् रश्मिम्, कृत्स्नजगत्प्रकाशनसमर्थम्, \textbf{कृणुते}= स्वात्मनः प्रकाशात् कुरुते। सैषा \textbf{ईयुषीणाम्}= गमनवतीनां पूर्वनिष्पन्नानाम् \textbf{शश्वतीनाम्}= बह्वीनामुषसाम्,  \textbf{उपमा}= सादृश्ययुता \textbf{विभातीनाम्}= विशेषेण प्रकाशमानानामागामिनीनाम् उषसां \textbf{प्रथमा}= आद्या उषा \textbf{व्यश्वैत्}= तेजसा प्रवृद्धासीत् ।

\customparagraph{\textbf{भावार्थः}}

प्राणिनां जीवनपोषणाय धनानि आनयन्ती उषा सर्वान् जनान् प्रबोधयति । विचित्रं प्रकाशमयं केतुं धारयति । प्राग्जातानामागमिनीनां च उषाणां सादृश्यभूता, प्रकाशमानानामागमिनीनां उषाणां प्रथमा इयम् उषा किरणैः प्रवृद्धासीत् ।

\customparagraph{\textbf{नेपाली भावार्थ}}

हाम्रो जिवनयापनका लागि आवश्यक पोषणरूपी सम्पत्ति लिएर आउने र सबै प्राणीलाई उद्बोधित गराउने उषा आफ्ना आश्चर्यजनक किरणहरूद्वारा सम्पूर्ण जगत्लाई प्रकाशित गर्न समर्थ छन् । पहिलेका अनेकौँ उषाहरूको सादृश्यका रूपमा रहेकी (पहिलेका उषाजस्ती) वर्तमान समयकी उषा प्रकाशशील आगामी उषाभन्दा पहिली र समृद्ध प्रकाशले युक्त छिन् ।

\begin{visheshshloka}[center]

\textbf{उदीर्ध्वं जीवो असुर्न आगादप प्रागात्तम आ ज्योतिरेति ।}

\textbf{आरैक्पन्थां यातवे सूर्यायागन्म यत्र प्रतिरन्त आयुः ॥१६॥}

\end{visheshshloka}

\customparagraph{\textbf{अन्वयः}}

(हे मनुष्याः !) उदीर्ध्वम् नः असुः जीवः आगात् । तमः अप प्रागात् । ज्योतिः आ एति । सूर्याय यातवे पन्थाम् आरैक् । अगन्म यत्र आयुः प्रतिरन्ते ।

\customparagraph{\textbf{शब्दार्थः}}

हे मनुष्याः ! \textbf{उदीर्ध्वम्}= शयनं परित्यज्य उत्तिष्ठत । \textbf{नः}= अस्माकम्, \textbf{असुः}= शरीरस्य प्रेरयिता, \textbf{जीवः}= जीवात्मा \textbf{आगात्}= आगतवान् । \textbf{तमः}= तिमिरम् अप \textbf{प्रागात्}= अपक्रान्तम्  दूरीभूतमित्यर्थः । \textbf{ज्योतिः}= प्रकाशः आ \textbf{एति}= आगच्छति । \textbf{सूर्याय}= भानवे \textbf{यातवे}= गमनाय \textbf{पन्थाम्}= मार्गम् \textbf{आरैक्}= विविक्तीकरोति, प्रकाशयतीत्यर्थः । तत्र \textbf{आगन्म}= गच्छामः, \textbf{यत्र}= यस्मिन् देशे \textbf{आयुः}= अन्नम्, \textbf{प्रतिरन्ते}= उदारा दानेन प्रवर्धयन्ति ।

\customparagraph{\textbf{भावार्थः}}

रात्रौ सुप्ता हे मनुष्या ! भवन्तः शयनं परित्यज्य उत्तिष्ठन्तु । अस्माकं जीवनस्य प्रेरकः जीवात्मा समागतः । रात्रिजन्यं तिमिरं दूरीभूतम् । प्रकाशवती उषा उदिता । सा भानवे मार्गं प्रकाशयति । वयं सर्वे तत्र गच्छामः, यत्र उदारा जीवनावधिं सोमादिहविर्लक्षणमन्नं वा प्रवर्धयन्ति ।

\customparagraph{\textbf{नेपाली भावार्थ}}

राती सुतेका हे मानिसहरू हो ! ओछ्यान छाडेर उठ । हाम्रो जीवनको प्रेरक जीवात्मा आइसकेको छ । अन्धकार हटेको छ । प्रकाशले पूर्ण उषाको उदय भइसकेको छ । उषाले सूर्य हिँड्ने बाटो प्रकाशित गरिरकेकी छन् । हामी सबै त्यतै जाऔं जहाँ हाम्रो आयु वा यज्ञका लागि अन्नको प्रवर्धन हुन्छ ।

\begin{visheshshloka}[center]

\textbf{स्यूमना वाच उदियर्ति वह्निः स्तवानो रेभ उषसो विभातीः ।}

\textbf{अद्या तदुच्छ गृणते मघोन्यस्मे आयुर्नि दिदीहि प्रजावत् ॥१७॥}

\end{visheshshloka}

\customparagraph{\textbf{अन्वयः}}

वह्निः रेभः उषसो विभातीः स्तवानः वाचः स्यूमना उदियर्ति । हे मघोनि ! अद्य गृणते तत् उच्छ । अस्मे प्रजावत् आयुः नि दिदीहि ।

\customparagraph{\textbf{शब्दार्थः}}

\textbf{वह्निः}= स्तोत्राणां बोढा \textbf{रेभः}= स्तोता \textbf{उषसो विभातीः}= तमसोऽपनोदनेन प्रकाशमानां उषां देवताम्, \textbf{स्तवानः}= स्तुवन् \textbf{वाचः}= वेदरूपायाः सम्बन्धिनी \textbf{स्यूमना}= स्यूमानि अनुस्यूतानि उक्थानि स्तोत्राणि, \textbf{उदियर्ति}= उद्गमयति उच्चारयति । अतो हे \textbf{मघोनि} ! =मघवति उषः ! \textbf{अद्य}= अस्मिन् समये, \textbf{गृणते}= स्तुवते तस्मै पुरुषाय, \textbf{तत्}= दृष्टिनिरोधकतया प्रसिद्धं नैशं तमः, \textbf{उच्छ}= दूरीकुरु ।  \textbf{अस्मे}= अस्मभ्यम् च \textbf{प्रजावत्}= प्रजाभिः पुत्रपौत्रादिभिर्युक्तम्, \textbf{आयुः}= चिरं जीवत्वम्, अन्नं वा \textbf{नि} \textbf{दिदीहि}= नितरां प्रकाशय, प्रयच्छेत्यर्थः ।

\customparagraph{\textbf{भावार्थः}}

स्तुतिवचनानां धारकः स्तुतिकर्ता उषां स्तुवन् वेदमन्त्ररूपं स्तुतिवचनं उच्चारयति । अतो हे धनवति उषे ! अद्य स्तुतिं कुर्वतो जनस्य नेत्रावरोधभूतं तमो दूरीकुरु । अस्मभ्यं पुत्रपौत्रादिभिर्युक्तम् चिरञ्जीवत्वम् अन्नं वा प्रयच्छ ।

\customparagraph{\textbf{नेपाली भावार्थ}}

स्तुतिवचनको राम्रो ज्ञान भएका स्तुतिवाचकहरू उषाको स्तुति गर्दै वेदसँग सम्बन्धित मन्त्रहरू उच्चारण गर्दछन् । हे धनवति उषा ! आज ती स्तुति गर्नेवालाको हेराइलाई अवरोध गर्ने अन्धकारको नाश गर । एवम् हामीलाई पनि छोरानातिसहितको लामो आयु वा अन्न प्रदान गर ।

\begin{visheshshloka}[center]

\textbf{या गोमतीरुषसः सर्ववीरा व्युच्छन्ति दाशुषे मर्त्याय ।}

\textbf{वायोरिव सूनृतानामुदर्के ता अश्वदा अश्नवत्सोमसुत्वा ॥१८॥}

\end{visheshshloka}

\customparagraph{\textbf{अन्वयः}}

दाशुषे मर्त्याय गोमतीः सर्ववीराः याः उषसः व्युच्छन्ति । वायोः इव सूनृतानाम् उदर्के अश्वदा ताः सोमसुत्वा अश्नवत् ।

\customparagraph{\textbf{शब्दार्थः}}

\textbf{दाशुषे}= हवींषि दत्तवते \textbf{मर्त्याय}= मनुष्याय यजमानाय \textbf{गोमतीः}= गोमत्यो बहुभिर्गोभिर्युक्ताः \textbf{सर्ववीराः}= सर्वैः सरणसमर्थैः वीरैः शूरैर्युक्ताः याः \textbf{उषसः}= उषादेव्यः, \textbf{व्युच्छन्ति}= तमो वर्जयन्ति । \textbf{वायोरिव}= वायुवच्छीघ्रं प्रवर्तमानानाम्, \textbf{सूनृतानाम्}= स्तुतिरूपाणां वाचाम्, \textbf{उदर्के}= समाप्तौ \textbf{अश्वदा}= अश्वानां दात्रीस्\textbf{ताः}= उषादेव्यः \textbf{सोमसुत्वा}= सोमानां कर्ता यजमानः, \textbf{अश्नवत्}= व्याप्नोतु, अर्थात् यजमानस्य कृते समागच्छतु इत्यर्थः ।

\customparagraph{\textbf{भावार्थः}}

यज्ञकर्त्रे मनुष्याय यज्ञार्थं गोभिर्युताः, विपरीतान् सर्वान् पराजेतुं शक्तैर्वीरैः समन्विताः, उषादेव्योऽन्धकारान् दूरीकुर्वन्ति । ताः उषादेव्यो वायुरिव तीव्रगत्या पठ्यमानानां स्तोत्राणामन्त्ये

अश्वदात्रीत्वेन यजमानस्य पुरः समागच्छन्तु।

\customparagraph{\textbf{नेपाली भावार्थ}}

यज्ञ गर्ने मनुष्यका लागि गाईहरूलाई साथमा लिएर आएका एवं विपरीत कार्य गर्ने (विरोधी)हरूलाई पराजित गर्ने सामर्थ्यले युक्त वीरहरूलाई पनि साथमा लिएर आएका उषाहरूले अन्धकारलाई हटाउँछन् । त्यस्ता उषाहरू वायुजस्तै तीव्रगतिमा पढिने स्तोत्रको अन्त्यमा अश्व आदि दक्षिणास्वरूप दान गर्ने पदार्थका साथमा यजमानका अगाडि प्रकट होऊन् ।

\begin{visheshshloka}[center]

\textbf{माता देवानामदितेरनीकं यज्ञस्य केतुर्बृहती वि भाहि ।}

\textbf{प्रशस्तिकृद् ब्रह्मणे नो व्युच्छा नो जने जनय विश्ववारे ॥१९॥}

\end{visheshshloka}

\customparagraph{\textbf{अन्वयः}}

देवानां माता, अदितेः अनीकम् यज्ञस्य केतुः बृहती वि भाहि । प्रशस्तिकृत् नः ब्रह्मणे व्युच्छ। हे विश्ववारे !  नः जने आ जनय ।

\customparagraph{\textbf{शब्दार्थः}}

(हे उषः ! त्वम्) \textbf{देवानाम्}= अमराणाम् \textbf{माता}= जननी असि, (उषायां सर्वे देवाः स्तुत्या प्रबोध्यन्ते अतः सा तज्जननीत्युच्यते) । अतः \textbf{अदितेः} = देवानां मातुः, \textbf{अनीकम्}= प्रत्यनीकं प्रतिस्पर्धिनी त्वमित्यर्थः । \textbf{यज्ञस्य}= यागकर्मणः \textbf{केतुः}= ज्ञापयित्री \textbf{बृहती}= महती सती वि \textbf{भाहि}= प्रकाशयस्व । \textbf{प्रशस्तिकृत्}= सम्यक् स्तुतमिति प्रशंसनं कुर्वती \textbf{नः=} अस्मदीयाय \textbf{ब्रह्मणे}= मन्त्ररूपाय स्तोत्राय \textbf{व्युच्छ}= विवासय, अज्ञानरूपं तिमिरमपनय। \textbf{हे विश्ववारे} ! सर्वैर्वरणीये उषः !, \textbf{नः}= अस्मान् \textbf{जने}= जनपदे \textbf{आ जनय=} आभिमुख्येन प्रादुर्भावय अवस्थापयेत्यर्थः ।

\customparagraph{\textbf{भावार्थः}}

हे उषः ! त्वं देवानां जननी अस्ति, तस्मात् देवमातुरदितेः प्रतिस्पर्धिनी, यज्ञस्य सूचयित्री, महती सती प्रकाशयस्व । तव स्तुतिं कुर्वन्तं जनं सम्यक् स्तुतमिति प्रशंसनं कुर्वती सती सम्यक् रूपेण स्तोत्रपाठाय अन्तःस्थितमज्ञानमयं तिमिरमपनय । सर्वैः सम्मानिता उषा, त्वम् अस्मान् रम्ये जनपदे स्थापय ।

\customparagraph{\textbf{नेपाली भावार्थ}}

हे उषा ! (उषाकालमा स्तुतिद्वारा उठाइने हुनाले) तिमी देवाताकी जननी हौ । त्यसैले तिमी देवमाता अदितिसँग प्रतिस्पर्धा गर्न समर्थ छौ । यज्ञलाई सूचित गर्ने तिमी महान् विशेषताले युक्त भएर प्रकाशित होऊ । तिम्रा स्तुति गर्नेहरूलाई ‘राम्रो स्तुति गर्‍यौ’ भनी प्रशंसा गर्दै तिमी शुद्ध स्तुतिमन्त्र उच्चारणका लागि हाम्रो अज्ञानरूपी अन्धकारलाई हटाऊ । हे सबैद्वारा मानिएकी उषा !  हामीलाई पवित्र जनपदमा स्थापित गर ।

\begin{visheshshloka}[center]

\textbf{यच्चित्रमप्न उषसो वहन्तीजानाय शशमानाय भद्रम् ।}

\textbf{तन्नो मित्रो वरुणो मामहन्तामदितिः सिन्धुः पृथिवी उत द्यौः ॥२०॥}

\end{visheshshloka}

\customparagraph{\textbf{अन्वयः}}

चित्रम् अप्नः यत् उषसः वहन्ति (तत्) ईजानाय शशमानाय भद्रं (भवति) तत् नः मित्रः वरुणः अदितिः सिन्धुः पृथिवी उत द्यौः मामहन्ताम् ।

\customparagraph{\textbf{शब्दार्थः}}

\textbf{चित्रम्}= चायनीयम्, महनीयम् \textbf{अप्नः}= आप्तव्यम् \textbf{यत्}= धनम् \textbf{उषसः}= उषादेव्यः, \textbf{वहन्ति}= आनयन्ति, तद् धनम्  \textbf{ईजानाय}= हविर्भिः इष्टवते \textbf{शशमानाय}= स्तुतिभिः संभजमानाय पुरुषाय \textbf{भद्रम्}= भजनीयं भवतीति शेषः । यदनेन सूक्तेनास्माभिः प्रार्थितम्, \textbf{तत्}= धनम्, \textbf{नः}= अस्मान्  \textbf{मित्रः}= सूर्यः, \textbf{वरुणः}= जलदेवः, \textbf{अदितिः}= देवमाता, \textbf{सिन्धुः}= समुद्रदेवः, \textbf{पृथिवी}= भूमाता, \textbf{द्युः}= द्युदेवश्चेति षड्‌देवता \textbf{मामहन्ताम्=}  स्वीकुर्वन्तु, दातुमुषसं प्रेरयन्तु इति।

\customparagraph{\textbf{भावार्थः}}

महनीयं प्राप्तव्यञ्च यत् धनम् उषा वहन्ती आगच्छन्ति, तत् धनं यज्ञकर्त्रे, स्तुतिकर्त्रे च यजमानाय कल्याणमयं भवति । (तत्तस्मै दातव्यम्) तादृशं धनम् अस्माकं कृते मित्रादयः षड्‌देवताः, स्वीकुर्वन्तु अर्थात् दातुमुषसं प्रेरयन्तु ।

\customparagraph{\textbf{नेपाली भावार्थ}}

अत्यन्तै महत्त्वपूर्ण र पाउनुपर्ने जुन धन उषादेवी आफ्ना साथमा लिएर आउँछिन्, त्यो धन यज्ञ गर्ने र स्तुति गर्ने यजमानकाे कल्याणका लागि हुन्छ । त्यस्तो कल्याणमय धन हाम्रा लागि मित्रादि छजना देवताले हामीलाई दिनका लागि स्वीकृति दिऊन् । अर्थात् यी देवताले हामीलाई धन दिन उषालाई प्रेरित गरून् ।

\begin{visheshshlokatwo}[center]{84}

\textbf{उषस्सूक्तं समाप्तम्}

\end{visheshshlokatwo}

\specialkhanda{\textbf{ऋग्वेदीयं पुरूरवा-उर्वशी-संवाद-सूक्तम्}}

\customparagraph{\textbf{संवादसूक्तपरिचयः}}

सुष्ठु=शोभनम्, उक्तम्= कथनं सूक्तम्, संवादाश्रितं सूक्तम्, संवादसूक्तमिति शब्दो निष्पद्यते । यानि सूक्तानि संवादाश्रितानि तानि संवादसूक्तानि कथ्यन्ते । लौकिकसाहित्यविकासबीजं वेदेषु प्राप्यते । तेषु ऋग्वेदे विविधविषयसम्बद्धानि  संवादसूक्तानि लभ्यन्ते । एतानि सूक्तानि नाट्यकाव्यस्याधाभूतानि मन्यन्ते । तत्र वर्णितयोः पात्रयोः संवाद एव संवादमूलस्य दृश्यकाव्यस्य प्रारम्भिकस्वरूपं प्राप्यते । ऋग्वेदे यमयमीसंवादः, सोमसूर्यासंवादः, इन्द्रवरुणसंवादः, पुरूरवा-उर्वशीसंवादप्रभृतयः संवादाः प्राप्यन्ते । तेेषु पुरूरवा-उर्वशीसंवादः शृङ्गाररसप्रधानो विद्यते । यमाधृत्य उत्तरवर्तिभिः कविभिर्नैकानि काव्यानि रचितानि दृश्यन्ते । तस्माद् वेदेषु स्थितानि संवादसूक्तानि लौकिकसाहित्यस्य कृते आधाररूपाणि स्वीक्रियन्ते ।

\customparagraph{\textbf{पुरूरवा-उर्वशीसूक्तम्}}

विपुलज्ञानगरिमगुम्फितस्य पौरस्त्यसाहित्यस्य समृद्धपरम्पराया आधार वैदिकसाहित्यमस्ति । संहिताब्राह्मणारण्यकोपनिषदादयो वैदिकग्रन्था लौकिकसंस्कृतसाहित्यस्य उपजीव्यभूताः सन्ति । तत्रस्थान् सूक्तभेदान् आधृत्य बहवो दृश्यश्रव्यकाव्यानि रचितानि । तेषु ‘पुरूरवा-उर्वशीसंवादः’ इत्याख्यं संवादसूक्तमपि वैशिष्ट्यपूर्णं विद्यते ।

सूक्तेऽस्मिन् उर्वशीपुरूरवसोः संवादो विद्यते । पौराणिकादिवाङ्मयेष्वनयोर्विस्तृता कथा विद्यते परमत्र अष्टादशपरिमितैः मन्त्रैः केवलं वियोगपश्चात्कालिकः संवाद एव वर्णितो विद्यते । विवाहानन्तरं पुत्रमपि प्रसूत्य स्वर्गीयत्वेन उर्वशी स्वर्गं गच्छति । कालक्रमेण एकदा सरोवरसमीपे तां दृष्ट्वा पुरूरवाऽऽत्मना सह मर्त्यलोक एव स्थातुमनुरुणद्धि, परमुर्वशी विविधैः कारणैस्तदनुरोधं नाङ्गीकरोति । एतादृशः संवादोऽस्मिन् सूक्ते प्रस्तुतोऽस्ति । सूक्तस्यास्य दैविकः सामाजिको राजनैतिकश्चेति विविधा अर्था उपलभ्यन्ते । तत्र दैविकपक्षे पुरूरवा मेधस्वरूपः विद्युत्स्वरूपा उर्वशी इति परिकल्प्य प्राकृतिकी व्याख्या क्रियते । तथा च राजनीतिकपक्षे पुरूरवा राजा, तस्य महिषी उर्वशी इत्येतादृशी व्याख्या विधीयते । सामाजिकदृष्ट्या तु पुरूरवा नायकः, उर्वशी च नायिका । अस्याः कथाया विस्तारः शतपथब्राह्मणग्रन्थे लभ्यते । कालिदासादिप्रणीतविक्रमोर्वश्यादिकाव्येषु नायकनायिकारूपेणानयोरुपस्थापनम् ।

\customparagraph{\textbf{वैशिष्ट्यम्}}

संवादेऽस्मिन् न केवलं प्रेमालापोऽपि तु गृहस्थाश्रमस्य वैशिष्ट्यम्, मातुः सन्ततिं प्रति कर्तव्यम्, स्त्रिषु स्थिता  दुर्मनोवृत्तिः, पतिपत्न्योः कर्तव्यम् इत्यादयो विषयाः समागताः सन्ति । तत्र उर्वशीं स्वगृहे स्थापयितुं पुरूरवाः पुत्रस्नेहं पुरस्करोति । ‘पुत्रस्य कृते मातुः स्नेह एव महत्त्वाधायको भवति । तवाभावे तव पुत्रो रोदिति। तस्मात् त्वमत्र तिष्ठ’ इति । एवमेव स्त्रियाश्चरित्रस्य विषये उर्वशी एव कथयति यत्- ‘हे पुरूरवः ! त्वं विलापं मा कुरु, स्त्रीजनानां पुरुषं प्रति स्नेहः चिरस्थायी न भवति। ताः कठोरहृदया भवन्ति । अतो मामपि तादृशमेव जानीहि ।’  तद्यथा-

\begin{shloka}[center]

पुरुरवो मा मृथा मा प्र पप्तो मा त्वा वृकासो अशिवास उ क्षन् ।

न वै स्त्रैणानि सख्यानि सन्ति सालावृकाणां हृदयान्येता ॥

\end{shloka}

एवमेव पुरूरवाः श्वसुरगृहं प्रति नार्यः कर्तव्यमपि प्रस्तौति । यत्-

‘पतिपत्न्योः संवादाभावे गृहस्थाश्रमस्य सौख्यं न लभ्यते, दीर्घतराणि दिनानि सुखप्रदायकानि न भवन्ति, अतस्त्वं कठोरहृदया मा भव, मया सह प्रेमालापं कुरु’ इति बोधयन् पुरूरवाः कथयति-

\begin{shloka}[left]

हये जाये मनसा तिष्ठ घोरे वचांसि मिश्रा कृणवावहै ।

न नौ मन्त्रा अनुदितास एते मयस्करन् परतरे चानाहन् ॥

\end{shloka}

विश्ववाङ्मयस्य प्राचीनतमे ग्रन्थे वर्णितमेतादृशं सरसं वृत्तान्तं सर्वथा महत्त्वाधायकमस्ति । अत्र पतिपत्न्योः स्नेहातिरेकः, गृहस्थाश्रमसञ्चालनाय दम्पत्योः कर्तव्यम्, सन्ततिं प्रति पित्रोः स्नेहः, राज्ञः कर्तव्यम् इत्यादयो विषया वर्णिताः सन्ति । सूक्तमिदं शतपथब्राह्मणे विद्यमानस्य पुरूरवउर्वश्युपाख्यानस्य लौकिकसाहित्यस्य च कृते उपजीव्यभूतमस्ति । इममेव वृत्तान्तमधिकृत्य कविकुलगुरुकालिदासेन ‘विक्रमोर्वशीयम्’ नाटकं प्रणीतम् । नाटकमिदं रसिकसमाजेनादृतञ्च । एतादृशैर्वैशिष्ट्यैः संवादस्यास्य महत्त्वं समधिकं विद्यते ।

\begin{visheshshloka}[center]

\textbf{हये जाये मनसा तिष्ठ घोरे वचांसि मिश्रा कृणवावहै नु}

\textbf{न नौ मन्त्रा अनुदितास एते मयस्करन्परतरे चनाहन् ॥ १ ॥}

\end{visheshshloka}

\customparagraph{संस्कृत व्याख्या}

\{उर्वशीं प्रति पुरूरवस उक्तिः-\} हये ! इति सम्बोधनम् । हे घोरे= कठोरगुणवती जाये= पत्नि ! मनसा= अस्मदुपर्यनुरागवता मनसा युक्ता सती तिष्ठ= सन्निधौ निवस । किमर्थमिति- वचांसि= वाक्यानि मिश्राणि= उक्तिप्रत्युक्तिरूपाणि नु= अद्य क्षिप्रं वा कृणवावहै= करवावहै । किमर्थं संवादः करणीय इत्युच्यते-  नौ= आवयोः, मन्त्राः= रहस्यार्थाः, एते=विवक्षिताः अनुदितासः= परस्परमभाष्यमाणा गुम्फिताः सन्तः परतरे= आगामिनि चनाहन्= अनेकेषु दिवसेषु मयः= सुखम्, न करन्= न कुर्वन्ति । तस्मात् परस्परं संवादं करवावहै ।

\customparagraph{\textbf{नेपाली}}

\{पुरूरवा उवर्शीलाई सम्बोधन गर्दै भन्दछन्-\} हे मेरी कठोर हृदय भएकी पत्नी ! एकछिन पर्ख, मसँग केही कुरा गर । हाम्रा मनका कति कुरा भित्रभित्रै गुम्सिएका छन् । ती अव्यक्त कुराहरूले पछि आउने दिनहरूमा पनि दुःख दिन सक्छन् । त्यसैले तिमी र म केही कुरा गरौं ।

\begin{visheshshloka}[center]

\textbf{किमेता वाचा कृणवा तवाहं प्राक्रमिषमुषसामग्रियेव}

\textbf{पुरूरव: पुनरस्तं परेहि दुरापना वात इवाहमस्मि ॥ २ ॥}

\end{visheshshloka}

\customparagraph{संस्कृत व्याख्या}

\{पुरूरवसम् उर्वशी प्रत्युवाच-\} एतया= प्रयोजनरहितया वाचा= वचनेन किं कृणव= किं करवाव । अहम्= उर्वशी त्विदानीं त्वत्सकाशात् प्राक्रमिषम्= अतिक्रान्तवत्यस्मि । अतिक्रमे दृष्टान्तः- बह्वीनामुषसां मध्येऽग्रे-अग्रे भवा उषाः पूर्वोषां यथा प्राक्रमीत्तथाहमपीति । तस्माद्धे पुरूरवः ! त्वं पुनः अस्मत्सकाशाद् अस्तम्= गृहम्, परेहि= परागच्छ । मयि मा अभिलाषं कार्षीः । स्वस्या दुर्ग्रहत्वमाह- अहम्= उर्वशी वात इव= वायुरिव दुरापना= दुष्प्रापा दुरापा वा अस्मि ।

\customparagraph{\textbf{नेपाली}}

\{पुरूरवालाई उर्वशी प्रत्युत्तरमा भन्छिन्-\} हे पुरूरवा ! यस्ता निष्प्रयोजन कुरा किन गर्छौ ? त्यसको अर्थ छैन। म त अहिले तिमीबाट धेरै टाढा भइसकेकी छु । जस्तैः- पूर्व-पूर्व गइसकेका उषालाई पछि-पछि आउने उषाले जसरी भेट्न सक्दैनन् त्यस्तै तिमीले मलाई भेट्न असम्भवप्राय छ । त्यसैले तिमी म नजिकबाट आफ्नो घर फर्क । म हावा जस्तै अप्राप्य छु ।

\begin{visheshshloka}[center]

\textbf{इषुर्न श्रिय इषुधेरसना गोषाः शतसा न रंहि:}

\textbf{अवीरे क्रतौ वि दविद्युतन्नोरा न मायुं चितयन्त धुनयः ॥ ३ ॥}

\end{visheshshloka}

\customparagraph{संस्कृत व्याख्या}

\{अनया ऋचा पुरूरवाः स्वस्य विरहजनितं वैक्लव्यं तां प्रति ब्रूते-\} हे उर्वशि ! तव विरहे इषुधे= इषवो धीयन्तेऽत्रेतीषुधिर्निषङ्गः, ततः सकाशात् इषुः=बाणः, असना= असनायै (शत्रूणामुपरि प्रक्षेप्तुम्), श्रियै= विजयार्थम्, न= न शक्नोमि । त्वद्विरहस्य कारणेनाहं रंहिः= वेगवान् (भूत्वा) गोषा= शत्रुगवां मध्ये शतसा= शतसङ्‌ख्यकाः स्वायत्तीकर्तुं न= न प्रभवामि । अवीरे= वीरवर्जिते क्रतौ= युद्धादिराजकर्मणि मत्सामर्थ्यं न विदविद्युतत्= न विद्योतते । धुनयः= कम्पयितारो- ऽस्मदीया भटाः, उरा= उरौ विस्तीर्णे सङ्‌ग्रामे मायुम्= शब्दम्, सिंहनादादिकम्, न चिन्तयन्त= न बुध्यन्ते । तव विरहे मम सैन्यबलमपि निष्क्रियमस्तीत्यर्थः ।

\customparagraph{\textbf{नेपाली}}

\{पुरूरवा उर्वशीसित आफ्नो विरहवेदना बताउँछन्-\} हे उर्वशी ! तिम्रो विरहले गर्दा बाण राख्ने ठोक्राबाट बाण निकालेर विजय प्राप्त गर्ने उद्देश्यले शत्रुमाथि प्रहार गर्ने सामर्थ्य ममा छैन । शत्रुका सयौं गाईलाई वेगवान् भई युद्ध गरी ल्याउने सामर्थ्य पनि ममा छैन । तिम्रा विरहमा मेरा सैनिकहरू पनि युद्ध प्रयोजनका लागि गरिने सिंहनादरहित छन् । अर्थात् तिम्रा विरहमा म सबै किसिमले निष्क्रिय एवं विह्वल छु ।

\begin{visheshshloka}[center]

\textbf{सा वसु दधती श्वशुराय वय उषो यदि वष्ट्यन्तिगृहात्}

\textbf{अस्तं ननक्षे यस्मिञ्चाकन्दिवा नक्तं श्नथिता वैतसेन ॥ ४ ॥}

\end{visheshshloka}

\customparagraph{संस्कृत व्याख्या}

\{इदमुर्वशीवाक्यम्, सा पुरात्मना कृतं कार्यमुषसे निवेदयति-\} हे उषः !  सा इयमुर्वशी वसु= वासकम्, सुगन्धयुतम्, वयः= अन्नम्, श्वशुराय= भर्तुः पुरूरवसः पित्रे दधती= प्रयच्छन्ती  तत्र गृहे स्थिता यदि पतिं  वष्टि= कामयते तदा अन्तिगृहात्= स्वभर्तृभोगगृहान्तिके यत् श्वशुरस्य भोजनगृहं तदन्तिगृहम्, तस्माद् गृहात् सा उर्वशी अस्तम्= पतिगृहम्, ननक्षे= व्याप्नोति । यस्मिन्= गृहे चाकन्= कामयत उर्वशी, दिवा, नक्तम्, अहनि, रात्रौ च वैतसेन= पुंस्प्रजननेन श्नथिता= तााडिता च भवति । एवमुर्वश्यात्मानं परोक्षेण निर्दिदेश ॥ ४॥

\customparagraph{\textbf{नेपाली}}

\{यस ऋचाद्वारा उर्वशी पहिलेको वृत्तान्त सत्यताका लागि उषालाई सुनाउँछिन्-\} हे उषा ! यो त्यो उर्वशी हो जो पहिले सुगन्धित अन्न ससुरा अर्थात् पति पुरूरवाका बुबालाई दिन्थी । त्यहाँ रहेका वेला पतिसँग सहवासको कामना भएमा भोजनगृहबाट नजिकै रहेको आफ्ना पति पुरूरवाको घरमा आउँथी । रात, दिन, साँझ आदि जस्तोसुकै समयमा पनि  पतिगृहमा पतिसँग रमाइलो गर्दथी । (यसरी उर्वशीले घुमाउरो पाराले आफ्नो अतीत बताउँछिन्)

\begin{visheshshloka}[center]

\textbf{त्रिः स्म माह्न: श्नथयो वैतसेनोत स्म मेऽव्यत्यै पृणासि}

\textbf{पुरूरवोऽनु ते केतमायं राजा मे वीर तन्व१स्तदासीः ॥ ५ ॥}

\end{visheshshloka}

\customparagraph{संस्कृत व्याख्या}

\{अनया ऋचा उर्वशी पुरूरवसमेव सम्बोध्य उक्तवती- \} हे पुरूरवः !  त्वं मा= माम्, अह्नः= अहनि वैतसेन= पुंव्यञ्जनेन दण्डेन त्रिः= त्रिवारम्, श्नथय स्म= अश्नथयः, अताडयः । उत= अपि च अव्यत्यै= सपत्निभिः सह पर्यायेण पतिमागच्छति सा व्यती, न तादृशी इति अव्यती तस्यै, निरन्तरभोगयुतायै मे= मह्यम्, पृणासि= पूरयसि । हे पुरूरवः ! ते= तव केतम्= गृहम्, अनु आयम्= अन्वगमम्  पूर्वमिति । हे वीर ! राजा ! त्वं च मे= मम तन्वः= शरीरस्य तत्= तदा आसीः= अभवः सुखयितेति । परमप्येवं मन्तव्यम्, किमिति कातरो भवसीत्युवाच  ॥ ५॥

\customparagraph{\textbf{नेपाली}}

\{यो ऋचाद्वारा उर्वशी पुरूरवालाई सम्बोधन गर्दै भन्छिन्- \} “हे पुरूरवा ! तिमीले मलाई प्रतिदिन तीन पटक पुंव्यञ्जित दण्डले भोग गर्दथ्यौ। मलाई तिम्रा अरू पत्नी अर्थात् सौताकले झैँ पालो पर्खिनु पर्दैनथ्यो किनकि मलाई तिमी औधी माया गर्थ्यौ । यसरी म पूर्व समयमा तिम्रा घरमा बसें । त्यतिबेला तिमी मेरो शरीरका सुखदाता थियौ । (यो सबै विगत हो । त्यसलाई सम्झेर अहिले किन दुःखी बन्छौ ?”)

\begin{visheshshloka}[center]

\textbf{या सुजूर्णिः श्रेणि: सुम्नआपिर्ह्रदेचक्षुर्न ग्रन्थिनी चरण्युः}

\textbf{ता अञ्जयोऽरुणयो न सस्रुः श्रिये गावो न धेनवोऽनवन्त ॥ ६ ॥}

\end{visheshshloka}

\customparagraph{संस्कृत व्याख्या}

\{पुरूरवसो वाक्यम्- \} या ‘सुजूर्णि’ इति नामिका उर्वशी श्रेणिः, सुम्नआपिः, ह्रदेचक्षुः, एतन्नामिकाभिरालिभिर्मानिनिभिः सहिता ग्रन्थिनी= ग्रन्थनवती, सन्दर्भवती चरण्युः= चरणशीला आजगाम । यद्वा- ग्रथिनीत्येतत्सखिभूताप्सरोनामधेयम् । या सुजूर्णिः= सुजवोर्वशी सा ताभिः सह पूर्वमाजगाम । ताः= अप्सरसः, अञ्जयः= आभरणोपेताः, अरुणयः= अरुणवर्णाः, न सस्नुः= पूर्ववन्न गच्छन्ति, अधुना न समागच्छन्तीत्यर्थः । ता अप्सरसः   श्रिये= श्रयणाय धेनवः= नवप्रसूताः गाव इव, न अनवन्त= अधुना न शब्दायन्ते,  (यथा गाव आश्रयार्थं शब्दायन्ते तथाधुनाप्सरसः न शब्दायन्तीति) ।

\customparagraph{\textbf{नेपाली}}

पुरूरवा भन्छन्-  जो सुजूर्णि उर्वशी, श्रेणि, सुम्नआपि, ह्रेदचक्षु आदि अभिमानले युक्त साथीहरूका साथमा आउँथिन् ती उर्वशी आभूषणद्वारा अलङ्‌कृत रक्त वर्णका अप्सराहरूका साथ अहिले यहाँ आउँदिनन् ।  जसरी गाईहरू आश्रयका लागि कराउँछन् (आवाज निकल्दछन्) त्यसरी नै उनीहरूले (उर्वशीसहित उनका साथी अप्सराहरूले) कुनै आवाज निकाल्दैनन्।

\begin{visheshshloka}[center]

\textbf{समस्मिञ्जायमान आसत ग्ना उतेमवर्धन्नद्य१: स्वगूर्ताः ।}

\textbf{महे यत्त्वा पुरूरवो रणायावर्धयन्दस्युहत्याय देवाः ॥ ७ ॥}

\end{visheshshloka}

\customparagraph{संस्कृत व्याख्या}

\{उवर्शीवाक्यम्-\} अस्मिन्= पुरूरवसि जायमाने= भूलोके सञ्जाते सति ग्नाः= अप्सरसोऽपि सम् आसत= सङ्गता अभवन् । अथवा- जायमाने=  यज्ञार्थं प्रवर्धमाने सति ग्नाः= देववत्न्योऽपि समासत= समभवन्। उत= अपि च ईम्= एनं पुरूरवसम्, स्वगूर्ताः= स्वयं गामिन्यः, नद्यः= तासामाश्रयभूता अवर्धयन्, संस्तुतवत्य इत्यर्थः । किञ्च हे पुरूरवः ! यत्= यस्मात् त्वा= त्वाम्, महे= महते रणाय= रमणीयाय सङ्‌ग्रामाय दस्युहत्याय= चौरहननाय देवाः= देवताः, अवर्धयन्, प्रेरितवन्त  इत्यर्थः।

\customparagraph{\textbf{नेपाली}}

\{उर्वशीका साथमा आएका अन्य अप्सराहरूको र उर्वशीको पनि तिमीले उपभोग गरेका थियो भन्ने उर्वशीको कथन-\} पृथ्वीलोकमा यी (तिमी) पुरूरवा जन्मिएको अवसरमा अप्सरा/देववेश्याहरू पनि नजिक आएका थिए । (अथवा यिनले याज्ञिक कार्यलाई बढाउँदा देवपत्नीहरू पनि यहाँ आएका थिए ।) स्वयं गमन गर्ने नदीहरूका आश्रय भएका (स्वतः आएका) अप्सराहरूले पुरूरवालाई बढाएका (सेवा गरेका) थिए ।  त्यस्तै हे पुरूरवा ! देवताहरूले रमणीय सङ्ग्रामका लागि र शत्रुहरूलाई संहार गर्नका लागि तिमीलाई प्रेरित गरेका थिए ।

\begin{visheshshloka}[center]

\textbf{सचा यदासु जहतीष्वत्कममानुषीषु मानुषो निषेवे ।}

\textbf{अप स्म मत्तरसन्ती न भुज्युस्ता अत्रसन्रथस्पृशो नाश्वा: ॥ ८ ॥}

\end{visheshshloka}

\customparagraph{संस्कृत व्याख्या}

\{क्रमशः त्रीणि पुरूरवसो वाक्यानि-\} यदा सचा= सहायभूतः पुरूरवाः अत्कम्= स्वकीयं रूपम्, जहतीषु= परित्यजन्तीषु अमानुषीषु= देवताभूतास्वप्सरःसु मानुषः सन् निषेवे= अभिमुखं गच्छति, तदानीं ताः= अप्सरसः मत्= मत्त अप= अपसृत्य अत्रसन्= अगच्छन् (पलायिताः) । पलायने दृष्टान्तः- यथा तरसन्ती= मृगस्त्री भुज्युः= भोगसाधनभूता सा व्याधाद्भीता पलायते, एवमेव रथस्पृशो नाश्वा= रथे नियुक्ता अश्वाः पलायन्ते, तद्वद् एता अप्सरसो मां दृष्ट्‌वा पलायन्ते इत्यर्थः । (पुरूरवा अनेकाभिरप्सरोभिः सह उर्वशी त्वया भुक्तेति उर्वशीकथनं प्रत्याचष्टे) ।

\customparagraph{\textbf{नेपाली}}

\{पुरूरवाले क्रमशः तीन मन्त्रमा आफ्नो धारणा राख्छन्-\} जुन बेला पुरूरवा आफ्ना सखीहरूका साथमा आफ्नो रूप परिवर्तन गरी परिभ्रमण गर्ने उर्वशीसहितका अप्सराहरूका अगाडि मनुष्य रूपमा पुग्थे तब ती अप्सराहरू व्याधाबाट त्रसित मृगीझैं अथवा रथमा जुटेका घोडाझैँ भाग्थे । (यस पद्यमा- ‘उर्वशीको र साथमा आएका उनका सखीहरूको पनि तिमीले उपभोग गरेका थियौ’ भन्ने उर्वशीका कथनको पुरूरवाले प्रतिवाद गरेका हुन् ।)

\begin{visheshshloka}[center]

\textbf{यदासु मर्तो अमृतासु निस्पृक्सं क्षोणीभि: क्रतुभिर्न पृङ्क्ते ।}

\textbf{ता आतयो न तन्व: शुम्भत स्वा अश्वासो न क्रीळयो दन्दशानाः ॥ ९ ॥}

\end{visheshshloka}

\customparagraph{संस्कृत व्याख्या}

यत्= यदा आसु = एतासु, अमृतासु= मरणधर्मरितासु अप्सरःसु, मर्तः= मनुष्यः पुरूरवाः, निस्पृक्= निःशेषेण स्पृशन्, क्षोणीभिः= वाग्भिः क्रतुभिर्न= कर्मभिश्च सं पृङ्क्ते= सम्पर्कं करोति (नकारः समुच्चयार्थः) । ताः= अप्सरसः, आतयः= आतिभूतास्तदानीं स्वाः तन्वः= आत्मीयानि रूपाणि न शुम्भत= न प्रकाशयन्ति । अश्वासो  न= अश्वा इव क्रीलयः= सङ्‌क्रीडमानाः, दन्दशानाः= दन्दशूकाः (सरीसृपाः) जिह्वाभिरात्मीयाः सृक्का= भक्षयन्तः । ते यथा चाल्पेन धावन्तः स्वरूपं न प्रयच्छन्ति रथिकाय तद्वदिति दुःखाद्‌ ब्रूते ॥ ९॥

\customparagraph{\textbf{नेपाली}}

जुन वेला मनुष्य रूपका पुरूरवा देवयोनिविशेष अप्सराहरूलाई स्पर्श गर्दै वाणी एवं कर्मका माध्यमले सम्पर्क बनाउन चाहन्थे त्यस वेला तिनीहरूले आफूलाई प्रकट गर्दैनथे (समर्पित हुँदैनथे)। क्रीडाकारी अश्वहरूले आफ्ना जिब्राले सृक्का (ओँठको छेउ) चाट्दै चञ्चलतापूर्वक भागेर रथीका समक्षमा आफ्नो स्वरूप प्रकट नगरेझैँ गर्दथे । (जसरी चञ्चलतापूर्वक भागेका र क्रीडामा निमग्न भएका घोँडाहरूले सर्पलेझैँ आफ्ना जिब्राले ओँठका प्रान्तभाग चाड्दै रथीका समक्षमा आफ्नो स्वरूप प्रकट गर्दैनन् त्यसरी नै उर्वशी र अन्य अप्सराहरूले मेरा समक्षमा आफ्नो समर्पणमूलक स्वरूप प्रकट गरेका थिएनन्) ।

\begin{visheshshloka}[center]

\textbf{विद्युन्न या पतन्ती दविद्योद्भरन्ती मे अप्या काम्यानि}

\textbf{जनिष्टो अपो नर्य: सुजात: प्रोर्वशी तिरत दीर्घमायु: ॥ १० ॥}

\end{visheshshloka}

\customparagraph{संस्कृत व्याख्या}

या= उर्वशी विद्युन्न= विद्युदिव पतन्ती= गच्छन्ती दविद्योत्= द्योतते । किं कुर्वती- अप्या= अप इत्यन्तरिक्षनाम, तत्सम्बन्धीनि व्याप्तानि वा, काम्यानि= अस्मदभिमतान्युदकानि वा मे= मह्यम्, भरन्ती= सम्पादयन्ती । यदा आगतायास्तस्याः सकाशात् अपः= व्याप्तः, कर्मसु नर्यः= नरेभ्यो हितः सुजातः= सुजननः पुत्रः जनिष्टः‍= अजनिष्टः, उत्पद्यते, तदानीं मम उर्वशी दीर्घमायुः प्र तिरत= प्रवर्धयति ।

\customparagraph{\textbf{नेपाली}}

यी उर्वशी विद्युत्‌समान गमन गर्दै प्रकाशमयी हुन्छिन् । अन्तरिक्षसित सम्बन्ध राख्ने अथवा व्यापक रूपमा विद्यमान मेरा सबै अभिलाषा (जलसम्बन्धी कामना) तिनले पूर्ण गराएकी थिइन् । मेरा समीपमा आएकी उर्वशीबाट जब मनुष्यहरूको कल्याण गर्ने पुत्र उत्पन्न हुनेछ तब उर्वशीले मलाई दीर्घजीवि बनाऊन् ।

\begin{visheshshloka}[center]

\textbf{जज्ञिष इत्था गोपीथ्याय हि दधाथ तत्पुरूरवो म ओज:}

\textbf{अशासं त्वा विदुषी सस्मिन्नहन्न म आशृणो: किमभुग्वदासि ॥ ११ ॥}

\end{visheshshloka}

\customparagraph{संस्कृत व्याख्या}

\{उर्वशीवाक्यम्-\} इत्था= इत्थम्, गोपीथ्याय= भूमे रक्षणाय जज्ञिसे= ‘आत्मा वै पुत्रनामा’ इति श्रुतेः पुत्ररूपेण जातः (गर्भस्थः) असि । पुनस्तदेवाह- हे पुरूरवः ! मे= ममोदरे मयि ओजः= अपत्योत्पादनसामर्थ्यम्, दधाथ= निहितवानसि । तत्= तथास्तु । अथापि स्थातव्यमिति चेत् तत्राह- अहं विदुषी भावि कार्यं जानती सस्मिन्नहन्= सर्वस्मिन्नहनि त्वया कर्तव्यं त्वा= त्वाम्, अशासम्= शिक्षितवत्यस्मि । त्वं मे= मम वचनं न आशृणोः= न शृणोषि । किं त्वम्, अभोक्ता= अपालयिता, विवाहप्राग्विहितं प्रतिज्ञातार्थमपालयन् वदासि= ‘हयेे जाये’ इत्यादिरूपं प्रलापं करोसि ।

\customparagraph{\textbf{नेपाली}}

\{उर्वशी भन्छिन्-\} हे पुरूरवा ! पृथ्वीको रक्षाका लागि तिमी मेरो पेटमा पुत्र रूपमा प्रकट भएका छौ । तिमीले मेरो गर्भमा पुत्र उत्पन्न गराउन समर्थ वीर्य स्थापित गरेका छौ । (अझै पनि तिमीले मलाई यहीँ बस भन्नु उचित हुँदैन, किनभने) पछि गर्नुपर्ने कर्तव्य जान्ने मैले सधैँ तिमीलाई कर्तव्य सम्बन्धी शिक्षा दिएकी थिएँ (कर्तव्य सम्झाएकी छु) । तिमीले भनाइ (कथन) सुनेनौँ। किन विवाहपूर्वको सर्त पालन नगरेर तिमी अहिले ‘हये जाये’ इत्यादि प्रलाप गरिरहेका छौ ।[^1]

\begin{visheshshloka}[center]

\textbf{कदा सूनुः पितरं जात इच्छाच्चक्रन्नाश्रु वर्तयद्विजानन्}

\textbf{को दम्पती समनसा वि यूयोदध यदग्निः श्वशुरेषु दीदयत् ॥ १२॥}

\end{visheshshloka}

\customparagraph{संस्कृत व्याख्या}

**\{**इदं पुरूरवसो वाक्यम्-\} कदा= कस्मिन् काले सूनुः= तवोदरजातः पुत्रः पितरं माम् इच्छात्= इच्छेत् । कदा वा विजानन्= पितरमधिगच्छन् चक्रन्= क्रन्दमानः, नाश्रु वर्तयत् (नेति चार्थे) । कः=किंविधः सन् सूनुः समनसा= समनसौ दम्पती= जायापती त्वां मां च वि यूथोत्= विश्लेषयत् । अध= अधुना यत्= यदा अग्निः= तव हृदयस्थितस्तेजोरूपो गर्भः श्वशुरेषु= मम पितृजनेषु सन्ततिसुखरूपेण दीदयत्= दीप्यते ।

\customparagraph{नेपाली}

\{पुरूरवाको कथन-\} कहिले त्यस्तो दिन आउला, जुन दिन मलाई तिम्रो पेटबाट जन्मिएको पुत्रले इच्छा गर्ला । कहिले मलाई चिनेर मेरा छेउमा आएर झगडा गर्दै (रुँदै) आँसु चुहाउला । कस्तो रूप लिएर छोराले एउटा आत्मा भएका दम्पती अर्थात् आफ्ना माता र पिता तिमी र मलाई परस्परमा एकाकार बनाउला । अहिले तिम्रो पेटमा रहेको अग्निरूप गर्भ कहिले तिम्रा ससुरा अर्थात् मेरा कुलका मान्यजनलाई आनन्द प्रदान गर्दै प्रकाशित होला।

\begin{visheshshloka}[center]

\textbf{प्रति ब्रवाणि वर्तयते अश्रु चक्रन्न क्रन्ददाध्ये शिवायै}

\textbf{प्र तत्ते हिनवा यत्ते अस्मे परेह्यस्तं नहि मूर माप: ॥ १३॥}

\end{visheshshloka}

\customparagraph{संस्कृत व्याख्या}

\{इदमुर्वशीवाक्यम्-\} हे पुरूरवः ! त्वां प्रति ब्रवाणि= प्रतिवच्मि । त्वदपत्यम् अश्रु= वाष्पम्, वर्तयते= वर्तयिष्यति । आध्ये= आध्याते वस्तुनि शिवायै= शिवे कल्याणे समुपस्थिते सति चक्रन्= रुदन्नश्रूणि विमुञ्चन् न क्रन्दत्= रोत्स्यति चेत्यर्थः (नकारश्चार्थे) । तत्= त्वदपत्यम्, ते= तुभ्यम्, हिनव= प्रहिणोमि । यत् अपत्यं ते= तव सम्बन्धि अस्मे= अस्मासु निहितमस्तीति त्वं परेह्यस्तम्= अस्तम्-स्वगृहं परेहि- प्रतिगच्छ । हे मूर ! मूढ, मा= माम्, नहि आप= न प्राप्नोषि ।

\customparagraph{\textbf{नेपाली}}

हे पुरूरवा ! तिमीलाई म स्पष्टसँग भन्छु सुन- जब तिम्रो छोरो जन्मिनेछ, त्यो रुँदै तिम्रो छेउमा आउनेछ । त्यसको म राम्ररी हेरविचार गर्छु । तिम्रो सन्तान जुन मेरो पेटमा हुर्किरहेको छ। भविष्यमा त्यसलाई तिमीले पाउनेछौ । अहिले तिमी आफ्नो घर फर्क । हे मूर्ख ! मलाई तिमीले कहिल्यै पाउन सक्नेछैनौ ।

\begin{visheshshloka}[left]

\textbf{सुदेवो अद्य प्रपतेदनावृत्परावतं परमां गन्तवा उ ।}

\textbf{अधा शयीत निॠतेरुपस्थेऽधैनं वृका रभसासो अद्युः ॥ १४ ॥}

\end{visheshshloka}

\customparagraph{संस्कृत व्याख्या}

\{अथ परिदूनः पुरूरवा उवाच-\}   सुदेवः= त्वया सह सुक्रीडः पतिः अद्य प्रपतेत्= अत्रैव पतेत् । अथवा अनावृत्= अनावृत्तः सन् परमां परावतम्= दूरादपि दूरदेशम्, गन्तवै= गन्तुं महाप्रस्थानगमनं कुर्यात् । अध= अथवा यत्र कुत्रापि गन्तुमसमर्थः, निर्ऋतेः= पृथिव्याः उपस्थे शयीत= शयनं कुर्यात् । यद्वा निर्ऋतिः= पापदेवता तस्या उपस्थे= उत्सङ्गे सन्निधौ म्रियतामित्यर्थः । अध= अथवा एनम्= पुरूरवसम्, वृकाः आरण्याः श्वानः, रभसासः= वेगवन्तः, अद्युः= भक्षयन्तु ।

\customparagraph{\textbf{नेपाली}}

\{उर्वशीको कुरा सुनेर दुःखी भएका पुरूरवा भन्छन्-\}  हे उर्वशी ! तिमीसँग पहिले प्रेमपूर्ण रमाइलो क्षण बिताएको तिम्रो पति यहीँ पृथिवीतलमा लड्न सक्छ । अथवा अनावृत (नाङ्गै) टाढाटाढा जान (अनन्त यात्रामा प्रस्थान गर्न) सक्छ । हिँड्‌न असमर्थ भई जहाकहीँ पृथिवीको काखमा सुत्न सक्छ (अथवा निर्ऋति- पापदेवताका काखमा मर्न पनि सक्छ) । त्यसबेला यो पुरूरवालाई आहाराका लागि झम्टिने ब्बाँसा आदि वन्यजन्तुले खानेछन् ।

\begin{visheshshloka}[center]

\textbf{पुरूरवो मा मृथा मा प्र पप्तो मा त्वा वृकासो अशिवास उ क्षन् ।}

\textbf{न वै स्त्रैणानि सख्यानि सन्ति सालावृकाणां हृदयान्येता ॥ १५ ॥}

\end{visheshshloka}

\customparagraph{संस्कृत व्याख्या}

(उर्वशी पुरूरवसं प्रत्युवाच-) हे पुरूरवः ! त्वं मा मृथाः= मृत्युं मा प्राप्नुहि । तथा मा प्रपप्तः= अत्रैव पतनं मा कार्षीः । तथा त्वा= त्वाम् अशिवासः= अशुभाः वृकासः= वृकाः मा उक्षन् (उ इत्येवकारार्थे) अक्षन्= माभ्यवहारयन्तु, मा भक्षयन्तु इत्यर्थः । (स्वस्नेहस्यासारतामाह-) स्त्रैणानि= स्त्रीणां कृतानि सख्यानि न वै सन्ति= न भवन्ति खलु । कथमिति- एता= एतानि सख्यानि सालावृकाणां हृदयानि= यथा- तेषां हृदयानि विश्वासापन्नानां घातकानि भवन्ति, तद्वत् स्त्रीणां सख्यानि कठोरपरिणामप्रदायीनि भवन्तीति ।

\customparagraph{\textbf{नेपाली}}

\{उर्वशी पुरूरवालाई भन्छिन्-\} हे पुरूरवा ! तिमी बेकारमा मर्ने कुरा नगर । तिमी कतै भुइँमा लड्ने काम पनि नगर । तिमीलाई अशुभ ब्वाँसा आदि जङ्गली ह्रिंसक जन्तुले पनि नखाऊन् । स्त्रीहरूसँगको मित्रता, वास्तविक मित्रता हुँदैन । यिनीहरूको मित्रता ब्बाँसाको हृदयजस्तो हुन्छ । (जसरी ब्बाँसा आदि ह्रिंसक जन्तुले पहिले विश्वास सिर्जना गरेर पछि विश्वासघात गर्छन् त्यस्तै स्त्रीसँगको मित्रताको परिणाम पनि चिरस्थायी नभई कठोर परिणाम दिने खालको हुन्छ ।) ।

\begin{visheshshloka}[center]

\textbf{यद्विरूपाचरं मर्त्येष्ववसं रात्री: शरदश्चतस्रः}

\textbf{घृतस्य स्तोकं सकृदह्न आश्नां तादेवेदं तातृपाणा चरामि ॥ १६ ॥}

\end{visheshshloka}

\customparagraph{संस्कृत व्याख्या}

यत्= यदाऽहं विरूपा= मनुष्यसम्पर्काद्विगतसहजभूतदेवरूपा पत्यानुकूल्येन नानारूपा वा मर्त्येषु= मनुष्येषु अचरम्,  तदानीं रात्रीः पूरयित्रीः चतस्रः शरदः अवसम्= न्यवसम् । तदानीं घृतस्य स्तोकं सकृदह्नः आश्नाम् । तादेव= तेनैव स्तोकेनाहम् इदं सम्प्रति तातृपाणा= तृप्ता सती चरामि ।

\customparagraph{\textbf{नेपाली}}

\{उर्वशी कथन-\} जहिले मनुष्यको सम्पर्कमा आएर देवरूप त्यागेँ अथवा अनुकूल पतिलाई पाएर विभिन्न रूप लिई मनुष्य (पुरूरवा) सँग विहार गरेँ, त्यतिबेला मैले चार वर्ष उनका साथमा बिताएकी थिएँ । त्यसबेला मैले प्रतिदिन एकपटक घिउ खाएकी थिएँ । त्यही घिउले तृप्त भई अहिलेसम्म म घुमिरहेकी छु ।

\begin{visheshshloka}[center]

\textbf{अन्तरिक्षप्रां रजसो विमानीमुप शिक्षाम्युर्वशीं वसिष्ठः ।}

\textbf{उप त्वा रातिः सुकृतस्य तिष्ठान्नि वर्तस्व हृदयं तप्यते मे ॥ १७ ॥}

\end{visheshshloka}

\customparagraph{संस्कृत व्याख्या}

\{पुरूरवसः कथनम्-\} अन्तरिक्षप्राम्= स्वतेजसान्तरिक्षस्य पूरयित्रीं तथा रजसः= रञ्जकस्योदकस्य विमानीम्= निर्मात्रीम् उर्वशीं वशिष्ठः= समानानां मध्येऽतिशयेन वासयिताहम् उप शिक्षामि= वशं नयामि । सुकृतस्य= शोभनकर्मणः, रातिः= दाता पुरूरवाः, त्वा= त्वाम् उपतिष्ठात्= उपतिष्ठतु । मे= मम, हृदयम्= मनः, तप्यते= सन्तप्तं भवति । अतो निवर्तस्व= प्रत्यागच्छ।

\customparagraph{\textbf{नेपाली}}

\{पुरूरवाको कथन-\} हे उर्वशी ! तिमी आफ्नो तेजद्वारा अन्तरिक्षलाई पूर्ण गराउँछ्यौ।  तिमी आनन्ददायक जल उत्पन्न गराउँछ्यौ । यस्ती तिमी उर्वशीलाई वशिष्ठ (इन्द्रियलाई अधीन बनाउन समर्थ) म आफ्ना वशमा राख्न चाहन्छु (तिमीलाई आफ्नी बनाउन चाहन्छु)। राम्रा कीर्तियुक्त काम गर्ने पुरूरवा तिम्रो सान्निध्यमा रहून्। मेरो हृदय अत्यन्तै सन्तप्त छ । अतः स्वर्गलोकबाट मेरा घरमा फर्क ।

\begin{visheshshloka}[center]

\textbf{इति त्वा देवा इम आहुरैल यथेमेतद्भवसि मृत्युबन्धुः ।}

\textbf{प्रजा ते देवान्हविषा यजाति स्वर्ग उ त्वमपि मादयासे ॥ १८ ॥}

\end{visheshshloka}

\customparagraph{संस्कृत व्याख्या}

\{उर्वशीकथनम्-\} हे ऐल= पुरूरवः ! त्वा= त्वाम्, इमे=देवाः, इति= एवम् आहुः= कथितवन्तः- त्वं मृत्युबन्धुः= मृत्योर्बन्धकः मृत्योर्बन्धुभूतो वा । मृत्युवशमप्राप्नुवन्स्त्वं यथे= यथा एतद्भवसि= भविष्यसि । प्रजा= प्रकर्षेण जायमानस्त्वं ते= तव सम्बन्धिनो यष्टव्यान् देवान्  हविषा यजासि= यजसि । उ स्वर्ग एव त्वमपि मादयासे, अस्माभिः सह मादयासे= मोदमनुभवसि ।

\customparagraph{\textbf{नेपाली}}

\{उर्वशीको कथन-\} हे पुरूरवा ! तिम्रा सम्बन्धमा देवताहरू यसो भन्छन्- तिमीले मृत्युलाई जित्नेछौ, अर्थात् तिम्रा लागि मृत्यु बन्धुसरह हुनेछन् । तिमी र तिम्रा सम्बन्धीले यज्ञ गरेर देवतालाई खुसी पार्नेछन् । (प्रसन्न भएका देवताका कृपाले) तिमी हामीसँगै स्वर्गमा आनन्दपूर्वक रहनेछौ ।

\begin{visheshshlokatwo}[center]{84}

\textbf{ऋग्वेदको दसौं मण्डलअन्तर्गत एक सय}

\textbf{पन्चानब्बे सूक्तमा वर्णित उर्वशी-पुरूरवा-संवादको}

\textbf{नेपाली भावानुवाद पूर्ण भयो ।}

\end{visheshshlokatwo}

\specialkhanda{\textbf{अथर्ववेदीयम् पृथ्वीसूक्तम्}}

\customparagraph{\textbf{अथर्ववेदपरिचयः}}

‘अथ लोकमङ्गलाय अर्व्यते प्रस्तूयते इति, आनन्तर्ये परब्रह्म गमयति प्रापयति वा’ इति विग्रहे अथपूर्वकात् ‘ऋ गतिप्रापणयोः’ इति धातोर्वनिप्प्रत्ययेन अथर्वशब्दस्य निष्पत्तिर्भवति । अस्यार्थो भवति यत्- सद्य एव परब्रह्मप्राप्तिर्येन भवति तद् ‘अथर्व’ इति । अथर्वसंहिताया शौनशाखायां विंशतिः कण्डाः, अष्टत्रिंशत् प्रपाठकाः, नवतिः अनुवाकाः, एकत्रिंशदधिकसप्तशतं सूक्तानि, तथा प्रायेण षट्‌सहस्रं मन्त्राश्च विद्यन्ते ।

ऋग्वेदादयस्त्रयो वेदा आमुष्मिकम् (पारलौकिकम्) फलं प्रयच्छन्ति, परम् अथर्ववेदः पारलौकिकेन सह  ऐहलौकिकमपि फलं प्रयच्छति । अथर्ववेदे शान्त्याः पौष्टिककार्यस्य च सम्पादनाय विविधानि सूक्तानि विद्यन्ते । अथर्ववेदसंहितायाः परिशिष्टे समुल्लिखितमस्ति यत्- ‘यस्य राज्ञो जनपदेऽथर्ववेदस्य विज्ञो निवसति, तस्य राज्ञो राष्ट्रं निरुपद्रवं भूत्वा वृद्धिं गच्छति ।’ अथर्ववेदसंहितायाः अध्येता ‘ब्रह्मा’ इति सम्बोध्यते। ‘ब्रह्मा’ इति पदस्यार्थं कुर्वता निरुक्तकारेण  उक्तं यत्- ‘ब्रह्मा सर्वविद्य सर्वं वेदितुमर्हति ।’ ‘सर्वविद्यः’ इत्यनेन पदेन त्रयीविद्यासंयुक्तः इत्यर्थो ज्ञायते । त्रयी विद्या नाम ऋग्यजुस्सामेति वेदत्रयी । अर्थादसौ ब्रह्मा सर्ववेदस्य ज्ञाता विद्यते । तस्माद् वैदिकसंहितासु अथर्ववेदसंहिता पृथग्वैशिष्ट्यमावहति ।

\customparagraph{\textbf{पृथ्वीसूक्तपरिचयः}}

पृथ्वीसूक्तमिदमथर्ववेदस्य शौनकीयसंहिताया द्वादशकाण्डान्तर्गतं प्रथमं सूक्तमस्ति । समानार्थकत्वाद् इदं सूक्तं भूमिसूक्तमपि कथ्यन्ते । पृथ्वीसूक्तस्य अथर्वा ऋषिः, पृथ्वी च देवता अस्ति । छन्दांसि नैकानि विद्यन्ते । ‘अथर्वा’ इति नाम्ना ऋषिणा दृष्टाः त्रिषष्ष्ठीसङ्‌ख्यका मन्त्राः अस्माकमावासभूतायाः पृथिव्याः स्तुतिरूपेण समागताः सन्ति । इयं स्तुतिरूपा रचना वैदिकसाहित्यस्य सौन्दर्यं विद्यते ।

सा पृथ्वी धर्मस्य कारणेन सुस्थिरा विद्यते । वाचिकं सत्यम्, धृतिः, ऋतम्, क्षत्रम्, दीक्षा, तपः, ब्रह्म, यज्ञः इत्येतदष्टकं तत्त्वम्; एते सर्वाणि तत्त्वानि अस्माकं निवासभूतां पृथिवीं धारयन्ति । प्राचीनकाले पृथिवी जलमध्ये निमग्ना आसीत् । तां वराहभगवान् समुद्धारयत् । अनन्तरं व्यक्ता पृथिवी जनैरावासस्थानत्वेन स्वीकृता । इयं पृथ्वी अस्माकं कृते माता अस्ति । तादृशी मातृतुल्या पृथिवी विचित्ररूपा विद्यते । अत्र नदी, पर्वतः, समुद्रः, अन्नमयी भूमिः इत्यादयः सन्ति । प्राणिनां जीवनधारकत्वेन फलप्रदातारो वनस्पतयोऽपि विद्यन्ते । द्विपदचतुष्पदयुता प्राणिनस्तत्र विचरन्ति । इन्द्रादिदेवगणैः सेविता इयं यज्ञस्य वेदितुल्या अस्ति । इयं यज्ञविधानेन देवमनुष्ययोः समागमभूभिरप्यस्ति। इयमनन्तकालात् प्राणिनो जन्मतो मृत्युपर्यन्तामवस्थां धारयति । इयं कुत्रचित् कठोरा कुत्रचित् कोमला सर्वेषां कृते समानरूपेणान्नदात्री अस्ति । तत्र पर्जन्यदेव जलवर्षणेन सिञ्चति अत एव स पतिरूपेण परिकल्पितो विद्यते । अत्र पृथिव्यां मनुष्याः श्रद्धया देवान्नुद्दिश्य हविः प्रददति । अस्यां सज्जना एव नापितु दुष्टजना अपि निवसन्ति । तस्मात् कविः स्तौति यत्- हे माता ! दुष्टजनस्य समृद्धिं मा कुरु, अस्मभ्यः तान् दूरीकुरु । मातृरूपेण वन्द्यमाना भवती अस्मान् सुखय । त्वं सर्वभावेन पूर्णा असि । सर्वपदार्थैः स्तुतिकारकान् अस्मान् संवर्धय । इत्यनेन प्रकारेण पृथिव्याः स्तुतिः पृथिवीसूक्ते विहिता विद्यते ।

\begin{visheshshloka}[center]

\textbf{सत्यं बृहदृतमुग्रं दीक्षा तपो ब्रह्म यज्ञः पृथिवीं धारयन्ति ।}

\textbf{सा नो भूतस्य भव्यस्य पत्न्युरुं लोकं पृथिवी नः कृणोतु ॥१॥}

\end{visheshshloka}

\customparagraph{\textbf{अन्वयः}}

सत्यम् बृहत् ऋतम् उग्रम् दीक्षा तपः ब्रह्म यज्ञः (इति) नः पृथिवीं धारयन्ति । सा भूतस्य भव्यस्य पत्नी पृथिवी नः उरुं लोकं कृणोतु ।

\customparagraph{\textbf{शब्दार्थः}}

सत्यम्= वाचिकं यथार्थभाषणम्, बृहत्= महत्तापादको गुणो धृतिः, गाम्भीर्यम्, ऋतम्= मानसं यथार्थसङ्कल्पनम्, उग्रम्= शाशकत्वम् दीक्षा= स्वाधिकारस्वकर्तव्यानुरूपः श्रद्धापूर्वको वाग्यमनादिनियमविशेषस्वीकारः, तपः= द्वन्द्वसहनपुरस्सरं वाङ्‌मनसोर्युक्तरूपं नियोजनम्, ब्रह्म= वेदः यज्ञः= हितकृदादानप्रदानात्मको व्यवहारः, इत्येतदष्टकम्, नः= अस्माकम् पृथिवीम्= आत्मवासभूताम् भूमिंम्, धारयन्ति= उपप्लवादिविरहेणावष्टम्भयन्ति । सा= सत्याद्यष्टकविधृता  पूर्वोक्ता, भूतस्य= पूर्वं विहितस्य कार्यस्य, भव्यस्य= भाविनश्च कार्यस्य, पत्नी= पालयित्री पृथिवी= भूमिः, नः= अस्माकं कृते, ऊरुम्= विस्त्तीर्णम्, लोकम्= वासस्थानम्, कृणोतु= विदधातु ।

\customparagraph{\textbf{भावार्थः}}

धर्मस्य बहुशाखित्वं शास्त्रेषु प्रसिद्धं तथाप्यत्र प्रमुखत्वेन अष्टावेव परिगणिता विद्यन्ते । ते च सत्यम् (भौतिकम्), बृहत् (गाम्भीर्यम्), ऋतम् (मानसिकम्), उग्रम् (शाशकगुणम्), दीक्षा (अधिकारः), तपः, ब्रह्म (वेदः), यज्ञश्चैता अष्टौ धर्मशाखाः । एता धर्मशाखा अस्माकं निवासभूतां पृथिवीं धारयन्ति । एतादृशी भूतस्य भाविनश्च सत्यस्य बृहतः ऋतस्य उग्रस्य दीक्षायास्तपसो ब्रह्मणो यज्ञस्य च पालयित्रीयंं पृथ्वी अस्माकं कृते विस्तीर्णम् आवासभूमिं ददातु ।

\customparagraph{\textbf{नेपाली भावार्थ}}

पृथिवीमा विभिन्न धर्मशाखाहरू रहेका छन्, तीमध्ये बोलीको सत्य, धैर्य, मानसिक सत्य, शासकपन, अधिकार र योग्यता अनुरूप कर्तव्य पालन, तप, वेद, यज्ञ जस्ता आठ प्रमुख धर्मका शाखाले पृथिवीलाई धारण गरेका छन् । यस्ती भूत र भविष्यका सत्यादिधर्मका स्वरूपलाई संरक्षण गर्ने पृथिवीले हामीलाई विस्तारसहितको निवासस्थान दिऊन्।

\begin{visheshshloka}[center]

\textbf{असम्बाधं बध्यतो मानवानां यस्या उद्वतः प्रवतः समं बहु ।}

\textbf{नानावीर्या ओषधीर्या बिभर्ति पृथिवी नः प्रथतां राध्यतां नः ॥२॥}

\end{visheshshloka}

\customparagraph{\textbf{अन्वयः}}

उद्वतः प्रवतः यस्याः बध्यतः मानवानाम् (कृते) असम्बाधं समं बहु ‘लोकं धारयन्ति’ । नानावीर्याः ‌ओषधीः बिभर्ति। (सा) पृथिवी नः प्रथताम् । (तथा) नः राध्यताम् ।

\customparagraph{\textbf{शब्दार्थः}}

उद्वतः= देवताः प्रवतः= प्रकृष्टगतयो महर्षयः सनकाद्या मनुष्याः, यस्याः= पूर्वोक्तायाः पृथिव्याः प्रदेशे बध्यतः= बध्यतां सासारिककर्मपाशे निबद्धानाम्, मानवानाम्= मनुष्याणाम् कृते, असम्बाधम्= बाधारहितम्, समम्= समानम् साधारणम् बहु= नानाप्रकारकं लोकम्= स्थानम् धारयन्ति स्थापयन्तीत्यर्थः। एवमेव या= पृथिवी, नानावीर्याः= विविधप्रभावसम्पन्नाः ‌ओषधीः= आमयनाशिनीर्हरितक्यादिनीः सत्त्वविवर्धिनीर्गोधूमधान्यादिनीश्चैता उभयविधा ओषधीः, बिभर्ति= धारयति । सा= पृथिवीः, नः= अस्मभ्यम्, प्रथताम्= विस्तृता भवतु । तथा सैषा पृथिवी नः= अस्माकम्, राध्यताम्= अभीप्सितं संसाधयतु ।

\customparagraph{\textbf{भावार्थः}}

देवैर्मुनिभिश्च पृथिव्यां सांसारिकेषु कर्मबन्धनेषु निबद्धानां साधारणमनुष्याणां कृते बाधारहितानि विविधानि पुण्यानि स्थानानि स्थापितानि सन्ति। या अनेकविधान् व्याधिनाशसत्त्ववर्धनसमर्थान्  ओषधिविशेषान् धारयति, तादृशी पृथिवी अस्माकं कृते विस्तृता भवतु, एवम् अस्माकम् अभीप्सितमपि संसाधयतु ।

\customparagraph{\textbf{नेपाली भावार्थ}}

देवता र दिव्य चरित्र भएका सनकादि मुनिहरूले पृथिवीमा  सांसारिक कर्मबन्धनमा बाँधिएका मानिसका लागि मुक्ति प्राप्त गर्ने विभिन्न पवित्र स्थानहरू स्थापित गरेका छन् । त्यस्तै पृथिवीले विभिन्न किसिमका रोग निको पार्ने औषधि र बलवर्धक अन्नहरू धारण गर्दछिन्। त्यस्ती पृथिवी हाम्रा लागि विस्तारवाली होऊन् र तिनले हाम्रा चाहनाहरू पूरा गरून् ।

\begin{visheshshloka}[center]

\textbf{यस्यां समुद्र उत सिन्धुरापो यस्यामन्नं कृष्टयः संबभूवुः ।}

\textbf{यस्यामिदं जिन्वति प्राणदेजत् सा नो भूमिः पूर्वपेये दधातु ॥३॥}

\end{visheshshloka}

\customparagraph{\textbf{अन्वयः}}

यस्यां समुद्रः उत सिन्धुः उत आपः संबभूवुः । यस्यां कृष्टयः अन्नं सम्बभूवुः । यस्याम् इदं प्राणत् एजत् जिन्वति । सा भूमिः नः पूर्वपेये दधातु ।

\customparagraph{\textbf{शब्दार्थः}}

यस्याम्= पृथिव्याम् समुद्रः= अब्धिः अन्तरिक्षं वा तत्र भवा या वृष्टेरापः, उत सिन्धुः= नदी, आपः= आधाराधेययोस्तादात्म्यनिर्देशात् जलाधिकरणानि वापीकूपतडागकानि सन्ति । तत्प्रभावेण यस्याम्= भूम्याम् कृष्टयः= मनुष्याः अन्नम्= गोधूमधान्यादिकम् संवभूबुः= उत्पादयन्तीत्यर्थः । यस्याम्= पृथिव्याम्, इदम्= दृश्यमदृश्यञ्च सर्वं जगत्, प्राणत्= जीवितम्, एजत्= कम्पमानञ्च जिन्वति= प्रीणाति, सा भूमिः नोऽस्मान् पूर्वपेये= पूर्वं यथा उत्पन्नस्य अन्नस्य रसमुपलभेमहि तथा दधातु= धारयतु, पोषयतु ।

\customparagraph{\textbf{भावार्थः}}

यस्यां पृथिव्यां समुद्रस्य माध्यमेेन वर्षायुता भूमयः सन्ति, एवमेव नदी, आपः जलपूर्णस्थलानि यथा वापीकूपतडागकानि सन्ति । एतेषां जलाधारस्थलानां प्रभावेण कृषिकर्मसंलग्नाः कृषका गोधूमधान्यादिकमन्नमुत्पादयन्ति । तथा च यस्यां पृथिव्यां स्थावरजङ्गमादिमयं जगत् प्रीणाति, सा भूमिरस्मान् भूमिवासिनः कृते अन्नरसोपभोगस्थाने स्थापयतु, अर्थात् दधातु ।

\customparagraph{\textbf{नेपाली भावार्थ}}

जुन पृथिवीमा समुुद्रका माध्यमले हुने वर्षाको पानी, नदीको पानी साथै कुवा, तलाउ, पोखरी आदिको पानी छ । यसलाई उपयोग गरी कृषकहरू गहूँ, धान आदि अन्न फलाउँछन् । जुन पृथिवीमा रहेर सम्पूर्ण स्थावरजङ्गमादि प्रसन्न हुन्छन् । त्यस्ती यी भूमिले हामीलाई जीवनका लागि अन्नरस प्रदान गरून् ।

\begin{visheshshloka}[center]

\textbf{यस्याश्चतस्रः प्रदिशः पृथिव्या यस्यामन्नं कृष्टयः संवभूबुः।}

\textbf{या बिभर्ति बहुधा प्राणदेजत् सा नो भूमिर्गोष्वप्यन्ने दधातु॥४॥}

\end{visheshshloka}

\customparagraph{\textbf{अन्वयः}}

यस्याः पृथिव्याः चतस्रः प्रदिशः (सन्ति) यस्यां कृष्टयः अन्नं संवभूबुः । या प्राणत् एजत् बहुधा विभर्ति, सा भूमिः नः गोषु अन्ने अपि दधातु ।

\customparagraph{\textbf{शब्दार्थः}}

यस्याः पृथिव्याः चतस्रः= चतुःसङ्‍ख्यकाः प्रदिशः= पूर्वदक्षिणपश्चिमोत्तरात्मका दिशः सन्ति । यस्यां पृथिव्याम्, कृष्टयः= कृषिकर्मसम्बद्धा मनुष्याः, अन्नम्= गोधूमधान्यादिकम्, संवभूबुः= संभावयन्ति, उत्पादयन्ति । या च पृथिवी प्राणत्= प्राणयुतम्, एजत्= स्पन्दमानञ्च जीवजातम्, बहुधा= बहुभिर्विधाभिः कन्दमूलफलसस्यादिभिर्जीवनोपायैः बिभर्ति= पोषयति । सा भूमिः नः=अस्मान् गोषु= गवोपलक्षणया सर्वेषु पशुषु अन्ने= भक्षणीये पदार्थजाते अपि, दधातु= धारयतु, स्थापयतु । पशुभिरन्नैश्च वयं स्यामेति भावः ।

\customparagraph{\textbf{भावार्थः}}

यस्याः पृथिव्याः पूर्वपश्चिमोत्तरदक्षिणात्मकाश्चतस्रो दिशः सन्ति, यस्यां पृथिव्यां कृषका अन्नमुत्पादयन्ति, या पृथिवी स्थावरजङ्गमादीन् सर्वान् विविधप्रकारेण धारयति । सा भूमिरस्मान् पशुभिरन्नैश्च पूर्णान् विदधातु ।

\customparagraph{\textbf{नेपाली भावार्थ}}

जुन पृथिवीका पूर्व, पश्चिम, उत्तर र दक्षिण गरी चार दिशा छन्; जुन पृथिवीमा कृषकहरू गहुँ, धान आदि अन्न उत्पादन गर्दछन्; जसले प्राणशील र गतिशील चराचरलाई धारण गर्दछिन्; त्यस्ती पृथिवीले हामीलाई पशुधन र अन्नधनले भरिपूर्ण गराऊन् ।

\begin{visheshshloka}[center]

\textbf{यस्यां पूर्वे पूर्वजना वि चक्रिरे यस्यां देवा असुरानभ्यवर्तयन्।}

\textbf{गवामश्वानां वयसश्च विष्ठा भगं वर्चः पृथिवी नो दधातु॥५॥}

\end{visheshshloka}

\customparagraph{\textbf{अन्वयः}}

यस्यां नः पूर्वे पूर्वजनाः वि चक्रिरे  । यस्यां देवाः असुरान् अभ्यवर्तयन् । गवाम् अश्वानां वयसः च विष्ठा पृथिवी नः भगं वर्चः दधातु ।

\customparagraph{\textbf{शब्दार्थः}}

यस्याम्= पृथिव्याम्, नः= अस्माकम्, पूर्वे= प्राचीनकाले, पूर्वजनाः= पूर्वजाः, वि चक्रिरे= विविधम्- आध्यात्मिकम्, आधिदैविकम्, आधिभौतिकम्, अभ्युदयसाधनभूतम्, आदर्शमयं कर्म कृतवन्तः । यस्याम्= पृथिव्याम् देवाः= सत्त्वगुणप्रधाना विभूतयः, असुरान्= तमःप्रधाना विभूतीः  अभ्यवर्तयन्= अभितोऽवर्तन्त, न्यगृह्णन्त, परिजिग्युरितिभावः । यस्यां च गवाम्= एतदुपलक्षणेन दृष्टसत्त्वगुणजीवानाम्, अश्वानाम्= दृष्टरजोगुणजीवानाम्, वयसः= दृष्टतमोगुणजीवानाम्, च विष्ठा= विशेषेण निवासस्थानमस्ति । सेयं पृथिवी नोऽस्मभ्यम् भगम्= सर्वप्रकारकमैश्वर्यम् वर्चः= तेजः (यतो दर्शनादेव विद्वान् वा वीरो वा धनी वा अन्यद् वा किञ्चन वरणीयगुणवत्वं विशेषतः प्रतीयते तत्तेजो वर्च इत्युच्यते) दधातु= स्थापयतु इति आशीः ।

\customparagraph{\textbf{भावार्थः}}

यस्यां पृथिव्याम् अस्माकं पूर्वजा विविधानि गौरवयुतानि कार्याणि कृतवन्तः । यत्र सात्त्विकविभूतियुता देवाः तामसिकविभूतियुतान् दानवान् उद्दण्डकार्येभ्यः संवार्य वशीकृतवन्तः। एवमेव इयं भूमिर्गवादिसत्त्वगुणप्रधानानां जीवानाम् अश्वादिरजोगुणप्रधानानां जीवानाम् वयस इति तमोगुणप्रधानानां पक्षिणाञ्च निवासस्थानमस्ति । एतादृशी गौरवगुणयुता पृथिवी, अस्मभ्यम् ऐश्वर्यं विशिष्टाकर्षणचिह्नमयं तेजश्च प्रदधातु ।

\customparagraph{\textbf{नेपाली भावार्थ}}

जुन पृथिवीमा हाम्रा पूर्वजहरूले विभिन्न प्रकारका गौरवपूर्ण महान् कार्यहरू गरेका थिए । जहाँ सात्त्विकगुणप्रधान देवताहरूले आसुरी शक्ति भएका दानवहरूलाई आफ्नो वशमा राखे । जुन पृथिवी सत्त्वगुणप्रधान गाई आदि प्राणी, रजोगुण प्रधान घोडा आदि प्राणी र तमोगुणप्रधान चराचुरुङ्गीको आवासभूमि छ । यस्ती गौरवपूर्ण गुण भएकी पृथिवीले हामीलाई ऐश्वर्य र आकर्षक तेज प्रदान गरून् ।

\begin{visheshshloka}[center]

\textbf{विश्वंभरा वसुधानी प्रतिष्ठा हिरण्यवक्षा जगतो निवेशनी।}

\textbf{वैश्वानरं बिभ्रती भूमिरग्निमिन्द्रऋषभा द्रविणे नो दधातु॥६॥}

\end{visheshshloka}

\customparagraph{\textbf{अन्वयः}}

विश्वम्भरा वसुधानी प्रतिष्ठा, हिरण्यवक्षाः जगतः निवेशनी वैश्वानरम् अग्निं बिभ्रती इन्द्र ऋषभा भूमिः नः द्रविणे दधातु ।

\customparagraph{\textbf{शब्दार्थः}}

विश्वम्भरा= (विश्वं बिभ्रतीति विश्वम्भरा) विश्वस्य भरणकर्त्री, वसुधानी= मणिरत्नादिधनानां धारयित्री, प्रतिष्ठा= स्थावराद्यात्मकविश्वस्याश्रयरूपा (गृहरूपत्वात्) हिरण्यवक्षाः= (हिरण्यं वक्षस्तद्रूपमस्या) अतिकठिनत्वात् आकरात्मना मध्यस्थत्वाच्च हिरण्यस्य वक्षस्त्वमुक्तम्, वक्षसि काठिन्यस्यैव प्रशस्तत्त्वात्। जगतः= स्थावरजङ्गमात्मकस्य निवेशिनी= आयातनभूता वैश्वानरम्= विश्वस्मिन्नपि प्रतिगतमग्निम्, बिभ्रती= अन्तर्धारयती, इन्द्रऋषभा= वर्षकर्मकर्ता इन्द्रेण सिक्ता  इयं भूमिः= पृथिवी नः= अस्मान् द्रविणे= बले धने वा दधातु= स्वामित्वेन स्थापयतु ।

\customparagraph{\textbf{भावार्थः}}

स्थावरजङ्गमात्मकस्य विश्वस्य भरणकर्त्री, मणिरत्नादिधनानां धारयित्री, सुवर्णखनिरूपं कठोरं वा वक्षं धारयित्री, स्थावरजङ्गमात्मकस्य जगतो वासस्थानभूता, अन्तर्वैश्वानरंं धारयित्री, इन्द्रेण वर्षाजलेन सिक्ता, एतादृशी ‘भूमिः’ अस्माकं कृते धनसामर्थ्यादिगुणान् विदधातु ।

\customparagraph{\textbf{नेपाली भावार्थ}}

सम्पूर्ण विश्वलाई भरणपोषण गर्ने, मणिरत्न आदि धनलाई धारण गर्ने, खानीका रूपमा सुनको छाती भएकी, स्थावरजङ्गमादिलाई आवासस्थान प्रदान गर्ने, सर्वत्र प्रसरणशील अग्निलाई धारण गर्ने, इन्द्रले वर्षाद्वारा सिञ्चन गरिएकी पृथिवीले हामीलाई बल, धन आदि सामर्थ्यले पूर्ण गरून् ।

\begin{visheshshloka}[center]

\textbf{यां रक्षन्त्यस्वप्ना विश्वदानीं देवा भूमिं पृथ्वीमप्रमादम्।}

\textbf{सा नो मधु प्रियं दुहामथो उक्षतु वर्चसा॥७॥}

\end{visheshshloka}

\customparagraph{\textbf{अन्वयः}}

विश्वदानीं भूमिं यां पृथ्वीम् अस्वप्नाः देवाः अप्रमादं रक्षन्ति । सा नः प्रियं मधु दुहाम् अथो वर्चसा उक्षतु ।

\customparagraph{\textbf{शब्दार्थः}}

विश्वदानीम्= सर्वस्याभिमतस्य दात्रीम्, भूमिम्= भूमगुणवतीम् (एकस्य बीजस्यानन्तफलदातृत्वात्) यां पृथ्वीम्= भूमिम्, अस्वप्नाः= स्वापो निद्रा तद्रहिताः, जागरूकाः, देवाः= अमराः, अप्रमादम्= सर्वदा सावधानाः सन्तः, रक्षन्ति= पालयन्ति । सा= पृथ्वी नोऽस्मभ्यम्, प्रियम्= अभीप्सितम्, मधुु= मधुरं फलम्, दुहाम्= दोग्धु । अथ= अपि च वर्चसा= तेजसा उक्षतु= सिञ्चतु । (वयं यथा वर्चस्विनो भूयास्म तथा कुर्यात्)

\customparagraph{\textbf{भावार्थः}}

भूमिवासिनः कृते ईप्सितफलप्रदात्रीम्, उत्पादकगुणवतीं पृथ्वीं स्वापरहिता जागरूकाः देवाः सर्वदा सावधाना सन्तः पालयन्ति । तादृशी पृथ्वी, अस्मभ्यमभीप्सितं मधुरं फलादिकं ददातु । एवमेव सास्मान् आकर्षकगुणमयेन तेजसा पूर्णं करोतु ।

\customparagraph{\textbf{नेपाली भावार्थ}}

पृथ्वीमा बसोबास गर्ने सम्पूर्ण लोकवासीलाई मनले चाहेको फल दिने, उत्पादकत्व शक्तिसम्पन्न पृथ्वीलाई कहिल्यै ननिदाउने देवताले सावधान भएर रक्षा गर्दछन् । त्यस्ती पृथ्वीले हामीलाई इच्छाएको फल दिऊन्, साथै वर्च (कर्तव्य शक्ति) ले पूर्ण गरून् ।

\begin{visheshshloka}[center]

\textbf{यार्णवेऽधिसलिलमग्र आसीद् यां मायाभिरन्वचरन् मनीषिणः।}

\textbf{यस्या हृदयं परमे व्योमन् सत्येनावृतममृतं पृथिव्याः।}

\textbf{सा नो भूमिस्त्विषिं बलं राष्ट्रे दधातूत्तमे॥८॥}

\end{visheshshloka}

\customparagraph{\textbf{अन्वयः}}

या अग्रे अर्णवे अधिसलिलम् आसीत् । यां मनीषिणः मायाभिः अन्वचरन् । यस्याः पृथिव्याः अमृतं हृदयं परमे व्योमन् सत्येन आवृतम् (अस्ति) सा भूमिः नः उत्तमे राष्ट्रे त्विषिं बलं दधातु ।

\customparagraph{\textbf{शब्दार्थः}}

या= पृथिवी, अग्रे= पूर्वम्, अर्णवे= समुद्रे, अधिसलिलम्= जलान्तर्निमग्ना आसीत् । याम्= पृथिवीम् मनीषिणः= मनसः ईशयितारो (ऐश्वर्यशालिनः) मनोनिग्रहितारस्तपश्चर्यान्विता हिरण्याक्षादयः, मायाभिः= प्रकृष्टज्ञानैः, अन्वचरन्= भूम्याः सविधे सलिल एवासीदन् (प्राप्तवन्तः) । यस्याः पृथिव्याः, अमृतम्= अविनाशीहृदयम् जीवभूतं देवयजनस्थानम्, परमे= उत्कृष्टे व्योमन्= व्योम्नि, सत्येन= सत्यस्वरूपेण देवजातेन, आवृतम्= सर्वतोभावेनाङ्गीकृतम् अस्ति । सा भूमिर्नोऽस्माकमुत्तमे= प्रशंसनीये, राष्ट्रे= देशे, त्विषिम्= दीप्तिम्, बलम्= वीर्यम्, दधातु= स्थापयतु ।

\customparagraph{\textbf{भावार्थः}}

इयं पृथिवी पुरा समुद्रजले आसीत् । यां हिरण्याक्षादयो विचक्षणाः सलिलान्तर एव प्राप्तवन्तः । अस्याः पृथिव्याः अमृतं जीवभूतं हृदयं द्युलोके (चन्द्रमसि) सत्त्वगुणयुतैर्देवैरावृतं विद्यते । तादृशी भूमिरस्माकं सर्वगुणसम्पन्ने देशे बलम्, वीर्यञ्च स्थापयतु ।

\customparagraph{\textbf{नेपाली भावार्थ}}

यी पृथिवी पहिले समुद्रजलमा अवस्थित थिइन् । त्यहाँ तपस्याबाट शक्ति प्राप्त गरेका हिरण्याक्ष आदिले यिनलाई जलमै आफ्नो नियन्त्रणमा राखेका थिए । जुन पृथिवीको अमृतस्वरूप जीवभूत हृदय आकाशमा स्थित चन्द्रमामा सत्त्वगुणी देवताहरूबाट संरक्षित छ । त्यस्ती पृथिवीले हाम्रो उत्कृष्ट देशलाई बल र वीर्य सम्पन्न गराऊन् ।

\begin{visheshshloka}[center]

\textbf{यस्यामापः परिचराः समानीरहोरात्रे अप्रमादं क्षरन्ति।}

\textbf{सा नो भूमिर्भूरिधारा पयो दुहामथो उक्षतु वर्चसा॥९॥}

\end{visheshshloka}

\customparagraph{\textbf{अन्वयः}}

यस्याम् अहोरात्रे परिचराः आपः समानीः अप्रमादं क्षरन्ति । भूरिधारा सा भूमिः नः पयः दुहाम्, अथो वर्चसा रक्षतु ।

\customparagraph{\textbf{शब्दार्थः}}

यस्याम्= पृथिव्याम् परिचराः= स्रोतस्यः (निर्झरिण्यः) आपः= वारीणि समानीः= समानपरिमाणेन, अप्रमादम्= अनवधानताविरहेण, अहोरात्रे= रात्रिन्दिवम्, क्षरन्ति= श्च्योतन्ते (प्रवहन्ति) । एवं स्रोतसां नानात्वात् भूरिधारा= प्रभूतधारावती, सा विश्वस्य धारयितृत्वादिना पूर्वमुक्ता भूमिः नोऽस्माकंं कृते पयः= वृष्टिलक्षणमुदकम्, दुहाम्= दोग्धु। अथो= अथ च वर्चसा= अन्नेन, अस्मान् उक्षतु= सिञ्चतु, समृद्धान् विदधातु ।

\customparagraph{\textbf{भावार्थः}}

यस्यां पृथिव्यां  निर्झरीणि वारीणि पर्वतमस्तकात् प्रमादराहित्येन शान्तरूपेण अधः पतन्ति । स्रोतसां विविधत्वात् प्रभूतधारावती पृथिवी अस्माकं कृते वृष्टिरूपमुदकं दोग्धु । अन्नादिविभूतिभिः अस्मभ्यं समृद्धं करोतु ।

\customparagraph{\textbf{नेपाली भावार्थ}}

जुन पृथिवीमा झरनाको जल रातदिन निरन्तर शान्तस्वभावले प्रवाहित हुन्छन्  । यस्ता पानीका झरना, नदी, ताल आदि स्रोतहरूलाई धारण गर्ने पृथिवीले हामीलाई वर्षाको पानी पनि उपलब्ध गराऊन् । साथै अन्नरूप सम्पत्तिले हामीलाई पूर्ण गराऊन् ।

\begin{visheshshloka}[center]

\textbf{यामश्विनावमिमातां विष्णुर्यस्यां विचक्रमे।}

\textbf{इन्द्रो यां चक्र आत्मनेऽनमित्रां शचीपतिः।}

\textbf{सा नो भूमिर्विसृजतां माता पुत्राय मे पयः॥१०॥}

\end{visheshshloka}

\customparagraph{\textbf{अन्वयः}}

याम् अश्विनौ अमिमाताम् । अस्यां विष्णुः विचक्रमे । यां शचीपतिः इन्द्रः आत्मने अनमित्रां चक्रे । सा नः भूमिः माता पुत्राय मे पयः विसृजताम् ।

\customparagraph{\textbf{शब्दार्थः}}

याम्= भूमिम् अश्विनौ= देवौ अमिमाताम्= इयत्तया परिमाणं निश्चिक्युः (निश्चितं कृतवन्तः)। यस्याम्= पृथिव्याम्, विष्णुः= नारायणः विचक्रमे= विशेषेण त्रिभिः पादैः क्रान्तवान् । यां= पृथिवीम्, शचीपतिः= कर्मणामधिपतिः, यद्वा शचीनामिकायाः पत्न्याः पतिः, इन्द्रः= देवराजः आत्मने= स्वस्य हिताय, अनमित्राम्= शत्रुरहिताम्, चक्रे= विदधे । सा= अश्विनीकुमाराभ्यां मिता, विष्णुपदैरतिक्रान्ता, इन्द्रेण अशात्रवीकृता नोऽस्माकमस्मत्सम्बन्धिनी भूमिः, माता= पृथिवी पार्थिवशरीरधारित्वात्, पुत्राय= तनुजाय मे= मह्यम् पयः‍‍= आप्यायनसाधनमन्नरसादि, विसृजताम्= प्रददातु ।

\customparagraph{\textbf{भावार्थः}}

यां पृथिवीमश्विनीकुमारौ कति परिणामिका पृथिवी इति जिज्ञासाशमनाय पूर्वम् अमिमाताम् । यां वामनावताररूपेणावतीर्णो भगवान् विष्णुः त्रिभिः पादैः क्रान्तवान् । यां देवराज इन्द्रः शत्रुरहितां कृतवान् । तादृशी मातृरूपा पृथिवी मह्यम् अन्नरसं ददातु।

\customparagraph{\textbf{नेपाली भावार्थ}}

जुन पृथिवीलाई अश्विनीकुमारले मापन गरे । जुन पृथिवीलाई विष्णुले वामन अवतार लिई तीन पाउले ढाके । जुन पृथिवीलाई देवराज इन्द्रले आफ्ना लागि शत्रुरहित बनाए । त्यस्ती आमाजस्ती पृथिवीले पुत्रतुल्य मलाई अन्नरस प्रदान गरून्।

\begin{visheshshloka}[center]

\textbf{गिरयस्ते पर्वता हिमवन्तोऽरण्यं ते पृथिवि स्योनमस्तु।}

\textbf{बभ्रुं कृष्णां रोहिणीं विश्वरूपां ध्रुवां भूमिं पृथिवीमिन्द्रगुप्ताम्।}

\textbf{अजीतोऽहतो अक्षतोऽध्यष्ठां पृथिवीमहम्॥११॥}

\end{visheshshloka}

\customparagraph{\textbf{अन्वयः}}

हे पृथिवि ! ते पर्वताः गिरयः हिमवन्तः ‘सन्तु’ । ते अरण्यं स्योनम् अस्तु । अहम् अजीतः अहतः अक्षतः (सन्) बभ्रुं कृष्णां रोहिणीं विश्वरूपां ध्रुवां भूमिं पृथिवीम् इन्द्रगुप्तां पृथिवीम् अध्यष्ठाम् ।

\customparagraph{\textbf{शब्दार्थः}}

हे पृथिवि ! ते= तव त्वत्सम्बन्धिनः पर्वताः= अद्रयः, गिरयः= अन्नवन्तः, मेघवन्तो वा हिमवन्तः= हिमयुताः शीतलाः सन्तु । ते= तव अरण्यम्= सिंहव्याघ्रचौरादिभिया रमणीयतामनाप्तं वनम्, स्योनम्= सुखकरम् अस्तु । अहम्= स्तुतिकर्ता अजीतः= परैरनभिभूतः शत्रुभिरपराजितः, अहतः= केनाप्यहिंसितः, अक्षतः= केनाप्युपकरणेन अहीनः सन्, बभ्रुम्= पिङ्गलवर्णाम्, कृष्णाम्= श्यामवर्णाम्, रोहिणीम्= लोहितवर्णाम्, विश्वरूपाम्= चित्ररूपाम्, ध्रुवाम्= कठिनाम्, भूमिम्= भूमगुणवतीं विप्रकीर्णां पृथिवीम्= पृथुत्वगुणसम्पन्नां विशालामित्यर्थः, इन्द्रगुप्ताम्= पर्जन्यस्वामिना देवराजेन रक्षिताम्, पृथिवीम्= भूमिम्, अध्यष्ठाम्= अधितिष्ठेयम्, स्वामित्वेनावस्थितो भूयासमिति याञ्चा ।

\customparagraph{\textbf{भावार्थः}}

हे पृथिवि ! तव पर्वताः हिमवन्तो जलवन्तोऽन्नवन्तश्च सन्तु । तव वनानि भयरहितानि सुखदायकानि सन्तु । अहं स्तुतिकारकः शत्रुभिरपराजितः परैरहिंसितः केनाप्युपकरणेनाहीनश्च सन् पिङ्गलाम्, कृष्णाम्, रक्तवर्णाम्, बालुकमयाञ्च भूमिम्; एवमेव पर्जन्येन इन्द्रेण पालितां कठोरां विशालाञ्च पृथिवीं स्वामित्वेनाप्नुयाम् (पृथिव्यां स्वामित्वेनावस्थितो भूयासम् ) ।

\customparagraph{\textbf{नेपाली भावार्थ}}

हे पृथिवी ! तिमीमा अवस्थित पर्वतहरू शीतल जल तथा अन्नले सम्पन्न बनून् । तिम्रा वनहरू बाघ, भालु, चोर आदिको भयबाट मुक्त होऊन् । म शत्रुहरूद्वारा अपराजित, कसैद्वारा अनाहत र विभिन्न उपकरणहरूले सम्पन्न भएर पहेँलो, कालो, रातो, बलौटेजस्ता विभिन्न स्वरूपको माटो भएको, इन्द्रद्वारा वर्षाका माध्यमले संरक्षित कठोर र विशाल पृथिवीको अधिपति बनूँ ।

\begin{visheshshloka}[center]

\textbf{यत्ते मध्यं पृथिवि यच्च नभ्यं यास्त ऊर्जस्तन्वः संबभूवुः।}

\textbf{तासु नो धेह्यभि नः पवस्व माता भूमिः पुत्रो अहं पृथिव्याः।}

\textbf{पर्जन्यः पिता स उ नः पिपर्तु॥१२॥}

\end{visheshshloka}

\customparagraph{\textbf{अन्वयः}}

हे पृथिवि ! यत् ते मध्यम्, च यत् ‘ते’ नभ्यम्, ‘च’ ते याः ऊर्जः तन्वः संबभूवुः । नः तासु धेहि । नः अभिपवस्व । उ भूमिः माता (अस्ति) अहं पृथिव्याः पुत्रः (अस्मि) । पर्जन्य पिता (अस्ति) । सः नः पिपर्तु ।

\customparagraph{\textbf{शब्दार्थः}}

हे पृथिवि ! यत् ते मध्यम्= मध्यस्थानीयो जम्बूद्वीपो विद्यते । च= तथा यत् ते नभ्यम्= नाभिस्थानीयो भारतवर्षः, अभयदातृत्वात्= अभयं नाभिः, इति श्रुत्यवगमात् भारतस्य अभयदातृत्वेन पुराणेषु स्मृतत्त्वाच्चेति । तथा ते= तव पृथिव्या याः ऊर्जः= बलदातॄणि, प्राणदातॄणि (‘ऊर्ज बलप्राणनयोः’ आध्यात्मिकधिदैविकाधिभौतिकाख्यतापत्रयोपशमनेन बलदातॄणि विशिष्टादृश्यशक्तिसञ्चारणेन प्राणदातॄणि) च तन्वः= शरीरावयवभूतानि विस्तीर्णानि वा पुण्यक्षेत्रतीर्थादीन्यायतनानि संबभूवुः= पूर्वतनानां देवतानां महर्षिणाञ्चानुग्रहेण सञ्जातानि सन्ति । नः‍= अस्मान् तासु= देशभूमिषु धेहि= निवासार्थं स्थापय । तथा नः= अस्मान् अभिपवस्व= सर्वविधरसान्नादिसमृद्ध्या स्थिरान्तवृत्तीन् रागद्वेषादिभिर्दोषैः संशोधय, दोषेष्वस्माकं प्रवृत्तिर्मा भवत्वित्यर्थः । यतः अस्माकंं भूमि माता अस्ति । तथा च अहम्= स्तोता पुत्रोऽस्मि, पुंनामनरकात् त्राता इति । तथा पर्जन्यः= जलदेवो नोऽस्माकम् पिता अस्ति, अन्नाद्युत्पादनद्वारा पालयितृत्वात् । सः= पर्जन्यो देवः, नोऽस्मान् पिपर्तु= पालयतु, यथाकालं वर्षणद्वारास्माकं मनोरथान् पूरयतु वा ।

\customparagraph{\textbf{भावार्थः}}

हे पृथिवि ! यत् ते मध्यभागे जम्बूद्वीपो विद्यते । अभयदातृत्वात् नाभिस्थले भारतवर्षो विद्यते । ते बलदातॄणि प्राणदातॄणि च शरीररूपाणि सन्ति । तत्र अस्मान् निवासार्थं स्थापय । अस्माकमवगुणान् दूरीकुरु । पुत्रवत् अन्नरसपानेन संवर्धनात् सा भूमि मम माता अस्ति । पुंनामनरकतारकत्वात् अहं पृथिव्याः पुत्रोऽस्मि । वर्षाजलसिञ्चनात् तेन पालनाच्च पर्जन्यः पिता अस्ति । अतः पर्जन्यपिता अस्मान् पालयतु पूरयतु वा ।

\customparagraph{\textbf{नेपाली भावार्थ}}

हे पृथिवी ! तिम्रो मध्यभागमा जम्बूद्वीप अवस्थित छ । भारतवर्ष नाभिस्थलका रूपमा रहेको छ । तिम्रा जुन अवयवहरू बल र प्राणप्रदायक छन्, त्यस स्थानमा हामीलाई राख । हाम्रा सबै दोषहरू हटाऊ किनकि अन्नरसादि पिलाएर संवर्धन गराउने तिमी मेरी माता हौ । स्तुतिकर्ता म तिम्रो पुत्र हुँ । वर्षा गराएर अन्नादि उत्पादन गराई तिमीलाई समर्थ बनाउने वर्षाका देवता पर्जन्य मेरा पिता हुन् । ती पितृरूप वर्षाका देवताले हामीलाई पालन गरून्, अथवा सर्वसामग्रीसम्पन्न बनाऊन् ।

\begin{visheshshloka}[center]

\textbf{यस्यां वेदिं परिगृह्णन्ति भूम्यां यस्यां यज्ञं तन्वते विश्वकर्माणः।}

\textbf{यस्यां मीयन्ते स्वरवः पृथिव्यामूर्ध्वाः शुक्रा आहुत्याः पुरस्तात्।}

\textbf{सा नो भूमिर्वर्धयद् वर्धमाना॥१३॥}

\end{visheshshloka}

\customparagraph{\textbf{अन्वयः}}

यस्यां भूम्यां वेदिं परिगृह्णन्ति । यस्यां विश्वकर्माणः यज्ञं तन्वते । यस्यां पृथिव्याम् आहुत्याः पुरस्तात् शुक्राः ऊर्ध्वाः स्वरवः मीयन्ते । सा वर्धमाना भूमिः नः वर्धयत् ।

\customparagraph{\textbf{शब्दार्थः}}

यस्यां भूम्यां विद्वांसः, वेदिम्= देवयजनस्थानम्, परिगृह्णन्ति= स्वीकुर्वन्ति । यस्यां च भूम्यां विश्वकर्माणः= नानाविधक्रियाकुशलाः श्रौतस्मार्तपौराणतान्त्रिककर्मविदः स्वस्वागमानुकूल्येन, यज्ञम्= देवोद्देश्येन द्रव्यत्यागात्मकं कर्म, तन्वते= विस्तारयन्ति । यस्यां पृथिव्याम्, आहुत्याः= देवताह्वानलक्षणाया आहूतेः (आहूतिरेव परोक्षवृत्त्या आहुतिशब्देन निर्दिश्यते इति ऐ. ब्राह्मणे स्पष्टमुक्तत्त्वात्) पुरस्तात्= पूर्वम्, शुक्राः= शुद्धाः पूताः, ऊर्ध्वाः= यूपस्य शिरःप्रदेशीयाः, स्वरवः= स्वरुनामकयज्ञसाधनानि (छिद्यमानस्य यूपस्य प्रथमं भूमौ पतितः काष्ठखण्डः स्वरुः कथ्यते, बहुत्वविवक्षया बहुवचनम्) मीयन्ते= निर्मीयन्ते । सा= एतादृशी वर्धमाना= विशाला भूमिः= पृथिवी नः= अस्मान् सर्वतोभावेन वर्धयत्= वर्धयतु ।

\customparagraph{\textbf{भावार्थः}}

यस्यां भूम्यां विद्वांसो यज्ञार्थं वेदिं स्वीकुर्वन्ति । यस्यां भूम्यां श्रौत्रस्मार्तपौराणतान्त्रिककर्मविदो विद्वांसः स्वस्वागमानुकूल्येन यज्ञं विस्तारयन्ति । यस्यां पृथिव्यां देवाह्वानपूर्वं शुद्धा यूपस्य शिरःप्रदेशीयाः स्वरवः (स्वरुनामकयज्ञसाधनानि) निर्मीयन्ते । एतादृशी विशालविस्तारयुता भूमिरस्मान् सर्वतोभावेन वर्धयतु ।

\customparagraph{\textbf{नेपाली भावार्थ}}

जुन भूमिमा विद्वान्‌हरू यज्ञका लागि वेदि बनाउँछन् । जुन भूमिमा विभिन्न शास्त्रवेत्ता विद्वान् होताहरू यज्ञ विस्तार गर्दछन् । जुन भूमिमा देवतालाई आह्वान गर्नु पूर्व यूपका लागि स्थापित स्तम्भको अघिल्लो भागबाट स्वरु (सुरो) बनाइन्छ । त्यस्ती गुणवती विशालविस्तार भएकी भूमिले हामीलाई सर्वगुणसम्पन्न बनाऊन् ।

\begin{visheshshloka}[center]

\textbf{यो नो द्वेषत्पृथिवि यः पृतन्याद्योऽभिदासान्मनसा यो वधेन।}

\textbf{तं नो भूमे रन्धय पूर्वकृत्वरि॥१४॥}

\end{visheshshloka}

\customparagraph{\textbf{अन्वयः}}

हे पृथिवि ! नः यः द्वेषत्, यः पृतन्यात्, यः मनसा अभिदासात्, यः वधेन ‘अभिदासात्’ । हे पूर्वकृत्वरि ! भूमे ! तं रन्धय ।

\customparagraph{\textbf{शब्दार्थः}}

हे पृथिवि ! नः= अस्मान् यः= प्राणी द्वेषत्= द्वेषं करोति । यश्च जनः पृतन्यात्= अस्माकं कृते पृतनां सेनामेकत्रीकर्तुमभिवाञ्छति । यश्च जनः मनसा= चित्तेन मानसिकव्यापारेण अभिदासात्= सर्वत उपक्षिणोति । यश्च जनः नः= अस्मान् वधेन= शारीरिकव्यापारेण अभिदासात्= हन्तुमिच्छति । हे पूर्वकृत्वरि ! अभिलषितपदार्थनामाकरभूतत्त्वात् प्रागेव सर्वस्य कार्यजातस्य विधात्रि भूमे ! पृथिवि ! तम्= द्वेषिणं जनम्, रन्धय= जहि । (विद्वेषिणो हिंसाभावस्य जागरणात्पूर्वमेव तादृशं दूषितमनोयुतं जनं जहीत्यर्थः)

\customparagraph{\textbf{भावार्थः}}

हे पृथिवि ! यो जनोऽस्माभिः सह द्वेषं करोति । यो जनोऽस्माकं तिरस्कारार्थं सैन्यसञ्चयं करोति । यो जनो मानसिकव्यापारेण तिरस्करोति । यो जनः शारीरिकव्यापारेण हन्तुमभिलषति । विलक्षणसामर्थ्यत्वात् पूर्वमेव कार्यसम्पादयित्रि हे पृथिवि ! दुरभिलाषायुतं जनं कार्यारम्भात् प्रागेव जहि ।

\customparagraph{\textbf{नेपाली भावार्थ}}

हे पृथिवि ! हामीसँग जसले वैरभाव राख्दछ । जसले हामीलाई पराजित गर्न सैन्यदल बनाउँछ । जसले मानसिकरूपले तिरस्कार गर्दछ। जसले भौतिक आक्रमण गरी मार्ने इच्छा राख्दछ । सामर्थ्य भएकाले पहिले नै कार्यसम्पादन गर्न सक्ने हे देवि पृथिवि ! त्यस्ता मान्छेलाई मनमा कुविचार आउन नपाउँदै (कुत्सित कार्य प्रारम्भ नगर्दै) संहार गर ।

\begin{visheshshloka}[center]

\textbf{त्वज्जातास्त्वयि चरन्ति मर्त्यास्त्वं बिभर्षि द्विपदस्त्वं चतुष्पदः।}

\textbf{तवेमे पृथिवि पञ्च मानवा येभ्यो ज्योतिरमृतं मर्त्येभ्य उद्यन्त्सूर्यो}

\textbf{रश्मिभिरातनोति॥१५॥}

\end{visheshshloka}

\customparagraph{\textbf{अन्वयः}}

हे पृथिवि ! त्वज्जाताः मर्त्याः त्वयि चरन्ति । त्वं द्विपदः बिभर्षि । त्वं चतुष्पदः ‘बिभर्षि’ । इमे पञ्चमानवाः तव (सन्ति) येभ्यः मर्त्येभ्यः उद्यन् सूर्यः रश्मिभिः अमृतं ज्योतिः आतनोति ।

\customparagraph{\textbf{शब्दार्थः}}

हे पृथिवि ! त्वज्जाताः= त्वय्युत्पन्नाः, मर्त्याः= मनुष्याः, त्वयि=  त्वत्प्रदेशेषु, चरन्ति= भ्रमन्ति । त्वं हि द्विपदः= पादद्वयवतो मुनष्यादिकान् प्रणिनः, बिभर्षि= धारयसि पोषयसि वा । एवम् त्वं चतुष्पदः= पादचतुष्टयुतः पश्वादिकांश्च बिभर्षि । इमे= दृग्बुद्धिगोचराः, पञ्च= पञ्चसङ्‍ख्यकाः, मानवाः= ब्राह्मणक्षत्रियवैश्यशूद्रनिषादाख्याः, गन्धर्वासुरपितृदेवरक्षांसि वा तव आश्रिताः सन्ति । येभ्यः= पञ्चभ्यो मर्त्येभ्यः, उद्यन्= उदयं गच्छन्नयं सूर्यः= भानुः, रश्मिभिः= किरणैः, अमृतम्= मृत्युविपरीतम्, ज्योतिः= प्रकाशम्, आतनोति= विस्तारयति ।

\customparagraph{\textbf{भावार्थः}}

हे पृथिवि ! त्वयि समुद्भूता मरणशीलाः प्राणिनस्त्वत्प्रदेशेषु सञ्चरन्ति । त्वं द्विपदयुताञ्चतुष्पदयुतान् प्राणिनो धारयसि पोषयसि वा । एवमेव दृग्बुद्धिप्रत्यक्षा ब्राह्मण-क्षत्रिय-वैश्य-शूद्र-निषादाख्याः,गन्धर्व-पित्र-सुर-देव-रक्षांसीति वा पञ्चमानवास्त्वयि समाश्रिताः सन्ति । येभ्यः पञ्चभ्यः उदीयमानोऽयं सूर्यो रश्मिभिरमृतमयं स्वास्थ्यवर्धकं प्रकाशं विस्तारयति।

\customparagraph{\textbf{नेपाली भावार्थ}}

हे पृथिवि ! तिमीमा उत्पन्न मरणशील प्राणीहरू तिम्रै प्रदेशहरूमा विचरण गर्दछन् । तिमी दुई पाउ भएका र चार पाउ भएका प्राणीहरूलाई धारण र पोषण गर्दछ्यौ । त्यस्तै आँखा वा ज्ञानबाट प्रत्यक्ष हुने ब्राह्मण, क्षत्रिय, वैश्य, शूद्र, निषाद अथवा गन्धर्व, पितृगण, देवता, असुर, राक्षस गरी पाँच मानवकी आश्रय दाता हौ । यी पाँच प्रकारका मानवलाई उदाउँदै गरेका सूर्यले आफ्ना किरणद्वारा अमृतमय स्वास्थ्यवर्धक प्रकाश छर्दछन् ।

\begin{visheshshloka}[center]

\textbf{ता नः प्रजाः सं दुह्रतां समग्रा वाचो मधु पृथिवि धेहि}

\textbf{मह्यम्॥१६॥}

\end{visheshshloka}

\customparagraph{\textbf{अन्वयः}}

हे पृथिवि ! ताः प्रजाः नः संदुह्रताम् । समग्राः वाचः मधु मह्यं धेहि।

\customparagraph{\textbf{शब्दार्थः}}

हे पृथिवि ! ताः= पूर्वनिर्दिष्टाः द्विपदचतुष्पदपञ्चजनाख्याः, प्रजाः= प्रकर्षेण जायन्ते इति प्रजाः सृष्टयः, नः= अस्मभ्यम्, संदुह्रताम्= दुग्धाम्, सर्वा अपि प्रजा मह्यं नमन्तामिति भावः । तथा समग्राः= समग्रायाः (षष्ठ्यर्थे प्रथमा) वाचः= वाण्याः, मधु= सारभूतं मधुरं तत्त्वम्, मह्यम्= मदर्थे, धेहि= स्थापय ।

\customparagraph{\textbf{भावार्थः}}

हे पृथिवि ! पूर्वपद्ये निर्दिष्टा द्विपदचतुष्पदपञ्चजनाख्याः प्रजा मम वशीभूता भवन्तु । एवमेव सर्वेषां भाषाणां सारभूतं तत्त्व मह्यं स्थापय । मां सम्पूर्णवाङ्मयतत्त्वज्ञानसम्पन्नं विदधातु इत्यर्थः ।

\customparagraph{\textbf{नेपाली भावार्थ}}

हे पृथिवि ! अघिल्लो पद्यमा वर्णन गरिएका दुई पाउ भएका चार पाउ भएका र पाँच किसिमका मानवहरू मलाई अभीष्ट प्रदान गरून् अर्थात् ती सबै मेरा अधीनमा होऊन् । तिमी सम्पूर्ण वाणीको सारभूत तत्त्व ममा स्थापित गर अर्थात् मलाई विश्वका सम्पूर्ण वाङ्मयको तत्त्ववेत्ता बनाऊ ।

\begin{visheshshloka}[center]

\textbf{विश्वस्वं मातरमोषधीनां ध्रुवां भूमिं पृथिवीं धर्मणा धृताम्।}

\textbf{शिवां स्योनामनुचरेम विश्वहा॥१७॥}

\end{visheshshloka}

\customparagraph{\textbf{अन्वयः}}

विश्वहा विश्वस्वम् ‍ओषधीनां मातरं धर्मणा धृतां ध्रुवां भूमिं शिवां स्योनां पृथिवीम् अनुचरेम ।

\customparagraph{\textbf{शब्दार्थः}}

(वयम्) विश्वहा= सर्वेषु अहसु, सर्वदा इत्यर्थः । विश्वस्वम्= सर्वेषामपि प्राणिनां जीवनदातृत्वाद्धनभूताम्, ओषधीनाम्= गोधूमधान्यादिकानाम्, मातरम्= उत्पादयित्रीम्, धर्मणा= सत्यवदनाग्निहोत्रादिवेदप्रतिपादितेन क्रियाकलापेन, धृताम्= अवष्टब्धाम्, ध्रुवाम्= पर्वताद्यात्मना कठिनाम्, भूमिम्= भावयित्रीम् अन्नाद्युत्पत्तये उर्वरात्वेन प्रसिद्धाम्, शिवाम्= कल्याणस्वरूपाम्, स्योनाम्= सुखप्रदाम्, पृथिवीम्= जगतीतलम्, अनुचरेम= परिभ्रमाम, अप्रतिहतगतयो वयं सर्वदा सर्वत्र भवेमेत्याशीः ।

\customparagraph{\textbf{भावार्थः}}

अप्रतिहतगतियुता वयं मानवाः सर्वेषां प्राणिनां धनभूताम्, गोधूमधान्यादीनामुत्पादयित्रीम्, सत्यादिवेदप्रतिपादितधर्मक्रियादिभिर्धृताम्, पर्वताद्यात्मना कठिनाम्, अन्नाद्युत्पत्तये उर्वरात्वेन भूगुणसम्पन्नाम्, सुखप्रदाम्, पृथिवीं सर्वत्र परिभ्रमाम, निर्वाधरूपेण सर्वतोभावेन सर्वत्र गन्तुं शक्नुमेत्यर्थः ।

\customparagraph{\textbf{नेपाली भावार्थ}}

हामी सधैँ प्राणीहरूको धनस्वरूप अन्नकी जननी, सत्य आदि वेदप्रतिपादित क्रियात्मक धर्मद्वारा धारण गरिएकी, पर्वतादिका कारण कठिन, उब्जाउरूपका कारण भूगुणसम्पन्न, कल्याणकारिणी, प्राणीहरूलाई सुखदिने, पृथिवीमा निर्वाधरूपले निरन्तर भ्रमण गरौँ ।

\begin{visheshshloka}[center]

\textbf{महत्सधस्थं महती बभूविथ महान्वेग एजथुर्वेपथुष्टे। महाँस्त्वेन्द्रो}

\textbf{रक्षत्यप्रमादम्। सा नो भूमे प्र रोचय हिरण्यस्येव सन्दृशि मा नो}

\textbf{द्विक्षत कश्चन॥१८॥}

\end{visheshshloka}

\customparagraph{\textbf{अन्वयः}}

हे भूमे ! महत् सधस्थं (इति) महती बभूविथ । ते वेगः, एजथुः, वेपथुः महान् (अस्ति) त्वा महान् इन्द्रः अप्रमादं रक्षति । सा नः हिरण्यस्य इव सन्दृशि प्ररोचय । नः कश्चन मा द्विक्षत ।

\customparagraph{\textbf{शब्दार्थः}}

हे भूमे ! महत्= द्रव्यदेशकालफलादिना बृहत्, सधस्थम्= सह तिष्ठन्ति ऋत्विजो यत्रेति सधस्थम्, यज्ञकर्म तत् कर्तव्यमिति त्वम्, महती= महत्त्वशालिनी विस्तृता वा, बभूविथ= जातासि । ते= तव वेगः= गमनसामर्थ्यं महान् अस्ति । (यावद्धि सूर्यः एकस्मिन् वर्षेे विचरति, तावदियं भूमिः स्वल्पेन कालेनातिक्रामतीति महावेगत्वमेतस्याः) एवम्  ते= तव, एजथुः= हर्षमूलकं कम्पनम्, (तच्च प्रभूतसस्यसम्पत्त्या दिशां प्रसादेन च लक्ष्यते) वेपथुः= भयशोकमूलकं कम्पनम्, (एतच्च भूचलनाद्यौत्पातिकतया लक्ष्यते) एतदुभयमेव ते महान् अस्ति । (प्राचूर्येण  लाभहानिसम्पादयितृत्वात्) अत एव त्वा= त्वाम्, महान्= कर्मणा महत्त्वशाली इन्द्रः, परमैश्वर्यसम्पन्नः सूर्यो देवः, अप्रमादम्= सावधानतया, रक्षति= स्वकीयेनाकर्षणेन धारयति । सा= एवंविधा त्वम्, नः= अस्मान्, हिरण्यस्येव= सुवर्णस्येव, सन्दृशि= दर्शने, प्ररोचय= कान्तिसम्पन्नान् कुरु । (यथा हिरण्यं दर्शनादेव मनः प्रलोभयत्येवमस्मानपि सार्वदेशिकानां मध्ये शरीर-गति-कर्म-बल-विद्यादिभिर्दर्शनीयतमानादर्शभूतान् विधेहि) नोऽस्मान् कश्चन= कोऽपि प्राणी, मा द्विक्षत= न्यायपरायणतया द्वेषं मा कुर्यादिति ।

\customparagraph{\textbf{भावार्थः}}

हे भूमे ! त्वं महतां यज्ञानां कृते विस्तारवती असि । तव वेगो महानस्ति । हर्षमयं शोकमयञ्च कम्पनादिकं महद् विद्यते । अत एव महैश्वर्यसम्पन्नः सूर्यो महत्त्वशाली इन्द्रश्च सावधानतया त्वां रक्षतः । एवंविधा त्वं अस्मान् सुवर्णवत् दर्शनीयान् कुरु । कोऽपि प्राणी अस्माभिः सह द्वेषं मा कुर्यादिति ।

\customparagraph{\textbf{नेपाली भावार्थ}}

हे भूमि ! तिमी ठुल्ठुला यज्ञका लागि विशाल बनेकी हौ । तिम्रो गति महान् छ । अतिशय धान्यादिको फलभारले हुने हर्षपूर्ण कम्पन र भूकम्पादिले हुने शोकदुःखादिले पूर्ण कम्पन पनि महान् छन् । त्यसैले महत्त्वशाली इन्द्र र ऐश्वर्यशाली सूर्य सावधान भएर तिम्रो रक्षा गर्दछन् । यस्ती बहुगुणसम्पन्न तिमी हामीलाई सुवर्णवत् दर्शनीय बनाऊ । त्यस अवस्थामा कसैले न्यायप्रिय हाम्रो डाहा नगरून् ।

\begin{visheshshloka}[center]

\textbf{अग्निर्भूम्यामोषधीष्वग्निमापो बिभ्रत्यग्निरश्मसु। अग्निरन्तः}

\textbf{पुरुषेषु गोष्वश्वेष्वग्नयः ॥१९॥}

\end{visheshshloka}

\customparagraph{\textbf{अन्वयः}}

भूम्याम् ओषधीषु अग्निः (अस्ति) आपः अग्निः बिभ्रति । अश्मसु अग्निः (अस्ति) पुरुषेषु अन्तः अग्निः (अस्ति) गोषु अश्वेषु अग्नयः (सन्ति)

\customparagraph{\textbf{शब्दार्थः}}

(पूर्वे हि षष्ठे मन्त्रे ‘वैश्वानरं बिभ्रती भूमिरग्निमि’त्युक्तम्, तदेवात्र विशद्यते) भूम्याम्= अन्नोत्पत्तिहेतुभूतायां पृथिव्याम्, ओषधीषु= आमयनाशिनीषु पिप्पलीशतावर्यादिषु (ओषध्यः फलपाकान्ता, इत्युक्तेः फलपाकान्तासु धान्यगोधूमाद्योषधीषु च) अग्निः= वह्निरस्ति । (अन्नभुक्त्यनन्तरमौष्ण्यस्य सर्वजनवेद्यत्वात्) आपः= जलानि च वडवानलत्वेन प्रख्यातम् अग्निं बिभ्रति= धारयन्ति । एवम् अश्मसु= पर्वतेषु तदीयखण्डेषु वा अग्निर्वर्तते । (परस्परसङ्‍घर्षणेन अग्निर्दर्शनात्) पुरुषेषु= मनुष्येषु अन्तः= मध्ये अग्निः प्रविष्टोऽस्ति । (अत एव यदा अग्निः= पुरुषान्निःसृतो भवति, तदा स शैत्यमाप्तो मृत इति उच्यते) एवमेव गोषु= अश्वेषु च अग्नयः सन्ति । (गवाश्वयोः सर्वपशूपलक्षकत्वम्, ‘अग्नयः’ इति बहुवचनञ्च नामान्तराभिप्रायेण बोध्यम्)

\customparagraph{\textbf{भावार्थः}}

पृथिव्यां स्थितासु ओषधीषु लक्षणया अन्नादिषु च भक्षणानन्तरमुष्णत्वानुभवाद् अग्नयः सन्ति । जलेऽपि वडवानलादिनाम्ना प्रख्यातोऽग्निरस्ति । परस्परघर्षणजन्यप्रमाणात् पर्वतखण्डेषु पाषाणेषु चाग्निरस्ति । उष्णत्वहीनेन मृत इति ज्ञानात् पुरुषेष्वपि अग्निरस्ति । एवमेव गोऽश्वप्रभृतिषु पशुष्वपि अग्नयः सन्ति । पृथिव्यां सर्वत्र सर्वेषु पदार्थेष्वग्निरस्तीति भावः ।

\customparagraph{\textbf{नेपाली भावार्थ}}

पृथिवीमा रहेका औषधि र औषधीय गुण भएका अन्नादिमा पनि अग्नि रहेको हुन्छ किनकि औषधि वा अन्नादि भक्षण गरेपछि उष्णताको अनुभव हुन्छ । पानीमा पनि वडवानलका रूपमा अग्नि रहेको हुन्छ। ढुङ्गामा पनि अग्नि हुन्छ किनकि परस्परमा घर्षण गर्दा आगो निस्किएको देखिन्छ । मानिस आदि प्राणीमा पनि अग्नि हुन्छ किनकि शरीरमा उष्णता रहेन भने मृत ठानिन्छ । त्यस्तै प्रकारले गाईवस्तु आदि प्राणीमा पनि अग्नि रहेको हुन्छ ।

\begin{visheshshloka}[center]

\textbf{अग्निर्दिव आतपत्यग्नेर्देवस्योर्व१न्तरिक्षम् ।}

\textbf{अग्निं मर्तास इन्धते हव्यवाहं घृतप्रियम् ॥ २०॥}

\end{visheshshloka}

\customparagraph{\textbf{अन्वयः}}

अग्निः दिवः आपतति । अग्नेः देवस्य ऊरु अन्तरिक्षम् (अस्ति) । मर्तासः हव्यवाहं घृतप्रियम् अग्निम् इन्धते ।

\customparagraph{\textbf{शब्दार्थः}}

अग्निः= बह्निरेव, दिवः= आकाशात्, सूर्यरूपेण आतपति= सर्वत्र तापं वितरति । एतादृशस्य अग्नेर्देवस्य ऊरु= विस्तीर्णम्, अन्तरिक्षम्= मध्यस्थानमस्ति । मर्तासः= मरणधर्माणो मनुष्याः, हव्यवाहम्= हव्यं वहतीति तम्, देवोद्देशेन प्रदत्तहविषो वोढारम् । घृतप्रियम्= घृतं प्रियमस्यास्ति तम्, तक्रनिर्मन्थनोत्पन्नदुग्धसारभूतं घृतम्, तस्य प्रियम् अग्निदेवं स्वकल्याणाय इन्धते= यज्ञादिषु प्रदीपयन्ति ।

\customparagraph{\textbf{भावार्थः}}

अग्निरेव सूर्यरूपेण आकाशात् तपति । तत्र अग्निदेवस्य विस्तीर्णम् अन्तरिक्षलोकं विद्यते । मानवा अपि स्वकल्याणाय हवनद्रव्यवोढारं घृतप्रियमग्निं यज्ञादिषु प्रदीपयन्ति ।

\customparagraph{\textbf{नेपाली भावार्थ}}

अग्नि देव नै आकाशबाट सूर्यका रूपमा तातो प्रकाश छर्छन् । त्यहाँ अग्निको विशाल अन्तरिक्ष लोक छ । मानिसहरू आफ्नो कल्याणका लागि हव्यपदार्थ बोकेर देवतासम्म पुर्‍याउने र घिउ मन पराउने अग्निलाई यज्ञका लागि प्रज्वलित गर्दछन् ।

\begin{visheshshloka}[center]

\textbf{अग्निवासाः पृथिव्यसितज्ञूस्त्विषीमन्तं संशितं मा कृणोतु ॥२१॥}

\end{visheshshloka}

\customparagraph{\textbf{अन्वयः}}

अग्निवासाः असितज्ञूः पृथिवी मा त्विषीमन्तं संशितं कृणोतु ।

\customparagraph{\textbf{शब्दार्थः}}

अग्निवासाः= अग्निरिव वासः स्थितिर्यस्याः सा (यथा ह्यग्निरङ्गीकृते पदार्थे सर्वत्रैवात्मसात्कृत्यावतिष्ठते, तथैवेयं पृथिव्यस्माकं सर्वेष्ववयवेषु विद्यते पार्थिवशरीरत्वादिति) असितज्ञूः= असितं संसारपर्यावर्तेऽपि रागद्वेषादिक्षुद्रजनसुलभदोषैरबद्धं पुरुषं जानातीत्यसितज्ञूः, क्षुद्रदोषज्ञा पृथिवी देवता, मा= माम्, त्विषीमन्तम्= दीप्तिमन्तम्, संशितम्= तीक्ष्णीकृतम्, सर्वत्रापि सद्विषये प्रथमपरिगणनीयमित्यर्थः, कृणोतु= विदधातु ।

\customparagraph{\textbf{भावार्थः}}

अग्निसमाना सर्वत्र विद्यमाना, रागद्वेषादिदुर्गुणविहिनस्य जनस्य विषये पूर्णतया ज्ञानवती पृथिवी मां सर्वानतिरिच्य सद्‍गुणवैभवशालिनं विदधातु ।

\customparagraph{\textbf{नेपाली भावार्थ}}

अग्निदेवझैं सबैतिर व्यापक रूपले विद्यमान; राग, द्वेष आदि दुर्गुण भएका मानिसलाई सजिलै चिन्ने देवी पृथिवीले मलाई अरूको मूल्याङ्कनमा सबैभन्दा अगाडि गनिने तीक्ष्ण बुद्धिवाला, उत्कृष्ट व्यक्तित्वले सम्पन्न बनाऊन् ।

\begin{visheshshloka}[center]

\textbf{भूम्यां देवेभ्यो ददति यज्ञं हव्यमरङ्‍कृतम् ।}

\textbf{भूम्यां मनुष्या जीवन्ति स्वधयान्नेन मर्त्याः ।}

\textbf{सा नो भूमिः प्राणमायुर्दधातु जरदृष्टिं मा पृथिवी कृणोतु ॥२२॥}

\end{visheshshloka}

\customparagraph{\textbf{अन्वयः}}

मनुष्याः भूम्याम् अरङ्‍कृतं हव्यं यज्ञं च देवेभ्यः ददति । भूम्यां मर्त्याः स्वधया अन्नेन जीवन्ति । सा भूमिः नः प्राणम् आयुः दधातु । पृथिवी मा जरदृष्टिं कृणोतु ।

\customparagraph{\textbf{शब्दार्थः}}

मनुष्याः= मननेन कार्यकारिणः पुरुषाः, भूम्याम्= पृथिव्यामधिकरणे, अरङ्‍कृतम्= अलङ्‍कृतं सुसम्पादितम्, हव्यम्= देवप्रदेयभागम्, यज्ञम्= पूजास्तुतिहोमादिकर्म च देवेभ्यः= द्युलोकवासिभ्यः, ददति= प्रवितरन्ति । तथा भूम्यामेव मर्त्याः= मरणधर्माणो मनुष्याः, स्वधया= पितृप्रदेयभागेन अन्नेन= मनुष्यपशुपक्ष्यादिप्रदेयभागेन च कृत्वा पितृ-नर-पशु-पक्ष्यादिकान् सर्वानपि जीवन्ति= आजीवयन्ति । सैषा भूमिर्नोऽस्माकम्, प्राणम्= जीवनहेतुभूतम्, आयुः= पूर्णावस्थां च दधातु= धारयतु । पृथिवी= देवता मा= माम्, स्वकीयरसप्रदानेन तदाप्यायनेन च जरदष्टिम्= जरन्तमप्यश्नुवानम्, वृद्धावस्थायामपि कार्यसामर्थ्ययुतमित्यर्थः, कृणोतु= करोतु ।

\customparagraph{\textbf{भावार्थः}}

सद्विचारयुताः पुरुषाः पृथिव्यां स्थिताः सन्तः सुसंस्कृतं हव्यं देवेभ्यः समर्पयन्ति। तथा भूम्यामेव स्थिता मनुष्याः पितृप्रदेयभागेन मनुष्य-पशु-पक्ष्यादिप्रदेयभागेन च पितृ-नर-पशु-पक्ष्यादिकान् आजीवयन्ति। तादृशी भूमिरस्माकं प्राणान् आयुश्च दधातु । सा भूमिदेवता मां स्वकीयान्नरसप्रदानेन तदाप्यायनेन च वृद्धावस्थायामपि कार्यसामर्थ्ययुतं करोतु ।

\customparagraph{\textbf{नेपाली भावार्थ}}

राम्रो ज्ञान भएका व्यक्तिले पृथिवीमा परिशोधित हवनीय द्रव्य देवताका लागि दिन्छन् । त्यस्तै पृथिवीमै मानिसहरु पितृलाई दिने र मनुष्य-पशु-पक्ष्यादिलाई दिने भागले पितृ, मानिस, पशु, पक्षी आदि प्राणीहरूको जीवन धारण गराउँछन् । त्यस्ती भूमिले हामीलाई प्राण र आयु दिऊन् । साथै वृद्धावस्थासम्म विभिन्न काम गर्ने सामर्थ्य प्रदान गरून्, अर्थात् निरोगी बनाऊन् ।

\begin{visheshshloka}[center]

\textbf{यस्ते गन्धः पृथिवि संबभूव यं बिभ्रत्योषधयो यमापः ।}

\textbf{यं गन्धर्वा अप्सरसश्च भेजिरेे तेन मा सुरभिं कृणु मा नो}

\textbf{द्विक्षत कश्चन ॥२३॥}

\end{visheshshloka}

\customparagraph{\textbf{अन्वयः}}

हे पृथिवि ! ते यः गन्धः संबभूव । यम् ओषधयः बिभ्रति । यम् आपः बिभ्रति । यं गन्धर्वाः अप्सरसः च भेजिरे । तेन मा सुरभिं कृणु । कश्चन मा नो द्विक्षत ।

\customparagraph{\textbf{शब्दार्थः}}

हे पृथिवि ! ते‍= तव यो गन्धः= परिचायको विशेषगुणो गन्धः, संबभूव= वैशिष्ट्येन प्रादुर्भूतः । यं गन्धम् ओषधयः= गोधूमधान्यादयः पदार्थाः, बिभ्रति= धारयन्ति । यं गन्धम् आपः= जलानि, बिभ्रति= धारयन्ति । यं च गन्धर्वाः अप्सरसः= इति देवयोनिविशेषाश्च भेजिरे= सेवन्ते (स्वस्य महात्म्यख्यापनाय) तेन= गन्धेन मा= माम् सुरभिम्= सुगन्धिम्, कृणु= कुरु । सुगन्धवत्वादेव कश्चन= कोऽपि प्राणी, नोऽस्मान्, मा द्विक्षत= द्वेषं न कुर्यात् । कस्यापि वयं घृणास्पदा न भवेमेति भावः।

\customparagraph{\textbf{भावार्थः}}

हे पृथिवि ! त्वयि यो विशेषो गन्धो विद्यते । (तत्र गन्धवती पृथिवी इति), यं गन्धम् ओषधयो गोधूमधान्यादयो धारयन्ति । यं गन्धं जलानि धारयन्ति। यं गन्धर्वाप्सरसो देवयोनिविशेषाः सेवन्ते। तेन विशिष्टगुणयुतेन सुगन्धेन मामपि सुरभिं कुरु, येन केऽपि जना मया सह द्वेषं न कुर्वन्तु ।

\customparagraph{\textbf{नेपाली भावार्थ}}

हे पृथिवि ! तिमीमा विशेष किसिमको सुगन्ध रहेको छ । जसलाई पृथिवीमा रहेका ओषधीय गुण भएका अन्नादिले धारण गर्दछन् । जसलाई पानीले पनि धारण गरेको छ । जसलाई गन्धर्व, अप्सरा आदि देवयोनिविशेषले पनि सेवन गर्दछन् । त्यस विशेष किसिमको सुगन्धले मलाई पनि पूर्ण गर, जसकरण कसैले पनि ममाथि द्वेष वा घृणाको व्यवहार नगरून् ।

\begin{visheshshloka}[center]

\textbf{यस्ते गन्धः पुष्करमाविवेश यं संजभ्रुः सूर्याया विवाहे ।}

\textbf{अमर्त्याः पृथिवि गन्धमग्रे तेन मा सुरभिं कृणु मा नो}

\textbf{द्विक्षत कश्चन ॥२४॥}

\end{visheshshloka}

\customparagraph{\textbf{अन्वयः}}

हे पृथिवि ! यः ते गन्धः पुष्करम् आविवेश । यं गन्धम् अमर्त्याः अग्रे सूर्यायाः विवाहे संजभ्रुः । तेन मा सुरभिं कृणु । नः कश्चन मा द्विक्षत।

\customparagraph{\textbf{शब्दार्थः}}

हे पृथिवि ! यः पूर्ववर्णितः, ते= तव विशेषगुणो गन्धः, पुष्करम्= कमलपुष्पम्, आविवेश= प्रविष्टोऽभूत् । यं च गन्धम्= पुष्करगन्धम्, अमर्त्याः= देवाः, अग्रे= सर्वतः प्रथमम्, सूर्यायाः= सूर्यपुत्र्याः, विवाहे= पाणिग्रहणकर्मणोऽवसरे, संजभ्रुः= अहरन्, सञ्चित्य स्वस्मिन् स्थापयामासुः । (सूर्याया विवाहोऽत्र चतुर्दशे काण्डे श्रूयते) तेन= पुष्करगन्धेन, मा= माम्, सुरभिम्= सुगन्धसम्पन्नम्, कृणु= विधेहि । नोऽस्मान् कश्चन= कोऽपि व्यक्तिः, मा द्विक्षित= द्वेषं मा विदध्यात् ।

\customparagraph{\textbf{भावार्थः}}

पूर्वश्लोके वर्णितो विशेषगुणसम्पन्नः सुगन्धः कमलपुष्पं प्रविष्टो जातः । यं कमलपुष्पस्थं गन्धं देवाः सर्वतः प्रथमं सूर्यपुत्र्या विवाहे नीतवन्तः । हे पृथिवि ! तादृशेन महत्त्वशालिना पुष्करगन्धेन मां परिपूर्णं कुरु । तस्मात् मां कोऽपि द्वेषं मा विदध्यात् ।

\customparagraph{\textbf{नेपाली भावार्थ}}

पृथिवीमा रहेको पूर्ववर्णित गन्ध कमलको फूलमा प्रवेश गर्‍यो । त्यो सुगन्धसम्पन्न कमलको फूल सर्वप्रथम देवताले सूर्यकी पुत्रीको विवाहका अवसरमा टिपेर लगे । हे पृथिवी ! त्यस्तो सर्वोत्कृष्ट कमलको सुगन्धले मलाई पूर्ण बनाऊ । जसकारण कसैले पनि ममाथि द्वेष नराखून्।

\begin{visheshshloka}[center]

\textbf{यस्ते गन्धः पुरुषेषु स्त्रीषु पुंसु भगो रुचिः ।}

\textbf{यो अश्वेषु वीरेषु यो मृगेषूत हस्तिषु ।}

\textbf{कन्यायां वर्चो यद् भूमे तेनास्माँ अपि संसृज मा नो}

\textbf{द्विक्षत कश्चन॥२५॥}

\end{visheshshloka}

\customparagraph{\textbf{अन्वयः}}

हे भूमे ! यः ते गन्धः पुरुषेषु स्त्रीषु (अस्ति) ‘तेन’ पुंसु भगः रुचिः (भवति)  उत यः अश्वेषु वीरेषु (अस्ति) यः मृगेषु हस्तिषु (अस्ति) यत् कन्यायां वर्च (अस्ति) तेन अस्मान् अपि संसृज । नः कश्चन मा द्विक्षत ।

\customparagraph{\textbf{शब्दार्थः}}

हे भूमे !‍= पृथिवि ! यस्ते विशेषगुणो गन्धः, पुरुषेषु= पुंव्यक्तिषु स्त्रीषु= स्त्रीव्यक्तिषु चास्ति । तेन च गन्धेन हेतुना पुंसु= पुरुषव्यक्तिषु भगः= स्त्रीणां सेव्यत्वम्, तत्र रुचिः= अभिलाषश्च भवति । उत= अपि च यो गन्धः, अश्वेषु= हयेषु वीरेषु= शौर्यप्रधानपुरुषेषु अस्ति । तथा यो गन्धः, मृगेषु= विविधेषु पशुषु हस्तिषु= करिषु चास्ति । यत्= यो गन्धः कन्यायाम्= कमनीयानां कुमार्याम्, वर्चः= लावण्यम्, तद्रूपेणास्ति । तेन= गन्धेन लावण्येन च अस्मानमपि संसृज= संयोजय । तस्मात् नोऽस्मान्, कश्चन= कोऽपि जनः, मा== द्विक्षत= मा विद्विषेदिति ।

\customparagraph{\textbf{भावार्थः}}

हे पृथिवि ! यस्ते गन्धरूपो विशेषगुणः स्त्रीपुरुषेष्वस्ति, तेन पुरुषव्यक्तिषु स्त्रीणां भोग्यत्वमभिलाषश्च भवति । यो गन्धोऽश्वेषु वीरपुरुषेषु च विद्यते । एवमेव आरण्यग्राम्यसिंहादिपशुषु हस्तिषु चास्ति । तथा च यो गन्धः कुमारिषु स्त्रीषु लावण्यरूपेणावस्थितो भवति । तादृशेन सर्वसामर्थ्यसम्पन्नेन गन्धेन अस्मानपि संयोजय, येन केऽपि जना अस्मान् प्रति विद्वेषं मा कुर्वन्तु ।

\customparagraph{\textbf{नेपाली भावार्थ}}

हे पृथिवि ! जुन तिम्रो विशेष गुण पुरुष र स्त्रीमा रहेको छ । जुन सुगन्धले पुरुषहरूमा महिलाहरूको भोग गर्ने सामर्थ्य र उनीहरूप्रति आकर्षण उत्पन्न गराउँछ । जुन गन्ध घोडामा, वीर पुरुषमा पनि रहेको छ । जुन गन्ध मृगादि पशु र हात्तीमा पनि अवस्थित छ । जुन गन्ध कुमारी सुन्दरी युवतीमा लावण्य रूपले रहन्छ । त्यस्तो विशिष्ट सुगन्धले हामीलाई पनि पूर्ण गराऊ; जसकारण हामीप्रति कसैले द्वेष गर्न नसकून् ।

\begin{visheshshloka}[center]

\textbf{शिला भूमिरश्मा पांसुः सा भूमिः सन्धृता धृता ।}

\textbf{तस्यै हिरण्यवक्षसे पृथिव्या अकरं नमः ॥२६॥}

\end{visheshshloka}

\customparagraph{\textbf{अन्वयः}}

सन्धृता धृता सा भूमिः शिला भूमिः अश्मा पांसुः (इति अस्ति) तस्यै हिरण्यवक्षसे पृथिव्यै नमः अकरम् ।

\customparagraph{\textbf{शब्दार्थः}}

पृथिव्यां तावत् सन्धृता धृता च द्विविधा भूमिर्भवति । यस्यामन्नादिनानाविधं जीवनधारकं गुरुपदार्थजातं सम्यगुत्पद्यते सा सन्धृता भूमिः । यस्यां च जीवनोपयोगीनि दारुपुष्पादीनि सम्भवन्ति लोह-रजतादिखनिर्वोपलभ्यते सा धृता भूमिः । सा= पूर्ववर्णिता भूमिः= पृथ्वी, शिला भूमिः अश्मा पांसुरिति चतुर्धा अस्ति । तत्र शिला= शिलाप्रायो भूप्रदेशः पर्वतीयः । भूमिः =  भूमगुणवत्वात् प्रभूतसस्योत्पादनशालिनी । अश्मा= शर्करा तत्प्रायो देशः शर्करिलः । पांसुः= सिकता तत्प्रायः प्रदेशः सिकतिलः । एतादृश्यै प्रभूतभेदशालिन्यै, तस्यै= पूर्वं स्तुतायै, हिरण्यवक्षसे= सुवर्णाकरं वक्षःस्थानीयमित्येतादृश्यै, पृथिव्यै= देव्यै, नमः= नमस्कारः, अकरम्= करोमि ।

\customparagraph{\textbf{भावार्थः}}

अन्नादिजीवनसाधनानामुत्पादयित्री सन्धृता भूमिः, जीवनोपयोगीनि वस्तूनि समुत्पादयित्री धृता भूमिरिति भूमिर्द्विविधा विद्यते । एवमेव सा भूमिः प्रस्तरयुता, अन्नोत्पादनशालिनी, अश्मवती, सिकतायुता च चतुर्विधा विद्यते । एतादृश्यैः नानारूपवत्यै सुवर्णखनिरूपवक्षस्थलयुतायै पृथिव्यै देव्यै नमस्करोमि ।

\customparagraph{\textbf{नेपाली भावार्थ}}

यी पृथिवी अन्न आदि प्राणीको जीवन बचाउने साधनलाई उत्पादन गर्ने ‘सन्धृता’ र जीवनका लागि उपयोगी फलाम, सुन आदि खनिज आदि उत्पादन गर्ने ‘धृता’ भूमि गरी दुई प्रकारकी छन् । त्यस्तै माटाको प्रकारका आधारमा ढुङ्‍ग्यान भूमि, मलिलो भूमि, कङ्कणमय भूमि र बालुवामय भूभि गरी चार प्रकारकी हुन्छिन् । यस्ती विभिन्न प्रकारकी र सुनको खानीरूपी छाती भएकी पृथिवीलाई म नमस्कार गर्दछु ।

\begin{visheshshloka}[center]

\textbf{यस्यां वृक्षा वानस्पत्या ध्रुवास्तिष्ठन्ति विश्वहा ।}

\textbf{पृथिवीं विश्वधायसं धृतामच्छा वदामसि ॥२७॥}

\end{visheshshloka}

\customparagraph{\textbf{अन्वयः}}

यस्यां वनस्पत्याः वृक्षाः ध्रुवाः विश्वहा तिष्ठन्ति । (तां) विश्वधायसं पृथिवीं धृताम् अच्छ वदामसि ।

\customparagraph{\textbf{शब्दार्थः}}

यस्याम्= भूम्याम्, वानस्पत्याः= पुष्पप्रसवानन्तरं फलवन्तः, वृक्षाः= तरवः, ध्रुवाः= स्थायित्वेन विश्वहा= सर्वेषु अहःसुु तिष्ठन्ति । तां विश्वधायसम्= विश्वं बहुलं दधाति फलप्रदानद्वारा सात्त्विकं जीवनमिति विश्वधायास्तां पृथिवीं ‘धृता’ इति नाम्ना अच्छ वदामसि= वदामो वयमिति । (मन्त्रैरेभिराकराणां लक्षणद्योतनं भवति) ।

\customparagraph{\textbf{भावार्थः}}

पृथिवी उत्पादयित्री अस्ति । तत्र सर्वादा वनस्पतयोऽन्नफलानि उत्पादयन्ति । तस्मात् पृथिवी सर्वदा फलभारयुता भवति । अतः वयं तां फलभारयुतां पृथिवीं धारणकर्तृत्वेन धृता इति नाम्ना कथयामः ।

\customparagraph{\textbf{नेपाली भावार्थ}}

पृथिवी सबै प्रकारका बीजकी उत्पादयित्री हुन् । उनमा सदाबहार फलपूर्ण वनस्पतिहरू शोभायमान हुन्छन् । अतः त्यस्ती सर्वदा फलभारले युक्त हुने पृथिवीलाई हामी धृता नामले पुकार्छौँ । (यहाँ प्रयुक्त ‘धृता’ शब्दले खनिजपदार्थयुक्त पृथिवीलाई पनि बुझाउँछ) ।

\begin{visheshshloka}[center]

\textbf{उदीराणा उतासीनास्तिष्ठन्तः प्रक्रामन्तः ।}

\textbf{पद्भ्यां दक्षिणसव्याभ्यां मा व्यथिष्महि भूम्याम् ॥२८॥}

\end{visheshshloka}

\customparagraph{\textbf{अन्वयः}}

भूम्याम् उदीरणाः उत आसीनाः ‘उत’ तिष्ठन्तः ‘उत’ प्रक्रामन्तः दक्षिणसव्याभ्यां पद्भ्यां ‘पृथिवीम्’ मा व्यथिष्महि ।

\customparagraph{\textbf{शब्दार्थः}}

भूम्याम्= पृथिव्याम्, उदीरणाः=गच्छन्तः, उतापि च आसीना= उपविष्टाः उत तिष्ठन्तः = ऊर्ध्वावस्थायिनः, प्रक्रामन्तः= धावन्तो वयम्, दक्षिणसव्याभ्याम्= द्वाभ्याम्, पद्भ्याम्= चरणाभ्याम्, मा व्यथिष्महि= पृथिवीं मा व्यथयेम।

\customparagraph{\textbf{भावार्थः}}

अस्यां पृथिव्यां दक्षिणवामाभ्यां चरणाभ्यां गच्छन्तः, स्थिररूपेण वसन्तः, धावन्तो वा वयं कदापि पृथिवीं मा व्यथयेम । श्रमेण दुःखयुता न भवेम, सदा शारीरिकैर्व्याधिभिरसङ्‌क्रान्ता निरोगिणो भवेमेति भावः ।

\customparagraph{\textbf{नेपाली भावार्थ}}

यो पृथिवीमा विभिन्न कामका सिलसिलामा हिँड्‌दा, बस्दा, दौडिँदा कहिल्यै पनि आफ्ना दाहिने र देब्रे दुबै पाउको आघातले पृथिवीलाई व्यथित नबनाऊँ।

\begin{visheshshloka}[center]

\textbf{विमृग्वरीं पृथिवीमा वदामि क्षमां भूमिं ब्रह्मणा वावृधानाम् ।}

\textbf{ऊर्जं पुष्टं बिभ्रतीमन्नभागं घृतं त्वाभि निषीदेम भूमे ॥ २९॥}

\end{visheshshloka}

\customparagraph{\textbf{अन्वयः}}

विमृग्वरीं पृथिवीम् आ वदामि- हे भूमे ! त्वा अभिनिषीदेम, क्षमां भूमिं ब्रह्मणा वावृधानाम्, ऊर्जं पुष्टं घृतम् अन्नभागं बिभ्रतीम् ।

\customparagraph{\textbf{शब्दार्थः}}

विमृग्वरीम्= गवेषणेन प्राप्याम्, पृथिवीम् भूभागम् आ वदामि कथयामि किमिति चतुर्थे पादे उच्यते- हे भूमे ! पृथिवी ! त्वा= त्वाम् अभिनिषिदेम= परितो वयं त्वदीयफलभोक्तृत्वेनावतिष्ठेम राजानमिव । कीदृशी भूमिर्विमृग्वरीत्युच्यते-  	क्षमाम्= सर्वस्याप्युप्तस्य बीजस्य  सहनशीलां सर्वविधान्नफलमूलजननीमिति भावः । भूमिम्= भूमगुणवतीम् अतिशयफलदात्रीम्, ब्रह्मणा= वेदेन यज्ञेन च वावृधानाम्= यस्यां वेदाध्ययने यज्ञादिविध्यनुष्ठाने शास्त्रत आदेशो मनसश्च प्रवृत्तिः समुदियात्ताम् । तथा ऊर्जम्= बलप्राणप्रदान्नरसयुक्तम् पुष्टं= दृढम्, घृतम्= स्नेहसमन्वितम् च अन्नभागम्= खाद्यवस्तुनः सेव्यमंशम्, बिभ्रतीमिति विमृग्वरीविशेषणानि । क्षमादिगुणयुता हि पृथिवी सर्वैरपि विमृग्यत इति विमृग्वरीत्युच्यते ।

\customparagraph{\textbf{भावार्थः}}

गवेषणेन प्राप्याम् (विमृग्वरीम्) पृथ्वीं प्रति कथयामि- हे भूमे ! त्वदीयफलभोक्तृत्वेन वयं तव परितः अवतिष्ठेम । कीदृशीं विमृग्वरीमित्युच्यते- सर्वबीजानां सहनशीलाम् अर्थादुत्पादयित्रीम्, अतिशयफलदात्रीम्, वेदेन यज्ञेन चानुष्ठानद्वारा बलप्राणसारभूतम् अन्नरसं बिभ्रतीम्  इत्यादिविमृग्वरी विशेषणानि । (क्षमादि गुणयुता हि पृथिवी सर्वैरपि विमृग्यत इति विमृग्वरीत्युच्यते) ।

\customparagraph{\textbf{नेपाली भावार्थ}}

खोजी गरेर प्राप्त भएकी (विमृग्वरी) पृथिवीलाई म भन्छु कि हे भूमि तिमीबाट उत्पन्न भएका अन्नादिफलको भोक्ता भएर हामी तिम्रा चारैतिर अवस्थित हुन्छौँ । विमृग्वरी यो विशेषणको भाव स्पष्ट पार्दै भनिन्छ कि सबै बीजहरूलाई ग्रहण गरेर सहनशक्तिले युक्त भई धारण गर्ने वेद र यज्ञद्वारा अनुशिष्ट मानिसको कर्मका कारण बल र प्राणका लागि सारभूत अन्नरसलाई धारण गर्ने विमृग्वरी पृथिवी । (अर्थात् क्षमादि गुणले युक्त भएकी पृथ्वीलाई सबैले खोजी गर्दछन् । अतः पृथ्वीलाई विमृग्वरी भनिएको हो) ।

\begin{visheshshloka}[center]

\textbf{शुद्धा न आपस्तन्वे क्षरन्तु यो नः सेदुरप्रिये तं नि दध्मः ।}

\textbf{पवित्रेण पृथिवि मोत् पुनामि ॥ ३०॥}

\end{visheshshloka}

\customparagraph{\textbf{अन्वयः}}

हे पृथिवि ! नः तन्वे शुद्धाः आपः क्षरन्तु, यः नः ‘अपां’ सेदु तम् अप्रिये निदध्मः, मा पवित्रेण उत्पुनामि ।

\customparagraph{\textbf{शब्दार्थः}}

हे पृथिवि ! नः= अस्माकम्, तन्वे= वपुरर्थे शरीरसन्धारणाय, शुद्धाः= निर्दुष्टाः, आपः= जलानि क्षरन्तु= श्चोच्यन्ताम् । यश्च नः= अस्माकम्, सेदुः= (षिध गत्याम्) गमनार्हः अग्राह्योऽहितकरत्वादीदृशो जलभागस्तम आदिहेयगुणविवर्धकः, तं अप्रिये= शास्त्रैरप्रशस्ते प्रदेशे निदध्मः= स्थापयामः । नन्वस्मदर्थमेव किमिति शुद्धजलाभ्यर्थनेति चेद्- यतो ह्ययं मा= माम् पवित्रेण= ब्रह्मणा वेदेन तदध्ययनद्वारा तदीयानुष्ठानकर्मणा वा उत्पुनामि उत्कृष्टतया स्वशरीरं पूतं करोमि ।

\customparagraph{\textbf{भावार्थः}}

हे पृथिवि ! रोगरहितं शुद्धं जलम् अस्माकं शरीराय आकाशात् पततु, वर्षारूपेणेत्यर्थः । अस्माकं कृते अहितकरो जलभागः शास्त्रैर्गर्हिते प्रदेशे स्थापयामः । अहं वेदाध्ययनानुष्ठानादिभिः शरीरं पवित्रं करोमि ।

\customparagraph{\textbf{नेपाली भावार्थ}}

हे पृथिवी ! रोगरहित शुद्ध जल हाम्रो शरीरको पुष्टिका लागि वर्षाका रूपमा आकाशबाट झरोस् । हाम्रा लागि अहितकर पानी शास्त्रद्वारा निन्दित अर्थात् आवासका लागि अयोग्य स्थानमा राख्छौँ । हामी वेदका मन्त्रको पाठ र वेदले बताएका अनुष्ठानहरूको सम्पादनद्वारा आफूलाई पवित्र गराउँछौँ ।

\begin{visheshshloka}[center]

\textbf{यास्ते प्राचीः प्रदिशो या उदीचीर्यास्ते भूमे अधराद्याश्च पश्चात् ।}

\textbf{स्योनास्ता मह्यं चरते भवन्तु मा नि पप्तं भुवने शिश्रियाणः॥३१॥}

\end{visheshshloka}

\customparagraph{\textbf{अन्वयः}}

हे भूमे ! भुवने शिश्रियाणः मा निपप्तम् । ते याः प्राचीः प्रदिशः याः उदीचीः ‘प्रदिशः’ याः ‘च’ ते अधरात् ‘प्रदिशः’ याः च पश्चात् ‘प्रदिशः’ (सन्ति) ताः चरते मह्यं स्योनः भवन्तु ।

\customparagraph{\textbf{शब्दार्थः}}

हे भूमे ! भुवने= लोके त्वदीये पृष्ठभागे इति भावः । शिश्रियाणः= सेवमानो यत्र कुत्रापि संवसन्नहमित्यर्थः । मा निपप्तम्= रोगादिना सम्मानतो वा निपतितो मा भूयासम् । हे भूमे ! ते= तव याः प्राचीः प्रदिशः= पूर्वदिगाश्रयभूताः प्राणिनः, याश्च उदीचीः प्रदिशः‍= औदिच्याः प्राणिनः, याश्च ते अधरात् प्रदेशः= दाक्षिणत्याः प्राणिनः, याश्च पश्चात् प्रदिशः= पाश्चात्त्याः प्राणिनः सन्ति । ते प्राणिनः (सर्वासु दिक्षु वसन्तः सर्वेऽपि प्राणिनः) अथवा ताः सर्वा अपि प्रदिशः पूर्वदक्षिणपश्चिमोत्तराभिधानाः प्रसिद्धा दिशः, चरते= कार्यवशात् दर्शनेच्छया वा इतस्ततो भ्राम्यते, मह्यं प्रार्थिने स्तोत्रे स्योनाः= सुखमयाः भवन्तु= जायन्ताम् ।

\customparagraph{\textbf{भावार्थः}}

हे भूमि ! तव पृष्ठभागे समाश्रितोऽहं कदापि केनापि कारणेन सम्मानात् रोगादेर्वा पतनं नाप्नुयाम् । तव पूर्वदक्षिणपश्चिमोत्तराभिधानेषु दिक्षु निवसन्तो जनाः पूर्वाद्या दिशो वा कार्यवशाद्दर्शनेच्छया वा इतस्ततो चरते स्तुतिपाठकाय मह्यं सुखमया भवन्तु ।

\customparagraph{\textbf{नेपाली भावार्थ}}

हे पृथिवी ! तिम्रा पृष्ठभागमा आश्रय लिएर क्रियाशील भएको म कुनै पनि कारणले पतित न होऊँ । तिम्रा पूर्व, पश्चिम, उत्तर, दक्षिण जस्ता चारै दिशामा आश्रय लिएर बसेका प्राणीहरू (अथवा पूर्व आदि दिशाहरू नै) विशेष कामले वा भ्रमणका सिलसिलामा त्यहाँ पुग्दा मेरा लागि अनुकूल, सुखकारक बनून् ।

\begin{visheshshloka}[center]

\textbf{मा नः पश्चान्मा पुरस्तान्नुदिष्ठा मोत्तरादधरादुत ।}

\textbf{स्वस्ति भूमे नो भव मा विदन् परिपन्थिनो वरीयो}

\textbf{यावया वधम् ॥३२॥}

\end{visheshshloka}

\customparagraph{\textbf{अन्वयः}}

‘हे भूमे !’ नः पश्चात् मा नुदिष्ठाः, पुरस्तात् मा ‘नुदिष्ठाः’ उत्तरात् मा ‘नुदिष्ठाः’ उत अधरात् ‘मा नुदिष्ठाः’ हे भूमे ! नः स्वस्ति भव । परिपन्थिनः ‘नः’ मा विदन् । ‘नः’ वरीयः वधं यावय ।

\customparagraph{\textbf{शब्दार्थः}}

हे भूमे ! नः= अस्मान् स्तोतॄन्, पश्चात्= पश्चप्रदेशात्, मा नुदिष्ठाः= मा प्रेरयेः ‘अस्मान् मा अपसारय इत्यर्थः’ (पश्चादित्यादिदिग्वाचिनः शब्दाः, परमत्र सामर्थ्यात्तत्तदिगवस्थितदेशपरतयोन्नेयाः,  दिशां विभुत्वात्ततो नोदनानुपपत्तेः) ततश्च पाश्चात्त्यप्रदेशेभ्यो मा प्रेरयेः । (देशतः प्रेरणं हि ततोऽपसारणं तन्मा कार्षीः) पुरस्तात्= पूर्वदेशेभ्यो मा नुदिष्ठाः= मा प्रेरयेः । उत्तरात्= औत्तरेयप्रदेशेभ्यो मा नुदिष्ठाः । उत= तथा अधरात्= दक्षिणप्रदेशेभ्यो मा नुदिष्ठाः । (सर्वेष्वपि चतुर्दिगवस्थितप्रदेशेषु गतानामस्माकं कदापि हीनतयापसारणं मा भवतु) हे भूमे ! सर्वत्रैव नः= अस्माकं कृते, स्वस्ति= कल्याणदायिनी भव । तथा नोऽस्मान् परिपन्थिनः= मार्गावरोधिनो विद्विषः, मा विदन्= अस्माकं महत्त्वं प्रभावं वा मा प्राप्नुवन्तु। क्षुद्रहृदयैः परिपन्थिभिस्तदीयदेशनिवासात् ईर्ष्याभावेन प्रयुक्तम् वरीयः= अत्यन्तं महत्, वधम्= हानिसम्पादकं कृत्यम्, यावय= अस्मत्तः पृथक्‍कुरु ।

\customparagraph{\textbf{भावार्थः}}

हे पृथिवि ! अस्मान् पूर्वदक्षिणपश्चिमोत्तराभिधानेभ्यः स्थानेभ्यो मा अपसारय । अस्माकं कृते त्वं कल्याणकारिणी भव । मार्गावरोधिनः शत्रवोऽस्माकं महत्त्वं प्रभावं वा मा प्राप्नुवन्तु । तैः परिपन्थिभिः प्रयुक्तं हानिसम्पादकं कृत्यम् अस्मत्तो दूरीकुरु ।

\customparagraph{\textbf{नेपाली भावार्थ}}

हे पृथिवि ! तिम्रा पूर्व, पश्चिम, उत्तर, दक्षिण, जुन दिशामा रहे पनि हामीलाई तिमीले नहटाउनू । तिमी हाम्रा लागि कल्याणदायिनी बन । विरोधी शत्रुहरूले हाम्रा महत्त्व, प्रभाव आदि केही पनि प्राप्त नगरून् । अर्थात् प्रतिस्पर्धामा हामीलाई नभेटून् । तिनीहरूले गर्ने घातक प्रहारलाई तिमीले हामीसम्म आउन नदिनू ।

\begin{visheshshloka}[center]

\textbf{यावत्तेऽभि विपश्यामि भूमे सूर्येण मेदिना ।}

\textbf{तावन्मे चक्षुर्मा मेष्टोत्तरामुत्तरां समाम् ॥३३॥}

\end{visheshshloka}

\customparagraph{\textbf{अन्वयः}}

हे भूमे ! उत्तराम् उत्तरां समां मेदिना सूर्येण यावत् ते अभिविपश्यामि तावत् मे चक्षुः मा मेष्ट ।

\customparagraph{\textbf{शब्दार्थः}}

हे भूमे ! उत्तराम् उत्तरां समाम्= उत्तरोत्तरेष्वागामिषु सर्वेषु वर्षेषु, (अत्यन्तसंयोगे द्वितीया) मेदिना= स्निग्धेन मित्रेण, सूर्येण= भानुना, तत्कृपया यावद् अहम्, ते= तव प्रदेशान्, अभिपश्यामि= अवलोकयामि, तावत्कालपर्यन्तम्, मे= मम स्तोतुः, चक्षुः= चक्षुरिन्द्रियम्, मा मेष्ट= मापगतं भवतु। पार्थिवाद्रोगादिदोषान्मम चक्षुरिन्द्रियं मा नश्येदिति याच्ञा ।

\customparagraph{\textbf{भावार्थः}}

हे भुमे ! आगामिषु वर्षेषु स्नेहयुतस्य सूर्यस्य कृपया जीवनपर्यन्तं यावत् कालं त्वां पश्यामि तावत् कालं लोकगतैरामयादिदोषैर्मम दर्शनेन्द्रियं दर्शनशक्तिहीनं मा भवेत् ।

\customparagraph{\textbf{नेपाली भावार्थ}}

हे पृथिवि ! मेरो जीवनमा आगामी जति वर्ष हुन्छन्, ती वर्षमा स्नेही भगवान् सूर्यका कृपाले मेरा आँखाको दर्शनशक्ति पृथिवीका रोगादिदोषका कारण क्षय नहोस् । अर्थात् म बाँचुन्जेल स्वस्थ आँखाले तिमीलाई हेर्न पाऊँ ।

\begin{visheshshloka}[center]

\textbf{यच्छयानः पर्यावर्ते दक्षिणं सव्यमभि भूमे पार्श्वम् ।}

\textbf{उत्तानास्त्वा प्रतीचीं यत् पृष्ठीभिरधिशेमहे ।}

\textbf{मा हिंसीस्तत्र नो भूमे सर्वस्य प्रतिशीवरि ॥३४॥}

\end{visheshshloka}

\customparagraph{\textbf{अन्वयः}}

हे भूमे ! यत् शयानः दक्षिणं सव्यं पार्श्वम् अभि पर्यावर्ते । यत् उत्तानाः त्वा प्रतीचीं (कृत्वा) पृष्ठिभिः अधिशेमहे । तत्र हे सर्वस्य प्रतिशीवरि ! हे भूमे ! नः मा हिंसीः ।

\customparagraph{\textbf{शब्दार्थः}}

हे भूमे ! यदहं शयानः‍ स्वपन्, दक्षिणम्= दक्षिणदिगवस्थितं भागम्, सव्यम्= वामं, पार्श्वम्= कक्षयोरधोदेशभागञ्च, अभि= अभिलक्ष्य पर्यावर्ते= परिवर्तनं करोमि । यच्च उत्तानाः= द्युमुखा भूत्वा, त्वा= त्वां मातरं पृथिवीम्, प्रतीचीम्= कृत्वा अनभिमुखीभूय, पृष्ठीभिः= स्वीयपृष्ठभागैर्वयमधिशेमहे= स्वपामः तत्र= अस्मिन्नपराधकदम्बके, हे सर्वस्य प्रतिशीवरि ! = समस्तस्यापि कार्यजातस्यावहेलनकारिणां कृते प्रतिशायनशीले कार्यनाशकारिणीति भावः । (यदि हि पृथिवी स्वपृष्ठवासिभ्यः क्रुध्वान्नादि नोत्पादयेत्तर्हि सर्वमपि प्रणश्येदेवेति) हे भूमे ! = भूगुणवति मातः पृथ्वी ! नः= अस्मान्, मा हिंसीः= मा विनाशय ।

\customparagraph{\textbf{भावार्थः}}

(पूर्वं दशमे मन्त्रे पृथिव्या मातृत्वमुक्तम्। मातुश्च श्रेयस्त्वात्तत्सविधे वयस्कव्यक्तेः शयनम्, पर्यावर्तः, पराङ्‍मुखीभावश्च नोचित इति क्षमापयति-) हे भूमे ! यदि अहं स्वपन् सन् दक्षिणवामभागे पार्श्वपरिवर्तनं करोमि, उन्मुखश्च स्वपिमि, तदा ‘अयं मम पुत्रः स्थानदोषेण मातरं तिरस्करोति’ इति विचार्य हे क्रोधेन सर्वनाशनसमर्था जननी ! मां  मा विनाशय । तेषु अपराधेषु क्षमस्व इति भावः ।

\customparagraph{\textbf{नेपाली भावार्थ}}

(पहिले वर्णन गरिएको दसौँ श्लोकमा पृथिवीलाई आमा मानिएको छ, आमा मान्या भएकाले सुतेको अवस्थामा विपरीत दिशा भएको खण्डमा अपराध हुने विचार गरी यहाँ त्यस्ता सम्भावित अपराधका लागि पृथ्वीसँग क्षमायाचना गरिएको छ-) हे पृथिवी ! जब म सुत्छु, तब यता उता फर्किन्छु, जब उत्तानो परेर सुती तिमीलाई पछाडि पार्छु त्यस अवस्थामा स्थान दोषका कारण यो छोराले अपमान गर्‍यो भनेर मेरो विनाश नगर्नू । तपाईं रिसाउनुभयो भने संसारै नाश हुन सक्छ, त्यसैले त्यस्ता अनजानमा भएका गल्तीमा क्षमा दिनुहोस् ।

\begin{visheshshloka}[center]

\textbf{यत्ते भूमे विखनामि क्षिप्रं तदपि रोहतु ।}

\textbf{मा ते मर्म विमृग्वरि मा ते हृदयमर्पिपम्॥३५॥}

\end{visheshshloka}

\customparagraph{\textbf{अन्वयः}}

हे भूमे ! यत् ते विखनामि तत् अपि क्षिप्रं रोहतु । हे विमृग्वरि ! ते मर्म मा ‘अर्पिपम्’ (तथा) ते हृदयं मा अर्पिपम् ।

\customparagraph{\textbf{शब्दार्थः}}

हे भूमे ! यत् ते= तव क्षेत्रम् अहम्, विखनामि= विशेषेण आकराद्युपलब्धये खनामि, तदपि क्षिप्रं रोहतु= शीघ्रं पूर्णतामाप्नोतु । (अपि शब्दोऽत्र वास्तुनिर्माणस्कन्दमूलादिनिमित्तेन क्रियमाणानां खननछेदनादिलघुहिंसादीनां सूचकः) हे विमृग्वरि ! = मार्गणीयेषु प्रशस्ये पृथिवि ! ते= तव मर्म= जीवनप्रदेशमवयवबन्धनहेतुभूतं स्थानम्, मा अर्पिपम्= व्यथितं मा कार्षम् । तथा ते= तव हृदयम्= केन्द्रप्रदेशम् च मा अर्पिपम्= न हिनसानि ।

\customparagraph{\textbf{भावार्थः}}

हे भूमे! आवश्यकतापरिपूर्तये गृहनिर्माणकन्दमूललाभार्थं यदहं तव क्षेत्रं खनामि तत् शीघ्रमेव पूर्ववद् भवतु । अहं तव पृष्ठभागे मर्मस्थले पीडाजन्यं कार्यं न करोमि । तव हृदयरूपे केन्द्रप्रदेशे हिंसादिकं दुःखजन्यं किमपि कार्यं न करोमि ।

\customparagraph{\textbf{नेपाली भावार्थ}}

हे पृथिवि ! म तिम्रो वक्षस्थल अर्थात् पृष्ठभूमिमा खनिज पदार्थ निकाल्न, घर बनाउन, कन्दमूल आदि प्राप्त गर्न खन्ने काम गर्दछु । मैले खनेका ती सबै ठाउँ छिटै पुरिऊन् । म तिम्रो मर्मस्थलमा पीडा दिने कुनै पनि कार्य नगरूँ। तिम्रो हृदयस्थलमा दुःखजन्य हिंसाका कार्य पनि नगरूँ ।

\begin{visheshshloka}[center]

\textbf{ग्रीष्मस्ते भूमे वर्षाणि शरद् हेमन्तः शिशिरो वसन्तः ।}

\textbf{ऋतवस्ते विहिता हायनीरहोरात्रे पृथिवि नो दुहाताम् ॥३६॥}

\end{visheshshloka}

\customparagraph{\textbf{अन्वयः}}

हे भूमे ! ग्रीष्मः वर्षाणि शरत् हेमन्तः शिशिरः वसन्तः (इति) हायनीः ऋतवः विहिताः (सन्ति) । हे पृथिवि ! ते अहोरात्रे (च) नः ते दुहाताम् ।

\customparagraph{\textbf{शब्दार्थः}}

हे भूमे ! ग्रीष्मः=निदाघः, वर्षाणि= वर्षाः, शरत्, हेमन्तः, शिशिरः, वसन्तः, इति षट् हायनीः= हायन्यः (‘वा छन्दसि’ पूर्वसवर्णदीर्घः) हायनम्= संवत्सरः, तत्सम्बधिनः ऋतवः, समेषां कृते विहिताः= विशेषेण हितकृतः सन्ति । हे पृथिवि ! ते= ऋतवः, अहोरात्रे च नः= अस्माकम्, ते= तानि पूर्वमन्त्रोक्तानि खात-मर्म-हृदयस्थानानि दुहाताम् (अत्र अचो व्यत्यय) दुहताम्= प्रपूरयन्तु ।

\customparagraph{\textbf{भावार्थः}}

हे भूमे ! ग्रीष्मः (ज्येष्ठाषाढौ) वर्षाः (श्रावणभाद्रौ) शरद् (आश्विनकार्तिकौ) हेमन्तः (मार्गशीर्षपौषौ) शिशिरः (माघफाल्गुनौ) वसन्तः (चैत्रवैशाखौ) इति द्वादशमासयुता षट्‍सङ्‍ख्यकाः ऋतवः, जगतो हितकारकाः सन्ति । हे पृथिवि ! ते ऋतवः अहोरात्रे च गृहनिर्माणादीनि कर्मणि अस्माभिर्विहितानि खातानि प्रपूरयन्तु । मर्मणि हृदयस्थाने च जातां पीडां विनाशयन्तु।

\customparagraph{\textbf{नेपाली भावार्थ}}

हे भूमि  ! ग्रीष्म, वर्षा, शरद्, हेमन्त, शिशिर, वसन्त गरी वर्षभरका बाह्र महिनाभित्र रहेका ६ ऋतु जगत्‌को हितकारक छन् । ती ६ ऋतु र ऋतुसम्बद्ध रातदिनले हामीले घर आदि बनाउँदा वा अन्य काम गर्दा पारेका खाल्डाहरू, छिटै पूरा गरून्, मर्मस्थानमा र हृदयमा गरिएका आधातहरू नष्ट गराऊन् अर्थात् तिमीलाई पूर्ववत् बनाऊन् ।

\begin{visheshshloka}[center]

\textbf{यापसर्पं विजमाना यस्यामासन्नग्नयो ये अप्स्वन्तः ।}

\textbf{परा दस्यून् ददती देवपीयूनिन्द्रं वृणाना पृथिवी न वृत्रम् ।}

\textbf{शक्राय दध्रे वृषभाय वृष्णे ॥३७॥}

\end{visheshshloka}

\customparagraph{\textbf{अन्वयः}}

या अपसर्पं विजमाना (आस्ते) यस्याम् अग्नयः आसन् । ये ‘अग्नयः’ अप्सु अन्तः (सन्ति तेभ्यः) देवपीयून् दस्यून् परा ददती, इन्द्रं वृणाना वृत्रं न ‘वृणाना’ पृथिवी शक्राय वृषभाय वृष्णे (रूपं) दध्रे ।

\customparagraph{\textbf{शब्दार्थः}}

या= पृथिवी अपसर्पम्= अविहितमार्गगामिनं पापिनं प्रति विजमाना= कम्पमाना आस्ते । यस्याम्= भूम्याम्, अग्नयः= वितानाग्नयः (आहवनीय-गार्हपत्य-दक्षिणाग्न्याख्याः, बहुत्वोपदेशाद् वितानाग्नयः) आसन्= सन्ति । तथा यस्यां पृथिव्याम् अप्सु= जलेषु अन्तः= मध्ये ये अग्नयः सन्ति। (जले वडवानालरूपेण प्रसिद्धत्वेऽपि बहुषु स्थलेषु यागकर्मण्येव जलाधिकरणत्वेनाग्नाविव हूयमानत्वाज्जलेऽग्निवासः इष्ट एव । अग्निजन्यत्वादपां कारणगुणानामनुवर्तमानत्वात् ।) तेभ्योऽग्निभ्यो देवपीयून्= देवविद्वेषिणः, दस्यून्= उपहिंसयितॄणि भूतानि, परा ददति= दानप्रतिलोम्यमाचरन्ती, (आयातनं साहाय्यं वा अनाचरन्ती) इन्द्रम्= ऐश्वर्यसम्पन्नं मुख्यं देवतत्त्वं वृणाना= अङ्गीकुर्वती । वृत्रम्= पाप्मानम्, न वृणाना= परित्यजन्ती पृथिवी अस्ति । सैषा पृथिवी, शक्राय= शुद्धाय शुचये, वृषभाय= पुंस्त्वसम्पन्नाय (पुंस्त्ववैशिष्ट्यद्योतनाय वृषभशब्दः) कर्मण्याय व्यक्तये वृष्णे= सेचकाय उपकर्त्रे (स्वीयं कामधेनुरूपम्) दध्रे‍‍‍= दधार । (अस्माभिरपि पृथिवीकृपाभिलाषिभिः शुचिभिः कर्मठैः सङ्‌घीभूतैश्च भाव्यमिति)

\customparagraph{\textbf{भावार्थः}}

या पृथिवी कुमार्गगामिनं पापिनं प्रति दुष्कार्यभयेन कम्पवती भवति । यस्यां पृथिव्याम् आहवनीय-गार्हपत्य-दक्षिणाख्यास्त्रिविधा अग्नयः सन्ति । यस्या जलेऽपि अग्नयः सन्ति । तेभ्योऽग्निभ्यो देवद्वेषिणः दुष्टान् प्रति दानप्रतिलोम्यमाचरन्ती, मुख्यं देवनायकमङ्गीकुर्वती, पापिनं परित्यजन्ती पृथिवी अस्ति । इयं पृथिवी शुक्लगुणसम्पन्नाय जनाय स्वीयं कामधेनुरूपं प्रकटयति ।

\customparagraph{\textbf{नेपाली भावार्थ}}

यी पृथिवी कुमार्गगामीदेखि दुष्कार्यका भयले काम्छिन् । यी पृथिवीमा आहवनीय, गार्हपत्य र दाक्षिण नाम गरेका अग्नि छन् । पृथिवीको जलमा पनि अग्नि छ । तिनै अग्निका सहायताले यी पृथिवी जनविरोधी, देवविरोधी दुष्टजनलाई फलबाट वञ्चित गराउँछिन् । देवराज इन्द्रलाई अपनाउँछिन् । पापीलाई परित्याग गर्दछिन् । यस्ती पृथिवी सर्वगुणसम्पन्न व्यक्तिका लागि आफ्नो सर्वदात्री कामधेनु रूप प्रकट गर्दछिन् ।

\begin{visheshshloka}[center]

\textbf{यस्यां सदोहविर्धाने यस्यां यूपो निमीयते ।}

\textbf{ब्रह्माणो यस्यामर्चन्त्यृग्भिः साम्ना यजुर्विदः ।}

\textbf{युज्यन्ते यस्यामृत्विजः सोममिन्द्राय पातवे ॥३८॥}

\end{visheshshloka}

\customparagraph{\textbf{अन्वयः}}

यस्यां सदोहविर्धाने ‘निमीयते’ । यस्यां यूपो निमीयते । यस्यां यजुर्विदो ब्रह्माणः ऋग्भिः साम्ना अर्चन्ति । यस्यां इन्द्राय सोमं पातवे ऋत्विजो युज्यन्ते ।

\customparagraph{\textbf{शब्दार्थः}}

यस्याम्= पृथिव्याम्, सदोहविर्धाने= सदोमण्डपः, हविर्धानमण्डपश्चेति द्वौ यज्ञीयमण्डपौ, निमीयते= निर्मीयेते । यस्यां च भूम्यां यूपः= पशुनियोजनीयो यज्ञीयः स्तम्भो, निमीयते= निर्मीयते । यस्यां च यजुर्विदो= यज्ञकल्पाम्नायवेत्तारः कुशला अध्वर्यवो वा ब्रह्माणा= ब्राह्मणाः, ऋग्भिः= नियताक्षरपादावसनैर्मन्त्रैः साम्ना= ऋगाधारत्वेन गीयमानैः स्वरविशेषवद्भिर्मन्त्रैः, अर्चन्ति= पूजयन्ति, देवानुपासन्ते । यस्यां च इन्द्राय= देवमुख्याय, सोमम्= सोमवल्लीरसम् पातवे= पातुं पाययितुमित्यर्थः । ऋत्विजः= यागकारयितृत्वेन वृता विद्वांसो ब्राह्मणाः, युज्यन्ते= युक्ता भवन्ति, ऐकमत्येन यथोपदेशं यथाधिकारं च कार्यं कुर्वन्ति।

\customparagraph{\textbf{भावार्थः}}

(‘यस्यामग्नयः’ इति पूर्वं यज्ञसाधनभूतानां वितानाग्नीनां निर्देशः कृत इति प्रसङ्गादन्यान्यपि पार्थिवानि यज्ञसाधनानि विव्रियन्ते, यस्यामिति-) यस्यां पृथिव्यां सदोमण्डपहविर्धानमण्डपौ निर्मीयेते । यस्यां यज्ञीयो यूपो निर्मीयते । यस्यां यज्ञविधिवेत्तारो ब्राह्मणा ऋग्वेदसामवेदमन्त्रैर्देवान् पूजयन्ति । यस्यां देवमुख्याय सोमरसं पातुम् ऋत्विजः संयुक्ता भवन्ति ।

\customparagraph{\textbf{नेपाली भावार्थ}}

पृथिवीमा सदोमण्डप र हविमण्डप गरी दुई किसिमका मण्डप बनाइन्छन् । जहाँ यज्ञका लागि यूपःस्तम्भ पनि निर्माण गरिन्छ । जहाँ यज्ञविधान जान्ने विद्वान् ब्राह्मणहरूबाट ऋग्वेदका र सामवेदका ऋचाहरूले देवताहरूको पूजा गरिन्छ । पृथिवीमा देवमुख्य इन्द्रलाई अमृत पिलाउन विद्वान् ब्राह्मणहरू एकीकृत भएर कार्य गर्दछन् ।

\begin{visheshshloka}[center]

\textbf{यस्यां पूर्वे भूतकृत ऋषयो गा उदानृचुः ।}

\textbf{सप्त सत्रेण वेधसो यज्ञेन तपसा सह ॥३९॥}

\end{visheshshloka}

\customparagraph{\textbf{अन्वयः}}

यस्यां पूर्वे भूतकृतः सप्त ऋषयः यज्ञेन तपसा सह वेधसः सत्रेण गा उदानृचुः ।

\customparagraph{\textbf{शब्दार्थः}}

यस्याम्= पृथिव्यामस्याम्, पूर्वे= प्राचीनाः, भूतकृतः= ऐेतिहासिकाः, सप्त= सप्तसङ्‍ख्यकाः, ऋषयः= गोत्रप्रवर्तका जमदग्नि-भरद्वाज-गौतम-अत्रि-वशिष्ठ-कश्यप-

विश्वामित्राख्या मुनयः, यज्ञेन= श्रुत्युक्तेन स्वोचितेनेज्याकर्मणा, तपसा= द्वन्द्वसहनेन च सह= युक्ताः, वेधसः सत्रेण = गोत्रप्रवर्तकस्वरूपप्रजापतित्वप्रापकेन यज्ञेन कृत्वा, गाः= ‘गोत्रप्रवर्तका इमः’ इति गिरः, उदानृचुः= उदानर्चुः, समपूजयन् ।

\customparagraph{\textbf{भावार्थः}}

यस्यां पृथिव्यां प्राचीनाः सप्तर्षयो यज्ञेन तपश्चर्यया च गोत्रप्रवर्तकत्वरूपप्रजापतित्वप्रापकेन सत्रेण ‘गोत्रप्रवर्तका इमः’ इति शब्देन सुपूजिताः ।

\customparagraph{\textbf{नेपाली भावार्थ}}

जुन पृथिवीमा प्राचीन जमदग्नि, भरद्वाज, अत्रि, गौतम, वशिष्ठ, कश्यप, विश्वामित्र जस्ता सप्त ऋषिहरू यज्ञ र तपस्याका साथमा गोत्रपवर्तक प्रजापतित्व प्राप्तिका लागि गरिने यज्ञद्वारा ‘यी ऋषि गोत्रप्रवर्तक हुन्’ भन्ने शब्दले प्रख्यात वा पूजित भए ।

\begin{visheshshloka}[center]

\textbf{सा नो भूमिरादिशतु यद्धनं कामयामहे ।}

\textbf{भगो अनुप्रयुङ्‌क्तामिन्द्र एतु पुरोगवः ॥४०॥}

\end{visheshshloka}

\customparagraph{\textbf{अन्वयः}}

सा भूमिः नः यत् धनं कामयामहे (तत्) आदिशतु । (तत्र) भगः अनुप्रयुङ्‌क्ताम्  । पुरोगवः इन्द्रः एतु ।

\customparagraph{\textbf{शब्दार्थः}}

सा= पूर्वमाख्याता भूमिः= पृथिवी नः= अस्मभ्यम् यत् धनम्= वसु, वयं कामयामहे= अभिलषामः, तत् आदिशतु= प्रयच्छतु । तदर्थं भगो= देवः अनुप्रयुङ्‌क्ताम्= भूम्या धनप्रदानेऽस्मदानुकूल्येनोद्योगं कुर्यात् । पुरोगवः= इन्द्रश्च एतु= आगच्छतु, भूम्या धनप्रधानानुमतये । (इन्द्रो हि स्तोतृभ्यो मुख्यतया गावो ददातीति पुरोगव उच्यते, इन्द्रो यथा स्वस्तोतृभ्यः उदारमनसा गावो ददाति, तथा स्तोतृभ्योऽस्मभ्यं भूमिर्दद्यादित्यर्थे इन्द्रोऽनुमन्यताम्, तदर्थं चात्र समेतु, बलधनयोर्दातृत्वमत्वात्)

\customparagraph{\textbf{भावार्थः}}

तादृशी पूर्ववर्णिता विविधगुणयुता पृथिवी धनमिभलषन्तमस्मान् धनं ददातु । अस्य कृते भगो देवता अनुकूलानि कार्याणि करोतु । पुरोगवो इन्द्रोऽपि पृथिवीं धनप्रदानाय प्रेरयितुमत्र समागच्छतु ।

\customparagraph{\textbf{नेपाली भावार्थ}}

पहिलेका श्लोकमा वर्णन गरिएकी पृथिवीले हामी धनको इच्छा गर्नेलाई धन दिऊन् । भग नामका देवताले पनि यसका लागि अनकूल कार्य गरून् । पृृथिवीलाई प्रेरित गर्न स्वयं गाई दिनेवाला इन्द्र आऊन् ।

\begin{visheshshloka}[center]

\textbf{यस्यां गायन्ति नृत्यन्ति भूम्यां मर्त्या व्यैलवाः ।}

\textbf{युध्यन्ते यस्यामाक्रन्दो यस्यां वदति दुन्दुभिः ।}

\textbf{सा नो भूमिः प्रणुदतां सपत्नानसपत्नं मा पृथिवी कृणोतु ॥४१॥}

\end{visheshshloka}

\customparagraph{\textbf{अन्वयः}}

व्यैलवाः मर्त्याः यस्यां भूम्यां गायन्ति नृत्यन्ति । यस्यां युध्यन्ते । यस्यां आक्रन्दः दुन्दुभिः वदति । सा भूमिः नः सपत्नान् प्रणुदताम् । पृथिवी मा असपत्नं कृणोतु ।

\customparagraph{\textbf{शब्दार्थः}}

इला= पृथिवी, तत्सम्बन्धिनः प्रदेशा ऐलाः, तेषु वसन्तीति ऐलवाः, विविधाश्च ऐलवाः, व्यैलवाः= विविधदेशवाशिनः, मर्त्याः= मनुष्याः, यस्यां भूम्याम्= भूप्रदेशे परस्परमेकीभूय हर्षोद्रेकेण गायन्ति= गायनं कुर्वन्ति, नृत्यन्ति= नृत्यं प्रदर्शयन्ति । यस्यां भूम्यामेव ते मनुष्याः परस्परं विमताः (विरुद्धमतयुताः) सन्तो द्वेषेण विजिगीषया, युध्यन्ते= सङ्‍ग्रामं कुर्वते । यस्यां पृथिव्यां विजिगीषयैवाक्रमणसंसूचनाय आक्रन्दः= आ समन्तात् क्रदति शब्दं करोत्यभिमुखं युद्धायेत्याक्रन्दः, रणसूचको दुन्दुभिः= वाद्यविशेषः, वदति= उद्घोष्यते । सा= एवंविधा सन्धि-विग्रह-यानाख्य- सन्धित्रयाधारभूता भूमिः, नः= अस्माकम्, सपत्नान्= शत्रून्, प्रणुदताम्= प्रकर्षेणापसारयतु । अपसार्य च समस्तानपि विद्विषः पृथिवी सर्वकालं मा= माम्, असपत्नम्= अविद्यमानशत्रुम्, अजातशत्रुमिति भावः, कृणोतु= करोतु ।

\customparagraph{\textbf{भावार्थः}}

यस्यां पृथिव्यां विविधप्रदेशवासिनो जना एकीभूय हर्षोत्कर्षेण प्रमोददायकानि गीतानि गायन्ति, दर्शनीयं नृत्यं प्रदर्शयन्ति, यस्याञ्च पृथिव्यां त एव विविधदेशवासिनो मर्त्याः कदाचित् स्वार्थवशेन परस्परं युद्धं कुर्वन्ति, यस्याञ्च युद्धारम्भसूचनाय दुन्दुभिनामको वाद्यविशेषो वाद्यते तादृशी भूमिरस्माकं शत्रून् दुरीकरोतु । अस्माकं सकलानपि शत्रून् अपसार्य मां शत्रुरहितम् ‘अजातशत्रुम्’ विदधातु ।

\customparagraph{\textbf{नेपाली भावार्थ}}

जुन पृथिवीमा विभिन्न स्थानमा बसोबास गर्ने मानिसहरू एकआपसमा मिलेर आनन्ददायक गीत गाउँछन्, आकर्षक नृत्य प्रदर्शन गर्दछन् । कहिलेकाहीँ स्वार्थवश एकआपसमा युद्ध गर्दछन् । युद्धको जानकारीका लागि दुन्दुभि बजाउँछन् । यस्तो सन्धि-विग्रह-यानजस्ता सन्धित्रयको आधारभूत पृथिवीले हाम्रा शत्रुहरू नाश गरून् । शत्रुहरूलाई नष्ट गराएर मलाई शत्रुविहीन अर्थात् अजातशत्रु बनाऊन् ।

\begin{visheshshloka}[center]

\textbf{यस्यामन्नं व्रीहियवौ यस्या इमाः पञ्च कृष्टय ।}

\textbf{भूम्यै पर्जन्यपत्न्यै नमोऽस्तु वर्षमेदसे ॥४२॥}

\end{visheshshloka}

\customparagraph{\textbf{अन्वयः}}

यस्यां व्रीहियवौ अन्नम् (अस्ति) । यस्याः इमाः पञ्चः कृष्टयः (सम्भवन्ति)। पर्जन्यपत्न्यै वर्षमेदसे भूम्यै नमः अस्तु ।

\customparagraph{\textbf{शब्दार्थः}}

यस्याम्= भूम्याम्, व्रीहियवौ धान्ये, विशेषतः अन्नम्= खाद्यम् अस्ति । यस्याः (पञ्चम्यर्थे षष्ठी) यस्मात्= अन्नात्, इमाः= प्रत्यक्षं दृश्यमानाः, पञ्च= पञ्चसङ्‍ख्यकाः कृष्टयः= मनुष्याः (ब्राह्मण-क्षत्रिय-वैश्य-शूद्र-निषादाख्याः, अथवा नृतिर्यक्‍पशुदेवासुरात्मकाः) पञ्चविधाः प्राणिनः संभवन्ति । तस्यै पर्जन्यपत्न्यै= पर्जन्यद्वारा पालितायै, तथा वर्षमेधसे= वर्षाभिवृद्धिमीयुष्यै, वर्षायाः कारणेन पुष्टियुतायै भूम्यै= पृथिव्यै नमोऽस्तु= नमस्काराणि संयुज्यन्ताम् ।

\customparagraph{\textbf{भावार्थः}}

यस्यां पृथव्यां व्रीहियवादीनि अन्नानि सञ्जायन्ते । यस्मादन्नात् पञ्च प्रकाराः (ब्राह्मण-क्षत्रिय-वैश्य-शूद्र-निषादाख्याः, अथवा नृतिर्यक्‍पशुदेवासुरात्मकाः) प्रणिनः सञ्जायन्ते जीवन्ति वा । एतादृश्यै पर्जन्यपालितायै, वर्षापरिपुष्टायै भूम्यै नमोऽस्तु ।

\customparagraph{\textbf{नेपाली भावार्थ}}

जुन पृथिवीमा धान, जौ आदि अन्नहरू उत्पादन हुन्छन् । जसमा पाँच प्रकारका (ब्राह्मण, क्षत्रिय, वैश्य, शूद्र, निषाद वा मानिस, पक्षी, पशु, देवता, असुर) प्राणी उत्पन्न भई जीवित रहन्छन् । त्यस्ती मेघद्वारा पालित, वर्षाका कारणले (सबै प्राणी र वनस्पतिका रूपमा) परिपुष्ट पृथ्वीलाई नमस्कार छ ।

\begin{visheshshloka}[center]

\textbf{यस्याः पुरो देवकृताः क्षेत्रे यस्या विकुर्वते ।}

\textbf{प्रजापतिः पृथिवीं विश्वगर्भामाशामाशां रण्यां नः कृणोतु ॥४३॥}

\end{visheshshloka}

\customparagraph{\textbf{अन्वयः}}

यस्याः देवकृताः पुरः (सन्ति) यस्याः क्षेत्रे (पापाः) विकुर्वते । (एवम्) विश्वगर्भां पृथिवीं प्रजापतिः नः आशामाशां रण्यां कृणोतु ।

\customparagraph{\textbf{शब्दार्थः}}

यस्याः= पृथिव्याः देवकृताः= देवैः स्वविशेषकार्येण कृताः= निर्मिताः, पुरः= काशीमथुरायोध्याकाञ्च्यादीनि नगराणि सन्ति । यस्याश्च क्षेत्रे= पुण्यप्रदेशे पुण्यक्षेत्रत्वेनाभिमते भूभागे (पापानि) विकुर्वते= नष्टानि भवन्ति । तदेवं विश्वगर्भाम्= सर्वविधाभीप्सितवस्त्वायतनभूतां पृथिवीम्= भूमिम्, प्रजापतिः देवः, नः= अस्माकं कृते,  आशामाशाम्= प्रतिदिशम्, रण्याम्= रमणीयाम्, कृणोतु= करोतु ।

\customparagraph{\textbf{भावार्थः}}

यस्यां पृथिव्यां दैवैः सर्वकल्याणाय निर्मितानि काशीमथुरायोध्याकाञ्च्यादीनि पुराणि सन्ति । यस्याः काश्यादीनां पुण्यक्षेत्राणां दर्शनपूजनादिकार्येण प्राणिनां पापानि नश्यन्ति । एतादृश्याः सर्वाभिलषितवस्त्वाधारभूतायाः पृथिव्याः  पूर्वपश्चिमादिदिशः प्रजापतिदेवता रमणीया विदधातु ।

\customparagraph{\textbf{नेपाली भावार्थ}}

जुन पृथिवीमा देवताहरूले सबैको कल्याणका लागि काशी, अयोध्या आदि पवित्र धार्मिक नगरहरू स्थापना गरेका छन् । जुन नगरको दर्शन, पूजन आदिले सबै पापहरू नष्ट हुन्छन्। त्यस्ती सबै अभिलषित वस्तुका आधारका रूपमा रहेकी पृथिवीका पूर्व, पश्चिम आदि सबै दिशा प्रजापति देवताले रमणीय बनाऊन् ।

\begin{visheshshloka}[center]

\textbf{निधिं बिभ्रती बहुधा गुहा वसु मणिं हिरण्यं पृथिवी ददातु मे ।}

\textbf{वसूनि नो वसुदा रासमाना देवी दधातु सुमनस्यमाना ॥४४॥}

\end{visheshshloka}

\customparagraph{\textbf{अन्वयः}}

गुहा बहुधा वसु निधिं बिभ्रती पृथिवी मे मणिं हिरण्यं ददातु । सुमनस्यमाना वसुदा रासमाना देवी नः वसूनि दधातु ।

\customparagraph{\textbf{शब्दार्थः}}

गुहा= गुहायाम् (सप्तम्या डादेशः) स्वस्यान्तः प्रदेशे बहुधा= बहुविधम्, वसु= वसूनां धनानाम्, निधिम्= आकारम् (पद्मादिनवनिधीन् वा) बिभ्रती= धारयन्ती पृथिवी मे= मह्यम्, मणिम्= वज्रवैडूर्यमरकताद्यश्मजातम्, मौक्तिकादि च (कुप्यम्), हिरण्यम्= सुवर्णरजतात्मकम्, अकुप्यं धनम्, ददातु= प्रवितरतु । मणिहिरण्यादिकमेव किमु अपि तु सुमनस्यमाना= प्रसन्नमानसा वसुदा= धनदायिनी, रासमाना= दानोन्मुखी देवी पृथिवी, नः= अस्माकं कृते वसूनि= बहूत्वोपदेशात् सर्वाणि अकुप्यातिरिक्तानि कुप्यात्मकान्यपि धनानि, दधातु= धारयतु, मामर्पयतु इति ।

\customparagraph{\textbf{भावार्थः}}

स्वान्तरप्रदेशे विविधानां धनानामाकरं धारयन्ती पृथिवी मह्यं मणिं हिरण्यं च ददातु । एवमेव प्रसन्नमानसा धनदायिनी दानशीला देवी पृथिवी अस्माकं सर्वाणि कुप्याकुप्यमयानि धनानि धारयतु प्रददातु वा ।

\customparagraph{\textbf{नेपाली भावार्थ}}

आफ्नो भित्रीभागमा (जमिनभित्र) मणि, सुवर्ण आदिको खानीलाई धारण गर्ने पृथ्वी देवीले ती सबै धन मलाई दिऊन् । साथै प्रसन्न चित्त भएकी धनदायिनी, दानशील पृथिवीले हामीलाई अकुप्य (खनिजजन्य सुवर्णादि बहुमूल्य धातु) र कुप्य (सुवर्ण र रौप्यभिन्न धातुका साथै) सबै प्रकारका धन प्रदान गरून् ।

\begin{visheshshloka}[center]

\textbf{जनं बिभ्रती बहुधा विवाचसं नानाधर्माणं पृथिवी यथौकसम् ।}

\textbf{सहस्रं धारा द्रविणस्य मे दुहां ध्रुवेव धेनुरनपस्फुरन्ती ॥४५॥}

\end{visheshshloka}

\customparagraph{\textbf{अन्वयः}}

विवाचसं नानाधर्माणं बहुधा जनं यथौकसं विभ्रती पृथिवी अनपस्फुरन्ती ध्रुवा धेनु इव द्रविणस्य सहस्रं धारा मे दुहाम् ।

\customparagraph{\textbf{शब्दार्थः}}

विवाचसम्= विविधभाषाभाषिणम्, नानाधर्माणम्= नैकविधमर्यादावन्तम्, बहुधा= बहुप्रकारकम्, जनम्= मनुष्यम्, यथौकसम्= यथास्थानम्, बिभ्रती= धारयन्ती पृथिवी= भूमिः, मे= मह्यम्, द्रविणस्य= धनस्य सहस्रं धारा= अनन्ताः शाखाः, दुहाम्= दोग्धु व्यवस्थापयत्वित्यर्थः। तत्र दृष्टान्तः - अनपस्फुरन्ती= अङ्‍गानि किञ्चिन्मात्रमपि न परिचालयन्ती, ध्रुवा= सुस्थिरा धेनुः= सवत्सा गौरिव। (यथा हि धीरा अकम्पमाना च गौः साधु दुग्धे तथैवेयं वसुधा सुस्थिरा सती स्वीयमन्तःस्थितं द्रविणं सर्वभाषाभाषिभ्यः सर्वधर्मानुयायिभ्यश्च साधु दुग्धाम्)

\customparagraph{\textbf{भावार्थः}}

विविधभाषायुतान्, विविधधर्मयुतान् च जनान् यथास्थानं धारयन्ती पृथिवी यथा शान्तस्वभावयुता सवत्सा गौः सारल्येन दुग्धस्य सहस्रधारां ददाति, तथैव शान्ता स्थिरा च सती सारल्येन मह्यं धनस्य सहस्रधारां व्यवस्थापयतु, निरन्तरमावश्यकं धनं ददातु इत्यर्थः।

\customparagraph{\textbf{नेपाली भावार्थ}}

अनेक भाषाभाषी र अनेक धर्मावलम्बी मानिसहरूलाई यथास्थानमा धारण गर्ने पृथ्वीले जसरी शान्तस्वभावको बाछोसहितको गाईले दुध दिन्छ त्यसरी नै शान्तस्वभावकी बनेर आवश्यकताअनुसार निरन्तर धन दिऊन् ।

\begin{visheshshloka}[center]

\textbf{यस्ते सर्पो वृश्चिकस्तृष्टदंश्मा हेमन्तजब्धो भृमलो गुहा शये ।}

\textbf{क्रिमिर्जिन्वत्पृथिवि यद्यदेजति प्रावृषि तन्नः सर्पन्मोप सृपद्यच्छिवं}

\textbf{तेन नो मृड ॥४६॥}

\end{visheshshloka}

\customparagraph{\textbf{अन्वयः}}

हे पृथिवि ! यः ते गुहा शये (सः) सर्पः, वृश्चिकः, तृष्टदंश्मा, हेमन्तजब्धः, भृमिलः, क्रिमिः, जिन्वत्, यत् यत् प्रावृषि एजति, तत् सर्पत् नः मा उप सृपत् । यत् शिवं तेन नो मृड ।

\customparagraph{\textbf{शब्दार्थः}}

हे पृथिवि ! यः ते= तव, गुहा= गुहयाम् (बिलरूपावकासे) शये= शेते, सः सर्पः= चक्षुश्रवाः, वृश्चिकः= प्रसिद्धो द्रुणः, तृष्टदंश्मा= यस्य दंशेन मनुष्यस्तृषार्तो भवत्येवं विधो जीवः, हेमन्तजब्धः= यस्य जग्ध्या भक्षणेन मनुष्योऽत्यन्तं हिमाक्रान्तो भवतीत्येवंविधो (मशकादि) जन्तुः, भृमलः= यस्य दंशनेन मनुष्यो घूर्णायते एवंविधो जीवः । यः किल क्रिमिः= क्रिमिभिः कीटैः (तृतीयाया सुः) । जिन्वत्= जिन्वति प्रीतो भवति । (कीटभक्षी विषधरो जीव इत्यर्थः), किं बहुना यद्यज्जीवजातं प्रावृषि= वर्षाकाले, एजति= कम्पते, प्रादुर्भवति परिदृश्यते वा तत् सकलमपि कृच्छ्रप्रदं (कष्टदायकम्) जीवजातम्, सर्पत्= इतस्ततः पृथिव्यां विचरत्, नः= अस्मान्, मा उपसृपत्= प्राप्तं न भवतु । (तेषां वयं दंशपात्रा मा भूमेति भावः) । अथ यत् शिवम्= सुखकारकं जीवजातं शुकसारिकादि, तेन नः= अस्मान् मृड= मृडय, सुखय ।

\customparagraph{\textbf{भावार्थः}}

हे पृथिवि ! सर्पाः, वृश्चिकाः, दंशेन तृष्णोत्पादकाः, भक्षणेन शीतोत्पादकाः, दंशनेन धूर्णयन्तः, कीटभक्षिणः, वर्षाकाले समुत्पद्यमानाश्चेति विविधा जीवास्तव गुहायां निवसन्ति । तेषां हिंस्रकजीवनां वयं दंशपात्रा मा भूम । यत् सुखकारकं जीवजातं शुकसारिकादयो जन्तवः सन्ति, त्वं तेभ्योऽस्मान् सुखय ।

\customparagraph{\textbf{नेपाली भावार्थ}}

हे पृथिवी ! सर्पहरू, विच्छीहरू, जसका टोकाइले मानिसहरू तिर्खाउँछन् त्यस्ता जीवहरू; जसका टोकाइले मानिसहरू जाडामा कामेझैं काम्छन्  त्यस्ता जीवहरू; जसको टोकाइले मानिसहरू रिंगटा लागेझैं फन्फनी घुम्छन् त्यस्ता जीवहरू;  अन्य कीटहरूलाई खाएर प्रसन्न हुने जीवहरू र वर्षामा उत्पन्न हुने जीवहरू तिमीभित्र रहेका छन् । (तिमीमा प्वाल बनाएर बसेका हुन्छन्) । तिनीहरू यताउति विचरण गर्दै हाम्रा नजिकमा नआऊन् (हामी तिनको टोकाइका पात्र नबनौं) । जुन सुँगा, मैना आदि सुखकारक जन्तुहरू छन् तिनीहरूबाट हामीलाई सुख देऊ ।

\begin{visheshshloka}[center]

\textbf{ये ते पन्थानो बहवो जनायना रथस्य वर्त्मानसश्च यातवे ।}

\textbf{यैः सञ्चरन्त्युभये भद्रपापास्तं पन्थानं जयेमानमित्रमतस्करं}

\textbf{यच्छिवं तेन नो मृड ॥४७॥}

\end{visheshshloka}

\customparagraph{\textbf{अन्वयः}}

‘हे पृथिवि !’ ते ये जनायनाः बहवः पन्थानः (सन्ति) । ‘य’ रथस्य अनसः च यातवे वर्त्म (सन्ति) । यैः भद्रपापाः उभये सञ्चरन्ति, तं पन्थानं जयेम । अनमित्रम् अतस्करम् यत् शिवं ‘वर्त्म’ तेन नो मृड ।

\customparagraph{\textbf{शब्दार्थः}}

हे पृथिवि ! ते= तव ये जनायनाः= जनानां गमनागमनसाधनभूताः, बहवः= समधिकाः पन्थानः= राजमार्गवीथीप्रभृतयः सन्ति । तथा रथस्य= अश्वयानादेः, अनसः= गवादिशकटादेश्च यातवे= यातुम्, वर्त्म= वर्त्मानि सन्ति। यैः= जनरथादिमार्गैः, भद्रपापाः= पुण्यकृतः पापकृतः, सभ्या असभ्या वा उभये= द्वयेऽपि सञ्चरन्ति= गच्छन्ति । तम्= सर्वविधम्, पन्थानम्= मार्गम्, जयेम= सर्वतः प्रथमं लभेम । (मार्गस्य प्रथमप्राप्तियोग्यतैव मार्गजयः) तथा हे पृथिवि ! अनमित्रम्= असपत्नं शत्रुरहितम्, अतस्करम्= चौरसाहसिकाद्युपद्रविविरहितम्, यत् शिवम्= शोभनं कल्याणमयं साधुजनसंसेवनीयम्, वर्त्म= मार्गः, अस्ति । तेन मार्गेण नः= अस्मान् मृड= सुखय । तेन पथास्माकं सर्वादा गमनं स्यादिति ।

\customparagraph{\textbf{भावार्थः}}

हे पृथिवि ! तव भूमौ जनानां गमनागमनाय मुख्यसहायकादयो विविधा मार्गाः सन्ति। एवमेव रथशकटादिगमनाय तादृशा एव योग्या मार्गाः सन्ति । तैर्मार्गैः सभ्या असभ्याश्च जनाः सञ्चरन्ति । तादृशेषु मार्गेषु गमनाय वयं प्रथमस्थानं प्राप्नुमः । यो मार्गः शत्रुरहितः, चौरतस्कराद्युपद्रवविहीनः कल्याणकरश्चास्ति तेनैव मार्गेणास्मान् नय ।

\customparagraph{\textbf{नेपाली भावार्थ}}

हे पृथिवी ! तिम्रा भूभागमा  मान्छेहरू आउन र जानका लागि राजमार्ग, गोरेटो, घोडेटो आदिका साथै रथ, बयलगाडा हिँड्ने बाटाहरू छन् । ती बाटामा पुण्यात्मा र पापी पनि हिँड्छन् । चोर र फटाहाहरू पनि हिँड्छन् । हामीलाई ती बाटामा प्राथमिकता मिलोस् । हाम्रा लागि हिँड्ने बाटाहरू शत्रु नभएका, चोर डाँकाहरू पनि नभएका, कल्याणमय र सुखप्रद होऊन् । हामीले हतारमा कहिल्यै कुबाटो हिँड्न नपरोस् ।

\begin{visheshshloka}[center]

\textbf{मल्वं बिभ्रती गुरुभृद् भद्रपापस्य निधनं तितिक्षुः ।}

\textbf{वराहेण पृथिवी संविदाना सूकराय विजिहीते मृगाय ॥४८॥}

\end{visheshshloka}

\customparagraph{\textbf{अन्वयः}}

मल्वं बिभ्रती गुरुभृद् भद्रपापस्य निधनं तितिक्षुः वराहेण संविदाना पृथिवी सूकराय मृगाय विजिहीते ।

\customparagraph{\textbf{शब्दार्थः}}

मल्वम्= अनावश्यकं वस्तु, निरर्थकं प्रणिजातं वा, बिभ्रती= धारयन्ती, अत एव गुरुभृत्= अत्यधिकभारवाहिनी भद्रपापस्य= सुकृतकारिणः पापकारिणश्च जनस्य निधनम्= पञ्चत्वम्, तितिक्षुः= मर्षयित्री । सुकृतकारिणो मृत्यौ न पृथिव्याः शोकः, न च पापकृतो मरणे हर्ष इति वराहेण= वर-आहारेण गोधूमधान्यादिना संविदाना= संपद्यमाना सम्पन्ना समृद्धेत्यर्थः सस्यशालिनी इयं पृथिवी सूकराय= सुकराय शोभनकृतये (सूकर इति दीर्घश्छान्दसः पूरुष इतिवत्) मृगाय ‍‍‍= गुणगुणिनोरभेदेनान्वेषकाय (‘मार्गणं मृगया मृगः’ इत्यमरः) जनाय विजिहीते= स्वस्वरूपगुप्तिं त्यजति, सर्वानपि कामदुग्धे इत्यर्थः । यद्वा उत्तरार्धस्य अन्यथा अर्थयोजना यथा- वराहेण= सूकराकृतिना विष्णुना संविदाना= संविन्ना प्राप्ता उद्धृता इति यावत्, पृथिवी सूकराय= मृगाय वराहमृगरूपिणे विष्णवे विजिहीते= विशेषतया स्वस्य ज्ञापनाय लज्जादिकं त्यजति, पत्ये योषिदिव । स्वकीयमान्तरिकभावं प्रकटयतीत्यर्थः ।

\customparagraph{\textbf{भावार्थः}}

पृथिव्यां निष्प्रयोजनानि वस्तूनि बहूनि सन्ति । अकार्यकारिणो जना अपि नैके विद्यन्ते । तेषां गुरुभारं वहन्ती, सत्कर्मकारिणः पापकर्मकारिणश्च निधनं मर्षयित्री, गोधूमधान्याद्यन्नैः सम्पन्ना इयं पृथिवी सत्कर्मकारिणमन्वेषकं प्रति स्वकीयमान्तरिकं कामधेनुरूपं प्रकटयति । तेषां कृतेऽभिलषितं वस्तु ददातीत्यर्थः ।  अथवा वराहावतारधारिणा विष्णुना समुद्धृता इयं पृथिवी वराहरूपिणे विष्णवे पत्ये योषिदिव सर्वस्वमान्तरिकभावं प्रकटयति ।

\customparagraph{\textbf{नेपाली भावार्थ}}

पृथिवीमा काम नलाग्ने वस्तु धेरै छन् । कतिपय मानिसहरू पनि निरर्थक जीवन बिताइरहेका छन् जसलाई हामी पृथिवीका भार भन्छौं। त्यस्ता प्रयोजनहीन वस्तुहरूको भारलाई पनि चुपचाप बहन गर्ने, अत एवं ठूलो भार वहन गरिरहेकी, धर्मात्मा होस् चाहे पापी सबैको मृत्युलाई सहन गर्ने, धान्यादि अन्नले भरिपूर्ण पृथिवी राम्रो काम गर्ने अन्वेषक व्यक्तिका लागि आफ्नो आन्तरिक कामधेनु रूप प्रकट गर्दछिन् । अथवा- वराह अवतारका समयमा विष्णुद्वारा उद्धार गरिएकी पृथिवी पतिका लागि पत्नीले झैं विष्णुका लागि आन्तरिक भाव प्रकट गर्दछिन् ।

\begin{visheshshloka}[center]

\textbf{ये त आरण्याः पशवो मृगा वने हिताः सिंहा व्याघ्रा पुरुषादश्चरन्ति ।}

\textbf{उलं वृकं पृथिवी दुच्छुनामित ऋक्षीकां रक्षो अपवाधयास्मत् ॥४९॥}

\end{visheshshloka}

\customparagraph{\textbf{अन्वयः}}

ह‍े पृथिवि ! ते ये वने हिताः आरण्याः मृगाः पशवः (उत) पुरुषादः सिंहाः व्याघ्राः चरन्ति । (तान्) उलं वृकं दुच्छुनां ऋक्षीकां रक्षः (च) इतः अस्मत् अप वाधय ।

\customparagraph{\textbf{शब्दार्थः}}

हे पृथिवि ! ते= तव ये= कीर्त्यमानाः, वने= कानने हिताः= हितकारिणः, आरण्याः= अरण्यचारिणः, मृगाः= (कृष्णसाररुरुन्यङ्‍कुरङ्‍कुशंवररोहिषाः। गोकर्णपृषतैणर्श्यरोहिताश्चामरो मृगाः॥) इत्युक्ता बहुविधा मृगजातीयाः पशवः= चतुष्पदो जीवाः सन्ति। वन एव पुरुषादः= मनुष्यभक्षिणः सिंहाः= केशरिणः व्याघ्राः= द्वीपिनश्च चरन्ति= परिभ्रमन्ति । एतान् सर्वान् तथा उलम्= आरण्यम् (जीवविशेषम्) वृकम्= ईहामृगम्, दुच्छुनाम्= दुष्टाश्च ते श्वानो दुच्छ्वानस्तेषाम् ‘कर्मणि षष्ठी’ दुष्टान् शुनः, ऋक्षीकाम्= भल्‍लुकसङ‍्घं तरक्षुं वा रक्षः= रक्षोजातीयं च इतः= अस्मान्निवासभूतात्स्थानात्, अस्मत्= अस्मत्सकाशाच्च अपवाधय= अपसारय ।

\customparagraph{\textbf{भावार्थः}}

हे पृथिवि ! तव वनक्षेत्रे मृगादयः चतुष्पदा जीवाः, सिंहव्याघ्रादयो मनुष्यभक्षिणो हिंस्रका जन्तवश्च वसन्ति। एवमेव दुष्टाः श्वानः, भल्लूकाः, रक्षोजातिविशेषा भयप्रदाश्चान्ये जन्तवोऽपि तत्र वसन्ति । तादृशान् भयप्रदान् जीवान् अस्मन्निवासस्थानाद् अस्मभ्यश्च दूरीकुरु ।

\customparagraph{\textbf{नेपाली भावार्थ}}

हे पृथिवी ! तिम्रा वनक्षेत्रमा मृग आदि गरिएका चारखुट्टे जीव, सिंह, बाघ, भालुजस्ता मानिस खाने ह्रिंसक जीवका साथै उल (मृगविशेषः) ब्बाँसो, दुष्ट कुक्कुर, भालु, राक्षस जातिविशेष जस्ता भयानक जीवहरू पनि हुन्छन् । त्यस्ता डर लाग्दा जीवलाई हामीबाट टाढै राख । अर्थात् त्यस्ता भयानक जीव हाम्रा सामु नआऊन् ।

\begin{visheshshloka}[center]

\textbf{ये गन्धर्वा अप्सरसो ये चारायाः किमीदिनः ।}

\textbf{पिशाचान्त्सर्वा रक्षांसि तानस्मद् भूमे यावय॥५०॥}

\end{visheshshloka}

\customparagraph{\textbf{अन्वयः}}

हे भूमे ! ये गन्धर्वाः अप्सरसः ये च अरायाः किमीदिनः पिशाचान् सर्वा रक्षांसि (सन्ति) तान् अस्मत् यावय ।

\customparagraph{\textbf{शब्दार्थः}}

हे भूमे ! ये गन्धर्वाः, अप्सरसः इति प्रसिद्धा देवयोनिविशेषाः, ये च अरायाः= धनहीना दरिद्राः, किमीदिनः= बुभुक्षिताः सन्तोऽद्य किं भुञ्जाम इति वाचो व्यवहर्तारः, पिशाचान्= पिशिताशनाः पिशाचाः ‘तानित्युक्तेः प्रथमान्तार्थः’ सर्वाः= सर्वाणि नवनवतिविधानि रक्षांसि= देवविद्विषो जातयः सन्ति । तान्= सर्वाण्यपि अस्मत्= अस्मत्सकाशात् यावय= अपसारय । एभिर्गृहीता वयं न कदापि भवेमेत्याशीः ।

\customparagraph{\textbf{भावार्थः}}

हे भूमे ! तव प्रदेशे ये गन्धर्वाप्सरसो देवजातिविशेषाः सन्ति, ये धनहीनाः सन्ति, ये अद्य किं भुञ्जेम इति वदन्तो दरिद्राः सन्ति, ये च मांसभोजिनो राक्षसाः सन्ति, तान् सर्वान् अस्मत् अपसारय । अर्थाद् वयं एभिः सङ्‍गताः कदापि न भवेम ।

\customparagraph{\textbf{नेपाली भावार्थ}}

हे भूमि ! तिम्रा प्रदेशमा गन्धर्व, अप्सरा आदि देवजाति विशेष रहन्छन् । सम्पत्ति नभएका कङ्गालीहरू र आज के खाऊँ ? भन्ने कुरा गरिरने दरिद्रीहरू पनि बस्छन् । साथै मांसभोजी राक्षसहरू पनि तिमीमा आश्रित छन् । तिनीहरूबाट हामीलाई टाढा राख । अर्थात् हामी कहिल्यै तिनीहरूका सामू नपरौं ।

\begin{visheshshloka}[center]

\textbf{यां द्विपादः पक्षिणः संपतन्ति हंसाः सुपर्णाः शकुना वयांसि ।}

\textbf{यस्यां वातो मातरिश्वेयते रजांसि कृण्वँश्च्यावयँश्च वृक्षान् ॥}

\textbf{वातस्य प्रवामुपवामनु वात्यर्चिः ॥५१॥}

\end{visheshshloka}

\customparagraph{\textbf{अन्वयः}}

यां सुपर्णाः हंसाः शकुनाः वयांसि द्विपादः पक्षिणः संपतन्ति । यस्यां रजांसि कृण्वन् वृक्षान् च्यावयन् च मातरिश्वा वातः ईयते । (तस्याः) अर्चिः वातस्य प्रवाम् उपवाम् अनु वाति ।

\customparagraph{\textbf{शब्दार्थः}}

याम्= पृथिवीम् सुपर्णाः= शोभनपतनाः सुगतयो हंसाः= मानससरोनिवासिनश्चक्राङ्‍गाः तथा शकुनाः= शुभाशुभशकुननिदानभूताः, वयांसि= प्रडीनोड्‍डीनसंडीनाख्यनानागतिमन्तो द्विपादः= पादद्वयवन्तः पक्षिणः= संपतन्ति= इतरस्माद्देशादन्यत्र देशे विचरन्ति, डीयन्ते । यस्यां च पृथिव्यां स्वप्रवहणेन रजांसि= पांसून् कृण्वन्= विदधानः, वृक्षान्= तरून् च्यावयन्= पातयन् मातरिश्वा= मातर्यन्तरिक्षे श्वसिति चेष्टत इति मातरिश्वा वातो वायुर्देवः, ईयते= प्रचलति । (एवं वायुकृतं पृथिव्याः अपकारमुक्त्वा तस्याः क्षमात्मकमौदार्यमाह-) तस्याः पृथिव्याः, अर्चिः= गन्धात्मकः शोचिः ‘अर्च्यते पूज्यते परिचीयतेऽनेनेति परिचायको विशेषगुणः’ वातस्य= वायोः प्रवाम्= स्वाभाविकीं प्रगतिम्, उपवाम्= मनुष्यैः स्वानुकूलतयापादितां कृत्रिमां समीपगतिम्, अनु= अनुलक्ष्यीकृत्यैव वाति= प्रवहतीति क्षमागुणप्रशस्तिः ।

\customparagraph{\textbf{भावार्थः}}

यस्यां पृथिव्यां सुन्दरगतियुता हंसाः, शकुनप्रकाशका नानागतिमन्त द्विपदगामिनः पक्षिणश्च उड्‍डीयन्ते । यत्र रजांसि नयन्, वृक्षान् पातयन् अन्तरिक्षगतो वायुः प्रचलति । तस्याः पृथिव्या अर्चिर्वातस्य गन्धवाहिवायोः प्रवाम् स्वभाविकीं गतिम्, मनुष्यैः कृतां कृत्रिमगतिम् अनुलक्ष्यीकृत्यैव प्रवहति ।

\customparagraph{\textbf{नेपाली भावार्थ}}

जुन पृथिवीमा राम्रो उड्ने कला भएका हाँस, शुभाशुभ शकुनसूचक, अनेक प्रकारका गति भएका दुईखुट्टे पक्षीहरू उड्छन् । जहाँ जमिनको धुलो उडाउँदै र रूखका हाँगाहरू भाँच्दै वायु वहन्छ । जहाँ गन्धवाहिनी वायुको स्वाभाविक गति मानिसहरूले सिर्जना गरेको कृत्रिमवायुको गतिभन्दा पछि बहन्छ ।

\begin{visheshshloka}[center]

\textbf{यस्यां कृष्णमरुणं च संहिते अहोरात्रे विहिते भूम्यामधि ।}

\textbf{वर्षेेण भूमिः पृथिवी वृतावृता सा नो दधातु भद्रया}

\textbf{प्रिये धामनि धामनि ॥५२॥}

\end{visheshshloka}

\customparagraph{\textbf{अन्वयः}}

यस्यां भूम्याम् अधि अहोरात्रे कृष्णम् अरुणं च संहिते विहिते । सा भूमिः वर्षेण भद्रया आवृता वृता पृथिवी नः प्रिये धामनि धामनि दधातु।

\customparagraph{\textbf{शब्दार्थः}}

यस्यां भूम्यामधि= उपरि अहोरात्रे= अहोरात्रयोः, कृष्णम्= कृष्णरूपम्, अरुणम्= अरुणरूपम् च संहिते= सन्ध्ये विहिते= निर्मिते भवतः । (प्रातः कालिकी अरुणोदयसामयिकी हि सन्ध्या अरुणसंहिता। सायंकालिकी रात्रिसमाना हि सन्ध्या कृष्णा संहितेति शोभनतमत्वं वसुधायाः) सा= एवं स्तुता भूमिः= सस्यानां भावयत्री, वर्षेण= संवत्सरात्मकेन कालेन, भद्रया= कल्याण्या आवृता= सूर्यस्य स्वस्य वा परिक्रमेण वृता= युक्ता पृथिवी= देवी नः= अस्मान् प्रिये= हितकृति धामनि धामनि= प्रदेशे प्रदेशे दधातु= स्थापयतु। सर्वस्मिन्नपि काले शोभन एव स्थाने वयं निवासमाप्नुवामेति भावः ।

\customparagraph{\textbf{भावार्थः}}

यस्यां भूम्याम् अहोरात्रयोः कृष्णारुणे द्विविधे सन्ध्ये विहिते । तस्मात् सा पृथिवी सौन्दर्यपूर्णा अस्ति । एतादृशी भूमिः द्वादशमासात्मकेन कालेन सूर्यस्य स्वस्य वा कल्याणमय्या आवृत्त्या युक्ता विद्यते । इयं पृथिवी अस्माकं कृते प्रियप्रदेशयुतां निवासभूमिं दधातु । वर्षपर्यन्तं सुरम्ये स्थाने वयं निवासमाप्नुम इति ।

\customparagraph{\textbf{नेपाली भावार्थ}}

जुन पृथिवीमा सायङ्‌कालको कृष्ण सन्ध्या र प्रातःकालको अरुण सन्ध्या गरी दुई किसिमका सन्ध्या निर्धारित छन् । जसकारण पृथिवी सौन्दर्यपूर्ण देखिन्छ । तिनै पृथिवी वर्षभरि सूर्यको वा आफ्नो कल्याणमय परिक्रमा गर्दछिन् । त्यस्ती पृथिवीले हामीलाई रमणीय स्थानमा राखून् । अर्थात् हामी राम्रा ठाउँमा बस्न पाऔं ।

\begin{visheshshloka}[center]

\textbf{द्यौश्च म इदं पृथिवी चान्तरिक्षं च मे व्यचः ।}

\textbf{अग्निः सूर्य आपो मेधां विश्वेदेवाश्च सं ददुः ॥५३॥}

\end{visheshshloka}

\customparagraph{\textbf{अन्वयः}}

मे इदं द्यौः च पृथिवी च अन्तरिक्षं च व्यचः । मे अग्निः सूर्यः आपः विश्वेदेवाः च मेधां सं ददुः ।

\customparagraph{\textbf{शब्दार्थः}}

मे= मम इदम्= आवासस्थानम्, द्यौः= तृतीयो लोकः, च= तथा पृथिवी= जगतीतलम्, च अन्तरिक्षम्= मध्यमो लोकश्च व्यचः= व्याप्नुवन्तु । तेषां व्याप्त्यात्रैव वयं लोकत्रयैश्वर्यमुपलभेमहि इत्याशयः। तथा मे‍= मह्यं स्तोत्रे, अग्निः= पृथिवीस्थानो देवः, सूर्यः= द्युस्थानो देवः, आपः= मध्यस्थाना जलदेवता, अन्ये विश्वेदेवाः= सर्वे देवाश्च मेधाम्= धारणवतीं बुद्धिं धनं वा, सं ददुः= सम्यक् दद्युरिति प्रार्थना ।

\customparagraph{\textbf{भावार्थः}}

मम इमामावासभूमिं द्युलोकः, पृथिवीलोकः, अन्तरिक्षलोकश्च व्याप्नुवन्तु । तेषां लोकानां व्याप्त्या त्रयाणामेव लोकानां सुखमत्र प्राप्नुयाम । मह्यं अग्निसूर्यजलदेवा बुद्धिं धनं वा प्रददतु ।

\customparagraph{\textbf{नेपाली भावार्थ}}

मेरो यो आवासभूमिलाई स्वर्गलोक, पृथिवीलोक र अन्तरिक्ष तिनै लोकले व्याप्त गरून् । जसबाट यहाँ तीनवटै लोकको सुख प्राप्त होस् । मलाई अग्नि सूर्य र जलदेवता तथा अन्य देवताले बुद्धि वा धन प्रदान गरून् ।

\begin{visheshshloka}[center]

\textbf{अहमस्मि सहमान उत्तरो नाम भूम्याम् ।}

\textbf{अभीषाडस्मि विश्वाषाडाशामाशां विषासहिः ॥५४॥}

\end{visheshshloka}

\customparagraph{\textbf{अन्वयः}}

(हे पृथिवि !) अहं भूम्यां नाम उत्तरः सहमानः अस्मि । अभिषाट् विश्वाषाट् आशामाशां विषासहिः अस्मि ।

\customparagraph{\textbf{शब्दार्थः}}

हे पृथिवि ! (देवैः प्रदत्ताया मेधयान्वितः कर्तव्यनिष्ठश्च तव प्रसादात् तव स्तोता) अहम्, भूम्याम्= सर्वस्मिन्नपि जगतीतले, नाम= नाम्ना उत्तरः= उत्कृष्टतरत्वेनाभिमतः सहमानः= मर्षयिता (परस्य गुणदोषोन्नत्यादिमनोविकृत्युत्पादकानामर्थानां सोढा धीर इति यावत्), अस्मि= भवामि । प्रतिकारासामर्थ्यादेवं स्यादिति चेदुच्यते- अभीषाट्= परितः सपत्नानामभिभविता, विश्वाषाट्= सर्वदेशीयस्यापि विरोधिनो जनस्याभिभविता, आशामाशाम्= प्रतिदिशम्, (विद्यमानानां विरोधीनाम्) विषासहिः= विशेषेणाभिभविता अस्मि ।

\customparagraph{\textbf{भावार्थः}}

हे पृथिवि ! तव स्तोता अहम्, तव कृपया अस्मिन् भूमण्डले परेषां गुणदोषादीनाम् उत्कृतष्टतरः सहिष्णुः अस्मि । किं प्रतिकारसामर्थ्याभावाद् इति शङ्‍कायामाह- नैव, भूमण्डले स्थितानां सर्वेषां विरोधीनां विजेता अहं ‘दिग्विजयी’ इत्युपनाम्ना ख्यातोऽस्मि ।

\customparagraph{\textbf{नेपाली भावार्थ}}

हे पृथिवी ! तपाईंको स्तुतिकारक म उत्तम सहनशील पुरुष हुँ । (म अरूका गुणदोष आदि सहन गर्न सक्दछु, अर्थात् तिनीहरूबाट म विकृत हुन्न)। प्रतिकार गर्ने सामर्थ्य नभएर म सहनशील भएको होइन । मैले आफ्ना चारैतिर भएका शत्रुहरूलाई दबाएको छु । सबै ठाउँमा रहेका आफ्ना विरोधीहरूलाई पनि दबाएको छु । विभिन्न दिशामा रहेका विपक्षीहरूको पनि दमन गरेको छु । (त्यसैले म दिग्विजयी हुँ)

\begin{visheshshloka}[center]

\textbf{अदो यद्देवि प्रथमाना पुरस्ताद् देवैरुक्ता व्यसर्पो महित्वम् ।}

\textbf{आ त्वा सुभूतमाविशत्तदानीमकल्पयथाः प्रदिशश्चतस्रः ॥५५॥}

\end{visheshshloka}

\customparagraph{\textbf{अन्वयः}}

हे देवि ! (त्वं) यत् पुरस्तात् देवैः उक्ता प्रथमाना अदः महित्वं व्यसर्पः । तदानीं त्वा सुभूतम् आ आविशत्, चतस्रः प्रदिशः अकल्पयथाः।

\customparagraph{\textbf{शब्दार्थः}}

हे देवि ! त्वं यत्= यदा पुरस्तात्= पूर्वम्, देवैः= अमरैः, उक्ता= प्रार्थिता सती प्रथमाना= विस्तारं लभमाना, अदः= एतावत् महित्वम्= पृथुत्वम्, व्यसर्पः= विशेषेणागच्छः । (देवप्रार्थनया त्वया स्वस्य महान् विस्तारः कृत इति भावः) तदानीम्= तदा त्वा= त्वाम् सुभूतम्= शोभनं सात्त्विकभावोपेतं भूतजातम्, आ आविशत्= सन्निविष्टम् । तदा च त्वम् चतस्रः= प्रदिशः, दिग्विभागान् अकल्पयथाः= अकार्षीः । (तदा हि विशालायां त्वयि वसुधायां निर्भया मनुष्या यथाभिलाषं सन्निवेशं चक्रिरे । अथ च प्राच्यावाच्यपाश्चात्त्यौदीच्यादिशब्दैर्दिशां नाम्ना मनुष्याँस्त्वं व्यभज इति भावः)

\customparagraph{\textbf{भावार्थः}}

हे पृथिवि ! प्राचीनकाले देवानामाग्रहेण त्वं विस्तृता जाता । तदा त्वयि सात्त्विकभावयुताः नानाविधाः प्राणिनो न्यवसन् । तदा त्वं प्राच्यादिचतस्रः प्रदिशः कल्पितवती । (तदा विशातायां त्वयि निर्भया मनुष्या यथाभिलाषं न्यवसन् । तासु दिशासु स्थिता जनाः पौरस्त्याः, पाश्चात्त्याः, उदीच्याः, अवाच्या इत्यभिधीयन्त) ।

\customparagraph{\textbf{नेपाली भावार्थ}}

हे पृथिवी ! प्राचीन कालमा देवताहरूको आग्रहले तिमीले आफूलाई विस्तार गर्‍यौ । त्यति वेला तिम्रा विस्तारित क्षेत्रमा सात्त्विक प्राणीहरू बस्न थाले । त्यसपछि तिमीले चार दिशाको कल्पना गर्‍यौ । (त्यस समयमा मानिसहरू निर्भय भएर चारै दिशामा बसे । त्यहाँ बस्ने मानिसहरू तिम्रा दिशाका आधारमा पूर्वीया, पश्चिमा इत्यादि भनिन थाले)।

\begin{visheshshloka}[center]

\textbf{ये ग्रामा यदरण्यं याः सभा अधि भूम्याम् ।}

\textbf{ये सङ्‍ग्रामाः समितयस्तेषु चारु वदेम ते ॥५६॥}

\end{visheshshloka}

\customparagraph{\textbf{अन्वयः}}

(हे पृथिवि !) भूम्याम् अधि ये ग्रामाः, यत् अरण्यम्, याः सभाः, ये सङ्‍ग्रामाः (याः) समितयः, तेषु ते चारु वदेम ।

\customparagraph{\textbf{शब्दार्थः}}

(हे जन्मभूमि !) अस्यां भूम्याम्= जगतीतले अधि= उपरि ये ग्रामाः= संवसथाः, जननिवासस्थानानि, यच्च अरण्यम्= वनम्, याः सभाः= शास्त्रीयचर्चास्थलानि, ये सङ्‍ग्रामाः= रणाङ्‍गणानि, सङ्‍ग्रामसाधनानि मल्लकलाशिक्षास्थलानि वा, या समितयः= समानगमनविधायिन्यः सामाजिकविचारपर्षदादयश्च सन्ति, तेषु= सर्वेष्वपि ते= अस्मन्निवासभूतायास्तव कृते, चारु= शोभनम्, वदेम= कथयेम । (ग्रामादिषु सर्वेष्वपि स्थानेषु जन्मभूम्यास्तव यशः प्रख्यापको हितकृच्च भूयासम्,  त्वत्कृते चारुवचनं ब्रूयां वेति भाव)

\customparagraph{\textbf{भावार्थः}}

हे जन्मभूमि ! तवोपरि जननिवासस्थानभूता ये ग्रामाः, यत् अरण्यम्, या विद्वद्विचारमन्थनाय समायोजिताः सभाः, ये परस्परश्रेष्ठतायाः परीक्षणे सञ्जाताः सङ्‍ग्रामाः, याश्च धर्मसभाः समाजिकसभा वा सन्ति, तेषु सर्वेषु स्थानेषु तव कृते वयं शोभनं प्रशस्तिमयं वचनं कथयेम ।

\customparagraph{\textbf{नेपाली भावार्थ}}

हे जन्मभूमि ! तपाईंका पृष्ठभूमिमा रहेका मानिस बस्ने गाउँ, जङ्गल, विद्वान्‌हरूका शास्त्रार्थसभा, युद्धस्थल, धर्मसभा (जनसभा) आदि ठाउँमा हामी तपाईंका बारेमा प्रशस्तिपूर्ण राम्रा कुरा मात्र बोलौँ ।

\begin{visheshshloka}[center]

\textbf{अश्व इव रजो दुधुवे वि ताञ्जनान्य आक्षियन् पृथिवीं यादजायत ।}

\textbf{मन्द्राग्रेत्वरी भुवनस्य गोपा वनस्पतीनां गृभिरोषधीनाम् ॥५७॥}

\end{visheshshloka}

\customparagraph{\textbf{अन्वयः}}

ये (मे) पृथिवीम् आक्षियन् तान् जनान् अश्वः रजः इव वि दुधुवे । मन्द्रा अग्रेत्वरी भुवनस्य गोपा वनस्पतीनां ओषधीनां गृभिः (भूमिः) यात् अजायत (तत् अपि अश्वः रजः इव दुधुवे)

\customparagraph{\textbf{शब्दार्थः}}

ये= जनाः, मम स्वामित्वेनाधिकृतां पृथिवीम्= (ग्रामारण्यसभासङ्‍ग्रामसमित्यात्मिकां पूर्वोक्ताम्) भूमिम्, आक्षियन्= मत्प्रतिकूल्येन निवसन्ति, तान् परिपन्थिनो जनान्= मनुष्यान् अश्वः= घोटकः, रजः= धूलिम् यथा धुनुते तथा विदुधुवे, अर्थाद् विशेषेण कम्पयामि अपसारयामीति भावः । किञ्च मन्द्रा= मादयित्री मदनशीला, अग्रेत्वरी= अग्रे त्वरयतीति अग्रेत्वरी मम प्रगतिप्रदात्री, भुवनस्य= लोकस्य गोपाः= पालयित्री, वनस्पतीनामोषधीनां च गृभिः= ग्रहणशीला पोषयित्री भूमिः यात्= यस्मात्कारणात् एवमन्यस्वामिका अजायत= सञ्जाता तदपि अश्वतनुसंलग्नधूलिकणिका इव विदुधुवे‍= धुनोमि ।

\customparagraph{\textbf{भावार्थः}}

ये जना मम स्वामित्वयुतां भूमिमधिक्रम्य निवसन्ति, तान् अहम् अपसारयामि । यथा अश्वः स्वशरीरे संलग्नां धूलिं शरीरकम्पेन दूरीकरोति तथैव तानपि दूरीकरोमि । मां हर्षदायिनी, प्रगतिदायिनी, भुवनरक्षिका, वनस्पतीः गोधूमादीन् च धारयित्री पृथिवी यस्य कारणेन अन्यस्वामिका जाता, तदपि तथैव दूरीकरोमि ।

\customparagraph{\textbf{नेपाली भावार्थ}}

जति पनि मानिस मेरो स्वामित्व भएको भूमिमा प्रतिकूल व्यवहार गर्दै जबरजस्ती बसेका छन्, तिनलाई ‘घोडाले शरीरमा लागेको धुलो शरीर हल्लाएर झारेझैं, हटाउँछु । मलाई हर्ष दिने, उन्नति दिने, लोकको रक्षा गर्ने वनस्पति र अन्न धारण गर्ने ती पृथिवी जुन कारणले तिनीहरूको स्वामित्वमा छिन्, त्यस कारणलाई पनि घोडाले शरीर हल्लाएर धुलालाई टाढा गराएझैँ हटाउँछु ।

\begin{visheshshloka}[center]

\textbf{यद्वदामि मधुमत्तद्वदामि यदीक्षे तद्वनन्ति मा ।}

\textbf{त्विषीमानस्मि जूतिमानवान्यान् हन्मि दोधतः ॥५८॥}

\end{visheshshloka}

\customparagraph{\textbf{अन्वयः}}

(हे पृथिवि !) यत् वदामि तत् मधुमत् वदामि । यत् ईक्षे तत् मा वनन्ति । त्विषीमान् जूतिमान् अस्मि । दोधतः अन्यान् अव हन्मि ।

\customparagraph{\textbf{शब्दार्थः}}

(हे पृथिवि !) त्वत्प्रसादादहं यत्= वाक्यम्, वदामि= कथयामि तत्= वाक्यम्, मधुमत्= मधुरं प्रियम्, वदामि= कथयामि । यत्= चराचरम्, ईक्षे= अवलोकयामि, तदीक्षिताः प्राण्यप्राणिजातयो मा= माम्, वनन्ति= संभजनते मामुपकुर्वन्ति । एवञ्चाहं त्विषीमान्= तेजस्वी जूतिमान्= प्रीतियुक्तश्च अस्मि= भवामि । तथा दोधतः= ममानिष्टचिन्तनद्वारा मामभिप्रकम्पयतः, अन्यान्= विद्विषः, अवहन्मि=नाशयामि ।

\customparagraph{\textbf{भावार्थः}}

हे पृथिवि ! तव प्रसादात् स्वकीयदोषत्यागकारणादहं यत् वचनं वदामि तत् सुमधुरं वदामि । यानि दृश्यानि अवलोकयामि, तानि सर्वाणि मां प्रति प्रीतियुक्तानि भवन्ति । अनेनाहं तेजस्वी सर्वप्रियश्च भवामि । मां विचालयितुमभिलषतः शत्रून् च नाशयामि ।

\customparagraph{\textbf{नेपाली भावार्थ}}

हे पृथिवि ! तपाईंका कृपाले आफ्ना दुर्गुणलाई त्याग गरेका कारण म जुन वचन बोल्छु, त्यो सुमधुर हुन्छ । जुन दृश्यलाई हेर्छु, ती सबै मप्रति सकारात्मक हुन्छन् । यसकारण म तेजस्वी र सर्वप्रिय छु । जसले मलाई विचलित बनाउँन खोज्छन् ती शत्रुहरूलाई परास्त गर्दछु ।

\begin{visheshshloka}[center]

\textbf{शन्तिवा सुरभिः स्योना कीलालोध्नी पयस्वती ।}

\textbf{भूमिरधिव्रवीतु मे पृथिवी पयसा सह ॥५९॥}

\end{visheshshloka}

\customparagraph{\textbf{अन्वयः}}

शन्तिवा सुरभिः स्योना कीलालोध्नी पयस्वती भूमिः पृथिवी पयसा सह मे अधिव्रवीतु ।

\customparagraph{\textbf{शब्दार्थः}}

शन्तिवा= शान्तिं वाति वापयति गमयतीति शान्तिप्रदा, सुरभिः= गन्धवती स्योना= सुखरूपा कीलालोध्नी= बीजस्याविनश्वरो भागः कीलालः स ऊधसि गर्भे विद्यते यस्या साः कीलालोध्नी- बीजधारिणी, पयस्वती= रसवती जलबहुला भूमिः प्रभूतसस्योत्पादयित्री, पृथिवी= भूमिः, पयसा साधर्म्यात् स्वान्तत्त्वेन सह मे= मह्यम् अधिव्रवीतु= संवदतात् । (सान्तस्तत्त्वायाः पृथिव्या अहमधिपो भूयासमिति पृथिव्याः संवदनम्)

\customparagraph{\textbf{भावार्थः}}

शान्तिप्रदा, गन्धवती, सुखप्रदा, बीजयुता, रसवती च भूमिः अन्तस्तत्त्वेन सह मह्यं संवदतात्, अर्थात् सान्तस्तत्त्वायाः पृथिव्या अहमधिपो भूयासम् ।

\customparagraph{\textbf{नेपाली भावार्थ}}

शान्ति दिने, गन्धवती, सबै प्राणीलाई सुख दिने, विभिन्न किसिमका बीज धारण गर्ने, रसिली एवम् उर्वरा भूमिले अन्तस्तत्त्वको साधर्म्यका कारण मसँग सादृश्य स्थापित गरून् । अर्थात् यस्ती पृथिवी मेरो स्वामित्वमा रहून् ।

\begin{visheshshloka}[center]

\textbf{यामन्वैच्छद्धविषा विश्वकर्मान्तरर्णवे रजसि प्रविष्टाम् ।}

\textbf{भुजिष्यं १ पात्रं निहितं गुहा यदाविर्भोगे अभवन्मातृमद्भ्यः॥६०॥}

\end{visheshshloka}

\customparagraph{\textbf{अन्वयः}}

अन्तरर्णवे रजसि प्रविष्टां यां विश्वकर्मा हविषा अन्वैच्छत् । यत् भुजिष्यं पात्रं गुहा निहितम् ‘अभवत्’ (तत्) मातृमद्भ्यः भोगे आविरभवत् ।

\customparagraph{\textbf{शब्दार्थः}}

अन्तरर्णवे‍= समुद्रमध्ये रजसि= जले प्रविष्टाम्= अन्तर्हिताम्, याम्= पृथिवीम्, विश्वकर्मा= नानाविधक्रियाकुशलो विष्णुर्देवः, हविषा= कर्मणा अन्वैच्छत्= अन्विष्टवान् समुद्दधारेति भावः । तथा यत् भुजिष्यम्= भोगसाधनम्, पात्रम्= स्नेहात्मकम्, गुहा= गुहायामन्तर्हृदयावकाशे निहितमभवत्= स्थापितं तिष्ठति, तत् मातृमद्भ्यः= समानामातृकेभ्यः, मात्रा चात्र पृथिवी । ततश्च तत्तद्देशीयपृथिवीजातेभ्यो भोगे= व्यवहारकाले आविरभवत्= प्रकाशमेति । (पृथिव्या हि देशविशेषस्य यो विशेषो रूपभाषागतिव्यवहारादिः सोऽन्तर्हितोऽपि समानदेशिकत्वादेकभूमिकमातृमद्भ्यः कृते कार्यकाले आविर्भवत्येव । सर्वोऽपि जनः स्वमातृदेशीयजनाय स्निह्यत्येवेति मानुषः स्वभावः पार्थिवेन महिम्ना सम्पद्यते)

\customparagraph{\textbf{भावार्थः}}

समुद्रजलान्तरे निमग्नां पृथिवीं जगत्पालको विष्णुर्वराहावतारं संधृत्य कर्तव्यपरायणतया समुद्दधार । यद् भोगप्राप्तेः साधनभूतं स्नेहं मानवस्यान्तस्करणे निहितं भवति, तद् जन्मभूमेः प्रेमरूपं मातृस्नेहं व्यवहारकाले यथावसरं प्रकाशमेति ।

\customparagraph{\textbf{नेपाली भावार्थ}}

समुद्रको जलभित्र डुबेको पृथिवीलाई कर्तव्यपरायण भगवान् विष्णुले वराह अवतार धारण गरी उद्धार गर्नुभयो । त्यो भोग प्राप्तिको साधन भएको पृथ्वीको माया मानिसका हृदयरूपी गुफामा रहेको हुन्छ, जुन प्रेम मातृभूमिको स्नेहका रूपमा व्यवहार कालमा यथा अवसरमा प्रकट हुन्छ ।

\begin{visheshshloka}[center]

\textbf{त्वमस्यावपनी जनानामदितिः कामदुघा पप्रथाना ।}

\textbf{यत्त ऊनं तत्त आपूरयति प्रजापतिः प्रथमजा ऋतस्य ॥६१॥}

\end{visheshshloka}

\customparagraph{\textbf{अन्वयः}}

(हे पृथिवि !) त्वम् जनानाम् आपवनी, अदितिः, कामदुघा पप्रथाना असि । यत् ते ऊनं तत् ते ऋतस्य प्रथमजाः प्रजापतिः आपूरयति ।

\customparagraph{\textbf{शब्दार्थः}}

हे पृथिवि ! त्वम्= भवती, एकस्मिन्नेव प्रदेशे नानादेशवासिनाम् जनानाम्=मनुष्याणाम्, आवपनी= आवपनस्थानीया सम्मीलनकारिणी असि= भवसि । तथा देशभेदेन पृथक्त्वेन ज्ञायमानापि त्वं स्वरूपतः अदितिः= अखण्डिता एकरूपा, कामदुघा= विना कमपि भेदभावं समेषां सर्वविधकामानां पूरयित्री, पप्रथना= विशालकाया चासि । आपवनरूपत्वाच्च तव यत् ते क्षुद्रविचारवतां मनुष्याणां क्षोभेण ऊनम्= न्यूनम् भवति, धारकं तत्त्वं सत्यादिप्रथममन्त्रज्ञापितं तदूनं ते= तव ऋतस्य= यज्ञस्य वेदस्य कर्मणो वा प्रथमजाः= मुख्यज्ञाता आदर्शपुरुषः, अग्रगामी वा प्रजापतिः= प्रजापतिरूपेण समाजस्य देशस्य राष्ट्रस्य वा निर्मातृत्वेनाभिमतो महान् व्यक्तिः, आपूरयति= रिक्तं पूरयति ।

\customparagraph{\textbf{भावार्थः}}

हे पृथिवि ! त्वं समस्तदेशवासिनाम् एकत्र संमीलनकारिणी असि (विभिन्नेषु स्थानेषु समुत्पना जनाः कस्मिन्नपि एकस्मिन् स्थले एकत्रिता भवितुं शुक्नुवन्ति) । वस्तुतस्त्वमखण्डिता, सर्वेषां सर्वविधकामानां पूरयित्री, विशालाकारयुता चासि। तव प्रदेशेषु क्षुद्रकर्मकारिणो जना अपि सन्ति तेषां कारणेन यत् तव गुणेषु न्यूनं जायते तत् सर्वविधिज्ञ आदर्शपुरुषः प्रजापतिः पूरयति ।

\customparagraph{\textbf{नेपाली भावार्थ}}

हे पृथिवि ! तिमी समस्त भूगोलका मानिसहरूलाई मिलाएर राख्छ्यौ । तिमी अखण्डित रूपले रही सबैका सबै किसिमका चाहाना पूरा गर्दछ्यौ । तिमी विशाल स्वरूपकी छ्यौ । तिम्रा केही स्थानमा क्षुद्रजनले तिम्रो गुणमा न्यूनता गराए पनि सर्वज्ञ, कार्यकारी, आदर्शपुरुष प्रजापतिले त्यस न्यूनतालाई हटाएर तिमीलाई पूरा गराउँछन् ।

\begin{visheshshloka}[center]

\textbf{उपस्थास्ते अनमीवा अयक्ष्मा अस्मभ्यं सन्तु पृथिवि प्रसूताः।}

\textbf{दीर्घं न आयुः प्रतिबुध्यमाना वयं तुभ्यं बलिहृतः स्याम ॥६२॥}

\end{visheshshloka}

\customparagraph{\textbf{अन्वयः}}

हे पृथिवि ! ते प्रसूताः उपस्थाः अनमीवाः अयक्ष्माः अस्मभ्यं सन्तु। वयं नः आयुः दीर्घं प्रतिबुध्यमानाः तुभ्यं बलिहृतः स्याम ।

\customparagraph{\textbf{शब्दार्थः}}

हे पृथिवि ! ते= तव सकाशात् प्रसृताः= उत्पन्नाः, उपस्थाः= उप समीपे तिष्ठन्तीति उपस्थाः। पृथिवीयं वसुधेति रत्नादिवसून्येव समीपवर्तिनः, अथवा उपस्थेन पालनीया आनन्ददायिनो वा रत्नहिरण्यादयः । अनमीवाः= अमीवा अभ्यमनवान् रोगभूतो भयदाता तद्रहितोऽनमीवाः । अयक्ष्मा= अयो लोहम्, तत् क्ष्मा= आधारभूमिर्येषां ते, अयक्ष्माः, आयसपेटिकायां निवासवन्तः। लोहपात्रे स्थापिता इत्यर्थः। अस्मभ्यम्= स्तोतृभ्यः सन्तु= जायन्ताम् । तथा वयं नः= अस्माकम् आयुः= वयः, दीर्घम्= समधिकसमयं यावत् प्रतिबुध्यमानाः= अवगच्छन्तः, तुभ्यम्= पृथिव्यै बलिहृतः स्याम= यज्ञपूजादिविधयोपासनकर्तारो भवेम । (‘गृहीत इव केशेषु मृत्युना धर्ममाचरेत्’ इत्युक्तेः, स्वस्य जरठत्वं मत्वा पृथिव्याः पोषकानां प्रथममन्त्रे समुल्लिखितानां सत्य-धैर्य-मानस-यथार्थभाषण-क्षत्र-दीक्षा- तपो-यज्ञेत्यष्टानामनुष्ठातारो भूयास्म येन इयमवष्टब्धा जगत उपकुर्यात्)

\customparagraph{\textbf{भावार्थः}}

हे पृथिवि ! तव सकाशात् समुत्पन्नाः, तव निकटस्थाः (आनन्ददायिनो वा), रोगभयादिहीनाः, लोहपात्रे स्थापिता रत्नहिरण्यादयोऽस्मभ्यं जायन्ताम् । वयं स्वकीयं दीर्घमायुरवगच्छन्तस्तुभ्यं यज्ञपूजादिविधयोपासनकर्तारो भवेम । अर्थात् अस्माकं दीर्घमायुः प्रबोध्य पृथिव्याः स्थिरतायाः कृते सत्यधैर्यादिपृथिवीपोषकानामनुष्ठानानां कर्तारो भवेम इत्यर्थः ।

\customparagraph{\textbf{नेपाली भावार्थ}}

हे  पृथिवि ! तिम्रा भूमिबाट उत्पन्न, तिम्रा समीपमा रहेका, रोग भय आदिले रहित र फलामका पात्रमा राखिएका जुन रत्न सुवर्ण आदि छन् ती हामी स्तुतिकर्तालाई प्राप्त होऊन् । हामी आफ्नो दीर्घ आयुलाई जानेर तिम्रो स्थायित्वका लागि यस सूक्तको पहिलो मन्त्रमा पोषक (अवष्टम्भक) भनेर बताए सत्य आदि आचरणको पालन गर्दछौं ।

\begin{visheshshloka}[center]

\textbf{भूमे मातर्नि धेहि मा भद्रया सुप्रतिष्ठितम् ।}

\textbf{संविदाना दिवा कवे श्रियां मा धेहि भूत्याम् ॥६३॥}

\end{visheshshloka}

\customparagraph{\textbf{अन्वयः}}

हे भूमे ! मातः !  दिवा संविदाना (त्वम्)भद्रया सुप्रतिष्ठतं नि धेहि । (अथ) ‘मा’ कवे श्रियां भूत्यां ‘नि’ धेहि ।

\customparagraph{\textbf{शब्दार्थः}}

हे भूमे ! मातः ! मातृवदाराधनीये ! पालयित्रि ! पृथिवि देवि ! दिवा= अहर्निदानभूतेन सूर्येण संविदाना= संमता सती मा= माम् भद्रया= कल्याण्या बुद्ध्या सुप्रतिष्ठितम्= समादरभाजनम्, निधेहि= संस्थापय । अथ मा= माम् स्तोतारं कवे= कुः पृथिवी ततश्चतुर्थी तदात्म्यात्पार्थिवाय भोगजाताय तत्प्राप्तय इत्यर्थः । श्रियाम्= लक्ष्म्याम्, भूत्याम्= ऐश्वर्ये च निधेहि । (धनपतित्वं हि श्रियां धारणम्, धनस्य फलभोक्तृत्वम्, ततो यशः प्राप्तिश्चेति भूत्यां धारणम्, तदुभयमपि कुर्विति भावः)

\customparagraph{\textbf{भावार्थः}}

हे भूमे !  हे मातः ! भगवता सूर्येण सह संमिल्य कल्याणप्रदया मनीषया आदरपात्रं निधेहि। मां बुद्धिमन्तं कुरु इत्यर्थः । एवमेव मां पार्थिवानां भोगादीनां प्राप्तये लक्ष्म्यामैश्वर्ये च स्थापय ।

\customparagraph{\textbf{नेपाली भावार्थ}}

\{बुद्धि, लक्ष्मी र ऐश्वर्यको प्रार्थना-\} हे पृथिवि ! हे माता ! सूर्यदेवसँग सहकार्य गरी मलाई कल्याणकारी बुद्धिमा स्थिर गराऊ । साथै पृथिवीका अनेक भोग भोग्नका लागि मलाई लक्ष्मी र ऐश्वर्यमा स्थापित गराऊ ।

\textit{प्रस्तुत अथर्ववेदीय पृथिवीसूक्तको शब्दार्थ, संस्कृत भावार्थ र नेपाली भावार्थ गोपालचन्द्रशास्त्रीद्वारा सम्पादित सूक्तरत्नसङ्‌ग्रहमा सङ्कलित ज्योत्स्नाविवृतिसहित पियूषकणिकाभाष्यमा आधारित छ ।}

\begin{visheshshlokatwo}[center]{84}

\textbf{पृथिवीसूक्तं समाप्तम्}

\end{visheshshlokatwo}

\begin{center}

\textbf{अभ्यासार्थं प्रश्नाः}

\end{center}

\textbf{अतिसङ्‌क्षिप्तोत्तरापेक्षिणः प्रश्नाः  (१)}

१. प्रथमतः किम् नाम वेद आसीत् ? तस्य विभागः केन कृतः ?

२. विभक्ताः संहिताः कति सन्ति ? तेषां नामनि कानि ?

३. केन हेतुना वेदः ‘त्रयी’ कथ्यते ?

४. ऋग्वेदस्य उपलब्धा शाखाः कति सन्ति ?

५. वेदं सम्यक् ज्ञातुं कति अङ्गानि पठ्यन्ते ? तेषां नामानि कानि ?

६, यज्ञे कति ऋत्विजो भवन्ति ? तेषां नामानि कानि ?

७. ऋत्विक्षु सर्ववेदज्ञः कः ?

८. वेदाङ्गानि कानि ? तेषां नामानि कानि ?

९. नासदीयसूक्तं कस्माद् ग्रन्थादुद्धृतमस्ति ?

१०. उषस्सूक्तस्याधारो ग्रन्थः कः ?

११. संवादसूक्तं कस्या लौकिकसाहित्यविधाया बीजम् ?

१२. वैदिकयागानां  व्याख्यात्मकानां ग्रन्थानां नाम किम् ?

१३ ‘अध्वर्युः’ कस्य वेदस्य ऋत्वग् अस्ति ?

१४. पाठ्यांशे निर्धारितस्य उषस्सूक्तस्याधारग्रन्थः कः ? तस्य कस्मिन् भागे तल्लभ्यते ?

१५. पृथ्वीसूक्तं कस्य ग्रन्थस्य कस्माद् भागादुद्‌धृतमस्ति ?

\textbf{सङ्‌क्षिप्तोत्तरापेक्षिणः प्रश्नाः (६)}

१. ऋग्वेदसंहितायाः सामान्यः परिचयो देय ?

२. संहिताग्रन्थपाठे निर्धारिता ऋत्विजः समुल्लिख्य तेषां कार्याणि उल्लेखनीयानि ।

३. उपनिषदा प्रकाश्यो विषयः कः ? तस्य सामान्यः परिचयः समुपस्थाप्यताम् ?

\textbf{विवेचनात्मकाः प्रश्नाः (१०)}

१ संस्कृतवाङ्‌मये वैदिकसाहित्यस्य महत्त्वं विवेचनीयम् ।

२ ‘ऋग्वेदः लौकिकसाहित्यस्योपजिव्यमस्ति’ इत्याभाणकं समर्थनीयम् ।

३ उषस्सूक्ते प्रकृतिवर्णनस्य सौन्दर्यं प्रकाशनीयम् ।

४ नासदीयसूक्तस्य रहस्यमयं दार्शनिकपक्षमुपस्थापनीयम् ।

\specialmahakhandatwo{\textbf{द्वितीयः पाठः}}{आख्यायिकाध्ययनम्}

\customparagraph{\textbf{ब्राह्मणग्रन्थपरिचयः}}

पौरस्त्यवाङ्‌मये संहिताग्रन्थानन्तरं ब्राह्मणग्नन्थानां क्रमः आयाति । तत्र वेदसंहितायां सूत्ररूपेण स्थिता यज्ञस्य विधानादिविषया विवेचिताः प्राप्यन्ते । ‘बृहि’ वृद्धौ धातोर्मनिन्‌प्रत्यये सति ‘ब्रह्मन्’शब्दो निष्पद्यते । अनेन वेदमन्त्रयज्ञादिकं सूच्यते । ‘ब्रह्मन्’ शब्दाद् व्याख्यानार्थे ‘अण्’ प्रत्ययसंयोगात् ‘ब्राह्मणः’ इति शब्दो निर्मीयते । शब्दोऽयं ग्रन्थार्थे ‘ब्राह्मणम्’ इति नपुंसके प्रयुज्यते । अनया व्युत्पत्त्या मन्त्राणां यज्ञीयविधीनाञ्च व्याख्यानप्रतिपादको ग्रन्थ इत्यर्थो ज्ञायते । अत्र व्याख्यानशब्देन केवलं पर्यायशब्दादिकमुपस्थाप्य कृता शब्दार्थव्याख्या नापितु वेदमन्त्राणां यज्ञाध्यात्मतत्त्वादिरहस्यमयगूढार्थप्रतिपादिका विवृतिर्बोध्या ।

ब्राह्मणग्रन्थस्य स्वरूपं प्रतिपादयन् तैत्तिरीयसंहिताया भाष्यकारो भट्‌टभास्कर उल्लिखति यत्- “ब्राह्मणं नाम कर्मणस्तन्मन्त्राणां च व्याख्यानग्रन्थः” इति । ब्राह्मणग्रन्थस्य प्रतिपाद्यविषये विद्वांसो नैकानि मतानि प्रस्तुवन्ति । तेषु विधिः, अर्थवादश्चेति द्वौ प्रसिद्धौ विद्येते । तत्र विधिविषयान्तर्गतम्- इदं  कार्यं कर्तव्यम्, इदं तु परित्यज्यम्; इदं यस्मात् कारणात् कर्तव्यम्, इदं यस्मात् कारणात् परित्यक्तव्यमित्यादीनि विधिनिषेधज्ञापकानि  वेदवाक्यानि भवन्ति । अर्थवादे तु विधिवाक्येन निर्दिष्टस्य कर्मणः प्रशंसा एवं विधिनिषिद्धार्थस्य निन्दा वर्ण्यते । ब्राह्मणग्रन्था अपि चतसृभिः संहिताभिः सम्बद्धाः सन्ति । तेषु ऋग्वेदस्य शाकलशाखाया ऐतरेयब्रह्मणम् शाङ्खायनशाखायाः कौषीतकी(शाङ्‌ख्यायन)ब्राह्मणमिति द्वौ ब्राह्मणग्रन्थौ स्तः । तयोरप्यस्माकं पाठ्यांशेन सम्बद्धं ‘शुनःशेपाख्यानम्’ ऐतरेयब्राह्मणस्य सप्तमपञ्चिकायास्त्रयत्रिंशेऽध्याये सङ्‌गृहीतं विद्यते ।

\specialkhanda{\textbf{ऐतरेयब्राह्मणान्तर्गतं शुनःशेपाख्यानम्}}

\customparagraph{\textbf{सारः}}

सूर्यवंशीयो राजा हरिश्चन्द्रः सन्तानहीन आसीत् । नारदस्य प्रेरणया तेन वरुणः प्रार्थितः । वरुणेन पुत्रः प्रदत्तः, परं ‘कतिचित् कालपश्चात् यज्ञीयविधिना मामेव प्रत्यर्पणं करणीय’मिति प्रस्तावोऽपि पुरस्कृतः। वरुणप्रसादाद् राज्ञः पुत्रः सञ्जातः। राजा हरिश्चन्द्रः तस्य नाम ‘रोहितः’ इति कृतवान्। पुत्रप्राप्त्यानन्तरं वरुणः प्राक्कृतपणानुसारं हरिश्चन्द्रं यागार्थं स्मारयति परं पुत्रप्रेम्णा हरिश्चन्द्रो यागाय पशोर्योग्यतानुसारं प्रथमतः जन्मतः दशदिनानन्तरं, ततो दन्तोत्पत्तेरनन्तरम्, ततो नवीनदन्तोत्पत्तेरनन्तरम्, ततः क्षत्रियसंस्कारानन्तरं पुत्रं यजामीति विलम्बयति । हरिश्चन्द्रप्रस्तावं वरुणो वारं वारं ‘तथास्तु’ इत्युक्त्वा स्वीकरोति । यदा पुत्रो रोहितो धनुर्धारणसमर्थः क्षत्रियसंस्कारयुतो जायते तदा स्वपितुः ‘मया त्वं वरुणाय दातव्यः’ इति कथनमस्वीकृत्य धनुर्धृत्वा वनं गच्छति । यज्ञार्थमागतेन वरुणेन शप्तो राजा हरिश्चन्द्र उदररोगेणाक्रान्तः सन् पीडितो जायते । एकवर्षानन्तरं पित्रुः वृत्तान्तमिदं श्रुत्वा गृहं प्रत्यावर्तनाय वनात् प्रस्थितं रोहितं ब्राह्मणरूपेण समागतः इन्द्रः पर्यटनस्य महत्त्वं प्रतिपाद्य गृहगमनान्निवारयत् । प्रतिवर्षान्ते इन्द्रेण निवारितो रोहितः षड्‌वर्षपर्यन्तं वनेऽवसत् । षष्ठे वर्षे तत्र वने अजीगर्ताख्यस्य कस्यचिद् ब्राह्मणस्य मध्यमपुत्रं शतगोपरिमितेेन मूल्येन क्रीत्वा गृहं प्रतिनिवृत्तः । तस्य मध्यमपुत्रस्य नाम शुनःशेप आसीत् । पुत्रस्य स्थाने शुनःशेपस्यार्पणं स्वीकारयोग्यमिति वरुणानुमोदनपश्चाद् पशुस्थाने शुनःशेपं निश्चित्य राजा यज्ञमारेभे। परं तत्र यज्ञे शुनःशेपं यूपे नियोक्ता अन्यः कोऽपि क्रूरः पुरुषो न प्राप्यते । तदा स्वेच्छानुसारं शतपरिमाणं गोधनं प्राप्य अजिगर्त एव स्वपुत्रं शुनःशेपं यूपे बध्नाति । अनन्तरं यज्ञविध्यनुसारं यूपपूजादिकं विधाय शुनःशेपस्य बलिसमयः समागतः । शुनःशेपबलिकार्येऽपि क्रूरजनो नाप्यते । पुनः स्वेच्छानुसारम् अजिगर्त एव शतगोपरिमितेन मूल्येन पुत्रवधकार्ये नियोज्यते । अनन्तरं शुनःशेपो ‘धनलुब्धः पिता मामवश्यमेव हन्ती’ति विचार्य प्रजापतेर्मन्त्रं जपति । प्रजापतिरग्निं स्तौतुं प्रेरयति । अग्निरपि सवितारं स्तौतुं कथयति । सविता च वरुण एव अस्य वृत्तान्तस्य मूलमिति विज्ञाप्य तमेव ध्यातुं वदति । वरुणोऽपि अग्निम्, अग्निश्च विश्वेदेवान्, ते च अश्विनीकुमारौ, तावपि उषादेवीं प्रार्थयितुं सूचयतः । यदा उषादेव्या मन्त्रं शुनःशेपः पठन्नासीत्, तदैव स बन्धनमुक्तो जातः । राजा हरिश्चन्द्रश्च उदररोगाद् मुक्तोऽभवत् । अनन्तरं ‘शुनःशेपः कस्य पुत्रो भवेत् ?’ इति यज्ञस्था ब्राह्मणाः अनिर्णीता जाताः । तैः सर्वैः ‘शुनःशेपो यं वृणोति स एवास्य पिता’ इति सिद्धान्तमुपस्थापयन्ति । एतद् वचनं श्रुत्वा शुनःशेपो विश्वामित्रस्याङ्कमाश्रयति । तदा अजीगर्तः प्राक्‌कृतं  स्वकीयं पापकर्म प्रति खेदं मन्यमानः शुनःशेपं पुत्रत्वेन स्वसमीपे समागन्तुमाह्वयति परं शुनःशेपस्तद् वचनमस्वीकरोति । विश्वामित्रोऽपि देवप्रसादेन प्राप्तं पुत्रं स्वीकृत्य तस्य नाम ‘देवरातः’ इति करोति । स च मम पुत्रेषु ज्येष्ठः श्रेष्ठश्च भवतीत्यपि समुद्‌घोषयति । तदा शुनःशेपः स्वकीयानां पुत्राणां विचारं ज्ञातुं विश्वामित्रं प्रेरयति । विश्वामित्रश्च शतपुत्रान् तत्राहूय शुनःशेपं ज्येष्ठत्वेन श्रेष्ठत्वेन च स्वीकर्तुं आज्ञापयति परं मधुच्छन्दसो ज्येष्ठाः पञ्चाशत् पुत्राः शुनःशेपं ज्येष्ठं भ्रातरं मन्तुमस्वीकुर्वन्ति । तान् विश्वामित्रः कुलकलङ्कत्वेन नीचजातौ पतनरूपं शापं ददाति । अन्यान् आज्ञापालकान् मधुच्छन्दसः कनिष्ठान् पुत्रान् वरेण उपकरोति । तदनन्तरं शुनःशेपो विश्वामित्रस्य पञ्चाशत् पुत्रैः सह विश्वामित्रकुलावतंसरूपेण सम्मानितो भवति ।  इति शुनःशेपाख्यानस्य सारो विद्यते ।

\customparagraph{\textbf{वैशिष्ट्यम्}}

लोके हरिश्चन्द्रोपाख्यानं प्रसिद्धमस्ति । पुराणादिषु हरिश्चन्द्रस्य करुणरसप्रवाहिणी कथा प्रचलिता विद्यते । लौकिकसाहित्येऽपि हरिश्चन्द्रस्योपाख्यानमाधृत्य श्रव्यदृश्यकाव्यानि प्रणीतानि विद्यन्ते । लोके मध्यमपुत्रस्य महत्त्वहीनत्वबोधार्थं शुनःशेपस्याख्यानं लोककथारूपेणापि श्रुतिपरम्परायां प्रचलितं विद्यते । एतेषां सर्वेषां मुख्यस्रोतस्तु ऋग्वेदान्तर्गतम् ऐतरेयब्राह्मणमस्ति । तत्र राज्ञो हरिश्चन्द्रस्य वृत्तान्तं विद्यते । तस्यैवान्तरे शुनःशेपाख्यानं संश्लिष्टं प्राप्यते । हरिश्चन्द्रो वरुणस्य प्रसादात् पुत्रं तु प्राप्नोति परं तस्य प्रस्तावानुसारं कर्तुमस्वीकृत्य पुत्रो रोहितो वनं गच्छति । षड्‌वर्षानन्तरं शतसङ्‌ख्यकेन गोमूल्येन अजीगर्तस्य मध्यमपुत्रं शुनःशेपं क्रीणाति । तदनन्तरं सञ्जातमाख्यानं शुनःशेपाख्यानरूपेण प्रसिद्धं विद्यते ।

अस्माकं लौकिकसमाजे साहित्ये च उपाख्यानस्यास्य महत्त्वं सन्दृश्यते । एतद् उपाख्यानं कथावस्तुत्वेनाङ्‌गीकृत्य प्रणितानि करुणरसप्रसृतानि काव्यानि भावकहृदयाकर्षकानि सन्ति । आख्यानेनानेन धनस्य लोभेन जनकोऽपि वात्सल्यं विस्मृत्य कीदृशं निकृष्टं कार्यं कर्तुमुद्गच्छतीति प्रदर्श्य एतादृशस्य वात्सल्यरहितस्य पितुरवस्था गर्हणीया भवतीत्यपि संबोध्य लोकान् ‘रामादिवत् प्रवर्त्तितव्यं न रावणादिवत्’ इति काव्यप्रयोजनरूपा शिक्षापि प्रकाश्यते । एवमेव नैरन्तर्येण कृता भगवद्भक्तिरसामान्यां बाधामपि दूरीकरोतीति आध्यात्मिकं विश्वासमपि दृढीकरोति ।

एतादृशीं भौतिकीमाध्यात्मिकीञ्च शिक्षां प्रकाशयन् ऐतरेयब्राह्मणाख्ये वैदिकग्रन्थे विद्यमानं शुनःशेपोपाख्यानं लोके साहित्यसमाजे च स्वकीयं सुश्रव्यं प्रेरकञ्च वैशिष्ट्यं प्रकाशयति।

\specialupakhanda{\textbf{प्रथमः खण्डः}}

\begin{visheshshloka}[center]

\textbf{हरिश्चन्द्रो ह वैधस ऐक्ष्वाको राजाऽपुत्र आस, तस्य ह शतं जाया बभुवुस्तासु पुत्रं न लेभे, तस्य ह पर्वतनारदौ गृह उषतुः, स ह नारदं प्रपच्छ ॥}

\end{visheshshloka}

\customparagraph{संस्कृतव्याख्या}

इक्ष्वाकुवंशीयस्य वैधसराज्ञः पुत्रो हरिश्चन्द्रोऽपुत्र आसीत् । स शतसङ्‌ख्यानां पत्नीनां मध्ये कस्याञ्चिदपि पुत्रं न लेभे । तस्य राज्ञो गृहे पर्वतनारदनामकौ मुनी निवासं चक्रतुः । तयोर्मध्ये नारदं मुनिम् एतया कथ्यमानया गाथया प्रपच्छ ।

\customparagraph{\textbf{नेपाली}}

इक्ष्वाकु वंशमा वेधस नाम गरेका राजाका छोरा हरिश्चन्द्र निःसन्तान थिए । तिनका सयवटी रानीमध्ये कुनैबाट पनि सन्तान भएनन् । ती राजाका घरमा पर्वत र नारद नाम गरेका मुनिद्वयले निवास गरे । तीमध्ये नारदलाई राजाले निम्न गाथाका माध्यमले सोधे-

तां गाथां दर्शयति-

\begin{visheshshloka}[center]

\textbf{यं न्विमं पुत्रमिच्छन्ति, ये विजानन्ति, ये च न ।}

\textbf{किंस्विद् पुत्रेण विन्दते, तन्म आचक्ष्व नारदेति ॥}

\end{visheshshloka}

\customparagraph{संस्कृतव्याख्या}

ये= विवेकयुता देवमनुष्यादयो विजानन्ति, ये=विवेकरहिता तिर्यग्‌जातयो न जानन्ति । ते= सर्वे इमं पुत्रमिच्छन्ति । तेन पुत्रेण किंस्वित्= किं नाम फलम्, विन्दते=  (पिता) लभतेे । हे नारद ! तत्= कारणम्= मे= माम् आचक्ष्व= वद इति ।

\customparagraph{\textbf{नेपाली}}

हे नारद ! देवता, मनुष्य आदि विवेकशील प्राणी र पशु, पक्षी आदि विवेकहीन सबै प्राणी पुत्र प्राप्तिको अभिलाषा राख्दछन् । सन्तानबाट के फाइदा हुन्छ ? पिताले पुत्रबाट के प्राप्ति गर्दछ? त्यसको फल मलाई भन्नुहोस् ।

तस्योत्तरमवतारयति-

\begin{visheshshloka}[center]

\textbf{स एकया पृष्टो दशभिः प्रत्युवाच ॥ इति ।}

\end{visheshshloka}

\customparagraph{\textbf{नेपाली}}

ती राजाले एउटा गाथाद्वारा सोधिएको प्रश्नको नारदले दश गाथाबाट उत्तर दिए ।

तत्र प्रथमामाह-

\begin{visheshshloka}[center]

\textbf{ऋणमस्मिन् संनयत्यमृतत्त्वं च गच्छति ।}

\textbf{पिता पुत्रस्य जातस्य पश्येच्चेज्जीवतो मुखम् ॥ इति ।}

\end{visheshshloka}

\customparagraph{संस्कृतव्याख्या}

पिता स्वकीयांशतः समुत्पन्नस्य, सुखेन जीवतः पुत्रस्य मुखं यदि पश्येत् तर्हि तस्मिन् पुत्रे स्वकीयं लौकिकं वैदिकञ्च ऋणं सम्यगवस्थापयति । अर्थात् स्मृतिप्रमाणानुसारम् ऋणमुक्तो भवति । पुत्रे ऋणभारादिकं न्यस्य स पिता अमृतत्त्वं प्राप्नोति ।

\customparagraph{\textbf{नेपाली}}

पिताले आफ्नो अंशबाट जन्मिएर सुखपूर्वक जीवनयापन गरिरहेको छोराको मुख हेर्ने सौभाग्य पाउँछ भने त्यसले लौकिक र वैदिक ऋण त्यही छोरामा स्थापित गर्दछ। अर्थात् आफ्नो ऋणबाट मुक्त हुन्छ । यसरी पुत्रवान् पिताले अमृतत्त्वलाई प्राप्त गर्दछ ।

\begin{visheshshloka}[center]

\textbf{यावन्तः पृथिव्यां भोगा यावन्तो जातवेदसि ।}

\textbf{यावन्तो अप्सु प्राणीनां भूयान्पुत्रे पितुस्ततः ॥ इति ।}

\end{visheshshloka}

\customparagraph{संस्कृतव्याख्या}

पृथिव्यां भोज्यवासादयो ये भोगाः सन्ति । ये अग्नौ दहनपाचनादयो भोगाः सन्ति । ये स्नानपानादयोऽप्सु भोगाः सन्ति । प्राणिनामेते यावन्तो भोगाः सन्ति । ततः= तावद्भ्यः सर्वेभ्यो भोगेभ्यः, भूयान्‌‍‍‍= समधिको भोगः पितुः पुत्रे विद्यते । (अत्यन्तसुखहेतुत्वात्) तथा चाहुः- ‘पुत्रोत्पत्तिविपत्तिभ्यां नापरं सुखदुःखयोः’ इति।

\customparagraph{\textbf{नेपाली}}

पृथिवीमा प्राणीहरूका लागि खानेकुरा वासस्थान आदि जति पनि भोगका पदार्थ छन्, त्यस्तै जति भोग अग्निमा दहन र पाचन क्रियामा छन्, जति भोग जलमा स्नान र पान आदिमा छन् । तीभन्दा धेरै भोग पिताका लागि पुत्रमा हुन्छ । किनकि पुत्रका माध्यमबाट असामान्य सुख पाइन्छ ।

तृतीयां गाथामाह-

\begin{visheshshloka}[center]

\textbf{शश्वत् पुत्रेण पितरोऽत्यायन् बहुलं तमः ।}

\textbf{आत्मा हि जज्ञ आत्मानः स इरावत्यतितारिणी ॥ इति ।}

\end{visheshshloka}

\customparagraph{संस्कृतव्याख्या}

पितरः‍= जनकाः स्वांशेन समुत्पन्नेन पुत्रेण शश्वत्= सर्वदा, सर्वत्र इहलोके परलोके च इत्यर्थः । बहुलम्= अभ्यधिकम्, तमः= ऐहिकम्, आमुष्मिकञ्च दुःखम् (नारकीयं दुःखम्) अत्यायन्= अतिक्रामन्ति । तथा च बौधायन आह-

\begin{shloka}[center]

पुदिति नरकस्याऽऽख्या दुःखं च नरकं विदुः ।

पुत्तस्त्राणात्ततः पुत्रमिहेच्छन्ति परत्र च ॥ इति ।

\end{shloka}

किञ्च ‘हि’ यस्मात्कारणात् पितुः पुत्र उत्पन्न इत्युक्ते सति पिता स्वस्मात् स्वयमेव जज्ञे= उत्पन्न इत्युक्तं भवति । ततः पिता स्वयं स्वकीयं दुःखं विनाशयति, तथा पुत्रोऽप्येतद्‍दुःखं विनाशयतीति द्रष्टव्यम् । तस्मात् सः= पुत्र इरावती= अन्नयुता अतितारिणी= नदीसमुद्रादेरतितारणहेतुनौरिति शेषः । यथा- नौर्दुर्घटं नद्यादिकं तारयति एवं पुत्रोऽप्यैहिकमामुष्मिकञ्च दुःखं तारयतीत्यर्थः ।

\customparagraph{\textbf{नेपाली}}

पुत्रवान् माता र पिता पुत्रद्वारा इहलोक र परलोक दुबैमा भोग्नु पर्ने अत्यन्त दुःखरूपी नरकलाई पार गर्दछन् किनकि पिता स्वयं आफूलाई आफ्नो पुत्रका रूपमा उत्पन्न गर्दछ । त्यसैले त्यो पुत्र नदी समुद्र आदि अप्ठ्यारा ठाउँ पार गर्न अन्नले भरिएको नाउझैं हुन्छ । अर्थात् लामो नदी वा समुद्र पार तर्नका लागि धेरै दिन लाग्न सक्छ त्यसैले खाद्यसामग्री सहितको नाउ आवश्यक हुन्छ त्यस्तैै भवरूपी सागरको दुःखबाट पार पाउनका लागि पुत्र अन्नसहितको नौकाजस्तै हुन्छ ।

चतुर्थीं गाथामाह-

\begin{visheshshloka}[center]

\textbf{किं नु मलं किमजिनं किमु श्मश्रुणि किं तपः ।}

\textbf{पुत्रं ब्राह्मण इच्छध्वं स वै लोकोऽवदावदः॥ इति ।}

\end{visheshshloka}

\customparagraph{संस्कृतव्याख्या}

अत्र मलाजिनश्मश्रुतपःशब्दैराश्रमचतुष्टयं विवक्षितम् । मलरूपाभ्यां शुक्रशोणिताभ्यां संयोगात् ‘मल’ शब्देन गार्हस्थ्यं विवक्षितम् । कृष्णाजिनसंयोगादजिनशब्देन ब्रह्मचर्यं विवक्षितम्। क्षौरकर्मराहित्याच्छ्‌मश्रुशब्देन वानप्रन्थं विवक्षितम् । इन्द्रियनियमसद्भावात् तपःशब्देन पारिव्राज्यं विवक्षितम् । मलम्= गार्हस्थ्यम्, किं नु= किं नाम सुखं करिष्यति ? न किञ्चिदित्यर्थः । एवमुत्तरत्रापि योज्यम् । हे ब्राह्मणाः= विप्राः, विप्रक्षत्रियाद्याः सर्वे यूयं सुखहेतुत्वात् पुत्रमिच्छध्वम् । स वै= स एव पुत्रः, अवदावदः= वदितुमयोग्यानि निन्दावाक्यान्यवदाः तैर्वाक्यैर्नोद्यते न कथ्यते इति अवदावदः । एवं प्रघट्टेन तेन कथ्यते इति- अवदावदः= दोषराहित्यान्निन्दानर्ह इत्यर्थः । तादृशो लोकः= भोगहेतुः पुत्रः । तस्मादाश्रमेभ्योऽप्याधिक्येन पुत्रेच्छा कर्तव्या । यद्यपि हरिश्चन्द्र एवात्र प्रष्टा, तथाऽपि तेन सह ऋषीणां बहूनां सभायामवस्थानात् ‘ब्राह्मणः’  इति सम्बोधनम् ।

\customparagraph{\textbf{नेपाली}}

(यहाँ मल, अजिन, श्मश्रु, तप जस्ता शब्दले चार आश्रमलाई बुझाएको छ)  शुक्र र रगतको सम्मिश्रण मल रूपको गृहर्थाश्रमले के सुख दिन सक्छ ? कालो मृगचर्म अर्थात् ब्रह्मचर्याश्रमले के सुख दिन्छ ? क्षौरकर्मरहित दारी र जुँगावाला वानप्रस्थ आश्रमले के सुख दिन्छ ? इन्द्रियनियमनरूप तपवाला सन्न्यासाश्रमले के सुख दिन्छ ? अर्थात् ती सबै आश्रमबाट केही लाभ हुँदैन । त्यसैले हे ब्राह्मणलगायत क्षत्रियादिहरू हो ! तपाईंहरू सबै सुखको हेतुभूत पुत्रको इच्छा गर्नुहोस् । किनकि त्यो पुत्र दोषरहित हुनाले अनिन्दनीय लोक अर्थात् सुखभोगको कारण हो । (यहाँ प्रश्न सोध्ने हरिश्चन्द्र एक्लै भए पनि ब्राह्मणहरू भनेर सभामा रहेका अन्य ब्राह्मणहरूलाई पनि सङ्‌केत गरिएको हो)

पञ्चमीं गाथामाह-

\begin{visheshshloka}[center]

\textbf{अन्नं ह प्राणः शरणं ह वासो रूपं हिरण्यं पशवो विवाहाः ।}

\textbf{सखा ह जाया कृपणं ह दुहिता ज्योतिर्ह पुत्र परमे व्योमन्॥}

\end{visheshshloka}

\customparagraph{संस्कृतव्याख्या}

अन्नादयो लोके सुखहेतुत्वेन प्रसिद्धाः । तथा हि शरीरे प्राणावस्थितिहेतुत्वादन्नमेव प्राणः । वासः= वस्त्रम्, शीतोपद्रवाद् रक्षकत्वेन शरणम्= गृहसमानम् । हिरण्यम्= कर्णाभरणादिकम्, दृष्टिप्रियत्वात्, रूपम्= रूपसम्पादकम् । पशवः= गवाश्वादयः, विवाहाः= विशेषेण निर्वाहकाः । जाया= पत्नी सुखभोगे सहकारित्वात् सखा ह= सखिस्वरूपा । एवमेते सुखहेतुत्वेन प्रसिद्धा अपि तात्कालिकमल्पमेव सुखं प्रयच्छन्ति । दुहिता ह= पुत्री कृपणम्= केवलदुःखकारित्वाद् दैन्यहेतुः । तथा च स्मर्यते-

संभवे स्वजनदुःखकारिता संप्रदानसमयेऽर्थहारिका ।

यौवनेऽपि बहुदोषकारिका दारिका हृदयदारिका पितुः॥

पुत्रो ह= पुत्रस्तु ज्योतिःस्वरूपम्= प्रकाशस्वरूपम्, तमोनिवारकत्वेन प्रकाशरूपत्वेन स हि पितरम्, परमे व्योमन्= उत्कृष्टे आकाशे परब्रह्मस्वरूपेऽवस्थापयति । (‘आकाशस्तल्लिङ्‍गात्’ इत्यनेन व्याससूत्रेण आकाशव्योमादिशब्दानां ब्रह्मपरत्वं निर्णीतम् । पुत्रस्य च ब्रह्मज्ञानहेतुत्वं पूर्वमेवामृतत्त्वं च गच्छतीत्यत्र प्रतिपादितम्)

\customparagraph{\textbf{नेपाली}}

शरीरमा प्राण रक्षा गर्ने हुनाले अन्नादिलाई प्राण भनिन्छ । जाडो आदि उपद्रवलाई छल्ने हुनाले कपडालाई घरसमान मानिन्छ। सुन आदिका गहना सौन्दर्य बढाउने हुनाले रूपसमान मानिन्छन् । गाईवस्तु विशेष रूपले वहन गर्ने हुँदा निर्वाहक हुन् । पत्नी सुखभोगमा सहायक हुनाले साथी हुन् । यी सबै सुखका कारणका रूपमा प्रसिद्ध छन् । यिनले तात्कालिक मात्र सुख प्रदान गर्दछन् । पुत्री केवल दुःखदायक भएकाले दैन्य स्वरूप हुन्छन्, तर पुत्र अन्धकारको निवारक ज्योतिस्वरूप हुन्छ जसले पितालाई परब्रह्मस्वरूप उत्कृष्ट आकाशमा स्थापित गराउँछ ।

षष्ठीं गाथामाह-

\begin{visheshshloka}[center]

\textbf{पतिर्जायां प्रविशति, गर्भो भूत्वा स मातरम् ।}

\textbf{तस्यां पुनर्नवो भूत्वा दशमे मासि जायते ॥ इति ।}

\end{visheshshloka}

\customparagraph{संस्कृतव्याख्या}

पत्युराकारद्वयमस्ति- वर्तमानपुरुषाकार एकः, रेतोरूपेण गर्भाकारो द्वितीयः । जायाया अप्याकारद्वयमस्ति-  पतिरूपमाकारं प्रति जाया भवति । गर्भारूपमाकारं प्रति माता भवति । अतः स= तादृशः पतिः स्वयं रेतोरूपेण गर्भो भूत्वा पूर्वमवस्थितां जायां भविष्यदाकारेण मातरं सतीं प्रविशति । तस्याम्= मातरि पुनर्नवो भूत्वा= पूर्वमन्यस्यां मातर्युत्पन्नो जरठः । इदानीं पुनर्नूतनबालो भूत्वा तस्याम्= इदानींतन्यामस्यां मातरि गर्भपाके सति दशमे मास्युत्पद्यते । तस्मात् पुत्रः स्वस्माद् अन्यो न भवति ।

\customparagraph{\textbf{नेपाली}}

पति नै रेत (वीर्य) रूपमा पत्नीमा प्रवेश गर्दछ, त्यो पति स्वयं वीर्यका माध्यमले गर्भरूप भएर भविष्यमा आमा बन्ने आफ्नै पत्नीबाट गर्भकाल पूृण भएपछि पुत्रका रूपमा दसौँ महिनामा उत्पन्न हुन्छ । त्यसैले पुत्र आफूभन्दा भिन्न होइन ।

सप्तमीं गाथामाह-

\begin{visheshshloka}[center]

\textbf{तज्जाया जाया भवति यदस्यां जायते पुनः ।}

\textbf{आभूतिरेषा भूतिर्बीजमेतन्निधीयते ॥ इति ।}

\end{visheshshloka}

\customparagraph{संस्कृतव्याख्या}

यत्= यस्मात् कारणात् अस्याम्= गर्भधारिण्याम् अयम्= पिता पुत्ररूपेण पुनर्जायते। तत्= तस्मात् कारणात् लोकप्रिसद्धा या जायाऽस्ति, सा जायतेऽस्यामिति व्युत्पत्त्या जायाशब्दवाच्या भवति । किं चैषा भूत्याभूतिशब्दाभ्यामभिधीयते । भवत्यस्यां पुत्ररूपेण पतिरित्येषा ‘भूति’ शब्दवाच्या, रेतोरूपेणाऽऽगत्यास्यां पुत्ररूपेण भवति ‘आभूति’ शब्दवाच्या । एतत्= एतस्यां स्त्रियाम्, बीजम्= रेतोरूपम्, निधीयते= प्रक्षिप्यते । तस्मादुक्ताः शब्दा उपपद्यन्ते ।

\customparagraph{\textbf{नेपाली}}

गर्भधारण गर्ने पत्नीमा पिता पुत्र रूपले उत्पन्न हुन्छन् । त्यसैले ती पत्नी ‘जायते अस्याम्’ यस व्युत्पत्तिअनुसार ‘जाया’ शब्दवाच्य हुन्छिन् । ‘भवति अस्यां पुत्ररूपेण पतिः’ यो व्युत्पत्तिअनुसार पत्नी ‘भूतिः’ शब्दवाच्य हुन्छेे । पति नै वीर्यरूपमा आएर पत्नीबाट पुत्रका रूपमा उत्पन्न हुने भएकाले ‘आगत्य भवतीति’ यस व्युत्पत्ति अनुसार पत्नीलाई ‘आभूति’ पनि भनिन्छ । पत्नीमा पुत्रका रूपमा उत्पन्न हुनका लागि वीर्यको स्थापना गरिने हुनाले उक्त ‘भूति’ र ‘आभूति’ शब्द पत्नी (जाया) अर्थमा निष्पन्न हुन्छन् ।

अष्टमीं गाथामाह-

\begin{visheshshloka}[center]

\textbf{देवाश्चैतामृषयश्च तेजः समभरन् महत् ।}

\textbf{देवा मनुष्यानब्रुवन्नेषा वो जननी पुनः ॥ इति ।}

\end{visheshshloka}

\customparagraph{संस्कृतव्याख्या}

एताम्= एतस्यां योषिति देवाश्च महर्षयश्च स्वकीयं महत्तेजः= रेतोरूपं सारम्, समभरन्= पुत्रोत्पादनाय सम्पादितवन्तः। स्वयमेव संपाद्य ततो मनुष्यानित्यब्रुवन्- हे मनुष्याः ! या इयं स्त्री इदानीं जायारूपेण वर्तते, सेयं पुनः वः=युष्माकम्, पुत्ररूपे जन्मनि जननी भवति ॥

\customparagraph{\textbf{नेपाली}}

यी स्त्रीमा देवता र ऋषिहरूले आफ्नो महान् रेतरूप सारलाई पुत्रोत्पादनका लागि नै स्थापित गरे । त्यसपछि मानवलाई भने- हे मानव हो ! जो यी स्त्री अहिले तिमीहरूकी पत्नी छन्, तिनै पत्नी पुत्ररूपमा उत्पन्न भएका तिमीहरूकी आमा हुनेछिन् ।

नवमीं गाथामाह-

\begin{visheshshloka}[center]

\textbf{नापुत्रस्य लोकोऽस्तीति तत्सर्वे पशवो विदुः ।}

\textbf{तस्मात्तु पुत्रो मातरं स्वसारं चाधिरोहति ॥ इति ।}

\end{visheshshloka}

\customparagraph{संस्कृतव्याख्या}

लोकः= लोकजन्यं सुखम्, अपुत्रस्य नास्ति । न हि पुत्रदर्शनेन यत्सुखं लभ्यते, तदन्यदर्शनेन क्वचिदपि दृश्यते । इति= यदस्ति तत् सर्वे पशवः= गोमहिष्यादयोऽपि जानन्ति यस्माद्, तस्मादेव कारणात्पशुजातौ जातः पुत्रः= वत्सः स्वकीयां मातरं भगिनीं वा पुत्रोत्पादनार्थमधिरोहति ।

\customparagraph{\textbf{नेपाली}}

जसको पुत्र (सन्तति) हुँदैनन् उसले इहलोकजन्य कुनै प्रकारको सुख पाउँदैन, किनकि पुत्रदर्शनजनित जुन सुख छ त्यो कतैबाट पनि पाइँदैन । यो कुरो सबै पशुले पनि बुझेका छन् । त्यसैले पशुयोनिमा उत्पन्न भएका प्राणीले पुत्र उत्पादनका लागि आफ्नै आमा वा बहिनीमा पनि रेतको आधान गर्दछन् ।

दशमीं गाथामाह-

\begin{visheshshloka}[center]

\textbf{एष पन्था उरुगायः सुशेवो यं पुत्रिण आक्रमन्ते विशोकाः।}

\textbf{तं पश्यन्ति पशवो वयांसि च तस्मात्ते मात्राऽपि मिथुनीभवन्ति॥}

\end{visheshshloka}

\customparagraph{संस्कृतव्याख्या}

पुत्रिणः= पुत्रवन्तो देवमनुष्यादयः, विशोकाः=शोकरहिताः सन्तः यम्= सुखानुभवरूपं पन्थानम्, आक्रमन्ते= प्राप्नुवन्ति । एष पन्थाः= पुत्रसुखानुभवरूपो मार्गः, उरुगायः= उरुभिर्महद्भिः शास्त्रज्ञै राजामात्यादिभिश्च गीयते । तथा सुशेवः=सुष्ठु सेवितं योग्यः, सुखाधिक्यस्य विद्यमानत्वात् । तम्= पुत्रसुखानुभवरूपं पन्थानम्, पशवः= गवादयः, वयांसि= पक्षिणः, पश्यन्ति= जानन्ति । तस्मात् ते= पशुपक्ष्यादयः पुत्रसुखार्थं मात्रा= जनन्या सह मिथुनीभवन्ति, किं किमुतान्यया स्त्रिया सहेत्यर्थः ।

\customparagraph{\textbf{नेपाली}}

देव मनुष्यादि पुत्रवान् व्यक्तिले शोकरहित भई जुन सुखको प्राप्तिरूप मार्गलाई सेवन गर्दछन्, त्यो ठूला व्यक्ति (राजा, मन्त्री, शास्त्रज्ञ आदि) द्वारा पनि प्रशंसा गरिएको मार्ग (सन्तानजन्य सुख) सबैद्वारा सेवनीय छ । पुत्रजन्य सुख पशुपक्षी आदि तिर्यक् जातिले पनि राम्ररी बुझेका हुन्छन् । त्यसैले त्यो सुख पाउनका लागि तिनीहरू माताका साथ पनि मैथुन गर्न उद्यत हुन्छन् भने अन्य स्त्रीका साथमा  गरिने मैथुनको त के कुरा गर्नु ।

नारदेनोक्तां दशगाथां निगमयति-

\begin{visheshshloka}[center]

\textbf{इति हास्मा आख्याय ॥१३॥}

\end{visheshshloka}

\customparagraph{संस्कृतव्याख्या}

इति ह= अनेनैव प्रकारेण अस्मै= तस्मै हरिश्चन्द्राय आख्याय= उत्तरमभिधाय अवस्थित इति शेषः ।

\customparagraph{\textbf{नेपाली}}

यस प्रकारले राजा हरिश्चन्द्रले राखेको जिज्ञासाको उत्तर दिएर मुनि नारदले विश्राम लिए ।

\begin{visheshshlokatwo}[alternating]{88}

\textbf{ऐतेरियब्राह्मणको  तेत्तिसौं अध्याय (सप्तपञ्चिकाको तेस्रो अध्यायको प्रथम खण्ड)}

\textbf{नेपाली भावानुवाद पूर्ण भयो ।}

\end{visheshshlokatwo}

\specialupakhanda{\textbf{द्वितीयः खण्डः}}

पुत्रनिमित्तलाभप्रतिपादिका नारदोक्ता दश गाथा उदाहृत्य पुत्रजन्मोपायोपदेशपरं नारवाक्यमवतारयति-

\begin{visheshshloka}[center]

\textbf{अथैनमुवाच- वरुणं राजानमुपधाव, पुत्रो मे जायतां तेन त्वा यजा इति ।}

\end{visheshshloka}

\customparagraph{संस्कृतव्याख्या}

अथ= पुत्रेच्छानिमित्तकथनानन्तरम्, एनम्= पुत्रार्थिनं हरिश्चन्द्रम्, नारद उवाच- हे हरिश्चन्द्र ! वरुणं राजानं उपधाव= प्रार्थयस्व । येन प्रकारेण प्रार्थनीयः सोऽभिधीयते । हे वरुण ! त्वत्प्रसादात् मे पुत्रो जायताम्, ततस्तेन= पुत्रेण त्वा यजा= त्वामुद्दिश्य यज्ञं करवाणीति ।

\customparagraph{\textbf{नेपाली}}

यो पुत्रको महिमा बताएपछि पुत्रार्थी हरिश्चन्द्रलाई नारदले भने- हे हरिश्चन्द्र ! तिमी वरुणसँग गएर पुत्रका लागि प्रार्थना गर कि हे वरुणदेव ! तपाईंको कृपाबाट मेरो छोरो भए त्यसैले तपाईंको यज्ञ गर्नेछु ।

प्रार्थनाप्रकारं नारदोपदेशमुक्त्वा हरिश्चन्द्रवृत्तान्तं दर्शयति-

\begin{visheshshloka}[center]

\textbf{तथेति; स वरुणं राजानमुपससार, पुत्रो मे जायतां तेन त्वा यजा इति;  तथेति; तस्य ह पुत्रो जज्ञे रोहितो नाम ॥ इति ।}

\end{visheshshloka}

\customparagraph{संस्कृतव्याख्या}

नारदोपदेशमङ्गीकृत्य हरिश्चन्द्रो वरुणम् उपससार= प्रार्थयामास। पुत्रो मे जायतां तेन त्वा यजा इति; इत्युक्तवान् । सः= वरुणोऽपि तथाऽस्त्विति तदीयपुत्रोत्पत्त्यै वरं दत्तवान् । तेन च वरेणोत्पन्नस्य पुत्रस्य ‘रोहितः’ इत्येतन्नामाभूत् ।

\customparagraph{\textbf{नेपाली}}

राजाले भने- “ठिक छ, म त्यसै गर्छु ।” नारदको उपदेशले प्रेरित भएर ती हरिश्चन्द्रले राजा वरुणसँग प्रार्थना गरे- “हे वरुणदेव ! मेरो पुत्र भयो त्यसैका माध्यमले तपाईंको यज्ञ गर्नेछु।” राजा वरुणले वर दिँदै भने- “ठिक छ, त्यस्तै हुनेछ ।” वरुणको यस्तो वरको प्रभावले हरिश्चन्द्रलाई ‘रोहित’ नामक पुत्र प्राप्त भयो ।

अथ बहुभिः पर्यायैर्वरुणस्योक्तिर्हरिश्चन्द्रस्य प्रत्युक्तिश्चेत्युभयं प्रतिपाद्यते ।

तत्र प्रथमं पर्यायं दर्शयति-

\begin{visheshshloka}[center]

\textbf{तं होवाचाजनि वै ते पुत्रो यजस्व माऽनेनेति; स होवाच, यदा वै पशुर्निर्दशो भवत्यथ स मेध्यो भवति, निर्दशो न्वस्त्वथ त्वा यजा इति; तथेति ॥ इति ।}

\end{visheshshloka}

\customparagraph{संस्कृतव्याख्या}

तम्= हरिश्चन्द्रम्, वरुण उवाच- हे हरिश्चन्द्र !  ते= तव पुत्रः, अजनि वै= उत्पन्न एव, अनेन पुत्रेण मामुद्दिश्य यागं कुरु’ इति । एवं वरुणनोक्ते हरिश्चन्द्रः पुनः प्रत्युवाच- हे वरुण ! यागार्थं पशुर्यदा निर्दश= दशाधिकदिनस्य भवति, तदा स= पशुः, मेध्य= यागयोग्यो भवति। निर्गतान्याशौचदिनानि दशसङ्‍ख्यकानि यस्मात् पशोः सोऽयं ‘निर्दशः’ । तस्मादयं नु क्षिप्रम्= शीघ्रम्, निर्दशोऽस्तु, अथः= अनन्तरं त्वा= त्वां प्रति अहम् ‘यजै’ इत्येतद्वाक्यम् । वरुणः ‘तथास्तु’ इत्यङ्‍गीचकार ।

\customparagraph{\textbf{नेपाली}}

ती हरिश्चन्द्रलाई वरुणले भने- “हे हरिश्चन्द्र ! तिम्रो पुत्र उत्पन्न भयो । अब तिमी यसले मेरो यज्ञ गर ।”  यसरी वरुणले भनेपछि  हरिश्चन्द्रले भने “जन्मिएको दश दिनपछि (आशौच सकिएपछि) पशु यज्ञका लागि योग्य हुन्छ, मेेरो छोरो पनि दश दिनको भएपछि म यज्ञ गर्नेछु” वरुणले भने- “ठिक छ, त्यसै गर्नू।”

द्वितीयमुक्तिप्रत्युक्तिपर्यायं दर्शयति-

\begin{visheshshloka}[center]

\textbf{स ह निदर्श आस तं होवाच- निदर्शो न्वभूद्यजस्व माऽनेनेति; स होवाच- यदा वै पशोर्दन्ता जायन्तेऽथ स मेध्यो भवति; दन्ता न्वस्य जायन्तामथ त्वा यजा इति; तथेति ॥ इति ।}

\end{visheshshloka}

\customparagraph{संस्कृतव्याख्या}

दशदिनाशौचापगमे शुद्धत्वाद् यथा यागयोग्यत्वम्, तथा दन्तोत्पन्नत्वादवयवसम्पूर्त्या यागयोग्यत्वमित्यभिप्रायः । स्पष्टमन्यत्।

\customparagraph{\textbf{नेपाली}}

जहिले त्यो बालक दश दिनको भयो, तब वरुणले हरिश्चन्द्रलाई भने- “यो बालक दश दिनको भयो, अब यसबाट मेरो यज्ञ गर ।” ती हरिश्चन्द्रले भने- “जब पशुका दाँत उम्रिन्छन् तब त्यो पूर्ण अवयवका कारण यज्ञका लागि योग्य हुन्छ । त्यसैले यस बालकको दाँत जहिले उम्रिन्छन् तब म तपाईंको यज्ञ गर्छु ।” वरुणले भने- “ठिकै छ, त्यसै गर।”

तृतीयं पर्यायं दर्शयति-

\begin{visheshshloka}[center]

\textbf{तस्य ह दन्ता जज्ञिरे; तं होवाचाज्ञत वा अस्य दन्ता यजस्व माऽनेनेति; स होवाच- यदा वै पशोर्दन्ताः पद्यन्तेऽथ स मेध्यो भवति; दन्ता न्वस्य पद्यन्तामथ त्वा यजा इति, तथेति ॥ इति ।}

\end{visheshshloka}

\customparagraph{संस्कृतव्याख्या}

अज्ञत वै= जाता एव । पद्यन्ते= पतन्ति । प्रथमोत्पन्नानां दन्तानामस्थायित्वेन मुख्यपश्ववयवत्वाभावात् तत्पाते सति पशोर्मेध्यत्वम् ॥

\customparagraph{\textbf{नेपाली}}

ती बालक रोहितका दाँत उम्रिएपछि वरुणले हरिश्चन्द्रलाई भने- “अब यस बालकका दाँत निस्किए, अतः मेरो यज्ञ गर ।” ती हरिश्चन्द्रले भने- “जब पशुका पहिलेका दाँत झर्छन् तब त्यो पशु यज्ञका लागि योग्य हुन्छ । त्यसकारण यस बालकका दुधे दाँत जब झर्छन् तब म तपाईंको यज्ञ गर्नेछु ।” वरुणले भने- “ठिकै छ, त्यसै गर ।”

चतुर्थं पर्यायं दर्शयति-

\begin{visheshshloka}[center]

\textbf{तस्य ह दन्ताः पेदिरे; तं होवाचापत्सत वा अस्य दन्ता, यजस्व माऽनेनेति; स होवाच- यदा वै पशोर्दन्ताः पुनर्जायन्तेऽथ स मेध्यो भवति; दन्ता न्वस्य पुनर्जातयन्तामथ त्वा यजा इति, तथेति ॥ इति ।}

\end{visheshshloka}

\customparagraph{संस्कृतव्याख्या}

अनपत्सत वै= पतिताः । पुनरुत्पन्नानां दन्तानां स्थिरत्वेन सम्पूर्णावयवत्वात् पशोर्मेध्यत्वम्॥

\customparagraph{\textbf{नेपाली}}

त्यस बालकका दाँत झरे । ती हरिश्चन्द्रलाई वरुणले भने- “अब यो बालकका दाँत झरे, अतः यसले मेरो यज्ञ गर।” ती हरिश्चन्द्रले भने- “जब पशुका दाँत झरेर फेरि उम्रन्छन् तब अवयव पूर्ण हुने भएकाले त्यो पशु यज्ञका लागि योग्य बन्दछ, त्यसैले यस बालकका पनि जब दाँत पुनः उम्रिन्छन् तब म तपाईंको यज्ञ गर्नेछु ।” वरुणले भने- “ठिक छ, त्यसै गर्नू ।”

पञ्चमं पर्यायं दर्शयति-

\begin{visheshshloka}[center]

\textbf{तस्य ह दन्ताः पुनर्जज्ञिरे; तं होवाचाज्ञत वा अस्य पुनर्दन्ताः, यजस्व माऽनेनेति; स होवाच- यदा वै क्षत्‌त्रियः सान्नाहुको भवत्यथ स मेध्यो भवति; सन्नाहं तु प्राप्नोत्वथ त्वा यजा इति, तथेति ।}

\end{visheshshloka}

\customparagraph{संस्कृतव्याख्या}

अज्ञत वै= जाता एव । पश्वन्तरस्य पुनर्दन्तोत्पत्तिमात्रेण मेध्यत्वेऽप्यस्य पशोः क्षत्त्रियत्वात् स्वजात्युचितधनुर्बाणकवचादिसन्नाहशीलित्वे सति जात्युचितव्यापारसम्पूर्तौ मेध्यत्वम्, तस्मात् ‘नु’ क्षिप्रमेवासौ सन्नाहं प्राप्नोतु, अनन्तरमेव यजा इत्युत्तरं वरुणोऽङ्‍गीचकार ॥

\customparagraph{\textbf{नेपाली}}

बालक रोहितका दाँत झरेर फेरि उम्रिए । त्यस वेला वरुणले हरिश्चन्द्रलाई भने- “अब यस बालकका फेरि दाँत पनि निस्किए, अब मेरो यज्ञ गर ।” हरिश्चन्द्रले भने- “जब क्षत्रिय बालक उपयुक्त धनुष, बाण एवं कवच आदि धारण गरी आफ्नो जातिको उचित कार्य गर्न समर्थ हुन्छ तब यज्ञका लागि योग्य मानिन्छ । अतः यो बालक छिटै शस्त्रादि धारण गर्न समर्थ हुनेछ त्यसपछि म तपाईंको यज्ञ गर्नेछु ।” वरुणले भने- “ठिकै छ ।”

षष्ठं पर्यायं दर्शयति-

\begin{visheshshloka}[center]

\textbf{स ह सन्नाहं प्रापत्, तं होवाच- सन्नाहं नु प्राप्नोद्यजस्व माऽनेनेति, स तथेत्युक्त्वा पुत्रमामन्त्रयामास; ततायं वै मह्यं त्वामददाद्धन्त त्वयाऽहमिमं यजा इति ॥}

\end{visheshshloka}

\customparagraph{संस्कृतव्याख्या}

सन्नाहप्राप्तेरूर्ध्वं सः= हरिश्चन्द्रो वरुणोक्तिमङ्‍गीकृत्य पुत्रमामन्त्र्यैवमुवाच । उपलालनार्थं पुत्रे पितृवाचि ततशब्दप्रयोगः । हे तत ! हे पुत्र ! अयम्=वरुणः, मह्यम्= मां हरिश्चन्द्रम्, त्वाम्= रोहितम्, पुत्रवरेण अददात्= दत्तवान् हन्त !  इति दुष्टोऽहम्, इमम्= वरुणम्, त्वया=पुत्रेण यजा= यागरूपां पूजां करवाणीति हरिश्चन्द्रस्योक्तिः ।

\customparagraph{\textbf{नेपाली}}

अब ती रोहित शस्त्रादि धारण गर्न समर्थ भए, अनि वरुणले हरिश्चन्द्रलाई भने- “अब तिम्रो छोरो शस्त्रधारी भयो, मेरो यज्ञ गर ।” हरिश्चन्द्रले भने- “ठिक छ ।” यति भनेर पुत्र रोहितलाई राजाले बोलाएर भने- “हे पुत्र ! वरुणद्वारा वरका रूपमा तिमी मलाई दिइएका हौ । अतः म तिमीद्वारा वरुणको यज्ञ गर्दछु ।”

अथ रोहितनाम्नः पुत्रस्य कृत्यं दर्शयति-

\begin{visheshshloka}[center]

\textbf{स ह नेत्युक्त्वा धनुरादायारण्यमुपातस्थौ; स संवत्सरमरण्ये चचार ॥१४॥ इति।}

\end{visheshshloka}

\customparagraph{संस्कृतव्याख्या}

स ह= स खलु रोहिताख्यः पुत्रः पितुर्वाक्यं निषिध्य स्वरक्षणार्थं धनुः स्वीकृत्यारण्यं प्रत्युपगतोऽभूत् । कस्मिंश्चिदरण्ये नैरन्तर्येण सः= रोहितः संवत्सरम्= वर्षपर्यन्तम्, चचार ।

\customparagraph{\textbf{नेपाली}}

ती पुत्र रोहितले पिताजीको कुरालाई अस्वीकार गरी आफ्नो रक्षाका निमित्त धनु धारण गरी निरन्तर एक वर्षसम्म कुनै वनमा घुमिरहे ।

\begin{visheshshlokatwo}[center]{84}

\textbf{ऐतेरियब्राह्मणको  तेत्तिसौं अध्याय (सप्तपञ्चिकाको तेस्रो अध्यायको द्वितीय खण्ड)}

\textbf{नेपाली भावानुवाद पूर्ण भयो ।}

\end{visheshshlokatwo}

\specialupakhanda{\textbf{तृतीयः खण्डः}}

रोहितस्यारण्ये संवत्सरचरणानन्तरभाविनं हरिश्चन्द्रादीनां वृत्तान्तं दर्शयति-

\begin{visheshshloka}[center]

\textbf{अथ हैक्ष्वाकं  वरुणो जग्राह, तस्य होदरं जज्ञे तदु ह रोहितः शुश्राव सोऽरण्याद् ग्राममेयाय}

\textbf{तमिन्द्रः पुरुषेणरूपेण पर्येत्योवाच-}

\textbf{नानाश्रान्ताय श्रीरस्तीति रोहित शुश्रुम ।}

\textbf{पापो नृषद्वरो जन इन्द्र इच्चरतः सखा चरैवेति ॥ इति ।}

\end{visheshshloka}

\customparagraph{संस्कृतव्याख्या}

अथ= रोहितस्यारण्ये संवत्सरवासानन्तरमेव ऐक्ष्वाकम्= इक्ष्वाकुवंशोत्पन्नं हरिश्चन्द्रम्, वरुणः= देवः, रोगरूपेण जग्राह । वरुणेन गृहीतस्य हरिश्चन्द्रस्य उदरं जज्ञे= जलेन पूरितमुच्छूनं  महोदरनामकं रोगस्वरूपमुत्पन्नम् । तदु ह= तदपि सर्वं वृत्तान्तमरण्ये स्थितो रोहितः= हरिश्चन्द्रपुत्रो मनुष्यमुखाच्छुश्राव । श्रुत्वा च सः=रोहितः पितरं द्रष्टुमरण्यात् ग्रामं प्रत्याजगाम । आगच्छन्तं रोहितं मार्गमध्य इन्द्रः केनचिद्‍ब्राह्मणपुरुषरूपेण प्राप्य इदमुक्तवान्- आ समन्ताच्छ्रान्तः आश्रान्तः, सर्वत्र पर्यटनेन श्रान्तिं प्राप्तः, तद्विपरीतोऽनाश्रान्तः= एकत्रैव निवासशीलः, तादृशाय तथाविधस्य पुरुषस्य कृते श्रीः= बहुविधा सम्पत्, नास्ति । यद्वा नानेति पदच्छेदः । श्रान्ताय= सर्वत्र पर्यटनेन श्रान्तस्य नाना श्रीः= बहुविधा सम्पदस्ति । इति= अनेन प्रकारेण हे रोहित ! वयं नीतिकुशलानां पुरुषाणां मुखाच्छुश्रुम । वरो जनः= विद्यादिभिः श्रेेष्ठोऽपि पुरुषो नृषत्पापः= नृषु मनुष्येषु सीदतीति नृ‍षत् । श्रेष्ठोऽपि बन्धुगृहेषु सर्वदाऽवस्थितस्तैरवज्ञातः पापः= तुच्छो भवेत् । अतस्तव पितृगृहे वासो न युक्तः । न चारण्ये चारतो मम सहायो नास्तीति शङ्कनीयम् । इन्द्र एव= परमेश्वर एव चरतः= भ्रमतस्तव सखा= मित्रम्, भविष्यति। तस्मात् चरैव= सर्वथाऽरण्ये चरस्वेत्येवमुवाच । एवं बहुष्वपि पर्यायेषु द्रष्टव्यम् ।

\customparagraph{\textbf{नेपाली}}

यसरी पिताजीको कुरालाई अस्वीकार गरी रोहित वन गएपछि वरुणले राजालाई रोग रूपले सङ्‍क्रमित गरे । उनलाई ‘जलोदर’ नामक रोग लाग्यो । यो कुरा जङ्‍गलमा रहेका रोहितले मानिसहरूका मुखबाट सुने । पिताको दुरावस्थाका बारेमा सुनेर पितालाई भेट्न गाउँतर्फ जान लागे । रोहितलाई गाउँ जाँदै गरेको देखेर कुनै पुरुषको रूप धारण गरेका इन्द्रले तिनका नजिक आई भने-

हे रोहित ! मैले नीतिविद् विद्वान्‌का मुखबाट यस्तो कुरा सुनेको छु कि अल्छी नगरी अविश्रान्त रहने अर्थात् सबैतिर घुम्ने मान्छेलाई अनेक प्रकारको सम्पत्ति प्राप्त हुन्छ । कुनै पनि मानिस विद्या आदिले श्रेष्ठ भए पनि गृहस्थ भएर बन्धुबान्धवका साथमा एकै स्थानमा बसिरहन्छ भने त्यो तुच्छ हुन जान्छ, केही प्रगति गर्दैन । अतः तिमी फर्किएर पिताका समीपमा बस्नु ठिक छैन । तिमीले यो सोच्नु पर्दैैन कि जङ्‍गलमा मेरा सहायक कोही छैनन् किनकि यत्रतत्र भ्रमणका क्रममा इन्द्र अर्थात् परमेश्वर नै तिम्रा सहायक हुन्छन् । त्यसैले तिमी वनमै भ्रमण गर ।

तत्रेन्द्ररोहितयोः संवादे प्रथमं पर्यायं दर्शयित्वा द्वितीयं पर्यायं दर्शयति-

\begin{visheshshloka}[center]

\textbf{चरैवेति वै मा ब्राह्मणोऽवोचविति ह द्वितीयं संवत्सरमरण्ये चचार, सोऽरण्याद्}

\textbf{ग्राममेयाय, तमिन्द्रः पुरुषरूपेण पर्येत्योवाच-}

\textbf{पुष्पिण्यौ चरतो जङ्‌घे भूष्णुरात्मा फलग्रहिः ।}

\textbf{शेरेऽस्य सर्वे पाप्मानः श्रमेण प्रपथे हतश्चरैवेति ॥ इति ।}

\end{visheshshloka}

\customparagraph{संस्कृतव्याख्या}

ब्राह्मणरूपस्येन्द्रस्य वाक्यं श्रुत्वा ‘ब्राह्मणोऽयमरण्ये चरैवेत्येवं मामुक्तवानि’ति मनसि ब्राह्मणवाक्ये महान्तमादरं कृत्वा पुनरेकं संवत्सरमरण्ये चरित्वा पश्चात् पितरं द्रष्टुं तमागच्छन्तं पुनरपीन्द्रो ब्राह्मणरूपेणाऽऽगत्यैवमुवाच- चरतः= पर्यटनं कुर्वतः पुरुषस्य जङ्‌घे पुष्पिण्यौ भवतः। यथा पुष्पयुक्ता वृक्षशाखा, लता वाऽथवा सुगन्धोपेता सेव्या भवति, एवं चरतो जङ्‍घे श्रमजयेन सेव्ये भवतः । तथैव आत्मा= मध्यदेहः, भूष्णुः= वर्धिष्णुः फलग्राहिः= आरोग्यरूपफलयुक्तो भवति । यथा वर्धमानो वृक्षः कालेन फलानि गृह्णाति, एवं चरतः पुरुषः पुरुषस्य बीजादिदीपनादिपाटवेन मध्यदेह आरोग्यरूपं फलं गृह्णाति । तथैव अस्य= चरतः पुरुषस्य सर्वे पाप्मानः= सर्वाणि पापानि प्रपथे= प्रकृष्टे तीर्थक्षेत्रादिमार्गे श्रमेण= तत्तद्देवतादिदर्शने तीर्थयात्रादिप्रयासेन हताः= विनाशिताः सन्तः शेरे= शेरते शयाना इव भवन्ति । यथा शयानाः पुरुषाः स्वकार्यं कृषिवाणिज्यादिकं कर्तुमशक्ताः, एवं पुण्येन विनष्टाः पाप्मानो नरकं दातुमसमर्था इत्यर्थः । तस्मात् सर्वथा अरण्ये चर, न पितुर्गृहेऽवतिष्ठस्व ।

\customparagraph{\textbf{नेपाली}}

रोहितले ब्राह्मण रूपका इन्द्रको वचन सुनेर धेरै आदरपूर्वक मनमा विचार गरे कि ‘यी ब्राह्मण हुन्, यिनले मलाई भने कि भ्रमण गर’ ब्राह्मणको कुरा मानेर रोहितले फेरि दोस्रो वर्ष पनि जङ्‌गलमै भ्रमण गर्ने विचार गरी त्यसै गरे । वर्ष पूरा भएपछि फेरि पितालाई भेट्न गाउँतर्फ के जान लागेका थिए, बाटैमा इन्द्र पुनः पुरुष रूप धारण गरी आएर रोहितलाई भन्न लागे- हे रोहित ! पर्यटन गर्नेवाला पुरुषका जाँघ (तिघ्रो) पुष्पयुक्त सुगन्धित वृक्ष वा लहराजस्तै सेव्य हुन्छ । त्यस्तै विचरण गर्नाले श्रमका कारण त्यो अत्यन्त सेव्य हुन्छ । त्यसरी नै आत्मारूप शरीरको मध्यभाग पनि संवर्धित भई आरोग्य रूप फलले युक्त हुन्छ । जसरी बढ्दै गरेको रुख समयले फल प्राप्त गर्दछ त्यसरी नै भ्रमण गर्ने पुरुषको मध्य शरीर बीजादिको पूर्णताले आरोग्य रूप फल प्राप्त गर्दछ । त्यस्तै भ्रमण गर्नेवाला पुरुषका सबै पाप उत्कृष्ट तीर्थ, क्षेत्र आदिका मार्गमा देवातादिदर्शन रूप तीर्थयात्राका श्रमले नष्ट भएर सुत्छन् अर्थात् सुतेको मानिस जसरी कृषि आदि कार्य गर्न असमर्थ हुन्छ त्यसैगरी सुतेका पापले कुनै हानि पुर्‍याउन सक्दैनन् । त्यसैले तिमी सर्वथा विचरण गर ।

तृतीयं पर्यायं दर्शयति-

\begin{visheshshloka}[center]

\textbf{चरैवेति वै मा ब्राह्मणोऽवोचदिति ह तृतीयं संवत्सरमरण्ये चचार, सोऽरण्याद् ग्राममेयाय}

\textbf{तमिन्द्रः पुरुषरूपेण पर्येत्योवाच-}

\textbf{आस्ते भग आसीनस्योर्ध्वस्तिष्ठति तिष्ठतः ।}

\textbf{शेते निपद्यमानस्य चराति चरतो भगश्चरैवेति ॥ इति ।}

\end{visheshshloka}

\customparagraph{संस्कृतव्याख्या}

भगः= सौभाग्यम्, आसीनस्य= उपविष्टस्य आस्ते= तथैव तिष्ठति, न तु वर्धते । अभिवृद्धिहेतोरुद्योगस्याभावात् । तिष्ठतः= उपवेशनं परित्यज्योत्थापनं कुर्वतः पुरुषस्य भगः ऊर्ध्वः= अभिवृद्धेरुन्मुखस्तिष्ठति । कृषिवाणिज्याद्युद्योगस्य संभावितत्त्वात् । निषद्यमानस्य= भूमौ शयानस्य भगः शेते= निद्रां करोति । विद्यमानधनरक्षादिचिन्ताया अप्यभावात् सर्वथैव विनश्यति । चरतः= तेषु तेषु देशेष्वर्जनाय पर्यटनं कुर्वतः पुरुषस्य भगः= सौभाग्यं चराति= दिने दिने वर्धते । तस्मात्त्वं चरैवेति, नत्वेकत्र तिष्ठ।

\customparagraph{\textbf{नेपाली}}

रोहितले ती ब्राह्मणको कुरा सुनेर आदरपूर्वक विचार गरे कि ब्राह्मणले मलाई भ्रमण गर्न भनेका छन् । त्यसकारण उनी तेस्रो वर्ष पनि वनमै घुमिरहे । वर्ष दिन बितेपछि जब पिताजीलाई भेट्ने मनसायले गाउँ फर्कन लागे तब पुरुष रूपधारी इन्द्र आएर उनलाई भने- हे रोहित ! बसेका व्यक्तिको भाग्य निष्क्रिय भई बस्छ, उठेको व्यक्तिको भाग्य समृद्धितर्फ उन्मुख हुन्छ । भूमिमा सुत्ने व्यक्तिको भाग्य सुत्छ । देशविदेशमा पर्यटन गर्ने व्यक्तिको भाग्य दिनदिनै समृद्ध हुन्छ । त्यसैले समृद्धि प्राप्तिका लागि तिमी पनि भ्रमण गर ।

चतुर्थं पर्यायं दर्शयति-

\begin{visheshshloka}[center]

\textbf{चरैवेति वै मा ब्राह्मणोऽवोचदिति ह चतुर्थं संवत्सरमरण्ये चचार; सोऽरण्याद् ग्राममेयाय}

\textbf{तमिन्द्रः पुरुषेण पर्येत्योवाच-}

\textbf{कलिः शयानो भवति संजिहानस्तु द्वापरः ।}

\textbf{उत्तिष्ठस्त्रेता भवति कृतं संपद्यते चरंश्चरैवेति॥ इति ।}

\end{visheshshloka}

\customparagraph{संस्कृतव्याख्या}

पुरुषस्य चतस्रोऽवस्था भवन्ति- निद्रा, तत्परित्यागः, उत्थानम्, सञ्चरणञ्चेति । ताश्चोत्तरोत्तरश्रेष्ठत्वात् कलिद्वापरत्रेतासत्ययुगैः समाना भवन्ति । अतः सञ्चरणस्य सर्वोत्तमत्वात् त्वं चरैवेति, सञ्चर इत्यर्थः ।

\customparagraph{\textbf{नेपाली}}

रोहितले ती ब्राह्मणलाई धेरै आदर गरेर सोचे कि सज्जन ब्राह्मणले मलाई भ्रमण गर्न भनेका छन् । त्यसैले रोहित चौथो वर्ष पनि वनमै घुमिरहे । जब वर्षान्तमा पितालाई सम्झेर भेट्ने विचार गरी जङ्‍लबाट गाउँतर्फ लागे तब पुरुष रूपधारी इन्द्र आएर उनलाई भने- हे रोहित ! कलियुग सुतिरहेको हुन्छ, द्वापर जागिरहेको हुन्छ, त्रेतायुग उठेको हुन्छ, त्यस्तै सत्ययुग चाहिँ सञ्चरण गरिरहन्छ । यी युगहरू क्रियाकलापका आधारमा उत्तरोत्तर श्रेष्ठ छन् । अतः कली, द्वापर, त्रेता र सत्य युगझैं उठेर. जागेर मात्रै नभएर सञ्चरण नै उत्तरोत्तर श्रेष्ठ हुनाले तिमी पनि श्रेष्ठताका लागि भ्रमण गर ।

पञ्चमं पर्यायं दर्शयति-

\begin{visheshshloka}[center]

\textbf{चरैवेति वै मा ब्राह्मणोऽवोचदिति ह पञ्चमं संवत्सरमरण्ये चचार, सोऽरण्याद् ग्राममेयाय}

\textbf{तमिन्द्रः पुरुषरूपेण पर्येत्योवाच-}

\textbf{चरन् वै मधु विन्दति चरन् स्वादुमुदुम्बरम् ।}

\textbf{सूर्यस्य पश्य श्रेमाणं यो न तन्द्रयते चरंश्चरैवेति ॥ इति।}

\end{visheshshloka}

\customparagraph{संस्कृतव्याख्या}

चरत्= भ्रमदेव पुरुषः क्वचिद् वृक्षाग्रे मधुम्= माक्षिकं लभते । क्वचित् स्वादु= मधुरम्, उदुम्बरम्= फलविशेषम्, लभते । अर्थात् भ्रमशीलः पुरुषः कुत्रचित् किमपि प्राप्नोति । अत्र सूर्यो दृष्टान्तरूपेण ज्ञातुं शक्यते । यः= सूर्यः सर्वत्र चरन्नपि न तन्द्रयते= आलस्यं नानुभवति, अतस्तस्य सूर्यस्य श्रेयाणम्= श्रेष्ठत्वम्, जगद्वन्द्यत्वं पश्य। तस्मात्त्वमपि चरैव ॥

\customparagraph{\textbf{नेपाली}}

यसरी सज्जन ब्राह्मणले मलाई भ्रमणको फइदाका बारेमा भने, अब पनि म उनको कुरा मान्छु भन्ने सोचेर पाँचौं वर्ष पनि रोहित वनमा घुम्न लागे । पाँचौं वर्ष पूरा भएपछि पिताजीलाई सम्झेर ती रोहित गाउँतर्फ प्रस्थान गर्नासाथ पुनः ती पुरुषरूपधारी इन्द्र आएर भने-

जो मानिस भ्रमणशील हुन्छ, त्यसले हिँड्‌दै जाँदा बाटामा मौरीको मह, डुम्री आदि मधुर फल भेट्न सक्छ । यसरी निरन्तर विचरण गर्दा यशरूपी फल पाउने व्यक्तिको उदाहरणमा सूर्य पनि पर्छन् । सूर्य आलस्य नमानी निरन्तर विचरण गर्छन्, त्यसैकारण संसारका पूज्य बनेका छन् । त्यसैले तिमी पनि निरन्तर भ्रमण गर, संसारले मान्ने बन्नेछौ ।

इत्थमिन्द्रकृतेन रोहितोपदेशेन चरतो रोहितस्य  स्वजीवने पितुरारोग्ये च कारणभूतं श्रेयोलाभं दर्शयति-

\begin{visheshshloka}[center]

\textbf{चरैवेति वै मा ब्राह्मणोऽवोचदिति ह षष्ठं संवत्सरमरण्ये चचार, सोऽजीगर्तं}

\textbf{सौयवसिमृषिमशयना परीतमरण्य उपेयाय-}

\end{visheshshloka}

\customparagraph{संस्कृतव्याख्या}

षष्ठे संवत्सरे पूर्ववदरण्यसञ्चारी स ह= रोहितः, तस्मिन्नरण्ये  कञ्चिदृषिम्, उपेयाय= प्राप्तवान् । कीदृशीमृषिम् ? अशनया परीतम्= अन्नस्यालाभेन क्षुत्पीडितम्, अजीगर्तनामकं सूयवसस्य पुत्रम्॥

\customparagraph{\textbf{नेपाली}}

पूर्ववत् सज्जन ब्राह्मणले मलाई भ्रमणको फइदाका बारेमा भने, अब पनि म उनको कुरा मान्छु भन्ने सोचेर छैटौं वर्ष पनि रोहित वनमा घुम्न लागे । यसपटक रोहितले वनमा भोकले पीडित सूवयसका पुत्र अजीगर्त नामका ऋषिलाई भेटे ।

अथाजिगर्तरोहितयोः संवादं दर्शयति-

\begin{visheshshloka}[center]

\textbf{तस्य ह त्रयः पुत्रा आसुः, शुनःपुच्छः शुनःशेपः शुनोलाङ्‍गुल इति; तं होवाच- ऋषेऽहं ते शतं ददाम्यहमेषामेकेनाऽऽत्मानं निष्क्रीणा इति, स ज्येष्ठं पुत्रं निगृह्णान उवाच- नन्विममिति; नो एवेममिति कनिष्ठं माता; तौ ह मध्यमे सम्पादयाञ्चक्रतुः शुनःशेपे, तस्य ह शतं दत्त्वा स तमादाय सोऽरण्याद् ग्राममेयाय ॥ इति।}

\end{visheshshloka}

\customparagraph{संस्कृतव्याख्या}

तस्य= अजीगर्तस्य शुनःपुच्छादिनामकास्त्रयः पुत्रा आसुः । तम्= पुत्रवन्तमृषिम्, रोहितः उवाच- हे ऋषे ! ते= तुभ्यम्, अहं गवां शतं ददामि । एषाम्= पुत्राणाम्,  एकेन अहम् आत्मानं निष्क्रिणा इति (अहमेतेषां पुत्राणां मध्ये एकेन पुत्रेणाऽऽत्मानं मद्देहं वरुणान्निष्क्रीणं मूल्यं दत्त्वाऽऽत्मानं मोचयामि) इति रोहित उक्तवान्। रोहितस्य वचनं श्रुत्वा सोऽजीगर्तो ज्येष्ठं पुत्रं शुनःपुच्छनामकं हस्तेन निगृह्णानः= स्वसमीपे समाकर्षन् रोहितं प्रत्येवमुवाच- तुभ्यमेकः पुत्रो दीयते, इमं नु शुनःपुच्छं तु न ददामि, मम प्रियत्वादिति । ततो माता कनिष्ठंं पुत्रं हस्तेन गृहीत्वैवमुवाच- इमं शुनोलाङ्‍गुलं तु मम प्रियं नो एव= सर्वथा न ददामीति । ततः तौ उभौ मातापितरौ मध्यमे पुत्रे शुनःशेपे दानं सम्पादयाञ्चक्रतुः । ततस्तस्य अजीगर्तस्य सः= रोहितो गवां शतं दत्त्वा तं शुनःशेपम् आदाय वनात् स्वकीयं ग्रामं प्रति प्रत्याजगाम ।

\customparagraph{\textbf{नेपाली}}

ती अजीगर्तका तीन छोरा थिए- शुनःपुच्छ, शुनःशेप र शुनलाङ्‍गुल । राजकुमार रोहितले अजीगर्तलाई भने- हे ऋषि ! यी छोराहरूमध्ये एक मलाई दिनुहोस्, त्यसबाट निष्क्रय गरेर म आफूलाई वरुणबाट मुक्त गराउँछु । यो सुनेर अजीगर्तले जेठो छोरो शुनःपुच्छलाई तानेर भन्यो- तिमीलाई एउटा पुत्र त दिन्छौँ तर यो होइन किनकि मेरो सबैभन्दा प्यारो जेठो छोरो छ । यसरी नै आमाले कान्छो छोरो शुनोलाङ्‌गुललाई तानेर भनिन्- यो मेरो प्यारो छ, यसलाई दिन्न । यसपछि ती दम्पती माइँलो छोरो शुनःशेपलाई दिन राजी भए। रोहितले सय गाई दिएर शुनःशेपलाई लिई जङ्‍गलबाट गाउँ फर्किए ।

तदागमदादूर्ध्वकालीनं वृत्तान्तं दर्शयति-

\begin{visheshshloka}[center]

\textbf{स पितरमेत्योवाच- तत हन्ताहमनेनाऽत्मानं निष्क्रीणा इति; स वरुणं राजानमुपससारानेन त्वा यजा इति; तथेति, भूयान् वै ब्राह्मणः क्षत्त्रियादिति वरुण उवाच, तस्मा एतं राजसूयं यज्ञक्रतुं प्रोवाच, तमेतमभिषेचनीये पुरुषं पशुमालेभे ॥१५॥ इति।}

\end{visheshshloka}

\customparagraph{संस्कृतव्याख्या}

सः= रोहितः पितरमागत्यैवमुवाच- हे तत=पितः ! हन्त= आवयोर्हर्षः सम्पन्नः । अहम्= रोहितः,  अनेन‍= शुनःशेपरूपेण मूल्येन आत्मानम्= मद्देहम्, निष्क्रीणं मूल्यं दत्त्वाऽऽत्मानं मोचयामीत्यर्थः । पुत्रेण तथोक्ते सति सः= हरिश्चन्द्रः, वरुणमुपेत्य अनेन= शुनःशेपेन ब्राह्मणेन त्वा= त्वाम्, यजा= यक्ष्यामीत्युक्तवान् । सः= वरुणोऽपि तथेत्यङ्‌गीकृत्यैवमुवाच- क्षत्त्रियात्= तव पुत्राद् रोहितादप्ययं ब्राह्मणो भूयान्= अभ्यधिक एव, मम प्रियः, इत्युक्त्वा तस्मै= हरिश्चन्द्राय कर्तव्यत्वेन राजसूयमुपदिदेश। सः= हरिश्चन्द्रो राजसूयं प्रक्रम्य तस्य मध्ये योऽयमभिषेचनीयाख्य एकाहः सोमयागः, तस्मिन् तमेतम्= शुनःशेपम्, पुरुषं पशुम्= पुरुषरूपं पशुम्, आलेभे= सवनीयपशुत्वेनालब्धुं निश्चितवान् ।

\customparagraph{\textbf{नेपाली}}

रोहितले पिताका नजिकै आएर भने- हे पिताजी ! खुसीको कुरा छ कि म यो शुनःशेप रूपको मूल्य दिएर वरुणबाट आफूलाई छुटाउँछु । रोहितको यस्तो कुरा सुनेर ती हरिश्चन्द्र वरुणका नजिक आई भन्छन्- हे वरुण ! म यो ब्राह्मण पुत्रबाट तिम्रो यज्ञ गर्छु । वरुणले भने- ठिक छ, किनकि तिम्रो छोरो क्षत्त्रियभन्दा ब्राह्मण बालक नै मलाई अधिक प्रिय छ । त्यसपछि वरुणले राजालाई राजसूय यज्ञको उपदेश दिए । ती हरिश्चन्द्रले त्यो एक दिने सोमयागको बीचमा पशुरूपमा शुनःशेपलाई स्थापित गर्न लागे ।

\begin{visheshshlokatwo}[center]{84}

\textbf{ऐतरेयब्राह्मणको  तेत्तिसौं अध्याय (सातौँ पञ्चिकाको तेस्रो अध्यायको तृतीय खण्ड)}

\textbf{नेपाली भावानुवाद पूर्ण भयो ।}

\end{visheshshlokatwo}

\specialupakhanda{\textbf{चतुर्थः खण्डः}}

हरिश्चन्द्रस्य राजसूयानुष्ठानं दर्शयति-

\begin{visheshshloka}[center]

\textbf{तस्य ह विश्वामित्रो होताऽऽसीज्जमदग्निरध्वर्युर्वशिष्ठो ब्रह्माऽयास्य उद्गाता; तस्मा उपाकृताय नियोक्तारं न विविदुः; स होवाचाजीगर्तः सौयवसिर्मह्यमपरं शतं दत्ताहमेनं नियोक्ष्यामीति; तस्मा अपरं शतं ददुस्तं स निनियोज ॥ इति ।}

\end{visheshshloka}

\customparagraph{संस्कृतव्याख्या}

तस्य= यज्ञस्य विश्वामित्रः- होता, जमदग्निः- अध्वर्युः, वशिष्ठः- ब्रह्मा, अयास्यः- उद्गाता चासीत् । तत्र जमदग्निरध्वर्युरभिषेचनीये सोमयागे तम्= शुनःशेपं सवनीयपशुत्वेनोपाकृतवान् । तस्मा= तस्मै उपाकृताय= मेध्यत्वेन संस्कृताय शुनःशेपाय नियोक्तारम्= यूपे बन्धनकर्तारं क्रूरं कञ्चिदपि पुरुषं न विविदुः= लेभिरे । तदानीं सुयवसस्य पुत्रः, शुनःशेपस्य पिता अजीगर्तः, उवाच- हे यज्ञस्थाः कर्तारः ! मह्यं पूर्वस्माच्छतादपरं गोशतं दत्त, ततोऽहम् एनम्= शुनःशेपम्, यूपे= यज्ञस्तम्भे नियोक्ष्यामि= रशनया कट्यां शिरसि पादयोर्बध्वा बन्धनं करिष्यामि । तस्मा= तस्मै अजीगर्ताय अपरं गोशतं ददुः । सः= अजीगर्तः तम्= शुनःशेपम्, निनियोज= यूपे बन्धनं चकार ।

\customparagraph{\textbf{नेपाली}}

हरिश्चन्द्रले गर्न लागेको त्यस यज्ञमा विश्वामित्र- होता, जमदग्नि- अध्वर्यु वशिष्ठ- ब्रह्मा र अयास्य ऋषि उद्गाता थिए । मेध्य वस्तुका रूपमा मन्त्रादि सुसंस्कृत त्यस शुनःशेपलाई यज्ञको खाबामा बाँध्ने क्रूर पुरुष कोही पाइएनन् । त्यसवेला शुनःशेपका पिता अजीगर्तले भने- हे यजमान र ब्राह्मणहरू हो ! यदि मलाई पहिले दिएजस्तै अरू सय गाई दिने हो भने यो काम पनि मै गर्छु । त्यो अजीगर्तलाई सय गाई फेरि दिइयो । त्यसले शुनःशेपलाई यज्ञीय खम्बामा बाध्यो ।

अथ विशसनप्रवृत्तिं दर्शयति-

\begin{visheshshloka}[center]

\textbf{तस्मा उपकृताय नियुक्तायाऽऽप्रीताय पर्यग्निकृताय विशसितारं विविदुः; स होवाच- अजीगर्तः सौयवसिर्मह्यमपरं शतं दत्ताहमेनं विशसिष्यामीति; तस्मा अपरं शतं ददुः, सोऽसिं निःशान एयाय ॥ इति ।}

\end{visheshshloka}

\customparagraph{संस्कृतव्याख्या}

उपाकृताय= मेध्यरूपेण सुसंस्कृताय, नियुक्ताय= यूपे स्थापिताय, आप्रीताय=  आप्रीसंज्ञिताभिरेकादशभिः प्रयाज-याज्याभिर्यद् यजनं तद् आप्रीकरणम्, पर्यग्निकृताय= दर्भरूपेणोल्मुकेन त्रिःप्रदक्षिणीकरणं तत् पर्यग्निकरणम्, तथाविधसंस्कारचतुष्टययुक्ताय तस्मा= तस्मै शुनःशेपाय हिंसितारं क्रूरपुरुषं न लेभिरे । ततः सोऽजीगर्तः पूर्ववदपरं गोशतं गृहीत्वा मारयितुम् असिम्= खड्गम्,  निःशानः= निशितं तीक्ष्णं कुर्वन्नेव एयाय= जगाम ।

\customparagraph{\textbf{नेपाली}}

मेध्य पशुका रूपमा विभिन्न उपाकरणले सुसंस्कृत बनाएर यज्ञीय खम्बामा बाँधेपछि ‘आप्री’ संज्ञक एघार मन्त्रको पाठ पूर्ण गरेर साथै कुशरूप उल्मुकले तीनपटक प्रदक्षिणारूप पर्यग्निकरण गरेपछि त्यसलाई बध गर्न कोही क्रूर पुरुष पाइएनन् । त्यसैवेला अजीगर्तले अरू सय गाई दिन्छौ भने यो काम पनि मै गर्छु भन्यो । उसको इच्छाअनुसार सय गाई दिएर अजीगर्तलाई नै बलि दिन लगाइयो । त्यो अजीगर्त खड्गमा धार लाएर शुनःशेपलाई प्रहार गर्न तयार भयो ।

पितुरजीगर्तस्य वृत्तान्तं दर्शयति-

\begin{visheshshloka}[center]

\textbf{अथ ह शुनःशेप ईक्षाञ्चक्रेऽमानुषमिव वै मा विशसिष्यन्ति, हन्ताहं देवता उपधावामीति स प्रजापतिमेव प्रथमं देवतानामुपससार, कस्य नूनं कतमस्यामृतानामित्येतयर्चा ॥ इति।}

\end{visheshshloka}

\customparagraph{संस्कृतव्याख्या}

अथ= पितुः पुत्रमारणोद्योगानन्तरं शुनःशेपः= पुत्रो मनस्येवम् ईक्षाञ्चक्रे= विचारितवान् । ‘अन्यत्र पर्यग्निकृतं पुरुषमारण्यांश्चोत्सृजन्त्यहिंसायै’ इति श्रुतेः पर्यग्निकरणादूर्ध्वं मनुष्यं परित्यजन्ति । एते तु मा= माम् अमानुषमिव= मनुष्यव्यतिरिक्तमजादिपशुमिव विशसिष्यन्ति= मारयिष्यन्ति । हन्त= हा कष्टमेतत्समापन्नम् । अहमितः परं रक्षार्थं देवता उपधावामि= भजामि, इति= एतद् विचार्य देवानां मध्ये प्रथमम्= मुख्यम्, प्रजापतिमेव= प्रजापतिदेवमेव ‘कस्य नूनम्–’ इत्यृचा उपससार= सेवितवान् ।

\customparagraph{\textbf{नेपाली}}

आफ्नै पिता आफूलाई मार्न उद्यत भएपछि पुत्र शुनःशेपका मनमा यस्तो विचार आयो- ‘पर्यग्निसंस्कार गरेर अहिंसाका लागि मनुष्यलाई वनमा छाडिदिन्छन्’ भन्ने श्रुतिकथन छ । तर यी यज्ञकर्ताहरूको व्यवहार देख्दा मलाई मार्ने कुरामा शङ्का छैन, यो मेरा लागि अत्यन्त कष्टप्रद छ, त्यसैले अन्तिम उपाय स्वरूप म देवताका शरणमा जान्छु । यस्तो विचार गरी शुनःशेपले सर्वप्रथम प्रजापति देवताको ‘कस्य नूनम्—’ यो मन्त्रद्वारा प्रार्थना गर्‍यो ।

तस्य प्रजापतेः सहकारित्वेनाग्नेः सेवां दर्शयति-

\begin{visheshshloka}[center]

\textbf{तं प्रजापतिरुवाचाग्निर्वै देवानां नेदिष्ठस्तमेवोपधावेति सोऽग्निमुपससाराग्नेर्वर्यं प्रथमस्यामृतानामित्येतयर्चा ॥ इति ।}

\end{visheshshloka}

\customparagraph{संस्कृतव्याख्या}

तम्= प्रार्थयितारं शुनःशेपम्, प्रजापतिरेवमुक्तवान्- अग्निः सर्वेषां देवानां नेदिष्ठः= अग्निरेव हविर्वहनेन देवानामतिसमीपवर्ती । अतस्तमेवोपास्स्वेत्युक्तः सः= शुनःशेपः ‘अग्नेर्वयम्—’ इत्यृचाऽग्निमुपासितवान् ।

\customparagraph{\textbf{नेपाली}}

ती आफ्नो प्रार्थना गरिरहेका शुनःशेपलाई प्रजापतिले भने- हविष्य पुर्‍याइदिने भएकाले सबै देवतासँग अग्नि नजिक छन् । त्यसकारण तिमी अग्निसँग जाऊ । त्यसपछि शुनःशेप ‘अग्नेर्वयं प्रथमस्यामृतानाम्’ यो मन्त्रपाठ गर्दै अग्निका शरणमा गए ।

तस्य चाग्नेः सहकारित्वेन सवितुरुपासनं दर्शयति-

\begin{visheshshloka}[center]

\textbf{तमग्निरुवाच- सविता वै प्रसवानामीशे तमेवोपधावेति; स सवितारमुपससाराभि त्वा देव}

\textbf{सवितरित्येतेन तृचेन ॥ इति ।}

\end{visheshshloka}

\customparagraph{संस्कृतव्याख्या}

प्रसवानाम्= सर्वेषु कार्येषु प्रेरणारूपानामनुज्ञानात् सविता  ईशे‍= स्वामी भवति । तस्मात्तमेवोपश्रावेत्यग्निनोपदिष्टः ‘अभि त्वा---’ इति तृचेन सवितारमुपासितवान् ।

\customparagraph{\textbf{नेपाली}}

आफ्नो प्रार्थना गर्दै गरेका शुनःशेपलाई अग्निले भने- प्रेरक रूपमा अनुज्ञा दिनाले सविता सबै कार्यका स्वामी हुन् । त्यसैले तिमी सविताको प्रार्थना गर । अग्निको कुरा सुनेर शुनःशेप ‘अभि

त्वा देव सवितः०’ यो ऋचाले स्तुति गर्दै सविताका शरणमा गए ।

तस्य च सवितुः सहकारित्वेन वरुणस्योपासनं दर्शयति-

\begin{visheshshloka}[center]

\textbf{तं सवितोवाच- वरुणाय वै राज्ञे नियुक्तोऽसि, तमेवोपधावेति; स वरुणं}

\textbf{राजानमुपससारात उत्तराभिरेकत्रिंशता ॥ इति ।}

\end{visheshshloka}

\customparagraph{संस्कृतव्याख्या}

सविता शुनःशेपमुवाच- हे शुनःशेप ! वरुणार्थं त्वं यूपे बद्धोऽसि । अतो वरुणमुपास्स्वेति । सवित्रोक्तः स पूर्वस्मात् सवितृविषयात् तृचाद् उत्तराभिरेकत्रिंशत् सङ्‌ख्याकाभिर्ऋग्भिर्वरुणमुपासितवान् । ‘न हि ते क्षत्त्रत्वम्०’ इत्याद्याः सूक्तविशेषभूता दशर्चो ‘यच्चिद्धि ते विश---’ इत्यादिकमेकविंशत्यृचं सूक्तम्, इत्येवमेकत्रिंशत्‌सङ्‍ख्यकाभिर्ऋग्भिः शुनःशेपो वरुणस्य स्तुतिं कृतवान् ।

\customparagraph{\textbf{नेपाली}}

प्रार्थना गर्दै गरेका शुनःशेपलाई सूर्यले भने- हे शुनःशेप ! तिमी वरुणका लागि बाँधिएका हौ । त्यसैले तिमी तिनैका शरणमा जाऊ । सविताको कुरा सुनेर शुनः शेपले पूर्वसूक्तका अवशिष्ट ‘नहि ते क्षत्त्रत्वम्०’ आदि (१/२४/६-१५) दश र ‘यच्चिद्धि ते विश०’ आदि (१/२५) सूक्तका सम्पूर्ण एक्काइस ऋचा गरी जम्मा एकतीस ऋचाबाट वरुणको स्तुति गरे ।

तस्य च वरुणस्य सहकारित्वेन पुनरप्यग्नेरुपासनं दर्शयति-

\begin{visheshshloka}[center]

\textbf{तं वरुण उवाचाग्निर्वै देवानां मुखं सुहृदयतमस्तं नु स्तुह्यथ त्वोत्स्रक्ष्याम इति; सोऽग्निं}

\textbf{तुष्टावात उत्तराभिर्द्वाविंशत्या ॥ इति ।}

\end{visheshshloka}

\customparagraph{संस्कृतव्याख्या}

शुनःशेपं वरुण उवाच-  अग्निद्वारेणैव सर्वैर्देवैर्हविःस्वीकाराद् हविर्वहनादतिशयेन प्रीत्या सुहृदयतमोऽयमग्निः सर्वेषां देवानां मुखस्थानीयोऽस्ति । तम्= अग्निम्, नु= क्षिप्रम्, स्तुहीति वरुणेनोक्तः पूर्वोक्ताभ्यः ऋग्भ्य उत्तराभिर्द्वाविंशतिसङ्‍ख्याकाभिर्ऋग्भिरग्निं तुष्टाव । ‘वसिष्वा हि--’ इत्यादिकं दशर्चं सूक्तम्, ‘अश्वं न त्वा-- इत्यादिकं त्रयोदशर्चं सूक्तम्’, तत्रान्त्यां परित्यज्य वसिष्वसूक्तद्वयगता ऋचो द्वाविंशतिसङ्‌ख्याकाः ।

\customparagraph{\textbf{नेपाली}}

आफ्नो शरणमा आएका शुनःशेपलाई वरुणले भने- “यज्ञको हवि वहन गरेर सबै देवतालाई पुर्‍याइदिनाले अग्नि सबै देवताका अत्यन्त प्रिय छन्, त्यसैले तिमी तिनैको स्तुति गर । अनि म तिमीलाई छाडिदिन्छु ।” त्यसपछि ती शुनःशेपले पूर्वोक्त ऋचाभन्दा पछि रहेका (१-२६ सूक्तका सम्पूर्ण)  १० ऋचा र (१-२७ सूक्तका अन्तिम १३ औँ मन्त्र छाडी) १२  ऋचा गरी बाइस ऋचा द्वारा अग्निको स्तुति गरे ।

अग्नेः सहकारित्वेन विश्वेषां देवानामुपासनं दर्शयति-

\begin{visheshshloka}[center]

\textbf{तमग्निरुवाच- विश्वान्नु देवान् स्तुह्यथ त्वोस्रक्ष्याम इति; स विश्वान् देवांस्तुष्टाव, नमो}

\textbf{महद्भ्यो नमो अर्भकेभ्य इत्येतयर्चा ॥ इति।}

\end{visheshshloka}

\customparagraph{संस्कृतव्याख्या}

(यद्यपि वरुणपाशेन बद्धत्वाद् वरुण एव शुनःशेपमुत्स्रष्टुं समर्थः। तथाऽप्यग्न्यादीनां सहकारित्ववचनं दार्ढ्यार्थं द्रष्टव्यम्) विश्वे देवा गणरूपा न भवन्ति तथापि सर्वे देवास्तान् ‘नमो महद्भ्यः--’ इत्येतयर्चा उपासितवान् ।

\customparagraph{\textbf{नेपाली}}

(यद्यपि वरुणका कारण पाशबद्ध भएकाले वरुण नै मुक्तिका लागि पर्याप्त छन्, तथापि उनले अन्यको स्तुति गर्न प्रेरित गर्नु भनेको सबै देवताको सम्मान र समर्थनका लागि हो भनी बुझ्नुपर्छ) ती शुनःशेपलाई अग्निले भने- “विश्वे देवको स्तुति गर अनि हामी तिमीलाई छाडिदिन्छौँ।” अग्निको आज्ञाअनुसार तिनले विश्वे देवको ‘नमो महद्भ्यः’ (१/२७/१३) भन्ने ऋचाले स्तुति गरे ।

तेषां  च विश्वेषां देवानां सहकारित्वेनेन्द्रस्योपासनं दर्शयति-

\begin{visheshshloka}[center]

\textbf{तं विश्वे देवा ऊचुः- इन्द्रो वै देवानामोजिष्ठो बलिष्ठः सहिष्ठः सत्तमः पारयिष्णुतमस्तं नु स्तुह्यथ त्वोत्स्रक्ष्याम इति; स इन्द्रं तुष्टाव यच्चिद्धि सत्य सोमपा इति चैतेन सूक्तेनोत्तरस्य च पञ्चदशभिः ॥ इति।}

\end{visheshshloka}

\customparagraph{संस्कृतव्याख्या}

इन्द्रस्य वैशिष्ट्यप्रकाशका ओजोबलादिशब्दाः पूर्वाचार्यैरेवं व्याख्याताः-

\begin{shloka}[center]

ओजो दीप्तिर्बलं दाक्ष्यं प्रसह्यकरणं सहः ।

सुजनः सन्पारयिष्णुरुपक्रान्तसमाप्तिकृत् ॥

\end{shloka}

\customparagraph{संस्कृतव्याख्या}

इष्ठप्रत्यय-तमप्रत्ययाभ्यां तत्र तत्रातिशय उच्यते । तादृशमिन्द्रं ‘यच्चिद्धि तस्य सोमपाः --’ इत्यनेन सप्तर्चेन सूक्तेन उत्तरस्मिन्नपि ‘आ व इन्द्रम्--’ इत्यादिके द्वाविंशत्यृचे सूक्ते पञ्चदशभिर्ऋग्भिश्च तुष्टाव ।

\customparagraph{\textbf{नेपाली}}

आफ्नै स्तुति गर्दै गरेका शुनःशेपलाई ती विश्वे देवले भने- देवताहरूमा इन्द्र सबैभन्दा धेरै ओजवान्, बलवान्, वीर्यवान्, सहिष्णु एवम् समस्याबाट पार लगाउन सक्ने सामर्थ्य भएका छन्। अतः तिमी तिनैको स्तुति गर अनि तिमीलाई छोडिदिन्छौं । विश्वेदेवले भने अनुसार तिनले ‘यच्चिद्धि तस्य सोमपाः --’ आदि गरिएका (१/२९ का सम्पूर्ण) सात ऋचा र ‘आ व इन्द्रम्--’ आदि  (२२ ऋचा भएका ३० /१ सूक्तका आरम्भका) १५ ऋचा गरी २२ ऋचाले  इन्द्रको स्तुति गरे।

पूर्वमप्यनुग्रहेण चरैवेत्येवमिन्द्रो रोहितं प्रत्युपदिदेश । तस्योपदेशस्य फलपर्यवसायित्वमनुस्मरतः तस्येन्द्रस्य शुनःशेपस्तुतौ प्रीत्यातिशयं दर्शयति-

\begin{visheshshloka}[center]

\textbf{तस्मा इन्द्रः स्तूयमानः प्रीतो मनसा हिरण्यरथं ददौ, तमेतया प्रतीयाय शश्वदिन्द्र इति ॥}

\end{visheshshloka}

\customparagraph{संस्कृतव्याख्या}

शुनःशेपेनेन्द्रः स्तूयमानः प्रीतो भूत्वा तस्मै= शुनःशेपाय सुवर्णमयं दिव्यं रथमारोहणार्थं स्वकीयेन मनसैव ददौ । शुनःशेपोऽपि तदीयमनुग्रहमवगत्य पूर्वोक्ताभ्यः पञ्चदशभ्य ‘उत्तरया’ ‘शश्वदिन्द्र’ इत्येतयर्चा तम्= रथम्, प्रतीयाय= मनसैव प्रतिजगाम ।

\customparagraph{\textbf{नेपाली}}

स्तुति गरिएका इन्द्रले प्रसन्न भएर ती शुनःशेपलाई सुवर्णमय रथ आरोहणका लागि आफ्नो मनैबाट दिए । शुनःशेपले पनि इन्द्रको कृपालाई जानेर पहिलेका पन्ध्र ऋचाभन्दा पछि ‘शश्वदिन्द्रः--’ इत्यादि (१/३०/१६) ऋचाले स्तुति गर्दै मनदेखि नै स्वीकार गरे ।

इन्द्रस्य सहकारित्वेनाश्वीनोरुपासनं दर्शयति-

\begin{visheshshloka}[center]

\textbf{तमिन्द्र उवाचाश्विनौ नु स्तुह्यथ त्वोत्स्रक्ष्याम इति; सोऽश्विनौ तुष्टावात उत्तरेण तृचेन ॥ इति}

\end{visheshshloka}

\customparagraph{संस्कृतव्याख्या}

इन्द्रकथनानुसारं शुनःशेपः ‘शश्वदिन्द्र’ इत्यस्या उत्तरेण ‘आश्विनावश्वावत्या--’ इति तृचेन अश्विनौ स्तुतवान् ।

\customparagraph{\textbf{नेपाली}}

ती शुनःशेपलाई इन्द्रले भने- तिमी अश्विनीकुमारको स्तुति गर, अनि तिमीलाई छाडिदिन्छौं । इन्द्रको भनाइ मानेर ती शुनःशेपले ‘आश्विनावश्वावत्या--’ इत्यादि तीन (१/३०/१७-१९) ऋचाद्वारा अश्विनीकुमारको स्तुति गरे ।

अश्विनोः सहकारित्वेनोषस उपास्तिं दर्शयति-

\begin{visheshshloka}[center]

\textbf{तमश्विना ऊचतुरुषसं नु स्तुह्यथ त्वोत्स्रक्ष्याम इति; स उषसं तुष्टावात उत्तरेण तृचेन ।}

\end{visheshshloka}

\customparagraph{संस्कृतव्याख्या}

अश्विनीकुमारकथनानुसारं स ‘कस्त उषः--’ इत्यादिकेन उत्तरेण तृचेन उषसं तुष्ठाव ।

\customparagraph{\textbf{नेपाली}}

ती शुनःशेपलाई अश्विनीकुमारले भने- तिमी उषाको स्तुति गर। अनि हामी तिमीलाई छाडिदिन्छौं । तिनले उषाको स्तुति ‘कस्त उष’ इत्यादि (१/३०/२०-२२) मन्त्रले गरे ।

अथोक्तानां सर्वासां देवतानामनुग्रहेण शुनःशेपस्य बन्धमोक्षं हरिश्चन्द्रस्याऽऽरोग्यं च दर्शयति-

\begin{visheshshloka}[center]

\textbf{तस्य ह स्मर्च्यृच्युक्तानां वि पाशो मुमुचे, कनीय ऐक्ष्वाकस्योदरं भवत्युत्तमस्यामेवर्च्युक्तायां वि पाशो मुमुचेऽगद ऐक्ष्वाक आस ॥ इति ।}

\end{visheshshloka}

\customparagraph{संस्कृतव्याख्या}

तस्य पूर्वोक्तस्य तृचस्य सम्बन्धिन्याम् ऋचि-ऋचि= एकैकस्यामृच्युक्तायां क्रमेण शुनःशेपस्य पाशो विमुमुचे= विशेषेण मुक्तोऽभूत् । ऐेक्ष्वाकस्य‍= हरिश्चन्द्रस्य यन्महोदरम्= व्याधिविशेषम्, तदपि क्रमेण कनीयः= अत्यल्पं भवति । उत्तमस्यामृच्युक्तायां पाशो विमुमुचे, एवं सर्वात्मना मुक्तोऽभूत् । ऐक्ष्वाकोऽपि अगदः= निःशेषेण रोगरहित आस ॥

\customparagraph{\textbf{नेपाली}}

ती पूर्वोक्त एक- एक ऋचा पाठ गर्दै जाँदा क्रमशः शुनःशेपका बन्धनहरू खुल्दै गए । हरिश्चन्द्रको पेट पनि सामान्य हुँदै गयो । अन्तिम मन्त्र पढ्दै गर्दा शुनःशेपका सबै बन्धन खुले साथै हरिश्चन्द्र रोगमुक्त भए ।

\begin{visheshshlokatwo}[center]{84}

\textbf{ऐतेरियब्राह्मणको  तेत्तिसौं अध्याय (सप्तपञ्चिकाको तेस्रो अध्यायको चतुर्थ खण्ड)}

\textbf{नेपाली भावानुवाद पूर्ण भयो ।}

\end{visheshshlokatwo}

\specialupakhanda{\textbf{पञ्चमः खण्डः}}

बन्धनान्मुक्तस्य शुनःशेपस्योत्तरकालीनं वृत्तान्तं दर्शयति-

\begin{visheshshloka}[center]

\textbf{तमृत्विज ऊचुस्त्वमेव नोऽस्याह्नः संस्थामधिगच्छेत्यथ हैतं शुनःशेपोऽञ्जः सर्वं ददर्श तमेताभिश्चतसृभिरभिसुषाव, यच्चिद्धि त्वं गृहे गृह इत्यथैनं द्रोणकलशमभ्यवनिनायोच्छिष्टं चम्वोर्भरेत्येतयर्चाऽथ हास्मिन्नन्वारब्धे पूर्वाभिश्चतसृभिः स स्वाहाकाराभिर्जुहवाञ्चकाराथैनमवभृथमभ्यवनिनाय, ‘त्वं नो अग्ने वरुणस्य विद्वानित्येताभ्यामथैनमत ऊर्ध्वमग्निमाहवनीयमुपस्थापयाञ्चकार’ शुनश्चिच्छेपं निदितं सहस्रादिति ॥ इति ।}

\end{visheshshloka}

\customparagraph{संस्कृतव्याख्या}

देवानुग्रहयुतं तम्= शुनःशेपं विश्वामित्रादयः सर्वे ऋत्विजः, एवमुचुः- हे शुनःशेप ! त्वमेव नः= अस्माकम्, अस्याह्नः= अभिषेचनीयाख्यस्य संस्थाम्= समाप्तिम्, अधिगच्छ= प्राप्नुहि, अनुष्ठापयेत्यर्थः । तैरेवमुक्ते सत्यनन्तरं शुनःशेपः एतम्= अभिषेचनीयाख्यं सोमयागम्, अञ्जःसवं ददर्श= अञ्जसा ऋजुमार्गेण सवः सोमभिषवो यस्मिन् यागे सोञ्जःसवः, तादृशं प्रयोगप्रकारं निश्चितवान् । निश्चित्य च तम्= सोमम्, ‘यच्चिद्धि त्वं गृहे गृहे-–’ इत्यादिभिश्चतसृभिर्ऋग्भिः अभिषुतं कृतवान् । अथ एनम्= अभिषुतं सोमम्, एतया ‘उच्छिष्टं चम्वोः–’ इत्यृचा द्रोणकलशमभिलक्ष्य अवनिनाय= द्रोणकलशे प्रक्षिप्तवान् । अथ= अनन्तरम्, अस्मिन्= हरिश्चन्द्रे अन्वारब्धे= शुनःशेपदेहमुपस्पृष्टवति सत्युक्ताभ्यः ऋग्भ्यः पूर्वाभिः= ‘यत्र ग्रावे’त्यादिभिश्चतसृभिर्ऋग्भिः स्वाहाकारसहिताभिः सोमं जुहवाञ्चकार । अथ= होमानन्तरमेव कर्तव्यमवभृथमभिलक्ष्य अवनिनाय=सर्वमवभृथसाधनं तद्देशे नीत्वा ‘त्वं नो अग्ने वरुणस्य विद्वान्—’ इत्यादिकाभ्यां ऋग्भ्यामवभृथयागं कृतवान् । अथ=  तथा कृत्वा तत ऊर्ध्वमेनमाहवनीयमग्निम् ‘शुनश्चिच्छेपं निदितं सहस्रात्--’ इत्यादिना उपस्थापयाञ्चकार= हरिश्चन्द्रमुपस्थापने प्रेरयामास ।

\customparagraph{\textbf{नेपाली}}

देवताहरूको कृपा प्राप्त गरेका ती शुनःशेपलाई विश्वामित्र आदि ऋषिहरूले भने- हे शुनःशेप ! तिमी हाम्रो यो अभिषेचनीय अह्नविधिलाई समापन गर । ऋत्विज्‌हरूले यति भनेपछि ती शुनःशेपले त्यो अभिषेचनीय नामक सोमयागको सरल मार्गले सोमरस निकाल्ने विशेष विधिको दर्शन गरे । त्यो सोमलाई ‘यच्चिद्धि त्वम्–’ इत्यादि (१/२८/५-८) चार मन्त्रद्वारा (अभिषुत) निचोरियो । यसपछि त्यो निचोरिएको सोमरसलाई ‘उच्छिष्टं चम्वोः--’ इत्यादि १/२८/९ ऋचाले द्रोण कलशमा अभिषेचन गरियो । यसपछि हरिश्चन्द्रले शुनःशेपको शरीरलाई स्पर्श गरेर ‘यत्र ग्राव--’ इत्यादि १/२८/१-४ चार मन्त्रले स्वाहाकारसहित सोमयाग गरे । यसरी हवन गरिसकेपछि ‘अभिमृथ कर्म’ अर्थात् यज्ञको अन्तिम कृत्यलाई अभिलक्षित गरी सम्पूर्ण यज्ञसाधनलाई त्यहाँ ल्याएर ‘त्वं नो अग्ने वरुणस्य-’ इत्यादि (४/ १/४-५) दुई ऋचाले अवभृथ अर्थात् समाप्ति यज्ञ गरियो । यसरी अन्तिम कृत्य गरिसकेपछि यो आहवनीय अग्निलाई ‘शुनश्चिच्छेपं निदितं--’ (५/२/७) इत्यादि ऋचाले हरिश्चन्द्र प्रति प्रेरित गरियो ।

विश्वामित्राजीगर्तयोः कञ्चित्संवाद दर्शयति-

\begin{visheshshloka}[center]

\textbf{अथ ह शुनःशेपो विश्वामित्रस्याङ्‍कमाससाद; स होवाचाजीगर्तः सौयवसिर्ऋषे पुनर्मे पुत्रं देहीति; नेति होवाच विश्वामित्रो, देवा वा इमं मह्यमरासतेति, स ह देवरातो वैश्वामित्र आस, तस्यैते कापिलेयबाभ्रवाः ॥ इति ।}

\end{visheshshloka}

\customparagraph{संस्कृतव्याख्या}

अथ= अभिषेचनीयसमाप्तेरनन्तरं हरिश्चन्द्रसहितेष्वृत्विक्षु विस्मितेषु स शुनःशेप इत ऊर्ध्वं कस्य पुत्रोऽस्त्विति विचारे सति तदीयेच्छैव नियामकेति महर्षीणां वचनं श्रुत्वा शुनःशेपः स्वेच्छया विश्वामित्रपुत्रत्वमङ्‌गीकृत्य सहसा तदीयमङ्कमाससाद । पुत्रो हि सर्वत्र पितुरङ्के निषीदति । तदानीं सूयवसपुत्रोऽजीगर्तो विश्वामित्रं प्रत्येवमुवाच । हे महर्षे ! मदीयपुत्रमेनं पुनरपि मह्यं देहीति स विश्वामित्रो नेति निराकृत्यैवमुवाच- “प्रजापत्यादयो देवा इमं शुनःशेपं मह्यम् अरासत=दत्तवन्तः, तस्मात् तुभ्यं न दास्यामीति ।” सः= शुनःशेपो देवैर्दत्तत्त्वाद् ‘देवरातः’ इति नामधारी विश्वामित्रपुत्र एव आस । तस्य देवरातस्य ये कापिलेयाः= कपिलगोत्रोत्‍पन्नाः, बाभ्रवाः= बभ्रुगोत्रोत्पन्नाश्च बन्धवोऽभवन् ।

\customparagraph{\textbf{नेपाली}}

यसपछि हरिश्चन्द्रसहित सम्पूर्ण ऋषिहरू शुनःशेप कसको छोरो हुने भन्ने कुरामा अनिर्णीत बनेपछि शुनःशेपले जसलाई पिता स्वीकार गर्छ, त्यही यसको पिता हुन्छ भन्ने सर्वसम्मत निर्णयमा पुग्नासाथ शुनःशेप विश्वामित्रलाई पिताका रूपमा स्वीकार गरी उनको काखमा बस्न पुगे । यो देखेर अजीगर्तले भने- “हे विश्वामित्र ! त्यो मेरो छोरो मलाई दिनुहोस् ।” विश्वामित्रले भने- हुन्न, देवताले मलाई दिएकाले यो मेरो छोरो भयो । यसरी ती शुनःशेप देवताले दिएकाले ‘देवरात’ नाम गरेका विश्वामित्रका छोरा भए । यिनीबाट कपिल गोत्रका र बभ्रु गोत्रका बान्धवहरू उत्पन्न भए ।

अथ शुनःशेपाजीगर्तयोः संवादं दर्शयति-

\begin{visheshshloka}[center]

\textbf{स होवाचाजीगर्तः सौयवसिस्त्वं वेहि विह्वयावहा स होवाचाजीगर्तः सौयवसिः -}

\textbf{आङ्गिरसो जन्मनाऽस्याजीगर्तिः श्रुतः कविः ।}

\textbf{ऋषे पैतामहात्तन्तोर्माऽपगाः पुनरेहि मामिति;}

\textbf{स होवाच शुनःशेपः -}

\textbf{अदर्शुस्त्वा शासहस्तं न यच्छूद्रेष्वलप्सत ।}

\textbf{गवां त्रीणि शतानि त्वमवृणीथा मदङ्गिरः ॥ इति ।}

\end{visheshshloka}

\customparagraph{संस्कृतव्याख्या}

विश्वामित्रेण निराकृतः सोऽजीगर्तः शुनःशेपं प्रत्येवमुवाच- हे पुत्र ! त्वं वा= त्वमेव विश्वामित्रादपगतः सन् एहि= अस्मद् गृहे आगच्छ। त्वदीयमाता चाहं चोभावावां विह्वयावहै= विशेषेण तवाऽऽह्वानं करवावहै । इत्युक्त्वा तूष्णीमवस्थितं तं शुनःशेपं प्रति पुनरपि कयाचिद् गाथयैवमुवाच- हे शुनःशेप ! त्वं जन्मना अङ्गिरसः= अङ्गिरोगोत्रोत्पन्नः, आजीगर्तिः= अजीगर्तस्य पुत्रः, कविः=विद्वान् श्रुतः= प्रसिद्धः, अतो हे ऋषे ! शुनःशेप ! पैतामहात्= पितामहेन प्रजापतिना सम्पादितात् तन्तोः= सन्तानात्, अङ्गिरोवंशात्, माऽपगाः= अपगतो मा भव । तस्मात् पुनरपि मामेहि= मद्‌गृहे समागच्छ, इत्युक्तः शुनःशेपः स्वयमपि गाथायाः प्रत्युत्तरं ददौ ।

शासः= विशसनहेतुः खड्गः, हेऽजीगर्त ! मद्वधाय शासहस्तम्= खड्‌गहस्तं त्वां सर्वेऽपि अदर्शुः= दृष्टवन्तः । यत्= क्रूरं कर्म शूद्रेषु= अत्यन्तनीचेष्वपि  नालप्सत= न लब्धवन्तो लौकिका जनाः, तादृशं क्रौर्यं कर्म त्वया कृतम् । हे अङ्गिर ! अङ्गिरगोत्रोत्पन्नाजीगर्त ! मत्= मां निमित्तीकृत्य त्वं गवां शतानि त्रीणि अवृणीथाः= वृतवानसि । तदिदमत्यन्तकष्टम्, इति उक्तवान् ।

\customparagraph{\textbf{नेपाली}}

विश्वामित्रबाट तिरस्कृत गरिएका सूयवसका पुत्र अजीगर्तले भने- हे पुत्र ! तिमी विश्वामित्रलाई छाडेर मसँग आऊ । हामी तिम्रा माता पिता दुबै तिमीलाई प्रेमपूर्वक बोलाउँछौं । यति भन्दा पनि चुप बसेका शुनःशेपलाई अजीगर्तले फेरि एक गाथा सुनाउँछन्- हे पुत्र ! तिमी जन्मले अङ्गिरा ऋषिका कुलमा उत्पन्न अङ्गिरस हौ । तिमी अजीगर्तका छोरा ऋषिका रूपमा प्रसिद्ध छौ । अतः हे पुत्र ! तिमी हजुरबुबाहरूले स्थापना गरेको अङ्गिरा वंशलाई नछोड । यसरी अजीगर्तले भनेपछि ती शुनःशेपले गाथारूपमै उत्तर दिए- हे तात ! मलाई मार्न तयार तपाईंका हातमा तरबार सबैले देखेका छन् । जुन कर्मलाई नीच शूद्रले पनि स्वीकार गर्दैन, अन्य लौकिक मानिसको त कुरै छाडौं । हे अङ्गिरा गोत्रमा उत्पन्न तात ! तपाईंले त मेरा सट्टामा तीन सय गाई पनि लिएको हो नि !

पुनरप्यजीगर्तशुनःशेपयोरुक्तिप्रत्युक्तीर्दर्शयति-

\begin{visheshshloka}[center]

\textbf{स होवाचाजीगर्तः सौयवसिः}

\textbf{तद्वै मा तात तपति पापं कर्म मया कृतम् ।}

\textbf{तदहं निह्नुवे तुभ्यं प्रतियन्तु शता गवामिति ॥ इति ।}

\textbf{स होवाच शुनःशेपः-}

\textbf{यः सकृत्पापकं कुर्यात् कुर्यादेनत्तत्तोऽपरम् ।}

\textbf{नापगाः शौद्रान्न्यायादसन्धेयं त्वया कृतमिति ॥ इति ।}

\end{visheshshloka}

\customparagraph{संस्कृतव्याख्या}

शुनःशेपेनोपालब्धोऽजीगर्तः स्वकीयमनुतापं दर्शयितुं काञ्चिद्गाथामेवमुवाच-  हे तात ! पितृवदुपलालनीय शुनःशेप ! मया यत्पापं कृतम्, तद्वै मा= माम्, तपति= सन्तापयति । अहं तत्पापं निह्नुवे= परिहरामि । गवां त्रीणि शतानि पूर्वं मया गृहीतानि, तानि तुभ्यम्= त्वदर्थम्, प्रतियन्तु= प्रत्येकमेव प्राप्नुवन्तु । इतरयोः पुत्रयोर्गावो मा भूवंस्तवैव सर्वाः सन्त्विति । ततः स शुनःशेपो गाथया प्रत्युत्तरं ददौ- यः= पुरुषः, धर्मशास्त्रभीतिरहितः सकृत्पापकं कुर्यात् स पुरुषस्ततः पापादन्यद् एनत्= पापम्, तदभ्यासवशात् कुर्यादेव । त्वं तु शौद्रान्न्यायात्= नीचजातिसम्बन्धिनः क्रूरादाचरणाद् नापागाः= अपगतो न भवसि । असन्धेयम्= प्रतिसन्धानरहितं पापं त्वया कृतमित्येषा शुनःशेपस्य प्रत्युक्तिः ॥

\customparagraph{\textbf{नेपाली}}

सुयवसका छोरा ती अजीगर्तले शुनःशेपलाई भने- “हे पितृवत् लालनपालन युक्त छोरा शुनःशेप ! मैले जुन पाप कर्म गरेको छु, त्यसले मेरो मनलाई अत्यन्तै सन्ताप दिइरहेको छ । म त्यो पापको निवारण गर्न चाहन्छु । जुन तीन सय गाई मैले लिएको छु त्यो सबै तिम्रै लागि हुनेछन् ।” यसरी अनुरोध गर्ने अजीगर्तलाई शुनःशेपले गाथाकै माध्यमबाट उत्तर दिए- “जो मानिस धर्मशास्त्रबाट भयभीत नभई एकपटक पाप गर्दछ, त्यो व्यक्ति अभ्यासवश अन्य पाप पनि गर्न सक्दछ। अतः नीच जातिसँग सम्बन्धित क्रूर कर्म गर्ने तपाईं त्यो कर्मबाट निवृत्त हुन सक्नुहुन्न किनकि तपाईंले केही नसोची क्रूर कार्य गर्नुभएको हो ।”

अथ विश्वामित्रकृत्यं दर्शयति-

\begin{visheshshloka}[center]

\textbf{असन्धेयमिति ह विश्वामित्र उपपपाद; स होवाच विश्वामित्रः-}

\textbf{भीम एव सौवयसिः शासेन विशिशासिषुः ।}

\textbf{अस्थान्मैतस्य पुत्रो भूर्ममैवोपेहि पुत्रतामिति ॥ इति ।}

\end{visheshshloka}

\customparagraph{संस्कृतव्याख्या}

असन्धेयम्= अप्रतिसमाधेयं पापमिति शुनःशेपेन यदीरितं तदेतदप्रतिसमाधेयत्वं विश्वामित्र उपपपाद= युक्तिभिरुपपादितवान्। तदुपपादनार्थमेव विश्वामित्र इमां गाथामुवाच- सुवयसस्य पुत्रोऽजीगर्तो भीम एव= भयहेतुरेव सञ्शासेन हस्तगतखड्‌गेन स विशिशासिषुरस्थात्= विशसनकर्तुमिच्छुरवस्थितवान् । अतो हे शुनःशेप ! एतस्य= पापिष्ठस्य पुत्रो मा भूः, ममैव पुत्रतामुपेहि ॥

\customparagraph{\textbf{नेपाली}}

शुनःशेपले भनेको ठिक छ कि यो अजीगर्तले गरेको पापको निवृत्ति हुन सक्दैन, भन्ने कुरालाई सही थप्दै विश्वामित्र भन्छन्- “जुन वेला हातमा खड्‍ग लिएर यो अजीगर्त तिमीलाई मार्न उद्यत भयो, त्यसवेला यो अत्यन्त भयानक देखिन्थ्यो । अतः हे शुनःशेप ! तिमी यो पापीको छोरो बन्न स्वीकार नगर, मेरै पुत्र बन ।”

अथ शुनःशेपविश्वामित्रयोरुक्तिप्रत्युक्ती दर्शयति-

\begin{visheshshloka}[center]

\textbf{स होवाच शुनःशेपः -}

\textbf{स वै यथा नो ज्ञपयाऽऽराजपुत्र तथा वद ।}

\textbf{यथैवाऽऽङ्गिरसः सन्नुपेयां तव पुत्रतामिति ॥}

\textbf{स होवाच विश्वामित्रः-}

\textbf{ज्येष्ठो मे त्वं पुत्राणां स्यास्तव श्रेष्ठा प्रजा स्यात् ।}

\textbf{उपेया दैवं मे दायं तेन वै त्वोपमन्त्रये इति ॥ इति ।}

\end{visheshshloka}

\customparagraph{संस्कृतव्याख्या}

विश्वामित्रेणैवं बोधितः शुनःशेपः पुनरपि गाथया विश्वामित्रं प्रत्येवमुवाच- अयं विश्वामित्रो जन्मना क्षत्त्रियः सन् स्वकीयेन तपोमहिम्ना ब्राह्मण्यं प्राप्तवानित्येवं तद्‌वृत्तान्तं सूचयितुं हे राजपुत्र ! इति सम्बोधितवान् । स वै= तथाविधो राजजातीय एव सन् यथा= येन प्रकारेण नः= अस्माभिः सर्वैरासमन्ताज्ज्ञपय= ब्राह्मणत्वेन ज्ञायसे तथैवास्मद्विषयेऽपि त्वं वद- कथं वदितव्यमिति । तदुच्यते- अहमिदानीमङ्गिरोगोत्रः संस्तत्परित्यागेन तव पुत्रत्वं येनैव प्रकारेणोपेयां तथैवानुगृहाणेति शेषः । एतद्वाक्याभिप्रायः पूर्वे सङ्‌क्षिप्य दर्शितः -

\begin{shloka}[center]

पुराऽत्मानं नृपं विप्रं तपसा कृतवानसि ।

एवमाङ्गिरसं मा त्वं वैश्वामित्रमृषे कुरु ॥ इति ।

\end{shloka}

\customparagraph{संस्कृतव्याख्या}

ततो विश्वामित्रो गाथया प्रत्युत्तरमुक्तवान्- हे शुनःशेप ! त्वं मे पुत्राणां मध्ये ज्येष्ठः स्याः= ज्येष्ठो भव । तव पुत्रादिरूपा प्रजाऽपीतरस्याः श्रेष्ठा स्यात् । मे= मह्यं विश्वामित्राय दैवम्= देवैः प्रसन्नैर्दत्तं दायम्= पुत्रत्वरूपं लाभम्, उपेयाः= प्राप्नुहि । तेन वै= तेनैव प्रकारेण त्वोपमन्त्रये= त्वां पुत्रत्वेन व्यवहरामि ॥

\customparagraph{\textbf{नेपाली}}

हे राजपुत्र ! (विश्वामित्र) भन्नुहोस् कि तपाईंले भनेजस्तो अर्थात् तपाईंले क्षत्रिय भएर जसरी ब्राह्मणत्व प्राप्त गर्नुभयो त्यसरी म अङ्गिरा गोत्रमा उत्पन्न भएर कसरी तपाईंको पुत्रत्व प्राप्त गर्न सक्छु ? यसपछि विश्वामित्रले भने- हे शुनःशेप ! तिमी मेरा पुत्रहरूमध्ये ज्येष्ठ हुनेछौ । तिम्रा सन्तान रूप प्रजा पनि श्रेष्ठ हुनेछन्। म विश्वामित्रका लागि देवताहरूद्वारा प्रसन्नतापूर्वक दिइएका तिमीले मेरो उत्तराधिरकार प्राप्त गर्नेछौ । त्यसरी नै म पुत्रत्वको रूपमा बोलाउने गर्छु ।

पुनरपि शुनःशेपविश्वामित्रयोर्वृत्तान्तं दर्शयति-

\begin{visheshshloka}[center]

\textbf{स होवाच शुनःशेपः-}

\textbf{संज्ञानानेषु वै ब्रूयात् सौहार्द्याय मे श्रियै ।}

\textbf{यथाऽहं भरतऋषभोपेयां तव पुत्रतामिति ॥}

\textbf{अथ ह विश्वामित्रः पुत्रानामन्त्रयामास-}

\textbf{मधुच्छन्दाः शृणोतन ऋषभो रेणुरष्टकः ।}

\textbf{ये के च भ्रातरः स्थ नास्मै ज्यैष्ठ्याय कल्पध्वमिति ॥ इति ।}

\end{visheshshloka}

\customparagraph{संस्कृतव्याख्या}

विश्वामित्रेण प्रलोभितः शुनःशेपः स्वकार्यदार्ढ्यार्थं गाथायैवमुवाच- संज्ञानानेषु मद्विषयैकमत्यं प्राप्तेषु त्वदीयपुत्रेषु सर्वोऽपि मां ब्रूयात्, ज्येष्ठभ्रातृत्वेन व्यवहरतु । तच्च मे= मम सौहार्द्याय=भ्रातृभिरितरैः स्नेहातिशयाय श्रीयै= धनलाभाय च सम्पद्यते । हे भरतऋषभः‍= भरतवंशश्रेष्ठो विश्वामित्र ! अहं तव पुत्रतां यथोपेयां तथैवैतेषां पुत्राणामग्रेऽनुगृहाणेति शेषः । ततो विश्वामित्र इतरान् पुत्रानाहूय गाथायैवमाज्ञापितवान्- ये मधुच्छन्दा नाम, यश्चर्षभः, योऽपि रेणुः, योऽप्यष्टकः, एते मुख्याः, हे पुत्रकाः ! शृणोतन= मदीयामाज्ञां शृणुत । ये केचिद् यूयं सर्वे भ्रातरः स्थ ते सर्वेऽपि अस्मै= शुनःशेपादस्मादपि ज्येष्ठाय= श्रेष्ठाय न कल्पध्वम्= विचारमपि मा कुरुत, असौ शुनःशेप एव युष्माकं मध्ये ज्येष्ठो भूत्वाऽवतिष्ठतामिति ।

\customparagraph{\textbf{नेपाली}}

विश्वामित्रद्वारा आश्वस्त शुनःशेप फेरि भन्छन्- “मेरो विषयमा भएका सबै कुरा तपाईंका पुत्रहरूलाई अवगत गराउनुहोस् । तिनले मलाई ज्येष्ठ भ्राताका रूपमा मानून् र आपसी प्रेम र धनका लागि पनि सहमत होऊन् । उनीहरूले यो कुरा मेरै अगाडि स्वीकार गरून्।”  शुनःशेपका कुरा सुनेर विश्वामित्रले आफ्ना सम्पूर्ण छोराहरूलाई बोलाए- हे मधुच्छन्दा, ऋषभ, रेणु, अष्टक आदि मुख्य रहेका छोराहरू हो ! मेरो आज्ञा सुन । आज उप्रान्त तिमीहरू सबैले यी शुनःशेपलाई ज्येष्ठ मान्ने कुराभन्दा अन्य केही नसोच्नू किनकि अब यी शुनःशेप नै तिमीहरूका ज्येष्ठ भ्राता हुन्।

\begin{visheshshlokatwo}[center]{84}

\textbf{ऐतेरियब्राह्मणको  तेत्तिसौं अध्याय (सप्तपञ्चिकाको तेस्रो अध्यायको पञ्चम खण्ड)}

\textbf{नेपाली भावानुवाद पूर्ण भयो ।}

\end{visheshshlokatwo}

\specialupakhanda{\textbf{षष्ठः खण्डः}}

अथ विश्वामित्रपुत्राणां वृत्तान्तमाह-

\begin{visheshshloka}[center]

\textbf{तस्य ह विश्वामित्रस्यैकशतं पुत्रा आसुः पञ्चाशदेव ज्यायांसो मधुच्छन्दसः पञ्चाशत्कनीयांसः ॥ इति ।}

\end{visheshshloka}

\customparagraph{संस्कृतव्याख्या}

विश्वामित्रस्य मधुच्छन्दोनामकः कश्चित्पुत्रो मध्यमस्ततो ज्येष्ठाः पञ्चाशत्, कनिष्ठाश्च पञ्चाशत् पुत्रा आसन् ।

\customparagraph{\textbf{नेपाली भावार्थ}}

विश्वामित्रका मध्यम छोरा मधुच्छन्द नामका थिए । तीभन्दा जेठा पचास र तीभन्दा कान्छा पचास थिए ।

तेषां मध्ये ज्येष्ठानां वृत्तान्तमाह-

\begin{visheshshloka}[center]

\textbf{तद्ये ज्यायांसो न ते कुशलं मेनिरे, ताननु व्याजहारान्तान् वः प्रजा भक्षीष्टेति, त एतेऽन्ध्राः पुण्ड्राः शबरा पुलिन्दा मूतिबा इत्युदन्त्या बहवो वैश्वामित्रा दस्यूनां भूयिष्ठाः ॥ इति ।}

\end{visheshshloka}

\customparagraph{संस्कृतव्याख्या}

तत्= तेष्वेकशतसङ्‌ख्याकेषु पुत्रेषु मध्ये ये ज्यायांसः= मधुच्छन्दसो ज्येष्ठाः पञ्चाशत्सङ्‌ख्याकाः सन्ति, ते शुनःशेपस्य विश्वामित्रपुत्रत्वं कुशलं न मेनिरे= इदं समीचिनमिति नाङ्‌गीकृतवन्तः । तान्= पञ्चादशसङ्‌ख्यकान् ज्येष्ठाननुलक्ष्य विश्वामित्रो व्याजहार= व्याहरणं  शापरूपं वाक्यमुक्तवान्- हे ज्येष्ठपुत्रा ! वः=युष्माकं मदीयाज्ञातिलङ्‌घिनां प्रजाः= पुत्रादिकाः, अन्तान्भक्षीष्ट= चाण्डालादिरूपान्नीचजातिविशेषान् भजतामिति । त एते शप्ताः सन्तोऽन्ध्रत्वादिपञ्चविधनीचजातिविशेषा भवन्ति । ‘इति’ शब्दस्य तत्प्रदर्शनार्थत्वादन्येऽपि नीचजातिविशेषाः सर्वे विवक्षिताः। उद्गतोऽन्त उदन्तोऽत्यन्तनीचजातिस्तत्र भवा ‘उदन्त्याः’। ते बहवः= अनेकविधा वैश्वामित्राः= विश्वामित्रसन्ततिजाः, दस्यूनाम्= चौराणाम्, मध्ये भूयिष्ठाः‍= समधिकाः सञ्जाताः ।

\customparagraph{\textbf{नेपाली}}

ती सय पुत्रमध्ये जति ‘मधुच्छन्द’ नामक पुत्रभन्दा जेठा थिए, ती शुनःशेपलाई दाजु मान्न तयार भएनन् । ती ज्येष्ठ पचासलाई विश्वामित्रले यस्तो सराप दिए- “हे ज्येष्ठपुत्रहरू हो ! तिमीहरूले मेरो आज्ञा उल्लङ्‌घन गर्‍यौ, त्यसकारण तिमीहरूका सन्तान चाण्डाल आदि गरिएका तुच्छ जातिका हुनेछन् ।” यसरी विश्वामित्रले सराप दिएपछि ती पचास भाइ छोराका सन्तान आन्ध्र पुण्ड्र शबर पुलिन्द मूतिब आदि नाम गरिएका चाण्डालजस्ता तुच्छ चोर जातिका भए । उक्त जातिका चोर, लुटाहा, फटाहाहरूमा विश्वामित्रका सन्तानहरूको बाहुल्य छ ।

ज्येष्ठानां वृत्तान्तमुक्त्वा मधुच्छन्दसस्तत्कनिष्ठानां च वृत्तान्तं गाथया दर्शयति-

\begin{visheshshloka}[center]

\textbf{स होवाच मधुच्छन्दाः पञ्चाशता सार्धम्-}

\textbf{यन्न पिता सञ्जानीते तस्मिंस्तिष्ठामहे वयम् ।}

\textbf{पुरस्त्वा सर्वे कुर्महे त्वामन्वञ्चो वयं स्मसीति ॥ इति ।}

\end{visheshshloka}

\customparagraph{संस्कृतव्याख्या}

कनिष्ठपुत्राणां पञ्चशता सह मधुच्छन्दोनामकः स मध्यमः पुत्रः शुनःशेपं प्रत्येवमुवाच- “हे शुनःशेप ! अस्माकं पिता विश्वामित्रो यत्कार्यं त्वदीयज्येष्ठपुत्रत्वरूपं सञ्जानीते= सम्यग्जानात्यङ्‌गीकरोति तस्मिन् कार्ये वयं तिष्ठामहे= तत्कार्यमङ्‌गीकुर्मः । सर्वे वयं त्वा= शुनःशेपनामानं त्वाम्, पुरस्कुर्महे= पुरस्कृत्य ज्येष्ठं कृत्वा वर्तामहे । त्वामन्वञ्च= शुनःशेपमनुगच्छन्तः स्मसि= भवाम”  इत्युक्तवान् ।

\customparagraph{\textbf{नेपाली}}

अन्य पचास कान्छा भाइहरूको नेतृत्व गर्दै मधुच्छन्दले भने-  हे शुनःशेप ! हाम्रा पिता विश्वामित्रले तपाईंलाई ज्येष्ठपुत्र मान्ने जुन कार्य राम्ररी विचार गरेर अङ्‌गीकार गर्नुभएको छ, त्यसलाई हामी पचास भाइ स्वीकार गर्दछौं । हामी तपाईंलाई जेठा दाजुका रूपमा अगाडि लगाएर हिँड्न तयार छौं । सर्वदा तपाईंकै अनुसरण गर्नेछौँ ।

विश्वामित्रवृत्तान्तं दर्शयति-

\begin{visheshshloka}[center]

\textbf{अथ ह विश्वामित्रः प्रतीतः पुत्रांस्तुष्टावः ॥ इति ।}

\end{visheshshloka}

अथ= मधुच्छन्दःसहितानां पञ्चाशत्कनिष्ठपुत्राणां शुनःशेपविषयज्येष्ठपुत्रत्वाङ्‌गी- कारानन्तरम्, विश्वामित्रः प्रतीतः= तेषु प्रत्ययं प्राप्तः, मदनुकूला इति विश्वस्य प्रीत्या एताभिर्गाथाभिस्तान् पुत्रान् त‍ुष्ठाव-

\customparagraph{\textbf{नेपाली}}

यसरी मधुच्छन्दसहित पचास पुत्रले आफ्नो कुरा स्वीकार गरेकाले खुसी भएर निम्नानुसारका गाथाका माध्यमबाट विश्वामित्रले तिनीहरूको प्रशंसा गरे-

प्रथमां गाथामाह-

\begin{visheshshloka}[center]

\textbf{ते वै पुत्राः पशुमन्तो वीरवन्तो भविष्यथ ।}

\textbf{ये मानं मेऽनुगृह्णन्तो वीरवन्तमकर्त मा ॥ इति ।}

\end{visheshshloka}

\customparagraph{संस्कृतव्याख्या}

हे मधुच्छन्दःप्रमुखाः कनिष्ठपुत्राः ! ये=यूयम्, मानम्= मदीयं मतम्, अनुगृह्णन्तः= आनुकूल्येन स्वीकृतवन्तः, माम्= विश्वामित्रम्, वीरवन्तम्= स्वधर्मशूरपुत्रयुतम्, अकर्त= कृतवन्तः । ते वै= तादृशा यूयं पुत्राः, पशुमन्तः= बहुविधपशुधनयुताः, वीरमन्तः= बहुविधानुकूलपुत्रयुक्ताश्च भविष्यथ ।

\customparagraph{नेपाली}

हे मधुच्छन्द आदि गरिएका पाचास कनिष्ठ छोरा हो ! तिमीहरूले मैले भनेको कुरा स्वीकार गरेर मलाई आफ्नो धर्ममा स्थिर वीरपुत्रहरूका पिताका रूपमा प्रसिद्ध बनाएका छौ । म तिमीहरूलाई यो आशीर्वाद दिन्छु कि तिमीहरू अनेक प्रकारका पशुधनले सम्पन्न र विभिन्न गुणयुक्त सन्तानले पनि पूर्ण हुनेछौ ।

द्वितीयां गाथामाह-

\begin{visheshshloka}[center]

\textbf{पुर एत्रा वीरवन्तो देवरातेन गाथिनाः ।}

\textbf{सर्वे राध्याः स्थ पुत्रा एष वः सद्विवाचनम् ॥ इति ।}

\end{visheshshloka}

\customparagraph{संस्कृतव्याख्या}

गाथिनशब्दो विश्वामित्रस्य पितरमाचष्टे । हे गाथिनाः, गाथिपौत्राः ! पुर एत्रा= युष्माकं पुरतो गन्त्रा मुख्येन देवरातेन (शुनःशेपेन) सह यूयं सर्वे वीरवन्तः=  श्रेष्ठपुत्रयुक्ताः, राध्याः स्थ= सर्वैः पुरुषैराराधनीयाः पूज्या भवथ । हे पुत्राः ! मधुच्छन्दप्रभृतयः, एष देवरातो वः=युष्माकम्, सद्विवाचनम्= सन्मार्गस्य विशेषतोऽध्यापनं करिष्यतीति शेषः ॥

\customparagraph{\textbf{नेपाली}}

हे गाथीका नातिहरू हो ! देवरात अर्थात् शुनःशेपले निर्धारण गरेको असल मार्गमा हिँड्ने भएकाले तिमीहरू सबै श्रेष्ठपुत्रले सम्पन्न भएर सबै मानिसहरूका आराध्य हुनेछौ । हे मधुच्छन्द आदि गरिएका कनिष्ठ पुत्र हो ! यी देवरातले तिमीहरूलाई विशेष गरेर सन्मार्गको ज्ञान गराउनेछन् ।

तृतीयां गाथामाह-

\begin{visheshshloka}[center]

\textbf{एष वः कुशिका वीरो देवरातस्तमन्वित ।}

\textbf{युष्मांश्च दायं म उपेता विद्यां यामु च विद्मसि ॥ इति ।}

\end{visheshshloka}

\customparagraph{संस्कृतव्याख्या}

हे कुशिकाः, कुशिकनाम्नो मत्पितामहस्य सम्बन्धिनो मधुच्छन्दःप्रभृतय !  एष देवरातो वः=युष्माकं ज्येष्ठभ्राता, तम्=देवरातम्, यूयमन्वित= अनुगता भवत । मे= मदीयम्, दायम्= धनम्, युष्मांश्च उपेता= प्राप्स्यति, चकाराद्देवरातं च । यामु च= यामपि कांश्चिद् वेदशास्त्रादिरूपां विद्यां विद्मसि= वयं जानीमः, साऽपि विद्या युष्मानुपेता=प्राप्स्यति ।

\customparagraph{\textbf{नेपाली}}

हे कुशिक नामका मेरा पितामहका सन्तानहरू हो ! यी तिमीहरूका ज्येष्ठभ्राता देवरात जो छन्, यिनको तिमीहरू सबैले अनुगमन गर्नू । मेरो जति सम्पत्ति छ, त्यो सबै तिमीहरूले पाउनेछौं । (देवरातले पनि पाउनेछन्) । म जति विद्या जान्दछु त्यो पनि तिमीहरूले (देवरातले समेत) पाउनेछौ ।

चतुर्थीं गाथामाह-

\begin{visheshshloka}[center]

\textbf{ते सम्यञ्चो वैश्वामित्राः सर्वे साकं सरातयः ।}

\textbf{देवराताय तस्थिरे धृत्यै श्रैष्ठ्याय गाथिनाः ॥ इति ।}

\end{visheshshloka}

\customparagraph{संस्कृतव्याख्या}

हे वैश्वामित्राः= मम पुत्राः, ये गाथिनाः= गाथिपौत्रास्ते यूयं सर्वेऽपि सम्यञ्चः= समीचीनबुद्धयो ये साकम्= देवरातेन सार्धम्, सरातयः= रातिर्धनसम्पत्तिस्तया युक्ताः सन्तो देवराताय= मदीयज्येष्ठपुत्रस्य देवरातस्य धृत्यै= धारणं युष्मत्पोषणम्, श्रैष्ठ्याय= युष्माकं मध्ये श्रेष्ठत्वं च तस्थिरे= अङ्गीकृतवन्तः ।

\customparagraph{\textbf{नेपाली}}

हे म विश्वामित्रका छोरा एवम् गाथीका नातिहरू हो ! तिमीहरू सबैले समीचीन बुद्धि भएकाहरूले देवरातका साथ धनसम्पत्ति प्राप्त गर्नेछौ किनभने तिमीहरूले मेरो श्रेष्ठपुत्र देवरातको संरक्षकत्व (लालनपालन) र आफ्नो श्रेष्ठत्व पनि  स्वीकार गरेका छौ ।

पञ्चमीं गाथामाह-

\begin{shloka}[center]

\textbf{अधीयत देवरातो रिक्थयोरुभयोर्ऋषिः ।}

\textbf{जह्‌नूनां चाऽऽधिपत्ये दैवे वेदे च गाथिनाम् ॥ इति ।}

\end{shloka}

\customparagraph{संस्कृतव्याख्या}

अधियत= स्मृतिकारैर्महर्षिभिः स्मर्यते । कथमिति तदुच्यते- अयं देवरातो द्व्यामुष्यायणत्वादुभयोरजीगर्तविश्वामित्रयोः सम्बन्धिनी ये रिक्थे= धने तयोर्ऋषिर्द्रष्टा तदुभयमर्हतीत्यर्थः । अजीगर्तस्य कूटस्थऋषिर्जह्‌नुसंज्ञकस्तस्य वंशे जाताः सर्वे जह्नवस्तेषां चाऽऽधिपत्ये स्वामित्वे देवरातो योग्यः । तथा दैवे= देवसम्बन्धिनि यागादिकर्मणि वेदे च मन्त्ररूपे समर्थः । गाथिनामस्मत्पितृवंशोत्पन्नानां च सर्वेषामाधिपत्ये योग्योऽस्ति देवरातः।

\customparagraph{\textbf{नेपाली}}

यसरी स्मृतिकारहरूद्वारा भनिन्छ- कि देवरात अजीगर्त र विश्वामित्र दुबैका धनको (ऋषिपरम्पराका) उत्तराधिकारी हुन्छन् । जह्‌नु ऋषिका वंशमा उत्पन्न हुनेहरूको स्वामित्वका लागि र गाथी वंशमा उत्पन्न हुनेहरूको विद्यारूप स्वामित्वका लागि यी देवरात योग्य छन् । साथै  देवसम्बन्धी यज्ञयागादि कर्म र मन्त्ररूप वेदका लागि पनि समर्थ छन् ।

उक्तमुपाख्यानमुपसंहरति-

\begin{visheshshloka}[center]

\textbf{तदेतत्परऋक्शतगाथां शौनःशेपमाख्यानम् ॥ इति ।}

\end{visheshshloka}

\customparagraph{संस्कृतव्याख्या}

‘कस्य नूनं–’ निधारयेत्यन्ताः सप्ताधिकनवतिसङ्‌ख्याका ऋचः, ‘त्वं न–’ ‘स त्वम्—’ इत्यादिकास्तिस्र ऋचः, एवमृचां शतम् । पूर्वोक्तादृक्शतात् परोऽधिका एकत्रिंशत्सङ्‌ख्याका ‘यं न्विमम्–’ इत्याद्या गाथाः, यस्मिन्नाख्याने तदेतत्= परऋक्शतगाथाम् ।

(शुनःशेपेन दृष्टाः सप्ताधिकनवतिसङ्‌ख्यायुता ऋचो याः सन्ति, अन्येन दृष्टास्तिस्र ऋचो याः सन्ति, ब्राह्मणे प्रोक्ता एकविंशद्‌गाथाविशेषा या सन्ति, तैः सर्वैरुपेतं ‘हरिश्चन्द्रो ह वैधसः’ इत्यादिकं सर्वं शुनःशेपविषयमाख्यानम्)

\customparagraph{\textbf{नेपाली}}

यस उपाख्यानमा ‘कस्य नूनं—’ भन्ने मन्त्रदेखि लिएर ‘निधारय०’ सम्म शुनःशेपद्वारा देखिएका सन्तानब्बे ऋचा छन् । ‘त्वं न’ र ‘स त्वं’ आदि अन्य ऋषिहरूद्वारा देखिएका तीन ऋचा छन् । यसरी समग्रमा सय ऋचा प्रयोग गएका छन् । यसअतिरिक्त ‘य न्विमं–’ आदि एकतीस गाथाहरू पनि यहाँ रहेका छन् । यी सबै मिलाएर शुनःशेपको आख्यान सय ऋचा र एकतीस गाथामा वर्णन गरिएको छ ।

तस्याख्यानस्य राजसूयक्रतौ विनियोगं दर्शयति-

\begin{visheshshloka}[center]

\textbf{तद्धोता राज्ञेऽभिषिक्तायाऽऽचष्टे ॥ इति ।}

\end{visheshshloka}

\customparagraph{संस्कृतव्याख्या}

राजसूयक्रतावभिषेचनीयाख्ये कर्मणि यदा राजाभिषिक्तो भवति, तदानीं तस्मै राज्ञे तत्= आख्यानम्, होता कथयेत् ।

\customparagraph{\textbf{नेपाली}}

राजसूय यज्ञमा राजालाई अभिषेक गर्ने समयमा होताले यो आख्यान राजालाई सुनाऊन्।

तत्रेति कर्तव्यं दर्शयति-

\begin{visheshshloka}[center]

\textbf{हिरण्यकशिपावासीन आचष्टे; हिरण्यकशिपावासीनः प्रतिगृह्णाति । यशो वै हिरण्यं यशसैवैनं तत्समर्धयति ॥ इति ।}

\end{visheshshloka}

\customparagraph{संस्कृतव्याख्या}

होता यदोपाख्यानं कथयति, तदानीं हिरण्यकशिपौ= सुवर्णनिर्मितसूत्रैर्निष्पादिते कशिपौ= आसने स होतोपविशेत् । तदाख्यानमध्येऽध्वर्युश्च हिरण्यकशिपावासीनो वक्ष्यमाणं प्रतिगरं ब्रूयात् । हिरण्यस्य यशोहेतुत्वात् यशस्त्वम् । तथा सति एनम्= राजानं  यशसैव समृद्धं करोति ।

\customparagraph{\textbf{नेपाली}}

होता जब यो आख्यान भन्छन् तब तिनले सुवर्णमय सूत्रबाट निर्मित आसनमा बसेर आख्यान भन्नुपर्छ । त्यस आख्यानका बीचमा अध्वर्यु पनि सुवर्णमय सूत्रबाट बनेको आसनमा बसेर प्रतिगरण गर्दछन्। वस्तुतः सुवर्ण नै यश हो । त्यसैले यसरी राजालाई यशले नै समृद्ध बनाउँछ।

अध्वर्युणा प्रयोक्तव्यं प्रतिगरविशेषं दर्शयति-

\begin{visheshshloka}[center]

\textbf{ओमित्युचः प्रतिगर; एवं तथेति गाथायाः; ओमिति वै देवम्, तथेति मानुषम्, दैवेन चैवैनं तन्मानुषेण च पापादेनसः प्रमुञ्चति ॥ इति ।}

\end{visheshshloka}

\customparagraph{संस्कृतव्याख्या}

होत्रा प्रयुक्ताया एकैकस्या ऋचोऽन्तेऽध्वर्योरोमित्यतादृशः प्रतिगरो भवति । ओमित्येतच्छन्दोरूपं दैवं देवैरङ्गीकारार्थे प्रयुज्यते । तत्तथेत्यन्तं (तथेति) मानुषम्= मनुष्या अङ्गीकारे तथेतिशब्दं प्रयुञ्जते । तत्= तेन प्रतिगरेण दैवेन मानुषेण वाध्वर्युः, एनम्= राजानम्, पापात्= ऐहिकादपकीर्तिरूपात्, एनसः= नरकहेतोश्च प्रमुञ्चति= प्रमुक्तं करोति ।

\customparagraph{\textbf{नेपाली}}

होताद्वारा प्रयुक्त एक एक ऋचाका अन्त्यमा अध्वर्युबाट ‘ओेम्’ यसरी प्रतिगर (प्रत्युच्चारण) हुन्छ । यसरी नै गाथाको प्रयोग गर्दा ‘तथा’ ‘यस्तै नै होस्’ यो प्रतिगर हुन्छ । वस्तुतः ‘ओम्’ यो दैवी हो। देवताहरूलाई अङ्गीकरणका लागि प्रयुक्त हुन्छ । साथै ‘तथा’ मनुष्यलाई अङ्गीकरणका लागि प्रयुक्त हुन्छ । त्यसैले ‘तथा’ मानुषी हो। यसरी त्यो दैवी र मानुषी प्रतिगरद्वारा अध्वर्यु यी राजालाई अपकीर्तिरूप पापबाट र नरकरूपी आमुष्मिक पापबाट प्रकृष्ट रूपले मुक्त गर्दछन् ।

क्रत्वर्थत्वेनोपाख्यानं विधाय क्रतुनिरपेक्षपुरुषार्थत्वेन विधत्ते-

\begin{visheshshloka}[center]

\textbf{तस्माद्यो राजा विजिती स्यादप्ययजमान आख्यापयेतैवैतच्छौनःशेपमाख्यानम्, न हास्मिन्नल्पं चनैनः परिशिष्यते ॥ इति ।}

\end{visheshshloka}

\customparagraph{संस्कृतव्याख्या}

यस्मादुपाख्यानं पापप्रशमनहेतुः, तस्माद् अयजमानोऽपि= राजसूयक्रतुरहितोऽपि राजा विजिती= यदि विजयोपेतः स्यात्, तदानीमेतच्छौनःशेपमाख्यानम्, आख्यापयेत्= तस्य राज्ञो यः कश्चिद् ब्राह्मण उपाख्यानं ब्रूयात् । तथा सति तस्मिन्= राजनि अल्पं चन= किञ्चिदपि एनः= पापम्, न ह वै परिशिष्यते= सर्वं पापं नश्यतीत्यर्थः ।

\customparagraph{\textbf{नेपाली}}

किनकि यो आख्यान पापनाशक छ । त्यसैले राजसूय यज्ञले रहित अयजमान जुन कुनै राजा यदि विजयी भयो भने त्यसलाई पनि यो शुनःशेप आख्यान कुनै ब्राह्मणका मुखबाट सुनाउनुपर्छ । यसरी सुनाए त्यस राजामा लेशमात्र पनि पाप रहँदैन । अर्थात् राजा सबै पापले मुक्त हुन्छ ।

आख्यानप्रयुक्तां दक्षिणां विधत्ते-

\begin{visheshshloka}[center]

\textbf{सहस्रमाख्यात्रे दद्याच्छतं प्रतिगरित्र एते चैवाऽऽसने श्वेतश्चाश्वतरीरथो होतुः ॥ इति ।}

\end{visheshshloka}

\customparagraph{संस्कृतव्याख्या}

योऽयमाख्याता होता विद्यते, तस्मै क्रत्वर्थदक्षिणामन्तरेणोपाख्यानप्रयुक्तां दक्षिणां गोसहस्ररूपां दद्यात्। प्रतिगरित्रे= अध्वर्यवे गोशतं दद्यात् । ये हिरण्यकशिपुरूपे द्वे आसने स्तः, एते अपि ताभ्यामेव दद्यात् । अश्वतरीभ्यां सङ्कीर्णज्ञातियुक्ताभ्यामधिकशक्तिभ्यां युक्तो रथः= अश्वतरीरथः, स च रजतेनाऽलङ्‌कृतत्त्वाच्छेतः । सोऽपि तादृशो होतुर्देयः ।

\customparagraph{\textbf{नेपाली}}

ती वाचक होताका लागि दक्षिणाका रूपमा सहस्र गाई दिनुपर्छ । साथै प्रतिगर गर्ने अध्वर्युलाई सय गाई दिनुपर्छ । तिनै होता र अध्वर्युलाई ती सुवर्णनिर्मित आसन पनि दिनुपर्छ । यसअतिरिक्त होतालाई चाँदीले सजाइएको सेतो अत्यन्त शक्तिशाली दुई घोडीले युक्त रथ पनि दिनुपर्छ ।

\begin{visheshshloka}[center]

पूर्वमयज्ञसंयुक्तः कल्पो विजयिनो राज्ञोऽभिहितः । इदानीं पुत्रकामानामपि सर्वेषामपि विधत्ते-

\textbf{पुत्रकामाः हाप्याख्यापयेरल्लँभन्ते ह पुत्राल्लँभन्ते ह पुत्रान्॥ इति ।}

\end{visheshshloka}

\customparagraph{संस्कृतव्याख्या}

ये पुत्रकामाः सन्ति, ते  ह= प्रसिद्धमेतदुपाख्यानम् आख्यापयेरन्= ब्राह्मणमुखाच्छृणुयुरित्यर्थः । ते पुत्राल्लँभन्ते । अभ्यासोऽध्यायसमाप्त्यर्थः ।

\customparagraph{\textbf{नेपाली}}

जो पुत्रको कामना राख्छन्, तिनले पनि शुनःशेपको कथा ब्राह्मणका मुखबाट सुन्नुपर्छ । यसबाट तिनलाई पुत्रप्राप्ति हुन्छ ।

\begin{visheshshlokatwo}[center]{84}

\textbf{ऐतेरियब्राह्मणको  तेत्तिसौं अध्याय (सप्तपञ्चिकाको तेस्रो अध्यायको छैटौँ खण्ड)}

\textbf{नेपाली भावानुवाद पूर्ण भयो ।}

\end{visheshshlokatwo}

\specialkhanda{\textbf{पुरूरवोर्वश्याख्यानम्}}

\begin{tcolorbox}[title=, colback=theme!5!white, colframe=theme!75!black]

चतुर्दशकाण्डात्मकस्य शतपथब्राह्मणस्य माध्यन्दिनशाखाया एकादशकाण्डस्य  पञ्चमाध्यायस्य प्रथमे ब्राह्मणे वर्णितं पुरुरवोर्वश्याख्यानम्  (तृतीयप्रपाठकस्य तृतीयं ब्राह्मणम् )

\end{tcolorbox}

\customparagraph{\textbf{शतपथब्राह्मणस्य परिचयः}}

ब्राह्मणग्रन्थेषु शतपथब्राह्मणस्य वैशिष्ट्यं सातिशयं विद्यते । शुक्लयजुर्वेदस्य ब्राह्मणग्रन्थस्य नाम शतपथब्राह्मणमस्ति । शतपथब्राह्मणनाम्ना द्वे ब्राह्मणे प्राप्येते । तयोर्नाम माध्यन्दिनशतपथब्राह्मणं काण्वशतपथब्राह्मणञ्च विद्यते । अनयोर्विषयसाम्येऽपि कण्डिकादिसङ्‌ख्यायामन्तरमस्ति । माध्यन्दिनशतपथब्राह्मणे १४ काण्डानि ६८ प्रपाठकाः १०० अध्यायाः ४३८ ब्राह्मणानि ६८०६ कण्डिकाश्च (कुत्रचित् ७६२४ कण्डिका उल्लिखिताः प्राप्यन्ते) सन्ति । काण्वशतपथे १७ काण्डानि १०४ अध्यायाः ४३५ ब्राह्मणानि ६८०६ कण्डिकाश्च विद्यन्ते । काण्वशतपथे प्रपाठकविभाजनं नोपलभ्यते ।

शताध्यायसमन्वितत्त्वादस्य ग्रन्थस्य नाम ‘शतपथम्’ इत्यभवत् । ब्राह्मणग्रन्थेषु शतपथब्राह्मणं प्राचीनतमं मन्यते । अयं ग्रन्थः प्रायशः समग्रयज्ञयागानां साङ्गोपाङ्गं विवरणं प्रस्तौति । वर्ण्यविषयदृष्ट्या शतपथब्राह्मणं सम्पूर्णब्राह्मणसाहित्यस्य शिरोभूषणं मन्यते । अत्र यागवर्णनप्रसङ्गे यत्र तत्र वर्णिता आख्यायिका यज्ञविधिनियमादिशुष्कवर्णनेष्वपि सरसतामाविष्कुवन्ति । यज्ञीयक्रियाणां विधिनियमादिनिर्देशः, प्राचीनाख्यानानां सरसविवेचनमस्य ब्राह्मणस्योत्कर्षं प्रकाशयति । अत्र समाविष्टा रसगर्भिता आख्यायिका लौकिकसाहित्ये गद्यकाव्यविकासदृष्ट्या बीजभूता मन्यन्ते । पुराणेषु वर्णितस्य मत्स्यावतारस्य बीजमस्यामेवाख्यायिकायामस्ति । एवम्प्रकारेण रमणीया नैका आख्यायिका शतपथब्राह्मणे समुपलभ्यन्ते । या लौकिकगद्यसाहित्यस्य विकासे पर्याधारत्वेन महत्त्वाधायिका विद्यन्ते । बृहत्कायात्मकस्य शतपथब्राह्मणस्यान्तिमभागे बृहदारण्यकोपनिषद् निबद्धा विद्यते ।

शतपथब्राह्मणे वर्णितासु कथासु पुरुरवउर्वश्याख्यानं प्रसिद्धं विद्यते । अत्र नायकः पुरुरवसः स्वर्गीयया नायिकया उर्वश्या सह संवादमयी सरसा शृ्ङ्गाररसवाहिनी कथा सङ्‌गृहीता विद्यते । सा कथा पौराणिकसाहित्यग्रन्थेषु लौकिकसाहित्यग्रन्थेषु च कविप्रतिभया सुसंस्कृत्य प्रस्तुता प्राप्यते । लौकिककाव्येषु वर्णिताः कथा रसानुकूल्येन स्थाने स्थाने परिवर्तिता दृश्यन्ते । अतः साहित्यिकदृष्ट्या अस्याः कथाया महत्त्वं विद्यते ।

\begin{visheshshloka}[center]

\textbf{उर्वशी हाप्सराः पुरूरवसमैडं चकमे । तं ह विन्दमानोवाच । त्रिः स्म माऽह्नो वैतसेन दण्डेन हतात् । अकामां स्म मा निपद्यासै । मो स्म त्वा नग्नं दर्शम् । एष वै नः स्त्रिणामुपचार इति ॥ १॥}

\end{visheshshloka}

\customparagraph{संस्कृतव्याख्या}

इडा नाम मनोर्दुहिता, तस्या पुत्रः= ऐडस्तमैडं पुरूरवसम्, अप्सरा देवजातिविशेषोत्पन्ना उर्वशी चकमे= कामितवती । तम्= पुरूरवसम्, विन्दमाना सा उर्वशी उवाच= उक्तवती । हे पुरूरवः ! अह्नः= दिवसस्य त्रिःस्म= त्रिवारम्, मा= माम्, दण्डेन= पुंस्प्रजननेन हतात्= उपभुक्ष्व । तथा अकामाम्= कामरहिताम्, सुरताभिलाषरहिताम्, मा स्म निपाद्यासै= निगृह्य न प्राप्नुयाः । तथा नग्नम्, वस्त्ररहितम्, त्वा= त्वाम्, मो स्म दर्शम्= नैव पश्येयम् । एतत् त्रितयं यावन्नियमेन करिष्यसि तावदहं तव जाया भविष्यामि । एष वै= एतादृश खलु, नोऽस्माकं स्त्रीणाम्, उपचारः= परिचर्या (सेवाविशेषः) इति ।

\customparagraph{\textbf{नेपाली}}

मनुपुत्री इडाका छोरा ऐड अर्थात् पुरूरवालाई देवयोनिविशेष अप्सरा उर्वशीले मन पराइन् । ती पुरूरवालाई चाहने उर्वशीले पुरूरवासमक्ष आफ्ना तीन सर्त राखिन्- १. मसँग दिनमा तीनपटकमात्र सहवास गर्न पाउँछौ । २. मलाई कामना नभएका वेला तिमीले सहवास गर्नेछैनौ। ३. तिमीलाई कुनै पनि वेला वस्त्र विना देख्नु नपरोस् । यी तीन सर्त जबसम्म तिमी पालन गर्छौ तबसम्म म तिम्री पत्नी बनेर पृथ्वीलोकमा बस्नेछु । अन्यथा तिमीले मलाई गुमाउनेछौ । यी तीन कुरा हामी स्त्रीजातिका सेवासर्त हुन् ।

\begin{visheshshloka}[center]

\textbf{सा हास्मिन् ज्योगुवास । अपि हास्माद्‌गर्भिण्यास तावत् । ज्योग्घास्मिन्नुवास । ततो ह गन्धर्वाः समूदिरे । ज्योग्वा इयमुर्वशी मनुष्येष्ववात्सीत् । उपजानीत । यथेयं पुनरागच्छेदिति । तस्यै हाविर्द्व्युरणा शयन उपवद्धाऽऽस । ततो ह गन्धर्वा अन्यतरमुरणं प्रमेयुः ॥ २॥}

\end{visheshshloka}

\customparagraph{संस्कृतव्याख्या}

इत्थं कृतनियमा सा= उर्वशी अस्मिन्= पुरूरवसि ज्योक्= चिरकालम्, उवास= निवासं कृतवती । अपि= खलु अस्मात्पुरूरवसः यावता कालेन गर्भिणी आस= गर्भयुक्ता बभूव । तावत्पर्यन्तं चिरकालमस्मिन्नुवासेत्यर्थः । एवम्प्रकारेण चिरकालं मनुष्यलोकसंस्थायिन्यामुर्वश्यां तद्विरहमसहमानैर्गन्धर्वैः कृतं तद्धरणोपायं दर्शययति- ततः= तस्मात् ऊर्वश्या मनुष्योपभोगाद्धेतोः, गन्धर्वाः=देवजातिविशेषाः, समूदिरे= संवादं मन्त्रालोचनं कृतवन्तः। इयमुर्वशी मनुष्येषु= मनुष्ययोनौ ज्योक्= चिरकालम्, अवात्सीत्= वासमकार्षीत् । यथा= येन प्रकारेण इयम्= उर्वशी अस्मान्पुनरागच्छेत्= पुनस्सहचारिणी भवेत्, तं प्रकारम् उपजानीत= पर्यालोचयत । एतस्मिन्नवसरे तस्याः खलु उर्वश्याः शयने= शयनप्रदेशे द्व्युरणा= द्वौ-उरणौ बालकौ मेषौ यस्या तादृशी अविर्मेषी उपबद्धाऽऽस= उपनिबद्धा बभूव । तौ द्वावुरणौ उर्वश्या पुत्रत्वेन पालितौ । ततः= संवादानन्तरम्, गन्धर्वाः तयोः= उरणयोः, अन्यतरम्= एकम् उरणम्, प्रमेथुः= अपजह्रुः ॥ २॥

\customparagraph{\textbf{नेपाली}}

यसरी सर्तसहित उर्वशी धेरै समयसम्म पुरूरवासँग मर्त्यलोकमा बसिन् । कालक्रमले उर्वशी गर्भवती भइन् । यति लामो समयसम्म मर्त्यलोकमा रहेकी उर्वशीलाई कसरी हाम्री सहगामिनी बनाउन सकिएला भन्ने विषयमा गन्धर्वहरूले विचारविमर्श गरे । तिनले पहिलेका तीन सर्त मध्ये कुनैबाट विचलित गराउने अठोट गरे ।  यही उद्देश्यका साथ यसै समयमा शयनकक्षमा उर्वशीसँग सुतिरहेका भेडाका दुई पाठा मध्ये एउटालाई गन्धर्वहरूले हरण गरे ।

\begin{visheshshloka}[center]

\textbf{सा होवाच अवीर इव बत मेऽजन इव पुत्रं हरन्ति । द्वितीयं प्रमेथुः  ।  सा ह तथैवोवाच ॥ ३॥}

\end{visheshshloka}

\customparagraph{संस्कृतव्याख्या}

तस्मिन्नुरणेऽपहृते सति सा= उर्वशी उवाच= उक्तवती । अवीरे= शौर्यरहिते इव= देशे अजने= जनरहिते इव=देशे मे=मम, मदीयं उरणरूपम् पुत्रम् = पुत्रवत् पालितमुरणम्, हरन्ति=नयन्ति, बत इति खेदे । तया एवमुक्तेऽपि पुरूरवास्तूष्णीं स्थितवान् । ततो गन्धर्वा द्वितीयम्= अपरम्, उरणं प्रमेथुः= अपजह्रुः । सा खलु उर्वशी तथैव= प्रागिव, उवाच= कथितवती।

\customparagraph{\textbf{नेपाली}}

उर्वशीले पुत्रवत् पालेका भेडाका पाठामध्ये एउटा हराएपछि उर्वशी भन्छिन्- “कोही साहसी वीर नभएको निर्जन ठाउँमा रहेझैं मेरो एउटा पाठो हरायो, यो अत्यन्तै दुःख लाग्दो कुरा भयो।” यति भन्दा पनि पुरूरवा मौन रहे । गन्धर्वले फेरि दोस्रो पाठो पनि चोरे । ती उर्वशीले पहिलो पाठो हराउँदाझैं पुरूरवालाई निन्दा गर्दै भनिन् ।

\begin{visheshshloka}[center]

\textbf{अथ हायमीक्षाञ्चक्रे । कथं नु तदवीरं कथमजनं स्यात्; यत्राहं स्यामिति । स नग्न एवानूत्पपात । चिरं तन्मेने; यद्वासः पर्यधास्यत । ततो ह गन्धर्वा विद्युतं जनयाञ्चक्रुः । तं यथा दिवैवं नग्नं ददर्श । ततो हैवेयं तिरोबभूव । पुनरैमीत्येत् । तिरोभूतां स आध्या जल्पन् कुरुक्षेत्रं समया चचार । अन्यतःप्लक्षेति । बिसवती । तस्यै हाध्यन्तेन वव्राज । तद्ध ता अप्सरस आतयो भूत्वा पुप्लुविरे ॥ ४॥}

\end{visheshshloka}

\customparagraph{संस्कृतव्याख्या}

एवमस्यां विलपन्त्यां सत्यां पुरूरवाः किं कृतवानिति तदाह- अथ= अनन्तरम्, अयम्= उर्वश्याः समीपे शयानः पुरूरवाः ईक्षाञ्चक्रे= पर्यालोचितवान्- “यत्र=यस्मिन् स्थाने अहम्= पुरूरवाः स्याम्= भवेयम्, तत्स्थानं कथं नु= केन खलु प्रकारेण अवीरम्= वीररहितम्,  तत् स्थानं कथं वा अजनम्= जनरहितम्, स्यात्=भवेत्।” एवं विचार्य सः=पुरूरवाः, नग्नः= विवसनः, एव= सन् अनूत्पपात= गन्धर्वैः सह योद्‌धुमुच्चक्राम । नग्नमेव इत्यस्मिन् कारणमाह- यद्= यदा वासः= वस्त्रम्, पर्यधास्यत= परिधानं क्रियते तत्= तावत् कालपर्यन्तम्, चिरं मेने= कालविलम्बः स्यादिति बुध्यते स्म । तस्मान्नग्न एवानूत्पपातेत्यर्थः । ततः= तदनन्तरम्, गन्धर्वाः, नग्नमागच्छन्तं पुरूरवसं दृष्ट्वा उर्वशीसमयभङ्गाय नग्नस्यैव तस्य दर्शनार्थम्, विद्युतम्= प्रकाशम्, जनयाञ्चक्रुः= आविष्कृतवन्तः । तत्प्रभया यथा दिवा= अहनि तथैव तम्=पुरूरवसम्, उर्वशी ददर्श= दृष्टवती, तदैव मम समयभङ्गो जात इत्यनन्तरमेव इयम्= उर्वशी तिरोबभूव= अन्तर्हितवती । पुनः= भूय ऐमि= आगच्छामि, इति ब्रुवाणा एत्= स्वर्गलोकमागच्छत् । पुनः पुरूरवसा कृतं दर्शयति- तिरोभूताम्= अन्तर्हिताम्, उर्वशीं सः= पुरूरवाः, आध्या= विरहजनिता मनःपीडा आधिस्तया जल्पन्= प्रलपन् कुरुक्षेत्रं समया= कुरुक्षेत्रस्य समीपे चचार= उर्वशीं द्रष्टुकामो भ्रमणं कृतवान्। तत्र= कुरुक्षेत्रे एकत्र देशे अन्यतःप्लक्षा इति नाम बिसवती= पद्मिनी सरसि तस्यै हाधि= तस्याः समीपप्रदेशे अन्तेन= अन्तिकेन तटप्रदेशेन वव्राज= गतवान् । अथोर्वश्या पुनः प्रादुर्भावमाह- तद्ध= तत्र खलु अन्यतःप्लक्षाख्ये सरसि ताः= उर्वशीप्रमुखा अप्सरसः, आतयः= जलचरपक्षिविशेषस्यैषा संज्ञा तद्रूपा भूत्वा पुप्लुविरे= परिप्लवनं जलक्रीडां कृतवत्यः ॥ ४॥

\customparagraph{\textbf{नेपाली}}

यसरी भेडाका पाठा अपहृत भएपछि पुरूरवा मनमनै विचार गर्छन्- “यो ठाउँ कसरी वीर पुरुष नभएको निर्जन हुन सक्छ, जहाँ मजस्तो पराक्रमी राजा विद्यमान छ ।” यस्तो सोचेर राजा कपडा लगाउन थाल्दा चोर धेरै टाढा पुग्छन् भन्ने ठानेर वस्त्रविहीन नै पाठा चोर्ने गन्धर्वसँग युद्ध गर्न निस्किए । त्यै वेला सर्त भङ्ग गराउने मनसाय राखी उर्वशीले पुरूरवालाई नग्न अवस्थामै देखून् भनेर गन्धर्वले बिजुली चम्काइदिए । बिजुलीको प्रकाशबाट पुरूरवाको नग्नता देख्ने बित्तिकै सर्त भङ्ग भएकाले पछि फेरि आउनेछु भन्दै उर्वशी अदृश्य भइन् । उर्वशीसँगको वियोगले विक्षिप्त बनी पुरूरवा कुरुक्षेत्रमा  घुमिरहे । कुनै दिन त्यही कुरुक्षेत्रको  धेरै कमलहरू फुलेको ‘अन्यतःप्लक्ष’ नाम गरेको तलाउमा ती उर्वशीसहित प्रमुख अप्सराहरू आएर जलक्रीडा गर्न लागे॥ ४॥

\begin{visheshshloka}[center]

\textbf{तं हेयं ज्ञात्वोवाच- अयं व स मनुष्यः, यस्मिन्नहमवात्समिति । ता  होचुः । तस्मै वा आविरसामेति । तथेति । तस्मै हाविरासुः ॥ ५॥}

\end{visheshshloka}

\customparagraph{संस्कृतव्याख्या}

एतस्मिन्नन्तरे तीरदेशं गच्छन्तं तम्= पुरूरवसम्, ज्ञात्वा= विज्ञाय इयम्= उर्वशी सखीभिः सह उवाच= उक्तवती । अयम्= दृश्यमानः पुरुषः,  स मनुष्योऽस्ति, यस्मिन्= पुरुषे अहम्= उर्वशी अवात्सम्= कतिचित् कालं यावद् जायात्वेन स्थिता । ताः= उर्वश्याः सख्यः, एतद्वाक्यं श्रुत्वा ऊचुः= उक्तवत्यः- तस्मै= मनुष्याय वयम् आविरसाम= प्रादुर्भवामेति । तथेति परस्परमङ्गीकृत्य तस्मै= पुरूरवसे आविरासुः= पक्षिरूपं विहाय स्वकीयेन रूपेण प्रादुर्बभूवुः॥५॥

\customparagraph{\textbf{नेपाली}}

यसैबेला कुरुक्षेत्रको तटप्रदेशमा घुम्दै पुगेका ती पुरूरवालाई चिनेर उर्वशीले साथीहरूसँग भनिन्- “हे हाथीहरू हो ! यी देखिएका पुरुषसँग मैले पत्नी भएर केही समय बिताएकी थिएँ ।” उर्वशीका कुरा सुनेर साथीहरू भन्छन्- त्यो परिचित मानिसका लागि हामी आफ्नो रूपमा प्रकट होऔं । यसरी परस्पर सल्लाह गरेर सबै अप्सराहरू पक्षीको रूप त्यागी वास्तविक रूपमा प्रकट भए।

\begin{visheshshloka}[center]

\textbf{तां हायं ज्ञात्वाऽभिपरोवाद । “हये जाये मनसा तिष्ठ घोरे वचांसि मिश्रा कृणवावहै नु । न नौ मन्त्रा अनुदिदास एते मयस्करन्परतरे च नाहन्”}- इति । \textbf{उप नु रम सं नुवदावहै; इति हैवैनां तदुवाच ॥ ६॥}

\end{visheshshloka}

\customparagraph{संस्कृतव्याख्या}

ततोऽयम्= पुरूरवा अप्सरोगणमध्ये तामुर्वशीं ज्ञात्वा= अभिलक्ष्य परावृतः सन्नुवाद- हये जाये= इति आवभाषे । तदेतत् उर्वशीपुरूरवसोः उक्तिप्रत्युक्तिरूपं पञ्चदशर्चं सूक्तम्-

हये जाये ! हे पत्नि ! मनसा= हृदयेन घोरे= क्रूरे, उर्वशि ! तिष्ठ= मा गाः । नु= अद्य वचांसि= परस्परवचनानि मिश्रा= मिश्रीकृतानि कृणवावहै= आवां करवावहै । नौ= आवयोर्मन्त्राः= रहस्योक्तिविशेषाः, अनुदितासः= अनुक्ताः सन्तः, मयः= सुखनामैतत्, सुखम्, न करन्= न कुर्वति । तथा परतरेऽहन् च= परस्मिन्नप्यहनि अनुक्ता एते= संवादविशेषाः सुखं न कुर्वति । तस्माद् मा गा इत्यर्थः । हे जाये ! नु= अद्य उपरम= गमनान्निवर्तस्व । नु= क्षिप्रम् आवाम्, संवदावहै= मिश्रभाषणं करवावहै, इति ह= एवमेव खलु एतया ऋचा  एनाम्= उर्वशीम्, पुरूरवा उवाच= कथितवान्॥६॥

\customparagraph{\textbf{नेपाली}}

त्यसपछि पुरूरवाले अप्सराहरूका बीच यी उर्वशी नै हुन् भन्ने जानेर फर्काउने उद्देश्यले यस्ता कुरा गर्दछन्- “हे कठोर हृदय भएकी पत्नी उर्वशी ! तिमी नजाऊ, यहीं बस । हामी एकआपसमा  आन्तरिक कुरा गरौं । हामीमा आपसी कुरा नहुँदा दुबैलाई दुःख भएको छ। केही कुरा गरिएन भने आउने दिनमा पनि दुःख नै हुन्छ। त्यसैले तिमी नजाऊ । जाने इच्छा छाड । मसँग सुखदुःखका कुरा गर ।” यसरी पुरूरवाले यो ऋचाका माध्यमले उर्वशीलाई अनुरोध गर्दै भने ।

\begin{visheshshloka}[center]

\textbf{तं हेतरा प्रत्युवाच- “किमेता वाचा कृणवा तवाहं प्राक्रमिषमुषसामग्रियेव पुरूरवः पुनरस्तं परेहि दुरापना वात इवाहमसि”- इति । न वै त्वं तदकरोः । यदहमब्रवम् । दुरापा वा अहं त्वया एतर्ह्यस्मि पुनर्गृहानिहि; इति हैवैनं तदुवाच ॥ ७॥}

\end{visheshshloka}

\customparagraph{संस्कृतव्याख्या}

इतरा= उर्वशी पुरूरवसं प्रत्युवाच= कथितवती, एता= त्वदीयया अर्थशून्यया वाचा= वचनेन किं कृणवा= अहं किं कुर्याम् । समयभङ्गात्निमित्ताद् अहं त्वत्सकाशात् प्राक्रमिषम्= प्रस्थानमकार्षम् । तत्र दृष्टान्तः- सूर्योदयेन विनिवर्त्यानाम् उषसां मध्ये अग्रिया= अग्रे भवा, पूर्वा उषा यथा प्रक्रमति, न पुनरावर्तते, एवं मया प्राक्रमिषमित्यर्थः । अतस्त्वं निराशो भूत्वा अस्तम्= गृहनामैतत्, स्वकीयं गृहं पुनः परेहि= पराङ्‌मुखो निवर्तस्व । अहं पुनर्दुरापना= आप्तुमशक्या । यस्मादहं देवस्त्री त्वं मनुष्यः । अतो वात इव= यथा वातो गन्तुं न शक्यते मत्प्राप्तिस्तव दुष्करेत्यर्थः। इमाम् ऋचं श्रुतिर्व्याकरोति- ‘मो स्म त्वा नग्नं दर्शम्’ इति यदहमब्रवम्, न खलु त्वं तदकरोः, अतस्त्वया एतर्हि= एतस्मिन् काले अहं दुरापा= आप्तुमशक्या अस्मि, अतस्त्वं गृहान्= आत्मीयान् पुनरिहि= गच्छ । इति ह= एवमेव खलु एनम्= पुरूरवसम्, तत्= तदानीम्, उर्वशी उवाच= कथितवती ।

\customparagraph{\textbf{नेपाली}}

पुरूरवाको आग्रह सुनेर उर्वशी भन्छिन्- यी निरर्थक कुरा सुन्नुको कुनै फाइदा छैन । सर्तमा हारेका कारण म तिमीबाट टाढा भएकी हुँ । जसरी सूर्योदय भएपछि उषा जान्छिन्, पुनः फर्किन्नन्, त्यस्तै म पनि गइसकेपछि फर्किनेवाला छैन । मैले पैले नग्न दर्शन नहोस् भनेकी थिएँ, तिमीले त्यो पालन गर्न सकेनौ । अब तिम्रा लागि म वायुझैं दुष्प्राप्य छु किनकि तिमी मनुष्य म देवस्त्री, हाम्रो मिलन सम्भव छैन ।  अतः तिमी फर्किएर आफ्ना घर जाऊ । यसरी उर्वशी पुरूरवालाई उत्तर दिन्छिन् ।

\begin{visheshshloka}[center]

\textbf{अथ हायं परिद्यून उवाच । “सुदेवो अद्य प्रपते दनावृत्परावतं परमां गंतवा उ । अधा शयीत निर्ऋतेरुपस्थेऽधैनं वृका रभसासो अद्युः”--इति । सुदेवोऽद्योद्वा बध्नीत । प्र वा पतेत् । तदेनं वृका वा श्वानो वाऽद्युः । इति हैव तदुवाच ॥ ९॥}

\end{visheshshloka}

\customparagraph{संस्कृतव्याख्या}

एवं प्रकारेण निराकृतस्य पुरूरवसो दुःखं दर्शयति- अथ ह= एवं प्रत्याख्यानानन्तरमेव अयम्= पुरूरवाः परिद्यूनः= प्राप्तपरिदेवनः सन् शोकातुरः उवाच= उक्तवान्- सुदेवः= सुष्ठु दीप्यमानः अशनिः (वज्रः) मम मस्तके प्रपतेत् । अथ अयं जनः अनावृत्= अनावर्तमानः सन् प्राप्तपरिदेवनः परमाम्= उत्कृष्टाम्, परावतम्= परलोकलक्षणं दूरदेशम्, गन्तवै= गन्तुं प्रभवति । अथानन्तरमशनिहतः सन् निर्ऋतेर्भूम्याः उपस्थे= उत्सङ्गे शयीत । अथानन्तरम् एनम्= मृतम्, रभसासः= वेगयुक्ता वृकाः= अरण्याश्वानः अद्युः= भक्षयेयुः । यद्वा- सुदेव= सुक्रीडोऽयं जनः अनावृतः सन् राज्यभोगं विहाय परावतम्= अत्यन्तदूरदेशं प्रपतेत्= प्रगच्छेत्। अस्या ऋचस्तात्पर्यं श्रुतिर्दर्शयति- सुदेवः= सुक्रीडोऽयं जनः अद्य= इदानीम्, हे उर्वशि ! त्वद्विरहनिमित्तात् उद्वध्नीत वा= देहत्यागार्थं पाशेन उद्बन्धनं वा कुर्यात् । प्रपतेद्वा= प्रपतनं महाप्रस्थानं कुर्याद्वा । तत्= ततः परम्, एनम्= पुरूरवसम्, सालावृका वा श्वानो वा अद्युः= भक्षयेयुः । इति ह= एवमेव खलु तत्= तेन परिदेवनेन मन्त्रेणोक्तवान् ॥ ८॥

\customparagraph{\textbf{नेपाली}}

यसरी उर्वशीले निराश बनाएपछि पुरूरवा अत्यन्तै दुःखी भई भन्छन्- “हे उर्वशी ! मेरो शिरमा वज्रपात होस् । यो दुःखी मान्छे पीडाले छटपटाएर दूरदेश जान पनि सक्छ । अथवा मूर्छा परेर पृथ्वीमै ढलेको मृतप्राय मलाई जङ्गली कुक्कुरले खानेछन् अथवा म राज्य छाडेर अनन्त यात्रामा जानेछु । तिम्रा अभावमा अब मलाई जीवनको कुनै मोह छैन ।”  यसरी हे उर्वशी ! तिम्रा अभावमा म पुरूरवा पासो लागेर वा कुनै पहाडबाट फाल हानेर प्राण त्याग गरेर मर्छु अनि मलाई जङ्गली कुक्कुर आदिले खानेछन् भन्ने भाव यस पीडा प्रधान मन्त्रबाट अभिव्यक्त छ ।

\begin{visheshshloka}[center]

\textbf{तं हेतरा प्रत्युवाच । “पुरूरवो मा मृथा मा प्रपप्तो मा त्वा वृकासो अशिवास उ क्षन् । न वै स्त्रैणानि सख्यानि सन्ति सालावृकाणां हृदयान्येता”--इति । मैतदादृथाः । न वै स्त्रैणं सख्यमस्ति । पुनर्गृहानिहि । इति हैवैनं तदुवाच ॥ ९॥}

\end{visheshshloka}

\customparagraph{संस्कृतव्याख्या}

एवं परिदेवने कृते तन्निवारणार्थम् इतरा= उर्वशी पुरूरवसं प्रति उवाच= उक्तवती- हे पुरूरवः ! वृथा मा मृथा= न म्रियस्व । मा प्रपप्ते= प्रपतनं मा कृथा । तथा अशिवासः= दुष्टाः, वृकासः= वृकाः, अरण्यश्वानश्च त्वा= त्वाम्, मा क्षन्= मा भक्षयन्तु । स्त्रैणानि= स्त्रिया सम्बन्धीनि सख्यानि= सखित्वानि न खलु सन्ति= चिरकालरूढान्यपि सख्यानि स्त्रीषु नावतिष्ठन्त इत्यर्थः । एता= एतदीयानि स्त्रीसम्बन्धीनि मनांसि सालावृकाणां हृदयानि= तद्वत् क्रूराणीत्यर्थः । एतत्= मत्प्राप्त्युपायचिन्तनं मा कृथा= तद्विषयमादरं मा कार्षीः । न खलु स्त्रैणम्= स्त्रीसम्बन्धि सख्यमस्ति= स्थिरतरं भवति । अतः पुनरात्मीयान्=गृहान् गच्छ । एवं प्रकारेण एनम्= पुरूरवसम्, उर्वशी उवाच= कथितवती ।

\customparagraph{\textbf{नेपाली}}

यसरी वियोगको पीडाले विक्षिप्त भएका पुरूरवालाई उर्वशी भन्छिन्- “हे पुरूरवा ! बेकारमा मर्ने कुरा नगर । कतैबाट हामफाल्ने कुरा पनि नगर । तिमीलाई जङ्गली कुक्कुरहरूले नखाऊन् । स्त्रीसँगको सम्बन्ध वा मित्रता लामो समयसम्म रहँदैन । स्त्रीहरूको हृदय ब्बाँसाको जस्तो कठोर हुन्छ । यो मलाई पाउने जुन आशा छ, त्यसलाई छाडिदेऊ । स्त्रीहरूसँगको सम्बन्ध

दीर्घकालीन हुँदैन भन्ने कुरा बुझ । त्यसैले अब आफ्नै घर फर्क ।” यसरी विक्षिप्त भएका पुरूरवालाई उर्वशीले भनिन् ।

\begin{visheshshloka}[center]

\textbf{“यद्विरूपाऽचरं मर्त्येष्ववसं रात्रीः शरदश्चतस्रः । घृतस्य स्तोकं सकृदह्नः आश्नां}

\textbf{ता देवेदं तातृपाणा चरामि”  इति । तदेतदुक्तप्रत्युक्तं पञ्चदशर्चं बह्‌वृचाः प्राहुः ।}

\textbf{तस्यै ह हृदयमाव्ययाञ्चकार ॥ १०॥}

\end{visheshshloka}

\customparagraph{संस्कृतव्याख्या}

इत्थं सहसा प्राप्तुकामं पुरूरवसं निराकृत्य उर्वशी स्वप्राप्त्युपायान्तरं तस्योपदेष्टुं तत्कृतमुपकारं आविष्करोति-

यद्विरूपा=विगतदेवरूपा सति मर्त्येषु= मनुष्येषु मध्ये अचरम्= अवर्तिषि । तदानीं रात्रीः= रमयित्रीः चतस्रः शरदः= संवत्सरचतुष्टयपर्यन्तम्, त्वत्समीपे अवसम्= कृतनिवासाऽस्मि । तस्मिन् वा समये अह्नः= एकस्य दिवसस्य सकृत्= एकवारम्, घृतस्य= सर्पिषः स्तोकम्= बिन्दुः आश्नाम्= अभुञ्जि । तव स्वभूतं घृतम् अपिबम् इत्यर्थः । इदम्= इदानीम्, देवे= देवत्वेऽपि तान् तस्मादेव घृतस्तोकात् तातृपाणा= तृप्ता सती चरामि= निवसामि । अतस्त्वां न विस्मरामीत्यर्थः । तदेतत् पञ्चदशर्चं सूक्तं ‘हये जाये मनसा’ इत्यादिकम् उक्तप्रत्युक्तम्= पुरूरवस उक्तिः उर्वश्याः प्रत्युक्तिः ताभ्यामुपेतं बह्‌वृचाः ऋक्शाखाध्यायिनः प्राहुः=प्रकर्षेणाधीयते । ‘हये जाये’ इत्युदाहृतव्यतिरिक्तास्तत्रत्या अन्या ऋचोऽपि अत्रानुसन्धेया इत्यर्थः । इत्थम्= पुरूरवसा स्नेहप्रवाहेन उक्ते सति तस्याः उर्वस्या हृदयम्= मनः, अव्यायाञ्चकार= विगतकाठिन्यं बभूव ॥ १०॥

\customparagraph{\textbf{नेपाली}}

यसरी सजिलै पाउन सकिन्छ भन्ने सोचेका पुरूरवालाई निराकरण गरेपछि प्रकारान्तरले कसरी पाउन सकिन्छ ? भन्ने कुरा सङ्केत गर्न उर्वशी तिनले गरेको उपकारलाई सम्झिन्छिन्-

म (शाप विशेषले) देवस्वरूप छाडेर मनुष्य रूपमा मर्त्यलोकमा आएको अवस्थामा तिम्रा समीपमा मैले रमाइला रात्रिहरूले सहित चार वर्ष बिताएँ । ती वर्षमा दिनको एकपटक तिम्रा घरको घिउ खाएकी थिएँ । अहिले त्यसैले तृप्त भएर घुमिरहेकी छु । त्यसैले तिमीलाई भुल्न त सक्दिनँ । यहाँ यस्ता उर्वशी र पुरूरवाका बीच भएको संवाद वर्णन गरिएका  ‘हये जाये’ जस्ता धेरै ऋचा ऋक्शाखाध्यायीले पढेका छन् । यीभन्दा अरू ऋचा पनि अनुसन्धेय छन् । यसरी पुरूरवाले भनेका स्नेहशील कुरा र उर्वशीले सम्झिएका प्रेमपूर्ण कुराले उर्वशीको मन पग्लियो ।

\begin{visheshshloka}[center]

\textbf{सा होवाच । संवत्सरतमीं रात्रिमागच्छतात् । तन्म एकां रात्रिमन्ते शयितासे । जात उ तेऽयं तर्हि पुत्रो भवितेति । स ह संवत्सरतमीं  रात्रिमाजगामेद्धिरण्यविमितानि । ततो हैनमेकमूचुः एतत्प्रपद्यस्वेति तद्धास्मै तामुपप्रजिग्घ्युः ॥ ११॥}

\end{visheshshloka}

\customparagraph{संस्कृतव्याख्या}

इत्यमुक्तिप्रत्युक्तिभ्यां विलीनहृदया उर्वशी उवाच= पुरूरवसं प्रत्युक्तवती- संवत्सरतमीम्= संवत्सरस्य पूरणीमन्तिमां रात्रिम् आगच्छतात्, तत्= तत्र मे= मम अन्ते= निकटे ताम्= एकाम्, रात्रिं शयितासे= शयनं करिष्यसि । तर्हि= तथा सति ते= त्वदीयांशोऽयं ममोदरे गर्भरूपेण वर्तमानः पुत्रो जातो भविता= भविष्यति । एवमुर्वश्योक्तः पुरूरवाः संवत्सरतमीं रात्रिम् आजगामेत्= संवत्सरकालं स्थित्वा तस्यान्तिमां रात्रिमवगम्य हिरण्यविमितानि= सुवर्णविनिर्मितानि सौधानि आजगामैव । ततस्तस्मादागमनादनन्तरम् एनम्= पुरूरवसम्, तत्रत्या जनाः इदमेकमूचुः= एतत्स्थानं प्रतिपद्यस्व= आरोहेति । स प्रपद्य स्थितः । तत्तत्र तस्मै पुरूरवसे तामुर्वशीम् उपप्रजिग्घ्युः= प्रेषितवन्तः ॥ ११॥

\customparagraph{\textbf{नेपाली}}

यसरी पुरूरवासँगको संवादले तरलित हृदय भएकी उर्वशीले पुरूरवालाई भनिन्- “प्रत्येक वर्षको अन्तिम रात्रि म आउनेछु । तिमी मसँग शयन गर्नेछौ । यो मेरो पेटमा रहेको पुत्र पनि जन्मिनेछ।” यसरी उर्वशीद्वारा भनिएका पुरूरवाले खुसी भएर सुवर्णमय भवन बनाए । त्यो अन्तिम रात्रिको समय आयो । त्यहाँ रहेका मानिसले पुरूरवालाई त्यस रमणीय महलमा आरोहण गर्न आग्रह गरे । उनी त्यही महलमा बसे । उर्वशीलाई पनि सबैले पुरूरवासँग पठाए ।

\begin{visheshshloka}[center]

\textbf{सा होवाच- गन्धर्वा वै ते प्रातर्वरं दातारः । तं वृणासा इति । तं वै त्वमेव वृणीष्वेति । युष्माकमेवैकोऽसानीति ब्रूतादिति । तस्मै ह प्रातर्गन्धर्वा वरं ददुः । स होवाच- युष्माकमेवैकोऽसानीति ॥ १२॥}

\end{visheshshloka}

\customparagraph{संस्कृतव्याख्या}

सा= उर्वशी पुनः पुरूरवसं प्राप्य एवमव्रवीत्-  हे पुरूरवः ! ते प्रातः= प्रातःकाले गन्धर्वा वरं दातारः= दास्यन्ति । तम्=वरम्, त्वदभिलषितमर्थं वृणासै= प्रार्थयस्वेति । अथ पुरूरवा आह- हे उर्वशि ! मे= मदीयम्, वरम् अहं न जानामि, त्वमेव मदर्थं वृणीष्वेति, येनोपायेन त्वत्प्राप्तिर्मम भवति तमुपायं कथय । सा प्रतिवक्ति- “हे गन्धर्वाः ! युष्माकं मध्ये अहमप्येकः असानि= भवानीति । एवम् मया प्रार्थनीयो वर इति गन्धर्वान् प्रति ब्रूतात्= ब्रूहीति ।” एवं प्रार्थनीयं वरमुपदिश्य सा निर्जगाम । अथ तस्मै= पुरूरवसे प्रातः= प्रभाते गन्धर्वा आगत्य वरं ददुः । स च उर्वश्योपदिष्टमेव वरमयाचतेत्याह-

\customparagraph{\textbf{नेपाली}}

ती उर्वशीले फेरि पुरूरवालाई पाएर यसो भनिन्- “हे पुरूरवा ! तिमीलाई प्रातःकालमा गन्धर्वहरू वर दिन आउनेछन् । त्यसबेला तिमीले आफूले चाहेको वर माग्नू ।” उर्वशीका कुरा सुनेर पुरूरवा भन्छन्- हे उर्वशी ! कस्तो वर माग्ने ? म जान्दिनँ, तिमी नै त्यो उपाय भन कि तिमीलाई मैले पाउन सकूँ । ती उर्वशीले उपाय भनिन्- “हे गन्धर्व हो ! तपाईंहरूमध्ये म पनि एक होऊँ । यस्तो वर चाहन्छु” भनी गन्धर्वहरूलाई भन्नू । यसरी वरका बारेमा भनेर उर्वशी गइन् । बिहान गन्धर्वहरूसँग उर्वशीले भनेझैँ वर मागे ।

\begin{visheshshloka}[center]

\textbf{ते होचुः - न वै सा मनुष्येषु अग्नेर्यज्ञिया तनूरस्ति, ययेष्ट्‌वाऽस्माकमेकः स्यादिति: तस्मै ह स्थाल्यामोप्याग्निं प्रददुः; अनेनेष्ट्‌वाऽस्माकमेको भविष्यसीति । तं च ह कुमारं चादायावव्राज। सोऽरण्य एवाग्निं निधाय कुमारेणैव ग्राममेयाय । पुनरैमीत्येत् । तिरोभूतं योऽग्निः अश्वत्थं तं, या स्थाली शमीं ताम् । स ह पुनर्गन्धर्वानेयाय ॥ १३॥}

\end{visheshshloka}

\customparagraph{संस्कृतव्याख्या}

एवं गन्धर्वरूपत्वे पुरूरवसा प्रार्थिते सति तत्प्राप्त्युपायफलं गन्धर्वैरुपदेश्यमनुक्रामति-  ते= गन्धर्वाः, उचुः= कथितवन्तः, न खलु अग्ने सम्बन्धिनी यज्ञिया= यज्ञार्हा सा तादृशी तनूः, मनुष्येष्वस्ति । यया= तन्वा इष्ट्‌वा= यागं कृत्वा अस्माकं मध्ये अयं पुरूरवा एकः स्यादिति विचार्य तस्मै= पुरूरवसे यज्ञार्हमग्निं स्थाल्यामोप्य= स्थापयित्वा प्रददुः= दत्तवन्तः । अनेन= अग्निना इष्ट्‌वा अस्माकमेकं त्वं भविष्यसीति । एवं गन्धर्वैर्दत्ते सति तं चाग्निम्, उर्वशीगर्भाज्जातं कुमारं चादाय आवव्राज= स्वस्थानमाजगाम । सः पुरूरवाः अरण्य एव स्थालीसहितमग्निं निधाय कुमारेणैव सह ग्राममेयाय= आजगाम । ततः सः पुरूरवाः पुनरैमि= आगच्छामि, इति अरण्ये अभिनिहितं तिरोभूतम्= अन्तर्निहितमग्निम्, आ ऐत्= आगच्छत् । तत्र योऽग्निर्निहितः, तमश्वत्थमपश्यत् । तत्र या तदाधारभूता स्थाली तां शमीमपश्यत् । अग्निरश्वत्थवृक्षात्मना परिणतः, स्थाली तु शमीवृक्षरूपेण परिणता । इत्यग्नौ तिरोहिते सति सः पुरूरवाः पुनर्गन्धर्वानेयाय= आजगाम॥१३॥

\customparagraph{\textbf{नेपाली}}

यसरी गन्धर्वको रूप पाऊँ भनेर पुरूरवाले प्रार्थना गरेपछि त्यो स्वरूप पाउने उपाय बताउँदै गन्धर्वहरू भन्छन्- अग्निजस्तो यज्ञका लागि योग्य अन्य शरीर मानव जातिमा छैन । यही अग्निको शरीरले यज्ञ गरेर पुरूरवा पनि हामी जस्तै एक गन्धर्व बन्न सक्छन् भन्ने सोचेर ती पुरूरवालाई यज्ञार्ह अग्नि थालीमा राखेर दिए । यो अग्निले यज्ञ गरेर तिमी हामीमध्ये एक बन्न सक्छौ । यसरी गन्धर्वले दिएको त्यो अग्नि र उर्वशीबाट जन्मिएको पुत्रलाई साथमा लिएर ती पुरूरवा आफ्नो स्थानमा गए । ती पुरूरवा जङ्गलमा नै थालीसहितको अग्निलाई  छाडेर पुत्र लिई गाउँ फर्किए । त्यसपछि पुरूरवा फेरि आउँछु भनी जङ्गल जहाँ थालीसहित अग्नि राखेको थियो त्यहाँ गए। जहाँ तिनले अग्नि राखेको थियो त्यहाँ अश्वत्थ वृक्ष उम्रिएको देखे र जहाँ थाली राखेका थिए  त्यहाँ शमीको वृक्ष उम्रिएको देखे । यसरी अग्नि तिरोहित भएपछि पुरूरवा फेरि गन्धर्वकहाँ गए ।

\begin{visheshshloka}[center]

\textbf{ते होचुः - संवत्सरं चातुष्प्राश्यमोदनं पच । स एतस्यैवाश्वत्थस्य तिस्रस्तिस्रः समिधो धृतेनान्वज्य समिद्वतीभिर्घृतवतीभिर्ऋग्भिरभ्याधत्तात् । स यस्ततोऽग्निर्जनिता, स एव स भवितेति ॥ १४॥}

\end{visheshshloka}

\customparagraph{संस्कृतव्याख्या}

आगमनप्रयोजनं जानानास्ते गन्धर्वाः ऊचुः- चातुष्प्राश्यम्= चतुर्भिर्ऋग्भिः प्राशनीयं ब्रह्मौदनं संवत्सरं पच= पचस्व । तत्रैतस्याश्वत्थस्य तिस्रतिस्रः समिधः गृहीत्वा घृतेनन्वज्य मूलादारभ्यान्तपर्यन्तम् आज्येनाक्त्वा समिद्वतिभिः=समित्पदयुक्ताभिः घृतवतीभिः= घृतशब्दयुताभिः ‘समिधाऽग्निं दुवस्यत–’ इत्यादिभिर्ऋग्भिः अभ्याधत्स्व । अस्मिन् व्रते कृते सति ततः= समिदाधानान्निमित्तात् योऽग्निर्जनिता= जनिष्यते । सोऽयं स एव भविता= भविष्यति । योऽस्माभिः पूर्वं दत्त इति ॥

\customparagraph{\textbf{नेपाली}}

पुरूरवा आउनुको प्रयोजन बुझेर गन्धर्व भन्छन्- चारवटा ऋग्वेदका मन्त्रले उपभोग गर्न योग्य ब्रह्म भोजन पकाऊ । त्यहाँ यो पिपलका तीनतीनवटा समिधा लिएर घिउले सिञ्चित गरी ‘समित्’ र ‘घृत’ शब्दयुक्त ‘समिधाग्निं दुवस्यत–’ इत्यादि ऋग्वेदका मन्त्रले हवन गर । यो व्रत सम्पन्न गरिसकेपछि त्यो समिधाबाट जुन अग्नि निस्किन्छ, त्यो नै हामीले पहिले दिएको अग्नि हुनेछ ।

\begin{visheshshloka}[center]

\textbf{ते होचुः - परोऽक्षमिव वा एतत् । आश्वत्थीमेवोत्तरारणिं कुरुष्व;}

\textbf{शमीमयीमधारणिम् । स यस्ततोऽग्निर्जनिता, स एव स भवितेति ॥ १५॥}

\end{visheshshloka}

\customparagraph{संस्कृतव्याख्या}

ते=गन्धर्वाः पुनर्विचार्य ऊचुः- यदेतत्= अश्वत्थसमिदाधानेनाग्नेः पुनः सम्पादनम् एतत्परोऽक्षमिव व्यवहितमिव खलु भवति । तस्मादन्य उपाय कथ्यते- हे पुरूरवः ! त्वं अश्वत्थविकृतिजामेवोत्तरारणिं कुरुष्व, अधरारणिं शमीमयीं तथा च अग्न्यात्मकस्य उपर्यवस्थानम्, स्थाल्यात्मिकायाः शम्या अधस्तात् स्थितिश्चेत्युभयमप्युपपद्यते । ततः= तस्मादरणिद्वयात् मन्थनेन यः अग्निर्जनिता सः प्रागस्माभिर्दत्त अग्नितुल्यो भवति ।

\customparagraph{\textbf{नेपाली}}

ती गन्धर्वहरूले फेरि विचार गरी भने- जुन यो अश्वत्थको समिधाले अग्निको पुनः सम्पादन गरिन्छ, यो परोक्ष नै हुन्छ । यसमा अर्को उपाय भनिन्छ- हे पुरूरवा ! तिमी अश्वत्थबाट बनेको उत्तर अरणि बनाऊ, अधर अर्थात् तलको चाहिँ शमीमय बनाऊ । यसरी प्रयोग गर्दा अग्न्यात्मक जुन पिपल हो त्यो माथि र जुन पात्रात्मक शमी हो त्यो तल पर्छ । यसरी दुबैको सही उपयोग हुन्छ । ती दुबै अरणिको मन्थनबाट जन्मिएको अग्नि पहिले हामीले दिएको अग्नि नै हुनेछ ।

\begin{visheshshloka}[center]

\textbf{ते होचुः - परोऽक्षमिव वा एतत् । आश्वत्थीमेवोत्तरारणिं कुर्वीत;}

\textbf{आश्वत्थीमधरारणिम् । स यस्ततोऽग्निर्जायते, स एव स भवितेति॥१६॥}

\end{visheshshloka}

\customparagraph{संस्कृतव्याख्या}

ते= गन्धर्वाः पुनरप्यूचुः- अधरारण्युत्तरारण्योः वृक्षद्वयजत्वम् । एतत्परोऽक्षमिव= अन्तर्हितमिव भवति । तस्मादरणिद्वयमपि अश्वत्थवृक्षजं कुर्यात् । स यस्ततोऽग्निरित्यादि पूर्ववत् ॥

\customparagraph{\textbf{नेपाली}}

ती गन्धर्वहरूले फेरि भने- जुन तल र माथिको अरणि (काठ) पिपल र शमीबाट बनेका हुन्छन्, ती परोक्षझैं हुन्छन्। त्यसैले दुबै अरणि पिपलकै वृक्षको बनाउनू । त्यसबाट जुन अग्नि उत्पन्न हुन्छ, त्यो पहिले हामीले दिएकै अग्नितुल्य हुन्छ ।

\begin{visheshshloka}[center]

\textbf{स आश्वत्थीमेवोत्तरारणिं चक्रे; आश्वत्थिमधरारणिम् । स यस्ततोऽग्निर्जज्ञे । स एव स आस। तेनेष्ट्‌वा  गन्धर्वाणामेक आस । तस्मादाश्वत्थीमेवोत्तरारणिं कुर्वीत; आश्वत्थीमधरारणिम् । स यस्ततोऽग्निर्जायते, स एव स भवति । तेनेष्ट्‌वा गन्धर्वाणामेको भवति ॥ १७॥}

\end{visheshshloka}

\customparagraph{संस्कृतव्याख्या}

इत्थं गन्धर्वैरुपदिष्टे सति पुरूरवाः आश्वत्थीमेवोत्तरारणिं चक्रे= कृतवान् । आश्वत्थीमेवाधरारणिं च चक्रे= कृतवान् । ततस्ताभ्यां मन्थनात् सः= तादृशोऽग्निर्जज्ञे, यः प्रागरण्ये निहितः । गन्धर्वैर्दत्तोऽग्निः स एव आस= बभूव। अग्निरूपात्केवलादश्वत्थात् अव्यवधानेन जायमानत्वात् । अथ सः पुरूरवाः तेन= मथितेनाग्निना इष्ट्‌वा गन्धर्वाणामेक आस= बभूव । इत्थमाख्यायिकासिद्धम् अरण्योरश्वत्थोपादानत्वम् अधुना यजमानस्यापि विधत्ते । यस्मादेवं पुरूरवाः अश्वत्थयोन्यग्निना इष्ट्‌वा स्वाभीष्टमालभत । तस्मादिदानींतनो यजमानोऽपि आश्वत्थीमेव अधरारणिमुत्तरारणिं च कुर्वीत । ततः सकाशाद्योऽग्निर्जायते स एव यज्ञार्हो  देेवेष्ववस्थितोऽग्निर्भवति । गतमन्यत् ॥ १७॥

\customparagraph{\textbf{नेपाली}}

यसरी गन्धर्वबाट उपदेश पाएर पुरूरवाले अश्वत्थबाट नै बनेको माथि पट्‌टिको र तलपट्‌टिको अरणि बनाएर मन्थन गरे । त्यसबाट व्यवधानरहित अग्नि निस्कियो । त्यो अग्नि पहिले गन्धर्वले दिएकै थियो । त्यसैबाट यज्ञ गरी पुरूरवा गन्धर्वहरूमध्ये एक भएर आफ्नी प्रेमिकासँग इच्छाअनुसार निरन्तर बस्न पाए ।

यसरी यो आख्यायिकाबाट अरणि मन्थनबाट निस्किएको अग्निले यज्ञ गर्दा पुरूरवाले झैं इच्छाएको फल पाइन्छ भन्ने प्रमाणित भएकाले अहिलेका यजमानहरूले पनि अश्वत्थको अरणि बनाएर यज्ञ गर्ने गर्नुपर्दछ ।

\begin{visheshshlokatwo}[center]{84}

\textbf{शतपथब्राह्मणअन्तर्गत एकादश काण्डको पाँचौं अध्यायको}

\textbf{प्रथम ब्राह्मणमा सङ्‌गृहीत पुरूरवोर्वश्याख्यानको नेपाली}

\textbf{भावानुवाद पूर्ण भयो ।}

\end{visheshshlokatwo}

\specialkhanda{\textbf{याज्ञवल्क्यमैत्रेयी-संवादाख्यानम्}}

\customparagraph{\textbf{परिचयः}}

वैदिकसाहित्यस्य विवृतिरूपमुपनिषत्साहित्यं विद्यते । उपनिषदमधीत्य वैदिकसाहित्यस्य ज्ञानं लभ्यते । नैकासु उपनिषत्सु बृहदारण्यकोपनिषदः स्वरूपं दीर्घकायात्मकं विद्यते । इयम् उपनिषद् माध्यन्दिनशतपथब्राह्मणस्य चतुर्दशकाण्डस्य चतुर्थाध्यायाद् नवमाध्यायपर्यन्तं निबद्धा अस्ति । तत्र चतुर्थब्राह्मणस्य द्वितीयेऽध्याये याज्ञवल्क्यमैत्रेयीसंवादो विद्यते । तत्र अद्वैतवेदान्तविषयः सुबोध्यशैल्या प्रस्तुतः प्राप्यते ।

याज्ञवल्क्यस्य कात्यायनी मैत्रेयी चेति द्वे पत्न्यौ आस्ताम् । एकदा मुनिना सन्न्यासाश्रमं प्रविश्य अवशिष्टं जीवनं यापयामीति मनसि निश्चित्य द्वयोः पत्न्योरग्रे धनादिकं विभज्य दातुं प्रस्तावः प्रस्तुतः । प्रस्तावोऽयं कात्यायन्या तु स्वीकृतः परं मैत्रेय्या पतिं प्रति जिज्ञासा प्रस्तुता । ‘हे स्वामिन् ! किमहं धनेनामरत्वं मोक्षं वा प्राप्तुं शक्नोमि ?’ महर्षियाज्ञवल्क्यः प्रत्युत्तरं ददाति यत्- ‘न हि, धनेन केवलं विषयसुखमेव लभ्यते । अमरत्वस्य कृते ब्रह्मज्ञानमावश्यकं भवति । तदर्थं सर्वाणि भौतिकानि साधनानि परित्यक्तुमावश्यकं भवति ।’ तदा मैत्रेयी कथयति यत्- ‘यदि भवान् दिव्यानन्दप्राप्तये सन्न्यासमाश्रयति चेदहं कथं भौतिकपदार्थान् इच्छामि ? मम धनादिके रुचिर्नास्ति । भवान् मोक्षार्थं धनधान्यादिसहितं गृहस्थाश्रमं त्यजति चेदहं तदेव कथमिच्छामि । अहमपि भवतो मार्गमनुसर्तुमिच्छामि। तस्माद् मां ब्रह्मज्ञानं वितर, अहमपि भवता सहैव गच्छामि ।’ तदनन्तरं स्वपत्न्या मैत्रेय्या विचारेण हृष्टो मुनिर्याज्ञवल्क्यस्तस्याः प्रश्नस्योत्तररूपेण अद्वैतवेदान्तस्य बीजम् `अहं ब्रह्मास्मि' इत्यस्य रहस्यं पत्नीं मैत्र्येयीं श्रावयति ।

\customparagraph{\textbf{सारः/वैशिष्ट्यम्}}

साधारणजनाः शरीरमेव आत्मा इति मन्यन्ते परं तेऽज्ञानेन सत्यतत्त्वाद् दूरं विद्यन्ते  । शरीराद् भिन्न एव आत्मा विद्यते । जगतः सर्वे पदार्था आत्माश्रिता विद्यन्ते, परं मूढाः शरीरमेव आत्मा इति मत्वा शरीरस्य विलासार्थं सर्वान् पदार्थान् योजयन्ति । भणितिरपि विद्यते यत्- `आत्मार्थं पृथिवीं त्यजेत्' अर्थात् आत्मनः मोक्षार्थं सर्वेऽपि पदार्थाः त्याज्या विद्यन्ते । तस्माद् आत्मनो ध्यानम्, मननम्, चिन्तनञ्चावश्यकं विद्यते । एतेनैव जनो वास्तविकरूपेण सन्तुष्टो भवति। आत्मा एव ब्रह्म विद्यते । आत्मा सर्वत्र विद्यते । आत्मानं परित्यज्य अन्यत्र ब्रह्मतत्त्वं न लभ्यते । ब्रह्मस्वरूपात्मा एव सर्वस्वमस्ति । लोके जनाः स्वकीये पुत्रादिके यत् स्नेहं कुर्वन्ति तत्रापि आत्मीयमेव प्रयोजनं ज्ञायते  । अतः आत्मनि एव मुख्यः स्नेहो विद्यते । तस्मादात्मा एव द्रष्टव्यो मन्तव्यो निदिध्यासितव्यश्च । आत्मतत्त्वत एव सर्वे जीवाः समुत्पन्नाः, पश्चात् तेऽत्रैवान्तर्लीयन्ते । यथा- लवणखण्डः समुद्रजलादेवोत्पद्यते । पुनः स लवणखण्डो जले स्थाप्यते चेत् तत्रैव लीनो जायते; तथैव ब्रह्मतत्त्वतः सञ्जाता नामोपाधिधारकाः पदार्थाः अन्ते तत्रैव लीयन्ते । ‘आत्मा’ ज्ञानात् परो विद्यते । यदि कोऽपि अहमात्मानं जानामि इति कथयति चेत् स किमपि न जानाति । स एव आत्मानम् अनन्तम् अवेद्यमस्तीति जानाति चेत् तेन आत्मतत्त्वं ज्ञातम् । अस्याध्ययनं नेति नेति इति शृङ्खलया भवति । तस्माद् आत्मा एव अद्वैतं ब्रह्मास्ति । अनन्तं अपरिमेयम् आत्मस्वरूपं ब्रह्मतत्त्वं विद्यते इति याज्ञवल्क्यमैत्रेयीसंवादस्य सारोऽस्ति ।

\begin{visheshshloka}[center]

\textbf{मैत्रेयीति होवाच याज्ञवल्क्य उद्यास्यन्वा अरेऽहमस्मात्स्थानादस्मि}

\textbf{हस्त तेऽनया कात्यायन्यान्तं करवाणीति॥१॥}

\end{visheshshloka}

\customparagraph{संस्कृतव्याख्या}

‘अरे’ इति सम्बोधनम् । मैत्रेयि != मम पत्नि ! अहम्= याज्ञवल्क्यः, अस्मात्स्थानात्= अस्माद् आश्रमाद् उद्यास्यम्= अपराश्रमम्, वानप्रस्थाश्रमम्, गन्तुं समुत्सुकोऽस्मि । अतो हन्त= अनुमतिं प्रार्थयामि। एवं ते= तव अनया कात्यायन्या= द्वितीयया भार्यया सह अन्तम्= धनस्य विभागम्, करवाणि इति ।

\customparagraph{\textbf{नेपाली भावार्थ}}

याज्ञवल्क्य पत्नी मैत्रेयीलाई भन्छन्- हे मैत्रेयि ! म यो भन्दा माथिल्लो आश्रम वानप्रस्थमा जान चाहन्छु । त्यसैले तिम्रो अनुमति चाहन्छु । यो जति सम्पत्ति छ, त्यो सबै तिमी र अर्की पत्नी कात्यायनीलाई भाग लगाइदिन्छु ।

\begin{visheshshloka}[center]

\textbf{सा होवाच मैत्रेयी । यन्नु मे इयं भगोः सर्वा पृथिवी वित्तेन पूर्णा स्यात्कथं तेनामृता स्यामिति नेति होवाच याज्ञवल्क्यो यथैवोपकरणवतां जीवितं तथैव ते जीवितं स्यादमृतत्त्वस्य तु नाशास्ति वित्तेनेति ॥२॥}

\end{visheshshloka}

\customparagraph{संस्कृतव्याख्या}

सा मैत्रेयी होवाच‍= एवमुक्तवती- भगोः= हे स्वामिन् ! यत्= इयं दृश्यमाना सर्वा= सम्पूर्णा पृथिवी मे= मम स्यात्= यदि भवेत् चेत् (नु इति वितर्के) कथम्= किम् तेन कारणेन अहम् अमृता= मृत्युरहिता स्याम् ? इति मैत्रेय्या जिज्ञासा । याज्ञवल्क्यो होवाच-  नैव नहि, अमृता नैव स्याः । किं तर्हि- यथा लोके‍= संसारे, उपकरणवताम्= धनादिसाधनवताम् जीवितम्= सुखोपायभोगसम्पन्नम्, तथैव= तद्वदेव तवापि जीवितं स्यात्, वित्तेन= वित्तसाध्येन कर्मणा मनसाऽपि= केवलं कल्पनयापि अमृतत्त्वस्य= अमरत्वस्य आशा= प्राप्तीच्छा नास्ति= न सम्पद्यते ।

\customparagraph{\textbf{नेपाली भावार्थ}}

मैत्रेयी भन्छिन्- हे स्वामि ! यो दृश्यमान पृथिवी मेरो स्वामित्वमा भयो भने के म अमर रहन सक्छु ? यस प्रश्नमा याज्ञवल्क्य उत्तर दिन्छन्- अहँ सक्दिनौ, तिमी अमर बन्न कदापि सक्दिनौ, किनकि संसारमा धन हुने मानिस केवल प्राण धारण गरेर निश्चित अवधिसम्म सुखपूर्वक बाँच्न सक्छन् तर अमर बन्न सक्दैनन् । तिम्रो पनि अवस्था त्यस्तै हुन्छ । धनसाध्य वस्तुबाट अमरत्वको कल्पनासम्म पनि गर्न सकिँदैन ।

\begin{visheshshloka}[center]

\textbf{सा होवाच मैत्रेयी येनाहं नामृता स्यां किमहं तेन कुर्यां यदेव भगवान्वेद तदेव मे ब्रूहि ॥३॥}

\end{visheshshloka}

\customparagraph{संस्कृतव्याख्या}

सा मैत्रेयी होवाच- यद्येवं चेत्, येन अहम् अमृता= मृत्युरहिता न स्याम्= न भवामि, तस्मात् अहं तेन= वित्तेन किम् कुर्याम्= मम कृते तस्य प्रयोजनं किमिति भावः। अतो भगवान्= याज्ञवल्क्यः, यदेव= अमृततत्त्वसाधनम्, वेद= जानाति तदेव= अमृतत्त्वसाधनम्, मे= मह्यम्, ब्रूहि= उपदिशतु ।

\customparagraph{\textbf{नेपाली भावार्थ}}

विदुषी मैत्रेयी भन्छिन्- यदि म सारा सम्पत्ति पाएर पनि मृत्युरहित हुन सक्दिनँ भने मेरा लागि त्यो सम्पत्तिको के प्रयोजन हुन्छ र ! त्यस कारण हे ऐश्वर्यवान् याज्ञवल्क्य ! हजुर जुन अमृतत्वका बारेमा जान्नुहुन्छ, त्यस अमृतत्वको मलाई उपदेश गर्नुहोस् ।

\begin{visheshshloka}[center]

\textbf{स होवाच याज्ञवल्क्यः प्रिया बतारे नः सती प्रियं भाषस एह्यास्स्व व्याख्यास्यामि ते व्याचक्षाणस्य तु मे निदिध्यासस्वेति॥४॥}

\end{visheshshloka}

\customparagraph{संस्कृतव्याख्या}

स याज्ञवल्क्यो होवाच- वत= इत्यनुकम्पायाम्- अरे मैत्रेयि ! नः= अस्माकं पूर्वमपि प्रिया = स्नेहशीला सती इदानीं प्रियम्= चित्तानुकूलमेव भाषसे= कथयसि, अत एहि= अत्रागच्छ, आस्स्व= उपविश (अहम्) व्याख्यास्यामि= यत्तव इष्टम् अमृतत्त्वसाधनम् आत्मज्ञानं कथयिष्यामि, व्याचक्षणस्य= व्याख्यानं कुर्वतो मे= मम निदिध्यासस्व= वाक्यान्यर्थतो निश्चयेन ध्यातुमिच्छ इति ।

\customparagraph{\textbf{नेपाली भावार्थ}}

पत्नी मैत्रेयीको जिज्ञासाको उत्तर दिँदै याज्ञवल्क्य भन्छन्- हे प्रिया ! मेरा लागि त‍िमी पहिलादेखि नै प्रिय थियौ । अहिले पनि मलाई प्रिय लाग्ने यति मीठो कुरा बोल्यौ । ठिक छ । यता आऊ । यहाँ बस । म तिमीले सोधेका कुराको व्याख्या गर्छु । तिमीले मैले व्याख्या गरेका वाक्यार्थको गहिरिएर चिन्तन गर्नू ।

\begin{visheshshloka}[center]

\textbf{स होवाच न वा अरे पत्युः कामाय पतिः प्रियो भवत्यात्मनस्तु कामाय पतिः प्रियो भवति । न वा अरे जायायै कामाय जाया प्रिया भवत्यात्मनस्तु कामाय जाया प्रिया भवति । न वा अरे पुत्राणां कामाय पुत्राः प्रिया भवन्त्यात्मनस्तु कामाय पुत्राः प्रिया भवन्ति । न वा अरे वित्तस्य कामाय वित्तं प्रियं भवत्यात्मनस्तु कामाय वित्तं प्रियं भवति । न वा अरे ब्रह्मणः कामाय ब्रह्म प्रियं भवत्यात्मनस्तु कामाय ब्रह्म प्रियं भवति । न वा अरे क्षत्रस्य कामाय क्षत्रं प्रियं भवत्यात्मनस्तु कामाय क्षत्रं प्रियं भवति । न वा अरे लोकानां कामाय लोकाः प्रिया भवन्त्यात्मनस्तु कामाय लोकाः प्रिया भवन्ति । न वा अरे देवानां कामाय देवाः प्रिया भवन्त्यात्मनस्तु कामाय देवाः प्रिया भवन्ति । न वा अरे भूतानां कामाय भूतानि प्रियाणि भवन्त्यात्मनस्तु कामाय भूतानि प्रियाणि भवन्ति । न वा अरे सर्वस्य कामाय सर्वं प्रियं भवत्यात्मनस्तु कामाय सर्वं प्रियं भवति । आत्मा वा अरे द्रष्टव्यः श्रोतव्यो मन्तव्यो निदिध्यासितव्यो मैत्रेय्यात्मनो वा अरे दर्शनेन श्रवणेन मत्या विज्ञानेेनेदं सर्वं विदितम् ॥५॥}

\end{visheshshloka}

\customparagraph{संस्कृतव्याख्या}

स होवाच= अमृतत्त्वसाधनं वैराग्यमुपदिदिक्षुर्याज्ञवल्क्यो जायापतिपुत्रादिभ्यो विरागमुत्पादयति यत् - अरे मैत्रेयि ! पत्युः कामाय= प्रयोजनाय जायायाः= पत्न्याः, पतिः प्रियो न भवति ।  किं तर्हि- आत्मनस्तु स्वकीयप्रयोजनाय एव भार्यायाः कृते पतिः प्रियो भवति । तथा न वा अरे जायायै इत्यादि समानमन्यत्, न वा अरे पुत्राणाम्, न वा अरे वित्तस्य, न वा अरे ब्रह्मणः, न वा अरे क्षत्रस्य, न वा अरे लोकानाम्, न वा अरे देवानाम्, न वा अरे भूतानाम्, न वा अरे सर्वस्य इत्यादिषु वाक्येषु पूर्वं पूर्वं यथा वर्णितम्, ते प्रीतिसाधने कृतम्, तत्र तत्रेष्टतरत्वाद्वैराग्यस्य। सर्वग्रहणमुक्तानुक्तार्थानां कृते ।

तस्माल्लोकप्रसिद्धमेतत्- आत्मैव प्रियः, नान्यः । तस्मादात्मप्रीतिसाधनत्वाद्गौणी अन्यत्र प्रीतिः, आत्मन्येव मुख्या । तस्मादात्मा वै अरे द्रष्टव्यः= दर्शनार्हः, दर्शनविषयतामापादयितव्यः; श्रोतव्यः= पूर्वमाचार्यत आगमतश्च; पश्चान्मन्तव्यः= ज्ञातव्यस्तर्कतः; ततो निदिध्यासितव्यः= निश्चयेन ध्यातव्यः; एवं प्रकारेण श्रवणमनननिदिध्यासनसाधनैर्निर्वर्तितैर्दृष्टः, साक्षात्कारो भवति। यदैकत्वमेतान्युपगतानि, तदा सम्यग्दर्शनं ब्रह्मैैकत्वविषयं प्रसीदति, नान्यथा श्रवणमात्रेण ।

यद् ब्रह्मक्षत्रादिकर्मनिमित्तं वर्णाश्रमादिलक्षणम् आत्मनि- अविद्याध्यारोपितप्रत्ययविषयं क्रियाकारकफलात्मकम्- अविद्याप्रत्ययविषयम्- रज्ज्वामिव सर्पप्रत्ययः । तस्माद् निवृत्तये श्रुतिः कथयति, यत्- आत्मनि खल्वरे मैत्रेयि ! दृष्टे श्रुते मते विज्ञाते इदं सर्वं विदितं विज्ञातं भवति ।

\customparagraph{\textbf{नेपाली भावार्थ}}

हे मैत्रेयि ! पतिको प्रयोजनका लागि पति प्रिय हुँदैन, आफ्नै प्रयोजनका लागि पत्नीलाई पति प्रिय हुन्छ । त्यस्तै पत्नीको प्रयोजनका लागि पत्नी प्रिय हुन्न, त्यहाँ पतिको आफ्नै प्रयोजन हुन्छ। पुत्रको प्रयोजनका लागि पुत्र प्रिय हुँदैन, त्यहाँ आफ्नै प्रयोजन रहेको हुन्छ । धनको प्रयोजनका लागि धन प्रिय हुँदैन, आफ्नै प्रयोजनका लागि धन प्रिय हुन्छ । ब्राह्मणको प्रयोजनका लागि ब्राह्मण प्रिय हुने होइन, त्यहाँं पनि आफ्नै स्वार्थ रहेको हुन्छ । क्षत्रियको प्रयोजनका लागि क्षत्रिय प्रिय हुँदैन, त्यहाँ आफ्नै प्रयोजन रहेको हुन्छ । लोकको प्रयोजनका लागि लोक प्रिय हुँदैन, त्यहाँ आफ्नै प्रयोजन रहन्छ । देवताका लागि देवता प्रिय हुँदैनन्, त्यहाँ आफ्नै कार्यसिद्धि रूप प्रयोजन रहन्छ । प्राणीहरूको प्रयोजनका लागि प्राणी प्रिय हुँदैनन् । त्यहाँ पनि आफनै स्वार्थ लुकेको हुन्छ । सबैको प्रयोजनका लागि सबै प्रिय हुँदैनन्, ती सबैमा आ-आफ्नै स्वार्थ रहेको हुन्छ ।

हे मैत्रेयि ! यो जुन आत्मा छ त्यो नै दर्शन गर्न योग्य, श्रवण गर्न योग्य, मनन गर्न योग्य र ध्यान गर्न योग्य छ । यही आत्माको दर्शन, श्रवण, मनन एवं गहिरो ज्ञानले यी सबै तिमीले बुझ्न चाहेका कुराको पूर्ण ज्ञान हुन्छ।

\begin{shloka}[left]

ननु कथमन्यस्मिन्विदितेऽन्यद्विदितं भवति ? नैष दोषः, न हि आत्मव्यतिरेकेणान्यत्किञ्चिदस्ति; यद्यस्ति न तद्विदितं स्यात्; न त्वन्यदस्ति; आत्मैव तु सर्वम्; तस्मात्सर्वमात्मनि विदिते विदितं स्यात् । कथं पुनरात्मैव सर्वमित्येतच्छ्रावयति-

\end{shloka}

\textbf{नेपाली भावार्थ}

शङ्का- तर अन्य कुराको ज्ञान हुँदा त्योभन्दा भिन्न वस्तुको ज्ञान कसरी हुन्छ ?

समाधान- यसमा कुनै दोष छैन । किनकि आत्मालाई छाडेर अन्य कुनै भिन्न वस्तु छैन । यदि हुन्थ्यो भने त्यसको ज्ञान पनि हुँदैनथ्यो । अन्य वस्तु छँदैछैन । आत्मा नै सबै हो । त्यसैले आत्माको ज्ञान भए सबै ज्ञातव्य कुराको ज्ञान हुन्छ । तथापि आत्मा नै सबै कसरी हुन्छ ? यी सबै वेदवाच्य छन् । यही कुरा याज्ञवल्क्य भन्छन् -

\begin{visheshshloka}[center]

\textbf{ब्रह्म तं परादाद्योऽन्यत्रात्मनो ब्रह्म वेद, क्षत्रं तं परादाद्योऽन्यत्रात्मनः क्षत्रं वेद, लोकास्तं परादुर्योऽन्यत्रात्मानो लोकान्वेद, देवास्तं परादुर्योऽन्यत्रात्मनो देवान्वेद, भूतानि तं परादुर्योऽन्यत्रात्मनो भूतानि वेद, सर्वं तं परादाद्योऽन्यत्रात्मनः सर्वं वेदेदं ब्रह्मेदं क्षत्रमिमे लोका इमे देवा इमानि भूतानीदं सर्वं यदयमात्मा ॥६॥}

\end{visheshshloka}

\customparagraph{संस्कृतव्याख्या}

ब्रह्म= ब्रह्मणजातिः, तम्= पुरुषम्, परादात्= परादध्यात्, पराकुर्यात्; कम् ? यो अन्यत्रात्मनः = आत्मस्वरूपव्यतिरेकेण आत्मैव न भवतीयं ब्राह्मणजातिरिति तां यो वेद, तं पराध्यात्  सा ब्राह्मणजातिरनात्मस्वरूपेण मां पश्यतीति; परमात्मा हि सर्वेषामात्मा ।

क्षात्रम् क्षत्रियजातिः, तथा लोका देवा भूतानि सर्वम् । येषां ‘इदं ब्रह्म’ ‘इदं क्षत्रम्’ इत्यादि रूपेण यानि अनुक्रान्तानि तानि सर्वाणि आत्मैव । यदयम् आत्मा= योऽयमात्मा द्रष्टव्यः श्रोतव्य इति प्रकृतः; यस्मादात्मनो जायते, आत्मन्येव लीयते, आत्ममयं च स्थितिकाले, आत्मव्यतिरेकेणाग्रहणात्, आत्मैव सर्वम् ।

\customparagraph{\textbf{नेपाली भावार्थ}}

जसले ब्राह्मणजाति आत्माभन्दा भिन्न भन्ने बुझ्छ, त्यसलाई ब्राह्मणजाति परास्त गरिदिन्छ । जसले क्षत्रिय जातिलाई आत्माभन्दा भिन्न मान्छ, त्यसलाई क्षत्रिय जातिले परास्त गर्दछ । जसले लोकलाई आत्माभन्दा भिन्न मान्छ, त्यसलाई लोकले नै परास्त गर्दछ। देवगणलाई जसले आत्माभन्दा भिन्न मान्छ त्यसलाई देवगणले नै परास्त गर्दछन् । जसले भूतगणलाई आत्माभन्दा भिन्न मान्छ, त्यसलाई भूतगणले नै परास्त गर्दछन् । अन्य सबै जेजति वस्तुहरू छन् तिनलाई जसले आत्माभन्दा भिन्न मान्छ तिनलाई ती सबै वस्तुले परास्त गर्छन् । यो ब्राह्मण, यो क्षत्रिय, यी लोकहरू, यी देवगण, यी भूतगण र अन्य जेजति छन्, ती सबै आत्ममय छन् । आत्माभन्दा भिन्न केही छैन ।

\begin{shloka}[left]

कथं पुनरिदानीमिदं सर्वमात्मैवेति ग्रहीतुं शक्यते ?

चिन्मात्रानुगमात्सर्वत्र चित्वरूपतैवेति गम्यते । तत्र दृष्टान्त उच्यते- यत्स्वरूपव्यतिरेकेणाग्रहणं यस्य, तस्य तदात्मत्वमेव लोके दृष्टम् ।

\end{shloka}

\customparagraph{\textbf{नेपाली भावार्थ}}

प्रश्न- किन्तु यस अवस्थामा ‘यो सबै आत्मा नै हो’ यस्तो विचार कसरी मान्न सकिन्छ ?

उत्तर- सबैतिर चिन्मात्राको अनुवृत्ति हुनाले सबै चिजको चित्स्वरूपता नै हुन्छ,  यस्तो बुझिन्छ । यस विषयमा एउटा दृष्टान्त बताइएको छ- जसलाई जसको शरीरबाट अलग गर्न सकिन्छ त्यो तद्रूप नै हुन्छ, यस्तो लोकमा देखिन्छ ।

\begin{visheshshloka}[center]

\textbf{स यथा दुन्दुभेर्हन्यमानस्य न बाह्याञ्शब्दाञ्शक्नुयाद्‍ग्रहणाय}

\textbf{दुन्दुभेस्तु ग्रहणेन दुन्दुभ्याघातस्य वा शब्दो गृहीतः ॥७॥}

\end{visheshshloka}

\customparagraph{संस्कृतव्याख्या}

सः दृष्टान्तः यथा- लोके यथा दुन्दुभेर्भेर्यादेः दण्डादिना  ताड्यमानस्य बाह्यान् बहिर्भूतान् शब्दान् दुन्दुभिशब्दसामान्यान्निकृष्टान् दुन्दुभिशब्दविशेषान् ग्रहीतुम् न शक्नुयात्= न शक्यते; दुन्दुभेर्जातस्य शब्दस्य सामन्यविशेषत्वेन ग्रहणेन दुन्दुभिशब्दा एते, इति शब्दविशेषा गृहीता भवन्ति, तेषां दुन्दुभिशब्दसामान्यव्यतिरेकेणाभावात् ।

दुन्दुभ्याघातस्य वा, दुन्दुभेराहननम् आघातः, दुन्दुभ्याघातविशिष्टस्य शब्दसामान्यस्य ग्रहणेन तद्गता विशेषा गृहीता भवन्ति, न तु त एव निर्भिद्य ग्रहीतुं शक्यन्ते, तेषां विशेषरूपेणाभावात् । तथा प्रज्ञानव्यतिरेकेन स्वप्नजागरितयोर्न कश्चिद्वस्तुविशेषो गृह्यते, तस्मात्प्रज्ञानव्यतिरेकेण अभावो युक्तस्तेषाम् ॥७॥

\customparagraph{\textbf{नेपाली भावार्थ}}

त्यो दृष्टान्त यस प्रकार छ- लोकव्यवहारमा जसरी दण्डादिले प्रहार गरिएका दुन्दुभि, भेरी आदि वाद्यविशेषबाट बाहिर फैलिएका शब्दविशेषलाई (दुन्दुभिको सामान्य शब्दबाट निकालिएको विशेष शब्दलाई) कसैले ग्रहण गर्न सक्दैन । दुन्दुभिको ग्रहणबाटै अर्थात् दुन्दुभिको सामान्य शब्दलाई विशेष रूपले ‘यो दुन्दुभिको शब्द हो ’ यस प्रकारले विशेष शब्द पनि ग्रहण हुन्छ । किनकि दुन्दुभिको सामान्य शब्दलाई छाडेर त्यसको सत्ता नै रहँदैन।

अथवा त्यो दुन्दुभिको आघात प्रहार रूप विशिष्ट शब्द, सामान्य शब्द ग्रहण हुँदा नै त्यस भित्र रहेको विशेष शब्दको पनि ग्रहण हुन्छ। त्यसलाई अलग गरेर छुट्टै त्यसको ग्रहण हुँदैन । किनकि विशेष रूपले त्यसको अभाव छ ।

यसरी नै स्वप्न र जागरितलाई कुनै वस्तु विशेषको प्रज्ञानबाट अलग ग्रहण गर्न सकिँदैन । त्यसैले प्रज्ञानबाट भिन्न त्यसको अभाव उचित नै हो ।

(त्यो दृष्टान्त यस्तो छ कि जसरी प्रहार गरिएको दुन्दुभिको बाह्य शब्दहरूलाई कसैले पनि ग्रहण गर्न सक्दैनन्, किनकि दुन्दुभि या दुन्दुभिको प्रहारलाई ग्रहण गर्दा नै त्यसको शब्द पनि ग्रहण गरिएको हुन्छ )

\begin{visheshshloka}[center]

\textbf{स यथा शङ्‍खस्य घ्मायमानस्य न बाह्याञ्शब्दाञ्शक्नुयाद्‍}

\textbf{ग्रहणाय शङ्‍खस्य तु ग्रहणेन शङ्‍खघ्मस्य वा शब्दो गृहीतः ॥८॥}

\end{visheshshloka}

\customparagraph{संस्कृतव्याख्या}

तथा सः= द्वितीयो दृष्टान्तो यथा- शङ्‍खस्य= वाद्यविशेषस्य घ्मायमानस्य= शब्देन संयोज्यमानस्य, आपूर्यमाणस्य न बाह्यान् शब्दाञ्छक्नुयात् (ग्रहीतुम्) इत्येवमादि पू्र्ववत्॥८॥

\customparagraph{\textbf{नेपाली भावार्थ}}

त्यसको दोस्रो उदाहरण यस प्रकार छ- जसरी फुकेर शब्दसंयोजन गरिएको शङ्‍खको बाह्य शब्दलाई छुट्याएर कसैले ग्रहण गर्न सक्दैन इत्यादि पूर्ववत् यसको अर्थ हो ।

(दोस्रो दृष्टान्त यस प्रकार छ- कोही व्यक्ति बजाउँदै गरेका शङ्खको बाह्य शब्द छुट्टै ग्रहण गर्न सक्दैन किनकि शङ्ख वा शङ्ख बजाइएको ग्रहण गर्दैमा त्यो बाह्य शब्दको पनि ग्रहण हुन्छ)

\begin{visheshshloka}[center]

\textbf{स यथा वीणायै वाद्यमानायै न बाह्याञ्शब्दाञ्शक्नुयाद् ग्रहणाय वीणायै तु ग्रहणेन वीणावादस्य वा शब्दो गृहीतः ॥९॥}

\end{visheshshloka}

\customparagraph{संस्कृतव्याख्या}

तथा (तृतीयदृष्टान्तः) वाद्यमानायै वीणायै= वाद्यमानाया वीणायाः, अनेकदृष्टान्तोपादानमिह सामान्यबहुत्वख्यापनार्थम्- अनेके हि विलक्षणाचेतनाचेतनरूपाः सामान्यविशेषाः, तेषां पारम्पर्यगत्या यथैकस्मिन्महासामान्येऽन्तर्भावः प्रज्ञाघने; कथं नाम प्रदर्शयितव्य इति; दुन्दुभिशङ्‍खवीणाशब्दसामान्यविशेषाणां यथा शब्दत्वेऽन्तर्भावः, एवं स्थितिकाले तावत्सामान्यविशेषव्यतिरेकाद् ब्रह्मैकत्वं शक्यमवगन्तुम् ॥९॥

\customparagraph{\textbf{नेपाली भावार्थ}}

अगाडिका दुई उदाहरणझैँ ‘वीणायै वाद्यमानायै’ को पनि बुझ्नुपर्छ । यहाँ अनेक दृष्टान्तहरूको ग्रहण सामान्यको बहुलता स्पष्ट गर्नका लागि हो । चेतन र अचेतन, सामान्य र विशेष अनेक एवं विलक्षण छन् । तिनीहरूको जसरी परम्पराअनुसार एक प्रज्ञानघन महासामान्यमा अन्तर्भाव हुन्छ, भन्ने तथ्यलाई कुनै न कुनै रूपले देखाउनु यी विभिन्न उदाहरणहरूको उद्देश्य हो । जसरी उदाहरणमा आएका दुन्दुभि, शङ्ख र वीणाको सामान्य एवं विशेष शब्दहरूको शब्दत्वमा अन्तर्भाव हुन्छ त्यसरी नै स्थितिकालमा सामान्य र विशेषले अभिन्न हुनाले ब्रह्मको एकत्वको ज्ञान हुन सक्छ।

(त्यो तेस्रो दृष्टान्त यस्तो छ- जस्तै कोही बजाइएको वीणाको बाह्य शब्दलाई ग्रहण गर्न समर्थ हुँदैन; तर वीणा या वीणाको स्वर ग्रहण हुँदा त्यो बाह्य शब्दको पनि ग्रहण हुन्छ)

\begin{shloka}[left]

\textit{\textbf{परमात्माको निश्वासभूत ऋग्वेदादिको परमात्माभन्दा अभिन्नत्व प्रतिपादन-}}

एवमुत्पत्तिकाले प्रागुत्पत्तेर्ब्रह्मैवेति शक्यमवगन्तुम् । यथाग्नेर्विस्फुलिङ्गधूमाङ्गारार्चिषां प्राग्विभागादग्निरेवेति भवत्यग्न्यैकत्वम्, एवं जगन्नामरूपविकृतं प्रागुत्पत्तेः प्रज्ञानघन एवेति युक्तं ग्रहितुमित्येतदुच्यते-

\end{shloka}

\customparagraph{\textbf{नेपाली भावार्थ}}

यसरी यो स्पष्ट हुन्छ कि उत्पत्तिकालमा उत्पत्तिभन्दा पहिले ब्रह्म नै थियो । जसरी अग्निको झिल्का, धुवाँ, कोइला र ज्वालाहरूको अस्तित्व विभाजन हुनुभन्दा पूर्व अग्नि नै थियो । त्यसैले अग्निको एकत्व सिद्ध हुन्छ । त्यसरी नै नाम, रूप विकारलाई प्राप्त गरेको जगत् उत्पत्तिभन्दा पहिले प्रज्ञानघन ब्रह्म मात्रै थियो । यस्तो जान्न उचित हुन्छ । यस सन्दर्भमा यस प्रकार भनिन्छ-

\begin{visheshshloka}[center]

\textbf{स यथार्द्रैधाग्नेरभ्याहितात्पृथग्धूमा विनिश्चरन्त्येवं वा अरेऽस्य महतो भूतस्य निश्वसितमेतद्यदृग्वदो यजुर्वेदः सामवेदोऽथर्वाङ्गिरस इतिहासः पुराणं विद्या उपनिषदः श्लोकाः सूत्राण्यनुव्याख्यानानि व्याख्यानान्यस्यैवतानि निश्वसितानि ॥१०॥}

\end{visheshshloka}

\customparagraph{संस्कृतव्याख्या}

सः दृष्टान्तः (चतुर्थो दृष्टान्तः) यथा- आर्द्रैधाग्नेः= आर्द्रैरेधोभिरिद्धोऽग्निः,  तस्मादभ्याहितात्= नियोजितात् पृथक् पृथक् धूमाः= धूमविस्फुल्लिङ्गादयः (धूमग्रहणं विस्फुल्लिङ्गादिप्रदर्शनार्थम्) विनिश्चरन्ति= विनिर्गच्छन्ति ।

एवम्- यथायं दृष्टान्तः, अरे मैत्रेयी ! अस्य परमात्मनः प्रकृतस्य महतो भूतस्य निश्वसितमेतत् । किं तन्निश्वसितमिव ततो जातमित्युच्यते, यत्- ऋग्वेदो यजुर्वेदः सामवेदोऽथर्वाङ्गिरसः चतुर्विधं मन्त्रजातम्, \textbf{इतिहासः}- उर्वशीपुरूरवसोः संवादादिः ‘उर्वशी हाप्सराः’ इत्यादिब्राह्मणमेव, \textbf{पुराणम्}- ‘असद्वा इदमग्र आसीत्’ (तै. उ. २।६।१) इत्यादि, \textbf{विद्याः}- देवजनविद्या वेदः सोऽयमित्याद्याः, \textbf{उपनिषदः}- ‘प्रियमित्येतदुपासीत’ (बृ.उ. ४।१।३) इत्याद्याः, \textbf{श्लोकाः}- ब्राह्मणप्रभवा मन्त्राः ‘तदेते श्लोकाः’ (बृ.उ. ४।३।११) इत्यादयः, \textbf{सूत्राणि}- वस्तुसङ्‍ग्रहवाक्यानि वेदे यथा- ‘आत्मेत्येवोपासीत’ (१।४।७) इत्यादीनि, \textbf{अनुव्याख्यानानि}- मन्त्रविवरणानि, व्याख्यानान्यर्थवादाः, अथवा वस्तुसङ्‍ग्रहवाक्यविवरणान्यनुव्याख्यानानि यथा- चतुर्थाध्याये ‘आत्मेत्येवोपासीत’ इत्यस्य, यथा वा- ‘अन्योऽसावन्योऽहमस्मीति न स वेद यथा पशुरेवम्’ (१।४।१०) इत्यस्यायमेवाध्यायशेषः, मन्त्रविवरणानि \textbf{व्याख्यानानि}, एवमष्टविधं ब्राह्मणम् । एतावतानि सर्वाण्येव ब्रह्मणो निश्वासजातानीत्यर्थः॥१०॥

\customparagraph{\textbf{नेपाली भावार्थ}}

(चौथो दृष्टान्त) जसरी भिजेका दाउराले आगो बाल्दा विभिन्न प्रकारका धुवाँ निस्कन्छन्, अर्थात् झिल्का, ज्वाला  आदिसहितको धुवाँ निस्कन्छ । (यहाँ धुवाँ शब्द झिल्का, ज्वाला आदिलाई बुझाउनका लागि प्रयुक्त छ) यसरी नै यहाँ दृष्टान्त छ कि हे मैत्रेयि ! परब्रह्म, जुन प्रकृत महद्भुतका रूपमा रहेको छ, यो चराचर जगत् त्यसैको निश्वास हो । जसरी विना प्रयत्न पुरुष निश्वास लिन्छ अरे ! मैत्रेयि ! त्यसरी नै यो जगत् त्यही महान् पुरुषको निश्वासबाट उत्पन्न भएको हो ।

त्यस महद्भुतको निश्वासबाट के के निस्किए त्यो बताइन्छ- लोकप्रसिद्ध यी ऋग्वेद, यजुर्वेद, सामवेद र अथर्ववेद गरी चार प्रकारका मन्त्र समुदाय छन्; \textbf{इतिहास}- जस्तै उर्वशी र पुरूरवाको संवाद आदि ‘उर्वशी हाप्सराः’ इत्यादि ब्राह्मण नै इतिहास हुन् । \textbf{पुराण}- ‘आरम्भमा यो असत् नै थियो’ इत्यादि; \textbf{विद्या}- ‘वेदः सोऽयम्’ इत्यादि देवजनविद्या; \textbf{उपनिषद्-} ‘प्रिय छ, यसलाई यसरी उपासना गर्नुपर्छ’ इत्यादि; \textbf{श्लोकः-} ‘तदेते श्लोकाः’ इत्यादि ब्राह्मण भागका मन्त्र; \textbf{सूत्र-} ‘वस्तुसङ्ग्रहवाक्यानि’ जस्तै वेदमा- ‘आत्मा छ, जसलाई यसरी उपासना गर्नुपर्छ’ इत्यादि मन्त्र छन् । \textbf{अनुव्याख्यान-} मन्त्रविवरण, \textbf{व्याख्यान-} अर्थवाद, अथवा वस्तुसङ्‍ग्रह वाक्यको विवरण नै अनुव्याख्यान हो । जसअनुसार चौथो अध्यायमा ‘आत्मेत्येवोपासीत’ यो वाक्यको व्याख्या छ, अथवा ‘अन्योऽसावन्योऽहमस्मीति न स वेद यथा पशुरेवम्’ यो वाक्यको व्याख्यान यो शेष अध्याय हो । मन्त्रविवरणको अर्थ मन्त्रव्याख्यान हो । यसरी इतिहास आदि पदले सङ्‌केत गरेका आठ प्रकारका ब्राह्मण भाग छन् । इत्यादि सबै त्यही महापुरुष परब्रह्मका निश्वासबाट उत्पन्न भएका हुन् ।

\customparagraph{\textit{\textbf{आत्मा नै सबैको आश्रय हो- यस सम्बन्धमा दृष्टान्त}}}

\begin{shloka}[left]

किञ्चान्यत्, न केवलं स्थित्युत्पत्तिकालयोरेव प्रज्ञानव्यतिरेकेणाभवाज्जगतो ब्रह्मत्वम्, प्रलयकाले च । जलबुद्‍बुदफेनादीनामिव सलिलव्यतिरेकेणाभावः, एवं प्रज्ञानव्यतिरेकेण तत्कार्याणां नामरूपकर्मणां तस्मिन्नेव लीयमानानामभावः । तस्मादेकमेव ब्रह्म प्रज्ञानघनमेकरसं प्रतिपत्तव्यमित्यत आह । प्रलयदर्शनाय दृष्टान्तः -

\end{shloka}

\customparagraph{\textbf{नेपाली भावार्थ}}

योभन्दा अर्को कुरा यो छ कि जगत्‌को ब्रह्मत्व केवल उत्पत्ति र स्थितिकालमा मात्र प्रज्ञानलाई छाडेर नरहनाको कारण होइन, अपितु प्रलयकालमा पनि हो । जसरी पानी, बुद्‍बुद र फिँज आदिको अस्तित्व पानीलाई छाडेर छैन, त्यसरी नै प्रज्ञानभन्दा भिन्न त्यसका कार्य र त्यसैमा लीन हुने नाम, रूप र कर्महरूको पनि सत्ता रहँदैन । त्यसैले एउटै प्रज्ञाघन एक रस ब्रह्म हो- यस्तो बुझ्नुपर्छ । यसैबाट श्रुति (निम्नाङ्कित मन्त्र) ले स्पष्ट पार्दछ । प्रलय प्रदर्शित गर्नलाई यहाँ दृष्टान्त दिइएको छ-

\begin{visheshshloka}[center]

\textbf{स यथा सर्वासामपां समुद्र एकायनमेवं सर्वेषां स्पर्शानां त्वगेकायनमेवं सर्वेषां गन्धानां नासिके एकायनमेवं सर्वेषां रसानां जिह्वैकायनमेवं सर्वेषां रूपाणां चक्षुरेकायनमेवं सर्वेषां शब्दानां श्रोत्रमेकायनमेवं सर्वेषां सङ्‍कल्पानां मन एकायनमेवं सर्वासां विद्यानां हृदयमेकायनमेवं सर्वेषां कर्मणां हस्तावेकायनमेवं सर्वेषामानन्दानामुपस्थ एकायनमेवं सर्वेषां विसर्गाणां पायुरेकायनमेवं सर्वेषामध्वानां पादावेकायनमेवं सर्वेषां वेदानां वागेकायनम् ॥११॥}

\end{visheshshloka}

\customparagraph{संस्कृतव्याख्या}

स इति दृष्टान्तः, यथा= येन प्रकारेण सर्वासाम्= नदीवापीतडागादिगतानामपाम्, समुद्रः= अब्धिः, एकायनम्= एकगमनमेकप्रलयोऽविभागप्राप्तिरित्यर्थः; यथायं दृष्टान्तः; एवं सर्वेषाम्= स्पर्शानां मृदुकर्कशकठिनपिच्छिलादीनां वायोरात्मभूतानां त्वगेकायनम्, त्वगिति त्वग्विषयं स्पर्शसामान्यमात्रम्, तस्मिन्प्रविष्टाः स्पर्शविशेषाः  आप इव समुद्रम्- तद्‍व्यतिरेकेणाभावभूता भवन्ति; तस्यैव हि ते संस्थानमात्रा आसन्।

तथा तदपि स्पर्शसामान्यमात्रं त्वक्छब्दवाच्यं मनःसङ्‍कल्पे मनोविषयसामान्यमात्रे, त्वग्विषय इव स्पर्शविशेषाः; प्रविष्टं तद्‍व्यतिरेकेणाभावभूतं भवति; एवं मनोविषयोऽपि बुद्धिविषयसामान्यमात्रे प्रविष्टस्तद्‍व्यतिरेकेणाभावभूतो भवति; विज्ञानमात्रमेव भूत्वा प्रज्ञानघने परे ब्रह्मण्याप इव समुद्रे प्रलीयते ।

\customparagraph{\textbf{नेपाली भावार्थ}}

त्यो दृष्टान्त यस प्रकार छ कि- सम्पूर्ण नदी, तलाउ, पोखरी आदिका जलको समुद्र नै एकमात्र प्रलय अर्थात् अभेद प्राप्तिको स्थान हो ।  त्यसरी नै हावाका विभिन्न स्वरूप जस्तै कोमल, कठोर, कर्कश, चिप्लो आदि समस्त स्पर्शहरूको एकमात्र प्रलय स्थान छाला हो । यहाँ छाला भन्नाले छालासम्बन्धी स्पर्शसामान्यमात्र सम्झनुपर्दछ । त्यसैमा समुद्रमा जल समान स्पर्श विशेष प्रविष्ट छ । समुद्र र छाला नहुँदा जल र स्पर्श सत्ताशून्य हुन्छन् किनकि ती तिनैका अवस्थितिमात्र भिन्न रहेका वस्तु हुन् ।

यसरी नै त्यो त्वक् शब्द वाच्य स्पर्श सामान्य, छालाको विषयमा स्पर्शविशेष समान मनको विषय सामान्यमात्र रूप मनको सङ्कल्पमा प्रविष्ट भएर त्योभन्दा भिन्न सत्ताशून्य हुन्छ । यसरी नै मनको विषय पनि बुद्धिको सामान्य विषयमात्रमा प्रवेश गरेर त्योभन्दा भिन्न रहँदैन एवम् विज्ञानमात्र नै भएर समुद्रमा जलझैं प्रज्ञानघन परब्रह्ममा लीन हुन्छ।

(त्यो दृष्टान्त जसरी समस्त पानीको समुद्र एक अयन (प्रलयस्थान) हो, यसरी नै समस्त स्पर्शहरूको छाला नै अयन स्थान हो, यसरी नै समस्त गन्धहरूको नासिका नै एक अयन स्थान हो, यसरी नै समस्त रूपहरूको चक्षु नै एक अयन स्थान हो, यसरी नै समस्त शब्दहरूको श्रोत्र एक अयन स्थान हो, यसरी नै समस्त सङ्कल्पहरूको मन एक अयन स्थान हो, यसरी नै समस्त विद्याहरूको हृदय एक अयन स्थान हो, यसरी नै समस्त कर्महरूको हात एक अयन हो, यसरी नै समस्त आनन्दहरूको उपस्थ एक अयन स्थान हो, यसरी नै समस्त विसर्ग (त्याग) हरूको पायु एक अयन स्थान हो, यसरी नै समस्त मार्गहरूको पाउ एक अयन स्थान हो, यसरी नै समस्त वेदहरूको वाणी एक अयन स्थान हो ।)

\customparagraph{\textit{\textbf{विवेकका माध्यमबाट देहादि विज्ञानघनस्वरूप हुन् भन्ने ज्ञानका लागि पानीमा लगाइएको नुनको दृष्टान्त-}}}

\begin{shloka}[left]

तत्र ‘इदं सर्वं यदयमात्मा’ (२।४।६।) इति प्रतिज्ञातम्, तत्र हेतुुरभिहितः - आत्मसामान्यत्वम्, आत्मजत्वम्, आत्मप्रलयत्वं च । तस्मादुत्पत्तिस्थितिप्रलयकालेषु प्रज्ञानव्यतिरेकेणाभावात् “प्राज्ञानं ब्रह्म” (ऐ. उ. ५।३।) “आत्मैवेदं सर्वम्” (छा. उ. ७।२५।२) इति प्रतिज्ञातं यत्, तत्तर्कतः साधितम् । स्वाभाविकोऽयं प्रलय इति पौराणिका वदन्ति । यस्तु बुद्धिपूर्वकः प्रलयो ब्रह्मविदां ब्रह्मविद्यानिमित्तः, अयमात्यन्तिक इत्याचक्षते- अविद्यानिरोधद्वारेण यो भवति; तदर्थोऽयं विशेषारम्भः -

\end{shloka}

\customparagraph{\textbf{नेपाली भावार्थ}}

त्यहाँ यो प्रतिज्ञा गरिएको छ कि ‘यो जो जे छ, त्यो सबै आत्मा हो’ यहाँ आत्मसामान्यत्व, आत्मजनितत्त्व र आत्मप्रलयत्व यी हेतु बताइएको छ । अतः उत्पत्ति, स्थिति र प्रलय कालमा प्रज्ञानभन्दा भिन्न केहीको पनि अस्तित्व नहुनाले जुन यस्तो प्रतिज्ञा गरिएको थियो कि “प्रज्ञान ब्रह्म हो” “यो सबै आत्मा नै हो” यसलाई तर्कले पनि सिद्ध गरिएको छ । ‘यो प्रलय स्वाभाविक हो’ यस्तो पौराणिक मानिसहरू भन्छन् । ब्रह्मवेत्ताहरूको जुन ब्रह्मविद्याजनित बुद्धिपूर्वक प्रलय हुन्छ, त्यो आत्यन्तिक हो- यस्तो भन्छन्, जुन अविद्याको निरोधबाट हुन्छ; त्यसैका लागि यो विशेष आरम्भ गरिन्छ-

\begin{visheshshloka}[center]

\textbf{स यथा  सैन्धवखिल्य उदके प्रास्त उदकमेवानुविलीयते न हास्योद्‍ग्रहणायेव स्यात् । यतो यतस्त्वाददीत लवणमेवैवं वा अर इदं महद्‍भुतमनन्तमपारं विज्ञानघन एव । एतेभ्यो भूतेभ्यः समुत्थाय तान्येवानु विनश्यति न प्रेत्य संज्ञास्तीत्यरे ब्रवीमीति होवाच याज्ञवल्क्यः ॥१२॥}

\end{visheshshloka}

\customparagraph{संस्कृतव्याख्या}

यथा= येन प्रकारेण जले प्रक्षिप्तं लवणखण्डं जलेन सह द्रवीभूत्वा लीनो भवति । तं लवणखण्डं पूर्ववत् समासादयितुं कोऽपि समर्थो न भवति । यस्माद् यस्मात् स्थानाज्जलं नीयते तत्र तत्र लवणमिश्रितमेव ज्ञायते । हे मैत्रेयि ! तेनैव प्रकारेण महद्‌भूतमनन्तमपारं विज्ञानघन एवास्ति । स एतेभ्यो भूतेभ्यः शरीरादिभ्यः समुत्थाय जलस्य फेनबुद्‌बुदादिवद् अन्ते तत्रैव विज्ञानघने/आत्मनि विलीनो भवति । ते प्रेत्यभावं प्राप्य संज्ञायुता न भवन्ति, तदा एक एव विज्ञानघनरूप आत्मा एवावतिष्ठते । हे मैत्रेयि ! एतादृशं रहस्यमहं त्वां कथयन्नस्मीति याज्ञवल्य उवाच ॥१२॥

\customparagraph{\textbf{नेपाली भावार्थ}}

जसरी पानीमा लगाइएको नुनको डल्लो केही समयपछि त्यही पानीमा विलीन हुन्छ । कसैले पनि पूर्ववत् प्राप्त गर्न सक्दैन । सबैतिर व्याप्त हुनाले पानी जहाँबाट लिए पनि नुनिलो नै हुन्छ । त्यस्तै पानीमा फिँज, फोका आदि समान त्यो विज्ञानघन आत्मतत्त्व यो सम्पूर्ण चराचर जगत्‌का नामरूपधारीका रूपमा प्रकट भएर अन्त्यमा त्यही आत्मामा नै विलीन हुन्छ । त्यसबेला केवल विज्ञानघन आत्मतत्त्वमात्र रहन्छ । हे मैत्रेयि ! तिमीलाई मैले यही गहिरो कुरा बताउने प्रयास गरिरहेको छु, यसरी याज्ञवल्क्य भन्छन्।

\customparagraph{(\textbf{विशेष स्पष्टीकरण}-}

यहाँ यो दृष्टान्त दिइन्छ- स यथा इत्यादि । सिन्धुको विकारको नाम सैन्धव हो । यहाँ सिन्धु शब्दले पानीलाई बुझाएको छ । स्यन्दन गर्नाले अर्थात् बहनाले यसलाई सिन्धु भनिएको हो । त्यसको विकार वा त्यसबाट उत्पन्न हुने वस्तुलाई सैन्धव भनिएको हो । जुन सैन्धव छ र खिल्य (डल्लो,ढिको) छ, त्यसैलाई सैन्धवखिल्य (नुनको ढिको) भनिन्छ । खिल शब्दमा स्वार्थमा यत् प्रत्यय भएर खिल्य बनेको हो । यो नुनको ढिको आफ्नै कारणभूत पानीमा लगाउँदा पानीको साथमा बिलेर त्यसैमा लीन हुन्छ। पार्थिव तेजसको सम्पर्कले जुन त्यो ढिकामा कठिनता प्राप्त भयो, त्यो आफ्नै कारणसँग संयोग हुँदा निवृत्त हुन्छ, यही पानीको घुलन हो, यसैका साथ त्यो नुनको ढिको पनि बिल्यो- यस्तो भनिन्छ। यसबाट यो भनिन्छ कि त्यो नुनको ढिको पानीका साथ लीन भयो ।

यो बिलाइसकेको नुनको ढिकालाई पूर्ववत् ग्रहण गर्न जतिसुकै विज्ञले पनि सक्दैन । ‘यहाँ प्रयुक्त इव शब्द अर्थहीन छ ।’ त्यसलाई ग्रहण गर्न समर्थ हुँदैन, किन ? किनकि जुनजुन ठाउँबाट त्यो व्यक्ति पानीलाई ग्रहण गर्छ, चाख्छ; त्यो पानी नुनकै स्वादको हुन्छ, त्यसमा ढिकोपन रहँदैन ।

यस प्रकार जस्तो यो दृष्टान्त छ, त्यस्तै हे मैत्रेयि ! यो परमात्मा नामको महद्‍भूत छ । जुन महद्‍भूतबाट तिमी अविद्याबाट भिन्न भई देहेन्द्रिय रूप उपाधिका सम्बन्धले खिल्यभावमा प्राप्त भएकी छ्यौ, एवम् मरणधर्मशील मान्छेका रूपमा जन्म, मृत्यु, क्षुधा र पिपासा आदि सांसारिक धर्मशील एवम् म नामरूप शरीरवाली अमुक वंशमा उत्पन्न भएकी छु- यस्तो भाववाली भएकी छ्यौ । देहेन्द्रियजनित उपाधिका सम्पर्कले भ्रान्तिका कारण उत्पन्न भएको तिम्रो त्यो खिल्यभाव आफ्नै कारणले महासमुद्र स्थानीय अजर, अमर, अभय, शुद्ध, सैन्धवघन तुल्य एकरस प्रज्ञाघन, अनन्त, अपार, अखण्ड एवम् अविद्या जनित भ्रान्तिमय भेदले रहित भई परमात्मामा प्रविष्ट गरिएको छ ।

त्यसमा प्रविष्ट हुनाले त्यो खिल्यभावको आफ्नै कारणले लीन गराइएपछि अविद्याजनित भेदभाव समाप्त हुनाले यो एक अद्वैत महद्‍भूत नै रहन्छ । महान् भूत हुनाले यो महद्‍भूत कहलिएको छ किनकि आकाशादिको पनि कारण हुनाले यो सबैभन्दा महान् छ । तीनवटै कालमा त्यसको स्वरूपको व्यभिचार हुँदैन । त्यो सधैँ जस्ताको त्यस्तै रहन्छ । त्यसैले त्यो ‘भूत’ हो। ‘भूत’ शब्दमा ‘त’ यो निष्ठाप्रत्यय त्रैकालिक छ ।

अथवा ‘भूत’ शब्द  परमार्थवाची छ । अर्थात् यो महत् हो र पारमार्थिक पनि हो । (त्यसैले महद्भूत हो) यद्यपि हिमालय आदि पर्वतहरू समान लौकिक वस्तु पनि महान् हुन्छन् तर ती स्वप्न वा मायाका समान हुन्छन्, परमार्थ वस्तु होइनन् । त्यसैले श्रुति यसलाई परिभाषित गर्दछ कि यो महत् र भूत पनि छ । अनन्त अर्थात् यसको अन्त्य छैन, त्यसैले यो अनन्त छ। कदाचित् यसको अनन्तता आक्षेपिक हुन सक्छ त्यसैले यसका लागि ‘अपारम्’ यो विशेषण दिइएको छ । विज्ञप्तिको नाम विज्ञान हो, जो विज्ञान र घन छ त्यसैलाई ‘विज्ञानघन’ भनिन्छ । यहाँ ‘घन’ शब्द अन्य जातिगत वस्तुको निषेधका लागि हो । जस्तै- सुवर्णघन, लोहघन आदि । ‘एव’ शब्द निश्चयार्थक हो- तात्पर्य यो हो कि यसभित्र कुनै अन्य विजातीय वस्तु छैन ।

यदि यो आत्मतत्त्व एक, अद्वैत, परमार्थतः शुद्ध र सांसारिक दुःखहरूले छोएको छैन भने आत्माको यो खिल्यभाव किन छ ? तथा म उत्पन्न भएँ, मरेँ, सुखी भएँ, दुःखी भएँ, अहम्, मम इत्यादि लक्षणवाला अनेक सांसारिक धर्मले दूषित किन छ ? यसका सम्बन्धमा भनिन्छ-

यी भूतहरूबाट जो यी देह र इन्द्रियरूप विषयका आकारमा परिणत वस्तुहरू छन् ती सबै पानीको फिँज र भुल्का वा फोका समान जलरूप स्वच्छ परमात्माका विकार हुन् । जसका सम्पूर्ण विषयसम्मका पारमार्थिक विवेक ज्ञानबाट प्रज्ञानघन ब्रह्ममा लीन हुनु बताइएको छ । यी सबैको हेतुभूत सत्य शब्दवाच्य भूतहरूबाट लवणखण्ड समान उत्पन्न भएर जसरी पानीबाट सूर्यचन्द्रादिको प्रतिबिम्ब अथवा जस्तै अलक्तक आदि उपाधिहरूका कारण स्वच्छ स्फटिकको रक्तादि भाव हुन जान्छ; त्यसरी नै देहेन्द्रियरूप उपाधिहरूका कारण विशेषात्मरूप खिल्यभावले उत्पन्न भएर जुन भूतबाट यो उत्पन्न भएको छ, त्यो देह र इन्द्रियका आकारमा परिणत एवम् आत्माको खिल्यभाव रूप विशेषत्वको हेतुभूत जुन समय, शास्त्र र आचार्यको ब्रह्मविद्यासम्बन्धी उपदेशले समुद्रमा नदी समान लीन हुँदै नाशलाई प्राप्त गर्दछ; पानीमा फिँज र भुल्का समान त्यसको नाश हुनासाथ नै यो विशेषात्मरूप खिल्यभाव पनि नष्ट हुन्छ । जसरी पानी र अलक्तक आदि हेतुहरू नाश हुँदा सूर्य, चन्द्र र स्फटिक आदिको प्रतिबिम्ब पनि नष्ट हुन्छन्; केवल चन्द्रादिको पारमार्थिक स्वरूप मात्रै रहन्छन्; त्यसरी नै (देहेन्द्रियको समाप्तिमा) अनन्त, अपार र स्वच्छ प्रज्ञानघन मात्र रहन्छ।

प्रेत्य-देहेन्द्रियभावले मुक्त भएर त्यसको विशेष संज्ञा रहँदैन, यसैबाट हे मैत्रेयि ! म यो भन्छु कि त्यसको ‘म अमुक  हुँ’ ‘म अमुकको छोरो हुँ’ ‘यो खेत र धन मेरो हो’ ‘म सुखी छु’ ‘म दुःखी छु’ इत्यादि प्रकारको विशेष संज्ञा रहँदैन किनकि त्यो त अविद्याजन्य ज्ञान हो र कारणसहित अविद्यालाई ब्रह्मविद्याले नाश गरिसकेको छ।  त्यसैले चैतन्यस्वरूप ब्रह्मवेत्ताको विशेष संज्ञा रहने सम्भावना कहाँ हुन्छ ? त्यसको त शरीरमा रहँदारहँदै पनि कुनै संज्ञा हुने सम्भावना रहँदैन । फेरि सबै प्रकारको देह र इन्द्रियहरूबाट मुक्त हुँदा त रहन कसरी सक्छ र !  यसरी याज्ञवल्क्यले आफ्नी पत्नी मैत्रेयीलाई परमार्थदृष्टिको निरूपण गर्नुभयो)

\customparagraph{\textit{\textbf{मैत्रेयीको शङ्का र याज्ञवल्क्यको समाधान}}}

एवं प्रवोधिता-

\begin{visheshshloka}[center]

\textbf{सा होवाच मैत्रेय्यत्रैव मा भगवानमूमुहन्न प्रेत्य संज्ञास्तीति स होवाच न वा}

\textbf{अरेऽहं मोहं ब्रवीम्यलं वा अर इदं विज्ञानाया॥१३}॥

\end{visheshshloka}

\customparagraph{संस्कृतव्याख्या}

सा मैत्रेयी होवाच- शरीरपातानन्तरं किमपि संज्ञादिकं न भवतीति कथयित्वा ऐश्वर्यवान् भवान् माम् अमूमुहत्= मोहञ्चकार । याज्ञवल्क्य उत्तरयति यत्- हे मैत्रेयि ! अहं त्वां मोहप्रवर्धकं वचनमुपदिशन् नास्मि, यद् मया उपदिष्टोऽर्थो विद्यते, स  महद्‌भूतस्य आत्मतत्त्वस्य विज्ञानाय पर्याप्तोऽस्ति । त्वमेव विपरीतं चिन्तयसीत्यर्थः ।

\customparagraph{\textbf{नेपाली भावार्थ}}

याज्ञवल्क्यको उपदेश सुनेपछि मैत्रेयी भन्छिन्- ‘शरीरको त्यागपछि कुनै संज्ञा रहँदैन, भनेर हजुरले मलाई मोहित बनाउनुभयो ।’ यसको उत्तरमा याज्ञवल्क्य भन्छन्- हे मैत्रेयि ! म तिमीलाई मोहको उपदेश गरिरहेको छैन, अरे ! यो त त्यो (महद्‍भूतको) विशेष ज्ञान गराउनका लागि पर्याप्त छ॥१३॥

( \textbf{विशेष स्पष्टीकरण}- मैत्रेयी भन्छिन्- यो एउटै वस्तुब्रह्ममा नै विरुद्धधर्मवत्ताको वर्णन गरेर हे श्रीमान्‌ ! हजुरले ममा झन मोह उत्पन्न गराउनुभयो । याज्ञवल्क्यले मोह उत्पन्न गराउने ब्रह्मको विरुद्धधर्मताको वर्णन कसरी गर्नुभयो, त्यो बताइन्छ- प्रथमतः ‘त्यो विज्ञानघन नै हो’ यस्तो प्रतिज्ञा गरेर फेरि ‘देहत्याग पश्चात् कुनै संज्ञा रहँदैन’ यस्तो भनिएको छ । त्यो कसरी विज्ञानघन हो र कसरी देहत्यागपछि त्यसको कुनै संज्ञा रहँदैन ? एउटै अग्नि उष्ण र शीतल दुबै प्रकारको हुन सक्दैन, त्यसैले ममा भ्रम उत्पन्न भयो ।

याज्ञवल्क्य भन्छन्- हे मैत्रेयि ! म मोहको उपदेश गरिरहेको छैन अर्थात् मोह उत्पन्न गर्ने कुरा मैले गरेको छैन । मैत्रेयी भन्छिन्- त्यसो हो भने त्यो विज्ञानघन हो र त्यसको कुनै संज्ञा छैन भने हजुरले त्यसका दुई विरुद्ध धर्म किन बताउनुभयो ? याज्ञवल्क्य भन्छन्- मैले यी दुई धर्म एउटै धर्मीमा छन् भनेर बताएको छैन, भ्रान्तिका कारण तिमीले नै एक वस्तुलाई विरुद्ध धर्मवाला बुझ्यौ तर मैले यसो भनेको छैन । मैले त यसो भनेको थिएँ- आत्माको जुन अविद्याद्वारा प्रस्तुत गरेको देहेन्द्रियसम्बन्धी खिल्यभाव छ, त्यसलाई विद्याद्वारा नाश गरेपछि त्यस खिल्यभावका कारण रहन गएको जुन शरीरादिसम्बन्धी अन्यत्वदर्शनरूप विशेष किसिमको संज्ञा हुन्छ, त्यो कार्यकारणसंघातरूप उपाधिलाई लीन गर्नाले कुनै हेतु न रहने हुनाले यही तरिकाबाट नष्ट हुन्छ- जस प्रकार जलादि आधारकको नाश हुँदा चन्द्रादिको प्रतिबिम्ब र त्यसबाट हुने प्रकाशादिको नाश हुन्छ । तर जसरी वास्तविक चन्द्रमा र सूर्यादिको स्वरूपको नाश हुँदैन त्यसरी नै असंसारी ब्रह्मको स्वरूप विज्ञानघनको पनि नाश हुँदैन; त्यसैलाई ‘विज्ञानघन’ यस नामले भनिएको छ; त्यही सम्पूर्ण जगत्‌को आत्मा हो एवं भूत पदार्थको नाश हुँदा पनि परमार्थतः त्यसको नाश हुँदैन । विनाशी त अविद्याजनित खिल्यभाव नै हो । जस्तो कि “विकार वाणीबाट आरम्भ हुने नाम मात्र हो” यो अन्य श्रुतिको भनाइबाट पनि सिद्ध हुन्छ; किन्तु यो त पारमार्थिक हो, हे मैत्रेयि ! यो आत्मा त अविनाशी हो; अतः जसरी यसको व्याख्या गरिएको छ, त्यसरी नै यो अनन्त, अपार महद्‍भूत जान्न सकिन्छ । “विज्ञाताको विज्ञान विशेष रूपले लोप हुँदैन, किनकि यो अविनाशी हो” यस्तो कुरा श्रुतिले अगाडि भन्छ पनि)

\customparagraph{\textit{\textbf{व्यवहार द्वैतमा छ, परमार्थ व्यवहारातीत हो -}}}

कथं तर्हि प्रेत्य संज्ञा नास्ति ? इत्युच्यते, शृणु-

\customparagraph{\textbf{नेपाली}}

शरीरको त्यागपछि त्यसको संज्ञा कसरी रहँदैन ? त्यस बारेमा भन्छु, सुन-

\begin{visheshshloka}[center]

\textbf{यत्र हि द्वैतमिव भवति तदितर इतरं जिघ्रति तदितर इतरं पश्यति तदितर इतरं शृणोति तदितर इतरमभिवदति तदितर इतरं मनुते तदितर इतरं विजानाति यत्र वा अस्य सर्वमात्मैवाभूत्तत्केन कं जिघ्रेत्तत्केन कं पश्येत्तत्केन कं शृणुयात्तत्केन कमभिवदेत्तत्केन कं मन्वीत तत्केन कं विजानीयात् । येनेदं सर्वं विजानाति तं केन विजानीयाद्विज्ञातारमरे केन विजानीयादिति ॥१४॥}

\end{visheshshloka}

\customparagraph{संस्कृतव्याख्या}

यत्र अविद्यावस्थायामेकस्यैवात्मानो द्वैतवत् ज्ञानं भवति । तस्मात् कारणात्  तत्= स एक एवापि इतर इतरं जिघ्रति = गन्धमाघ्रापयति । इतर इतरं शृणोति= श्रुतिविषयं करोति । इतर इतरमभिवदति= सम्मानभावं दर्शयति । इतर इतरं मनुते= चिन्तयति । इतर इतरं विजानाति= बोधयति । परं प्रेत्यानन्तरं तु एतादृशानां नामरूपाणां सर्वेषामात्मा एव अवशिष्टश्चेत् स तु एक एव, तदा स केन= करणेन कं पश्येत्, कं शृणुयात्, कम् अभिवदेत्, कं मन्वीत, कं विजानीयात्, अर्थात् तत्र विभिन्नेन्द्रियाभावाद् दर्शनादिक्रिया न भवति । परब्रह्म येनेदं सर्वं विजनाति तं विज्ञातारं केन करणेन को वान्यो विजानीयात्, अरेे मैत्रेयि ! तं सर्वविज्ञातारं कोऽन्य केन करणेन विजानीयात् । अर्थाद् अविद्यावस्थायां भ्रमात्मकरूेण ज्ञातं नामरूपादिकं वैभिन्न्यं ब्रह्मविद्याया ज्ञानेन तान् सर्वान् विनश्य केवलम् आत्मारूपेणैवावशिष्यते । तत्र प्रेत्यभावानन्तरं  संज्ञादिकं द्वैतं नावशिष्यते इति मया विहितोपदेशस्य सारः । इति याज्ञवल्यो विषयसारं प्रकाशयति ।

\customparagraph{\textbf{नेपाली भावार्थ}}

जहाँ अविद्या अवस्थामा द्वैत-जस्तो भान हुन्छ, त्यहीँ अन्य अन्यलाई सुँघ्छ, अन्‍य अन्यलाई देख्छ, अन्य अन्यलाई सुन्छ, अन्य अन्यलाई अभिवादन गर्दछ, अन्य अन्यलाई मनन गर्छ तथा अन्य अन्यलाई जान्दछ । किन्तु जहाँ यसका लागि सबै आत्मा नै भएको छ; त्यहाँ कसद्वारा कसलाई सुघ्ने, कसद्वारा कसलाई हेर्ने, कसद्वारा कसलाई सुन्ने, कसद्वारा कसलाई अभिवादन गर्ने, कसद्वारा कसलाई मनन गर्ने र कसद्वारा कसलाई जान्ने ? जसले यी सबैलाई जान्दछ त्यसलाई कसद्वारा जान्ने ? हे मैत्रेयि ! स्वयं विज्ञातालाई कसद्वारा जान्ने ? वा जान्न सकिन्छ । त्यसैले अविद्याबाट प्राप्त जुन नामरूपधारी पदार्थको ज्ञान छ त्यो ब्रह्मविद्याको ज्ञानबाट नष्ट हुन्छ । नाश भइसकेपछि केवल परमात्मा मात्रै बाँकी रहन्छ । त्यसको नामरूपादि केही हुँदैन । हे मैत्रेयि ! ‘मैले तिमीलाई यही वास्तविकता बुझाउन खोजेको हुँ ।’ यसरी याज्ञवल्क्यले मैत्रेयीमा उत्पन्न भएको भ्रमलाई निवारण गरे ।

\begin{visheshshlokatwo}[center]{84}

\textbf{शतपथ ब्राह्मणअन्तर्गत याज्ञवल्क्यमैत्रेयीसंवादाख्यानको}

\textbf{नेपाली अनुवाद समाप्त भयो ।}

\end{visheshshlokatwo}

\begin{center}

\textbf{अभ्यासार्थं प्रश्नाः}

\end{center}

\textbf{अतिसङ्‌क्षिप्तोत्तरापेक्षिणः प्रश्नाः (१)}

१. ब्राह्मणग्रन्था विशेषतः कस्य कृते उपयोगिनः सन्ति ?

२. शुनःशेपाख्यानं कस्माद् ग्रन्थादुद्‌धृतमस्ति ?

३. कस्य कृपया हरिश्चन्द्रः पुत्रं लभते ?

४. ‘चरैवेति चरैवेति’ कः कं कथयति ?

५. शुनःशेपः कस्य पुत्रासीत् ?

६. रोहितः कं कति मूल्येन क्रीणाति ?

७. पित्रा विक्रीतं शुनःशेपं पुत्रत्वेन कः स्वीकरोति ?

८. शुनःशेपाख्याने कति खण्डाः सन्ति ?

९. शतपथब्राह्मणस्य नामकरणबीजं किम्  ?

१०. पुरूरवोर्वश्याख्यानस्योपजीव्यो ग्रन्थः कः ?

११. बृहदारण्यकोपनिषद् कस्य ग्रन्थस्यांशभूता विद्यते ?

१२. पुरूरवसा सह उर्वश्या विहितं पणत्रयं किम् ?

१३. गन्धर्वाः केन विधिना पुरूरवसः पणभङ्गं कुर्वन्ति ?

१४. उर्वश्याः कथनानुसारं स्त्रियो मनः केन सह तुल्यं भवति ?

१५. पुरूरवसः प्रार्थनां स्वीकृत्य उर्वशी आत्मना सह समागमाय कीदृशम् उपायं भणति ?

१६. केन विधिना पुरूरवा गन्धर्वाणां वासस्थानं प्राप्तवान् ?

१७. याज्ञवल्क्यमैत्रेयीसंवादः कस्मिन् ग्रन्थे समाविष्टो विद्यते ?

१८. याज्ञवल्क्यस्य पत्न्योर्नामनी के ?

१९. याज्ञवल्क्यो मैत्रेय्या कीदृशेन वचनेन हृष्टः सन् उपदिशति ?

२०. याज्ञवल्क्यमैत्रेयीसंवादस्य प्रतिपाद्यो विषयः कः ?

\textbf{सङ्‌क्षिप्तोत्तरापेेक्षिणः प्रश्नाः (६)}

१‍. शुनःशेपाख्याने प्रतिपादितं पर्यटनस्य महत्त्वं प्रस्तूयताम् ।

२. अजीगर्तस्य मध्यमपुत्रविक्रयमुचितं किम् ? स्वमतमुपस्थाप्यताम् ।

३‍. मैत्रेय्या चरित्रं वर्ण्यताम् ।

५. धनज्ञानयोः कस्य महत्त्वमधिकम् ? याज्ञवल्यमैत्रेयीसंवादानुसारं समर्थ्यताम् ।

\textbf{विवेचनात्मकाः प्रश्नाः (१०)}

१. शुनःशेपाख्यानस्याध्ययनेन का शिक्षा लभ्यते ? विशदीक्रियताम् ।

२. याज्ञवल्क्यमैत्रेयीसंवादे वर्णितं ब्रह्मविद्यायाः स्वरूपं व्याख्यायताम् ।

३. पुरूरवोर्वश्याख्यानस्य द्वयोर्नायकनायिकयोः संवादात् का शिक्षा लभ्यते ? प्रकाश्यताम् ।

\specialmahakhandatwo{\textbf{तृतीयः पाठः}}{\textbf{उपनिषदध्ययनम्}}

\customparagraph{\textbf{प्रस्थानत्रयीपरिचयः}}

प्रस्थानशब्दस्य सामान्यार्थाः प्रयाणम्, गमनम्, प्रशस्तमार्गः, पद्धतिप्रभृतयो भवन्ति । एतदाधार एव ‘प्रस्थानत्रयी’ इत्यस्य शब्दस्याशयः परमार्थज्ञानप्राप्तिरूपप्रशस्तमार्गो भवति । प्रस्थानत्रयीत्यत्र उपनिषद्-श्रुतिप्रस्थानम्, गीता-स्मृतिप्रस्थानम्, ब्रह्मसूत्रम्- न्यायप्रस्थानमिति ज्ञानार्जनरूपाणि त्रीणि प्रस्थानानि समाविशन्ति ।

\customparagraph{\textbf{उपनिषद्प्रस्थानम्}}

उपनिषद् वेदस्य चरमो भागो मन्यते । उपनिषद्‌ग्रन्थेषु ब्रह्मवेत्तारो मुनयो ब्रह्मविद्याया गभीररहस्यमयार्थानां सुस्पष्टं विवेचनं कृतवन्तः । यत्र पौरस्त्यज्ञानस्य अपारो भण्डारः सन्निहितो  विद्यते । उपनिषदो जिज्ञासुजनं ब्रह्मविद्यायामुन्मुखीकृत्य मोक्षलाभाय प्रेरयन्ति । विशेषतः परब्रह्मणो जीवात्मनो जगतश्च चिन्तने उपनिषदः समर्पिताः सन्ति । अनन्तज्ञानदृष्ट्या इमाः उपनिषदो ब्रह्मविद्याया अनन्तपाराः पारावारसदृशाः सन्ति । तस्मात् प्रस्थानत्रय्यामुपनिषदः सर्वश्रेष्ठस्थानमावहन्ति ।

\customparagraph{\textbf{श्रीमद्भगवद्गीताप्रस्थानम्}}

अपरं प्रस्थानं श्रीमद्भगवद्गीता विद्यते । यत्र उपनिषदि वर्णितानामध्यात्मतत्त्वानां सारभूतोऽर्थः सन्निहितोऽस्ति । गीयते आत्मविद्योपदेशात्मिका ब्रह्मतत्त्वोपदेशमयी कथा यत्र सा गीता (शब्दकल्पद्रुमः) । गीतेति नाम्ना ‘रामगीता’ ‘गणेशगीता’  इत्याद्या नैकाः शास्त्रसम्बद्धा गीता दृश्यन्ते । तासु गीतासु सकलशास्त्रसारभूता ज्ञानकर्मभक्तियोगत्रयपरिलसिता श्रीमद्भगवद्गीता साधारणजनानां कृते पवित्रीकरणसमर्था ज्ञानगङ्‌गा एव । सर्वतत्त्वज्ञस्य व्यासस्य ब्रह्माण्डस्थितसम्पूर्णपदार्थवर्णनपरं महाभारतं ज्ञानसागरमेव । तत्रैव वैदिकमतप्रतिष्ठापनाय आध्यात्मिकज्ञानप्रतिपादिका सर्ववेदान्तसारयुता गीता सङ्‌गृहीताऽस्ति । गीता सर्वेषु प्रस्थानेषु सर्वजनसंवेद्यतया सर्वैरध्यात्मविद्याजिज्ञासुजनैरादृता विद्यते ।

\customparagraph{\textbf{ब्रह्मसूत्रप्रस्थानम्}}

बादरायणविरचितं ब्रह्मसूत्रं तृतीयं प्रस्थानं मन्यते । तस्मिन् आपाततः परस्परविरोधवद् भासमानानामुपनिषद्वाक्यानां समन्वयप्रदर्शनपुरःसरं पर्यवसानं प्रदर्शितं विद्यते । एवमेव श्रुतिविपरीततर्कसमूहस्य निराकरणञ्च विहितं प्राप्यते । बादरायणश्चतुर्भिरध्यायैर्वेदान्तशास्त्रस्य सूत्राणि रचयामास । ते चाध्यायाः समन्वय-अवरोध-साधन-फलाख्या विद्यन्ते । तत्र अथातो ब्रह्मजिज्ञासा इत्यारभ्य वैदिकज्ञानकाण्डविचारः प्रकाशितः ।

संस्कृतवाङ्‌मये प्रस्थानत्रयी एव वैदिकधर्मदर्शनाधारभूता मन्यते तथापि प्राधान्येन उपनिषदेव वेदानामन्तिमभागत्वात्  श्रीमद्भगवद्गीताब्रह्मसूत्रयोरवलम्बनभूता स्वीक्रियते ।

\customparagraph{\textbf{उपनिषदो विशिष्टपरिचयः}}

‘उपनिषद्यते प्राप्यते, ज्ञायते ब्रह्मविद्या अनया’ उपनिषदिति उपनिषदः सामान्योऽर्थः। ‘उपनिषद्’ इति शब्दे पदत्रयं विद्यते- उप+नि+षद् । तत्र उप, नि इति उपसर्गद्वयपूर्वात् ‘षद्’ धातोः उपनिषच्छब्दस्य निष्पत्तिर्भवति । ‘उप’ इत्यस्य समीपमित्यर्थो भवति । तस्मादस्य शब्दस्य व्युत्पत्तिलभ्योऽर्थो भवति यत्- ‘गुरोः समीपमुपविश्य यत्र ब्रह्मविद्याया ज्ञानं प्राप्यते, सा उपनिषद्’ इति ।

उपनिषदो वेदस्यांशभूता एव । सर्वा अपि उपनिषदश्चतुर्षु वेदेषु -अन्यतमे वेदे सम्बद्धाः । उपनिषदि प्रधानप्रतिपाद्यविषयो ब्रह्मविद्या एव । यद्यपि उपनिषदः सङ्‌ख्याः शताधिका अस्ति । तथापि पौरस्त्यसाहित्ये प्रामुख्येन दश उपनिषदः प्रसिद्धा मन्यन्ते । तदुक्तं यत् -

\begin{shloka}[center]

ईश-केन-कठ-प्रश्न-मुण्ड-माण्डुक्य-तित्तिरीः ।

ऐतेरियं च छान्दोग्यं बृहदारण्यकं दश ॥ इति ।

\end{shloka}

‘एता दश एव उपनिषदः’ इति न मन्तव्यम् । श्वेताश्वतर-जाबाल- कौशितक्याद्युपनिषदोऽपि महत्त्वाधायिकाः सन्ति । तेष्वपि अस्माकं पाठ्यक्रमे ईशावास्यकेनोपनिषदौ स्थपितौ विद्येते । अतः तयोरेवात्र परिचयपुरस्सरं पाठश्च प्रस्तूयते ।

\specialkhanda{शुक्लयजुर्वेदस्था ईशावास्योपनिषद्}

\customparagraph{\textbf{सारः/वैशिष्ट्यञ्च}}

ईशावास्योपनिषद् शुक्लयजुर्वेदान्तर्गता मन्त्रोपनिषद् । शुक्लयजुर्वेदसंहिताया अन्तिमाध्याय एव ईशोपनिषद् अस्ति । माध्यन्दिनसंहितान्तर्गतायामीशोपनिषदि सप्तदश काण्वसंहितायान्तर्गतामीशोपनिषदि चाष्टदश मन्त्राः सन्ति । अस्या उपनिषदो नाम ‘ईशावास्यमिदं सर्वम् …..’ इति प्रथममन्त्रस्यांशमुद्धृत्य निर्मितमस्ति । अस्या उपनिषदो नाम ‘ईशावास्योपनिषद्’ वा ‘ईशोपनिषदि’ति प्रसिद्धमस्ति ।

अस्या उपनिषदो मुख्योपदेशो विद्यते यत्-  चराचरजगतः कश्चिन्नियामकोऽस्ति। तेन नियामकेन इदं स्थावरजङ्गमात्मकं सर्वं जगद् व्याप्तमस्तीति । ईशावास्यमिदं सर्वमित्यस्य इदं सर्वं चराचरमयं यत्किञ्चित् दृश्यते, तत्सर्वं ईशा= ईशेन नियामकेन परब्रह्मणा वास्यम्= व्याप्तमित्यर्थः। अत्र जनाः सर्वम् ईशावास्यमिति मनसि निधाय परोपकारमयं कर्म कुर्वन्ति तदा तेषां कृते कर्मणो बन्धनं न भवति इत्यपि तत्रोपदिष्टम्। यथा-

\begin{shloka}[center]

कुर्वन्नेवेह कर्माणि जिजीविषेच्छतं समाः ।

एवं त्वयि नाऽन्यथेतोऽस्ति न कर्म लिप्यते नरे ॥

\end{shloka}

यो नरः केवलं स्वार्थपूर्णं कार्यं करोति सोऽत्र किञ्चित् सुखमनुभवति, परं शरीरपातानन्तरमन्धकारवृतं लोकं प्राप्नुवन्ति। अतः परोपकारः करणीय इति तत्रोपदिष्टं विद्यते यत्-

असुर्या नाम ते लोका अन्धेन तमसा वृत्ताः ।

तान् ते प्रेत्याभिगच्छन्ति ये के चात्महनो जनाः ॥

अत्र विद्याविद्याया विषये सम्यक् चर्चा विद्यते । ये जना भौतिकेच्छापूर्तिरूपामविद्यामुपासते ते अन्धं तमः प्रविशन्ति । आध्यात्मिकज्ञानरूपाया विद्यायाः सेवनेनैव अमृतमाश्वाद्यते । बुद्धिमन्तो जना अविद्यामेव नापितु विद्यामपि साधयित्वा अक्षय्यं फलं प्राप्नुवन्ति । तस्मादत्र ज्ञानकर्मणोः समुच्चयः उपदिष्टो विद्यते, यथा-

विद्यां चाविद्यां यस्तद् वेदोभयं सह ।

अविद्यया मृत्युं तीर्त्वा विद्ययाऽमृतमश्नुते ॥

एवं प्रकारेणास्यामुपनिषदि चराचरजगति व्यापकस्य नियामकस्य ईशस्य वर्णनम्, अद्वैतसिद्धान्तस्य बीजम्, ब्रह्मविद्याया महत्त्वम्, ज्ञानकर्मणो समुच्चयः, दिव्याः प्रार्थना इत्यादिकं समुपदिष्टं विद्यते । तस्माद् वैदिकसाहित्ये अस्याः उपनिषदो महत्त्वमन्यूनं ज्ञायते ।

\begin{center}

\textbf{शान्तिपाठः}

\end{center}

\begin{visheshshloka}[center]

\textbf{पूर्णमदः पूर्णमिदं पूर्णात्पूर्णमुदच्यते ।}

\textbf{पूर्णस्य पूर्णमादाय पूर्णमवेवावशिष्यते॥}

ॐ शान्तिः शान्तिः शान्तिः !

\end{visheshshloka}

\customparagraph{संस्कृतव्याख्या}

अदः= इदं परब्रह्म पूर्णमस्ति, एवमिदं कार्यब्रह्मरूपं जीवमयं जगदपि  पूर्णमेवास्ति। (प्रलयकाले) पूर्णात्= पूर्णब्रह्मणः पूर्णमेव= पूर्णरूपेणैव उदच्यते= जगदुत्पत्तिर्भवति। पूर्णस्य= पूर्णब्रह्मणः पूर्णम्= पूर्णत्वम्, सर्वव्यापकत्वम्, आदाय= गृहीत्वा पूर्णमेव= पूर्णपरब्रह्मस्वरूपयुतोऽवशिष्यते= अवशिष्टो भवति, पूर्णब्रह्मरूपो भवति,  जीव इति शेषः।

\customparagraph{\textbf{नेपाली भावार्थ}}

ब्रह्मतत्त्व सर्वत्र परिपूर्ण एवं सर्वव्यापक छ। जीव पनि कार्यब्रह्मस्वरूपकै भएकाले ऊ पनि स्वयंंमा पूर्ण छ । सम्पूर्ण जगत् पनि पूर्ण ब्रह्म नै व्यापक रूपमा रहेकाले त्यो पनि पूर्ण छ । अज्ञानका कारण अपूर्ण महसुस गरेको जीवले ब्रह्मतत्त्वको पूर्णता र सर्वव्यापकतालाई बुझ्न सक्यो भने ऊ पनि पूर्ण परब्रह्म स्वरूपकै हुन सक्छ ।

\begin{visheshshloka}[center]

\textbf{ईशा वास्यमिदं सर्वं यत्किञ्च जगत्यां जगत् ।}

\textbf{तेन त्यक्तेन भुञ्जीथा मा गृधः कस्यस्विद्‌धनम्॥}

॥ १॥

\end{visheshshloka}

\customparagraph{संस्कृतव्याख्या}

परमेश्वरः संसारस्य सर्वेषु वस्तुषु व्याप्तोऽस्ति । ते सर्वे पदार्थाः परमात्मनः कारणेनैव अस्तित्वयुताः सन्ति । अतो यस्य वस्तुनो जीवनयापनार्थम् उपयोगः क्रियते तस्य स्वकीयावश्यकतानुरूपं त्यागपूर्वकमेव उपयोगः कर्तव्यः । त्यागपूर्वकमिति-    भोगसमये वस्तुषु लिप्सा न भवेत् ।  वस्तुभोगः केवलं जीवनसञ्चालनाय एव भवति । ईश्वरः सर्वत्र सर्वेषु वस्तुष्वस्ति, अतः किमर्थं लिप्सा कर्तव्या । तस्मात् त्यागपूर्वकमेव भोगः करणीयः । भोगसमये कस्यापि अन्यस्याधिकारगतं धनं न ग्रहणीयम् । अन्यस्य धनमित्युक्ते स्वपरिश्रमेण यत् धनं प्राप्यते तस्य धनस्य एव उपयोगः करणीयः ।

\customparagraph{\textbf{नेपाली भावार्थ}}

यो संसारमा जे-जति पदार्थहरू छन्, ती सबैमा परमात्मा व्याप्त छन् । अर्थात् ती सबै आत्मतत्त्वकै चैतन्यलाई प्राप्त गरेर अस्तित्वमा रहेका छन् । त्यसैले यो सम्पूर्ण चराचर जगत् नै परमात्माको धन हो। व्यक्तिले जीवनयापनका लागि उपयोग गर्दा स्वार्जित धनको मात्र ग्रहण गर्नुपर्दछ । अरूको हिस्साको धनमा लोभ राख्नु हुँदैन । आफूले आर्जन गरेकै धन पनि अधिक छ भने परोपकारको भावना राखी अन्यमा वितरण गर्नुपर्छ । उपभोग गर्दा लोभको नभई त्यागको भावना राख्नुपर्छ ।

\begin{visheshshloka}[center]

\textbf{कुर्वन्नेवेह कर्माणि जिजीविषेच्छतं समाः ।}

\textbf{एवं त्वयि नान्यथेतोऽस्ति न कर्म लिप्यते नरे ॥}

॥ २॥

\end{visheshshloka}

\customparagraph{संस्कृतव्याख्या}

अस्मिन् जगति शास्त्रविहितकर्माणि कुर्वन् सर्वनियन्त्रकं परमेश्वरं स्मरन् शतं वर्षं यावत् जीवितुम् इच्छां कुर्यात् । शास्त्रेषु शतं वर्षं यावत् मनुष्यस्य जीवनं पूर्णं मन्यते । अतः शतं यावत् जीवितुम् अभिलषेत् । मनुष्येभ्य एतदतिरिक्तोऽन्यः कश्चिद्‌ हितकरो मार्गो नास्ति । निष्कामभावनया शास्त्रविहितं कर्म कुर्वन् नरः कदापि कर्मणा लिप्तो न भवति । पुनर्जन्मतो मुक्तिलाभार्थम् एतदतिरिक्तो नास्ति अन्यो मार्गः ।

\customparagraph{\textbf{नेपाली भावार्थ}}

यो संसारमा शास्त्रनिर्दिष्ट कर्महरू गर्दै निष्काम भावले सर्वसञ्चालक परमेश्वरलाई सम्झँदै सय वर्ष जिउने इच्छा गर्नुपर्दछ । शास्त्रअनुसार मानवका लागि सय वर्ष पूर्ण आयु मानिन्छ । त्यसैले मानवले सय वर्ष बाँच्ने इच्छा राख्नुपर्दछ । मानवका लागि योभन्दा अर्को कल्याणपरक मार्ग छैन । यो शास्त्रविहित कर्म गर्ने मानिस कदापि कर्ममा लिप्त हुँदैन । अतः त्यागपूर्वक सांसारिक वस्तुको उपभोग वा निष्काम कर्मभन्दा पुनर्जन्मको दुःखबाट मुक्ति पाउने अन्य मार्ग छैन ।

\begin{visheshshloka}[center]

\textbf{असुर्या नाम ते लोका अन्धेन तमसाऽऽवृताः ।}

\textbf{ताँस्ते प्रेत्याभिगच्छन्ति ये के चात्महनो जनाः ॥}

॥ ३॥

\end{visheshshloka}

\customparagraph{संस्कृतव्याख्या}

असूर्या नामकाः केचन लोकाः सन्ति । यत्र विभिन्नाधमयोनिगता जीवा निवसन्ति । तेऽसुरलोका अन्धेन= अतिशयेन तमसा= अज्ञानावरणेन, दुःखक्लेशेन आवृताः= आच्छादिताः सन्ति । आत्महनो जनाः‍= येऽज्ञानमग्नाः सन्तोऽऽत्मानं घ्नन्ति, ते मनुष्याः प्रेत्य= मृत्युं प्राप्य तान्= असुर्याख्यान् लोकान् अभिगच्छन्ति, वारं वारं तानेव लोकान् प्राप्नुवन्ति । अतोऽज्ञानं स्वीकृत्यात्महननं न कर्तव्यम्, अन्यथा तेषु लोकेषु तुच्छयोनिरेव प्राप्यते ।

\customparagraph{\textbf{नेपाली भावार्थ}}

कतिपय असुर्या नामका लोक छन् जहाँ ज्ञानहीन निम्न कोटिका प्राणीहरूको बसोबास छ । ती लोक अज्ञानरूपी आवरणले ढाकिएका छन् । जोजो मानवले ज्ञानबाट विमुख भएर आत्मघात गर्छन्, तिनीहरू मृत्युपछि पटकपटक तुच्छ योनिमा जन्म लिएर तिनै लोकमा दुःख भोग गर्न पुग्छन् । त्यसकारण सचेत मानिसले ज्ञानको मार्गलाई अनुसरण गर्नुपर्छ ।

\begin{visheshshloka}[center]

\textbf{अनेजदेकं मनसो जवीयो}

\textbf{नैनद्देवा आप्नुवन् पूर्वमर्षत् ।}

\textbf{तद्धावतोऽन्यानत्येति तिष्ठत्}

\textbf{तस्मिन्नपो मातरिश्वा दधाति ॥}

॥ ४॥

\end{visheshshloka}

\customparagraph{संस्कृतव्याख्या}

अनेजत्= अचलम्, एकम्= अद्वितीयम्, मनसो जवीयः= सर्वान्नतिरिच्य तीव्रगतियुक्तम्, पूर्वम्= प्राचीनकालादेव अर्षत्= सर्वत्र व्याप्तम्, एनत्= ब्रह्मतत्त्वम्, देवाः‍= इन्द्रादयोऽपि न प्राप्नुवन् । अर्थात् ब्रह्मतत्त्वं केऽपि कदापि न प्राप्तवन्तः । तत् ब्रह्म तिष्ठत्= अचलदेव धावतः= धवमानान्, अन्यान्= वायुप्रभृतीन् देवान् अत्येति‍= अतिक्रामति । तस्मिन् भूते सति= तस्य ब्रह्मणः कारणेनैव मातरिश्वा= वायुसदृशो देवः, अपः‍= वर्षारूपेण जलानि दधाति । जलदेवोऽपि तस्यैव साहाय्येन जलवृष्टिं कर्तुं समर्थो भवति । अतस्तद् ब्रह्म सर्वशक्तिमदस्ति ।

\customparagraph{\textbf{नेपाली भावार्थ}}

अचल, अद्वितीय, मनभन्दा पनि तीव्र वेगशाली, प्राचीन कालदेखि नै सर्वत्र व्याप्त, त्यस ब्रह्मतत्त्वलाई इन्द्रादि अथवा इन्द्रियरूपी देवताले पनि प्राप्त गर्न सकेनन्। निरन्तर तीव्र वेगमा दौडिने वायु आदि गतिशील तत्त्वलाई पनि जो एकै ठाउँमा बसीबसी जित्छ, एवम् त्यही ब्रह्मतत्त्वकै सहायताले जलदेवता पनि जलवर्षा गराउँन सक्षम छन् ।

\begin{visheshshloka}[center]

\textbf{तदेजति तन्नैजति तद्दूरे तद्वन्तिके ।}

\textbf{तदन्तरस्य सर्वस्य तदु सर्वस्यास्य बाह्यतः॥}

॥ ५॥

\end{visheshshloka}

\customparagraph{संस्कृतव्याख्या}

तत्= ब्रह्म एजति= चलति, गतिशीलमस्ति । तत्= ब्रह्म न एजति= न चलति, गतिहीनमपि वर्तते । तद् ब्रह्म दूरेऽस्ति, अस्मादत्यन्तं दूरमस्ति । तत् ब्रह्म अन्तिके च= हृदयेऽप्यवस्थितत्त्वादत्यन्तं निकटेऽपि विद्यते । तद् ब्रह्म सर्वस्य जगतोऽन्तःस्थितमस्ति, अर्थात् सर्वेषु वस्तुषु पदार्थेषु च निहितमस्ति । तथा च तत् ब्रह्म सर्वतः सर्वस्य पदार्थस्य बाह्यतोऽपि स्थितमस्ति । तात्पर्यमिदं यत्-  ब्रह्म सर्वेषां पदार्थानां मध्येऽपि अस्ति । पदार्थाद् बहिरप्यस्ति ।

\customparagraph{\textbf{नेपाली भावार्थ}}

त्यो ब्रह्मतत्त्व अचल भए पनि प्राणीहरूका रूपमा गतिशील पनि छ । इन्द्रियबाट अगम्य भएकाले टाढा र हृदयमा अवस्थित भएकाले अत्यन्त नजिक पनि छ। त्यो ब्रह्मतत्त्व सम्पूर्ण चराचरको सञ्चालकका रूपमा भित्र छ भने यो ब्रह्माण्डको सीमा बाहिर पनि व्याप्त छ। अर्थाद् ब्रह्म तत्त्व भित्र, बाहिर सर्वत्र व्याप्त छ ।

\begin{visheshshloka}[center]

\textbf{यस्तु सर्वाणि भूतान्यात्मन्येवाऽनुपश्यति ।}

\textbf{सर्वभूतेषु चात्मानं ततो न विजुगुप्सते ॥}

॥ ६॥

\end{visheshshloka}

\customparagraph{संस्कृतव्याख्या}

यो पुरुषः (मनुष्यः) परमेश्वरविषये- ‘ईश्वरः सर्वेषु वस्तुष्वस्ति, सर्वाणि वस्तूनि सर्वे प्राणिनश्च आत्मनि ईश्वरे एव सन्ति’ एवं चिन्तयति, तथा च ‘सर्वभूतेषु आत्मानमीश्वरं पश्यति’ स केनापि सार्द्धं कदापि न विजुगुप्सते । अर्थाद् यो मनुष्यः ‘ईश्वरः सर्वत्र सर्वेषु पदार्थेषु च वर्तते, सर्वे पदार्थाः, मनुष्याः, जीवाश्च ईश्वरे सन्ति’ इति चिन्तयति, स महापुरुषः सर्वत्र अहमेवेति विचार्य कुत्रापि केनापि सार्द्धं कदापि द्रोहम्= ईर्ष्यादिकं न करोति ।

\customparagraph{\textbf{नेपाली भावार्थ}}

जसले कुनै पनि पदार्थ, वा चराचर जगत्‌मा ईश्वर निहित छ भन्ने अनुभूति गर्दछ, त्यस मानवले सर्वत्र आफ्नै स्वरूप देख्दछ, आफूभन्दा भिन्न कोही नभएको महसुस गर्नाले उसले कसैसँग घृणा वा द्वेष आदि नकारात्मक भाव राख्दैन ।

\begin{visheshshloka}[center]

\textbf{यस्मिन् सर्वाणि भूतान्यात्मैवाऽभूत् विजानतः ।}

\textbf{तत्र को मोहः कः शोक एकत्वमनुपश्यतः ॥}

॥ ७॥

\end{visheshshloka}

\customparagraph{संस्कृतव्याख्या}

यस्मिन् परिस्थितौ काले वा मनुष्यः, ‘ईश्वरस्य सम्पूर्णस्वरूपं विजानतः; सर्वाणि भूतानि आत्मा एव अभूत्’ इति चिन्तयति,  अर्थात् ईश्वरः सर्वत्र सर्वेषु पदार्थेष्वस्तीत्यनुभवति, सर्वत्र एकत्वं पश्यति, एवम् ईश्वरीयपदार्थस्य च एकत्वं पश्यन् यो जीवति, तस्य पार्श्वे को मोहः, कः शोकः; अर्थात् ईश्वरपदार्थयोरेकत्वं पश्यन् मनुष्यः शोकरहितो मोहरहितश्च भवति ।

\customparagraph{\textbf{नेपाली भावार्थ}}

जुन समयमा मानिस ईश्वरका सम्पूर्ण रूप जान्दछ; सबै नामरूपधारी पदार्थमा आत्मरूप ईश्वर विद्यमान छ भन्ने चिन्तन गर्दछ;  सर्वत्र ईश्वरकै अनुभूति गर्दै जीवनयापन गर्दछ; त्यस समयमा मानिसलाई मोह, शोकजस्ता विकारहरूले प्रभावित पार्न सक्दैनन् ।

\begin{visheshshloka}[center]

\textbf{स पर्यगाच्छुक्रमकायमव्रण-}

\textbf{मस्नाविरं शुद्धमपापविरुद्धम् ।}

\textbf{कविर्मनीषी परिभूः स्वयम्भूर्याथातथ्यतो-}

\textbf{ऽर्थान् व्यदधाच्छाश्वतीभ्यः समाभ्यः ॥}

॥ ८॥

\end{visheshshloka}

\customparagraph{संस्कृतव्याख्या}

शुक्रम्= परमतेजोमयम्, अकायम्= शरीररहितम्, अव्रणम्= अक्षतम्, अस्नाविरम्= स्थूलपाञ्चभौतिकपदार्थरहितम्, शुद्धम्= दिव्यसच्चिदानन्दस्वरूपम्, अपापविद्धम्= शुभाशुभकर्मसम्पर्कशून्यं परमात्मतत्त्वम्;  कविः= सर्वद्रष्टा, मनीषी= सर्वज्ञः सर्वनियन्ता च स्वयम्भूः= स्वेच्छया उत्पन्नः सः= आत्मा, पर्यगात्= सर्वत्र व्याप्तो जातः। स एव शाश्वतिभ्यः समाभ्यः= अनादिकालात्, याथातथ्यतः= यथार्थरूपेण यथायोग्यम्, अर्थान्= कर्तव्यपदार्थान् व्यदधात्= चक्रे, यथानुरूपं व्यभजदित्यर्थः । अर्थात् आदिकालादेव परमेश्वरः सर्वान् कर्तव्याकर्तव्यपदार्थान् संसृज्य सर्वेषु प्राणीषु वितरतीत्यर्थः।

\customparagraph{\textbf{नेपाली भावार्थ}}

परमतेजोमय, निराकार, नसा, हाड, छाला आदिको बन्धनयुक्त पाञ्चभौतिक शरीरले रहित, शुद्ध, कर्मका विकारबाट अस्पृश्य परमात्मतत्त्व छ । सर्वद्रष्टा सर्वज्ञ स्वेच्छाले उत्पन्न त्यो आत्मतत्त्व सर्वत्र व्याप्त भयो । त्यसैले अनादिकालदेखि यथार्थरूपले कर्तव्याकर्तव्यपदार्थको सिर्जना गरी सबैमा विभाजन गर्दछ ।

\begin{visheshshloka}[center]

\textbf{अन्धन्तमः प्रविशन्ति येऽविद्यामुपासते ।}

\textbf{ततो भूय इव ते तमो य उ विद्यायां रताः ॥}

॥ ९॥

\end{visheshshloka}

\customparagraph{संस्कृतव्याख्या}

ये= जनाः, अविद्यामुपासते= केवलम् अज्ञानादिभिरुत्पन्नस्य कर्मण उपासनां कुर्वन्ति; ते‍= जनाः अन्धतमः= अन्धकारावृतं स्थानम्, प्रविशन्ति= गच्छन्ति । तथा च ये विद्यायां रताः= देवाराधनारूपस्य ज्ञानमार्गस्यानुसरणं कुर्वन्ति; ते ततो भूय इव= तस्मादप्यधिकम् अन्धाकारावृतं स्थानं गच्छन्ति । तात्पर्यमिदं यत्- कर्ममार्गस्य उपासको जनोऽज्ञानादन्धकारे पतति, ज्ञानमार्गस्य उपासकश्च ततोऽप्यधिकमन्धकारं प्राप्नोति । अत एकस्यैव मार्गस्योपासना न कर्तव्या । उभयोरविद्याविद्ययोर्मार्गयोरुपासना कर्तव्या इति।

\customparagraph{\textbf{नेपाली भावार्थ}}

कुनै कामना राखेर गरिने सकाम कर्महरू सबै अज्ञानमय छन् । किनकि सकाम कर्म गर्ने मानिसहरू बाह्यशरीरादिलाई आत्मा ठानेर त्यसको सुखका लागि मरिमेट्‌छन् । त्यस्तो कर्म अर्थाद् अविद्याको उपासना गर्ने व्यक्तिहरू घोर अन्धकारमा निमग्न हुन्छन् । त्यस्तै विद्या अर्थाद् देवाराधनादिमा मात्र लागिरहने उपासकहरू पनि ‘म नै ईश्वर हुँ’ भन्ने भावले ग्रस्त रहन्छन्, त्यसैले तिनीहरू त झन् महाभयङ्‌कर अन्धकारको भुमरीमा पर्दछन् । अतः विद्या (देवोपासना) र अविद्या (सकाम कर्म) दुबैको अनुसरण गर्नुपर्दछ ।

\begin{visheshshloka}[center]

\textbf{अन्यदेवाहुर्विद्ययाऽन्यदाहुरविद्यया} ।

\textbf{इति शुश्रुम धीराणां ये  नस्तद्विचचक्षिरे ॥}

॥ १०॥

\end{visheshshloka}

\customparagraph{संस्कृतव्याख्या}

विवेकशीला विद्वांसो विद्यया अन्यदेव फलं प्राप्यते इत्याहुः= कथितवन्तः । एवमेव अविद्यया च भिन्नमेव फलं प्राप्यत इत्याहुः । इति= एवं प्रकारेण धीराणां विदुषां मुखेन वयं श्रुतवन्तः । ये धीराः सन्ति ते न अस्मान् तदुक्तं तथ्यं विचचक्षिरे= उक्तवन्तः । अस्याशयो यत्- ये ब्रह्मविषयकचिन्तकाः, ईश्वरज्ञानसम्पन्नाश्च सन्ति, ते ‘ज्ञानेन देवभावरूपेण, कर्मणा पदार्थभोगरूपेण च फलं प्राप्तुं शक्नुवन्ति ।’ इति वदन्ति ।

\customparagraph{\textbf{नेपाली भावार्थ}}

विवेकशील विद्वान्‌हरू ‘विद्याबाट भिन्नै प्रकारको र अविद्याबाट भिन्नै प्रकारको फल प्राप्त हुन्छ’ भन्दछन् । यस प्रकार गहिरो ज्ञान भएका विद्वान्‌का मुखबाट हामीले सुनेका हौँ । विद्वान्‌हरूले हामीलाई यो तथ्य बताएका हुन् । अर्थात्- ब्रह्मज्ञ विद्वान्‌हरू ‘विद्या या ज्ञानले देवभाव रूपमा र अविद्या वा कर्मले पदार्थ भोग रूपमा फल पाउन सकिन्छ’ भन्दछन् ।

\begin{visheshshloka}[center]

\textbf{विद्याञ्चाविद्याञ्च यस्तद्वेदोभयं सह ।}

\textbf{अविद्यया मृत्युं तीर्त्वा विद्ययाऽमृतमश्नुते ॥}

॥ ११॥

\end{visheshshloka}

\customparagraph{संस्कृतव्याख्या}

यो मनुष्यो विद्यां जानाति= देवोपासनां करोति तथा  अविद्यामपि जानाति= कर्म अपि करोति । एवं प्रकारेण यो जनो विद्याऽविद्यारूपमुभयविधं ज्ञानं धारयति, स उभयविधज्ञानेन युक्तो भूत्वा अविद्यायाः कारणेन मृत्युं तरति‍= मृत्युसंसारात् पारं याति। तदनन्तरं विद्यया अमरत्वं प्राप्नोति ।

\customparagraph{\textbf{नेपाली भावार्थ}}

जुन मानिसले अविद्या= कर्म र विद्या= देवोपासना दुबैलाई साथसाथै अङ्गीकार गर्छ, त्यस व्यक्तिले अविद्यारूप कर्मद्वारा दुःख, दरिद्रताजस्ता मृत्युतुल्य अवस्थाबाट पार पाउँछ भने विद्यारूप देवोपासनाबाट दिव्यतारूप अमृतत्त्वलाई प्राप्त गर्दछ ।

\begin{visheshshloka}[center]

\textbf{अन्धन्तमः प्रविशन्ति येऽसम्भूतिमुपासते ।}

\textbf{ततो भूय इव तमो य उ सम्भूत्यां रताः ॥}

॥ १२॥

\end{visheshshloka}

\customparagraph{संस्कृतव्याख्या}

ये मनुष्याः, असम्भूतिम्= अव्यक्तप्रकृतितत्त्वम् उपासते= उपासनां कुर्वन्ति । ते अन्धन्तमः= अन्धकारयुक्तामवस्थां प्रविशन्ति= गच्छन्ति । ये च केवलं सम्भूत्याम्= कार्यब्रह्मणि हिरण्यगर्भाख्ये रता भवन्ति, ते जनास्तस्माज्जातेन सिद्ध्यादिना साहङ्काराः सन्तस्ततो भूयः= तस्मादधिकं तमसावृत्तमवस्थां गच्छन्ति= प्राप्नुवन्ति। आशयो यत्- ये केवलमव्यक्तप्रकृतिमुपासते, ते प्रकाशरहितमज्ञानमयं लोकं गच्छन्ति, तथा च ये मनुष्याः केवलं कार्यब्रह्म भजन्ते, ते ततो विलक्षणसामर्थ्यं प्राप्य अहङ्कारनिमग्नाः सन्तस्ततोऽपि अधिकं ज्ञानहीनमवस्थां प्राप्नुवन्ति ।

\customparagraph{\textbf{नेपाली भावार्थ}}

जो मानिस असम्भूति अर्थात् अव्यक्त प्रकृतिको उपासना गर्दछ त्यसले बाहिरी कुराको सन्तुष्टि मात्र चाहने हुँदा त्यो अन्धकारको अवस्थामा पुग्छ । त्यस्तै जसले कार्यब्रह्मको मात्र उपासना गर्छ, त्यसमा विभिन्न सामर्थ्य आदिको प्राप्ति भएर त्यो ‘म सर्वेसर्वा हुँ’ भन्ने अहङ्कारले ग्रस्त हुन्छ र झन् ठुलो अन्धकारको अवस्थामा प्रवेश गर्दछ ।

\begin{visheshshloka}[center]

\textbf{अन्यदेवाहुः सम्भवादन्यदाहुरसम्भवात् ।}

\textbf{इति शुश्रुम धीराणां येनस्तद्विचचक्षिरे ॥}

॥ १३॥

\end{visheshshloka}

\customparagraph{संस्कृतव्याख्या}

सम्भवात्= कार्यब्रह्मोपासनात् अन्यदेव= भिन्नमेव (अणिमाद्यैश्वर्यरूपम्) फलमाहुः= कथयन्ति । अर्थात् सम्भवस्य कार्यब्रह्मण उपासनायाः फलम् अन्यदस्ति । एवम् असम्भवात्= अव्याकृतोपासनादन्यदेव (प्रकृतिलयरूपम्) फलं प्राप्यते, इत्याहुः = कथयन्ति । इति= एतादृशं वचनम्, वयं धीराणां मुखात् शुश्रुमः= शृणुमः, ये ब्रह्मविषयकचिन्तकाः ईश्वरज्ञानसम्पन्ना विद्वांसस्ते नोऽस्मान् तत्= तस्मिन् विषये विचचक्षिरे= व्याकृताऽव्याकृतोपासनाफलं व्याख्यातवन्तः ।

\customparagraph{\textbf{नेपाली भावार्थ}}

सम्भव अर्थात् कार्यब्रह्मतत्त्वको उपासनाबाट भिन्नै र असम्भव अर्थात् अव्यक्त प्रकृतिको उपासनाबाट भिन्नै किसिमको फल प्राप्त हुन्छ । यसरी विद्वान्‌हरूले दुई खाले उपासनाको रहस्य बताएको हामीले सुनेका छौं ।

\begin{visheshshloka}[center]

\textbf{सम्भूतिञ्च विनाशञ्च यस्तद्वेदोभयं सह ।}

\textbf{विनाशेन मृत्युं तीर्त्वा सम्भूत्याऽमृतमश्नुते ॥}

॥ १४॥

\end{visheshshloka}

\customparagraph{संस्कृतव्याख्या}

यः मनुष्यः सम्भूतिम्= विनाशरहितमव्यक्तप्रकृतिं जानाति । तस्य उपासनां करोति, तथा च यो विनाशशीलं देवपितृमनुष्यरूपमसम्भूतिरूपं कार्यब्रह्म जानाति अर्थात् विनाशशीलस्य उपासनां करोति  । अर्थात् सम्भूतिं विनाशञ्च उभयविधं स्वरूपं जानाति, स उभयविधज्ञानेन युक्तो भूत्वा कार्यब्रह्मण उपासनया मृत्यं तरति, पुनर्जन्ममृत्युं न प्राप्नोति, तदनन्तरमव्यक्तप्रकृतेरुपासनया अमरत्वं प्राप्नोति । (शाङ्करभाष्यानुसारं सम्भूतिमित्यत्र अकारलोपः)

\customparagraph{\textbf{नेपाली भावार्थ}}

जुन मानिसले विनाशशील कार्यब्रह्म हिरण्यगर्भको र अविनाशी अव्यक्त प्रकृतिको एकसाथ उपासना गर्दछ, त्यसले विनाशशील कार्यब्रह्मको उपासनाबाट सिद्धि सामर्थ्य आदि प्राप्त गरी दुःखादि मृत्युतुल्य अवस्थाबाट पार पाउँछ भने प्रकृतिको उपासनाले अमृतलाई प्राप्त गर्दछ ।

\begin{visheshshloka}[center]

\textbf{हिरण्मयेन पात्रेण सत्यस्यापिहितं मुखम् ।}

\textbf{तत्त्वं पूषन्नपावृणु सत्यधर्माय दृष्टये ॥}

॥ १५॥

\end{visheshshloka}

\customparagraph{संस्कृतव्याख्या}

हे पूषन् ! सूर्यमण्डलाभिमानी ईश्वरः, हिरण्मयेन= ज्योतिर्मयेन पात्रेण= अपिधानभूतेन मायया सत्यस्य= आदित्यमण्डलस्थस्य ब्रह्मणो मुखम्, अपिहितम्= आच्छादितं विद्यते । अतो  सत्यधर्माय= सत्यधर्मणि स्थिताय मह्यं दृष्टये= दर्शनाय त्वम्=ईश्वरः, तत्= आवरणम्, अपावृणु= उद्‌घाटय । येन अहं परमेश्वरस्य परब्रह्मणः सम्यक् स्वरूपं द्रष्टुं शक्ष्ये ।

\customparagraph{\textbf{नेपाली भावार्थ}}

हे ईश्वरः ! ज्योतिर्मय मायारूप पात्रले सत्यस्वरूप तपाईंको वास्तविक मुख छोपेको छ । अर्थात् संसारमा सच्चिदानन्दस्वरूप आत्मा नभई त्यसको अनन्त ऐश्वर्यमात्र देखिएको छ । म सत्यलाई हेर्न चाहने उपासक हुँ । तसर्थ मलाई आफ्नो वास्तविक स्वरूपको दर्शनका लागि त्यो मायारूप आवरणलाई हटाइदिनुहोस् ।

\begin{visheshshloka}[center]

\textbf{पूषन्नेकर्षे यम सूर्य प्राजा-}

\textbf{पत्य व्यूह रश्मीन् समूह ।}

\textbf{तेजो यत्ते रूपं कल्याणतमं तत्ते पश्यामि}

\textbf{योऽसावसौ पुरुषः सोऽहमस्मि॥}

॥ १६॥

\end{visheshshloka}

\customparagraph{संस्कृतव्याख्या}

हे पूषन् ! सर्वेषां पालनकर्तृदेव !, हे एकर्षे‍ ! = अद्वितीयः सन् अन्तरिक्षे विचरणशील !, यम ! = सर्वनियन्त्रक !, प्राजापत्य ! = प्रजापतिपुत्र !, सूर्य ! = दिनकर !, ते= तव रश्मीन्= किरणसमूहान् व्यूह= विगमय, समूह=एकीकुरु, सर्वप्रकाशं स्वरूपे संस्थापयतु । यत्- ते= तव कल्याणतमम्= परमकल्याणमयम्, रूपम्= वास्तविकं स्वरूपमस्ति, तत्= रूपम्, अहं पश्यामि। यः=एतादृशः, असौ=दूरे स्थितः ‘असौ’ इत्यनेन सङ्केतितो दिव्यपुरुषः प्रकाशमण्डले स्थितेन येन सर्वत्र प्रकाशः प्रसारितः, सोऽहमस्मि । अहमपि परमपुरुष एवास्मि । अहं ब्रह्मास्मि इति भावः ।

\customparagraph{\textbf{नेपाली भावार्थ}}

सबै प्राणीलाई परिपोषित गर्ने, एक्‌लै अन्तरिक्षमा विचरण गर्ने, सबैका नियन्ता, प्रजापतिका छोरा हे सूर्य ! म तिम्रो प्रकाशमय बाहिरी आवरण होइन, भित्री वास्तविक रूपलाई हेर्न चाहन्छु। त्यसकारण तिमी आफ्ना किरणहरूलाई आफूबाट अलग गराऊ र तिनलाई एकत्रित गर । सारा मानिसहरू सूर्यमण्डलाभिमानी पुरुषलाई अर्को कुनै रूप बताउँछन्, तर त्यो सूर्यमण्डलाभिमानी चेतन र यो शरीरभित्र अवस्थित चेतन म दुबै एउटै हौँ ।

\begin{visheshshloka}[center]

\textbf{वायुरनिलममृतमथेदं भस्मान्तं शरीरम् ।}

\textbf{ॐ क्रतो स्मर कृतं स्मर क्रतो स्मर कृतं स्मर ॥}

॥१७॥

\end{visheshshloka}

\customparagraph{संस्कृतव्याख्या}

ईश्वरस्य नश्वररूपम्, अनश्वररूपञ्च ज्ञात्वा ज्ञानी मनुष्यः स्वस्य अन्तिमं लक्ष्यं प्राप्तुं प्रार्थयति, यत्- अथ= शरीरत्यागानन्तरम्, वायुः = शरीरान्तर्गतप्राणरूपवायुः, तस्मिन् अमृतत्त्वेऽनिले‍= समष्टिवायौ प्रविशेत्= प्रवेशं कुुर्यात्, इदं नश्वरं शरीरं भस्मान्तं भूयात्, अग्नौ प्रविश्य नष्टं भवेत् । ॐ= ओम्‌कारमुच्चार्य उपासितो हे क्रतो ! यज्ञमयपरमेश्वर !, मां स्मर= स्मरणं करोतु, कृतं स्मर‍= मम कृतं सर्वविधं कार्यं स्मरतु । हे क्रतो ! सर्वस्रष्टा ! जगन्नियन्त्रक ! मम सर्वविधं कृतं कार्यं स्मर ।

\customparagraph{\textbf{नेपाली भावार्थ}}

यो भौतिक शरीर त्यागेपछि मेरो प्राण वायु यही समष्टि वायुमा मिलोस् एवम् शरीर पनि भस्म भई यहीँ समाप्त होस् । ॐकारको उच्चारणद्वारा उपासना गरिएका हे परमेश्वर ! अब मप्रति अनुग्रह गरी आफ्ना कर्तव्यलाई सम्झ । हे ब्रह्मात्मा ! मेरो तिमीप्रतिको साधना र मप्रतिको तिम्रो कर्तव्य दुबैलाई कहिल्यै नबिर्स ।

\begin{visheshshloka}[center]

\textbf{अग्ने नय सुपथा रायेऽअस्मान्}   \textbf{विश्वानि देव वयुनानि विद्वान् ।}

\textbf{युयोध्यस्मज्जुहुराणमेनो}   \textbf{भूयिष्ठां ते नमउक्तिं विधेम ॥}

॥१८॥

\end{visheshshloka}

\customparagraph{संस्कृतव्याख्या}

हे अग्ने ! सर्वात्मनिष्ठ देव !, अस्मान्= मादृशान् मनुष्यान्, राये= उपासनाजन्यफलभोगाय सुपथा= शोभनेन मार्गेण अन्तिमं लक्ष्यं प्राप्तुं नय= प्रेरय । हे देव ! दिव्यदानादिगुणयुक्त भगवन् !,  विश्वानि= सम्पूर्णानि वयुनानि= कर्माणि विद्वान्= त्वं ज्ञानवानसि ।  अतोऽस्मत्= अस्माकं जुहुराणम्= मार्गावरोधकम्, एनः= यत् किमपि पापादिकमस्ति, तत् सर्वं युयोधि= दूरं करोतु । विघ्नं दूरीकरणार्थं ते= भवते भूयिष्ठाम्= बारं बारं नम उक्तिम्= नमस्कारवचनम्, विधेम‍= वदाम, कथयाम।

\customparagraph{\textbf{नेपाली भावार्थ}}

हे अग्निदेव ! जीवनभर गरेका उपासनाका फल भोग्न  हामीलाई राम्रा मार्गबाट दिव्यलोक लैजाऊ । हे अग्निदेव ! तिमी हामीले गरेका असल, खराब सबै कर्मका ज्ञाता छौ । हाम्रो दिव्यता प्राप्तिको मार्गलाई अवरोध गर्ने पापलाई नाश गरिदेऊ । हामी तिमीलाई निरन्तर नमन गर्दछौँ ।

\begin{visheshshlokatwo}[center]{84}

\textbf{शुक्लयजुर्वेदसँग सम्बद्ध ईशावास्योपनिषद् पूरा भयो ।}

\end{visheshshlokatwo}

\specialkhanda{\textbf{तवल्कारब्राह्मणस्था}  \textbf{केनोपनिषद्}}

\customparagraph{\textbf{परिचयः}}

इयम् उपनिषद् सामवेदसम्बद्धास्ति । सामवेदान्तर्गते तवलकारब्राह्मणे चतुर्थाध्याये यो दशमोऽनुवाक् विद्यते, तत्र केनोपनिषद् सङ्‌गृहीताऽस्ति । अस्यामुपनिषदि चत्वारः खण्डाः सन्ति। परमात्मा एव सर्वकर्ता, तस्यैव इच्छाशक्त्या, ज्ञानशक्त्या, क्रियाशक्त्या च सर्वं प्रवर्तते इति अस्या उपनिषदः प्रधानो विषयः । ‘केनैषितं पतति प्रेषितं मनः—’ इति प्रथममन्त्रस्य केन इति प्रथमशब्देन अस्य नाम ‘केन’ इति सञ्जातमिति ज्ञायते । इयं ब्रह्मणोपनिषद् विद्यते।

\customparagraph{\textbf{सारः/वैशिष्ट्यञ्च}}

केनोपनिषदानुसारं सर्वं जगत् परब्रह्म सञ्चालयति । परब्रह्मैव सर्वज्ञः सर्वकर्ता सर्वनियामकश्च विद्यते। परब्रह्म एव देवादीन् स्वकीयेषु कार्येषु संयोजयति । लोके जना चर्मचक्षुषा दृश्यं तत्त्वं ब्रह्मत्वेन स्वीकृत्य पूजयन्ति, परं तानज्ञानसागरे निमग्नान् जानीहि । ब्रह्मतत्त्वं त्वभौतिकमनुभूतिबोध्यं भवति । ब्रह्म एव अदृश्यरूपेण चराचरं जगत् सञ्चालयति । ‘अहं ब्रह्म वेद्मि’ इति यः चिन्तयति, स ब्रह्मविषये किमपि न जानाति परं ‘ज्ञानविषयात् परमपारं ब्रह्म, तदात्मरूपेण सर्वत्र व्याप्य तिष्ठति, तस्य विषये पूर्णरूपेण ज्ञानमसम्भवप्रायमिति’ चिन्तयति, एष चिन्तक एव ब्रह्मवेत्ता इति मन्यते । तस्योपरि नामरूपधारकाः केऽपि न विद्यन्ते, देवा अपि तस्यैव निर्देशानुसारं कार्यं कुर्वन्ति । अतः  परब्रह्म एव सर्वशक्तिमान् इति विषयस्य प्रबोधाय तत्र एका आख्यायिका वर्णिता विद्यते, यत्-

एकदा देवैर्दानवेषु विजयो लब्धः । स विजयः परमात्मकृपया एवासीत् परं देवा अहङ्कारेण अस्माकमेव विजयोऽयमिति कथितवन्तः । अस्माकं सामर्थ्येनैव विजयः प्राप्त इति तेषामभिमान आसीत् । एतादृशं देवानामहङ्कारं शमयितुं परमात्मा यक्षस्वरूपं गृहीत्वा दूरतः तिष्ठति। देवा एनं यक्षरूपं परमात्मानं पश्यन्ति परं न जानन्ति । ततो देवेषु एष दृश्यमानो यक्षः कः ? इति जिज्ञासा जायते । अतः प्रथमतो देवा यक्षं ज्ञातुम् अग्निं प्रेषयन्ति ।

‘ते अग्निमब्रूवन्-  ‘हे जातवेदः ! एतद् विजानीहि कोऽयं ‘यक्षः’ इति । सोऽपि ‘तथेति’ स्वीकृत्य यक्षसमीपं गच्छति । तदा प्राक् यक्ष एव अग्निं पृच्छति, यत्- ‘त्वं कोऽसि’ इति। अग्निर्वदति यत्- अहमग्निरस्मि । ‘जातवेदाः’ इति नाम्ना अहं प्रसिद्धः । तदनन्तरं यक्षेन पृष्टं यत्- ‘त्वयि किं वीर्यम् ? त्वयि कीदृशी शक्तिर्वर्तते ?’ तदा अग्निर्वदति यत्- ‘यदस्यां पृथिव्यां यत्किञ्चिद् वस्तु वर्तते तत् सर्वं दहेय’मिति । तदा यक्षोऽग्नेः पुरः तृणखण्डं स्थापयित्वा वदति यत्- ‘एतं दह’ इति । किन्तु अग्निस्तद्दहनेऽसमर्थो भवति । निराशोऽग्निः प्रत्यावर्त्य देवानां पुरो वदति यत्- ‘अहं न जानामि कोऽयं यक्षः’ इति । तदनन्तरं देवैः प्रेषितो वायुः समागच्छति । तदा प्राक् अग्निना सह यादृशो वार्तालापोऽभवत्, वायुना सहापि तादृश एव संवादो भवति । सोऽपि यक्षस्य परीक्षायामनुत्तीर्णो भवति । प्रत्यावर्त्य वदति यत्- ‘कोऽयं यक्षः’ इति ज्ञातुमहं न समर्थः’ इति । तदनन्तरं देवराजस्तत्र गच्छति, परं तदा स यक्षोऽदृश्यो भवति । तत्स्थाने ‘उमा’ नाम्नी एका स्त्री उपस्थिता भवति । सा देवान् कथयति यत्- ‘भो देवाः ! भवन्तोऽहङ्कारेण अमन्यन्तो यत्- अस्माकमेव विजयोऽयम्, सर्वो महिमा चास्माकमेव इति । परं तत्सर्वं परमात्मनः कृपया सञ्जातमिति बोधाय स्वयं परमात्मा यक्षरूपं धृत्वा समागत आसीत् । तस्माद् भवद्भिरहङ्कारो न कर्तव्यः’ इति । एवमत्र आख्यायिकायाः सारो विद्यते यत्- देवानामपि प्रेरकं परबह्म चेत् किं नाम अन्येषाम्, अतो  ब्रह्माण्डेषु सर्वधारकः सर्वसञ्चालकश्च परमात्मा एव इति उपदेशः ।

यो जनः सम्यक् रूपेण ब्रह्माण्डसञ्चालकस्य परमेश्वरस्य ज्ञानगर्भितं ब्रह्मतत्त्वं जानाति, स भवसागरे स्थितात् पापबन्धनात् मुक्तिं प्राप्नोति । जगति परब्रह्म एव अद्वैता शक्तिरिति सारतया अद्वैतवादस्य बीजमत्र वर्णितं विद्यते ।

\begin{center}

\textbf{शान्तिपाठः}

\end{center}

\begin{visheshshloka}[center]

\textbf{ॐ अप्यायन्तु ममाङ्गानि वाक् प्राणश्चक्षुः श्रोत्रमथो बलमिन्द्रियाणि च सर्वाणि। सर्वं ब्रह्मोपनिषदं माऽहं ब्रह्म निराकुर्याम्, मा मा ब्रह्म निराकरोदनिराकरणमस्त्वनिराकरणं मे अस्तु । तदात्मनि निरते य उपनिषत्सु धर्मास्ते मयि सन्तु ते मयि सन्तु॥}

ॐ शान्तिः शान्तिः शान्तिः !

\end{visheshshloka}

\customparagraph{संस्कृतव्याख्या}

मम अङ्गानि- वाक्, प्राणः, चक्षुः, श्रोत्रमित्यादीनि पञ्चज्ञानेन्द्रियाणि अथ=अनन्तरं बलम्, एवं पञ्चकर्मेन्द्रियाणि चेति सर्वाणि इन्द्रियाणि चाप्यायन्तु= परिपुष्टानि भवन्तु, उपासनाध्यानादिकार्यार्थमित्यर्थः । सर्वमुपनिषदम्= उपनिषदि प्रतिपाद्यं ब्रह्म एवास्ति । अहं ब्रह्म मा निराकुर्याम् । माञ्च ब्रह्म मा निराकरोत् । अनिराकरणम्= निराकरणाभावात्मक आवयोः सम्बन्धोऽस्तु= भवेत् । मे= मां प्रति ब्रह्मणोऽनिराकरणमस्तु । तत्= एतत्प्रकारेण आत्मनि= ब्रह्मणि निरते= निमग्ने मयि, उपनिषत्सु= उपनिषद्‌ग्रन्थेषु ये धर्माः= उपासनायामावश्यकाः शमदमादिधर्माः सन्ति ते मयि सन्तु= भवन्तु, ते मयि सन्तु= अवश्यमेव भवन्तु ।

\customparagraph{\textbf{नेपली भावार्थ}}

मेरा सम्पूर्ण अङ्गहरू सबल बनून्;  पञ्चज्ञानेन्द्रिय, पञ्चकर्मेन्द्रिय सामर्थ्यले पूर्ण बनून् । तिनीहरू ब्रह्मतत्त्वको चिन्तन गर्न योग्य बनून् । यो सबै उपनिषद् प्रतिपाद्य ब्रह्म नै हो । म विचलित भएर ब्रह्मलाई तिरस्कार नगरूँ, मलाई पनि ब्रह्मले तिरस्कार नगरोस् । यसप्रकारले दुबैमा आदानप्रदानको सुमधुर सम्बन्ध बनिरहोस् । निरन्तर ब्रह्ममा समर्पित ममा उपनिषद्‌मा व्याख्या गरिएका शम, दम आदि सद्‌गुणहरू विकसित होऊन् । ती सद्‌गुणहरू निरन्तर ममा रहून् ।

\specialupakhanda{\textbf{प्रथमः खण्डः}}

\begin{visheshshloka}[center]

\textbf{ॐ केनेषितं पतति प्रेषितं मनः}   \textbf{केन प्राणः प्रथमः प्रैति युक्तः ।}

\textbf{केनेषितां वाचमिमां वदन्ति}   \textbf{चक्षुः श्रोत्रं क उ देवो युनक्ति ॥}

॥ १॥

\end{visheshshloka}

\customparagraph{संस्कृतव्याख्या}

केन= कर्तृतत्त्वेन इषितम्=प्रेरितं प्रेषितं मनः पतति= विषयं प्रति गच्छति । केन= कर्त्रा युक्तः= प्रेरितः प्राणः प्रथमः= इन्द्रियव्यापारात् प्राक् प्रैति= गमनागमनं करोति । केन= कर्त्रा इषिताम्= प्रेरिताम्, इमां वाचम्= वाणीम्, वदन्ति= उच्चारयन्ति। चक्षुः= दर्शनेन्द्रियम्, श्रोत्रम्= श्रवणेन्द्रियम्, क उ देवः= किं नाम देवः, युनक्ति= प्रेरयति ।

\customparagraph{\textbf{नेपाली भावार्थ}}

कुन तत्त्वद्वारा प्रेरणा प्राप्त गरेर मन विषयतिर उन्मुख हुन्छ ? कसद्वारा प्रेरित भई इन्द्रियव्यापारभन्दा पहिले भित्रबाहिर गर्दछ ? कसबाट प्रेरित भई बोली उच्चारण गरिन्छ ? हेराइ सुनाइ, आदि क्रियालाई कुन देवताले प्रेरित गर्छन् ?

\begin{visheshshloka}[center]

\textbf{श्रोत्रस्य श्रोत्रं मनसो मनो यद्}  \textbf{वाचो ह वाचं स उ प्राणस्य प्राणः ।}

\textbf{चक्षुषश्चक्षुरतिमुच्य धीराः}   \textbf{प्रेत्यास्माल्‌लोकादमृता भवन्ति ॥}

॥ २॥

\end{visheshshloka}

\customparagraph{संस्कृतव्याख्या}

यत्= तत्त्वविशेषं श्रोत्रस्य= श्रवणेन्द्रियस्यापि श्रोत्रमस्ति । मनसोऽपि मनोऽस्ति । वाचः = वाण्या अपि वाचमस्ति ।  प्राणस्य= जीवसारस्यापि प्राणोऽस्ति । चक्षुषः= दर्शनेन्द्रियस्यापि चक्षुरस्ति । धीराः= विवेकशीला जनाः, अतिमुच्य= श्रोत्राद्यात्मभावं परित्यज्य अस्माल्लोकात्= पुत्रबान्धवादिव्यवहारलक्षणाद् इहलोकात् । प्रेत्य= गत्वा अमृताः= मृत्युरहिता भवन्ति ।

\customparagraph{\textbf{नेपाली भावार्थ}}

जुन महान् तत्त्व कानमा सुन्ने शक्तिका रूपमा, मनमा सोच्नेे शक्तिका रूपमा, वाणीमा बोल्ने शक्तिका रूपमा, प्राणमा जिउने शक्तिका रूपमा र आँखामा हेर्ने शक्तिका रूपमा रहेको छ, त्यो आत्मतत्त्व हो । अर्थात् यी सबैलाई सञ्चालन गर्ने प्रेरक बनेर त्यही आत्मतत्त्व रहेको हुन्छ । विवेकशील मानिसहरू जब श्रोत्रादि इन्द्रियहरूमा आत्मत्वभावलाई त्यागी तिनलाई प्रकाशित गर्ने चेतन तत्त्व नै आत्मा हो भन्ने अवस्थामा पुग्दछन् तब तिनीहरू मरणरहित हुन्छन् ।

\begin{visheshshloka}[center]

\textbf{न तत्र चक्षुर्गच्छति न वाग्गच्छति नो मनो न विद्मो न विजानिमो यथैतदनुशिष्यादन्यदेव तद्विदितादथो अविदितादधि । इति शुश्रुम पूर्वेषां ये नस्तद्विचचक्षिरे ॥ ३॥}

\end{visheshshloka}

\customparagraph{संस्कृतव्याख्या}

न तत्र चक्षुर्गच्छति= तद् ब्रह्मतत्त्वं चर्मचक्षुःप्रत्यक्षं न भवति । न वाग्गच्छति= वाण्या आह्वयितुं वर्णयितुं वा नैव शक्यते । तत्र मनो न गच्छति= गन्तुं समर्थो न भवति । (ब्रह्म मनसो विषयो न भवति) यथा एतत्= ब्रह्मतत्त्वम्, अनुशिष्यात्= उपदेष्टुं शक्यते न वा इति वयं न विद्मः। न विजानीमः= तस्य वैशिष्ट्यमपि ज्ञातुं न शक्नुमः । अथ= एवं प्रकारेण तत्= ब्रह्म विदितात्= ज्ञानविषयात्, अन्यदेव= भिन्न एव । अविदितात्= ज्ञानाविषयभूताया मायायाः, अधि= उपरि विद्यते । इति= इत्यनेन प्रकारेण पूर्वेषाम्= प्राग्जातानां विदुषां वचांसि वयं शुश्रुम=शृणुम । ये= पूर्वाचार्याः, नः= अस्मान् तत्= ब्रह्मविषये विचचक्षिरे= कथितवन्तः ।

\customparagraph{\textbf{नेपाली भावार्थ}}

ब्रह्मतत्त्वलाई चर्मचक्षुले देख्न सक्दैन, त्यो वाणी र मनको पहुँचमा पनि छैन । त्यसका बारेमा वास्तविक उपदेश गर्न पनि सकिँदैन । त्यसको यथार्थ वैशिष्ट्य पनि जानिँदैन । त्यसैले त्यो ब्रह्मतत्त्व ज्ञानको विषयभन्दा पर छ । ज्ञानको विषयभन्दा पर रहेको मायातत्त्वभन्दा पनि माथि छ । यसप्रकारका पूर्वजका वाणीहरू हामीले सुनेका थियौं । ती ब्रह्मतत्त्ववेत्ता विद्वान् आचार्यहरूले त्यो ब्रह्मको बारेमा हामीलाई भनेका थिए ।

\begin{visheshshloka}[center]

\textbf{यद्वाचाऽनभ्युदितं येन वागभ्युद्यते ।}

\textbf{तदेव ब्रह्म त्वं विद्धि नेदं यदिदमुपासते॥}

॥ ४॥

\end{visheshshloka}

\customparagraph{संस्कृतव्याख्या}

यत्= ब्रह्मतत्त्वम्, वाचा= वाण्या अनभ्युदितम्= वर्णितुं न शक्यते । येन= तत्त्वेन वाक्= वाणी अभ्युद्यते= वक्तुं शक्यते । त्वं तदेव तत्त्वं ‘ब्रह्म’ इति विद्धि= जानीहि । यत्= इन्द्रियज्ञानात्मकं पदार्थमेव ब्रह्मात्मतत्त्वमिति उपासते= आराध्यते, इदं न ब्रह्म ।

\customparagraph{\textbf{नेपाली भावार्थ}}

जुन ब्रह्मतत्त्वलाई वाणीद्वारा स्वरूपादि बताउन सकिंदैन तर जसका सामर्थ्यले प्राणी बोलीद्वारा अनेक कुरा व्यक्त गर्न सक्छ, त्यो नै ब्रह्म हो भनी बुझ । इन्द्रियादिको पहुँचमा रहेको पदार्थलाई यही नै ब्रह्म हो भनी जसको मानिसहरू आराधना गर्दछन्, त्यो ब्रह्म होइन ।

\begin{visheshshloka}[center]

\textbf{यन्मनसा न मनुते येनाहुर्मनो मतम् ।}

\textbf{तदेव ब्रह्म त्वं विद्धि नेदं यदिदमुपासते॥}

॥ ५॥

\end{visheshshloka}

\customparagraph{संस्कृतव्याख्या}

यत्= तत्त्वम्, मनसा= अन्तःकरणेन न मनुते= ज्ञातुं न शक्यते परं यस्य कारणेन प्राणी मनसा विविधं कर्तव्याकर्तव्यं चिन्तयितुं समर्थो भवति तदेव ब्रह्म । त्वं तदेव विलक्षणसामर्थ्ययुतं तत्त्वं ब्रह्मतत्त्वमिति विद्धि= जानीहि । लोके नेत्रप्रत्यक्षं यत् तत्त्वं ब्रह्मतत्त्वमिति जनैराराध्यते तद् ब्रह्मतत्त्वं नैव ।

\customparagraph{\textbf{नेपाली भावार्थ}}

जुन तत्त्वलाई अन्तस्करणले जान्न सक्दैन तर जसका कारण प्राणीको मन विभिन्न थरी विचार गर्न सक्ने हुन्छ, त्यही सर्वसामर्थ्ययुक्त तत्त्वलाई तिमी ब्रह्मतत्त्व हो भनी बुझ । जुन इन्द्रिय प्रत्यक्ष पदार्थलाई ब्रह्म हो भनी मानिसहरू आराधना गर्छन्, त्यो ब्रह्म होइन । (त्यही नै ब्रह्म हो भनेपछि पुनः ‘नेदं ब्रह्म’ भनी आत्मेतर पदार्थको ब्रह्मत्व निराकरण गरिएको हो)

\begin{visheshshloka}[center]

\textbf{यच्चक्षुषा न पश्यति येन चक्षूंषि पश्यति ।}

\textbf{तदेव ब्रह्म त्वं विद्धि नेदं यदिदमुपासते ॥}

॥ ६॥

\end{visheshshloka}

\customparagraph{संस्कृतव्याख्या}

यत्= तत्त्वम्, चक्षूषा= चर्मचक्षुषा न पश्यति= द्रष्टुं न शक्यते परं येन नेत्रं दर्शनसामर्थ्ययुतं भवति । त्वं तदेव तत्त्वं ब्रह्म इति जानीहि । लोके नेत्रप्रत्यक्षं तत्त्वमेव इदं ब्रह्म इति जनैरुपास्यते तद् ब्रह्म नैव ।

\customparagraph{\textbf{नेपाली भावार्थ}}

जुन तत्त्वलाई चर्मचक्षुद्वारा देख्न सकिँदैन तर त्यही चर्मचक्षुले जसबाट हेर्ने सामर्थ्य प्राप्त गर्दछ, त्यस तत्त्वलाई तिमी ब्रह्म हो भन्ने बुझ। लोकमा जुन इन्द्रियसम्बद्ध पदार्थलाई यो नै ब्रह्म हो भनी आराधना गर्दछन्, त्यो वास्तविक ब्रह्म होइन ।

\begin{visheshshloka}[center]

\textbf{यत्छ्रोत्रेण न शृणोति येन श्रोत्रमिदं श्रुतम् ।}

\textbf{तदेव ब्रह्म त्वं विद्धि नेदं यदिदमुपासते ॥}

॥ ७॥

\end{visheshshloka}

\customparagraph{संस्कृतव्याख्या}

यत्तत्त्वं श्रोत्रेण= श्रवणेन्द्रियेण न श्रूयते परं येन इदं प्राणिनः श्रोत्रं श्रुतम्= श्रुतिगुणविशिष्टं भवति । त्वं तदेव तत्त्वं ब्रह्मतत्त्वमिति जानीहि । लोके जना यद् भौतिकं तत्त्वं इदमेव ब्रह्म इति मत्वा पूजयन्ति तद् ब्रह्म नैव ।

\customparagraph{\textbf{नेपाली भावार्थ}}

जुन तत्त्वलाई कानद्वारा सुन्न सकिँदैन तर त्यही कानमा सुन्ने क्षमताको जसले प्रतिपादन गरेको छ, तिमी त्यसैलाई ब्रह्म भनेर बुझ । लोकमा अज्ञानी मानिसहरूले भौतिक वस्तुलाई यो नै ब्रह्म हो भनी मान्दछन् तर त्यो वास्तविक ब्रह्म होइन ।

\begin{visheshshloka}[center]

\textbf{यत्प्राणेन न प्राणिति येन प्राणः प्रणीयते ।}

\textbf{तदेव ब्रह्म त्वं विद्धि नेदं यदिदमुपासत् ॥}

॥ ८॥

\end{visheshshloka}

\customparagraph{संस्कृतव्याख्या}

यत्तत्त्वं प्राणेन= प्राणवायुना न प्राणिति = न चेष्टते, न चलति परं येनैव कारणेन वायुः प्राणीषु श्वासप्रश्वासरूपेण चेष्टागुणविशिष्टो भवति  । तदेव त्वं ब्रह्म इति जानीहि । यत्= लोके प्रत्यक्षं पदार्थं ब्रह्म इति मत्वा आराध्यते, इदं ब्रह्म नैव ।

\customparagraph{\textbf{नेपाली}}

जुन तत्त्व प्राणवायुद्वारा श्वासप्रश्वासरूप चेष्टायुक्त हुँदैन तर जसले प्राणवायुलाई श्वासप्रश्वासरूप चेष्टायुक्त गराउँछ त्यसै विलक्षण तत्त्वलाई ब्रह्म हो भनी बुझ । इन्द्रियद्वारा प्रत्यक्ष पदार्थलाई देखाएर ‘ब्रह्म’ हो भनी लोकमा जसको पूजा गर्दछन्, त्यो वास्तविक ब्रह्म होइन ।

\begin{visheshshlokatwo}[center]{84}

\textbf{इति प्रथमः खण्डः}

\end{visheshshlokatwo}

\specialupakhanda{\textbf{द्वितीयः खण्डः}}

\begin{visheshshloka}[center]

\textbf{यदि मन्यसे सुवेदेति दभ्रमेवापि}

\textbf{नूनं त्वं वेत्थ ब्रह्मणो रूपम् ।}

\textbf{यदस्य त्वं यदस्य देवेष्वथ नु}

\textbf{मीमांस्यमेव ते मन्ये विदितम् ॥}

॥ १॥

\end{visheshshloka}

\customparagraph{संस्कृतव्याख्या}

(हे शिष्य !) यदि त्वं ‘अहं सुवेद’ इति मन्यसे तर्हि नूनमेव ब्रह्मणो दभ्रम्= स्वल्पमेव रूपं वेत्थ= जानासि । अथ= तस्मात् अस्य ब्रह्मणो यत्= प्रसिद्धं त्वम्= स्वकीयमात्मतत्त्वम्, एवमस्य ब्रह्मणः यत्= शास्त्रप्रसिद्धम्, देवेषु= इन्द्रादिषु स्थितं स्वरूपं विद्यते तद् मीमांस्यमेव= चिन्तनीयमेव । तस्मादेवाहं ब्रह्मतत्त्वं ते‍= तुभ्यम्, विदितम्= ज्ञातम्, मन्ये= विचारयामि । ‘स्वकीयमात्मतत्त्वम्, इन्द्रादिदेवात्मकमात्मतत्त्वञ्च एकमेवेति यदा नैरन्तर्येण चिन्तयसि तदाऽऽत्मतत्त्वं जानासि, तत एवाहं ‘मम शिष्येण ब्रह्मतत्त्वं ज्ञात’मिति मन्ये इति।’

\customparagraph{\textbf{नेपाली भावार्थ}}

हे शिष्य ! यदि तिमी ‘म सबै राम्ररी जान्दछु’ भन्ने सोच्छौ भने तिमीले ब्रह्मको थोरै अंशमात्र बुझेका छौ । शास्त्रप्रतिपादित इन्द्रादिदेवतामा रहेको अधिदैव चेतन तत्त्व र आफ्नो शरीरभित्र रहेको अध्यात्म चेतन तत्त्वलाई एउटै रूपमा निरन्तर चिन्तन गरेमा मात्र आत्मतत्त्वको वास्तविक स्वरूपको बोध हुन्छ, अनि मात्र म तिमीलाई ‘यी मेरा शिष्यले ब्रह्मलाई बुझे’ भन्ने ठान्दछु ।

\begin{visheshshloka}[center]

\textbf{नाहं मन्ये सुवेदेति नो न वेदेति वेद च ।}

\textbf{यो नस्तद्वेद तद्वेद नो न वेदेति वेद च ॥}

॥ २॥

\end{visheshshloka}

\customparagraph{संस्कृतव्याख्या}

(हे गुरो !) अहं ‘सुवेद= सर्वज्ञानयुत’ इति न मन्ये, न वेद= ब्रह्मतत्त्वं न जानामीत्यपि न चिन्तये, किञ्च वेद च= ब्रह्मतत्त्वं जानाम्यपि । नः= अस्माकं मध्ये यः= शिष्यः तत्= तादृशं रहस्यवाक्यम्, वेद= जानाति । तत्= ब्रह्ममेव वेद= जानाति। न वेद इति= ब्रह्मतत्त्वं न जानातीत्यपि नो= नहि वेद च= जानाति एव ।

\customparagraph{\textbf{नेपाली भावार्थ}}

हे गुरुवर ! म आफूलाई ‘पूर्ण रूपले ब्रह्मज्ञ छु’ भन्दिनँ, त्यस्तै मैले ब्रह्मलाई जानेकै छैन भन्ने पनि होइन; थोरबहुत  जानेको पनि छु । हामी शिष्यहरू मध्ये जसले मैले भनेको कुराको गहिराइ बुझ्दछ, त्यसले नै ब्रह्मतत्त्वको रहस्यलाई बुझ्दछ। उसले ब्रह्मलाई जान्दैन भन्ने होइन, चिन्तशील भए अवश्य जान्दछ ।

\begin{visheshshloka}[center]

\textbf{यस्यामतं तस्य मतं मतं यस्य न वेद सः ।}

\textbf{अविज्ञातं  विजानतां विज्ञातमविजानताम् ॥}

॥ ३॥

\end{visheshshloka}

\customparagraph{संस्कृतव्याख्या}

यस्य= अध्येतृजनस्य कृते अमतम्= ब्रह्मतत्त्वं ज्ञानस्य विषयो भवितुं न शक्नोति इति निश्चयो भवति, तस्य= जनस्य मतम्= ब्रह्मविषये सकारात्मकं मतमर्थात् स ब्रह्म वेत्ति । यस्य मतम्= ब्रह्म ज्ञानस्य विषय इति निश्चयो विद्यते, स न वेद= तादृशो जनो ब्रह्म न जानाति । विजानताम्= ‘अहं ब्रह्मतत्त्वं वेद्मि’ इति निश्चयवतां जनानां कृतेऽविज्ञातम्= ब्रह्मतत्त्वमज्ञातवद् भवति, एवम् अविजानताम्= अपरिमितं ब्रह्मतत्त्वं ज्ञातुमशक्यमिति निश्चयवतां कृते विज्ञातम्= तैर्ब्रह्मतत्त्वं ज्ञातं भवति ।

\customparagraph{\textbf{नेपाली भावार्थ}}

जुन अध्येताले ब्रह्मलाई ज्ञानको विषयमा अटाउने बुझ्दछ, त्यसले ब्रह्मतत्त्वलाई पटक्कै बुझेको हुँदैन । जसले ब्रह्मतत्त्व कुनै परिधिमा नअटाउने अपरम्पार भएकाले ज्ञानको विषयभन्दा पर छ भन्ने जान्दछ, त्यसले ब्रह्मलाई बुझेको मानिन्छ । ब्रह्मतत्त्वलाई इन्द्रियको विषय हो भनी निश्चय गरेर बुझ्न खोज्ने अध्येताहरूका लागि ब्रह्मतत्त्व अज्ञेय छ भने सारा ब्रह्माण्डको सिर्जनाभन्दा पहिले अवस्थित तत्त्व

नै ब्रह्म हो, जो ज्ञानको परिधिभन्दा बाहिर छ’ भनी बुझ्नेहरूका लागि मात्र ब्रह्म विशेषतः प्रकशित हुने स्वभावको हुन्छ।

\begin{visheshshloka}[center]

\textbf{प्रतिबोधविदितं मतममृतत्त्वं हि विन्दते ।}

\textbf{आत्मना विन्दते वीर्यं विद्यया विन्दतेऽमृतम् ॥}

॥ ४॥

\end{visheshshloka}

\customparagraph{संस्कृतव्याख्या}

हि= निश्चितार्थे, प्रतिबोधविदितम्= पूर्णबुद्धिवृत्त्या प्रत्येकेषु बोधेषु यद् ज्ञानं भवति, तदेव मतम्= ब्रह्मज्ञानं मन्यते । तेनैव अमृतत्त्वम्= मोक्षम्, विन्दते= प्राप्नोति । आत्मना= स्थिरेण चित्तेन  वीर्यम्= विद्यारूपं सामर्थ्यम्, विन्दते= प्राप्नोति। एतादृशो जन एव विद्यया= गुरुणा दत्तेनात्मज्ञानेन अमृतम्= विनाशरहितमात्मानं विन्दते= प्राप्नोति, मोक्षं लभते ।

\customparagraph{\textbf{नेपाली भावार्थ}}

प्रत्येक अध्ययनमा गरेको बोधको साक्षीको रूपमा जान्दा नै ब्रह्मलाई वास्तविक रूपले बुझेको मानिन्छ । त्यही बुझाइबाट मोक्ष प्राप्त हुन्छ । वास्तवमा आत्मज्ञानमा मोक्ष प्रदान गर्न सक्ने जुन सामर्थ्य छ त्यो सामर्थ्य आफ्नै एकाग्र चित्तबाट प्राप्त हुने हो, किनकि एकाग्र चित्त भएको व्यक्तिले नै गुरुद्वारा उपदिष्ट ज्ञानलाई ग्रहण गर्न सक्छ र त्यसैबाट त्यसले अविनाशी ब्रह्मात्मालाई प्राप्त गर्दछ।

\begin{visheshshloka}[center]

\textbf{इह चेदवेदीदथ सत्यमस्ति}   \textbf{न चेदिहावेदीन्महती विनष्टिः ।}

\textbf{भूतेषु भूतेषु विचित्य धीराः}   \textbf{प्रेत्यास्माल्लोकादमृता भवन्ति ॥}

॥ ५॥

\end{visheshshloka}

\customparagraph{संस्कृतव्याख्या}

इह= अस्मिन्नेव जन्मनि अवेदीत्, चेत् = यद्यात्मतत्त्वं ज्ञातम्, अथ= अनन्तरं सत्यमस्ति= जीवनं सार्थकं भवति । इह= अत्रैव न अवेदीत् चेत्= यद्यात्मतत्त्वं न ज्ञातम्, तदा महती= असामान्या विनष्टिः= क्षतिर्भवति । धीराः= धैर्ययुता जनाः, भूतेषु भूतेषु= चराचरप्राणिषु विचित्य= ब्रह्मतत्त्वमेव निश्चित्य अस्मात् लोकात्= अवस्थितात् शरीरात् प्रेत्य= आत्मत्वं परित्यज्य, अमृताः= मृत्युहीनाः, भवन्ति= जायन्ते ।

\customparagraph{\textbf{नेपाली भावार्थ}}

यदि मान्छेले यही जन्ममा ब्रह्मतत्त्वलाई चिन्न सक्यो भने त्यसको जीवनले सार्थकता पाउँछ । यसै जन्ममा ब्रह्मलाई चिनेन भने त्यसले ठुलो क्षति व्यहोर्नुपर्दछ । त्यसैले धैर्ययुक्त मानिसहरू चराचर जगत्‌का पदार्थहरूमध्ये ब्रह्मतत्त्वकै पहिचान गरी यो जन्मबाट शरीरमा जुन आत्मज्ञान छ त्यसलाई छाडी अमरत्वलाई प्राप्त गर्दछन् ।

\begin{visheshshlokatwo}[center]{84}

\textbf{इति द्वितीयः खण्डः समाप्तः}

\end{visheshshlokatwo}

\specialupakhanda{\textbf{तृतीयः खण्डः}}

संस्कृतव्याख्या

\begin{visheshshloka}[center]

\textbf{ब्रह्म ह देवेभ्यो विजिग्ये तस्य ह ब्रह्मणो विजये देवा अमहीयन्त। त ऐक्षन्तास्माकमेवायं विजयोऽस्माकमेवायं महिमेति॥ १॥}

\end{visheshshloka}

\customparagraph{संस्कृतव्याख्या}

ह= इति निश्चयार्थे । ब्रह्म देवेभ्यः= इन्द्रादिदेवेभ्यः, विजग्ये= विजयं प्राप्तवान् । ह= निश्चयार्थे । तस्य ब्रह्मणो विजये= सफलतायां देवाः‍= इन्द्रादयः, अमहीयन्तः= महिमायुता जाताः । परं ते देवा ऐक्षन्तः= अचिन्तयन्तो यत्- अयं विजयोऽस्माकमेवास्ति । एवमयं महिमाऽपि अस्माकमेवेति ।

\customparagraph{\textbf{नेपाली भावार्थ}}

निश्चय नै यी ब्रह्मले देवताहरूका लागि विजय प्राप्त गरे । ब्रह्मको विजयका कारणले नै देवताहरू सर्वविजेता, सर्वपूज्य र महिमाशाली बने । त्यसपछि ती देवताहरूले विचार गरे कि ‘यो विजय हामीले आफ्नै सामर्थ्यले प्राप्त गरेका हौं । यो विजयबाट प्राप्त ऐश्वर्य पनि हाम्रै हो ।’

\begin{visheshshloka}[center]

\textbf{तद्धैषां विजज्ञौ तेभ्यो ह प्रादुर्बभूव तन्न व्यजानत}

\textbf{किमिदं यक्षमिति ॥ २॥}

\end{visheshshloka}

\customparagraph{संस्कृतव्याख्या}

तत्= ब्रह्म, एषाम्= इन्द्रादिदेवानाम्, विजज्ञौ= अस्माकमेव विजयः, वयमेव सर्वशक्तिसम्पन्नाः, इत्यहङ्कारो जातः । तेषामहङ्कारमपनेतुं ब्रह्म तेभ्यः= देवताभ्यः, अग्रे प्रादुर्बभूव‍= प्रकटितो जातः, यक्षरूपेण इत्यर्थः । तत्= अदृष्टपूर्वं विचित्राकृतिमयं ब्रह्मतत्त्वम्, इन्द्रादिदेवता इदम्= पुरोदृश्यमाणम्, यक्षम्= विचित्राकृतिमयं जीवतत्त्वम्, किम्= किं नामधेयम् ? इति न व्यजानत= ज्ञानं न प्राप्तवन्तः ।

\customparagraph{\textbf{नेपाली भावार्थ}}

यसरी अज्ञानवश देवताहरूले आफू नै सर्वेसर्वा भएको अनुभव गर्न थाले, भन्ने कुरा परब्रह्मले थाहा पाएपछि तिनीहरूको अभिमानलाई शान्त गराउन यक्षको रूप धारण गरी प्रकट भए । देवताहरू तिनलाई देखेर छक्क पर्दै विचार गर्न थाले कि यो विचित्र स्वरूप भएको जीव को हो ? तिनले परब्रह्मलाई चिन्न सकेनन् ।

\begin{visheshshloka}[center]

\textbf{तेेऽग्निमब्रुवञ्जातवेद एतद्विजानीहि किमिदं यक्षमिति तथेति॥३॥}

\end{visheshshloka}

\customparagraph{संस्कृतव्याख्या}

ते= इन्द्रादिदेवा अग्निम्= अग्निदेवम्, अब्रुवन्= कथितवन्तो यत्- हे जातवेदः ! अग्ने !  इदम्= विचित्राकृतिरूपं यक्षं किम्= कोऽस्ति ? एतत्= वृत्तान्तम्, विजानीहि= विशेषेण परिचयं जानीहि। अग्निदेवोऽपि ‘तथा’ इति स्वीकृतिसूचकं शब्दमुक्तवान् ।

\customparagraph{\textbf{नेपाली भावार्थ}}

यक्षस्वरूप परब्रह्मलाई नचिनेर बिलखबन्दमा परेका देवताले अग्निदेवलाई भने- “हे जातवेदा, अग्नि ! यो विचित्र आकृति भएको जीव को हो ? यसका बारेमा यथार्थ पत्ता लगाऊ ।”  इन्द्रादिदेवताले यसो भनेपछि अग्निदेवले पनि ‘हुन्छ’ भनी प्रतिउत्तर दिए ।

\begin{visheshshloka}[center]

\textbf{तदभ्यद्रवत् तमभ्यवदत्कोऽसीत्यग्निर्वा अहमस्मीत्यब्रवीज्जातवेदा वा अहमस्मीति ॥ ४॥}

\end{visheshshloka}

\customparagraph{संस्कृतव्याख्या}

तत्= यक्षस्थितं स्थानमग्निः, अभ्यद्रवत्= गतवान् । तम्= अग्निं त्वं कोऽसीति, अभ्यवदत्= परब्रह्म अपृच्छत् । वै= निश्चयेनाहम् अग्निः‍= अहं देवेषु मध्ये एकोऽग्निः, जातवेदाः= जन्मना एव सम्पूर्णज्ञानयुतोऽहमस्मीति सर्वे मां वदन्तीति, अभ्यवदत्= अकथयत् ।

\customparagraph{\textbf{नेपाली भावार्थ}}

अग्निदेव शरीरधारी ब्रह्म भएका स्थानमा गए । तिनलाई यक्षरूप परब्रह्मले ‘तिमी को हौ’ भनी  सोधे । अग्निले उत्तर दिए- “म अग्निदेव हुँ, जन्मजात सर्वज्ञ भएकाले मलाई ‘जातवेदा’ पनि भन्छन्” ।

\begin{visheshshloka}[center]

\textbf{तस्मिँस्त्वयि किं वीर्यमिति । अपीदं सर्वं दहेयम्, यदिदं पृथिव्यामिति ॥ ५॥}

\end{visheshshloka}

\customparagraph{संस्कृतव्याख्या}

परब्रह्म अग्निं पृच्छति यत्- “तस्मिन्= तादृशी त्वयि किम्= कीदृशं वीर्यम्= पराक्रमः ?” इति । अग्निः प्रत्युत्तरति यत्- “पृथिव्यां यत्किञ्चिद् विद्यतेे, तत् सर्वं दहेयमिति ।”

\customparagraph{\textbf{नेपाली भावार्थ}}

ती परब्रह्मले अग्निलाई सोधे- तिमीमा त्यस्तो के क्षमता छ ? अग्निले उत्तर दिए- “संसारका जे-जति पदार्थ छन्, म ती सबै भस्म पारिदिन सक्छु ।”

\begin{visheshshloka}[center]

\textbf{तस्मै तृणं निदधावेतद्दहेति । तदुपप्रेयाय सर्वजवेन तन्न शशाक दग्धुं स तत एव निववृते, नैतदशकं विज्ञातुं यदेतद् यक्षमिति ॥ ६॥}

\end{visheshshloka}

\customparagraph{संस्कृतव्याख्या}

तस्मै= अग्नये एतत्= इदं दृश्यमानम्, दह= भस्मसात् कुरु, इत्युक्त्वा तृणम्= शुष्कघासखण्डम्, निदधौ= अग्नेरग्रभागे स्थापयामास । सर्वजवेन= सर्वसामर्थ्येन तत्= तृणं भस्मसात् कर्तुम्, उपप्रेयाय= तत्र गतवान् । परं तत्= तृणम्, दग्धुं न शशाक= न समर्थोऽभवत् । स= अग्निः तत एव= तस्मात् स्थानाद् निववृते= प्रत्यावृतः । स देवानामग्रे गत्वाऽऽह-  यद् एतद् यक्षम्= यक्षवत् पुरोदृश्यमाणं जीवम्, विज्ञातुम्= विशेषेण परिचयं प्राप्तुम्, न अशकम्= नाशक्नुवम् ।

\customparagraph{\textbf{नेपाली भावार्थ}}

ती यक्षले अग्निलाई एउटा सुकेको घाँसको त्यान्द्रो राखेर भने= “यसलाई जलाऊ ।”  अग्निले पनि नजिकै गएर आफ्नो पूर्ण बल लगाएर त्यसलाई जलाउने प्रयास गरे, तर सकेनन् । निराश भई फर्किएर अग्निले देवताहरूलाई भने= हे देवता हो ! मैले त्यस यक्षका बारेमा जान्न सकिनँ ।

\begin{visheshshloka}[center]

\textbf{अथ वायुमब्रुवन् वायवेतद्विजानीहि किमेतद् यक्षमिति तथेति॥७॥}

\end{visheshshloka}

\customparagraph{संस्कृतव्याख्या}

अथ= अनन्तरम्, वायुम्= वायुदेवम् अब्रुवन्= देवाः कथितवन्तः, यत्- हे वायो ! वायुदेव ! एतद् विजानीहि । एतद् यक्षं किमिति । वायुरपि ‘तथा’ इति स्वीकारं कृतवान् ।

\customparagraph{\textbf{नेपाली भावार्थ}}

यसरी अग्नि निराश भएर फर्किएपछि देवताहरूले वायुदेवलाई भने- “हे वायुदेव ! यिनका बारेमा राम्ररी बुझ्नुहोस्, यी यक्ष को हुन् ? यिनका बारेमा जानकारी लिनुहोस्” वायुदेवले पनि ‘हुन्छ’ भनी देवताहरूको अनुरोध स्वीकार गरे ।

\begin{visheshshloka}[center]

\textbf{तदभ्यद्रवत् तमभ्यवदत्कोऽसीति ।}

\textbf{वायुर्वा अहमस्मीत्यब्रवीन्मातरिश्वा वा अहमस्मीति ॥ ८॥}

\end{visheshshloka}

\customparagraph{संस्कृतव्याख्या}

तत्= यक्षं प्रति, वायुदेवोऽभ्यद्रवत्= गतवान् । परब्रह्म तं वायुदेवं “त्वं कोऽसि” इत्यभ्यवदत्= अपृच्छत् । अहं वायुरस्मीत्यब्रवीत्= अकथयत्, अहं ‘मातरि- अन्तरिक्षे, श्वयति- वर्द्धते’ इति ‘मातरिश्वा’ नामधेयोऽप्यस्मि ।

\customparagraph{\textbf{नेपाली}}

देवताहरूको कुरा मानेर ती वायु यक्ष नजिकै गए । यक्षले तिनलाई सोधे- “तिमी को हौ” वायुदेवले उत्तर दिए- “म वायु हुँ, आकाशमा वृद्धि प्राप्त गर्ने हुँदा मेरो नाम ‘मातरिश्वा’ पनि रहेको छ ।”

\begin{visheshshloka}[center]

\textbf{तस्मिन् त्वयि किं वीर्यमिति ? अपीदं सर्वमाददीयं}

\textbf{यदिदं पृथिव्यामिति ॥ ९॥}

\end{visheshshloka}

\customparagraph{संस्कृतव्याख्या}

तस्मिन्= तादृशे त्वयि वायुदेवे किम्= कीदृशम्, वीर्यम्= सामर्थ्यम् ? इति वायुर्यक्षेण पृष्टः । पृथिव्याम्= लोकेऽस्मिन् इदं यत्= यावन्ति वस्तूनि विद्यन्ते, मया इदं सर्वम् आददीयम्= स्वायत्तीकर्तुं योग्यमिति ।

\customparagraph{\textbf{नेपाली भावार्थ}}

ती वायुदेवलाई यक्षले सोधे- “त्यस्ता प्रसिद्ध तिमीमा के क्षमता छ ?” वायु भन्छन्- “म पृथ्वीमा भएका जे-जति पदार्थ छन्, ती सबैलाई आफ्नो वशमा गर्न सक्छु अर्थात् उडाइदिन सक्छु ।”

\begin{visheshshloka}[center]

\textbf{तस्मै तृणं निदधावेतदादत्स्वेति । तदुपप्रेयाय सर्वजवेन}

\textbf{तन्न शशाकादातुं स तत एव निववृते, नैतदशकं विज्ञातुं}

\textbf{यदेतद् यक्षमिति ॥ १०॥}

\end{visheshshloka}

\customparagraph{संस्कृतव्याख्या}

तस्मै= वायुदेवाय एतत्= मया पुरो न्यस्तम्, आदत्स्व= गृहाण इत्युक्त्वा तस्याग्रे पूर्ववत् तृणं निदधौ= स्थापयामास । सर्वजवेन= सर्वसामर्थ्येन तत्= तृणं स्वायत्तीकर्तुम् उपप्रेयाय= गतवान्, परं तत्= तृणम्, आदातुम्=ग्रहीतुम्, न शशाक= असमर्थो जातः । तत एव= तस्मात्स्थानादेव निववृते= प्रतिनिवृत्तः । तत्र देवसमीपे गत्वा वायुः कथयति- “यत् एतद् दृश्यमाणं यक्षं किमिति विज्ञातुं न अशकम्= नाशक्नुवम् ।”

\customparagraph{\textbf{नेपाली भावार्थ}}

ती परब्रह्मले वायुदेवताका अगाडि एउटा झारको त्यान्द्रो राखेर भने- “हे वायुदेव ! यसलाई उडाऊ ।” वायुदेवले पनि नजिक गएर आफ्नो सारा शक्ति लगाएर त्यस त्यान्द्रालाई उडाउन खोजे तर सकेनन् । त्यसपछि त्यहाँबाट फर्केर आई वायुदेवले देवताहरूलाई भने- “हे देव हो ! मैले तिनलाई चिन्न सकिन”

\begin{visheshshloka}[center]

\textbf{अथेन्द्रमब्रूवन् मघवन्नेतद् विजानीहि किमेतद् यक्षमिति ।}

\textbf{तथेति । तदभ्यद्रवत् । तस्मात् तिरोदधे ॥ ११॥}

\end{visheshshloka}

\customparagraph{संस्कृतव्याख्या}

अथ= अनन्तरम् देवा इन्द्रमब्रूवन्= कथितवन्तो यत्- हे मघवन् ! एतद् यक्षं किमिति, एतत्= रहस्यम्, विजानीहि । इन्द्रोऽपि तथा इति स्वीकृत्य तत् स्थानम् अभ्यद्रवत्= ययौ, परं यक्षरूपधारकं ब्रह्म तिरोदधे= अदृश्यमभवत् ।

\customparagraph{\textbf{नेपाली भावार्थ}}

वायुले पनि जान्न नसकेपछि देवताहरूले देवराज इन्द्रलाई नै ‘हे इन्द्र ! यी को हुन् ? यी यक्षका बारेमा जानकारी लिनुहोस्’ भने । इन्द्र पनि देवताहरूको कुरा स्वीकार गरी यक्ष नजिकै गए तर तुरुन्तै यक्षरूपधारी ब्रह्म अदृश्य भए ।

\begin{visheshshloka}[center]

\textbf{स तस्मिन्नेवाकाशे स्त्रियमाजगाम बहुशोभमानामुमां}

\textbf{हैमवतीं तां होवाच किमेतद् यक्षमिति ॥ १२॥}

\end{visheshshloka}

\customparagraph{संस्कृतव्याख्या}

सः= अत्यर्थं ब्रह्म विजिज्ञासुः शान्तभिमानः देवमुख्यः इन्द्रः, तस्मिन्= यत्र ब्रह्मणः प्रादुर्भावः, तिरोधानञ्चासीत् तत्र आकाशे= गगनक्षेत्रे बहुशोभमानाम्= कमनीयसौन्दर्ययुताम्, हैमवतीम्= हिमालयपुत्रीम्, स्त्रियम्= स्त्रीरूपाम्, उमाम्= पार्वतीम्, ददर्श । स इन्द्रः तस्या निकटम् आजगाम= गतवान् । एतत्= प्राग् दृश्यमाणं यक्षं किम् ? इति= अनेन प्रकारेण इन्द्रस्ताम्= भवानीम्, उवाच= अपृच्छत् ।

\customparagraph{\textbf{नेपाली भावार्थ}}

यक्ष अदृश्य भएपछि त्यही ठाउँको आकाशमा इन्द्रले हिमालयपुत्री पार्वतीलाई देखे । इन्द्र नजिक गएर पार्वतीलाई सोध्छन्- “हे भवानी ! ती अघि हामीले देखेका यक्ष को थिए ?”

\begin{visheshshlokatwo}[center]{84}

\textbf{इति तृतीयः खण्डः}

\end{visheshshlokatwo}

\specialupakhanda{\textbf{चतुर्थः खण्डः}}

\begin{visheshshloka}[center]

\textbf{सा ब्रह्मेति होवाच । ब्रह्मणो वा एतद्विजये महीध्वमिति ।}

\textbf{ततो हैव विदाञ्चकार ब्रह्मेति ॥ १॥}

\end{visheshshloka}

\customparagraph{संस्कृतव्याख्या}

सा= ब्रह्मरूपा उमा, उवाच- ब्रह्म इति= भवद्भिः प्राग्दृष्टं यक्षरूपं ब्रह्म आसीदिति। ब्रह्मणः= परमेश्वरस्य एतद्विजये= देवासुरसङ्‌ग्रामे प्राप्ते विजये, महीध्वम्= भवद्भिः स्वकीयं सामर्थ्यमेव कारणत्वेन स्वीकृतम्, परं तत्तु ब्रह्मणः कृपया लब्धमासीत् । तत एव= उमाया वचनानन्तरमेव प्रथमतो देवराजो ब्रह्म इति= ब्रह्म एवासीदिति विदाञ्चकार=ज्ञातवान् इति ।

\customparagraph{\textbf{नेपाली भावार्थ}}

इन्द्रले यक्षका बारेमा सोधेपछि उमा भन्छिन्- ती ब्रह्म थिए । ब्रह्मकै कारण प्राप्त विजयमा तिमीहरूले आफ्नै सामर्थ्यको विजय मानी सर्वशक्तिमान् रहेको अहङ्कार व्यक्त गर्‍यौ, तसर्थ तिमीहरूलाई आफ्नो वास्तविकता बोध गराउन आएका ती ब्रह्म थिए । उमाले यसो भनेपछि सर्वप्रथम देवराजले ती ब्रह्म रहेछन् भन्ने कुरा बुझे ।

\begin{visheshshloka}[center]

\textbf{तस्माद्वा एते देवा अतितरामिवान्यान् देवान् यदग्निर्वायुरिन्द्रस्ते ह्येनन्नेदिष्ठं}

\textbf{पस्पृशुस्ते ह्येनत्प्रथमो विदाञ्चकार ब्रह्मेति  ॥ २॥}

\end{visheshshloka}

\customparagraph{संस्कृतव्याख्या}

तस्मात्= ब्रह्मसाक्षात्कारस्य कारणेन, अग्निः, वायुः, इन्द्रश्चेते त्रयो देवाः, अन्यान्= इतरेभ्यो देवेभ्यः, अतितरामिव= श्रेष्ठाः सञ्जाताः । यत्= यस्मात्, ते= त्रयो देवाः, एनत्= ब्रह्म नेदिष्ठम्=  अतिसामीप्येन पस्पृशुः‍= पृष्टवन्तः, संवादमाध्यमेन ससम्बन्धा जाता इत्यर्थः । ते एनद् ‘ब्रह्म’ इति प्रथम विदाञ्चकार= ज्ञातवन्त इत्यर्थः ।

\customparagraph{\textbf{नेपाली भावार्थ}}

यक्षका रूपमा सर्वप्रथम ब्रह्मलाई साक्षात्कार गरेकाले अग्नि, वायु र इन्द्र यी तीन देवता अरू देवता भन्दा श्रेष्ठ मानिए । यिनीहरूले नै सर्वप्रथम ब्रह्मको स्पर्श गरे, अर्थात् संवादका माध्यमबाट ब्रह्मसँग अत्यन्त नजिकको सम्बन्ध स्थापित गरे । यिनीहरूले नै यक्षलाई ‘ब्रह्म’ रहेछ भनी सर्वप्रथम बुझेका थिए ।

\begin{visheshshloka}[center]

\textbf{तस्माद्वा इन्द्रोऽतितरामिवान्यान् देवान् स ह्येनन्नेदिष्ठं}

\textbf{पस्पर्श, स ह्येनत्प्रथमो विदाञ्चकार ब्रह्मेति ॥ ३॥}

\end{visheshshloka}

\customparagraph{संस्कृतव्याख्या}

यस्मात् अग्निवाय्वपेक्षया इन्द्र एव पूर्वं ब्रह्म ज्ञातवान् तस्मात् कारणात् इन्द्रोऽन्यान्= इतरान् देवान् अतितराम् इव = अतिशेते इव । अग्निवाय्वपेक्षया श्रेष्ठः सञ्जातः । सः= देवमुख्यः, एनत्= ब्रह्म नेदिष्ठम्= सामीप्येन पस्पर्श= पृष्टवान्, संवादमाध्यमेन सन्निकटसम्बन्धं स्थापितवान् । स एनद् ब्रह्म इति प्रथमः= सर्वेषु मध्ये प्रथमः, प्राग् विदाञ्चकार= ज्ञातवान् ।

\customparagraph{\textbf{नेपाली भावार्थ}}

त्यसैले देवराज इन्द्र अरू देवताभन्दा श्रेष्ठ मानिए किनकि उनले ब्रह्मसँग नजिकबाट संवादका माध्यमले सम्बन्ध राख्न सके । यिनले अरू देवताभन्दा पहिले ब्रह्मलाई बुझे ।

\begin{visheshshloka}[center]

\textbf{तस्यैष आदेशो यदेतद्विद्युतो व्यद्युतदा}

\textbf{इतीन्न्यमीमिषदा इत्यधिदैवतम् ॥ ४॥}

\end{visheshshloka}

\customparagraph{संस्कृतव्याख्या}

यदेतत् ब्रह्म विद्युतः = विद्युत्प्रकाशः, व्यद्युतत्=प्रकाशितः, आ= एतादृश एव ब्रह्म इति देवान् बोधयितुम् तस्य एष आदेशः‍= सदृशाभिव्यक्तिरासीत्, तथा च न्यमीमिषत् आ= नेत्रनिमेषवत् दर्शनम्, इति= एतादृश एव अधिदैवतम्= देवसम्बन्धिज्ञानमिति ज्ञातव्यम् ।

\customparagraph{\textbf{नेपाली भावार्थ}}

इन्द्रादि देवताका अघि ‘ब्रह्म’ बिजुली चम्किएजस्तै वा आँखा झिम्किएजस्तै झुलुक्क देखिएको थियो । बिजुली चम्किएझैँ र आँखा झिम्क्याएझैँ झुलुक्क देखिएको क्षणिक अनुभूति नै ब्रह्मको आधिदैविक स्वरूप हो । यही नै ब्रह्मको उपदेश हो ।

\begin{visheshshloka}[center]

\textbf{अथाध्यात्मं यदेतद् गच्छतीव  च}

\textbf{मनोऽनेन चैतदुपस्मरत्यभीक्ष्णं सङ्कल्पः ॥ ५॥}

\end{visheshshloka}

\customparagraph{संस्कृतव्याख्या}

अथ= अनन्तरम्, अध्यात्मम्= आध्यात्मिकं ज्ञानमुपदिश्यते । यत् मनः= अन्तस्करणम्, एतत्= ब्रह्म प्रति गच्छति इवानुभूयते । अनेन च= मनसा एतत्= ब्रह्म उपस्मरति= सामीप्येन स्मर्यते । अभीक्ष्णम्= निरन्तरम्, सङ्कल्पः च= ब्रह्मस्वरूपं विचारयति मन इति ।

\customparagraph{\textbf{नेपाली भावार्थ}}

यसरी ब्रह्मको आधिदैविक स्वरूपको उपदेश गरेपछि आध्यात्मिक स्वरूपको (प्रत्यगात्मको) उपदेश गरिन्छ- जुन मानिसको अन्तर्मन छ, त्यो ब्रह्मतर्फ उन्मुख भएझैं लाग्दछ । यही मनद्वारा यो ब्रह्मलाई नजिकबाट सम्झिइन्छ । साथै मनले नै सधैं ब्रह्मस्वरूपको निरन्तर चिन्तन गरिरहन्छ ।

\begin{visheshshloka}[center]

\textbf{तद्ध तद्वनं नाम  तद्वनमित्युपासितव्यं स य}

\textbf{एतदेवं वेदाभि हैनं सर्वाणि भूतानि संवाञ्छन्ति ॥ ६॥}

\end{visheshshloka}

\customparagraph{संस्कृतव्याख्या}

तत्= ब्रह्म तद्वनं नाम= ‘तद्वनम्’ इति नामधेयमस्ति । तद्वनमिति= पूजनीयत्वेन (वन-सम्भक्तौ, वन्यते, पूज्यते इति वनम्) उपासितव्यम्= आराधितव्यम् ।  सः= तादृशः, यः= उपासकः, एतत्= ब्रह्म एवम्= प्रकारेण वेद= जानाति । एनम्= ब्रह्मोपासकम्, सर्वाणि भूतानि= जगति सञ्जातानि, अभि संवाञ्छन्ति= सर्वतः सम्मानयन्ति ।

\customparagraph{\textbf{नेपाली भावार्थ}}

सकल ब्रह्मको उपासनाको फल बताइन्छ- त्यो ब्रह्म ‘तद्वन’ नामक छ । त्यो सम्पूर्ण प्राणीहरूद्वारा आराधनीय छ । जसले ब्रह्मलाई संसारका प्राणीहरूको आत्माका रूपमा आराधना गर्दछ, त्यस आराधकलाई (अथवा जसले आफूलाई प्रत्यगभिन्न ब्रह्म मान्दछ त्यसलाई) संसारका सबै प्राणीले सम्पूर्ण प्रकारले आदरका साथ आफ्नो आत्मीय ठान्दछन् ।

\begin{visheshshloka}[center]

\textbf{उपनिषदं भो ब्रूहीत्युक्ता त उपनिषद्}

\textbf{ब्राह्मीं वाव त उपनिषदमब्रूमेति ॥ ७॥}

\end{visheshshloka}

\customparagraph{संस्कृतव्याख्या}

भो= हे गुरो ! भवान् उपनिषदम्= ब्रह्मविद्याम्, मां ब्रूहि= उपदेशं करोतु । इति= अनेन प्रकारेण संस्पृष्टे सति गुरुरवोचत्- हे शिष्य ! ते = तुभ्यम्, उपनिषद् उक्ता= प्राक्कथितेन रहस्यमयेन वर्णनेन त्वमुपदिष्टोऽसीत्यर्थः । ते= तुभ्यम्, वाव= निश्चयेन ब्राह्मीम्= ब्रह्मसम्बन्धिनीम् उपनिषदम्= ब्रह्मविद्याम्, वयमब्रूम= वक्ष्याम ।

\customparagraph{\textbf{नेपाली भावार्थ}}

यति सुनिसकेपछि शिष्य भन्छन्- हे गुरुवर ! अब मलाई उपनिषद्‌मा वर्णन गरिएको ब्रह्मविद्याको उपदेश दिनुहोस् । शिष्यको कुरा सुनेर गुरु भन्नुहुन्छ- हे शिष्य ! तिमी  आत्मविद्याले उपदिष्ट भइसक्यौ । पूर्वोक्त प्रकारले तिमीलाई मैले आत्मविद्याको उपदेश गरिसकेँ । अब ब्रह्मविद्याको उपदेश गर्दछ ।

\begin{visheshshloka}[center]

\textbf{तस्यै तपो दमः कर्मेति प्रतिष्ठा वेदाः}

\textbf{सर्वाङ्गाणि सत्यमायतनम् ॥ ८ ॥}

\end{visheshshloka}

\customparagraph{संस्कृतव्याख्या}

तस्यै= ब्रह्मविद्यायै तपः= तपस्या, दमः= इन्द्रियनिग्रहः, कर्म= वर्णाश्रमनिर्धारितानि श्रौतस्मार्तादीनि कर्माणि, वेदाः= ऋगादयश्चत्वारो वेदाः, सर्वाङ्गाणि= शिक्षादीनि षड्‌वेदाङ्गानि, इति= एतादृशाः, प्रतिष्ठाः= ब्रह्मविद्याप्राप्तये साधनानि सन्तीत्यर्थः । (एतेषु हि सत्सु ब्राह्म्युपनिषत्प्रतिष्ठिता भवति ।) सत्यम्= अवितथम्, आयातनम्= आश्रयभूतम्, अर्थात् तानि साधनानि ब्रह्मविद्याया निवासस्थानरूपाणीति ।

\customparagraph{\textbf{नेपाली भावार्थ}}

तपस्या, इन्द्रियनिग्रह, श्रौत, स्मार्त आदि कर्महरू, चार वेद, ६ वेदाङ्ग, यी सबै ब्रह्मविद्याको प्राप्तिका लागि साधन हुन् । यिनीहरूको अवस्थितिमा नै ब्रह्मविद्या प्रतिष्ठित हुन्छ । साँच्चै भन्नु पर्दा ती कर्महरू ब्रह्मविद्याको आश्रयस्थल अर्थात् घर नै हुन् ।

\begin{visheshshloka}[center]

\textbf{यो वा एतामेवं वेदापहृत्य पाप्मानमनन्ते}

\textbf{स्वर्गे लोके ज्येये प्रतितिष्ठति प्रतितिष्ठति ॥ ९॥}

\end{visheshshloka}

\customparagraph{संस्कृतव्याख्या}

यः= मुमुक्षु जनः, एताम्= ब्रह्मविद्याम्, एवम्= अनेन प्रकारेण वेद= जानाति । स पाप्मानम्= अज्ञानरूपं पापम्, अपहृत्य= दूरीकृत्य अनन्ते= विनाशरहिते ज्येये= सर्वश्रेष्ठे स्वर्गे=सुखमये लोके= ब्रह्मलोके प्रतितिष्ठति, प्रतितिष्ठति= निश्चितरूपेणावस्थितो भवति । (प्रतितिष्ठति- सर्ववेदान्तवेद्यं ब्रह्म आत्मत्वेनावगम्य तदेव ब्रह्म प्रतिपद्यत इत्यर्थः)

\customparagraph{\textbf{नेपाली}}

यसप्रकारका कर्म गर्दै मुक्तिको चाहना गर्ने पुरुषले ब्रह्मविद्या जान्दछ । त्यसपछि ब्रह्मविद्या जानेको साधक अज्ञानरूपी पापलाई पखालेर विनाशरहित सर्वश्रेष्ठ सुखमय ब्रह्मलोकमा प्रतिष्ठित हुन्छ । (वेदान्त प्रतिपादित ब्रह्मलाई आत्माका रूपमा जानेर स्वयमेव ब्रह्म बन्दछ) ।

\begin{visheshshlokatwo}[center]{84}

\textbf{इति केनोपनिषदश्चतुर्थः खण्डः समाप्तः}

\end{visheshshlokatwo}

\begin{center}

\textbf{अभ्यासार्थं प्रश्नाः}

\end{center}

\textbf{अतिसङ्‌क्षिप्तोत्तरापेक्षिणः प्रश्नाः (१)}

१. उपनिषदि कस्य चर्चा लभ्यते ?

२. ईशावास्योपनिषद् कस्यां संहितायां विद्यते ?

३. ईशोपनिषदनुसारं सर्वव्यापकः कोऽस्ति ?

४. काण्वशाखान्तर्गतायामीशोपनिषदि कति मन्त्राः सन्ति ?

५. ईशोपनिषदि वर्णिताया विद्याया अविद्याया चार्थः कः ?

६. ‘विद्ययाऽमृतमश्नुते’ इत्यस्यार्थः कः ?

७. ईशकेनोपनिषदोर्नामकरणस्याधारः कः ?

८. सृष्टिक्रमं कः सञ्चालयति ?

९.   ब्रह्मणो रूपं को जानाति ?

१०. देवेषु कीदृशोऽहङ्कारः समागतः ?

११. देवानां परीक्षणार्थं परमात्मा केन रूपेण प्रदुर्बभूव ?

१२. अग्निः परीक्षायां कथं पराजयमनुभवति ?

१३. का इन्द्रं परब्रह्मणः सर्वश्रेष्ठत्वं श्रावयति ?

१४. ब्रह्मविद्याया ज्ञानस्य फलं कीदृशं भवति ?

\textbf{सङ्‌क्षिप्तोत्तरापेक्षिणः प्रश्नाः (६)}

१. ईशोपनिषदः ‘ब्रह्ममयं जगत्’ इति विषयं समर्थ्यताम् ।

२ .केनोपनिषदा प्रस्तुतां ब्रह्मविद्याया ज्ञानयुक्तिं प्रकाशयत ।

\specialmahakhanda{परिशिष्टः}

\specialupakhanda{\textbf{पाठ्यक्रमः}}

\begin{center}

नेपालसंस्कृतविश्विवद्यालयस्य साहित्यविषयस्य

शास्त्रिस्तरीयषष्ठपत्रस्य वैदिकसाहित्यांशस्य

पाठ्यांशविषयाः

प्रकरणम् - ‘क’      						              पूर्णाङ्कः - ३५

\customparagraph{प्रथमः खण्डः}

\end{center}

१.   वेदवाङ्‌मयस्य सामान्यः परिचयः ।

२.   प्राचीनतमाया ऋक्संहिताया विशेषता, साहित्यिकदृष्ट्या ऋग्वेदस्य विश्लेषणम् ।

३.   दार्शनिकसूक्तवैशिष्ट्यपरिचयप्रदानपूर्वकं नासदीयसूक्तान्तर्गतानां

(ऋक् १०/ १२९) सप्तानां मन्त्राणां विशिष्टमध्ययनम् ।

४.   ऋग्वेदगतप्रकृतिवर्णनवैशिष्ट्यप्रतिपादनपूर्वकं उषःसूक्तान्तर्गतानां

(ऋक् १/ ११३) विंशतिपद्यानां विशिष्टाध्ययनम् ।

५.  नाटकाद्याधारभूतसंवादसूक्तपरिचयपूर्वकं पुरूरवउर्वशीसूक्तान्तर्गतानाम्

(ऋक् १०/९५) अष्टादशानां पद्यानां विशिष्टाध्ययनम् ।

६.   अथर्ववेदस्य परिचयः ।

७.   अथर्ववेदीयस्य पृथ्वीसूक्तस्य (१२ /१ /१) अर्थात् शौनकशाखीयस्याथर्ववेदस्य

द्वादशकाण्डान्तर्गतस्य प्रथमानुवाकरूपस्य प्रथमसर्गस्य पदपदार्थविज्ञानपूर्वकं

विशिष्टमध्ययनम्।

\begin{center}

\customparagraph{द्वितीयः खण्डः}

\end{center}

१.    वैदिकवाङ्‌मयान्तर्गतानां ब्राह्मणग्रन्थानां सामान्यः परिचयः ।

२.    ऋग्वेदीयैतरेयब्राह्मणान्तर्गतस्य शुनःशेपाख्यानस्य

(३३ तमोऽध्यायः)  विश्लेषणात्मकमध्ययनम्,

३.    गद्यकाव्यविकासदृष्ट्या शतपथब्राह्मणस्याध्ययनम्,

४.    शतपथब्राह्मणान्तर्गतस्य पुरुरवोर्वश्याख्यानस्य  विशिष्टाध्ययनम्,

५.    बृहदारण्यकोपनिषदन्तर्गतस्य याज्ञवल्क्यमैत्रेय्याख्यानस्य च विशिष्टमध्ययनम् ।

\begin{center}

\customparagraph{खण्डः- तृतीयः}

\end{center}

१.     प्रस्थानत्रय्यन्तर्गतानां ब्रह्मसूत्रगीतोपनिषदां परिचयात्मकमध्ययनम्,

२.     वेदस्यान्तिमभागभूतानां (वेदान्तानां) उपनिषदां विशिष्टः परिचयः ।

३.     ईशावास्योपनिषदन्तर्गतानां मन्त्राणां विशिष्टाध्ययनम् ।

४.     केनोपनिषदन्तर्गतानां (चतुर्ष्वेव खण्डेषु विद्यमानानां) मन्त्राणां  विशिष्टाध्ययनम् ।

\specialupakhanda{\textbf{वैदिकशब्दज्ञानम्}}

वेदः = शास्त्राणि विदन्ति प्राप्नुवन्ती अस्मिन्

ऋक् = ऋच्यन्ते स्तूयन्ते देवा अनया ।

यजुः= इज्यते अनेन ।

साम= ऋचामुपरि गानं भवति स साम।

अथर्वः = अथ लोकमङ्गलाय अर्व्यते प्रस्तूयते इति

संहिता= सम्यक् धीयते स्मर्यते वा ।

ऋत्विक्= ऋतो यजति इति ।

होता = जुहोति इति होता (होमकर्ता)

अध्वर्युः= अध्वरं यज्ञं यौति सम्पादयति  । (यजुर्वेदस्य ऋत्विक्)

उद्गाता= उद्गायति साम यः । (सामवेदस्य ऋत्विक्)

इष्टम् = इज्यते इष्यते वा यज्ञादि कर्म ।

उपनिषद् = उपनिषद्यते प्राप्यते ब्रह्मविद्या अनया

तैत्तरीयः= तैतिरं तित्तिरीप्रोक्तं यजुर्वेदीयशाखामधीते वेत्ति वा ।

ऐतरेय = इतरस्य तन्नाम्न मुनेरपत्यम्

धर्मः= धरति लोकान्, ध्रियते पुण्यात्मभिरिति वा ।

निरुक्तिः= निर्वचनम्, निश्चयकथनम् ।

ब्राह्मणः= ब्रह्म वेदः, तदधीते वेत्ति वा सः ।

शतपथ = शतं पन्थानो यत्र

आरण्यकः = अरण्ये भवः

मन्त्रः= मन्त्र्यते गुप्तं परिभाष्यते, मननात् त्रायते वा ।

वेदव्यासः= वेदं व्यस्यति पृथक्करोति इति वेदव्यासः ।

विद्युत्= विशेषेण द्योतते ।

विप्रः= विशेषेण प्रतिपूरयति षट्‌कर्माणि इति ।

शाकलः= शाकलेन प्रोक्तम् अधीयते शाकलाः, तेषां सङ्‌ग्रहः घोषो वा ।

शाखा= शाखेति गगनं व्याप्नोति  इति गुरुशिष्यपरम्पराध्ययनम् ।

प्रस्थानत्रयम् = प्रस्थानानां ज्ञानमार्गाणां त्रयम्

बालखिल्यः = अङ्‌गुष्ठपरिमिताः षष्टिसहस्रसङ्‌ख्यका मुनयः

माध्यन्दिनः = मध्यन्दिनेन प्रोक्ता अधीता वा

शाकलः = शकलेन मुनिना प्रोक्तमधीयते शाकलाः, तेषां सङ्‌घोऽअङ्को घोषो वा शाकलः

आश्वलायनः = अश्वलः अश्वग्राहको मुनिभेदस्तस्यापत्यम्

कर्मकाण्डः = कर्मणां कर्तव्यताप्रतिपादकं काण्डम्

कल्प = कल्पते समर्थो भवति स्वक्रियायै विरुद्धलक्षणया असमर्थो भवति वा

व्याकरणम् = व्याक्रियन्ते व्युत्पाद्यन्ते अर्थवत्तया प्रतिपाद्यन्ते शब्दा येन

ज्योतिषम् = ज्योतिः सूर्य्यादीनां ग्रहाणां गत्यादिकं प्रतिपाद्यतया अस्त्यस्येति

आस्तिकः = अस्ति परलोक इति मतिर्यस्य

वशिष्ठः = वशवतां वशिनां (जितेन्द्रियाणां) श्रेष्ठः

ब्रह्मचर्यम् = ब्रह्मणे वेदार्थं चर्य्यम् आचरणीयम्

बृहस्पतिः = बृहतां वाचां पतिः

सालोक्यम् = समानो लोको यस्य सलोकः, सलोकस्य भावः

विष्णुधर्मोत्तरपुराणस्य    काव्यशास्त्रसम्बद्धांशः

\specialupakhanda{\textbf{विष्णुधर्मोत्तरपुराणसारः}}

\begin{center}

\customparagraph{शास्त्रिकक्षायास्तृतीयखण्डे निर्धारितः}

\customparagraph{विष्णुधर्मोत्तरपुराणस्य    काव्यशास्त्रसम्बद्धांशः}

\end{center}

१. अनुप्रासयमकाद्यलङ्कारनिरूपणम्

२. शास्त्रेतिहासकाव्यमहाकाव्यानां लक्षणवर्णनम्

३. प्रहेलिकालक्षणभेदकथनम्

४. मन्त्रब्राह्मणकल्पपुराणकथनपुरस्सरं नाटिकादिद्वादशरूपकलक्षणकथनम्

५. नायकनायिकाभेदशृङ्गारादिरसवर्णनम्

६. नृत्तगीतवाद्यवैशिष्ट्यम्

७. अङ्गसञ्चालनादिवैशिष्ट्यम्

८. वाचिकाहार्याङ्गिकसात्त्विकभेदेन

चतुःप्रकारकाभिनयालङ्काराङ्गरचनासज्जीवादिनाट्योपकरणवर्णनम्

१०. इन्द्रियार्थाभिनयप्रकारः

११. इष्टानिष्टमध्यस्थभेदेन गात्रप्रह्लादनादिना चेष्टादिनिरूपणम्

१२. प्रभातगगनरात्रिप्रदोषदिवसर्तुघनधराजलाशयदिङ्‌नक्षत्राद्यभिनयवर्णनम्

१३. शृङ्गारादिनवरसकथनप्रसङ्गेन शृङ्गाराद्धास्यस्य रौद्रात्करुणस्य वीरादद्भुतस्य

बीभत्साद्भयानकस्योत्पत्तिमभिधाय तत्तद्देवतावर्णभेदव्यञ्जनादिनिरूपणम्

\customparagraph{\textbf{विष्णुधर्मोत्तरपुराणस्थकाव्यशास्त्रसम्बद्धांशस्य सारांशः}}

पौरस्त्यकाव्यशास्त्रस्य कृते पौराणिका ग्रन्था महत्त्वाधिकं भजन्ते । प्रायशः सर्वेऽपि पौराणिकग्रन्थाः श्रव्यदृश्यकाव्ययोः कृते कथातत्त्वस्योपजीव्यरूपेण स्वीकृता विद्यन्ते । इत्येव नापितु केषुचिद् ग्रन्थेषु काव्यशास्त्रीयविषयस्य महत्त्वाधायिका चर्चा प्राप्यते । सन्दर्भेऽस्मिन् अग्निपुराणे काव्यशास्त्रीयविषयाणां विवेचनं लभ्यते । एवमेव विष्णुधर्मोत्तरपुराणस्य तृतीये खण्डे मार्कण्डेयवज्रयोः संवादेऽपि काव्यशास्त्रीयविषयाः समुल्लिखिता विद्यन्ते ।

विष्णुधर्मोत्तरपुराणकारेण रूपकरसयोः सन्दर्भे नाट्यशास्त्रस्यानुसरणं कृतं विद्यते । अन्यत्रापि केषुचित् स्थलेषु प्राग्जातानां काव्यशास्त्रिणामनुकरणम्, कुत्रचिद् नवीनमतस्थापनमपि प्राप्यते ।

तृतीयखण्डस्य द्वितीयाध्यायादारभ्य शास्त्रीयविषया विवेचिता विद्यन्ते । तेषु चित्रसूत्रम्, नृत्तशास्त्रम्, गीतम्, शब्दशास्त्रम्, अलङ्कारः, महाकाव्यम्, प्रहेलिका, द्वादशरूपकम्, आहार्याभिनयः, सामान्याभिनयः, रसः, एकोनपञ्चाशद् भावाश्च निरूपिताः सन्ति ।

चित्रसूत्रमित्यत्र चित्रकलाया वर्णनं विद्यते । अस्य पुराणस्य चित्रसूत्रं पौरस्त्यकलासाहित्ये सुप्रसिद्धं विद्यते । अस्य ग्रन्थस्य चित्रसूत्रखण्डस्य कलाविषयकं नवीनं महत्त्वाधायकञ्च मतं विलोक्य प्राज्ञिकैर्विविधेषु विश्वविद्यालयेषु कलाविषयकपाठ्यक्रमे सम्बद्धांशः समाविष्टो विद्यते । तत्र चित्रसूत्रस्य महत्त्वविषये उल्लिखितं विद्यते, यत्-

\begin{shloka}[center]

चित्रसूत्रं न जानाति यस्तु सम्यङ् नराधिपः ।

प्रतिमालक्षणं वेत्तुं न शक्यस्तेन कर्हिचित् ॥

\end{shloka}

एवम् चित्रसूत्रज्ञाने नृत्तशास्त्रस्य ज्ञानमावश्यकमिति मार्कण्डेयः कथयति, यत्-

\begin{shloka}[center]

विना तु नृत्तशास्त्रेण चित्रसूत्रं सुदुर्विदम् ।

जगतोऽनुक्रिया कार्या द्वयोरपि यतो नृपः ॥

\end{shloka}

नृत्तशास्त्रस्य ज्ञानं गीतशास्त्रेण विना अपूर्णमेवेति मुनिः कथयति, यत्-

\begin{shloka}[center]

न गीतेन विना शक्यं ज्ञातुमातोद्यमच्युत ।

गीतशास्त्रविधानज्ञः सर्वं वेत्ति यथाविधि ॥

\end{shloka}

एवं प्रकारेण मार्कण्डेयो मुनी राजानं वज्रं गीतशास्त्रमुपदिशति । प्रसङ्गेऽस्मिन् काव्यशास्त्रीयविषयाः समागताः सन्ति । तत्र गद्यपद्ययोश्चर्चाक्रमे छन्दसो विवेचनम्, अभिधानकोशलिङ्गादिप्रदर्शनञ्च विद्यते । मूलतश्चतुर्दशाध्यायात् काव्यशास्त्रीयविषया वर्णिताः सन्ति ।

\customparagraph{\textbf{चतुर्दशाध्यायः}}

चतुर्दशाध्यायेऽनुप्रासयमकादीनां शब्दालङ्काराणाम्, रूपकादीनामर्थालङ्काराणां केवलं स्वरूपवर्णनं विद्यते । तत्र सप्तदशालङ्काराणामेव सामान्यरूपेण स्वरूपचर्चा प्राप्यते । तेषु १. अनुप्रासः २. यमकम् ३.  व्यतिरेकः ४. श्लेषः ५. उत्प्रेक्षा ६. अर्थान्तरन्यासः ७. उपन्यासः ८. विभावना ९. अतिशयोक्तिः १०. वार्ता  ११. यथासङ्‌ख्यम् १३. विशेषोक्तिः १४. विरोधः १५. निन्दास्तुतिः १६. निदर्शनम् १७. अनन्वयश्च सन्ति । स्थाने स्थाने पुरातनैः कथितमित्युल्लेखनं विद्यते । अतः पूर्वप्रसिद्धिमवलम्ब्य अलङ्कारस्य सामान्यपरिचयप्रकाशनमेव पुराणकारस्योद्देश्यमासीदिति मन्यते ।

विष्णुधर्मोत्तरपुराणकार उपमया परिचित आसीत् । तेन पूर्वाचार्यैः स्पष्टतया वर्णितत्वादुपमालङ्कारस्य चर्चा न कृता इत्यनुमीयते । तत्रोपमा नामोल्लेखनपुरस्सरं चरमश्लोके मार्कण्डेयमुनिरन्यैः काव्यशास्त्रज्ञैः कथिता अलङ्कारा लेशमात्रमेव मया कथिता इत्युल्लेखितवान्, यत्-

\begin{shloka}[center]

विना तथा स्यादुपमा तु यत्र तेनैव तस्यैव भवेन्नृवीर ।

तमन्वयाख्यं कथितं पुराणैरेतावदुक्तं तव लेशमात्रम् ॥

\end{shloka}

\customparagraph{\textbf{पञ्चदशोऽध्यायः}}

अनन्तरं पञ्चदशेऽध्याये शास्त्रेतिहासकाव्यमहाकाव्यानां लक्षणमुल्लिखितं विद्यते । \textbf{शास्त्रम्-} यत् धर्म-अर्थ-काम-मोक्षाणां (चतुर्वर्गस्य) प्राप्तये साक्षादुपदेशं ददाति, तत् `शास्त्रम्' इति कथ्यते । एतत् पूर्वाचार्यैर्निरन्तरं पठितं सुरक्षितं च अस्ति । \textbf{इतिहासः-} इतिहासे प्राग्जातानां सत्पुरुषाणाम् आचरणस्य धर्मार्थकामसाधकत्वस्य च वर्णनं भवति । अत्र विशेषरूपेण मोक्षस्य चर्चा क्रियते। \textbf{काव्यम्-} काव्यमपि पुरुषार्थचतुष्टयसाधकं भवति, किन्तु अत्र शास्त्रवत् शुष्क-उपदेशो न भवति । कस्यचिद् एकस्य नायकस्य चरित्राधृतं सत्, काव्यं सरसगत्या मानवं सन्मार्गं प्रति प्रेरयति । \textbf{महाकाव्यम्-} विष्णुधर्मोत्तरपुराणे विद्यमाने महाकाव्यलक्षणे उल्लिखिता कतपये विषया अन्येषां काव्यशास्त्रिणां महाकाव्यलक्षणे विद्यमानैर्विषयैस्सह साम्यं भजन्ते । अस्य मते महाकाव्ये नायक-प्रतिनायकयोश्चरित्रं वर्णितं भवति । तद् युद्ध-प्रयाणाभ्युदयादिभिरवान्तरविषयवर्णनैर्युक्तं भवति । तत्र देश-नगर-ऋतु-पर्वत-नद्यादीनां मनोहरं चित्रणं भवति । महाकाव्यस्य नायको धर्मविजयी सदाचारी च भवति । प्रतिनायकश्च लोकविजयी भवति ।  प्रतिनायकस्य घातः काव्ये वर्णनीयः, परं नायकस्य मरणं कदापि न वर्णनीयम् । काव्यस्य भाषा शुद्धा, छन्दोबद्धा, श्रुतिमधुरा च भवति; तत्र व्याकरणविरुद्धानाम्, कष्टाक्षराणाम्, अश्लीलानाञ्च शब्दानां प्रयोगो वर्जनीयो भवति । पुनरुक्तिदोषो विस्मय-हर्ष-भय-शोकादिषु भावेषु ग्राह्यो भवति । महाकाव्यं शृङ्गार-वीर-शान्तादिभिर्नवरसैः समन्वितं भवति । एवम्प्रकारेण महाकाव्यलक्षणं समुल्लिख्य  चरमश्लोके साररूपेण महाकाव्ये आवश्यकीयानि वस्तूनि निर्दिष्टानि विद्यन्ते, यथा-

\begin{shloka}[center]

काव्यं कलाकौशलसंप्रयुक्तं धर्मेण चार्थेन तथोपपन्नम् ।

धर्मेण युक्तस्य च नायकस्य कार्यं तथा चाभ्युदयान्तमेव ॥

\end{shloka}

\customparagraph{\textbf{षोडशोऽध्यायः}}

षोडशेऽध्याये प्रहेलिका इति मनोरञ्जनात्मकं काव्यभेदं सलक्षणं प्रतिपादितं विद्यते । राजसभामारभ्य ग्रामीणजनपर्यन्तं मनसो परिष्कारार्थं प्रहेलिकायाः प्रयोगस्य परम्परा प्राचीनकालादेव प्रारब्धा प्राप्यते । कादम्बर्यादिग्रन्थे प्रहेलिकायाः प्रयोगपरम्परा सङ्केतिता विद्यते यत्- “कदाचिदक्षरच्युतक-मात्राच्युतकबिन्दुमतीगूढचतुर्थपादप्रहेलिकाप्रदानैः….. दिवसमेवमारब्धविविधक्रिडापरिहासचतुरैः सुहृद्भिरूपेतो निशामनैषीत्” (कादम्बरी ३६) विशिष्टसामान्यजनैरनुमोदितो मनोरञ्जनप्रदश्च काव्यभेदोऽन्यत्र सादरं स्वीकृतो नोपलभ्यते । । प्रहेलिका एकपद्यात्मकं द्विपद्यात्मकं वा भवेत् । बहुपद्यात्मिका प्रहेलिका न कर्तव्या  इत्युल्लिख्य वर्णनसीमानमपि निर्धारयति । जनसाधारणे प्रचलितासु प्रहेलिकासु अश्लीलदोषो भवति परं काव्ये तद्वर्जनमावश्यकमिति निर्दिशति  यथा-

\begin{shloka}[center]

अश्लीलमेतास्वपि नैव कार्यं तद्वर्ज्यमाहुः कवयो नरेन्द्र ॥

उद्वेजनीयं तु सतां तदुक्तमश्लीलबन्धं सुखदं न काव्यम् ॥

\end{shloka}

\customparagraph{\textbf{सप्तदशोऽध्यायः}}

तदनन्तरं सप्तदशेऽध्याये ब्राह्मणकल्पग्रन्थलक्षणपुरस्सरं दशरूपकविवेचनञ्च विद्यते ।

\textbf{‘मन्त्रब्राह्मणयोर्वेदनामधेयम्’ इत्युक्तेः} सब्राह्मणाः मन्त्राः (संहिताः) वेदपदवाच्याः । मन्त्राणां व्याख्यानात्मका ग्रन्था ब्राह्मणनामधेयाः । कल्पनया (व्यवस्थितयोजनया) युक्तो ग्रन्थः कल्पोऽभिधीयते । पुराणस्य विषये अत्र द्विप्रकारकं विवेचनं दृश्यते । कथावस्तु, उपोद्घातः अनुषङ्गः, संहारश्चेति पुराणं चतुष्पादात्मकं भवति । सर्गः (सृष्टिः), प्रतिसर्गः (प्रलयः पुनस्सृष्टिश्च), वंशः (देवर्षीणां वंशावली), मन्वन्तराणि (कालचक्रं मनूनां कालश्च), तथा च वंशानुचरितम् (राजवंशानां विस्तृतं चरितम्) चेति सुप्रसिद्धं पुराणपञ्चलक्षणञ्चोल्लिखितमस्ति। अत्र चतुष्पदात्मकं पुराणमिति पुराणविषयिणी नूतना धारणा दृश्यते ।

\customparagraph{\textbf{द्वादशरूपकम्}}

नाट्यशास्त्रे दश रूपकाणि अष्टादश उपरूपकाणि च वर्णितानि सन्ति । परमत्र उपरूपकवर्गीकरणं विना रूपकाणां द्वादशभेदाः स्वरूपप्रयोजनादिवर्णनपुरस्सरं वर्णिताः सन्ति । ते च नाटकम्, नाटिका, प्रकरणम्, प्रकरणी, उच्छिष्टाङ्कः (अस्य लक्षणमेव विद्यते), प्राणः, समवतारः, ईहामृगः, व्यायोगः, वीथी, डिमः, प्रहसनञ्चेति द्वादश । तेषु मुख्यस्य नाटकस्य आवश्यकतत्त्वादिसहिता विस्तृता चर्चा विहिता, परमन्येषां रूपकभेदानां तद्वदित्युल्लिख्य सामान्यं विवेचनं कृतं लभ्यते । एवमेव विष्कम्भकप्रवेशकप्रभृतीनां नाट्याङ्गानामपि परिचयः प्राप्यते । पात्राणां शारीरिकस्वरूपम्, भाषास्वरूपम्, सम्बोधनञ्च उल्लिखितं विद्यते । मुखम्, प्रतिमुखम्, गर्वः, विमर्शः, निर्वहणञ्चेति पञ्चसन्धीन् उल्लिख्य वाससज्जा, विरहोत्कण्ठिता, स्वाधीनभर्तृका, कलहान्तरिता, खण्डिता, विप्रलब्धा, प्रोषितभर्तृका (अभिसारिका) चेति अष्टौ नायिकाभेदा अपि वर्णिताः सन्ति । (अष्टौ इत्युल्लिख्य सप्तभेदानामेव लक्षणमत्र निरूपितं विद्यते) द्वादशरूपकभेदविवेचनान्ते तेषु रसस्यानिर्वार्यत्वं सूच्यते । नाट्यशास्त्रे अष्टावेव रसभेदा उक्ताः । परमत्र नाट्येऽपि नवरसा भवन्तीति स्वमतमुपस्करोति पुराणकारः । यथा-

\begin{shloka}[center]

शृङ्गारहास्यकरुणवीररौद्रभयानकाः ॥

बीभत्साद्भुतशान्ताख्या नव नाट्यरसाः स्मृताः ॥

बन्धो रसानुगः कार्यः सर्वेष्वेतेषु यत्नतः ।

रसप्रधानमेवैतत्सर्वनाट्यं नराधिपः ॥

\end{shloka}

द्वादशरूपकानां स्वरूपवर्णनानन्तरं पुराणकारस्ते सर्वे भेदाः कलाकौशलसमन्विता लोकविधानप्रतिपादका, लोकहिताय धर्मादिचतुष्टयोपदेशकाः कार्याः’ इत्युल्लिखति । यथा-

\begin{shloka}[center]

एते कलाकौशलशीलयुक्ताः कार्यास्तथा लोकविधानयुक्ताः ॥

धर्मार्थकामाद्युपदेशगाश्च हिताय लोकस्य नरेन्द्रचन्द्र ॥

\end{shloka}

\customparagraph{\textbf{अष्टादशोऽध्यायः}}

अष्टादशेऽध्याये गीतलक्षणं प्रस्तुतं विद्यते । अस्याध्यायस्य अधिको भागो गद्यशैल्यां विद्यते । अत्र मार्कण्डेयमुनिर्गीतशास्त्रस्य गभीरं विवेचनं करोति । तत्र गीतस्योत्पत्तये अङ्गरूपाणि स्थानानि, ग्रामादीनि गीताङ्गानि, रसस्वरलयानां परस्परसम्बन्धादीनि च वर्णितानि विद्यन्ते । अध्यायान्ते नाट्ये गीतस्य वैशिष्ट्यं प्रतिपाद्यते । गीतं केवलं मनोरञ्जनस्य साधनमेव नापितु ईश्वरप्राप्ते सुलभं साधनमप्यस्ति । यदि कश्चन कुशलो गीतज्ञः सङ्गीतसाधनया साक्षात् परमपदं (मोक्षम्) यदि नाप्नोति तर्ह्यपि तस्य साधना व्यर्था न भवति । स मृत्योरनन्तरं देवस्यानुचरः (पार्षदः) भूत्वा दिव्यलोके भगवता सह शाश्वतमानन्दमनुभवति । तदुक्तम्-

\begin{shloka}[center]

सङ्‌क्षेपेण मया प्रोक्तं सुराराधनकारणात् ।

गीतज्ञो यदि गीतेन नाप्नोति परमं पदम् ।

देवस्यानुचरो भूत्वा तेनैव सह मोदते ॥

\end{shloka}

\customparagraph{\textbf{नवदशोऽध्यायः}}

गद्यशैल्यां रचितेऽस्मिन् अध्याये आतोद्यविधानवर्णनं विद्यते । ‘आतोद्यम्’ इत्यस्यार्थो वाद्यं विद्यते । तत्र वाद्यस्य ततम्, सुषिरम्, घनम्, अनबद्धञ्चेति चत्वारो भेदा भवन्ति । तत्र ततममित्यनेन वीणादि, सुषिरमित्यनेन वंशादि, घनमित्यनेन तालादि, अनबद्धमित्यनेन मुरजादि च ज्ञायन्ते । अध्यायेऽस्मिन् एतेषां  वाद्यसाधनानां वर्णनपुरस्सरं रसैस्सह तेषां सम्बन्धोऽपि वर्णितो विद्यते।

\customparagraph{\textbf{विंशोऽध्यायः}}

नाट्यज्ञैः परस्यानुकृतिर्नाट्यमिति कथितम् । नाट्यस्य शोभायै नृत्तं भवति । तस्य द्वौ भेदौ (आभ्यन्तरं बाह्यञ्च) भवतः । नृत्तस्य प्रयोगो नियमानुसारं मण्डपे करणीय इत्युल्लिख्य मण्डपनिर्माणविधिरप्यत्र वर्णितो विद्यते । एवमेव नृत्ताभिनयसम्बन्धादीन्यावश्यकानि तत्त्वानि च वर्णितानि सन्ति । अध्यायन्ते नृत्तस्य महत्त्वं प्रकाश्य पुराणकारो लिखति– नृत्तं न केवलं रसेन भावेन च युक्तं भवति, अपितु शृङ्गारादिकाव्यरसानामपि पूर्णतयानुकूलं भवति । गीतस्य तालस्य चानुगामी भवति । एतादृशं नृत्तं प्रेक्षकेभ्यः परमं सुखं ददाति, समाजस्य धर्मं कल्याणं च वर्धयति। यथा-

\begin{shloka}[center]

रसेन भावेन समन्वितं च तालानुगं काव्यरसानुगं च ।

गीतानुगं नृत्तमुशन्ति धन्यं सुखप्रदं धर्मविवर्धनं च ॥

\end{shloka}

\customparagraph{\textbf{एकविंशोऽध्यायः}}

अध्यायेऽस्मिन् नटादीनां शैय्यासनस्थानादीनां वर्णनं विद्यते । समम्, कुञ्चितम्, प्रसारितम्, विवर्तितम्, उद्वहितम्, नतञ्चेति षड् शैय्यास्थानानि वर्णितानि सन्ति । येषां प्रयोगोऽभिनयकाले अवस्थाभेदेन भवति।

\customparagraph{\textbf{द्वाविंशोऽध्यायः}}

एकविंशेऽध्याये वर्णितस्य विषयस्य विस्तारो द्वाविंशतितमे अध्याये विद्यते । अत्र उपवेशनस्य (आसनस्य) विविधाः पद्धतयः, यथा- जानुगत-स्वस्तिक-मण्डलासनादयो वर्णिताः सन्ति, ये चिन्ता-निर्वेद-शोक-नृपकर्म-प्रभृतीनाम् अवस्थानां प्रकर्षाभिनयार्थमुपयुज्यन्ते। एवमेव नृपामात्यादीनामासनस्वरूपञ्च वर्णितं विद्यते । अध्यायान्ते केनापि शिष्येण गुरुणा सह कदापि एकस्मिन् आसने न उपवेष्टव्यमिति गुरुशिष्ययोरासनस्वरूपं वर्णितं विद्यते । यथा-

\begin{shloka}[center]

नैकासनं स्याद् गुरुणा तु राजन् गजे रथे वापि शिलातले च ।

नैकासनं स्याद् गुरुणा तु दोषोऽथ वातिदीर्घे फलके च नित्यम् ॥

\end{shloka}

\customparagraph{\textbf{त्रयोविंशोऽध्यायः}}

त्रयोविंशतितमे अध्याये वैष्णवम्, समपादम्, वैशाखम्, मण्डलम्, प्रत्यालीढम्, आलीढञ्चेति षट् उत्थित-स्थानानि (शरीरस्योत्थितावस्था) वर्णितानि, यत्र पाद-विन्यासस्य ताल-प्रमाणस्य च माध्यमेन वीर्य-धैर्य-स्वभावादीनां चित्रणं भवति।

\customparagraph{\textbf{चतुर्विंशोऽध्यायः}}

अङ्गवर्णननामके चतुर्विंशतितमे अध्याये शिरसो विविधा गतयः, (यथा- उद्वहित-नतादयः) अन्येषां शारीरिकाङ्गानां सञ्चालनविधयः, क्रोध-आनन्द-विस्मयादि-भावाश्च निर्दिष्टाः सन्ति।

\customparagraph{\textbf{पञ्चविंशोऽध्यायः}}

पञ्चविंशतितमे अध्याये उपाङ्गानां (नेत्रकटाक्षादीनाम्) द्वात्रिंशत्प्रकाराणां दृष्टीनां वर्णनमस्ति। शृङ्गार-वीर-करुणादि-रसान् तथा अन्य-स्थायिभावान् प्रकटयितुम् अस्य उपाङ्गवर्णनस्य महत्त्वं विद्यते । विशेषतो नृत्तमुपाङ्गाभिनयाश्रितं भवतीत्यपि उल्लिख्यते । यथा-

\begin{shloka}[center]

उपाङ्गकर्माभिहितं तथैतन्मया समासेन नरेन्द्रचन्द्र ।

कार्यः प्रयत्नोऽत्र सदा बुधेन यस्मादुपाङ्गाश्रयमेव नृत्तम् ॥

\end{shloka}

\customparagraph{\textbf{षड्‌विंशोऽध्यायः}}

अध्यायेऽस्मिन् अभिनयस्य मुख्यसाधनत्वेन हस्तानां विविधा मुद्राः (यथा- पताक-मुष्ट्यादयः) शास्त्रोक्तरीत्या निर्दिष्टाः सन्ति।

\begin{shloka}[center]

हस्ता मयैते कथिता नृवीर सर्वं करायत्तमिदं हि नॄणाम् ।

यत्नश्च कार्यस्तु करेषु तस्माच्चातुर्यलास्याभिनयोपपन्नम् ॥

\end{shloka}

\customparagraph{\textbf{सप्तविंशोऽध्यायः}}

अभिनयस्य वाचिकः, आहार्यः, आङ्गिकः, सात्त्विकश्चेति चत्वारः प्रकाराः सन्ति । वचसा क्रियमाणोऽभिनयो वाचिकः कथ्यते । आहार्याभिनयः प्रस्तः, अलङ्कारः, अङ्गरचना, सज्जीवञ्चेति चतुर्विधः । तेषु मृदा, दारुणा, वस्त्रेण, चर्मणा, लोहैर्वा देवदानवयक्षगजाश्वमृगपक्षिणाम् अनुकृतिर्यत्र बिधीयते, तादृशोऽभिनयः प्रस्तः कथ्यते । माल्याभरणवाससां सज्जीकरणेन कृतोऽभिनयोऽलङ्कारः कथ्यते । शारीरिकसंस्कारसम्बद्धाभिनयोऽङ्गरचना कथ्यते।  सज्जीवाभिनये मानवेतराणां जीवानामभिनयो भवति ।

\begin{shloka}[center]

हेतिप्रमोक्षो न तु रङ्गमध्ये कार्यो भवेत् पार्थिववंशमुख्यः ।

आहार्यमेतत्कथितं समासाद्वक्ष्याम्यतश्चाङ्गिकमेव तुभ्यम् ।

\end{shloka}

\customparagraph{\textbf{अष्टाविंशोऽध्यायः}}

अभिनयस्य आङ्गिकाख्यस्तृतीयप्रकारोऽत्र वर्णितो विद्यते । अत्र आङ्गिकाभिनयस्य कृते सामान्य इति शब्दः प्रयुक्तोऽस्ति । अभिनयोऽयं शारीरिकाङ्गभागस्याभिनयो विद्यते। सामान्याभिनयोऽयं नाट्यकलाया आधारभूतो मन्यते, अतोऽस्मिन् प्रयत्नो विधेयः ।  सात्त्विकाभिनयस्तु सत्त्वजन्याभिनयो विद्यते । तस्याभिनयस्य वर्णनं रसभाववर्णनप्रसङ्गे कृतं विद्यते । तत्र सत्त्वजन्यस्वेदरोमाञ्चाद्यभिनयो भवति ।

अभिनयवर्णनप्रकरणान्ते मार्कण्डेयमुनिः ‘अभिनयप्रकारा अनन्ता विद्यन्ते । तेषु केचन मुख्या एव मया श्राविताः’ इति कथयति । यथा-

\begin{shloka}[center]

एतावदेवाभिनये तु तुभ्यं शक्यं मया पार्थिवमुख्य वक्तुम् ।

नाट्यं हि विश्वस्य यतोऽनुकारं कृत्स्नं ततो वक्तुमशक्यमीश ।

\end{shloka}

\customparagraph{\textbf{एकोनत्रिंशत्तमोऽध्यायः}}

अध्यायेऽस्मिन् नाट्यपात्राणां गमनगतिर्वर्णिता विद्यते ।

\customparagraph{\textbf{त्रिंशोऽध्यायः}}

अध्यायेऽस्मिन् शृङ्गार-हास्य-करुण-रौद्र-वीर-भयानक-बीभत्स-अद्भुत-शान्ताश्चेति नवरसा वर्णिता विद्यन्ते । अत्र स्वतन्त्ररूपेण स्वीकृतः शान्तो रसः पृथगेव व्यवस्थापितोऽस्ति । अन्येषां रसानां परस्परं सम्बन्धं दर्शयन् आचार्यः कथयति यत्- शृङ्गाराद् हास्यस्य, रौद्रात् करुणस्य, वीराद् अद्भुतस्य तथा बीभत्साद् भयानकस्य उत्पत्तिर्भवति। एतेषां रसानां विशिष्टा वर्णा अपि निर्दिष्टा विद्यन्ते । यथा- शान्तस्य स्वाभाविकवर्णः, शृङ्गारस्य श्यामः, रौद्रस्य रक्तः, हास्यस्य श्वेतः (सितः), भयानकस्य कृष्णः, वीरस्य गौरः, अद्भुतस्य पीतः, करुणस्य कापोतः (धूसरः) तथा बीभत्सस्य नीलश्चेति रसवर्णा निर्धारिता सन्ति। एवमेवात्र नवरसानां अधिष्ठातृ-देवानां स्पष्टो निर्देशो विद्यते। तद्यथा- हास्यरसस्य देवता प्रमथः (शिवगणः), शृङ्गारस्य भगवान् विष्णुः, रौद्रस्य रुद्रः, करुणरसस्य यमः, बीभत्सस्य महाकालः, भयानकस्य कालदेवः, वीरस्य महेन्द्रः (इन्द्रः), अद्भुतस्य ब्रह्मा तथा शान्तरसस्य परपुरुषः (परमात्मा) इति रसदेवा स्वीकृताः सन्ति।

अनन्तरमत्र नवरसानां स्वरूपं क्रमशो विवेचितं विद्यते । अध्यायान्ते रसस्य महत्त्वं प्रकाशितमस्ति । यथा-

\begin{shloka}[center]

नाट्यस्य मूलं तु रसः प्रदीष्टो रसेन हीनं न हि नृत्तमस्ति ।

तस्मात्प्रयत्नेन रसाश्रयस्य नृत्तस्य यत्नः पुरुषेण कार्यः ॥

\end{shloka}

\customparagraph{\textbf{एकत्रिंशोऽध्यायः}}

अस्मिन् भावाख्येऽध्याये \textbf{एकोनपञ्चाशत् (४९)} भावानां विवेचनं कृतम् अस्ति (स्थायिभावाः ९, सञ्चारिभावाः ३३ सात्त्विकभावाः ८= ४९) । तत्र प्रमुखभावानां स्वरूपम्, तेषाम् अभिनयविधिश्च वर्णितोऽस्ति।

\customparagraph{\textbf{नाट्ये रसभावमहत्त्वम्}}

रसानां भावानां च महत्त्वं प्रतिपादयन् पुराणकारः कथयति, यत्- हे नरेन्द्रचन्द्र ! संसारे कोऽपि नाट्यप्रयोगोऽथवा काव्यं केवलम् एकेन रसेन पूर्णं न भवति। अत्र रसाः, भावाः, वृत्तयश्चेतेषां सर्वेषां सम्मिश्रणं भवति। अत्र एकोनपञ्चाशत् (४९) भावाः तिसृषु अवस्थासु (स्थायि-व्यभिचारि-सात्त्विकभेदेन) वर्णिताः । एते सर्वे भावाः केन प्रकारेण युक्तिपूर्वकं रसेषु प्रयोक्तव्या येन प्रेक्षका आनन्दं प्राप्नुयुः, इत्यादिसर्वमुपदिष्टम्। एतेषां सर्वेषामङ्गानां प्रयोगो नाट्यकाव्ययोरात्मभूतं रसं विचार्यैव कार्यः । तदा एव साफल्यं प्राप्यते ।

\begin{shloka}[center]

न ह्येकरसजं काव्यं किञ्चिदस्ति प्रयोगजम् ।

भावो वापि रसो वापि प्रवृत्तिर्वृत्तिरेव च ॥

एकोनपञ्चाशदिमे समग्रा भावास्त्र्यवस्था गदिता मया ते ।

युक्त्या च ये यत्र रसे नियुक्ताः प्रोक्तं तदप्यत्र नरेन्द्रचन्द्र ॥

\end{shloka}

अत्रोपर्युल्लिखितानि काव्यशास्त्रे प्रसिद्धानि काव्यतत्त्वानि एव विद्यन्ते । विष्णुधर्मोत्तरपुराणस्य तृतीयखण्डस्य १४ अध्यायतः ३१ अध्यायपर्यन्तं काव्यशास्त्रेषु नाट्यशास्त्रेषु वर्णितानां विषयाणामेव नापितु कतिपयेषां नवीनविषयाणामपि सूक्ष्मरूपेण विवेचनं प्राप्यते । पुराणेऽस्मिन् शब्दशास्त्रम्, गीतम्, महाकाव्यम्, अलङ्कारः, प्रहेलिका, नाट्यम्, अभिनयम्, इत्यादिषु काव्यशास्त्रीयविषयेषु सुस्पष्टा चर्चा विद्यते । पुराणेनानेन कुत्रचित् नवीनशैल्या नवीनतथ्यञ्चोद्घाट्य लक्षणशास्त्रेतिहासे विशिष्टं योगदानं कृतं विद्यते ।

\begin{visheshshlokatwo}[center]{84}

\textbf{विष्णुधर्मोत्तरपुराणस्य काव्यशास्त्रीयांशसारः समाप्तः}

\end{visheshshlokatwo}

\specialupakhanda{\textbf{अग्निपुराणस्य काव्यशास्त्रीयांशसारः}}

\begin{center}

\textbf{शस्त्रिकक्षायाः साहित्यविषयस्य षष्ठपत्रे समाविष्टः अग्निपुराणस्य}

\textbf{काव्यशास्त्रीयांशसम्बद्धपाठ्यक्रमः (अध्यायः ३३७-३४७)}

\end{center}

\textbf{१. काव्यादिलक्षणम्, २. नाटकनिरूपणम्, ३. शृङ्गारादिनिरूपणम्, ४. रीतिनिरूपणम्, ५.नृत्यादावङ्गकर्मनिरूपणम्, ६. अभिनयादिनिरूपणम्, ७. शब्दालङ्काराः, ८. अर्थालङ्काराः, ९. शब्दार्थालङ्काराः, १०. काव्यगुणविवेकः, ११. काव्यदोषविवेकः ।}

\begin{center}

\textbf{अग्निपुराणे वर्णितः काव्यशास्त्रीयांशः}

\end{center}

\customparagraph{\textbf{विषयप्रवेशः}}

संस्कृतवाङ्मयस्य पौराणिकसाहित्यं महत्ताधायकं विद्यते । साहित्यिकरचनायां महर्षिव्यासप्रणीतानि पुराणानि उपजीव्यग्रन्थरूपेण स्वीक्रियन्ते । तत्र केषुचिद् ग्रन्थेषु काव्यशास्त्रसम्बद्धा विषया अपि प्रदर्शिताः प्राप्यन्ते । शास्त्रीयविषयवर्णनपरेषु पुराणेषु  अग्निपुराणं प्रसिद्धमस्ति  । तत्र ३३७ तमाद् अध्यायात् ३४७ तममध्यायं यावद् अग्निदेवः महर्षिकाश्यपाय काव्यशास्त्रीयतत्त्वं सम्यगुपदिशति । तत्र काव्यस्वरूपम्, रसः, गुणः, दोषः, रीतिः, अलङ्कार इत्यादिकाव्यशास्त्रीयतत्त्वानां विस्तारपूर्विका चर्चा प्राप्यते। एतदेव नापितु नाट्यतत्त्वविषयेऽपि तत्र साङ्गोपाङ्गं समीक्षणं विद्यते । अग्निपुराणस्य ३३७ तमादध्यायात् ३४७ तममध्यायं यावद् विवेचिता विषयाः साररूपेणात्र समुल्लिख्यन्ते ।

\customparagraph{\textbf{अध्याय ३३७}}

अध्यायेऽस्मिन् काव्यस्य परिभाषा, तस्य महत्त्वम्, गद्यपद्यकाव्ययोश्च विस्तरेण विवेचनं विद्यते । पद्यकाव्यस्य भेदान् समुल्लिख्य सर्गबन्धात्मकस्य महाकाव्यस्य विस्तृता चर्चाऽपि विहिता दृश्यते । काव्यस्य स्वरूपादिविवेचनं प्रायः काव्यशास्त्रिभिः स्व-स्वकाव्यशास्त्रेषु वर्णितैर्विवेचनैस्सह साम्यं भजते। सालङ्कारस्य शब्दार्थस्य काव्यम्, गद्यपद्यकाव्ययोर्भेदाः, महाकाव्यस्य लक्षणम्, महाकाव्ये काव्यात्मभूतस्य रसस्यानिर्वार्यता इत्यादीनां विषयाणां विवेचनं प्रायशः काव्यशास्त्रानुगतमेवावलोक्यते ।

\customparagraph{\textbf{वाङ्‌मयशास्त्रेतिहासादिपरिचयः}}

ध्वनिवर्णपदानां समुच्चयमेव वाङ्‌मयं कथ्यते । तेन च शब्देन शास्त्रम्, इतिहासः काव्यञ्च सङ्‌गृह्यते । शास्त्रेतिहासयोरभिधायाः प्राधान्यं भवति । काव्ये तु गभीरभावप्रकाशनसामर्थ्ययुतस्य व्यङ्‌ग्यार्थस्य प्राधान्यं भवति । अतः शास्त्रेतिहासाभ्यां काव्यं भिद्यते ।

\begin{shloka}[center]

ध्वनिर्वर्णाः पदं वाक्यमित्येतद् वाङ्‌मयं मतम् ॥

शास्त्रेतिहासवाक्यानां त्रयं यत्र समाप्यते ।

शास्त्रे शब्दप्रधानत्वमितिहासेषु निष्ठता ॥

अभिधायाः प्रधानत्वात्काव्यं ताभ्यां विभिद्यते ।

\end{shloka}

\customparagraph{\textbf{शास्त्रमहत्त्वम्}}

शास्त्रमिति सिद्धान्तप्रधानं मन्यते । सिद्धान्तनिर्माणे सर्वविषयबोधमपेक्ष्यते । अतः सर्वेषु वाङ्‌मयभेदेषु शास्त्रनिर्माणं कठिनं भवति । तस्माद् विलक्षणप्रतिभाराहित्येन साधकेन मृग्यमाणमपि  शास्त्रं न सिद्ध्यतीति पुराणकारस्य मतं विद्यते ।

\begin{shloka}[center]

नरत्वं दुर्लभं लोके विद्या तत्र सुदुर्लभा ।

कवित्वं दुर्लभं तत्र शक्तिस्तत्र च दुर्लभा ।

व्युत्पत्तिर्दुर्लभा तत्र विवेकस्तत्र दुर्लभः ॥

सर्वं शास्त्रमविद्वद्भिर्मृग्यमाणं न सिध्यति ।

\end{shloka}

\customparagraph{\textbf{काव्यलक्षणम्}}

सूत्रात्मकशैल्या व्यङ्ग्यभावेन ईष्टार्थव्यञ्जकपदावलीयुतं वाक्यमेव काव्यं भवति । तत् प्रस्फुरद्भिरलङ्कारैरलङ्‌कृतं गुणयुतं दोषवर्जितञ्च भवेत् । तस्य योनिः (उपजीव्यः) वेदः लोकश्च भवति। मन्त्रादिसिद्धं काव्यं तु अयोनिजमपि कथ्यते । (उपजीव्यं विना कल्पनया निर्मितं काव्यमयोनिजं कथ्यते)

\begin{shloka}[center]

काव्यं स्फुरदलङ्कारं गुणवद्दोषवर्जितम् ।

योनिर्वेदश्च लोकश्च सिद्धमर्थादयोनिजम् ॥

\end{shloka}

\customparagraph{\textbf{काव्यभाषा, काव्यभेदाश्च}}

काव्ये संस्कृतभाषा प्राकृतभाषा च प्रयुज्यते । देवराजाब्राह्मणविदूषकादीनां भाषा संस्कृतं भवति । सामान्यतः त्रिविधं प्राकृतं भवति। विशेषतः रूपकेषु भाषाभेदो दृश्यते । संस्कृतादिभाषानिर्मेयस्य काव्यादेर्गद्यं पद्यं मिश्रञ्चेति त्रयो भेदा भवन्ति ।

\begin{shloka}[center]

देवादीनां संस्कृतं स्यात् प्राकृतं त्रिविधं नृणाम्

गद्यं पद्यं च मिश्रं च काव्यादि त्रिविधं मतम् ।

\end{shloka}

\customparagraph{\textbf{गद्यलक्षणं तद्भेदाश्च}}

पादरहितेन पदसन्तानेन रचितं गद्यं कथ्यते । तस्य चूर्णिका, उल्कलिका, वृत्तगन्धिरिति त्रयो भेदा भवन्ति । कोमलपदावलीयुतम्, दीर्घसमाससहितम्, चूर्णिकाख्यं गद्यं भवति । दीर्घसमासयुतम् उत्कलिकाख्यं गद्यं भवेत् । अतिदीर्घसमासरहितं मध्यमस्वरूपकं पद्याभासयुतम् अनुत्कटं गद्यं वृत्तगन्धि इति कथ्यते ।

\begin{shloka}[center]

अपदः पदसन्तानो गद्यं तदपि गद्यते ।

चूर्णिकोत्कलिकावृत्तगन्धिभेदात् त्रिरूपकम् ॥

\end{shloka}

\customparagraph{\textbf{गद्यकाव्यभेदाः}}

गद्यकाव्यस्य आख्यायिका, कथा, खण्डकथा, परिकथा, कथानिका चेति पञ्चभेदा भवन्ति।

\begin{shloka}[center]

आख्यायिका कथा खण्डकथा परिकथा तथा ।

कथानिकेति मन्यन्ते गद्यकाव्यं च पञ्चधा ॥

\end{shloka}

यत्र स्रष्टुर्वंशप्रशस्तिवर्णनम्, कन्याहरणसङ्ग्रामविप्रलम्भादिविपत्तीनां वर्णनम्, वैदर्भ्यादिरीतीनां नागरिकादिवृत्तीनाञ्च संयोजनम्, उच्छ्वासपरिच्छेदैः खण्डविभाजनम्, चूर्णिकागद्यशैल्यां प्रणयनम्, मध्ये मध्ये वक्त्रापरवक्त्रादिच्छन्दसां सम्मिश्रणञ्च भवति सा आख्यायिका कथ्यते ।

\customparagraph{\textbf{कथा}}

यत्र कविः कतिपयैः श्लोकैः स्ववंशं प्रशंसति, मुख्यार्थप्रकाशनाय कथान्तरं योजयति, परिच्छेदादिभिः काव्यं न विभजति, खण्डीकरणस्यावश्यकतायां सत्यां लम्बकादिभिः खण्डयति सा कथा भवति ।

\customparagraph{\textbf{खण्डकथा/परिकथा/कथानिका}}

कथाया गर्भे चतुष्पदी (पद्यम्) संयोज्यते चेत् सा खण्डकथा कथ्यते । खण्डकथायां परिकथायाञ्च अमात्यो वणिग् ब्रह्मणो वा नायको भवति । तयोः करुणो रसश्चतुर्विधो विप्रलम्भो वा मुख्यो रसो भवति । अनयोः काव्यभेदयोः कथाप्रसङ्गस्य समापनं न भवत‍ि अपितु मुख्यकथाया अनुसरणं करोति । कथाख्यायिकयोर्मिश्रणरूपा परिकथा कथ्यते। यस्यां प्रारम्भे भयानकरसः, गर्भभागे करुणरसः, अन्त्ये चाद्भुतरसो भवति सा कथानिका कथ्यते । इयं कथानिका उत्तमा न मन्यते ।

\customparagraph{\textbf{पद्यलक्षणम्}}

हे काश्यप ! (वशिष्ठ) चतुष्पादयुतं पद्यं वृत्तं जातिरिति द्विधा विभज्यते । तयोर्वृत्तं अक्षरैर्गणनीयं भवति । तच्च उक्थं (वैदिकम्) जातिशेषजं (लौकिकम्) चेति द्विविधम् । यत्र मात्राभिर्गनणा क्रियते सा जातिः कथ्यते । छन्दस्शास्त्रज्ञः पिङ्गलो वृत्तस्य समम्, अर्धसमम्, विषमञ्चेति प्रकारत्रयमुल्लिखति । इयं छन्दोविद्या गभीरं काव्यसागरं तितीर्षुणां कृते नौरूपा मन्यते ।

\begin{shloka}[center]

पद्यं चतुष्पदी तच्च वृुत्तं जातिरिति द्विधा ।

वृत्तमक्षरसङ्‌ख्येयमुक्थं तत्कृतिशेषजम् ॥

मात्राभिर्गणना यत्र सा जातिरिति काश्यप !

सममर्धसमं वृत्तं विषमं पैङ्गलं त्रिधा ॥

\end{shloka}

\customparagraph{\textbf{पद्यकाव्यभेदाः}}

चतुष्पादयुतस्य पद्यकाव्यस्य पञ्चभेदा भवन्ति । महाकाव्यम्, कलापः, पर्यायबन्धः, विशेषकम्, कुलकम्, मुक्तकम्, कोषश्चेति ।

\begin{shloka}[center]

महाकाव्यं कलापश्च पर्याबन्धो विशेषकम् ।

कुलकं मुक्तकं कोष इति पद्यकुटुम्बकम् ॥

\end{shloka}

\customparagraph{\textbf{महाकाव्यलक्षम्}}

पौरस्त्यकाव्यशास्त्रे महाकाव्यस्वरूपनिर्धारणस्य परम्परा प्राचीना विद्यते । तस्यां परम्परायां अग्निपुराणस्य महाकाव्यलक्षणमपि प्रसिद्धं विद्यते । तद्यथा-

परिष्कृतया संस्कृतभाषया आरब्धं महाकाव्यमनेकैः सर्गैः सुबद्धं भवति । तस्य स्वरूपमङ्गीकृत्य अन्यं किमपि रचितं चेत् तदपि दुष्टं न मन्यते । तस्य कथा इतिहाससम्बद्धा लौकिकसच्चरित्राश्रिता वा भवति । तत्र स्थानानुकूलं गोप्यमन्त्रणा, दूतप्रयाणम्, सङ्‌क्षेपरूपेण युद्धवर्णनञ्च भवति । शक्वरी-अतिजगती-अतिशक्वरी-त्रिष्टुभ्-पुष्पिताग्राप्रभृतिभिरेवञ्च चारुभिर्वक्त्राभिः समैर्वृत्तैश्च युतं भवति । विभिन्नवृत्तान्तयुताः सर्गा नातिसङ्‌क्षिप्ता भवति । सर्गान्ते वृत्तपरिवर्तनं विधीयते । एकत्र अतिशक्वरिकाष्टिभ्यां वृत्तिभ्यामेकेन सङ्‌‌कीर्णं पद्यं भवति । अपरत्र मात्रावृत्तेन युक्तं भवति । प्रागतिरिच्य पश्चसर्ग उत्कृष्टो भवति । महाकाव्ये प्रकृतिवर्णनमपरिहार्यं भवति । वस्तुवर्णनं सर्ववृत्तिसर्वभावसर्वरीतिसर्वगुणालङ्कारैश्च पुष्टं भवति । एतादृशं विविधगुणगर्भितं महाकाव्यं भवति तस्मात् तस्य रचयिता अपि महाकविः कथ्यते । महाकाव्ये मूलतो वाग्वैदग्ध्यप्रधाने सत्यपि तत्र प्राणरूपेण रसो भवति । तद् अपृथग्यत्ननिर्वर्त्यं वाग्वैचित्र्यरूपं महाकाव्यशरीरं रससाध्यतां याति । तस्य चतुर्वर्गफलप्राप्तिरेव फलं भवति । अग्निपुराणे वर्णितं महाकाव्यलक्षणं दण्डिना कृतेन महाकाव्यलक्षणेन सह कुत्रचिद् भावसाम्यं कुत्रचित्तु पद्यसाम्यमपि भजते ।

\customparagraph{\textbf{अध्याय ३३८}}

अध्यायेऽस्मिन् रूपकस्य सप्तविंशतिभेदान्नुल्लिख्य नाटकस्वरूपम्, अर्थप्रकृतिः, नाटकसन्धिः इत्यादयो विषया वर्णिताः सन्ति । अध्यायान्ते श्रेष्ठनाटकस्य गुणा अपि समुल्लिखिताः विद्यन्ते । एतेषां नाट्याङ्गानां विवेचनं काव्यशास्त्रानुरूपमेवाभिलक्ष्यते ।

\customparagraph{\textbf{सप्तविंशतिरूपकाणि}}

अग्निपुराणे अभिनेयकाव्यवर्णनप्रसङ्गे प्रथमतो रूपकस्य भेदा वर्णिताः सन्ति । नाट्यशास्त्रीयेतिहासे रूपकस्य दश, उपरूपकस्य अष्टादश भेदा लभ्यन्ते, परमग्निपुराणे उपरूपकाणामुल्लेखं विनैव रूपकस्य सप्तविंशतिर्भेदा निरूपिताः सन्ति । नाट्यशास्त्रे वर्णितानां रूपकाणामुपरूपकाणाञ्चात्र वर्णितै रूपकैः सह स्वरूपनामादिषु साम्यं विलोक्यते । तेनानुमीयते यद्- उपरूपकभेदान् अपि रूपकभेदेषु सङ्‌गृह्य सप्तविंशतिरूपकाण्यत्र वर्णितानीति । नाटकम्, प्रकरणम्, डिमः, ईहामृगः, समवकारः, प्रहसनम्, व्यायोगः, भाणः, वीथिः, अङ्कः, त्रोटकम्, नाटिका, सट्टकम्, शिल्पकः, कर्णा, दुर्मल्लिका, प्रस्थानम्, भाणिका, भाणी, गोष्ठी, हल्लीशकम्, काव्यम्, श्रीगदितम्, नाट्यरासकम्, रासकम्, उल्लाप्यम्, प्रेङ्‌खणञ्चेत्यग्निपुराणोक्तानि सप्तविंशतिरूपकाणि । एतेषां रूपकभेदानां स्वरूपमपि उल्लिखितं विद्यते ।

अनन्तरमत्र आमुखः, कथानकभेदाः, अर्थप्रकृतिः, पञ्चसन्धिप्रभृतीनि नाट्यकाव्याङ्गानि वर्णितानि सन्ति । अध्यायस्यान्तिमभागे उत्कृष्टनाटकस्य कृते नाटकगुणा वर्णिता विद्यन्ते ।

\customparagraph{\textbf{अध्याय ३३९}}

अध्यायेऽस्मिन् नाट्यकाव्यस्यात्मभूतो रसो वर्णितो विद्यते । तत्र रसभावादयो निरूपिताः सन्ति ।

अनन्तरं नायिकाभेदा अपि समाविष्टा विद्यन्ते ।

रसविषयेऽपि परम्परामतिरिच्य रसप्रभेदान् दर्शयति । नाट्यशास्त्रे रूपकाणां कृते अष्टावेव रसा निर्धारिताः, परमत्र नवैव रसाः प्रतिपादिताः । प्रथमतः शृङ्‌गाररौद्रवीरबीभत्सा इति चतुरः रसान् जायन्ते । एते परब्रह्मणः सत्त्वादिगुणादुत्पद्यते । चतुर्भ्यो रसेभ्योऽन्ये रसाः समुत्पद्यन्ते । अतः शृङ्गाररौद्रवीरबीभत्सा इति चत्वाररसाः स्वभावजन्याः, अन्ये तु परजन्या इति ।

\begin{shloka}[center]

सत्त्वादिगुणसन्तानाज्जायते परमात्मनः ।

रागाद्भवति शृङ्गारो रौद्रस्तैक्ष्ण्यात्प्रजायते ॥

वीरोऽवष्टम्भजः सङ्कोचभूर्बीभत्स इष्यते  ।

शृङ्गाराज्जायते हासो रौद्रात्तु करुणो रसः ।

वीराच्चाद्भुतनिष्पत्तिः स्याद् बीभत्साद् भयानकः ॥

शृङ्गारहास्यकरुणारौद्रवीरभयानकाः ॥

बीभत्साद्भुतशान्ताख्याः स्वभावाश्चतुरो रसाः ।

\end{shloka}

\customparagraph{\textbf{रसस्वरूपम्}}

सहजानन्दो मानवस्य आत्मनि सदा विद्यमानो भवति। यदा कोऽपि सहृदयः काव्यं शृणोति वा नाटकं प्रेक्षते, तदा विभावादिभिः कारकैः तस्य हृदयस्थितः  सुसुप्त आनन्दो व्यज्यते। इयमेव अलौकिकी आनन्दमयी स्थितिः `रसः' इति नाम्ना प्रसिद्धा अस्ति। एषा रसप्रतीतिर्वस्तुतश्चैतन्यस्य चमत्कार एव, यत्र मानवः सांसारिकप्रपञ्चं विस्मृत्य केवलं ब्रह्मानन्दसहोदरं सुखम् अनुभवति।

\begin{shloka}[center]

आनन्दः सहजस्तस्य व्यज्यते स कदाचन ।

व्यक्तिः सा तस्य चैतन्यचमत्काररसाह्वया  ॥

\end{shloka}

\customparagraph{\textbf{रसमहत्त्वम्}}

यथा दानं विना लक्ष्मीर्न शोभते तथैव रसैर्विना काव्यं न शोभते । काव्यसंसारम् अपारं विद्यते । तस्य कविरेव प्रजापतिरस्ति । तस्य कृते काव्यसंसारं यथा रोचते स तथैव निर्माति । काव्यकर्ता कविर्यदि शृङ्गारी चेत् काव्यजगद् रसमयं भवति । स एव वीतरागश्चेत् काव्यमपि नीरसं जायते । रसभावयोः सामीप्यसम्बन्धो भवति ।  भावा रसान् विभावयन्ति । भावैश्च रसा भाव्यन्ते ।

\begin{shloka}[center]

लक्ष्मीरिव विना त्यागान्न वाणी भाति नीरसाः ॥

अपारे काव्यसंसारे कविरेव प्रजापतिः ।

यथास्मै रोचते विश्वं तथेदं परिवर्तते ॥

शृङ्गारी चेत्कविः काव्ये जातं रसमयं जगत् ।

स चेत्कविर्वीतरागो नीरसं व्यक्तमेव तत्

\end{shloka}

\customparagraph{\textbf{नायकनायिकाभेदाः}}

तत्र प्रथमतो नायकभेदाश्चत्वारः- धीरोदात्तः, धीरोद्धतः, धीरललितः, धीरप्रशान्तश्चेति । एतेषां प्रत्येकस्य अनुकूलः, दक्षिणः, शठः, धृष्टश्चेति चतुरो भेदान् समेत्य नायकस्य षोडशो भेदाः । पीठमर्दः, विटः, विदूषकश्चेति त्रयः शृङ्गारिकनायकस्यानुगामिनः । एतेषु नायकतुल्यैश्वर्यबलाभ्यां समन्वितः पीठमर्दः, नायकस्यान्तरङ्गसहचरो विटः, परिहासप्रदायको विदूषकश्चेति । एवमेव कौशिकाचार्यमते स्वकीया, परकीया, पुनर्भूश्चेति तिस्रो नायिका भवन्ति । केषाञ्चिदाचार्याणां मते सामन्यैव नायिका भवति, न पुनर्भूः । तासां नायिकानामप्यन्येऽनेके भेदा भवन्ति ।  अनन्तरमत्र स्थायिसञ्चारादिभावानां स्वरूपवर्णनम्, हावभावादिनाञ्च वर्णनं विद्यते ।

\customparagraph{\textbf{अध्याय ३४०}}

अध्यायेऽस्मिन् रीतिवृत्तयोः परिचयो विद्यते । अग्निपुराणकारो रीतिवृत्तयोर्महत्त्वं स्वीकृत्य भेदप्रदर्शनपूर्वकं तयोर्विवेचनं करोति।

वाग्विद्यायाः सम्यग् ज्ञानाय रीतेर्निरूपणं क्रियते । प्रथमतो रीतेः पाञ्चाली, गौडदेशीया (गौडी), वैदर्भी, लाटजा (लाटी) इति चत्वारो भेदा भवन्ति । तासु उपचारयुता कोमलपदावलीभिर्लघुसमासैश्च संरचिता पाञ्चाली, क्षीणसन्दर्भयुता दीर्घसमासा गौडीया, नात्यल्पैर्नाप्यधिकैर्लाक्षणिकतत्त्वैर्युता नातिकोमलशब्दार्थसहिता मुक्तसमासा वैदर्भी, सरलवाक्यसङ्‌घटिता अस्फुटसमासा अधिकलाक्षणिकतत्त्वैर्विहीना लाटीया चेति चतस्रो रीतयो भवन्ति । क्रियासु अविषमा वाक्यरचना वृत्तिः कथ्यते । तस्या भारती, आरभटी, कौशिकी, सात्वती चेति चत्वारो भेदा भवन्ति ।

\customparagraph{\textbf{अध्याय ३४१}}

अध्यायेऽस्मिन् नृत्त्यादौ समावश्यकं अङ्गकर्मनिरूपणं विद्यते । नायकनायिकानां चेष्टादियुतमङ्गकर्म आङ्गिकमभिनयं कथ्यते । तस्याभिनयस्य द्वादश प्रकाराः सन्ति ।

अनन्तरं शिरहस्तवक्षःपार्श्वकटिपादप्रभृतीनामङ्गानाम्, भ्रूलतादीनां प्रत्यङ्गानाञ्च चेष्टावर्णनं विद्यते । तत्र शिरसस्त्रयोदश, भ्रुवः सप्त, नेत्रयोर्नव, ओष्ठयोः षोडश, चिबुकस्य सप्त, मुखस्य षोडश, ग्रीवाया नव इति विविधानामङ्गोपाङ्गानां चेष्टावर्णनं विद्यते ।

\customparagraph{\textbf{अध्याय ३४२}}

अध्यायेऽस्मिन् सत्त्वाश्रयः, वागाश्रयः, अङ्गाश्रयः, आहरणाश्रयश्चेति चत्वारोऽभिनयप्रकारा वर्णिता विद्यन्ते । अनन्तरं शृङ्गारादिरसलक्षणम्, अलङ्कारस्य परिभाषा, शब्दालङ्कारस्य दशभेदा वर्णिता विद्यन्ते ।

\customparagraph{\textbf{अभिनयप्रकाराः}}

नटः प्रेक्षकाणामाभिमुख्याय यत् कार्यं विदधाति स बुधैरभिनयो ज्ञेयः । तस्य स्तम्भाद्याश्रितः सात्त्विकः, वागाश्रितो वाचिकः, शरीराश्रित आङ्गिकः, बुद्ध्याश्रित आहार्यश्चेति चत्वारो भेदा भवन्ति ।  अनन्तरमभिनयप्रकरणे प्रस्तुतं नृत्यविषयकं मुद्रादिकमन्यत्र दुर्लभमेव । तद्यथा-

\begin{shloka}[center]

शिरः पाणिरुरः पार्श्वं कटिरङ्‌घ्रिरिति क्रमात् ॥

अङ्गानि भ्रूलतादीनि प्रत्यङ्गान्यभिजायते ।

अङ्गप्रत्यङ्गयोः कर्म प्रयत्नजनितं विना ॥

भ्रूकर्म सप्तधा ज्ञेयं पातनं भृकुटी मुखम् ।

दृष्टिस्त्रिधा रसस्थायिसञ्चारिप्रतिबन्धना ॥

\end{shloka}

\customparagraph{\textbf{अलङ्कारलक्षणम्}}

अग्निपुराणकारेण काव्यशास्त्रप्रतिपादितलक्षणानुकूलमलङ्कारलक्षणं विहितं विद्यते । अर्थालङ्कारसन्दर्भे अलङ्कारैर्विना वाणी विधवेव सौन्दर्यहीना जायत इति कथयन् पुराणकारः काव्येऽलङ्कारस्य महत्त्वं प्रदर्शयति । परं रसप्रकरणे रसतत्त्वस्यानिवार्यत्वमपि स्वीकरोति । अतोऽयं मध्यमार्गी इति मन्तुं युज्यते । अयं काव्यशोभाविधायका धर्मा एव अलङ्कारा भवन्तीति अलङ्कारलक्षणं करोति ।

\begin{shloka}[center]

काव्यशोभाकारान्धर्मानलङ्कारान् प्रचक्षते ॥

\end{shloka}

\customparagraph{\textbf{अलङ्कारभेदाः}}

शब्दालङ्कारः अर्थालङ्कारः शब्दार्थालङ्कारश्चेति त्रिविधा अलङ्कारा भवन्ति ।

\customparagraph{\textbf{अध्याय ३४३}}

अध्यायेऽस्मिन् शब्दालङ्कारभेदानां सस्वरूपं निरूपणं विद्यते । तेषु अनुप्रासयमकश्चित्रबन्धप्रभृतयः शब्दालङ्काराः प्रमुखाः सन्ति ।

पदवाक्ययोर्वर्णानामावृत्तिरनुप्रासः कथ्यते । अस्य एकवर्णगतावृत्तिः अनेकवर्णगतावृत्तिश्चेति द्वौ भेदौ भवतः । एकवर्णगतानुप्रासस्य मधुरा, ललिता, प्रौढा, भद्रा, परुषा चेति पञ्च प्रकारा भवन्ति ।

यत्र भिन्नार्थप्रतिपादकानाम् अनेकवर्णानामावृत्तिर्भवति तत्र यमकं भवति । तस्यापि अव्यपेतं व्यपेतञ्चेति द्वौ भेदौ भवतः । व्यवधानयुता वर्णावृत्तिर्व्यपेतं व्यवधानरहिता वर्णावृत्तिरव्यपेतं यमकं भवति । अनयोः स्थानपादयोराधारे चत्वारो भेदा भवन्ति ।

सहृदयगोष्ठ्यां कुतूहलाविष्कारको वाग्बन्धः चित्रं कथ्यते । नानार्थानुयोगादस्य प्रश्नः, प्रहेलिका, गुप्तम्, च्युतम्, दत्तम्, च्युतदत्ताक्षरञ्चेति पञ्च भेदा भवन्ति । अग्निपुराणकारेण विहितः शब्दालङ्कारस्य चित्रबन्धभेदस्य विस्तारः काव्यशास्त्रीयग्रन्थेषु दुर्लभ एव दृश्यते। अनेकधावृत्तवर्णविन्यासैर्विहिता चित्रादिवस्तूनां शिल्पकल्पना एव बन्धः कथ्यते । तस्य गोमुत्रिकाबन्धः, अर्धभ्रमणबन्धः, सर्वतोभद्रबन्धः, अम्बुजबन्धः, चक्रबन्धः, चक्राब्जबन्धः, दण्डबन्धः, मुरजबन्धश्चेति अष्टौ भेदा सलक्षणं निरूपिताः सन्ति ।

\customparagraph{\textbf{अध्याय ३४४}}

अध्यायेऽस्मिन्नर्थालङ्कारा विवेचिता विद्यन्ते । तेषु प्रथमतः स्वरूपहेतुसादृश्योत्प्रेक्षातिशयविभावनाविरोधहेतुसम इत्येता अष्टावलङ्कारा वर्णिताः सन्ति । तदनन्तरं सादृश्याधृता उपमा रूपकम् सहोक्तिरर्थान्तरन्यास इत्येते चत्वारोऽलङ्कारा विवेचिताः । तत्पश्चाद् उपमाभेदानां निरूपणम्, विरोधमूलकालङ्काराणां निरूपणम्, हेतुमूलकालङ्काराणां निरूपणञ्च विद्यते ।

\customparagraph{\textbf{अर्थालङ्कारस्वरूपं महत्त्वं भेदाश्च}}

अर्थस्यालङ्करणमर्थालङ्कारः कथ्यते । अर्थालङ्कारं विना शब्दसौन्दर्यमपि मनोहरं न भवति । यथा कटकाद्यलङ्कारैर्विना विधवा कमनीया न भवति तथैवोपमाद्यर्थालङ्कारैर्विना सरस्वती (काव्यम्) अपि कमनीया । अर्थालङ्कारस्य स्वरूपम्, सादृश्यम्, उत्प्रेक्षा, अतिशयोक्तिः, विभावना, विरोधः, हेतुः, समञ्चेति अष्टौ भेदा भवन्ति ।

\begin{shloka}[center]

अलङ्कारमर्थानामर्थालङ्कार इष्यते ।

तं विना शब्दसौन्दर्यमपि नास्ति मनोहरम् ॥

अर्थालङ्काररहिता विधवेव सरस्वती ।

स्वरूपमथ सादृश्यमुत्पेक्षातिशयावपि ॥

विभावना विरोधश्च हेतुश्च सममष्टधा ।

\end{shloka}

\customparagraph{\textbf{अध्याय ३४५}}

अध्यायेऽस्मिन् शब्दार्थसौन्दर्यप्रदायकाः शब्दार्थालङ्कारा (उभयालङ्काराः) वर्णिता विद्यन्ते। अत्र प्रथमतः प्रशस्तिकान्त्यादिषडलङ्कारविवेचनम्, तदनन्तरमाक्षेप- समासोक्तिपर्यायोक्त्यपह्नुत्यलङ्काराणां विवेचनं विद्यते । एतेऽलङ्कारा गम्यमानार्थप्रधाना भवन्ति ।

अलङ्कारवर्गीकरणेऽप्यस्य मते नावीन्यं प्रतिभाति । काव्यशास्त्रपरम्परायां प्रायशः शब्दालङ्कारः, अर्थालङ्कारश्चेति प्रामुख्येन प्रकारद्वयं प्राप्यते, परमग्निपुराणकारः सभेदं शब्दार्थालङ्कारान् निरूपयति । काव्यशास्त्रकारैरर्थालङ्कारे परिगणिताः समाधिः, आक्षेपः, समासोक्तिः, अपह्नुतिः, पर्यायोक्तिश्चेत्यलङ्कारान् शब्दार्थालङ्कारवर्गे स्थापयति । अयमलङ्कारसम्प्रदायस्य पक्षपाती मन्यते । तथापि ध्वनेर्महत्त्वं स्वीकरोति । अलङ्कारप्रकरणे ध्वनेर्विषयम् उल्लिखति यत्-

\begin{shloka}[center]

श्रुतेरलभ्यमानोऽर्थो यस्माद्भाति सचेतनः ।

स आक्षेपो ध्वनिः स्याच्च ध्वनिना व्यज्यते यतः ॥

\end{shloka}

तत्र शब्दार्थालङ्कारसन्दर्भेऽपि

\begin{shloka}[center]

एषामेकतमस्यैव समाख्या ध्वनिरित्यतः  ॥

\end{shloka}

अस्य मते काव्येऽलङ्कारस्यातीवावश्यकता विद्यते । अयं प्रस्तौति, यत्- `अर्थालङ्काररहिता विधवेव सरस्वती ।' अर्थात् अर्थालङ्कारेण रहिता वाणी विधवेव सौन्दर्यहीना भवति । तस्मात् काव्येऽलङ्कारस्यावश्यकता विद्यते । उत्तरवर्त्तिनोऽलङ्कारवादिनो विद्वांसोऽस्यालङ्कारमतं सम्मानेन स्वीकुर्वन्ति ।

\customparagraph{\textbf{अध्याय ३४६}}

अध्यायेऽस्मिन् काव्यगुणविवेचनं विद्यते । अत्र गुणस्य महत्त्वम्, तस्य भेदोपभेदाः, सप्त शब्दगुणाः, षडर्थगुणाः, षट् शब्दार्थगुणाश्च  विवेचिता विद्यन्ते ।

गुणस्य स्वरूपम्, तस्य वर्णीकरणञ्च नवीनमाभाति। अलङ्‌कृतमपि निर्गुणं काव्यं लावण्यरहिताया युवत्याः शरीरे हारमिव भारायते इत्युल्लिख्य गुणमहत्त्वं प्रकाशयति । काव्यगुणं प्राक् सामान्यविशेषाविति द्विधा विभज्य सामान्यस्य शब्दगुणः, अर्थगुणः, उभयगुणश्चेति प्रकारत्रयमुल्लिखति । तत्र सप्त शब्दगुणाः, षड् अर्थगुणाः, षड् उभयगुणाश्चेति । विशेषस्तु स्वलक्षणगोचर एक एव । एतादृशं गुणस्य विंशतिप्रकारसमन्वितं नूतनं वर्गीकरणं पश्चवर्त्ती काव्यशास्त्री भोजः स्वकीये सरस्वतीकण्ठाभरणेऽनुकरोति ।

\customparagraph{\textbf{अध्याय ३४७}}

अध्यायेऽस्मिन् काव्यदोषविवेचनं विद्यते । अत्र दोषाणां भेदोपभेदान् विविच्य परिहारप्रकारश्च निरूप्यते । अध्यायान्ते कविसमयख्यातेर्भेदोपभेदवर्णनमपि समावेश्यते ।

सभ्यानामपि उद्वेगजनको दोषो भवति । स च वक्तृवाचकवाच्यानामेकद्वित्रिनियोगात् सप्तधा विभज्यते । यथा- वक्तृनियुक्तः, वाचकनियुक्तः, वाच्यनियुक्तः, वक्तृवाचकनियुक्तः, वाचकवाच्यनियुक्तः, वक्तृवाच्यनियुक्तः, वक्तृवाचकवाच्यनियुक्तश्चेति ।

पौरस्त्यकाव्यशास्त्रस्य विकासे पौराणिकग्रन्थानामपि महनीयं योगदानं विद्यते। तत्र प्रथमक्रमे ‘अग्निपुराणम्’ समायाति । अत्र प्रायशः काव्यशास्त्रीयतत्त्वानां चर्चा विद्यते । काव्यलक्षणम्, काव्यभेदाः, महाकाव्यलक्षणम्, रूपकभेदाः, नाटकलक्षणम्, अभिनयप्रकाराः, काव्यालङ्काराः, गुणाः, दोषा इत्यादयः सर्वेऽपि विषया एकादशसु अध्यायेषु सङ्‌गृहीता विद्यन्ते । कस्मिंश्चित् स्थले भामहदण्ड्यादिभिरलङ्कारशास्त्रिभिरुल्लिखितानि पद्यानि यथावत् स्थापितानि विद्यन्ते । कुत्रचित्तु नावीन्यमपि दृश्यते । अस्य पुराणस्य रचना कदा जाता इति विषये विदुषां मतैक्यं नास्ति । केचिद् विद्वांसो भामहसमयात् प्राग् भरतमुनिसमयाच्च पश्चादस्य रचना जातेति मन्यन्ते चेत् केचित्तु काव्यशास्त्रस्य प्रारम्भिकविकासानन्तरं भामहदण्डिनोः समयानन्तरमस्य रचना जतेति स्वीकुर्वन्ति । केचित्तु पुराणं प्रग् रचितम्, काव्यशास्त्रीयखण्डः पश्चात् प्रक्षिप्त इत्यपि वदन्ति । पुराणस्य रचनाशैल्यादिकमवलोक्य पुराणमिदम् अन्यपुराणेभ्योऽर्वाचीनं प्रतिभाति ।



\end{document}